% Preamble for the swarm-transcribed Weber volumes.
% Deliberately close to weber_v1_ge.tex so the verified p1-148 region and the swarm
% region can live in one document.
\documentclass[11pt]{article}
\usepackage[ngerman,english]{babel}
\usepackage[a4paper,margin=0.9in]{geometry}
\usepackage[T1]{fontenc}
\usepackage[utf8]{inputenc}
\usepackage{lmodern}
\usepackage[leqno]{amsmath}
\usepackage{amssymb,mathtools,mathrsfs}
\usepackage{microtype}
\usepackage{fancyhdr}
\usepackage{array,longtable,enumitem,multicol}
\usepackage{tikz}
\usepackage{graphicx}
\graphicspath{{../figures/b1/}{../figures/b2/}{../figures/b3/}}
\setlength{\parindent}{1.5em}
\setlength{\parskip}{0.18em}
\emergencystretch=6em
\setlength{\headheight}{14pt}
\pagestyle{fancy}
\fancyhf{}
\cfoot{\thepage}

% --- Weber conventions -------------------------------------------------------
% Sperrung (letter-spacing) is Weber's emphasis. There is no bold in his body text.
\providecommand{\gesperrt}[1]{\textls[130]{#1}}
% Printed-page marker: a no-op in the PDF. The machine-readable anchor is the
% "% !PRINTED n" comment in the source, which the collector and Codex both read.
\providecommand{\weberpage}[1]{}
% Displayed equation numbered in the LEFT margin, as Weber prints them.
\providecommand{\weq}[2]{\begin{equation}\tag{#1}#2\end{equation}}
% Figure placeholder. A later pass replaces it with TikZ or a cleaned scan crop.
% A marked gap is honest; a fabricated diagram is not.
\providecommand{\weberfigure}[2]{%
  \begin{center}\fbox{\parbox{0.78\linewidth}{\centering\small [Figur, S.~#1: #2]}}\end{center}}
% Source-faithful figure crop.  The asset search path works both from the live
% swarm_build directory and from a self-contained release bundle with assets/.
\providecommand{\weberfigureimage}[2]{%
  \begin{center}\includegraphics[width=#2\linewidth]{#1}\end{center}}

% Text symbols agents reach for that are not always defined. Providing them in the preamble
% fixes every occurrence at once and leaves the generated fragments untouched, which is the
% right trade whenever the fix is "define the macro" rather than "change the text".
\usepackage{textcomp}
% amsmath caps matrix environments at 10 columns. Weber prints permutation matrices far wider
% than that -- Bd III has an 11-column pmatrix of the w_nu -- and the failure is the opaque
% "Extra alignment tab has been changed to \cr".
\setcounter{MaxMatrixCols}{30}
\providecommand{\textquotedblbase}{\glqq}
\providecommand{\textquotedblleft}{``}
\providecommand{\textquotedblright}{''}
\providecommand{\textquotesinglbase}{\glq}

% --- Greek capitals are SLANTED in this fount -------------------------------------------
% RULING 2026-08-07, after the same observation was reported from five separate pages
% (printed 271, 275, 285, 333 and the nu/v cousin on 172).
%
% Verified at 8x zoom on printed 333 (b1 idx 358): every capital Delta on that page --
% "f(x) - f(alpha) = Delta", "Delta_alpha = f(alpha+delta) - f(alpha)" -- is the INCLINED
% form, which LaTeX calls \varDelta. \Delta, LaTeX's default, is upright and is not what
% Vieweg set. The same inclination is visible on capital Phi (printed 171, 282).
%
% This is a property of the TYPEFACE, not of any occurrence -- the same kind of fact as the
% varkappa ruling -- so it holds wherever the letter appears, and the right fix is one line
% here rather than several thousand edits across the fragments. \let leaves every fragment
% untouched, applies to the whole corpus at once, and is reversible by deleting these lines.
%
% Delta and Phi are the two verified by eye. The rest follow from its being a single fount;
% if any capital is later found upright on the page, delete its line and the fragments still
% say exactly what they said before.
\let\Gamma\varGamma
\let\Delta\varDelta
\let\Theta\varTheta
\let\Lambda\varLambda
\let\Xi\varXi
\let\uprightPi\Pi
\let\Pi\varPi
\let\Sigma\varSigma
\let\Upsilon\varUpsilon
\let\Phi\varPhi
\let\Psi\varPsi
\let\Omega\varOmega
% Lower-case rho is the OPEN-TAILED form in this fount. Verified at 10x on printed 419
% ("wenn also überhaupt ein \varrho vorhanden ist"): the tail curves down and left rather than
% descending straight, which is \varrho, not \rho. Same kind of fact as the capitals above, and
% the same one-line treatment. (kappa is already handled: the fragments write \varkappa
% explicitly, on the ruling recorded in resolve_class.py.)
\let\rho\varrho

% --- Fractions written in running text -------------------------------------------------
% Gauss's German prose says things like "von 6\tfrac{1}{2} Jahren" and "\tfrac{1}{2} Grad":
% a fraction as an ordinary word, not inside math. Bare \tfrac there is
%     ! Missing $ inserted.
% and the error points at the sentence, not at the macro.
%
% Wrapping the three fraction macros in \ensuremath fixes every occurrence at once and leaves
% the fragments untouched -- the right trade whenever the fix is "define the macro" rather than
% "change the text", as with the Greek capitals above. Inside math \ensuremath is a no-op, so
% the thousands of legitimate in-math uses are unaffected.
\let\swarmfrac\frac
\let\swarmtfrac\tfrac
\let\swarmdfrac\dfrac
\renewcommand{\frac}[2]{\ensuremath{\swarmfrac{#1}{#2}}}
\renewcommand{\tfrac}[2]{\ensuremath{\swarmtfrac{#1}{#2}}}
\renewcommand{\dfrac}[2]{\ensuremath{\swarmdfrac{#1}{#2}}}

% --- Unicode vulgar fractions -----------------------------------------------------------
% Agents sometimes type the single Unicode glyph the page shows -- "erstere ⅔" -- rather than a
% macro. utf8 inputenc has no default binding for U+2153..U+215E and stops the build with
%     ! LaTeX Error: Unicode character ⅔ (U+2154)
% Declaring them costs nothing, keeps the fragment reading exactly as the agent saw the page,
% and is the same "define it in the preamble" trade as the fraction macros above. Typing the
% glyph is not wrong; it is arguably closer to the page than \tfrac would be.
\DeclareUnicodeCharacter{2153}{\ensuremath{\swarmtfrac{1}{3}}}
\DeclareUnicodeCharacter{2154}{\ensuremath{\swarmtfrac{2}{3}}}
\DeclareUnicodeCharacter{2155}{\ensuremath{\swarmtfrac{1}{5}}}
\DeclareUnicodeCharacter{2156}{\ensuremath{\swarmtfrac{2}{5}}}
\DeclareUnicodeCharacter{2157}{\ensuremath{\swarmtfrac{3}{5}}}
\DeclareUnicodeCharacter{2158}{\ensuremath{\swarmtfrac{4}{5}}}
\DeclareUnicodeCharacter{2159}{\ensuremath{\swarmtfrac{1}{6}}}
\DeclareUnicodeCharacter{215A}{\ensuremath{\swarmtfrac{5}{6}}}
\DeclareUnicodeCharacter{215B}{\ensuremath{\swarmtfrac{1}{8}}}
\DeclareUnicodeCharacter{215C}{\ensuremath{\swarmtfrac{3}{8}}}
\DeclareUnicodeCharacter{215D}{\ensuremath{\swarmtfrac{5}{8}}}
\DeclareUnicodeCharacter{215E}{\ensuremath{\swarmtfrac{7}{8}}}
% Prime and double prime: arcminutes and arcseconds. Gauss's geodesy and astronomy volumes are
% full of "43° 37' 6\"741", and agents type the true Unicode prime characters rather than the
% ASCII apostrophe -- which is right, since that is what the page prints.
\DeclareUnicodeCharacter{2032}{\ensuremath{{}^{\prime}}}
\DeclareUnicodeCharacter{2033}{\ensuremath{{}^{\prime\prime}}}
\DeclareUnicodeCharacter{2034}{\ensuremath{{}^{\prime\prime\prime}}}
\DeclareUnicodeCharacter{00BD}{\ensuremath{\swarmtfrac{1}{2}}}
\DeclareUnicodeCharacter{00BC}{\ensuremath{\swarmtfrac{1}{4}}}
\DeclareUnicodeCharacter{00BE}{\ensuremath{\swarmtfrac{3}{4}}}

% Macros agents invent for a text fraction. \textfrac is not a real LaTeX command, but the
% intent is unambiguous and providing it is cheaper and safer than rewriting the fragment --
% the alternative is editing a table cell by hand and risking its alignment.
\providecommand{\textfrac}[2]{\ensuremath{\swarmtfrac{#1}{#2}}}

% Common shorthands inherited from the volume I file, so fragments that use them compile.
\providecommand{\eps}{\varepsilon}
\providecommand{\ii}{\mathrm{i}}
\providecommand{\tg}{\operatorname{tg}}
% Gauss writes the trigonometric functions in the period's Latin abbreviations -- "tang" for
% tangent, "cotang" for cotangent -- rather than Weber's "tg"/"cotg". Same function, different
% century's shorthand; both are reproduced as the author sets them.
\providecommand{\tang}{\operatorname{tang}}
\providecommand{\cotang}{\operatorname{cotang}}
\providecommand{\arctang}{\operatorname{arc\,tang}}
\providecommand{\ctg}{\operatorname{cotg}}
\providecommand{\cotg}{\operatorname{cotg}}
\providecommand{\sn}{\operatorname{sn}}
\providecommand{\cn}{\operatorname{cn}}
\providecommand{\dn}{\operatorname{dn}}
\providecommand{\am}{\operatorname{am}}
\providecommand{\dd}{\,d}
\providecommand{\tht}{\vartheta}
\providecommand{\vark}{\varkappa}
\providecommand{\Norm}{\operatorname{N}}
\providecommand{\ord}{\operatorname{ord}}
\providecommand{\ind}{\operatorname{ind}}
\providecommand{\abs}[1]{\left|#1\right|}
\providecommand{\sect}[2]{\section*{\S.\,#1. #2}}
\providecommand{\frp}{\mathfrak p}
\providecommand{\frP}{\mathfrak P}
\providecommand{\frG}{\mathfrak G}
\providecommand{\frA}{\mathfrak A}
\providecommand{\frB}{\mathfrak B}

% --- reader additions ---------------------------------------------------------
\usepackage{url}
\usepackage[unicode,bookmarksnumbered=false,bookmarksopen=true,bookmarksopenlevel=1,colorlinks=true,
  linkcolor=black,urlcolor=blue]{hyperref}
\usepackage{bookmark}
\renewcommand{\headrulewidth}{0pt}
\renewcommand{\weberpage}[1]{\phantomsection\label{wp:#1}}
\newcommand{\inhaltkopf}[1]{\begin{center}\large\textbf{#1}\end{center}\medskip}
\newcommand{\inhaltbuch}[2]{\par\medskip\begin{center}\if\relax\detokenize{#1}\relax\else#1\\[2pt]\fi{\large\textbf{#2}}\end{center}}
\newcommand{\inhaltabschnitt}[2]{\par\smallskip\begin{center}\small\if\relax\detokenize{#1}\relax\else#1\\\fi\textbf{#2}\end{center}}
\newcommand{\inhaltpar}[3]{\par\noindent\hangindent=3.4em\hbox to 3.4em{\S~#1.\hfil}#2\dotfill\,\hyperref[wp:#3]{#3}\par}
\newcommand{\inhaltzeile}[2]{\par\noindent #1\dotfill\,\hyperref[wp:#2]{#2}\par}

\begin{document}
\selectlanguage{english}
\pagenumbering{roman}
\begin{center}
{\Large Heinrich Weber}\\[6pt]
{\LARGE\textbf{Lehrbuch der Algebra}}\\[4pt]
{\large (Textbook of Algebra)}\\[8pt]
{\large Volume I}\\[4pt]
Braunschweig: Friedrich Vieweg und Sohn, 1895\\[24pt]
{\large English translation}\\[4pt]
with the page numbers of the original
\end{center}
\clearpage
\noindent Translated from the scan \url{https://archive.org/details/lehrbuchderalgeb01weberich} (Internet Archive).
The German text was transcribed from the scan and translated into English by AI models (Anthropic Claude and
OpenAI Codex), and both were checked against the scan by AI. Weber's text is in the public domain; the
transcription, translation and typesetting are dedicated to the public domain under CC0~1.0 (\url{https://creativecommons.org/publicdomain/zero/1.0/}).
Source files: \url{https://github.com/KokunoYumeto/weber-lehrbuch-der-algebra}.
\clearpage

% ===== front matter, scan leaf 10 =====
% not printed: faint show-through of the following title page
\begin{center}
\vspace*{6cm}
{\Large LEHRBUCH}\\[1.2cm]
{\small DER}\\[1cm]
{\Huge \gesperrt{ALGEBRA}.}\\[0.6cm]
{\large [Textbook of Algebra]}\\[1.2cm]
\rule{3.5cm}{0.4pt}
\end{center}

\clearpage
% ===== front matter, scan leaf 12 =====
% not printed: faint mirror-reversed library stamp ("MATH.-STAT.") between the title and "VON"
\begin{center}
\vspace*{2cm}
{\LARGE \gesperrt{LEHRBUCH}}\\[1.5em]
{\small DER}\\[1.5em]
{\Huge \gesperrt{ALGEBRA.}}\\[0.8em]
{\large [TEXTBOOK OF ALGEBRA]}\\[2.5em]
{\small BY}\\[1.5em]
{\LARGE HEINRICH WEBER,}\\[0.4em]
{\footnotesize PROFESSOR OF MATHEMATICS AT THE UNIVERSITY OF GÖTTINGEN.}\\[2em]
\rule{3cm}{0.4pt}\\[3em]
{\large \gesperrt{IN TWO VOLUMES.}}\\[2em]
{\large \textbf{FIRST VOLUME.}}\\[2em]
\rule{2.5cm}{0.4pt}\\[2.5em]
{\small \gesperrt{WITH 28 FIGURES PRINTED IN THE TEXT.}}\\[3em]
\rule{\textwidth}{1pt}\\[1.5em]
{\large BRAUNSCHWEIG,}\\[0.5em]
{\small PRINTED AND PUBLISHED BY FRIEDRICH VIEWEG AND SON.}\\[0.5em]
{\large \gesperrt{1895.}}
\end{center}

\clearpage
% ===== front matter, scan leaf 13 =====
% not printed: library stamp "MATH-STAT." with handwritten initials; handwritten note "Cat. for Math-Stat Lib. Gift of M. W. Haskell"; show-through of the title page from the other side of the leaf
\vspace*{0.35\textheight}
\begin{center}
\rule{0.6\textwidth}{0.4pt}\\[6pt]
\gesperrt{All rights reserved}.\\[2pt]
\rule{0.6\textwidth}{0.4pt}
\end{center}

\clearpage
% ===== front matter, scan leaf 14 =====
% not printed: handwritten call number "QA153 W35 v.1"; library stamps "MATH.-STAT. LIBRARY" and "M777593"
\begin{center}
{\Large \gesperrt{PREFACE.}}\\[1ex]
\rule{4em}{0.4pt}
\end{center}

In view of the development that algebra has undergone in recent decades, a comprehensive presentation and linking of the various theoretical considerations and manifold applications may well be useful, even after the textbook of \gesperrt{Serret}, which was excellent for its time.

For years I have cherished the plan of such an undertaking, which indeed appeared vast and far-reaching and required a good deal of preparatory work. Only after I had several times traversed the field as a whole in university lectures, and had treated individual parts in greater detail, did I resolve to set about the execution of the work, of which the first volume now lies before the reader in completed form.

It was my intention to provide a textbook which, without presupposing much prior knowledge, would introduce the reader to modern algebra and also lead him on to the higher and more difficult parts, in which interest in the subject first becomes truly alive. At the same time, the necessary tools, the elementary as well as the higher ones, were to be derived from the course of the development itself, in order to make the presentation as independent of other textbooks as possible.

There are two things that have become of quite special importance for the most recent development of algebra; these are, on the one hand, group theory, which is gaining ever greater dominance and whose ordering and clarifying influence can be felt everywhere, and then the intervention of number theory.

\clearpage
% ===== front matter, scan leaf 15 =====
% !PRINTED VI
\weberpage{VI}
\noindent Although algebra in part goes beyond the theory of numbers, and reaches over into other fields, e.g.\ the theory of functions, or, in its applications, into geometry as well, the science of numbers nevertheless always remains the foremost example for all algebraic considerations, and the questions of number theory that stand today in the foreground of interest are predominantly of an algebraic nature. This marked out the path that I had to follow in my work.

The extensive material has been divided into two volumes. The first volume contains the elementary part of algebra, which, by a traditional expression, may be called \gesperrt{calculation with letters}, then the rules for the numerical computation of the roots of equations, and the beginnings of \gesperrt{Galois}'s theory.

The second volume, which will, I hope, follow the first within a short time, is to present the general theory of finite groups, the theory of linear substitution groups, and applications to various individual problems, and is to conclude with the theory of algebraic numbers, where an attempt is made to unite the various points of view from which this theory has hitherto been considered.

Finally, the second volume is to provide an alphabetical index to both volumes.

As cannot be otherwise in a discipline that is in the course of rapid development and on which work is being done from the most diverse quarters, in algebra, too, the use of language and the notation are very varied and frequently not in agreement. This makes a unified presentation, as well as penetration into the various works, very difficult.

I have therefore endeavoured to adhere uniformly to a mode of expression that is as appropriate as possible, and in doing so I have frequently taken counsel with colleagues. I may express the hope of having thereby contributed to the consolidation of a uniform terminology.

The references to the literature and the historical notes given in the book make no claim whatsoever to completeness, although I have endeavoured

\clearpage
% ===== front matter, scan leaf 16 =====
% !PRINTED VII
\weberpage{VII}

\noindent to mention, as far as possible, the most important sources and literary aids at the appropriate places.

It is a pleasant duty for me to express my thanks here to many a friend and colleague who has taken a lively and active interest in the progress of my work. These thanks go first of all to my friend \gesperrt{Dedekind} for his faithful help with the proofreading. And even though he has exercised no influence on the plan and the execution of my work, I should nevertheless not like to leave unmentioned something from many years ago. Through a set of notes of a lecture course on higher algebra, in particular on the theory of \gesperrt{Galois}, which he gave in Göttingen in the winter of 1857/58, I gained an even livelier interest in this theory. Until then it had probably never been presented at our universities with such completeness.

I must also make mention here of the manifold stimulation and instruction that I owe to my friend and colleague F.~\gesperrt{Klein}. He has followed the progress of the work with the liveliest interest, and his expert advice, always most readily given, has been of great influence on many parts of the book.

I cannot name here all the colleagues and friends who, as the reader will find in the relevant places, have supported me with their advice. But I must still make mention here of Messrs.\ E.~\gesperrt{Hess} in Marburg, Fr.~\gesperrt{Meyer} in Clausthal and R.~\gesperrt{Fricke} in Braunschweig. Through their expert and careful execution of the laborious correction of the proof sheets, they have furthered the accuracy and correctness of the text.

Finally, my thanks go to the publishing house, which, by readily accommodating my wishes and by its care in printing and production, has contributed substantially to the success of the whole.

\bigskip

\gesperrt{Göttingen}, November 1894.

\medskip

\begin{flushright}
\textbf{\gesperrt{The Author}.}
\end{flushright}

\clearpage
% ===== front matter, scan leaf 18 =====
\inhaltkopf{CONTENTS OF THE FIRST VOLUME.}
\inhaltzeile{Introduction}{1}
\inhaltbuch{First Book.}{The Foundations.}
\inhaltabschnitt{First Section.}{Rational Functions.}
\inhaltpar{1}{Integral functions}{23}
\inhaltpar{2}{A theorem of \gesperrt{Gauss}}{25}
\inhaltpar{3}{Division}{27}
\inhaltpar{4}{Division by a linear function}{30}
\inhaltpar{5}{Fractional functions; divisibility}{32}
\inhaltpar{6}{Greatest common divisor}{34}
\inhaltpar{7}{Products of linear factors}{39}
\inhaltpar{8}{The binomial theorem}{42}
\inhaltpar{9}{Interpolation}{44}
\inhaltpar{10}{Solution of the interpolation problem by means of the differences}{46}
\inhaltpar{11}{Arithmetic series of higher order}{47}
\inhaltpar{12}{The multinomial theorem}{50}
\inhaltpar{13}{Derived functions}{51}
\inhaltpar{14}{Derivative of a product}{54}
\inhaltpar{15}{Integral functions of several variables: forms}{56}
\inhaltpar{16}{The derivatives of functions of several variables}{59}
\inhaltpar{17}{\gesperrt{Euler}'s theorem on homogeneous functions}{62}
\inhaltabschnitt{Second Section.}{Determinants.}
\inhaltpar{18}{Permutations of $n$ elements}{64}
\inhaltpar{19}{Permutations of the first and second kind}{65}
\inhaltpar{20}{Determinants}{68}
\inhaltpar{21}{Principal theorems on determinants}{70}

\clearpage
% ===== front matter, scan leaf 19 =====
% !PRINTED X
\weberpage{X}
\inhaltpar{22}{Minors}{72}
\inhaltpar{23}{Minors in the wider sense}{77}
\inhaltpar{24}{Linear homogeneous equations}{81}
\inhaltpar{25}{Elimination from linear equations}{87}
\inhaltpar{26}{Inhomogeneous linear equations}{89}
\inhaltpar{27}{Multiplication of determinants}{92}
\inhaltpar{28}{Determinants of the minors}{95}
\inhaltpar{29}{Interpolation}{98}
\inhaltabschnitt{Third Section.}{The roots of algebraic equations.}
\inhaltpar{30}{Concept of the roots. Multiple roots}{101}
\inhaltpar{31}{Continuity of integral functions}{103}
\inhaltpar{32}{Change of sign of $f(x)$. Roots of equations of odd degree and of pure equations}{107}
\inhaltpar{33}{Solution of pure equations by trigonometric functions}{111}
\inhaltpar{34}{Removal of the second term of an equation}{114}
\inhaltpar{35}{Cubic equations. \gesperrt{Cardano}'s formula}{116}
\inhaltpar{36}{\gesperrt{Cayley}'s expression for \gesperrt{Cardano}'s formula}{118}
\inhaltpar{37}{The biquadratic equations}{119}
\inhaltpar{38}{Proof of the fundamental theorem of algebra}{121}
\inhaltpar{39}{Algorithm for computing the roots}{127}
\inhaltpar{40}{Continuity of the roots}{132}
\inhaltabschnitt{Fourth Section.}{Symmetric functions.}
\inhaltpar{41}{Concept of the symmetric functions. Fundamental symmetric functions}{138}
\inhaltpar{42}{The power sums}{140}
\inhaltpar{43}{Proof of the main theorem for two variables}{143}
\inhaltpar{44}{General proof of the main theorem}{144}
\inhaltpar{45}{Second proof of the theorem on symmetric functions}{147}
\inhaltpar{46}{Discriminants}{150}
\inhaltpar{47}{Discriminants of the forms of the third and fourth order}{153}
\inhaltpar{48}{Resultants}{156}
\inhaltpar{49}{Elimination. \gesperrt{Bezout}'s theorem}{159}
\inhaltpar{50}{Elimination from several equations}{161}
\inhaltpar{51}{Reducible and irreducible functions}{164}
\inhaltpar{52}{Tschirnhausen transformation}{170}
\inhaltpar{53}{Application to the cubic and biquadratic equations}{173}
\inhaltpar{54}{The Tschirnhausen transformation of the equation of the 5\textsuperscript{th} degree}{175}
\inhaltabschnitt{Fifth Section.}{Linear transformation. Invariants.}
\inhaltpar{55}{Linear transformation}{178}
\inhaltpar{56}{Quadratic forms}{179}

\clearpage
% ===== front matter, scan leaf 20 =====
% !PRINTED XI
\weberpage{XI}
\inhaltpar{57}{Transformation of quadratic forms into a sum of squares}{181}
\inhaltpar{58}{Law of inertia of quadratic forms}{183}
\inhaltpar{59}{Transformation of forms of the $n^{\text{th}}$ degree}{185}
\inhaltpar{60}{Invariants and Covariants}{186}
\inhaltpar{61}{Linear transformation of binary forms}{189}
\inhaltpar{62}{Binary cubic forms}{192}
\inhaltpar{63}{The complete system of forms of the binary cubic form}{196}
\inhaltpar{64}{Biquadratic forms}{199}
\inhaltpar{65}{Solution of the biquadratic equation}{201}
\inhaltpar{66}{The covariants}{203}
\inhaltpar{67}{The complete system of invariants of the binary biquadratic form}{206}
\inhaltabschnitt{Sixth Section.}{Tschirnhausen Transformation.}
\inhaltpar{68}{\gesperrt{Hermite}'s form of the Tschirnhausen transformation}{210}
\inhaltpar{69}{Invariant property of the Tschirnhausen transformation}{212}
\inhaltpar{70}{Elaborations on \gesperrt{Hermite}'s theorem}{215}
\inhaltpar{71}{Transformation of the cubic equation}{219}
\inhaltpar{72}{General execution of the transformation}{223}
\inhaltpar{73}{The Bezoutiant}{225}
\inhaltpar{74}{Transformation of the equation of the fifth degree}{230}
\inhaltpar{75}{Normal form of the equation of the fifth degree}{233}
\inhaltbuch{Second Book.}{The Roots.}
\inhaltabschnitt{Seventh Section.}{Reality of the Roots.}
\inhaltpar{76}{Generalities on the reality of roots of equations and on discriminants}{241}
\inhaltpar{77}{Discussion of the quadratic and cubic equations}{243}
\inhaltpar{78}{Discussion of the biquadratic equation}{246}
\inhaltpar{79}{The Bezoutiant and its significance for the reality of the roots}{252}
\inhaltpar{80}{The inertia of forms of the second degree}{255}
\inhaltpar{81}{Quadratic forms with vanishing determinant}{257}
\inhaltpar{82}{Quadratic forms with non-vanishing determinant}{260}
\inhaltpar{83}{Number of the positive and negative squares}{261}
\inhaltpar{84}{Application to the Bezoutiant}{265}
\inhaltabschnitt{Eighth Section.}{Sturm's Theorem.}
\inhaltpar{85}{\gesperrt{Sturm}'s problem}{270}
\inhaltpar{86}{\gesperrt{Sturm}'s chains}{271}

\clearpage
% ===== front matter, scan leaf 21 =====
% !PRINTED XII
\weberpage{XII}

\inhaltpar{87}{First example: spherical harmonics}{273}
\inhaltpar{88}{Second example}{276}
\inhaltpar{89}{\gesperrt{Sturm}'s functions}{279}
\inhaltpar{90}{\gesperrt{Hermite}'s solution of \gesperrt{Sturm}'s problem}{280}
\inhaltpar{91}{Determination of \gesperrt{Hermite}'s form $H$}{282}
\inhaltpar{92}{The determinant of \gesperrt{Hermite}'s form}{283}
\inhaltpar{93}{Outline of the theory of characteristics}{285}
\inhaltpar{94}{Characteristic of a system of three functions}{287}
\inhaltpar{95}{Relation of the characteristics to the points of intersection}{290}
\inhaltpar{96}{Application of the characteristics to the localization of the complex roots of an equation}{292}
\inhaltpar{97}{Determination of the characteristic}{294}
\inhaltpar{98}{\gesperrt{Gauss}'s first proof of the fundamental theorem of algebra}{295}

\inhaltabschnitt{Ninth Section.}{Estimation of the Roots.}

\inhaltpar{99}{The \gesperrt{Budan-Fourier} theorem}{299}
\inhaltpar{100}{\gesperrt{Newton}'s rule}{304}
\inhaltpar{101}{\gesperrt{Descartes}'s theorem}{308}
\inhaltpar{102}{\gesperrt{Jacobi}'s criterion}{311}
\inhaltpar{103}{\gesperrt{Klein}'s geometric comparison of the various criteria}{312}
\inhaltpar{104}{Determination of an upper bound for the roots}{316}
\inhaltpar{105}{Estimation of the imaginary roots}{318}
\inhaltpar{106}{\gesperrt{Rolle}'s theorem}{319}
\inhaltpar{107}{\gesperrt{Laguerre}'s theorems for equations with only real roots}{322}

\inhaltabschnitt{Tenth Section.}{Approximate Calculation of the Roots.}

\inhaltpar{108}{Interpolation. Regula falsi}{331}
\inhaltpar{109}{\gesperrt{Newton}'s method of approximation}{335}
\inhaltpar{110}{\gesperrt{Daniel Bernoulli}'s method of approximation and related methods}{341}
\inhaltpar{111}{\gesperrt{Gräffe}'s method of approximation}{344}
\inhaltpar{112}{Trigonometric solution of cubic equations}{349}
\inhaltpar{113}{\gesperrt{Gauss}'s method for the solution of trinomial equations}{352}
\inhaltpar{114}{Calculation of the imaginary roots of a trinomial equation}{355}

\inhaltabschnitt{Eleventh Section.}{Continued Fractions.}

\inhaltpar{115}{Conversion of rational fractions into continued fractions}{358}
\inhaltpar{116}{Continued-fraction expansion of irrational numbers}{361}
\inhaltpar{117}{The convergents}{362}
\inhaltpar{118}{Solution of indeterminate equations in two unknowns}{365}

\clearpage
% ===== front matter, scan leaf 22 =====
% !PRINTED XIII
\weberpage{XIII}
\inhaltpar{119}{Convergence of the approximating fractions}{369}
\inhaltpar{120}{Equivalent numbers}{371}
\inhaltpar{121}{Expansion of equivalent numbers in continued fractions}{374}
\inhaltpar{122}{Quadratic irrational numbers}{377}
\inhaltpar{123}{Reduced numbers with negative discriminant}{379}
\inhaltpar{124}{Reduced numbers with positive discriminant}{383}
\inhaltpar{125}{Expansion of real quadratic irrational numbers in continued fractions}{386}
\inhaltpar{126}{Examples}{392}
\inhaltpar{127}{The \gesperrt{Pell} equation}{395}
\inhaltpar{128}{Derivation of all solutions of the \gesperrt{Pell} equation from the smallest positive one}{399}
\inhaltpar{129}{Approximate computation of the real roots of a numerical equation by continued fractions}{401}
\inhaltpar{130}{Rational roots of equations with integer coefficients. Reducible equations}{403}
\inhaltabschnitt{Twelfth Section.}{Theory of the Roots of Unity.}
\inhaltpar{131}{The roots of unity}{408}
\inhaltpar{132}{Primitive roots of unity}{411}
\inhaltpar{133}{Equations for the primitive $n$th roots of unity}{414}
\inhaltpar{134}{Irreducibility}{417}
\inhaltpar{135}{The discriminant of the cyclotomic equation}{421}
\inhaltpar{136}{Primitive congruence roots}{425}
\inhaltpar{137}{Multiplication and division of the trigonometric functions}{432}
\inhaltpar{138}{Determination of the sign. Quadratic residues}{439}
\inhaltbuch{Third Book.}{Algebraic Quantities.}
\inhaltabschnitt{Thirteenth Section.}{The Galois Theory.}
\inhaltpar{139}{The concept of a field}{449}
\inhaltpar{140}{Adjunction}{451}
\inhaltpar{141}{Functions in a field}{452}
\inhaltpar{142}{Algebraic fields}{455}
\inhaltpar{143}{Simultaneous adjunction of several algebraic quantities}{457}
\inhaltpar{144}{Primitive and imprimitive fields}{460}
\inhaltpar{145}{Normal fields. \gesperrt{Galois} resolvent}{464}
\inhaltpar{146}{The substitutions of a normal field}{467}
\inhaltpar{147}{Composition of substitutions}{470}
\inhaltpar{148}{Permutation groups}{472}
\inhaltpar{149}{Galois group}{476}

\clearpage
% ===== front matter, scan leaf 23 =====
% !PRINTED XIV
\weberpage{XIV}

\inhaltpar{150}{Transitive and intransitive groups}{481}
\inhaltpar{151}{Primitive and imprimitive groups}{483}

\inhaltabschnitt{Fourteenth Section.}{Application of the permutation groups to equations.}
\inhaltpar{152}{Action of the permutation groups on functions of independent variables}{488}
\inhaltpar{153}{Decomposition of permutations into transpositions and cycles}{492}
\inhaltpar{154}{Divisors of the group, cosets and conjugate groups}{501}
\inhaltpar{155}{Reduction of the \gesperrt{Galois} resolvent by adjunction. Normal divisors of a group}{507}
\inhaltpar{156}{The group of the resolvents}{511}
\inhaltpar{157}{Reduction of the \gesperrt{Galois} group by adjunction of arbitrary irrationalities}{513}
\inhaltpar{158}{Imprimitive groups}{516}

\inhaltabschnitt{Fifteenth Section.}{Cyclic equations.}
\inhaltpar{159}{Cubic equations}{522}
\inhaltpar{160}{Permutation groups of four elements}{524}
\inhaltpar{161}{Solution of the biquadratic equations}{528}
\inhaltpar{162}{\gesperrt{Abel}ian equations}{533}
\inhaltpar{163}{Reduction of the \gesperrt{Abel}ian equations to cyclic ones}{537}
\inhaltpar{164}{Resolvents of \gesperrt{Lagrange}}{542}
\inhaltpar{165}{Solution of the cyclic equations}{546}
\inhaltpar{166}{Division of the angle}{551}

\inhaltabschnitt{Sixteenth Section.}{Cyclotomy.}
\inhaltpar{167}{The cyclotomic periods and the period equations}{554}
\inhaltpar{168}{\gesperrt{Gauss}'s method for the computation of the resolvents}{560}
\inhaltpar{169}{Reduction of the cyclotomic equation to pure equations. Division into seventeen parts}{564}
\inhaltpar{170}{Properties of the numbers $\psi$}{570}
\inhaltpar{171}{The \gesperrt{Gauss}ian sums}{574}
\inhaltpar{172}{The periods of ${}^{1}\!/_{3}\,(n-1)$ and ${}^{1}\!/_{4}\,(n-1)$ terms}{579}
\inhaltpar{173}{The complex numbers of \gesperrt{Gauss}}{585}
\inhaltpar{174}{The field of the third roots of unity}{592}

\inhaltabschnitt{Seventeenth Section.}{Algebraic solution of equations.}
\inhaltpar{175}{Reduction of the groups by pure equations}{595}
\inhaltpar{176}{Metacyclic equations}{597}

\clearpage
% ===== front matter, scan leaf 24 =====
% !PRINTED XV
\weberpage{XV}
\inhaltpar{177}{Simplicity of the alternating group}{600}
\inhaltpar{178}{Non-metacyclic equations in the field of rational numbers}{603}
\inhaltpar{179}{Solution by real radicals}{606}
\inhaltpar{180}{Metacyclic equations of prime degree}{609}
\inhaltpar{181}{Application to the metacyclic equations of the 5\textsuperscript{th} degree}{621}
\inhaltpar{182}{The group of the resolvent}{627}

\inhaltabschnitt{Eighteenth Section.}{Roots of metacyclic equations.}

\inhaltpar{183}{Statement of the problem. Lemma}{630}
\inhaltpar{184}{Theorems on the resolvents}{633}
\inhaltpar{185}{Roots of metacyclic equations}{638}
\inhaltpar{186}{Removal of the restrictive assumptions}{641}
\inhaltpar{187}{Conditions of reality}{647}
\inhaltpar{188}{Metacyclic equations of the 5\textsuperscript{th} degree}{648}

\clearpage
\pagenumbering{arabic}

% ===== scan leaf 26, printed 1 =====
% !PRINTED 1
\weberpage{1}
\pdfbookmark[1]{Introduction}{bm1}
\begin{center}
{\large\gesperrt{INTRODUCTION.}}

\rule{0.18\linewidth}{0.4pt}
\end{center}

We presuppose, in our considerations, the \gesperrt{natural numbers}
$1,2,3,\ldots$ and the rules according to which calculations are made
with these numbers as known and given. The fundamental arithmetical
operations, the so-called \gesperrt{four species}, are \gesperrt{addition},
\gesperrt{multiplication}, to which, as repetition, \gesperrt{exponentiation}
belongs, \gesperrt{subtraction}, and \gesperrt{division}. The first two are
called the \gesperrt{direct operations of calculation}; they are
distinguished by the fact that in the realm of the natural numbers they
can be carried out without restriction. The first of the indirect or
inverse operations, subtraction, can be carried out only when the
minuend is greater than the subtrahend.

The problem of division can be understood in two ways. In the first
elementary interpretation one asks how often the divisor is contained
in the dividend. A number is not contained in a smaller one. But if the
dividend is equal to or greater than the divisor, then the answer to the
question gives a \gesperrt{quotient} and, in most cases, a remainder
which is smaller than the divisor. If no remainder is left, one says
that the division comes out exactly, or that the \gesperrt{dividend} is
divisible by the \gesperrt{divisor}, or that the divisor is a
\gesperrt{factor} or \gesperrt{divisor} of the number which forms the
dividend.

This problem, which already presents itself at the first stages of the
art of calculation, leads to a more deeply lying distinction among
numbers, which is the foundation of all number theory.

Since a factor can never be greater than the number of which it is a
factor, every number has only a finite number of factors. Every number
is divisible by $1$ and by itself, and a number which has no other
divisor is called a \gesperrt{prime-}
%CONT>>


% ===== scan leaf 27, printed 2 =====
% !PRINTED 2
\weberpage{2}
%<<CONT
\gesperrt{number}. The number $1$ itself is customarily not designated
a prime number, for reasons of expediency. If two numbers are divisible
by a third, then the sum and the difference of the first two are also
divisible by the third; and if one number is divisible by a second and
this second by a third, then the first is also divisible by the third.

Two numbers always have the common divisor $1$. If they have no other
common divisor, as for example the pairs of numbers $5$ and $7$, or
$21$ and $38$, then the two numbers are called \gesperrt{relative prime
numbers}, or also \gesperrt{numbers prime to one another}.

Among the common divisors of any two given numbers one will be the
greatest, and it is a very important problem to find this greatest
common divisor. This is accomplished by a procedure known under the
name \gesperrt{algorithm of the greatest common divisor}, which already
occurs in Euclid.\footnote{Elements, Book VII, II, vol. II of Heiberg's
edition.}% Source prints Roman II (proposition 2), not Arabic 11.

Let $a,a_1$ be the two numbers; take the larger of them, say $a$, as
dividend and the smaller $a_1$ as divisor, and determine the quotient
$q_1$ and, if the division does not come out exactly, the remainder
$a_2$, so that
\[
 a=q_1a_1+a_2,
\]
and $a_2<a_1$. Every common divisor of $a$ and $a_1$ is then also a
common divisor of $a_1$ and $a_2$, and conversely. If one proceeds with
$a_1,a_2$ just as with $a,a_1$, and puts, when the division does not
come out exactly,
\[
 a_1=q_2a_2+a_3,
\]
then again $a_3<a_2$, and every common divisor of $a_1$ and $a_2$ is
also a common divisor of $a_2$ and $a_3$, and conversely. If one
continues dividing in this way, then, since the numbers
$a,a_1,a_2,a_3,\ldots$ continually decrease, the division must
necessarily come out exactly after a finite number of steps; that is,
at last a pair of equations
\[
 a_{\nu-2}=q_{\nu-1}a_{\nu-1}+a_\nu,
 \qquad
 a_{\nu-1}=q_\nu a_\nu
\]
must occur. Then $a_\nu$ is a common divisor of all preceding $a$'s,
hence also of $a$ and $a_1$, and every common divisor of $a$ and $a_1$
is a divisor of $a_\nu$. Thus $a_\nu$ is the greatest common divisor of
$a$ and $a_1$, and at the same time we have arrived at the theorem that
every common divisor of two numbers must divide their greatest common
divisor.


% ===== scan leaf 28, printed 3 =====
% !PRINTED 3
\weberpage{3}

If $a$ and $a_1$ are relatively prime numbers, then the last divisor is
$a_\nu=1$. Multiplying, under this assumption, the preceding equations
by an arbitrary number $b$, it follows that $b$ is the greatest common
divisor of $ab$ and $a_1b$, and hence that every common divisor of $ab$
and $a_1$, or of $a$ and $a_1b$, must be a divisor of $b$. Thus if $a_1$
is relatively prime to $a$ and to $b$, it is also relatively prime to
$ab$; and from the special assumption that $a_1$ is a prime number it
follows that a product can be divisible by a prime number only when at
least one of the factors is divisible by it.

Thus, if the product $ab$ is divisible by $a_1$ and $a_1$ is relatively
prime to $a$, then $b$ must be divisible by $a_1$.

Let $a,b$ be any two numbers with greatest common divisor $d$, and let
$a=da'$, $b=db'$. Then $a'$ and $b'$ are relatively prime to one
another. Every number $m$ which is at once a multiple of $a$ and of $b$
has the form $m=am'=da'm'$, and in this $a'm'$ must be a multiple of
$b'$; therefore $m'$ too must be a multiple of $b'$. That is, every
number which is at once a multiple of $a$ and of $b$ is divisible by
$a'b'd$. This number $a'b'd$, which itself is a common multiple of $a$
and $b$, is therefore called the \gesperrt{least common multiple} of
$a$ and $b$.

If a number is divisible by two relatively prime numbers, then it is
also divisible by their product; and if, therefore, a number $m$ is
divisible by several numbers of which each two are relatively prime to
one another, then it is also divisible by the product of all these
numbers.

On this rests the proof of the important theorem that a number $m$ can
always be represented, and in only one way, as a product of prime
numbers. For since a divisor cannot be greater than the dividend, $m$
can certainly be divisible by only a finite number of prime numbers. If
$a$ is one of these prime numbers and $a^\alpha$ the highest power of
$a$ which divides $m$, then $\alpha$ too is a definite number. Let
likewise $b^\beta,c^\gamma,\ldots$ be the highest powers of the remaining
prime numbers occurring in $m$, namely $b,c,\ldots$, by which $m$ can
be divided. Then, since any two of the
numbers $a^\alpha,b^\beta,c^\gamma,\ldots$ are relatively prime, $m$
must also be divisible by the product
$a^\alpha b^\beta c^\gamma\cdots$; and $m$ must be equal to this product,
since otherwise another
%CONT>>


% ===== scan leaf 29, printed 4 =====
% !PRINTED 4
\weberpage{4}
%<<CONT
prime number, or a higher power of one of the prime numbers
$a,b,c,\ldots$, would have to divide $m$.

\begin{center}
\rule{0.18\linewidth}{0.4pt}
\end{center}

We now give an overview of the extensions of the concept of number
which are necessary in mathematics and which were gradually introduced.

By a \gesperrt{manifold or set}, or by the abbreviated sign
$\mathfrak M$, we understand a system of objects or elements of any
kind, which is so delimited and complete in itself that, for every
arbitrary object, it is completely determined whether it belongs to
the system or not, regardless of whether we are able, in each
particular case, actually to make the decision or not.

A set is called \gesperrt{ordered} when, of any two different elements,
one that is fully determined in itself always counts as the greater,
and in such a way that from $a>b$, $b>c$ one always has $a>c$. If
$a>b$ and $b>c$, or $a>b>c$, we say that $b$ lies between $a$ and
$c$.% [source-original sic] The print repeats the descending case; the logically expected ascending alternative is not silently substituted.

The natural numbers form an ordered set; between two successive
elements of it there lies no further element. Such a manifold is called
\gesperrt{discrete}. An ordered set with the property that between any
two elements further elements can always be found is called
\gesperrt{dense}. One can form a dense set by combining the natural
numbers in pairs and regarding these pairs as elements of a set. These
pairs shall be called \gesperrt{fractions} and denoted by $m:n$ or
$m/n$, and two such fractions $m:n$ and $m':n'$ are set equal when
$mn'=nm'$. If all fractions equal to one another are combined into one
element, one obtains a manifold which is ordered once one further
stipulates that $m:n$ is greater than $m':n'$ when $mn'>nm'$. That this
manifold is dense is seen as follows: if $\mu=m:n$, $\mu'=m':n'$ are
two fractions and $\mu>\mu'$, then, with $h$ arbitrary, one can put
\[
 \mu={hmn'\over hnn'},
 \qquad
 \mu'={hm'n\over hnn'}
\]
%CONT>>


% ===== scan leaf 30, printed 5 =====
% !PRINTED 5
\weberpage{5}
%<<CONT
and therein $hmn'>hm'n$. One can always take $h$ so large that
further numbers lie between $hmn'$ and $hm'n$, and if $p$ is such a
number, then $p:hnn'$ lies between $\mu$ and $\mu'$.

The points of a straight line can also be regarded as an ordered set
if by greater and smaller one understands some indication of position,
for instance farther right and farther left, or higher and lower.

\begin{center}
\rule{0.18\linewidth}{0.4pt}
\end{center}

A division of an ordered set $\mathfrak M$ into two parts $A,B$ of
such a kind that every element $a$ of $A$ is smaller than every
element $b$ of $B$ is called a \gesperrt{cut} in $\mathfrak M$ and is
suitably denoted by $(A,B)$. Such a cut arises when one selects an
element $\mu$ of $\mathfrak M$, assigns all smaller elements to $A$,
all larger elements to $B$, and assigns $\mu$ itself at will to $A$
or to $B$. Strictly speaking, according as one does the one or the
other, two cuts arise, but we shall always regard them as equal. If,
in a cut $(A,B)$, either $A$ contains a greatest element or $B$ a
least element $\mu$, then we say that $\mu$ produces the cut $(A,B)$.
But it may also occur that neither $A$ has a greatest nor $B$ a least
element.

If every cut in a dense set is produced by a definite element $\mu$,
then the set is called \gesperrt{continuous}.

Continuity as well as density are properties which, by the nature of
the matter, are inaccessible to our sense perception; consequently
they also cannot be demonstrated with rigor in things of the external
world, in spatial magnitudes, intervals of time, or masses, however
much they may seem to lie in the essence of our intuition. But pure
conceptual systems can very well be constructed to which density
without continuity, or also density and continuity, belong\footnotemark[1].
\footnotetext[1]{This was proved by \gesperrt{Dedekind}, to whom we owe
in general the definition of continuity given above. Compare the
writings of \gesperrt{Dedekind}, \emph{Stetigkeit und irrationale
Zahlen}, Brunswick 1872, 1892, and \emph{Was sind und was sollen die
Zahlen?}, Brunswick 1888, 1893. Other}

The \gesperrt{rational fractions} furnish an example of a dense set.
This manifold, which we shall denote by $\mathfrak R$,
%CONT>>


% ===== scan leaf 31, printed 6 =====
% !PRINTED 6
\weberpage{6}
%<<CONT
is not continuous. For take any rational fraction $\mu=m:n$, in which
$m$ and $n$ have no common divisor and are not both square numbers.
Then $\mu$ is not the square of a rational fraction. For if
$\mu=p^2:q^2$, then $mq^2=np^2$ would follow, and from this $m=p^2$,
$n=q^2$, contrary to the assumption. Therefore if we form a cut
$(A,B)$ in the manifold $\mathfrak R$ by assigning to $A$ every
element $a$ of $\mathfrak R$ whose square is smaller than $\mu$, and
to $B$ every element $b$ whose square is greater than $\mu$, then
neither $A$ contains a greatest element nor $B$ a least element, and
the cut $(A,B)$ is not produced by an element of $\mathfrak R$.

For suppose that $a=p:q$ is any element of $A$, so that
$p^2:q^2<m:n$, or $np^2<mq^2$. Let $y$ be any natural number and
choose another natural number $x$ so that $x>y$ and
$x(mq^2-np^2)>ny(2p+1)$. Then
\[
 x^2(mq^2-np^2)>nxy(2p+1)>n(2pxy+y^2),
\]
and hence also
\[
 mq^2x^2>n(px+y)^2.
\]
Thus, if we put $a'=(px+y):qx$, then $a'>a$ and $a'^2<\mu$;
hence $a'$ too is contained in $A$. In exactly the same way one can
show that there is no least element in $B$.

The manifold $\mathfrak R$ can, however, serve as our starting point
for constructing a continuous set. The totality of all cuts in
$\mathfrak R$ is certainly a manifold, which may be denoted by
$\mathfrak S$. Consider two different elements of it,
$\alpha=(A,B)$ and $\alpha'=(A',B')$. Then either $A$ is a part of
$A'$ or $A'$ is a part of $A$. For if some

\par\vfill
\noindent\rule{0.22\linewidth}{0.4pt}

{\footnotesize\noindent
manifolds to which continuity belongs were formed by
\gesperrt{Weierstrass} and G.\ \gesperrt{Cantor}.

The definition of continuity, as we here lay it down after
\gesperrt{Dedekind}, is exhaustive in the following respect: a set
continuous in this sense, if it also possesses the property of
measurability presently to be discussed, cannot be a part of a richer
continuous set. I do not know whether this property has already been
proved anywhere, and hope to return to it on another occasion. I
remark, however, that such a property can be proved only for measurable
sets. A merely ordered set can always, however dense it may be, be
regarded as a part of one still denser.}
%CONT>>


% ===== scan leaf 32, printed 7 =====
% !PRINTED 7
\weberpage{7}
%<<CONT
element $a$ belongs to $A$, then every smaller element of $\mathfrak R$
also belongs to $A$. If $A$ is a part of $A'$, we shall call $\alpha$
smaller than $\alpha'$, and thereby the set $\mathfrak S$ has become
ordered.

If we regard the cuts produced by rational fractions as equivalent to
those rational fractions themselves and call them, briefly,
\gesperrt{rational cuts}, then the set $\mathfrak S$ contains the set
$\mathfrak R$, and $\mathfrak S$ is therefore in any case a
\gesperrt{dense set}. But the set $\mathfrak S$ is also
\gesperrt{continuous}. For let us denote the cuts by the large German
letters $\mathfrak A,\mathfrak B,\ldots$, and consider any cut in the
manifold of cuts, $(\mathfrak A,\mathfrak B)$. We can determine an
element of $\mathfrak S$, say $\alpha=(A,B)$, by putting into $A$ every
rational fraction that produces one of the cuts of $\mathfrak A$, and
putting into $B$ all the remaining rational fractions, which therefore
produce the rational cuts in $\mathfrak B$. This cut $\alpha$ in
$\mathfrak R$ produces the cut $(\mathfrak A,\mathfrak B)$ in
$\mathfrak S$. This will be proved when it is shown that every element
$\alpha'$ of $\mathfrak S$ which is smaller than $\alpha$ belongs to
$\mathfrak A$, and every element $\beta'$ of $\mathfrak S$ which is
greater than $\alpha$ belongs to $\mathfrak B$.

Let $\alpha'=(A',B')$ and $\alpha'<\alpha$. Then there are rational
fractions in $A$ which are not contained in $A'$. There is therefore a
rational $\mu$ such that $\alpha'<\mu<\alpha$; this $\mu$ produces a
cut contained in $\mathfrak A$, hence itself belongs to $\mathfrak A$.
Since $\alpha'<\mu$, $\alpha'$ also belongs to $\mathfrak A$. In
exactly the same way one shows that every $\beta'$ greater than
$\alpha$ belongs to $\mathfrak B$, and the continuity of $\mathfrak S$
is thereby proved.

This very abstract way of looking at the matter gives us the assurance
that the assumption of a continuous set contains no contradiction, and
that such sets at least exist in the realm of thought. Geometry, and
also analysis, which has always willingly attached itself to
geometrical intuition, long tacitly assumed the \gesperrt{existence} of
continuous sets, for example in the points of a straight line or of any
other connected line-segment, as a kind of axiom. The
\gesperrt{difference between a dense and a continuous set}, which lies
at the basis of the distinction between commensurable and
incommensurable segments, was also not missed by the ancients.
\setcounter{footnote}{0}\footnote{Euclid, Elements, Book X.}


% ===== scan leaf 33, printed 8 =====
% !PRINTED 8
\weberpage{8}

We too shall in what follows not renounce the aid of geometrical
intuition, and shall, for example, without hesitation regard the points
of a straight line as a continuous set.

\begin{center}
\rule{0.18\linewidth}{0.4pt}
\end{center}

An ordered set $\mathfrak M$ is called \gesperrt{measurable} under the
following assumptions: addition and multiplication are generally
executable in $\mathfrak M$, and likewise the subtraction of a smaller
element from a greater one; that is, from any two elements $a,b$ (which
may also be identical) a new element $a+b$ of $\mathfrak M$ can be
derived according to a definite rule, in such a way that $a+b$ is
greater than $a$ and greater than $b$, and the familiar laws of addition
expressed in the formulae
\[
 a+b=b+a,
 \qquad (a+b)+c=a+(b+c)
\]
hold. Moreover, for two elements $a,c$, of which the second is the
greater, a third element $b$ can be found such that $a+b=c$, which is
also expressed by the sign of subtraction $b=c-a$. Thus two unequal
elements of $\mathfrak M$ always have a definite difference. From these
assumptions it follows that a sum becomes greater when one of its
summands becomes greater. The repeated, say $m$-fold, addition of the
same element $a$ is called multiplication, and its result is denoted by
$ma$.

Finally there is added one more assumption, namely that, for every
given $a$, a sufficiently high multiple $ma$ is greater than any other
element $b$ given in advance. Among the elements of a measurable set
there is therefore no greatest element.

In a dense measurable set there is also no least element. For if $a$
were the least and $b$ an arbitrary element, then between $b$ and $b+a$
no elements could lie; for if $b<c<b+a$, then it would follow from the
definition of measurability that $a'=c-b$ would be smaller than $a$.
Conversely it also follows that a measurable set in which no least
element occurs is dense. For if $a$ and $a+b$ are any two elements, one
need only add to $a$ an element smaller than $b$ in order to obtain an
element between $a$ and $a+b$.

If a set is continuous, then the assumptions for measurability can be
simplified still further, because then subtraction is a consequence of
addition. For if $a$ and $c$ are two ele-
%CONT>>


% ===== scan leaf 34, printed 9 =====
% !PRINTED 9
\weberpage{9}
%<<CONT
ments of a continuous ordered set $\mathfrak M$ in which addition
exists, and if $c>a$, then one obtains a cut $(A,B)$ in $\mathfrak M$
by assigning to $A$ all elements $x$ for which $a+x\le c$, and to $B$
those for which $a+x>c$. This cut is produced by an element $b$, for
which $a+b=c$ must hold.

The natural numbers form, according to our definition, a measurable set
in which a least element, namely $1$, occurs. The manifold of rational
fractions is likewise measurable when addition and subtraction are
explained according to the known rules of fractional calculation.
Especially important, and as it were typical for the measurable sets,
is the manifold of rectilinear segments, or lengths of a line, which
are added simply by laying them end to end. Quantities of material,
compared by means of the balance, and intervals of time, measured by
the clock, also furnish examples of measurable sets. The manner of
measuring does not lie in the nature of the manifold itself, but is put
into it by the thinking observer; thus, for example, it would be just as
admissible to understand by the sum $a+b$ of two segments $a$ and $b$
the hypotenuse of a right triangle with legs $a,b$, instead of, as is
usually assumed, the segment composed from $a$ and $b$ by laying them
end to end.

In order to make the continuous set $\mathfrak S$ of cuts in the
manifold $\mathfrak R$ of rational fractions into a measurable one,
observe first the following.

If $\alpha=(A,B)$ is an element in $\mathfrak S$ and $\mu$ is an
arbitrarily given rational fraction, then one can always choose an
element $a'$ in $A$ so that $a'+\mu=b'$ is contained in $B$. For if one
chooses two arbitrary elements $a,b$, then one can determine the
natural number $m$ so that $m\mu>b-a$, hence $a+m\mu$ is contained in
$B$. If then $h$ is the smallest integer for which $a+h\mu$ is
contained in $B$, then $a+(h-1)\mu=a'$ is in $A$ and hence
$a'+\mu=b'$ is in $B$.

We now understand by the sum $(A,B)+(A',B')=(A'',B'')$, or
$\alpha+\alpha'=\alpha''$, the cut in $\mathfrak R$ which is obtained
by assigning a rational fraction $a''$ to $A''$ only when an $a$ in
$A$ and an $a'$ in $A'$ exist such that $a''\le a+a'$. In fact
$(A'',B'')$ is a cut; for if $a''$ is contained in $A''$, the same
holds for every smaller fraction, and there are fractions which are in
$A''$ and fractions which are not in
%CONT>>


% ===== scan leaf 35, printed 10 =====
% !PRINTED 10
\weberpage{10}
%<<CONT
$A''$, namely fractions of the forms $a+a'$ and $b+b'$. Further,
$\alpha''$ is greater than $\alpha$ and than $\alpha'$. For $A''$
first contains all of $A$, and if $a'$ is an arbitrary fraction in
$A'$, one can choose $a$ in $A$ so that the element $a+a'$ of $A''$ is
contained in $B$; hence $A''$ is more comprehensive than $A$. The cuts
produced by rational fractions $\mu,\mu'$ give, by addition, the cut
produced by $\mu+\mu'$.

\begin{center}
\rule{0.18\linewidth}{0.4pt}
\end{center}

We now pass to the definition of \gesperrt{ratios}, which since
antiquity have been regarded as the foundation of the theory of
numbers, and here follow Euclid first of all.
\setcounter{footnote}{0}\footnote{Elements, Book V.}

If one combines the elements of a measurable set $\mathfrak M$ into
pairs, and regards these pairs themselves as elements, then a new
manifold arises. We denote such a pair by $a:b$, or also by $a/b$; but,
when $a$ and $b$ are different elements, we distinguish $a:b$ from
$b:a$, and call $a$ the \gesperrt{numerator} and $b$ the
\gesperrt{denominator} of $a:b$. We shall call these pairs ratios, and
shall now order and make measurable this new set.

Assume first that two whole numbers $m,n$ exist such that $na=mb$, as
is always the case, for example, when $\mathfrak M$ is the system of
natural numbers, or when $a,b$ are two commensurable segments. Then,
if $p,q$ are two other whole numbers, $qa=pb$ holds if and only if
$mq=np$. The pair of numbers $p,q$ is completely determined by this
requirement if the condition is added that $p,q$ are to be relatively
prime. Then, if $h$ is an arbitrary whole number, $m=hp$, $n=hq$. In
this case we call the ratio $a:b$ \gesperrt{rational} and set it equal
to the rational fraction $m:n$ or $p:q$. Accordingly these rational
fractions may be regarded as ratios of whole numbers.

All rational ratios equal to one another form a \gesperrt{rational
number}, and the rational numbers form, like the rational fractions,
an ordered, dense, and measurable manifold.


% ===== scan leaf 36, printed 11 =====
% !PRINTED 11
\weberpage{11}

The natural numbers themselves are inserted into the manifold of
rational numbers if by the natural number $m$ one understands the
ratio $m:1$.

We now return to any measurable manifold $\mathfrak M$ and take from
it any two elements $a$ and $b$. If one chooses, which is always
possible, two natural numbers $m,n$ so that $na>mb$, then the ratio
$a:b$ is called greater than the rational ratio $m:n$, or
\[
 {a\over b}>{m\over n},
\]
and if $m:n>p:q$, then also $a:b>p:q$. Similarly, if $n'a<m'b$, then
\[
 {a\over b}<{m'\over n'}.
\]

If $a:b>m:n$, then one can insert one, and consequently also
arbitrarily many, rational numbers $m_1:n_1$ such that
\[
 {a\over b}>{m_1\over n_1}>{m\over n}
\]
hold. That is, one can insert arbitrarily many rational ratios
\gesperrt{between} $a:b$ and $m:n$. To show this, choose an arbitrary
whole number $k$ and determine the whole number $h$ so that
$h(na-mb)>kb$, which is always possible. Then
\[
 {a\over b}>{hm+k\over hn}>{m\over n}.
\]
And in exactly the same way, if $n'a<m'b$ and
$h'(m'b-n'a)>k'a$, one obtains
\[
 {a\over b}<{h'm'\over h'n'+k'}<{m'\over n'}.
\]

If now $a:b$ and $\alpha:\beta$ are any two ratios, which for brevity
we shall also denote by $e,\epsilon$, and whose elements belong to the
same or to different manifolds, then two cases are possible: 1) no
rational ratio $\mu$ lies between $e$ and $\epsilon$, or 2) a rational
ratio lies between $e$ and $\epsilon$.

In case 1) the two ratios $e$ and $\epsilon$ are called
\gesperrt{equal to one another}. One sees that two ratios which are
equal to a third are also equal to one another; for if
$e<\mu<\epsilon$, then every other ratio $e'$ is either smaller
%CONT>>


% ===== scan leaf 37, printed 12 =====
% !PRINTED 12
\weberpage{12}
%<<CONT
than, equal to, or greater than $\mu$. If $e'$ is equal to $\mu$, then
rational ratios lie both between $e$ and $e'$ and between $e'$ and
$\epsilon$. If, however, $e'$ is smaller than $\mu$, then $\mu$ lies
between $e'$ and $\epsilon$, and if $e'$ is greater than $\mu$, then
$\mu$ lies between $e$ and $e'$; hence $e'$ cannot at the same time be
equal to $e$ and equal to $\epsilon$.

In case 2) the ratios $e,\epsilon$ are called \gesperrt{unequal}. Thus
either $\alpha)$ $e<\mu<\epsilon$ or $\beta)$
$e>\mu'>\epsilon$ may hold, and these two cases exclude one another,
because from $e<\mu<\epsilon$, $\epsilon<\mu'$ it follows that
$\mu<\mu'$, and consequently $e<\mu'$ results.

The relation of magnitude $2\alpha)$ or $2\beta)$ also remains
unchanged when $e$ or $\epsilon$ is replaced by an element equal to it.
For if $\mu<\epsilon$ and $\mu\gtrless\epsilon'$, then a rational ratio
lies between $\epsilon$ and $\epsilon'$, and $\epsilon,\epsilon'$ are
not equal.

In case $2\alpha)$, $e$ is called \gesperrt{smaller than} $\epsilon$;
in case $2\beta)$, $e$ is called \gesperrt{greater than} $\epsilon$.
Between two unequal ratios one can insert an arbitrary number of
rational ratios.

If now we collect all ratios equal to one another, we obtain a generic
concept which we designate as a \gesperrt{number} in the general sense
of the word. A number is therefore a name or sign for a certain
manifold whose elements are precisely the ratios equal to one of them.
\setcounter{footnote}{0}\footnote{The natural numbers too can be traced back in a simple and
consistent way to the generic concept.} Under this concept of number
are included the rational ratios, and hence also the natural numbers as
the ratios $m:1$; they form the \gesperrt{rational numbers}. Numbers
which do not arise from rational ratios are called \gesperrt{irrational
numbers}. According to what has been developed so far, the numbers form
an \gesperrt{ordered set}, and their order can be fixed if, for each
number, any one of the ratios contained under it is chosen as
representative.

\begin{center}
\rule{0.18\linewidth}{0.4pt}
\end{center}

Of two ratios with the same denominator and unequal numerators, that
one is greater whose numerator is greater; and
%CONT>>


% ===== scan leaf 38, printed 13 =====
% !PRINTED 13
\weberpage{13}
%<<CONT
of two ratios with the same numerator and unequal denominators, that
one is smaller whose denominator is greater. For if $a,a',b$ are
arbitrary elements of a measurable set and $a'>a$, choose first a whole
number $n$ such that $na>b$ and $n(a'-a)>b$. Then take $m$ to be the
least whole number satisfying $mb>na$; then $na<mb$, but $mb<na'$.
% [source-original sic] The print gives mb\ge na'\ge na+n(a'-a); the second relation is algebraically equality and is preserved verbatim.
For if $mb\ge na'\ge na+n(a'-a)$, then $mb>na+b$; hence, contrary to
the supposition, already $(m-1)b>na$. Thus
\[
 {a\over b}<{m\over n}<{a'\over b},
\]
and hence $a:b<a':b$. In exactly the same way one proves that if
$b'>b$, then $a:b>a:b'$. This theorem is also expressed by saying that
a ratio increases with the numerator and decreases with increasing
denominator.

From this one easily obtains the following theorem. If $a,b,c,d$ are
elements of the same measurable set and $a:b=c:d$, then also
$a:c=b:d$. For suppose that $a:c<b:d$; then there would have to be two
whole numbers $m,n$ such that
\[
 na<mc,
 \qquad nb>md.
\]
Then, however, by the theorem just proved, $na:nb<mc:md$, and hence
$a:b<c:d$, contrary to the supposition.

\gesperrt{To this is now attached the following principal theorem. If
of four quantities} $a,b,c,d$ \gesperrt{any three are given arbitrarily
from a continuous measurable set, then the fourth can be determined in
the same set so that} $a:b=c:d$.

The theorem is an immediate consequence of the assumed continuity. For
if, for example, an element $x$ of $\mathfrak M$ is put into $A$ or
into $B$ according as $x:b$ is smaller or greater than $c:d$, then a
cut is obtained, which is produced by an element $a$ satisfying
$a:b=c:d$. But this theorem also holds in certain non-continuous sets,
for example for the rational fractions.

It follows that as representatives of two numbers one can always
choose two ratios whose elements belong to the same
%CONT>>


% ===== scan leaf 39, printed 14 =====
% !PRINTED 14
\weberpage{14}
%<<CONT
manifold and have the same arbitrarily chosen denominator. Addition is
then defined so that
\[
 {a\over c}+{b\over c}={a+b\over c}
\]
holds.
This rule contains as a special case the addition of rational
fractions. To justify it generally it remains only to show that, if
$a:c=a':c'$ and $b:c=b':c'$, then also
$(a+b):c=(a'+b'):c'$ must hold. We prove this by showing that, if
$a:c=a':c'$ and $(a+b):c>(a'+b'):c'$, then also $b:c>b':c'$ must
hold. Let therefore, with $m,n$ two whole numbers,
\[
 {a+b\over c}>{m\over n}>{a'+b'\over c'},
\]
then $n(a+b)>mc$, and hence
\[
 {b\over c}>{mc-na\over nc},
 \qquad {mc'-na'\over nc'}>{b'\over c'}.
\]
On the other hand, from the supposition $a:c=a':c'$ it follows easily
that
\[
 {mc-na\over nc}={mc'-na'\over nc'}
\]
holds, and therefore $b:c>b':c'$, as was required to be proved.

It is thus shown that the numbers, as we have defined them, also form a
measurable set. If $a$ and $c$ are taken from a continuous manifold,
then, with fixed $c$, the ratios $a:c$ also form a continuous set;
hence, since continuous sets exist at all, it follows that the numbers
too form a continuous set.

If $\alpha,\beta,\gamma,\delta$ are now numbers, then from the
proportion $\alpha:\beta=\gamma:\delta$ any one of the four numbers
can be determined by the three others given. If one puts $\delta=1$
and seeks $\alpha$, one obtains the \gesperrt{multiplication}
$\alpha=\beta\gamma$; the commutativity of the factors is a consequence
of the theorem that from $\alpha:\beta=\gamma:\delta$ there follows
$\alpha:\gamma=\beta:\delta$. If one seeks $\gamma$, one obtains
division; and from the definition of addition given above
there follows the fundamental formula
$\alpha(\beta+\gamma)=\alpha\beta+\alpha\gamma$. The four fundamental
operations are therefore executable in the domain of numbers, with the
sole restriction that in subtraction the subtrahend must be smaller
than the minuend.

\begin{center}
\rule{0.18\linewidth}{0.4pt}
\end{center}


% ===== scan leaf 40, printed 15 =====
% !PRINTED 15
\weberpage{15}

The construction of a cut in the sequence of numbers always supplies
the proof for the existence of a number satisfying prescribed
requirements. Thus one obtains a cut $(A,B)$ by putting into $A$ all
and only those numbers whose square is smaller than a prescribed number
$\alpha$; to this cut there corresponds a definite number whose square
is equal to $\alpha$ and which is denoted by $\sqrt\alpha$. In this way
the existence of square roots is proved.

The numerical sequences introduced by G.\ \gesperrt{Cantor} for the
definition of irrational numbers can also be referred back to cuts.
\setcounter{footnote}{0}\footnote{Cantor, \emph{Ueber die Ausdehnung eines Satzes aus der
Theorie der trigonometrischen Reihen}. Mathematische Annalen, vol.~5
(1872); compare also Heine, \emph{Elemente der Functionenlehre}.
Journal für Mathematik, vol.~74 (1872).}

According to \gesperrt{Cantor}, by a numerical sequence one is to
understand any unbounded system of numbers ordered in a definite way:
\[
 S=x_1,x_2,x_3,x_4,\ldots
\]
having the property that there is a definite number $g$ below which
none of the numbers of $S$ sinks, and that, if $\delta$ is any chosen
number and $x_n,x_m$ are different elements of $S$, then $x_n-x_m$ or
$x_m-x_n$ always remains smaller than $\delta$ when $m$ and $n$ have
exceeded a sufficiently large number.

There are always numbers which are not exceeded by the numbers of such
a sequence; for, once $\delta$ has been chosen and $n$ suitably
determined, $x_m$, whatever the value of $m$, never exceeds the
greatest of the numbers $x_1,x_2,\ldots,x_n,x_n+\delta$. One now
obtains a cut $(A,B)$ by throwing into $B$ those numbers which, when
$n$ has a sufficiently high value, are no longer exceeded by any
$x_n$, and all other numbers, which are therefore exceeded by
infinitely many $x_n$, into $A$. If this cut is produced by the number
$\alpha$, then, however small $\varepsilon$ may be, there are always
infinitely many numbers $x_n$ between $\alpha-\varepsilon$ and
$\alpha$; and one may say that this number $\alpha$, completely
determined by $S$, is produced by the numerical sequence $S$. According
to \gesperrt{Cantor}, the numerical sequence $S$ is directly the
definition of the number $\alpha$. Of course one and the same
%CONT>>


% ===== scan leaf 41, printed 16 =====
% !PRINTED 16
\weberpage{16}
%<<CONT
number $\alpha$ can be produced by very different numerical sequences.
These numerical sequences, however, are all to be regarded as equal to
one another. Among other things, one may take all numbers of $S$ to be
rational fractions. Conversely, it is also easily shown that for every
given number $\alpha$ one can always give numerical sequences $S$ by
which $\alpha$ is produced, so that the totality of the numerical
sequences $S$ likewise forms a continuous set.

\begin{center}
\rule{0.18\linewidth}{0.4pt}
\end{center}

In the explanation of the fundamental operations an inconvenient
restriction has appeared in subtraction, from which we can free
ourselves by introducing \gesperrt{zero and the negative numbers}.

Let $x$ denote any element of the number system defined so far, which
we shall now call the system of \gesperrt{positive numbers}. We take
this number system a second time and, to distinguish it, denote in this
second system, which is to be called the system of \gesperrt{negative
numbers}, every element by $-x$. This second system we now order in
exactly the reverse way from the first, so that wherever in the system
$x$ ``greater'' stands, in the system $-x$ ``smaller'' is put, and
conversely. Addition and subtraction in $-x$ are explained just as in
$x$, so that $(-x)+(-y)=-(x+y)$, $(-x)-(-y)=-(x-y)$ shall hold.

But we wish to order these two number systems together in such a way
that every $-x$ is to be smaller than every $x$. We thus obtain an
ordered set in which there is no greatest and no least element. This
system is also in general continuous; only the single cut $(-x,x)$ is
produced by no element, and here, where the two systems meet, a
violation of continuity is still present. To restore continuity we
must therefore add, corresponding to the cut $(-x,x)$, a number
\gesperrt{zero, or 0}, which is precisely defined by this cut. Then
we have an ordered continuous set unbounded on both sides, the complete
series of \gesperrt{real numbers}.

In this enlarged number domain we now explain addition generally by
putting, by definition:
%CONT>>


% ===== scan leaf 42, printed 17 =====
% !PRINTED 17
\weberpage{17}
%<<CONT
\[
 x+(-x)=0,
 \qquad x+0=x,
\]
\[
 x+(-y)=x-y,\quad \text{if }x>y,
 \quad =-(y-x),\quad \text{if }y>x.
\]
With this definition of addition, if $z_1,z_2,z_3$ are any three
numbers of the whole number domain, the laws
\[
 z_1+z_2=z_2+z_1,
 \qquad (z_1+z_2)+z_3=z_1+(z_2+z_3),
\]
hold; they are called the commutative and associative laws. The sum of
any number of summands is formed by successively combining any two
summands into a sum. Subtraction no longer needs separate
consideration if $z_1-z_2$ is defined as $z_1+(-z_2)$ and
$-(-z)=z$ is set.

The number line is represented intuitively by points: from a fixed
point of a straight line, denoted by $0$, the positive numbers are laid
off as segments on one side, say the right, and the negative numbers
on the other (left) side. The image of the sum of two segments
$z_1+z_2$ is obtained by laying off from the point $z_1$ the segment
of length $\pm z_2$ to the right or left, according as $z_2$ is
positive or negative.

Multiplication and division are introduced into the enlarged number
domain by the equations
\[
 x(-y)=(-x)y=-xy,
\]
\[
 (-x)(-y)=xy,
 \qquad 0x=0
\]
are defined. Division, as the inverse of multiplication, is always possible except
when the divisor is zero.

\begin{center}
\rule{0.18\linewidth}{0.4pt}
\end{center}

A further extension of the concept of number consists in introducing
\gesperrt{complex quantities}. We combine every two numbers of the
complete number line into pairs $(x,y)$, and regard two such pairs
$(x,y)$ and $(a,b)$ as equal only when $x=a$, $y=b$. These pairs of
numbers form a manifold whose elements cannot be ordered, but with
which addition, subtraction, multiplication, and division are to be
performed according to the following rules. Let
\[
 (x,y)+(a,b)=(x+a,y+b),
\]
\[
 (x,y)(a,b)=(xa-yb,xb+ya),
\]
%CONT>>


% ===== scan leaf 43, printed 18 =====
% !PRINTED 18
\weberpage{18}
%<<CONT
\begingroup
\setlength{\abovedisplayskip}{5pt plus 2pt minus 1pt}
\setlength{\belowdisplayskip}{5pt plus 2pt minus 1pt}
\setlength{\abovedisplayshortskip}{3pt plus 1pt}
\setlength{\belowdisplayshortskip}{4pt plus 1pt}
and we further stipulate that $(x,0)=x$, which does not contradict
these equations. The pair $(x,y)$ is $=0$ only when $x$ and $y$
are both zero. For abbreviation we denote $(0,1)$ by $i$. The
equations above then give the consequences:
\[
 (x,0)(0,1)=(0,x)\quad\text{or}\quad =ix,
\]
\[
 (x,0)+(0,y)=(x,y)=x+yi,
\]
\[
 x+yi+a+bi=(x+a)+(y+b)i,
\]
and the inverse of addition:
\[
 x+yi-(a+bi)=(x-a)+i(y-b),
\]
and further multiplication:
\[
 (x+yi)(a+bi)=xa-yb+i(xb+ya),
 \qquad i^2=-1,
\]
\[
 (x+yi)(x-yi)=x^2+y^2,
\]
\[
 x+yi=(a+bi){(a-bi)(x+yi)\over a^2+b^2}
\]
or
\[
 {x+yi\over a+bi}={ax+by+i(ay-bx)\over a^2+b^2},
\]
which defines division except when $a+bi=0$.

Numbers of the form $ix$ are called \gesperrt{purely imaginary
numbers}, and $i$ the imaginary unit. The numbers $a+bi$ are called
\gesperrt{imaginary} or \gesperrt{complex}. The systems of
\gesperrt{real} and \gesperrt{purely imaginary numbers} are contained
among them as special cases.

Thus, in the domain of \gesperrt{complex numbers} $x+yi$, the
fundamental operations can be carried out without restriction (except
for division by zero), and calculation with real numbers is a special
case.

Following \gesperrt{Gauss}, one represents the complex numbers
$z=x+yi$ geometrically by the points of a plane, laying down a
rectangular coordinate system and regarding the point with coordinates
$x,y$ as the image of the numerical value $z$. The points of the
$x$-axis represent the real numbers $x$ in the way discussed above.
The points of the $y$-axis are the images of the purely imaginary
numbers $yi$. The origin is the image of the number $0$. The radius
vector from the origin to the point $z$ has the numerical value
$\varrho=\sqrt{x^2+y^2}$, which is called the \gesperrt{absolute
value}, or the magnitude, or, according to older terminology, the
modulus of the complex number $z$.
\endgroup


% ===== scan leaf 44, printed 19 =====
% !PRINTED 19
\weberpage{19}

The only number having absolute value $0$ is $0$. Every positive
number occurs as the absolute value of infinitely many complex
numbers, and the image points of all numbers with the same absolute
value lie on a circle whose center is at the origin.

Two imaginary numbers which differ only by the sign of $i$, namely
$x+yi$ and $x-yi$, are called conjugate imaginary. Their product is
the square of the absolute value of either one.

If one of two conjugate imaginary numbers is zero, the other is also
zero; hence in every correct numerical equation one may replace $i$
by $-i$ without disturbing its correctness.

We shall also state the frequently used theorem that the absolute
value of the sum of two nonzero complex numbers is never greater than
the sum of the absolute values of the summands, and is equal to it only
when the ratio (the quotient) of the two summands is real and positive.
Namely, let
\[
 z=x+yi,
 \qquad c=a+bi,
 \qquad Z=(x+a)+(y+b)i,
\]
\[
 \varrho^2=x^2+y^2,
 \qquad r^2=a^2+b^2,
 \qquad R^2=(x+a)^2+(y+b)^2,
\]
then
\[
 (r+\varrho-R)(r+\varrho+R)
   =(r+\varrho)^2-R^2=2(r\varrho-ax-by);
\]
this is certainly positive when $ax+by\leqq 0$. If, however,
$ax+by>0$, then from
\[
 r^2\varrho^2-(ax+by)^2=(ay-bx)^2,
\]
it follows that $r\varrho\geqq ax+by$, with equality only when
$ay-bx=0$. Hence $r+\varrho>R$, except in the special case
$r+\varrho=R$ when $ay-bx=0$ and $ax+by>0$, which proves the
assertion.

In the geometric representation this theorem says that in a triangle
one side is smaller than the sum of the other two. The theorem has the
further consequence that the absolute value of a sum cannot be smaller
than the difference of the absolute values of the summands. For,
applying the preceding theorem to the sum $c=Z-z$, one obtains
$r\leqq R+\varrho$, or
\[
 R\geqq r-\varrho.
\]
%CONT>>


% ===== scan leaf 45, printed 20 =====
% !PRINTED 20
\weberpage{20}
%<<CONT
Equality occurs here only when the quotient $z:c$ is real and
negative.

\begin{center}
\rule{0.18\linewidth}{0.4pt}
\end{center}

It remains to say a word about the most important aid of algebra,
literal calculation. This aid is so generally used that the term
literal calculation is sometimes employed outright as a synonym for
algebra. We presuppose the rules for calculating with such literal
expressions. Equations between literal expressions can be of two
kinds: either they are so-called identities, that is, the two
expressions set equal can be transformed by the rules of calculation
until they agree exactly. Such literal equations yield correct
numerical equations when the letters are replaced by any real or
complex numbers, provided that division by zero is not required. The
letters in such equations are often called variables, because one may
imagine, without ever reaching a contradiction, that different
numerical values are successively substituted for them.

Another kind of equation between literal expressions does not have the
character of an identity. It instead contains a requirement imposed on
the numbers that may be substituted for the letters without error.
Algebra has the task of determining numerical values that satisfy such
a requirement, that is, of solving the equation. In these equations
the letters are also called ``\gesperrt{unknowns}.'' Very often one
and the same equation contains letters of two kinds: some for which
arbitrary numerical values are to be substituted, and others whose
numerical value is to be determined.

\begin{center}
\rule{0.24\linewidth}{0.4pt}
\end{center}


% ===== scan leaf 46, printed 21 =====
% !PRINTED 21
\weberpage{21}
\thispagestyle{empty}
\vspace*{0.31\textheight}
\begin{center}
{\Large\gesperrt{FIRST BOOK.}}

\vspace{1.8em}
\rule{0.14\linewidth}{0.4pt}

\vspace{2.6em}
{\LARGE\gesperrt{THE FOUNDATIONS.}}

\vspace{3.5em}
\rule{0.24\linewidth}{0.4pt}
\end{center}
\vfill


% ===== scan leaf 47, printed 22 =====
% !PRINTED 22
\weberpage{22}
\thispagestyle{empty}
\vspace*{\fill}
\mbox{}
\vspace*{\fill}


% ===== scan leaf 48, printed 23 =====
% !PRINTED 23
\weberpage{23}
\pdfbookmark[0]{First Book. The Foundations.}{bm2}
\pdfbookmark[1]{First Section. Rational Functions.}{bm3}
\pdfbookmark[2]{§ 1. Integral functions}{bm4}
\thispagestyle{empty}
\begin{center}
{\large\gesperrt{First Section.}}

\vspace{0.8em}
{\Large\gesperrt{Rational Functions.}}

\vspace{0.8em}
\rule{0.16\linewidth}{0.4pt}

\vspace{0.8em}
{\large \S.\ 1.}

\vspace{0.5em}
{\large\gesperrt{Whole Functions.}}
\end{center}

The next object of consideration is \gesperrt{whole rational}
functions, or briefly \gesperrt{whole functions of one variable}. By
these we understand expressions of the following form:
\begin{equation}
 f(x)=a_0x^n+a_1x^{n-1}+a_2x^{n-2}+\cdots+a_{n-1}x+a_n,
\tag{1}
\end{equation}
where $n$ is an integral exponent, the \gesperrt{degree} of the function $f(x)$.
The degree is a \gesperrt{natural number}. Sometimes, however, it is
also useful to speak of whole rational functions of degree 0, by
which a constant is understood. $x$ is called the \gesperrt{variable},
$a_0,a_1,a_2,\ldots,a_{n-1},a_n$ the \gesperrt{coefficients}. Both $x$
and $a_0,a_1,\ldots,a_n$ are symbols for indeterminate quantities, with
which one proceeds according to the rules of literal calculation, and
for which definite numerical values may also be substituted under
some circumstances (cf. the Introduction).

If the function $f(x)$ is written in the way shown in (1), we call it
\gesperrt{ordered according to descending powers of $x$}. The summands
may be arranged in any other order; for example, the function may also
be ordered according to ascending powers of $x$:
\[
 f(x)=a_n+a_{n-1}x+\cdots+a_1x^{n-1}+a_0x^n.
\]
By addition (subtraction) and multiplication of whole rational
functions there again arise whole rational functions. The rules of
literal calculation directly give the laws of formation of these new
functions.

In addition, the coefficient of any power $x^\nu$ in the sum is equal
to the sum of the coefficients of $x^\nu$ in the individual summands.
The degree of the function that arises is
%CONT>>


% ===== scan leaf 49, printed 24 =====
% !PRINTED 24
\weberpage{24}
%<<CONT
\begingroup
\setlength{\abovedisplayskip}{5pt plus 2pt minus 1pt}
\setlength{\belowdisplayskip}{5pt plus 2pt minus 1pt}
\setlength{\abovedisplayshortskip}{3pt plus 1pt}
\setlength{\belowdisplayshortskip}{4pt plus 1pt}
\linespread{0.94}\selectfont
equal to the highest of the degrees of the summands, and can decrease
only in the special case in which the highest degree occurs in several
summands and the sum of the coefficients of the highest powers is
zero.

In the multiplication of two or more whole rational functions there
arises a function whose degree is equal to the sum of the degrees of
the factors.

To survey the law of formation of the coefficients for the product,
put
\begin{equation}
\begin{aligned}
 A(x)&=a_0x^m+a_1x^{m-1}+a_2x^{m-2}+\cdots,\\
 B(x)&=b_0x^n+b_1x^{n-1}+b_2x^{n-2}+\cdots,
\end{aligned}
\tag{2}
\end{equation}
\begin{equation}
 A(x)B(x)=C(x)=c_0x^{m+n}+c_1x^{m+n-1}+c_2x^{m+n-2}+\cdots
\tag{3}
\end{equation}
and obtain
\[
\begin{aligned}
 c_0&=a_0b_0,\\
 c_1&=a_0b_1+a_1b_0,\\
 c_2&=a_0b_2+a_1b_1+a_2b_0,\\
 &\cdots,
\end{aligned}
\]
or, in general notation, if $\nu$ denotes one of the numbers
$0,1,2,\ldots,m+n$:
\begin{equation}
 c_\nu=a_0b_\nu+a_1b_{\nu-1}+a_2b_{\nu-2}+\cdots
       +a_{\nu-1}b_1+a_\nu b_0,
\tag{4}
\end{equation}
where all $a$ whose index is greater than $m$, and all $b$ whose index
is greater than $n$, are to be set equal to zero. For $c_\nu$ is the
coefficient of $x^{m+n-\nu}$, and hence $c_\nu x^{m+n-\nu}$ arises by
multiplying all terms of the form $a_\mu x^{m-\mu}$ by all terms of the
form $b_{\nu-\mu}x^{n-\nu+\mu}$ and then adding. If $a_0$ and
$b_0=1$---that is, if both are equal to one---then also $c_0=1$.%
% The source's compressed “a_0 and b_0=1” is retained symbol-for-symbol,
% while the translation makes its mathematically intended joint scope explicit.

\gesperrt{If in all factors of such a product the highest power of
$x$ has coefficient $1$, then in the product too the coefficient of
the highest power is $=1$.}

From these rules for calculating with whole functions it follows that
the rules of calculation expressed in formulas such as $ab=ba$,
$(ab)c=a(bc)$, $(a+b)c=ac+bc$ and the like also hold when $a,b,c$ are
whole functions of $x$, and in the sense that whole functions are to
be regarded as equal only when equal powers have equal coefficients.

By whole functions of several variables we understand sums of products
of whole, positive powers of the variables $x,y,z\ldots$ with arbitrary
coefficients that count as constant. One can think of these functions
as arising
\endgroup
%CONT>>


% ===== scan leaf 50, printed 25 =====
% !PRINTED 25
\weberpage{25}
\pdfbookmark[2]{§ 2. A theorem of Gauss}{bm5}
%<<CONT
from functions of one variable $x$ when the coefficients
$a_0,a_1\ldots a_n$ occurring in them are themselves regarded as
whole rational functions of other variables $y,z\ldots$. Thus
functions of $m$ variables arise from functions of $m-1$ variables.
All terms of such a function are of the form $x^r y^s z^t\cdots$,
multiplied by a coefficient; and if the function is properly
collected, then each combination of the exponents $r,s,t\ldots$
occurs only once. Two whole functions so ordered are regarded as equal
only when they contain the same products $x^r y^s z^t\cdots$ with the
same coefficients.

\begin{center}
{\large \S.\ 2.}

\vspace{0.5em}
{\large\gesperrt{A Theorem of Gauss.}}
\end{center}

We shall at once apply the multiplication rule for two whole rational
functions in order to prove a theorem of Gauss which will be useful to
us later, but which here is to serve as an introduction to the mode of
calculation and as an example.\footnote{Gauss, \emph{Disquisitiones
arithmeticae}, Art.~42.}

We consider here the case in which the coefficients in the functions
$A(x),B(x)$ are whole numbers, so that, by (4), the coefficients of
$C(x)=A(x)B(x)$ are also whole numbers.

If all integral coefficients $a_0,a_1,\ldots,a_m$ of a function
$A(x)$ have no common divisor, then the function $A(x)$ is called an
\gesperrt{original or primitive} function, and the theorem we wish to
prove is:

\gesperrt{If $A(x)$ and $B(x)$ are primitive functions, then their
product $C(x)$ is also a primitive function.}

The proof follows almost immediately from looking at formula (4).

For if all coefficients $c_0,c_1,c_2\ldots c_{m+n}$, as given in
(4), had a common divisor greater than $1$, then there would also have
to be at least one prime number which divides all these coefficients.
Let $p$ be such a prime number. By the assumption that $A(x),B(x)$ are
primitive, this prime can divide neither all $a_0,a_1\ldots a_m$ nor
all $b_0,b_1\ldots b_n$.


% ===== scan leaf 51, printed 26 =====
% !PRINTED 26
\weberpage{26}

Suppose now that $p$ divides
\[
 a_0,a_1\ldots a_{r-1},\quad \hbox{but not }a_r,
\]
and divides
\[
 b_0,b_1\ldots b_{s-1},\quad \hbox{but not }b_s,
\]
Then, by assumption, $r\leqq m$, and may also be zero if $p$ already
does not divide $a_0$. Likewise $s\leqq n$.

If now, by (4), we form $c_{r+s}$ and arrange it in the following way,
\begin{equation}
\begin{aligned}
 c_{r+s}={}&a_rb_s+a_{r-1}b_{s+1}+a_{r-2}b_{s+2}+\cdots\\
            &+a_{r+1}b_{s-1}+a_{r+2}b_{s-2}+\cdots,
\end{aligned}
\tag{1}
\end{equation}
one sees at once that $c_{r+s}$ cannot be divisible by $p$, as had
been assumed; for the first term $a_rb_s$ is not divisible by $p$,
while all other terms, since they are multiplied by one of the
coefficients $a_0,a_1\ldots a_{r-1},b_0,b_1\ldots b_{s-1}$, are
divisible by $p$. Thus the assumption that $C(x)$ is not primitive
while $A(x)$ and $B(x)$ are primitive leads to a contradiction.

This theorem can be transferred to functions of several variables. We
call a whole rational function of $m$ variables with integral
coefficients \gesperrt{original or primitive} if the coefficients have
no common divisor, and we state the theorem that \gesperrt{the product
of two primitive functions is again a primitive function}.

To see its truth, we need only take, in the argument carried out above,
the coefficients $a_0,a_1,a_2\ldots$, $b_0,b_1,b_2\ldots$ not as whole
numbers but as whole rational functions of $m-1$ variables
$y,z\ldots$, and call such a function divisible by a prime number $p$
when all its coefficients are divisible by $p$. If we then suppose
that the theorem to be proved has already been proved for functions of
$m-1$ variables, formula (1) gives its correctness for functions of
$m$ variables, hence its general validity by complete induction, or by
the conclusion from $m-1$ to $m$.

An imprimitive function is one whose integral coefficients all have a
common divisor. We call this divisor the \gesperrt{divisor of the
function}. Then we can also express the theorem just proved as follows:
%CONT>>


% ===== scan leaf 52, printed 27 =====
% !PRINTED 27
\weberpage{27}
\pdfbookmark[2]{§ 3. Division}{bm6}
%<<CONT
\gesperrt{The divisor of a product of two whole functions is equal to
the product of the divisors of the two functions.}

For if $PA$ and $QB$ are two whole functions with divisors $P$ and
$Q$, then $A$ and $B$ are primitive functions. Hence $AB=C$ is also a
primitive function, and $PQ$ is the divisor of the function
$PA\cdot QB=PQC$.

Insofar as it refers to functions of one variable, we can give the
theorem the following form, without changing its content essentially;
in this form it is especially useful. Let
\[
 \varphi(x)=x^m+\alpha_1x^{m-1}+\alpha_2x^{m-2}+\cdots+\alpha_m,
\]
\[
 \psi(x)=x^n+\beta_1x^{n-1}+\beta_2x^{n-2}+\cdots+\beta_n
\]
\gesperrt{be two whole rational functions in which the highest powers
of $x$ have coefficient $1$, while the remaining coefficients are
rational numbers. Then, in the product}
\[
 \varphi(x)\psi(x)=x^{m+n}+\gamma_1x^{m+n-1}
 +\gamma_2x^{m+n-2}+\cdots+\gamma_{m+n}
\]
\gesperrt{the coefficients $\gamma$ cannot all be whole numbers if the
coefficients $\alpha,\beta$ in $\varphi(x)$ and $\psi(x)$ are not all
whole numbers.}

For let the smallest common denominator of the coefficients $\alpha$
of $\varphi$ be denoted by $a_0$, and that of the coefficients $\beta$
of $\psi$ by $b_0$. Then $a_0\varphi(x)=A(x)$ and
$b_0\psi(x)=B(x)$ are primitive functions of $x$. If the coefficients
$\gamma$ were whole numbers, the product $a_0b_0\varphi(x)\psi(x)$
would have divisor $a_0b_0$ and, if $a_0$ and $b_0$ were not both
equal to $1$, would therefore not be primitive; but this would
contradict the theorem proved above.

\begin{center}
{\large \S.\ 3.}

\vspace{0.5em}
{\large\gesperrt{Division.}}
\end{center}

Let, as before,
\begin{equation}
\begin{aligned}
 A&=A(x)=a_0x^m+a_1x^{m-1}+\cdots,\\
 B&=B(x)=b_0x^n+b_1x^{n-1}+\cdots
\end{aligned}
\tag{1}
\end{equation}
be two whole rational functions of $x$; but now suppose
%CONT>>


% ===== scan leaf 53, printed 28 =====
% !PRINTED 28
\weberpage{28}
%<<CONT
\begingroup
\setlength{\abovedisplayskip}{5pt plus 2pt minus 1pt}
\setlength{\belowdisplayskip}{5pt plus 2pt minus 1pt}
\setlength{\abovedisplayshortskip}{3pt plus 1pt}
\setlength{\belowdisplayshortskip}{4pt plus 1pt}
\linespread{0.94}\selectfont
that $m\geqq n$ and that $a_0$ and $b_0$ are different from zero. Then
the difference
\begin{equation}
 A-{a_0\over b_0}x^{m-n}B
\tag{2}
\end{equation}
is also a whole rational function of $x$, but its degree is
\gesperrt{smaller} than $m$, since the highest power in the two terms
of the difference has the same coefficient and therefore cancels. We
put this difference equal to
\begin{equation}
 A'=A'(x)=a'_0x^{m'}+a'_1x^{m'-1}+\cdots,
 \qquad m'<m.
\tag{3}
\end{equation}
If now $m'$ is still $\geqq n$, then in (2) we can put $A'$ in place of
$A$ and repeat the same reasoning.

This gives a chain of equations:
\begin{equation}
\begin{aligned}
 A-{a_0\over b_0}x^{m-n}B&=A',\\
 A'-{a'_0\over b_0}x^{m'-n}B&=A'',\\
 A''-{a''_0\over b_0}x^{m''-n}B&=A''',\\
 &\cdots,
\end{aligned}
\tag{4}
\end{equation}
and this chain can be continued until the degree of the function that
has arisen is smaller than $n$. Since in the sequence of functions
$A,A',A''\ldots$ the degree of each following function is lowered by
\gesperrt{at least} one unit, the chain of equations (4) consists of
\gesperrt{at most} $m-n+1$ members; but it may also contain fewer
members when \gesperrt{several} powers cancel simultaneously. If we add
all equations (4), denote the last of the functions $A,A'\ldots$ by
$C$, and, for abbreviation, set
\begin{equation}
 Q={a_0\over b_0}x^{m-n}+{a'_0\over b_0}x^{m'-n}+\cdots,
\tag{5}
\end{equation}
so that $Q$ also is a whole rational function of $x$, then there
follows
\begin{equation}
 A=QB+C.
\tag{6}
\end{equation}
The operation just described, by which the functions $Q,C$ are found
from $A,B$, is called \gesperrt{division}. $A$ is the
\gesperrt{dividend}, $B$ the \gesperrt{divisor}, $C$ the
\gesperrt{remainder}, and $Q$ the \gesperrt{quotient}. \gesperrt{The
degree of the remainder is always lower than the degree of the
divisor.}
\endgroup


% ===== scan leaf 54, printed 29 =====
% !PRINTED 29
\weberpage{29}

The coefficients of the functions $Q$ and $C$ are composed from the
coefficients $a$ and $b$ by addition, subtraction, multiplication,
and division. In the denominator, however, only powers of $b_0$
occur; hence, if $b_0=1$, the coefficients of $Q$ and $C$ are whole
functions of the $a$ and $b$. The highest power of $b_0$ that can
occur in the denominator is the $(m-n+1)$st, since in the chain (4)
the factor $b_0$ is added once in the denominator in each succeeding
equation. In special cases, however, the highest power of $b_0$ in
all denominators may be lower.

Take, for example, for $A$ a function of the third degree,
\begin{equation}
 f(x)=a_0x^3+a_1x^2+a_2x+a_3
\tag{7}
\end{equation}
and for $B$ the so-called first derivative of $f(x)$, which is of the
second degree,
\begin{equation}
 f'(x)=3a_0x^2+2a_1x+a_2,
\tag{8}
\end{equation}
and hence:
\begin{equation}
 Q={1\over 3}x+{a_1\over 9a_0},
\tag{9}
\end{equation}
\begin{equation}
 C={6a_0a_2-2a_1^2\over 9a_0}x
   +{9a_0a_3-a_1a_2\over 9a_0}.
\tag{10}
\end{equation}

How the calculation prescribed in equations (4) is expediently
arranged may here be assumed as known from the elements. We draw
attention to the analogy between this calculation and division in the
decimal number system. A function $f(x)$ represents a number written
decimally if the coefficients $a_0,a_1,\ldots$ are whole numbers
between zero (inclusive) and $10$ (exclusive), and $x=10$ is set. If
numbers greater than $10$ are also permitted among the coefficients,
a number can be represented by $f(x)$ in different ways. This is used
in the division procedure to avoid fractional and negative
coefficients in the results in the simplest way.


% ===== scan leaf 55, printed 30 =====
% !PRINTED 30
\weberpage{30}
\pdfbookmark[2]{§ 4. Division by a linear function}{bm7}

\begingroup
\setlength{\abovedisplayskip}{5pt plus 2pt minus 1pt}
\setlength{\belowdisplayskip}{5pt plus 2pt minus 1pt}
\setlength{\abovedisplayshortskip}{3pt plus 1pt}
\setlength{\belowdisplayshortskip}{4pt plus 1pt}
\linespread{0.94}\selectfont

\begin{center}
{\large \S.\ 4.}

\vspace{0.5em}
{\large\gesperrt{Division by a Linear Function.}}
\end{center}

We shall apply the general rule for division to one special case. Let
the dividend be
\begin{equation}
 f(x)=a_0x^n+a_1x^{n-1}+a_2x^{n-2}+\cdots+a_n
\tag{1}
\end{equation}
arbitrary, while the divisor is of the first degree, or, as one also
says, linear. We shall also suppose the coefficient of the first power
of $x$ in the divisor to be $=1$, and therefore take the divisor in the
simple form $(x-\alpha)$, where $\alpha$ remains arbitrary. In this
case the quotient $Q$ is of degree $(n-1)$, and the remainder $C$ is of
degree $0$, that is, independent of $x$. Thus set
\begin{equation}
 f(x)=(x-\alpha)Q+C,
\tag{2}
\end{equation}
so that $C$ no longer contains $x$; and if we set
\begin{equation}
 Q=q_0x^{n-1}+q_1x^{n-2}+\cdots+q_{n-2}x+q_{n-1}
\tag{3}
\end{equation}
then there follows from (2):
\begin{equation}
\begin{aligned}
 f(x)={}&q_0x^n+q_1x^{n-1}+\cdots+q_{n-2}x^2+q_{n-1}x+C\\
 &-\alpha q_0x^{n-1}-\cdots-\alpha q_{n-3}x^2
   -\alpha q_{n-2}x-\alpha q_{n-1},
\end{aligned}
\tag{4}
\end{equation}
and by comparison with (1):
\begin{equation}
\begin{aligned}
 q_0&=a_0,\\
 q_1-\alpha q_0&=a_1,\\
 q_2-\alpha q_1&=a_2,\\
 &\cdots,\\
 q_{n-1}-\alpha q_{n-2}&=a_{n-1},\\
 C-\alpha q_{n-1}&=a_n.
\end{aligned}
\tag{5}
\end{equation}
From this one obtains
\begin{equation}
\begin{aligned}
 q_0&=a_0,\\
 q_1&=a_0\alpha+a_1,\\
 q_2&=a_0\alpha^2+a_1\alpha+a_2,\\
 &\cdots,\\
 q_{n-1}&=a_0\alpha^{n-1}+a_1\alpha^{n-2}+\cdots+a_{n-1},\\
 C&=a_0\alpha^n+a_1\alpha^{n-1}+\cdots+a_{n-1}\alpha+a_n=f(\alpha).
\end{aligned}
\tag{6}
\end{equation}
$C$ arises from $f(x)$ by setting $x=\alpha$, and therefore can also
be denoted by $f(\alpha)$.
\endgroup


% ===== scan leaf 56, printed 31 =====
% !PRINTED 31
\weberpage{31}

Accordingly we also have the formula
\begin{equation}
 {f(x)-f(\alpha)\over x-\alpha}=Q(x),
\tag{7}
\end{equation}
where $Q(x)$ is a whole function of degree $n-1$.

If in expressions (6) we put the sign $x$ in place of the indeterminate
quantity $\alpha$, there arises a sequence of whole rational functions
of $x$, which, if for simplicity we set $a_0=1$, we write as follows:
\begin{equation}
\begin{aligned}
 f_0&=1,\\
 f_1&=x+a_1,\\
 f_2&=x^2+a_1x+a_2,\\
 &\cdots,\\
 f_{n-1}&=x^{n-1}+a_1x^{n-2}+a_2x^{n-3}+\cdots+a_{n-1}.
\end{aligned}
\tag{8}
\end{equation}
These functions $f_0,f_1,\ldots,f_{n-1}$ will later render us good
service. For the present we add the following remarks.

According to (8), the powers $1,x,x^2,\ldots,x^{n-1}$ of $x$ can be
expressed linearly in terms of the functions
$f_0,f_1,\ldots,f_{n-1}$, in such a way that only whole rational
combinations of the $a$ occur in the coefficients, for example
\[
\begin{aligned}
 1&=f_0,\\
 x&=f_1-a_1f_0,\\
 x^2&=f_2-a_1f_1+(a_1^2-a_2)f_0,\\
 &\cdots,
\end{aligned}
\]
and it follows that every whole rational function of $x$ whose degree
is not greater than $n-1$ can likewise be expressed linearly in terms
of $f_0,f_1,\ldots,f_{n-1}$ in the form
\begin{equation}
 y_0f_0+y_1f_1+\cdots+y_{n-1}f_{n-1},
\tag{9}
\end{equation}
where the coefficients $y_0,y_1,\ldots,y_{n-1}$ are independent of
$x$.

Thus, if $F(x)$ is any whole rational function of $x$, then by \S~3,
regarding $f(x)$ as divisor, one can set
\begin{equation}
 F(x)=Qf(x)+y_0f_0+y_1f_1+\cdots+y_{n-1}f_{n-1}
\tag{10}
\end{equation}
where $Q$ is also a whole rational function of $x$.

For the recursive calculation of the functions $f_r(x)$, relation (8)
gives
\begin{equation}
 f_r(x)-xf_{r-1}(x)=a_r.
\tag{11}
\end{equation}


% ===== scan leaf 57, printed 32 =====
% !PRINTED 32
\weberpage{32}
\pdfbookmark[2]{§ 5. Fractional functions; divisibility}{bm8}

\begin{center}
{\large \S.\ 5.}

\vspace{0.5em}
{\large\gesperrt{Fractional Functions; Divisibility.}}
\end{center}

If $F(x)$ and $f(x)$ are two whole rational functions of $x$, then
the fraction
\begin{equation}
 {F(x)\over f(x)}
\tag{1}
\end{equation}
is called a \gesperrt{fractional rational}, or briefly a
\gesperrt{fractional} or \gesperrt{rational function} of $x$.

If the degree of the numerator is lower than the degree of the
denominator, the function is called a \gesperrt{proper fraction}; in
the opposite case it is called an \gesperrt{improper fraction}.

According to \S~3, the whole rational functions $Q$ and $\varphi(x)$
can be determined so that
\begin{equation}
 F(x)=Qf(x)+\varphi(x),
\tag{2}
\end{equation}
and the degree of $\varphi(x)$ is smaller than the degree of $f(x)$;
consequently
\begin{equation}
 {F(x)\over f(x)}=Q+{\varphi(x)\over f(x)}
\tag{3}
\end{equation}
and hence the theorem:

\gesperrt{Every fractional function can be decomposed into the sum of
a whole function and a proper fractional function.}

If $n$ is the degree of $f(x)$, then the degree of $\varphi(x)$ is at
most $n-1$; it may also be lower. In particular, the case may occur in
which $\varphi(x)$ \gesperrt{vanishes identically}.

\gesperrt{In this case the function $F(x)$ is said to be divisible by
$f(x)$.} The function
\[
 {F(x)\over f(x)}
\]
is in this case only apparently fractional; in reality it is equal to
the whole function $Q$.

The formula
\[
 {x^m-1\over x-1}=x^{m-1}+x^{m-2}+x^{m-3}+\cdots+1
\]
gives a simple example.


% ===== scan leaf 58, printed 33 =====
% !PRINTED 33
\weberpage{33}

For divisibility of functions the same laws hold as for divisibility of numbers, in particular the following:

\gesperrt{1. If the function $F(x)$ is divisible by the function $f(x)$, and $f(x)$ is divisible by a third function $\varphi(x)$, then $F(x)$ is also divisible by $\varphi(x)$.}

For if
\[
 F=Qf,
 \qquad
 f=q\varphi,
\]
where $Q$ and $q$ are whole rational functions, then
\[
 F=Qq\varphi,
\]
and since $Qq$ is a whole rational function, $F$ is divisible by $\varphi$.

\gesperrt{2. If $F(x)$ is divisible by $f(x)$ and $Q$ is any whole rational function, then $QF(x)$ is also divisible by $f(x)$.}

\gesperrt{3. If $F(x)$ and $f(x)$ are divisible by $\varphi(x)$, then $F(x)\pm f(x)$ is also divisible by $\varphi(x)$, or more generally:}

\gesperrt{4. If $F_1,F_2\ldots$ are divisible by $f(x)$ and $Q_1,Q_2\ldots$ are arbitrary whole rational functions, then}
\[
 Q_1F_1+Q_2F_2+\cdots
\]
\gesperrt{is also divisible by $f(x)$.}

The last theorem includes the two preceding ones and is proved simply as follows.

If $F_1,F_2\ldots$ are divisible by $f$, then the whole rational functions $\Phi_1,\Phi_2\ldots$ can be determined so that
\[
 F_1=\Phi_1f,
 \qquad
 F_2=\Phi_2f,\ldots
\]
and consequently
\[
 Q_1F_1+Q_2F_2+\cdots=(Q_1\Phi_1+Q_2\Phi_2+\cdots)f.
\]
Since $Q_1\Phi_1+Q_2\Phi_2+\cdots$ is a whole rational function, the theorem is proved.

\gesperrt{5. Every function is divisible by itself.}

With respect to divisibility or indivisibility of functions nothing is changed if the functions are multiplied by arbitrary factors independent of $x$.

A non-zero quantity independent of $x$ may be regarded as a function of degree zero. If we call such a quantity a constant, we may say:

\gesperrt{6. Every function is divisible by every constant.}


% ===== scan leaf 59, printed 34 =====
% !PRINTED 34
\weberpage{34}
\pdfbookmark[2]{§ 6. Greatest common divisor}{bm9}

If a function is divisible by another, then the degree of the quotient is equal to the difference between the degree of the dividend and the degree of the divisor. The former therefore cannot be smaller than the latter. If the degrees are equal, the quotient is a constant, and hence:

\gesperrt{7. If, of two whole rational functions of the same degree, one is divisible by the other, then they differ from one another only by a constant factor, and the second is also divisible by the first.}

\gesperrt{8. By \S~4, the necessary and sufficient condition that the function $f(x)$ be divisible by the linear function $x-\alpha$ is that $f(\alpha)=0$.}

\begin{center}
{\large \S.\ 6.}

\vspace{0.5em}
{\large\gesperrt{Greatest Common Divisor.}}
\end{center}

It is a problem of fundamental importance to decide whether two whole rational functions have, besides the constants, any further common divisor. The solution is found by the \gesperrt{algorithm of the greatest common divisor} in exactly the same way as the corresponding question concerning divisibility of integers is answered. See the Introduction.

Let
\begin{equation}
 f(x)=A,
 \qquad
 \varphi(x)=A'
\tag{1}
\end{equation}
be two given functions; the degree of $\varphi(x)$ is supposed to be lower, or at least not higher, than that of $f(x)$.

We divide $A$ by $A'$ and denote the remainder, whose degree is lower than the degree of $A'$, by $A''$, thus:
\[
 A=Q'A'+A'';
\]
then we divide $A'$ by $A''$ and denote the remainder by $A'''$, and continue in this way in forming the sequence of functions
\begin{equation}
 A,A',A'',A''',\ldots,
\tag{2}
\end{equation}
whose degrees $n,n',n'',n'''\ldots$ always decrease and therefore, after a finite number of divisions, necessarily go down to zero. Let the last remainder of degree zero, and therefore a constant, be $A^{(\nu)}$. Then we have the chain of equations:
%CONT>>


% ===== scan leaf 60, printed 35 =====
% !PRINTED 35
\weberpage{35}
%<<CONT

\begin{equation}
\begin{aligned}
A&=Q'A'+A'',\\
A'&=Q''A''+A''',\\
&\cdots\\
A^{(\nu-3)}&=Q^{(\nu-2)}A^{(\nu-2)}+A^{(\nu-1)},\\
A^{(\nu-2)}&=Q^{(\nu-1)}A^{(\nu-1)}+A^{(\nu)},
\end{aligned}
\tag{3}
\end{equation}
where $A^{(\nu)}$ is a constant.

If now $A,A'$ have any common divisor, then by the theorems of the preceding paragraph, as is seen from the equations (3), it is also a divisor of $A'',A''',A''''$ and so on down to $A^{(\nu)}$. If $A^{(\nu)}$ is a non-zero constant, then no divisor of $f(x)$ and $\varphi(x)$ depending on $x$ can exist. Such functions are called \gesperrt{prime to one another}, or \gesperrt{relatively prime}. One also says, taking account only of divisors depending on $x$, that the functions have no common divisor.

The condition that the two functions have a common divisor is therefore
\begin{equation}
 A^{(\nu)}=0.
\tag{4}
\end{equation}
In this case $A^{(\nu-1)}$ is a divisor of $A^{(\nu-2)}$, as the last equation (3) shows, and by the preceding equation it is a divisor of $A^{(\nu-3)}$, and so on, hence also a divisor of $A$ and $A'$, that is, of $f$ and $\varphi$. Conversely, every common divisor of $A,A'$ is also a divisor of $A^{(\nu-1)}$; accordingly $A^{(\nu-1)}$ is called the \gesperrt{greatest common divisor} of $f$ and $\varphi$. The algorithm (3) shows that $A^{(\nu)}$ or $A^{(\nu-1)}$ can be derived from the coefficients of $f$ and $\varphi$ by the \gesperrt{rational operations of arithmetic}, and in such a manner that division is always only by the coefficients of the highest powers of $x$ in the functions $A,A',A'',\ldots$ which are not zero.

We shall apply these considerations to two examples.

First let
\begin{equation}
\begin{aligned}
 A&=a_0x^2+a_1x+a_2,\\
 B&=b_0x^2+b_1x+b_2
\end{aligned}
\tag{5}
\end{equation}
be two functions of the second degree, so that $a_0,b_0$ are non-zero.


% ===== scan leaf 61, printed 36 =====
% !PRINTED 36
\weberpage{36}

The first step is to form the equation
\begin{equation}
 A={a_0\over b_0}B+C,
\tag{6}
\end{equation}
where
\begin{equation}
 C=c_0x+c_1
\tag{7}
\end{equation}
and
\begin{equation}
 c_0={a_1b_0-b_1a_0\over b_0},
 \qquad
 c_1={a_2b_0-b_2a_0\over b_0}.
\tag{8}
\end{equation}
If $c_0=0$, so that $C$ is constant, then $A,B$ have a common factor only if also $c_1=0$; and then $A$ is divisible by $B$, that is, $A$ and $B$ differ only by a constant factor. If however $c_0$ is different from zero, we continue the calculation by putting
\[
 B=QC+D
\]
where, according to \S~4, if there $\alpha=-c_1:c_0$ is put,
\begin{equation}
 D={b_0c_1^2-b_1c_0c_1+b_2c_0^2\over c_0^2}.
\tag{9}
\end{equation}
If $D$ is non-zero, then $A,B$ have no common divisor. If however $D=0$, then $C$ is the greatest common factor of $A$ and $B$.

Substituting the values $c_0,c_1$ from (8) into (9), and suppressing the denominator $b_0c_0^2$, one obtains the condition for the existence of a common divisor $C$ of $A$ and $B$ in the form
\begin{equation}
\begin{aligned}
&a_0^2b_2^2+a_2^2b_0^2-2a_0a_2b_0b_2-a_1a_2b_0b_1-a_0a_1b_1b_2\\
&\qquad{}+a_0a_2b_1^2+a_1^2b_0b_2=0
\end{aligned}
\tag{10}
\end{equation}
or
\begin{equation}
 (a_0b_2-b_0a_2)^2+(a_0b_1-a_1b_0)(a_2b_1-a_1b_2)=0.
\tag{11}
\end{equation}
This condition is also fulfilled when $c_0$ and $c_1$ are both zero, and is therefore the necessary and sufficient condition that $A,B$ have a common divisor.

The left hand side of (10) or (11) is called the \gesperrt{resultant of the functions $A$ and $B$}.

As a second example we take the one already chosen in \S~3. Put
\begin{equation}
\begin{aligned}
 f(x)&=a_0x^3+a_1x^2+a_2x+a_3=A,\\
 f'(x)&=3a_0x^2+2a_1x+a_2=B,
\end{aligned}
\tag{12}
\end{equation}
and suppose $a_0$ different from zero. Then by \S~3 we have:
%CONT>>


% ===== scan leaf 62, printed 37 =====
% !PRINTED 37
\weberpage{37}
%<<CONT

\begin{equation}
A=QB+C,
\tag{13}
\end{equation}
\begin{equation}
 C=c_0x+c_1,
\tag{14}
\end{equation}
where
\begin{equation}
 c_0={6a_0a_2-2a_1^2\over 9a_0},
 \qquad
 c_1={9a_0a_3-a_1a_2\over 9a_0}.
\tag{15}
\end{equation}
If $c_0$ is zero, the algorithm is already closed. If $c_1$ is non-zero, then $A$ and $B$ have no common divisor; but if $c_1$ is also zero, then $B$ itself is the greatest common divisor of $A$ and $B$, that is, $A$ is divisible by $B$. The conditions for this are therefore
\begin{equation}
 3a_0a_2-a_1^2=0,
 \qquad
 9a_0a_3-a_1a_2=0.
\tag{16}
\end{equation}
If $c_0$ is not zero, we go one step further and put
\begin{equation}
 B=PC+D,
\tag{17}
\end{equation}
where $D$ becomes constant and has the expression
\begin{equation}
 D={a_2c_0^2-2a_1c_0c_1+3a_0c_1^2\over c_0^2}.
\tag{18}
\end{equation}
If this expression is non-zero, then $A$ and $B$ are \gesperrt{prime to one another}; if it is zero, then $A$ and $B$ have the greatest common divisor $C$. If for $c_0,c_1$ one substitutes the values (15), omits the denominator, still removes the non-vanishing factor $9a_0$, and changes the sign, this condition, after a simple calculation, assumes the form
\begin{equation}
 a_1^2a_2^2+18a_0a_1a_2a_3-4a_0a_2^3-4a_1^3a_3-27a_0^2a_3^2=0.
\tag{19}
\end{equation}
It is, as one readily sees either by calculation or from (18), also fulfilled when the conditions (16) hold, and is therefore the necessary and sufficient condition that $f(x)$ and $f'(x)$ have a common divisor. The left hand side of (19), which is a whole rational and homogeneous function of the coefficients of $f(x)$, is called the \gesperrt{discriminant of the function $f(x)$}.

We derive one more theorem from the algorithm (3), which is often used.

From the first of these equations it follows that
\begin{equation}
 A''=A-Q'A',
\tag{20}
\end{equation}
and if this value of $A''$ is substituted in the following equation, one obtains
\[
 A'''=(1+Q'Q'')A'-Q''A,
\]
%CONT>>


% ===== scan leaf 63, printed 38 =====
% !PRINTED 38
\weberpage{38}
%<<CONT

thus, denoting by $p,p'$ whole rational functions,
\begin{equation}
 A'''=pA+p'A'.
\tag{21}
\end{equation}
Substituting the expressions (20), (21) in the third equation (3), one obtains for $A''''$ again an expression of the form (21), and so one may continue, obtaining finally
\begin{equation}
 A^{(\nu)}=PA+P'A',
\tag{22}
\end{equation}
where $P,P'$ are whole rational functions whose coefficients are composed from the coefficients of $A$ and $A'$ by rational operations.

In the formula (22), $A^{(\nu)}$ is a constant. This theorem is especially important in the case where $A^{(\nu)}$ is non-zero, so that $A,A'$ are relatively prime. If in this case we put
\[
 P=A^{(\nu)}F(x),
 \qquad
 P'=A^{(\nu)}\Phi(x),
\]
then, after suppressing the factor $A^{(\nu)}$, the theorem can be expressed as follows:

\gesperrt{I. If $f(x)$ and $\psi(x)$ are two whole functions without a common divisor, then two other whole functions $F(x)$ and $\Phi(x)$ can be determined which satisfy the equation}
\begin{equation}
 F(x)f(x)+\Phi(x)\psi(x)=1
\tag{23}
\end{equation}
\gesperrt{identically.}

The theorem can still be generalized. Multiplying (23) by an arbitrary whole function $\chi(x)$ gives
\begin{equation}
 F(x)\chi(x)f(x)+\Phi(x)\chi(x)\psi(x)=\chi(x),
\tag{24}
\end{equation}
and by \S~3 we can put
\begin{equation}
 \Phi(x)\chi(x)=Q(x)f(x)+\varphi(x)
\tag{25}
\end{equation}
in such a way that $Q(x),\varphi(x)$ are whole rational functions of $x$ and the degree of $\varphi(x)$ is lower than that of $f(x)$. Substituting this in (24), and putting in place of $F(x)\chi(x)+Q(x)\psi(x)$ again $F(x)$, we obtain
\[
 F(x)f(x)+\varphi(x)\psi(x)=\chi(x).
\]
We can therefore generalize the preceding theorem as follows:

\gesperrt{II. If $f(x),\psi(x),\chi(x)$ are given whole rational functions, and $f(x)$ and $\psi(x)$ are without a common divisor, then the}
%CONT>>


% ===== scan leaf 64, printed 39 =====
% !PRINTED 39
\weberpage{39}
\pdfbookmark[2]{§ 7. Products of linear factors}{bm10}
%<<CONT

\gesperrt{whole rational functions $F(x),\varphi(x)$ can be so determined, with $\varphi(x)$ of lower degree than $f(x)$, that the equation}
\begin{equation}
 F(x)f(x)+\varphi(x)\psi(x)=\chi(x)
\tag{26}
\end{equation}
\gesperrt{is identically satisfied.}

\begin{center}
{\large \S.\ 7.}

\vspace{0.5em}
{\large\gesperrt{Products of Linear Factors.}}
\end{center}

According to \S~1, by multiplication of whole functions we again obtain such functions of higher degree, and the degree of the product is determined by the sum of the degrees of the individual factors.

If therefore we multiply $n$ linear factors by one another, there arises a function of the $n$th degree, whose manner of formation we must examine somewhat more closely.

We shall take the linear factors in the simple form $x-\alpha_1$, $x-\alpha_2$, $\ldots$, $x-\alpha_n$, and put, since the coefficient of the highest power is (by \S~1) $=1$,
\begin{equation}
\begin{aligned}
 f(x)&=(x-\alpha_1)(x-\alpha_2)\cdots(x-\alpha_n)\\
 &=x^n+a_1x^{n-1}+a_2x^{n-2}+\cdots+a_n.
\end{aligned}
\tag{1}
\end{equation}
A slight consideration shows the following law of formation of the coefficients $a_1,a_2\ldots a_n$:

$-a_1$ is equal to the sum of the $\alpha$'s; $a_2$ is equal to the sum of the products of every two of the $\alpha$'s; $-a_3$ is the sum of the products of every three of the $\alpha$'s; in general $(-1)^\nu a_\nu$ is the sum of the products of every $\nu$ of the quantities $\alpha$, or, in formulas,
\begin{equation}
\begin{aligned}
-a_1&=\sum \alpha_1,\\
+a_2&=\sum \alpha_1\alpha_2,\\
&\cdots\\
(-1)^\nu a_\nu&=\sum \alpha_1\alpha_2\cdots\alpha_\nu,\\
&\cdots\\
(-1)^n a_n&=\alpha_1\alpha_2\cdots\alpha_n.
\end{aligned}
\tag{2}
\end{equation}
But to see still more clearly the correctness of this law of formation, one uses \gesperrt{complete induction}.

One first verifies the correctness in the first cases by actually carrying out the multiplication
%CONT>>


% ===== scan leaf 65, printed 40 =====
% !PRINTED 40
\weberpage{40}
%<<CONT

\[
(x-\alpha_1)(x-\alpha_2)=x^2-(\alpha_1+\alpha_2)x+\alpha_1\alpha_2,
\]
\[
\begin{aligned}
 (x-\alpha_1)(x-\alpha_2)(x-\alpha_3)&=x^3-(\alpha_1+\alpha_2+\alpha_3)x^2\\
&\quad +(\alpha_1\alpha_2+\alpha_1\alpha_3+\alpha_2\alpha_3)x-\alpha_1\alpha_2\alpha_3.
\end{aligned}
\]
If we assume the law of formation proved for $n-1$ factors, and put
\[
 (x-\alpha_1)(x-\alpha_2)\cdots(x-\alpha_{n-1})=x^{n-1}+a'_1x^{n-2}+\cdots+a'_{n-1},
\]
then multiplication by $x-\alpha_n$ gives
\begin{equation}
\begin{aligned}
 a_1&=a'_1-\alpha_n,\\
 a_2&=a'_2-\alpha_n a'_1,\\
 a_3&=a'_3-\alpha_n a'_2,\\
&\cdots\\
 a_\nu&=a'_\nu-\alpha_n a'_{\nu-1},\\
&\cdots\\
 a_n&=-\alpha_n a'_{n-1},
\end{aligned}
\tag{3}
\end{equation}
and from this the correctness of the formulas (2) is immediately evident.

It is important to determine the number of terms which occur in each of the sums (2). We denote the number of terms occurring in the sum $(-1)^\nu a_\nu$ by $B_\nu^{(n)}$; it is the number of combinations without repetition of $n$ elements to the $\nu$th class, that is, of every $\nu$ distinct ones among the $n$ elements. To determine it, suppose first that the $B_{\nu-1}^{(n)}$ combinations to the $(\nu-1)$st class have been formed. From each of these combinations, by adding one at a time of the missing $n-\nu+1$ elements, one can derive $n-\nu+1$ combinations to the $\nu$th class. In this way however every combination is formed $\nu$ times, namely by addition of each of its elements; hence one obtains the relation
\begin{equation}
 B_\nu^{(n)}=B_{\nu-1}^{(n)}{n-\nu+1\over \nu}
\tag{4}
\end{equation}
whereas $B_1^{(n)}$ plainly has the value $n$. Therefore, if one multiplies the equations following from (4),
\begin{equation}
\begin{aligned}
 B_1^{(n)}&=n,\\
 B_2^{(n)}&=B_1^{(n)}{n-1\over 2},\\
&\cdots\\
 B_\nu^{(n)}&=B_{\nu-1}^{(n)}{n-\nu+1\over \nu}
\end{aligned}
\tag{5}
\end{equation}
it follows that:
%CONT>>


% ===== scan leaf 66, printed 41 =====
% !PRINTED 41
\weberpage{41}
%<<CONT

\begin{equation}
B_\nu^{(n)}={n(n-1)(n-2)\cdots(n-\nu+1)\over 1\cdot2\cdot3\cdots\nu}.
\tag{6}
\end{equation}
In what follows we shall often, when $m$ is an arbitrary positive integer, use the symbol
\begin{equation}
 \Pi(m)=1\cdot2\cdot3\cdots m,
 \qquad
 \Pi(0)=1
\tag{7}
\end{equation}
so that
\begin{equation}
 \Pi(m)=m\Pi(m-1).
\tag{8}
\end{equation}
With the help of this symbol the expression for $B_\nu^{(n)}$ can be displayed more conveniently as
\begin{equation}
 B_\nu^{(n)}=B_{n-\nu}^{(n)}={\Pi(n)\over \Pi(\nu)\Pi(n-\nu)},
\tag{9}
\end{equation}
which, if $B_0^{(n)}=1$ is put, is also valid for $\nu=0$ and $\nu=n$, and makes evident the invariance of $B_\nu^{(n)}$ under the interchange of $\nu$ with $n-\nu$.

The derivation of the formula (4), which we have just given according to the rules of the theory of combinations, is indeed perfectly correct and clear, but for its exact foundation it requires some reflection which cannot easily be compressed into a few words. We shall therefore subsequently prove its correctness once more by the method of complete induction.

The formulas (3), namely, give the following recurrence formulas:
\begin{equation}
\begin{aligned}
B_1^{(n)}&=B_1^{(n-1)}+1,\\
B_2^{(n)}&=B_2^{(n-1)}+B_1^{(n-1)},\\
&\cdots\\
B_\nu^{(n)}&=B_\nu^{(n-1)}+B_{\nu-1}^{(n-1)},\\
&\cdots\\
B_n^{(n)}&=B_{n-1}^{(n-1)}.
\end{aligned}
\tag{10}
\end{equation}
Now the formula (6) can be confirmed very easily by counting for the first values of $n$. If one assumes it correct for $n-1$, then (10) gives
\[
 B_\nu^{(n)}={\Pi(n-1)\over \Pi(\nu)\Pi(n-\nu-1)}+{\Pi(n-1)\over \Pi(\nu-1)\Pi(n-\nu)},
\]
whence, by (8),
\[
 B_\nu^{(n)}={\Pi(n)\over \Pi(\nu)\Pi(n-\nu)},
\]
and by this the general validity of formula (9) is proved.


% ===== scan leaf 67, printed 42 =====
% !PRINTED 42
\weberpage{42}
\pdfbookmark[2]{§ 8. The binomial theorem}{bm11}

From the meaning of $B_\nu^{(n)}$ it follows, and it is also proved from the formulas (10) by complete induction, that the $B_\nu^{(n)}$ are positive integers.
% SOURCE-ORIGINAL-SIC: The English states the mathematical sense of the source's historical expletive construction; the German leaf preserves it verbatim.

\begin{center}
{\large \S.\ 8.}

\vspace{0.5em}
{\large\gesperrt{The Binomial Theorem.}}
\end{center}

If in formula (1) of the preceding paragraph we set the previously arbitrary quantities $\alpha_1,\alpha_2\ldots\alpha_n$ equal to one another, we obtain the \gesperrt{binomial theorem}, which consists in nothing other than arranging the $n$th power of a binomial $x+y$ according to powers of $x$. Namely, if we put
\[
 \alpha_1=\alpha_2=\cdots=\alpha_n=-y
\]
then formula (2) gives
\[
 a_\nu=(-1)^\nu\sum \alpha_1\alpha_2\cdots\alpha_\nu=y^\nu B_\nu^{(n)},
\]
where, according to the definition of the preceding paragraph, $B_\nu^{(n)}$ means the number of terms of the sum, which all become equal to one another and equal to $(-1)^\nu y^\nu$.

Accordingly one obtains
\begin{equation}
 (x+y)^n=x^n+B_1^{(n)}x^{n-1}y+B_2^{(n)}x^{n-2}y^2+\cdots+B_n^{(n)}y^n,
\tag{1}
\end{equation}
or, written out,
\[
\begin{aligned}
(x+y)^n
&=x^n+nx^{n-1}y+{n(n-1)\over1\cdot2}x^{n-2}y^2+\cdots+y^n\\
&=\Pi(n)\sum_{\nu=0}^{n} {x^{n-\nu}y^\nu\over \Pi(n-\nu)\Pi(\nu)}
 =\Pi(n)\sum {x^\alpha y^\beta\over \Pi(\alpha)\Pi(\beta)},
\end{aligned}
\]
the last sum being extended over all non-negative integer values $\alpha,\beta$ which satisfy the condition $\alpha+\beta=n$.

Accordingly the coefficients $B_\nu^{(n)}$ are called the \gesperrt{binomial coefficients}.

For reference we put here a small table of the first values of the binomial coefficients:
\[
\begin{array}{crrrrrrrr}
 n=1&1&1\\
 n=2&1&2&1\\
 n=3&1&3&3&1\\
 n=4&1&4&6&4&1\\
 n=5&1&5&10&10&5&1\\
 n=6&1&6&15&20&15&6&1\\
 n=7&1&7&21&35&35&21&7&1
\end{array}
\]


% ===== scan leaf 68, printed 43 =====
% !PRINTED 43
\weberpage{43}

Among the various properties of the binomial coefficients we shall derive two, of which we shall afterwards make an interesting application.

From (1), replacing $x,y$ by $1,x$ and taking $n=0,1,2,3\ldots$, the formulas
\begin{equation}
\begin{aligned}
1&=B_0^{(0)},\\
1+x&=B_0^{(1)}+B_1^{(1)}x,\\
(1+x)^2&=B_0^{(2)}+B_1^{(2)}x+B_2^{(2)}x^2,\\
&\cdots\\
(1+x)^n&=B_0^{(n)}+B_1^{(n)}x+B_2^{(n)}x^2+\cdots+B_n^{(n)}x^n.
\end{aligned}
\tag{2}
\end{equation}
are obtained.

We now use the well-known summation formula for the geometric series:
\[
1+(1+x)+(1+x)^2+\cdots+(1+x)^n={(1+x)^{n+1}-1\over x},
\]
where, if $(1+x)^{n+1}$ is again expanded by the binomial formula, replacing $n$ by $n+1$ in the last formula (2) and setting $B_0^{(n+1)}=1$,
\begin{equation}
{(1+x)^{n+1}-1\over x}=B_1^{(n+1)}+B_2^{(n+1)}x+\cdots+B_{n+1}^{(n+1)}x^n.
\tag{3}
\end{equation}
Comparing this with the sum of the right hand sides of (2), and, according to the rules of \S~1, setting equal the coefficients of corresponding powers of $x$, one obtains
\begin{equation}
\begin{aligned}
 B_0^{(0)}+B_0^{(1)}+B_0^{(2)}+\cdots+B_0^{(n)}&=B_1^{(n+1)},\\
 B_1^{(1)}+B_1^{(2)}+\cdots+B_1^{(n)}&=B_2^{(n+1)},\\
&\cdots\\
 B_n^{(n)}&=B_{n+1}^{(n+1)},
\end{aligned}
\tag{4}
\end{equation}
or in general
\begin{equation}
 B_\nu^{(\nu)}+B_\nu^{(\nu+1)}+\cdots+B_\nu^{(n)}=B_{\nu+1}^{(n+1)}.
\tag{5}
\end{equation}
If however the equations (2) are multiplied in order by
\[
 B_0^{(n)},
 \quad -B_1^{(n)},
 \quad +B_2^{(n)},
 \ldots,
 \quad \pm B_n^{(n)}
\]
where the upper sign is to be taken for even and the lower for odd $n$, the sum of the left hand sides gives
\[
B_0^{(n)}-B_1^{(n)}(1+x)+B_2^{(n)}(1+x)^2-\cdots\pm B_n^{(n)}(1+x)^n=[1-(1+x)]^n=(-x)^n
\]
and equating the coefficients of equal powers on the right and left hand sides yields the system of formulas:
%CONT>>


% ===== scan leaf 69, printed 44 =====
% !PRINTED 44
\weberpage{44}
\pdfbookmark[2]{§ 9. Interpolation}{bm12}
%<<CONT

\begin{equation}
\begin{aligned}
0&=B_0^{(0)}B_0^{(n)}-B_0^{(1)}B_1^{(n)}+B_0^{(2)}B_2^{(n)}-\cdots\pm B_0^{(n)}B_n^{(n)},\\
0&=-B_1^{(1)}B_1^{(n)}+B_1^{(2)}B_2^{(n)}-\cdots\pm B_1^{(n)}B_n^{(n)},\\
&\cdots\\
\pm 1&=\pm B_n^{(n)}B_n^{(n)},
\end{aligned}
\tag{6}
\end{equation}
and this system of formulas can also be written in the reversed order, using the relation $B_\nu^{(n)}=B_{n-\nu}^{(n)}$, as
\begin{equation}
\begin{aligned}
1&=B_0^{(n)}B_0^{(n)},\\
0&=B_1^{(n)}B_0^{(n)}-B_0^{(n-1)}B_1^{(n)},\\
0&=B_2^{(n)}B_0^{(n)}-B_1^{(n-1)}B_1^{(n)}+B_0^{(n-2)}B_2^{(n)},\\
&\cdots\\
0&=B_n^{(n)}B_0^{(n)}-B_{n-1}^{(n-1)}B_1^{(n)}+B_{n-2}^{(n-2)}B_2^{(n)}-\cdots\pm B_0^{(0)}B_n^{(n)},
\end{aligned}
\tag{7}
\end{equation}
or, in a summarizing notation,
\begin{equation}
 \sum_{\nu=0}^{\mu}(-1)^\nu B_{\mu-\nu}^{(n-\nu)}B_\nu^{(n)}
 =\begin{cases}
 1,&\mu=0,\\
 0,&\mu=1,2,\ldots,n.
 \end{cases}
\tag{8}
\end{equation}


\begin{center}
{\large \S.\ 9.}

\vspace{0.5em}
{\large\gesperrt{Interpolation.}}
\end{center}

We shall make an application of the last obtained formulae to a problem from the theory of interpolation.

\gesperrt{A whole rational function of degree $n$ is to be determined which, for $n+1$ given values of the variable $x$, has prescribed values.}

Thus the problem is the determination of the $n+1$ coefficients $a_0,a_1,\ldots,a_n$ in
\begin{equation}
 f(x)=a_0x^n+a_1x^{n-1}+\cdots+a_{n-1}x+a_n
\tag{1}
\end{equation}
from the $n+1$ equations
\begin{equation}
 f(\alpha_0)=A_0,\quad f(\alpha_1)=A_1,\quad \ldots,\quad f(\alpha_n)=A_n,
\tag{2}
\end{equation}
when $\alpha_0,\alpha_1,\ldots,\alpha_n$ are the given values of $x$, and $A_0,A_1,\ldots,A_n$ are the corresponding values of $f(x)$.

Equations (2) form a system of equations of the first degree for the unknowns $a_0,a_1,\ldots,a_n$, which in general can be solved by determinants, as we shall see in the next section. We shall later return to this problem,
%CONT>>


% ===== scan leaf 70, printed 45 =====
% !PRINTED 45
\weberpage{45}
%<<CONT

but treat it now under the special supposition that the given values of $x$ are the first $n+1$ whole numbers $0,1,2,\ldots,n$.

The binomial coefficients $B_\nu^{(n)}$, as they were defined by \S~7, (6), keep their meaning even when $n$ is not a whole number, as was previously presupposed, but an arbitrary variable quantity. Then
\begin{equation}
 B_\nu^{(x)}={x(x-1)\cdots(x-\nu+1)\over 1\cdot2\cdot3\cdots\nu}
\tag{3}
\end{equation}
is a whole rational function of degree $\nu$ in $x$.
\footnote{The meaning of these generalized binomial coefficients for the expansion of powers of a binomial does not belong here, but to analysis.}
% SOURCE-ORIGINAL-SIC: The English regularizes the source footnote's misspelling "Binominalcoefficienten"; the German leaf preserves it verbatim.
The coefficient of $x^\nu$ is different from zero, and therefore $x^\nu$ can also be expressed by $B_\nu^{(x)}$ and by lower powers of $x$. Thus every whole function $f(x)$ of degree $n$ in $x$ can also be linearly expressed by $B_0^{(x)},B_1^{(x)},\ldots,B_n^{(x)}$ in the form
\begin{equation}
 f(x)=M_0B_0^{(x)}+M_1B_1^{(x)}+\cdots+M_nB_n^{(x)},
\tag{4}
\end{equation}
where $M_0,M_1,\ldots,M_n$ are constants; and the function $f(x)$ is determined when these constants are determined.

Now, by our supposition, $f(0),f(1),f(2),\ldots,f(n)$ are given. Since, by (3), $B_\nu^{(x)}$ always vanishes when $x$ has one of the values $0,1,2,\ldots,\nu-1$, the following linear equations for the unknowns $M$ are obtained from (4):
\begin{equation}
\begin{aligned}
 f(0)&=M_0B_0^{(0)},\\
 f(1)&=M_0B_0^{(1)}+M_1B_1^{(1)},\\
 f(2)&=M_0B_0^{(2)}+M_1B_1^{(2)}+M_2B_2^{(2)},\\
 &\quad\cdots\\
 f(n)&=M_0B_0^{(n)}+M_1B_1^{(n)}+M_2B_2^{(n)}+\cdots+M_nB_n^{(n)}.
\end{aligned}
\tag{5}
\end{equation}
These equations are now to be solved with respect to $M_0,M_1,\ldots,M_n$, which is done very easily with the help of the final equations (8) of the preceding paragraph. The first equation (5) gives directly
\begin{equation}
 M_0=f(0).
\tag{6}
\end{equation}


% ===== scan leaf 71, printed 46 =====
% !PRINTED 46
\weberpage{46}
\pdfbookmark[2]{§ 10. Solution of the interpolation problem by means of the differences}{bm13}

If the first equation (5) is multiplied by $B_0^{(1)}$, the second by $-B_1^{(1)}$, and they are added, one obtains, by the formula system mentioned (applied to $n=1$),
\[
 -M_1=B_0^{(1)}f(0)-B_1^{(1)}f(1),
\]
and so in general, if the first $\nu$ equations (5) are multiplied in order by
\[
 B_0^{(\nu)},\quad -B_1^{(\nu)},\quad +B_2^{(\nu)},\quad \ldots,\quad \pm B_\nu^{(\nu)}
\]
and then added,
\begin{equation}
 \pm M_\nu=B_0^{(\nu)}f(0)-B_1^{(\nu)}f(1)+B_2^{(\nu)}f(2)-\cdots\pm B_\nu^{(\nu)}f(\nu),
\tag{7}
\end{equation}
whereby, according to (4), the function $F(x)$ is determined and the problem is solved. It is clear that, so long as we make no special supposition about the values $f(0),f(1),\ldots,f(n)$, every arbitrary whole rational function of $x$ can be represented in this form.
% SOURCE-ORIGINAL-SIC (mathematical notation): The scan prints uppercase F(x) here although equations (1), (4), and the surrounding argument define lowercase f(x); both lanes remain source-faithful.

\begin{center}
{\large \S.\ 10.}

\vspace{0.5em}
{\large\gesperrt{Solution of the Interpolation Problem by Differences.}}
\end{center}

The definition of the whole rational function
\begin{equation}
 B_\nu^{(x)}={x(x-1)\cdots(x-\nu+1)\over 1\cdot2\cdot3\cdots\nu}
\tag{1}
\end{equation}
gives the relation, already proved earlier for the special case of a positive integer $x$,
\begin{equation}
 B_\nu^{(x+1)}=B_\nu^{(x)}+B_{\nu-1}^{(x)},\qquad B_0^{(x+1)}=B_0^{(x)}=1.
\tag{2}
\end{equation}
From this we obtain for the coefficients $M_0,M_1,\ldots,M_n$ a mode of determination which is much more convenient for practical calculation than the application of the formulae of the preceding paragraph.

Namely, by the formulae (4) and (6) of \S~9,
\begin{equation}
 f(x)=f(0)+M_1B_1^{(x)}+M_2B_2^{(x)}+\cdots+M_nB_n^{(x)},
\tag{3}
\end{equation}
and if in this expression we replace $x$ by $x+1$ and form the difference
\begin{equation}
 \Delta_x=f(x+1)-f(x)
\tag{4}
\end{equation}
then, with regard to (2),
\begin{equation}
 \Delta_x=M_1+M_2B_1^{(x)}+M_3B_2^{(x)}+\cdots+M_nB_{n-1}^{(x)},
\tag{5}
\end{equation}
from which follows (since (5) is an equation of the same kind as (3), except that $n-1$ has taken the place of $n$):
%CONT>>


% ===== scan leaf 72, printed 47 =====
% !PRINTED 47
\weberpage{47}
\pdfbookmark[2]{§ 11. Arithmetic series of higher order}{bm14}
%<<CONT

\[
M_1=\Delta_0=f(1)-f(0).
\]
Put
\begin{equation}
\begin{aligned}
 \Delta_{x+1}-\Delta_x&=\Delta_x',\\
 \Delta_{x+1}'-\Delta_x'&=\Delta_x'',\\
 &\cdots\\
 \Delta_{x+1}^{(n-2)}-\Delta_x^{(n-2)}&=\Delta_x^{(n-1)}.
\end{aligned}
\tag{6}
\end{equation}
then, accordingly,
\[
 M_0=f(0),\quad M_1=\Delta_0,
 \quad M_2=\Delta_0',\quad\ldots,\quad M_n=\Delta_0^{(n-1)},
\]
and formula (3) gives
\begin{equation}
 f(x)=f(0)+\Delta_0B_1^{(x)}+\Delta_0'B_2^{(x)}+\Delta_0''B_3^{(x)}+\cdots+\Delta_0^{(n-1)}B_n^{(x)}.
\tag{7}
\end{equation}
Likewise the functions $\Delta_x,\Delta_x',\Delta_x'',\ldots$ can be expressed:
\begin{equation}
\begin{aligned}
 \Delta_x&=\Delta_0+\Delta_0'B_1^{(x)}+\Delta_0''B_2^{(x)}+\cdots+\Delta_0^{(n-1)}B_{n-1}^{(x)},\\
 \Delta_x'&=\Delta_0'+\Delta_0''B_1^{(x)}+\cdots+\Delta_0^{(n-1)}B_{n-2}^{(x)}.
\end{aligned}
\tag{8}
\end{equation}
and so on.

The $\Delta_x,\Delta_x',\Delta_x'',\ldots,\Delta_x^{(n-1)}$ are whole rational functions of degrees $n-1,n-2,\ldots,0$; the last of them is therefore constant. In order to represent them all, one needs only the values
\[
 f(0),\quad \Delta_0,\quad \Delta_0',\quad \Delta_0'',\quad\ldots,\quad \Delta_0^{(n-1)},
\]
which one computes most easily by arranging a table. For the case $n=3$, for example, the table has the form:
\[
\begin{array}{c|c|c|c}
f(0) &&&\\
f(1)&\Delta_0&&\\
f(2)&\Delta_1&\Delta_0'&\\
f(3)&\Delta_2&\Delta_1'&\Delta_0''
\end{array}
\]
When $f(0),f(1),\ldots$ are given, the table is filled by simple subtractions.

\begin{center}
{\large \S.\ 11.}

\vspace{0.5em}
{\large\gesperrt{Arithmetic Series of Higher Order.}}
\end{center}

The foregoing is the foundation for the theory of arithmetic series of higher order.

We denote by
\begin{equation}
 u_0,u_1,u_2,u_3,\ldots
\tag{1}
\end{equation}
any infinite series of quantities and by
%CONT>>


% ===== scan leaf 73, printed 48 =====
% !PRINTED 48
\weberpage{48}
%<<CONT

\begin{equation}
\Delta_0=u_1-u_0,
 \Delta_1=u_2-u_1,
 \Delta_2=u_3-u_2,\ldots
\tag{2}
\end{equation}
the series of its differences; and
\begin{equation}
 \Delta'_0=\Delta_1-\Delta_0,
 \Delta'_1=\Delta_2-\Delta_1,\ldots
\tag{3}
\end{equation}
the series of its second differences, and so on.

It is clear that the series (1) is completely determined when its first term and the series of its first differences are given; for
\[
 u_1=u_0+\Delta_0,
 \quad u_2=u_0+\Delta_0+\Delta_1,
 \quad\ldots,
\]
\[
 u_m=u_0+\Delta_0+\Delta_1+\cdots+\Delta_{m-1}.
\]
Likewise the series (1) is completely determined when the first two terms and the series of its second differences are given, and so forth. The series (1) is called an \gesperrt{arithmetic series of $n$th order} when the series of its $n$th differences is constant, and hence the series of the $(n+1)$st differences consists entirely of zeros.

One obtains an arithmetic series of $n$th order by substituting the numbers $0,1,2,3,\ldots$ for $x$ in a whole function $f(x)$ of degree $n$. For if one puts
\[
 \Delta_x=f(x+1)-f(x),
\]
\[
 \Delta'_x=\Delta_{(x+1)}-\Delta_x,
\]
\[
 \cdots\cdots\cdots\cdots\cdots
\]
\[
 \Delta_x^{(n-1)}=\Delta_{(x+1)}^{(n-2)}-\Delta_x^{(n-2)},
\]
then $\Delta_x$ is of degree $(n-1)$, $\Delta'_x$ of degree $(n-2)$, and so on with respect to $x$, and therefore $\Delta_x^{(n-1)}$ is constant.

If now
\[
 u_0,u_1,u_2,\ldots
\]
is an arithmetic series of $n$th order, then the whole series is completely determined when the $n$ values $u_0,u_1,u_2,\ldots,u_{n-1}$ and, besides these, the constant $n$th difference are given. But this latter is determined if also the $(n+1)$st term $u_n$ is given. We can therefore state the theorem:

\gesperrt{An arithmetic series of $n$th order is completely determined when its first $n+1$ terms are given.}

Since a function $f(x)$ of degree $n$ is likewise completely determined by the arbitrarily given values
\[
 f(0),f(1),f(2),\ldots,f(n)
\]
%CONT>>


% ===== scan leaf 74, printed 49 =====
% !PRINTED 49
\weberpage{49}
%<<CONT

\gesperrt{it follows that from the whole rational functions \(f(x)\) of degree \(n\) all arithmetic series of \(n\)th order are produced when in them the series of natural numbers is substituted for \(x\).}

The expression for the general term is then given by formula (7) of the preceding paragraph.

The sums of the first $m+1$ terms of an arithmetic series of $n$th order,
\[
 s_m=u_0+u_1+\cdots+u_m,
\]
form an arithmetic series of $(n+1)$st order, since their first differences
\[
 s_{m+1}-s_m=u_{m+1}
\]
form an arithmetic series of $n$th order.

Thus, with the help of formula (7) of \S~10, the sum $s_m$ can be determined in general when $s_0,s_1,\ldots,s_{n+1}$ are assumed known.

To find the generating function $F(x)$ of $s_m$, when $f(x)$ is the generating function of $u_m$, put
\[
 F(0)=f(0),
\]
\[
 F(1)=f(0)+f(1),
\]
\[
 F(2)=f(0)+f(1)+f(2),
\]
\[
 \cdots\cdots\cdots\cdots
\]
and in formula (7), \S~10,
\[
 F(x)=F(0)+D_0B_1^{(x)}+D'_0B_2^{(x)}+\cdots
\]
one has to put
\[
\begin{aligned}
D_0&=F(1)-F(0)=f(1),& D'_0&=f(2)-f(1)=\Delta_1,\\
D_1&=F(2)-F(1)=f(2),& D'_1&=f(3)-f(2)=\Delta_2,\\
 D_2&=F(3)-F(2)=f(3),& D''_0&=\Delta_2-\Delta_1=\Delta'_1.
\end{aligned}
\]
Thus one obtains
\begin{equation}
 F(x)=f(0)+f(1)B_1^{(x)}+\Delta_1B_2^{(x)}+\Delta'_1B_3^{(x)}+\cdots.
\tag{4}
\end{equation}
For example, take $f(x)=x^2$; then $F(m)$ gives the sum of the first $m$ square numbers. Here
\[
 \Delta_x=2x+1,\qquad \Delta'_x=2,
\]
and therefore
%CONT>>


% ===== scan leaf 75, printed 50 =====
% !PRINTED 50
\weberpage{50}
\pdfbookmark[2]{§ 12. The multinomial theorem}{bm15}
%<<CONT

\[
 F(x)=x+3{x(x-1)\over 1\cdot2}+2{x(x-1)(x-2)
        \over 1\cdot2\cdot3}
      ={x(x+1)(2x+1)\over 6}.
\]
For $f(x)=x^3$ the same calculation gives
\[
 F(x)=\left({x(x+1)\over2}\right)^2.
\]
Thus the sum of the first $m$ cubes is equal to the square of the $m$th triangular number.

\begin{center}
{\large \S.\ 12.}

\vspace{0.5em}
{\large\gesperrt{The Polynomial Theorem.}}
\end{center}

In \S~8 the following form of the binomial theorem was derived:
\begin{equation}
 (x+y)^n=\Pi(n)\sum^{\alpha,\beta}{x^\alpha y^\beta\over \Pi(\alpha)\Pi(\beta)},
\tag{1}
\end{equation}
where the sum extends over all combinations of two numbers $\alpha,\beta$, neither of which is negative, and which satisfy the condition
\begin{equation}
 \alpha+\beta=n.
\tag{2}
\end{equation}

This form permits, first by induction, a generalization to the $n$th power of a polynomial:
\begin{equation}
 (x+y+z+\cdots)^n=\Pi(n)\sum^{\alpha,\beta,\gamma,\ldots}
 {x^\alpha y^\beta z^\gamma\cdots
\over \Pi(\alpha)\Pi(\beta)\Pi(\gamma)\cdots},
\tag{3}
\end{equation}
with the determination that $\alpha,\beta,\gamma,\ldots$ run through all positive or zero integer values which satisfy
\begin{equation}
 \alpha+\beta+\gamma+\cdots=n.
\tag{4}
\end{equation}
But to prove the correctness of this formula in general, we assume that it is proved when the polynomial contains one term fewer, as is indeed the case when the polynomial contains only two terms.

We then put
\begin{equation}
 u=y+z+\cdots
\tag{5}
\end{equation}
and apply formula (1) to $(x+u)$, from which there follows
%CONT>>


% ===== scan leaf 76, printed 51 =====
% !PRINTED 51
\weberpage{51}
\pdfbookmark[2]{§ 13. Derived functions}{bm16}
%<<CONT

\begin{equation}
 (x+y+z+\cdots)^n=\Pi(n)\sum^{\alpha,\nu}{x^\alpha u^\nu
\over \Pi(\alpha)\Pi(\nu)},
\tag{6}
\end{equation}
with the restriction
\begin{equation}
 \alpha+\nu=n.
\tag{7}
\end{equation}
By the assumption, however, it has already been proved that
\begin{equation}
 u^\nu=\Pi(\nu)\sum^{\beta,\gamma,\ldots}
 {y^\beta z^\gamma\cdots
\over \Pi(\beta)\Pi(\gamma)\cdots},
\tag{8}
\end{equation}
where
\begin{equation}
 \beta+\gamma+\cdots=\nu.
\tag{9}
\end{equation}
When this is inserted in (6), formula (3) is obtained immediately, and (7) passes over into (4).

The coefficients
\begin{equation}
 P_{\alpha,\beta,\gamma,\ldots}^{(n)}={\Pi(n)\over
 \Pi(\alpha)\Pi(\beta)\Pi(\gamma)\cdots},
\tag{10}
\end{equation}
which, according to their meaning, are whole numbers, are called the \gesperrt{polynomial coefficients}.

For example, for the third power of the trinomial one obtains
\begin{equation}
\begin{aligned}
(x+y+z)^3={}&x^3+y^3+z^3+3x^2y+3xy^2+3x^2z+3xz^2\\
&+3y^2z+3yz^2+6xyz.
\end{aligned}
\tag{11}
\end{equation}


\begin{center}
{\large \S.\ 13.}

\vspace{0.5em}
{\large\gesperrt{Derived Functions.}}
\end{center}

Let
\begin{equation}
 f(x)=a_0x^n+a_1x^{n-1}+a_2x^{n-2}+\cdots+a_n
\tag{1}
\end{equation}
be a whole rational function of order $n$.

If in it we replace $x$ by a binomial $x+y$, then we can apply the binomial theorem to each separate term, and can arrange the result according to descending or according to ascending powers of $x$ or of $y$. We shall carry out the arrangement according to ascending powers of $y$. The highest power of $y$ which occurs is the $n$th, and the coefficient of the zeroth power of $y$ is the function $f(x)$ itself, as one sees by putting $y=0$. We therefore put, denoting the other coefficients by
\[
 f'(x),\qquad {f''(x)\over \Pi(2)},\qquad {f'''(x)\over \Pi(3)},\ldots:
\]
%CONT>>


% ===== scan leaf 77, printed 52 =====
% !PRINTED 52
\weberpage{52}
%<<CONT

\begin{equation}
 f(x+y)=f(x)+yf'(x)+{y^2\over1\cdot2}f''(x)+\cdots
       =\sum_{\nu=0}^{n}{y^\nu\over\Pi(\nu)}\,f^{(\nu)}(x).
\tag{2}
\end{equation}

The functions $f'(x),f''(x),f'''(x),\ldots$ are called the \gesperrt{first, second, third, \ldots{} derivatives or derived functions of $f(x)$}. They are whole functions of $x$, and $f^{(\nu)}(x)$ cannot exceed the degree $n-\nu$, since the sum of the exponents of $x$ and $y$ in no term exceeds the degree $n$.

\gesperrt{The first derivative, which is thus the coefficient of the first power of $y$ in the expansion of $f(x+y)$ according to ascending powers of $y$}, is obtained by applying the binomial theorem to (1):
\begin{equation}
 f'(x)=na_0x^{n-1}+(n-1)a_1x^{n-2}+(n-2)a_2x^{n-3}+\cdots .
\tag{3}
\end{equation}

The principal theorem on derived functions follows from (2) if we transform $x$ into $x+z$ or $y$ into $y+z$:
\begin{equation}
 f(x+y+z)=\sum_{\nu=0}^{n}{y^\nu\over\Pi(\nu)}f^{(\nu)}(x+z)
 =\sum_{\nu=0}^{n}{(y+z)^\nu\over\Pi(\nu)}f^{(\nu)}(x).
\tag{4}
\end{equation}
If by $f^{(\nu,\mu)}(x)$ we denote the $\mu$th derivative of $f^{(\nu)}(x)$, then by (2)
\begin{equation}
 f^{(\nu)}(x+z)=\sum_{\mu=0}^{n-\nu}{z^\mu\over\Pi(\mu)}f^{(\nu,\mu)}(x),
\tag{5}
\end{equation}
and by the binomial theorem
\begin{equation}
 { (y+z)^\nu\over\Pi(\nu)}
 =\sum^{\beta,\gamma}{y^\beta z^\gamma\over\Pi(\beta)\Pi(\gamma)},
 \qquad \beta+\gamma=\nu.
\tag{6}
\end{equation}
Substituting this in (4), one obtains
\begin{equation}
 \sum_{\nu=0}^{n}\sum_{\mu=0}^{n-\nu}
 {y^\nu z^\mu\over\Pi(\nu)\Pi(\mu)}f^{(\nu,\mu)}(x)
 =
 \sum^{\beta,\gamma}
 {y^\beta z^\gamma\over\Pi(\beta)\Pi(\gamma)}f^{(\beta+\gamma)}(x).
\tag{7}
\end{equation}
The last sum is to be extended over all non-negative numbers $\beta,\gamma$ whose sum does not exceed the degree $n$ of $f(x)$. But the same number-combinations are also run through by the exponents $\nu,\mu$ on the left-hand side, and comparison of the coefficients of equal powers and products gives (by \S~1)
\begin{equation}
 f^{(\nu,\mu)}(x)=f^{(\nu+\mu)}(x),
\tag{8}
\end{equation}
and therefore the theorem:
%CONT>>


% ===== scan leaf 78, printed 53 =====
% !PRINTED 53
\weberpage{53}
%<<CONT

\gesperrt{The $\mu$th derivative of the $\nu$th derivative is the $(\nu+\mu)$th derivative of the original function.}

Thus all the higher derivatives are obtained by forming from each preceding one, according to the rule (3), the first derivative:
\begin{align}
 f(x)&=a_0x^n+a_1x^{n-1}+a_2x^{n-2}+a_3x^{n-3}+\cdots,\notag\\
 f'(x)&=na_0x^{n-1}+(n-1)a_1x^{n-2}+(n-2)a_2x^{n-3}+\cdots, \tag{9}\\
 f''(x)&=n(n-1)a_0x^{n-2}+(n-1)(n-2)a_1x^{n-3}+\cdots. \notag
\end{align}

These derivatives assume a somewhat simpler form if one uses another notation, which is often in use and for certain purposes almost indispensable, and which we shall discuss in connection with this.

Because of the indeterminacy of the coefficients $a_0,a_1,\ldots,a_n$, it is clearly no restriction if we write a whole rational function of degree $n$ in the form
\begin{equation}
 f(x)=a_0x^n+B_1^{(n)}a_1x^{n-1}+B_2^{(n)}a_2x^{n-2}+\cdots+a_n,
\tag{10}
\end{equation}
or explicitly
\begin{equation}
 f(x)=a_0x^n+na_1x^{n-1}+{n(n-1)\over1\cdot2}a_2x^{n-2}+\cdots.
\tag{11}
\end{equation}
When a function $f(x)$ is represented in this way, we shall say that it is ``\gesperrt{written with the binomial coefficients}.''

A greater agreement is already shown by the formulae (9), which then read:
\begin{align}
 f(x)&=a_0x^n+na_1x^{n-1}+{n(n-1)\over1\cdot2}a_2x^{n-2}+\cdots,\notag\\
 {1\over n}f'(x)&=a_0x^{n-1}+(n-1)a_1x^{n-2}+{(n-1)(n-2)\over1\cdot2}a_2x^{n-3}+\cdots, \tag{12}\\
 {1\over n(n-1)}f''(x)&=a_0x^{n-2}+(n-2)a_1x^{n-3}+{(n-2)(n-3)\over1\cdot2}a_2x^{n-4}+\cdots,\notag
\end{align}
wherein the right-hand sides all likewise appear written with the binomial coefficients. We shall later become more acquainted with the usefulness of this notation, but already here we must emphasize that the choice of the one or the other
%CONT>>


% ===== scan leaf 79, printed 54 =====
% !PRINTED 54
\weberpage{54}
\pdfbookmark[2]{§ 14. Derivative of a product}{bm17}
%<<CONT

manner of representation is nevertheless not wholly indifferent, and that the first often deserves preference. Especially in those cases in which the coefficients are numbers and the number-theoretic behavior of these coefficients is at issue, one must not overlook that through the binomial coefficients a numerical element foreign to the matter is introduced. That Gauss in the theory of quadratic forms (in the Disq. ar.) uses the notation with binomial coefficients when he represents the quadratic forms by $ax^2+2bxy+cy^2$, and that this notation has come into general use, has led in number theory to an unnecessary and very regrettable complication.

\begin{center}
{\large \S.\ 14.}

\vspace{0.5em}
{\large\gesperrt{Derivative of a Product.}}
\end{center}

The derived functions which we have considered here are no others than the differential quotients known from differential calculus; however, here, where only whole rational functions are involved, we have obtained the concept without using infinitesimal calculus, from the coefficients of development of the powers of $y$ in the function $f(x+y)$. If we denote the $\nu$th derivative of $f(x)$ by $D_\nu f$, then by (2), \S~13:
\begin{equation}
 f(x+y)=f(x)+yD_1f+{y^2\over1\cdot2}D_2f+{y^3\over1\cdot2\cdot3}D_3f+\cdots,
\tag{1}
\end{equation}
and from this follow at once the two fundamental rules expressed in the formulae
\begin{equation}
 D_\nu(Cf)=C D_\nu f;
\tag{2}
\end{equation}
\begin{equation}
 D_\nu(f+\varphi)=D_\nu f+D_\nu\varphi,
\tag{3}
\end{equation}
where $C$ is a constant, and $\varphi$ a second whole rational function of $x$.

A generalization of formula (2) gives the representation of the derivatives of the product $f\varphi$. Namely, if after (1) one writes, for abbreviation,
\begin{align}
 f(x+y)&=u_0+yu_1+y^2u_2+\cdots+y^n u_n,\notag\\
 \varphi(x+y)&=v_0+yv_1+y^2v_2+\cdots+y^m v_m,
\tag{4}
\end{align}
so that
\begin{equation}
 u_\nu={D_\nu f\over\Pi(\nu)},\qquad
 v_\nu={D_\nu\varphi\over\Pi(\nu)},
\tag{5}
\end{equation}
%CONT>>


% ===== scan leaf 80, printed 55 =====
% !PRINTED 55
\weberpage{55}
%<<CONT

then carrying out the multiplication of the two formulae (4), and seeking in the product the coefficient of $y^\nu$, gives
\begin{equation}
 {D_\nu(f\varphi)\over\Pi(\nu)}
 =u_\nu v_0+u_{\nu-1}v_1+u_{\nu-2}v_2+\cdots+u_0v_\nu,
\tag{6}
\end{equation}
or
\begin{equation}
 D_\nu(f\varphi)=\varphi D_\nu f+B_1^{(\nu)}D_{\nu-1}fD_1\varphi
 +B_2^{(\nu)}D_{\nu-2}fD_2\varphi+\cdots,
\tag{7}
\end{equation}
where $B_1^{(\nu)},B_2^{(\nu)},\ldots$ are the binomial coefficients. In a similar way one can, by using the polynomial coefficients, form the derivatives of a product of more than two factors.

We shall form the \gesperrt{first derivative}, which we now denote by $D$ instead of by $D_1$, for a product of $n$ factors. For two factors we obtain from (7)
\[
 D(f\varphi)=\varphi Df+fD\varphi=\varphi f'(x)+f\varphi'(x),
\]
and generally, if we denote the factors by $u_1,u_2,\ldots,u_n$ and their derivatives by $u'_1,u'_2,\ldots,u'_n$,
\begin{equation}
\begin{aligned}
D(u_1u_2\cdots u_n)
&=u'_1u_2\cdots u_n+u_1u'_2\cdots u_n+\cdots+u_1u_2\cdots u'_n,
\end{aligned}
\tag{8}
\end{equation}
or more briefly,
\[
 {D(u_1u_2\cdots u_n)\over u_1u_2\cdots u_n}
 =\sum {D u_\nu\over u_\nu}.
\]

Thus if we consider a product of linear factors
\begin{equation}
 f(x)=(x-\alpha_1)(x-\alpha_2)\cdots(x-\alpha_n),
\tag{9}
\end{equation}
we obtain, since the first derivatives of $x-\alpha_1$, $x-\alpha_2,\ldots,x-\alpha_n$ are all equal to $1$,
\begin{equation}
\begin{aligned}
f'(x)&=(x-\alpha_2)(x-\alpha_3)\cdots(x-\alpha_n)\\
&\quad +(x-\alpha_1)(x-\alpha_3)\cdots(x-\alpha_n)+\cdots\\
&\quad +(x-\alpha_1)(x-\alpha_2)\cdots(x-\alpha_{n-1}),
\end{aligned}
\tag{10}
\end{equation}
for which one can also write
\begin{equation}
 f'(x)={f(x)\over x-\alpha_1}+{f(x)\over x-\alpha_2}+\cdots+{f(x)\over x-\alpha_n}.
\tag{11}
\end{equation}
A very important result follows from this if $x$ is put equal to one of the values $\alpha_1,\alpha_2,\ldots,\alpha_n$, namely
%CONT>>


% ===== scan leaf 81, printed 56 =====
% !PRINTED 56
\weberpage{56}
\pdfbookmark[2]{§ 15. Integral functions of several variables: forms}{bm18}
%<<CONT

\begin{equation}
\begin{aligned}
 f'(\alpha_1)&=(\alpha_1-\alpha_2)(\alpha_1-\alpha_3)\cdots(\alpha_1-\alpha_n),\\
 f'(\alpha_2)&=(\alpha_2-\alpha_1)(\alpha_2-\alpha_3)\cdots(\alpha_2-\alpha_n),\\
 &\quad\cdots\\
 f'(\alpha_n)&=(\alpha_n-\alpha_1)(\alpha_n-\alpha_2)\cdots(\alpha_n-\alpha_{n-1}).
\end{aligned}
\tag{12}
\end{equation}

\begin{center}
{\large \S.\ 15.}

\vspace{0.5em}
{\large\gesperrt{Whole Functions of Several Variables: Forms.}}
\end{center}

Up to now we have chiefly considered whole rational functions of one variable; but we cannot always restrict ourselves to this, and indeed we have already used functions of several variables above. By a whole rational function of degree $n$ in several variables $F(x,y,z,\ldots)$ we understand a sum of terms
\[
 \sum A_{\alpha,\beta,\gamma,\ldots}x^\alpha y^\beta z^\gamma\cdots,
\]
where $\alpha,\beta,\gamma,\ldots$ are non-negative integral exponents whose sum $\alpha+\beta+\gamma+\cdots$ does not exceed the value $n$, and in at least one term actually attains it. Thus the degree is determined by the greatest value assumed by the sum $\alpha+\beta+\gamma+\cdots$.

If the sum of the exponents $\alpha+\beta+\gamma+\cdots$ has the same value in all terms, then the function is called \gesperrt{homogeneous}.

\gesperrt{A fundamental property of homogeneous functions of degree $n$ is this: if all variables are multiplied by the same factor, the result is the same as if the function were multiplied by the $n$th power}; in signs, if $t$ denotes an arbitrary variable,
\begin{equation}
 F(tx,ty,tz,\ldots)=t^nF(x,y,z,\ldots);
\tag{1}
\end{equation}
for if in the product $x^\alpha y^\beta z^\gamma\cdots$ the variables are replaced by $tx,ty,tz,\ldots$, it receives the factor
\[
 t^{\alpha+\beta+\gamma+\cdots}.
\]
If now $\alpha+\beta+\gamma+\cdots$ has in all terms the same value $n$, the factor $t^n$ can be taken out before the sum. But if the sum $\alpha+\beta+\gamma+\cdots$ has different values in the separate terms, no such common factor can be taken out, at least
%CONT>>


% ===== scan leaf 82, printed 57 =====
% !PRINTED 57
\weberpage{57}
%<<CONT

not without leaving different powers of $t$ in the separate terms.

By increasing the number of variables one can transform every non-homogeneous function into a homogeneous one of the same degree. If $m-1$ is the number of variables in a non-homogeneous function of degree $n$, then we put
\[
 x={x_1\over x_m},\qquad y={x_2\over x_m},\qquad z={x_3\over x_m},\ldots,
\]
and obtain in
\[
 x_m^n F\!\left({x_1\over x_m},{x_2\over x_m},{x_3\over x_m},\ldots\right)
\]
a whole homogeneous function of degree $n$ in the variables $x_1,x_2,\ldots,x_m$, which we denote by
\[
 \Phi(x_1,x_2,\ldots,x_m).
\]

It is sometimes advisable to write homogeneous functions of several variables \gesperrt{with the polynomial coefficients}. We therefore put
\begin{equation}
 \Phi(x_1,x_2,\ldots,x_m)=
 \sum{\Pi(n)\over\Pi(\alpha_1)\Pi(\alpha_2)\cdots\Pi(\alpha_m)}
 A_{\alpha_1,\alpha_2,\ldots,\alpha_m}
 x_1^{\alpha_1}x_2^{\alpha_2}\cdots x_m^{\alpha_m},
\tag{2}
\end{equation}
where the sum is extended over all non-negative numbers satisfying the condition
\begin{equation}
 \alpha_1+\alpha_2+\cdots+\alpha_m=n.
\tag{3}
\end{equation}
This notation, without the restriction (3), is also applicable to non-homogeneous functions.

But the homogeneous function can also be represented as
\begin{equation}
 \Phi(x_1,x_2,\ldots,x_m)=
 \sum A_{\nu_1,\nu_2,\ldots,\nu_n}x_{\nu_1}x_{\nu_2}\cdots x_{\nu_n},
\tag{4}
\end{equation}
where each of the indices $\nu_1,\nu_2,\ldots,\nu_n$, independently of the others, runs through the range of values $1,2,\ldots,m$. Thus the sum (4) consists of $m^n$ terms, but not all of them are different from one another. The product $x_{\nu_1}x_{\nu_2}\cdots x_{\nu_n}$ remains unchanged if the indices $\nu_1,\nu_2,\ldots,\nu_n$ are permuted among themselves in any way. The number of permutations of $n$ elements is $\Pi(n)$. If among these elements $\alpha_1,\alpha_2,\ldots$ respectively are equal to one another, then the number of permutations is reduced to
\[
 {\Pi(n)\over \Pi(\alpha_1)\Pi(\alpha_2)\cdots},
\]
%CONT>>


% ===== scan leaf 83, printed 58 =====
% !PRINTED 58
\weberpage{58}
%<<CONT

whence it follows that in (4) any product $x_1^{\alpha_1}x_2^{\alpha_2}\cdots$ occurs exactly
\[
 {\Pi(n)\over \Pi(\alpha_1)\Pi(\alpha_2)\cdots}
\]
times. If one further fixes that $A_{\nu_1,\nu_2,\ldots,\nu_n}$ shall not change when the indices are permuted in any way, then the notations (2) and (4) prove to be identical if, by collecting equal factors,
\[
 x_{\nu_1}x_{\nu_2}\cdots x_{\nu_n}=x_1^{\alpha_1}x_2^{\alpha_2}\cdots x_m^{\alpha_m}
\]
and
\[
 A_{\nu_1,\nu_2,\ldots,\nu_n}=A_{\alpha_1,\alpha_2,\ldots,\alpha_m}
\]
are put.

If the number of terms that occur in the function $\Phi$ [according to (2)] is denoted by $(m,n)$, then, by first counting the terms which have the factor $x_1$ and then the remaining ones, which form a homogeneous function of order $n$ of the remaining $m-1$ variables, one finds the recursion formula
\begin{equation}
 (m,n)=(m,n-1)+(m-1,n),
\tag{5}
\end{equation}
with whose aid one proves by complete induction the expression
\begin{equation}
 (m,n)={m(m+1)\cdots(m+n-1)\over 1\cdot2\cdots n}
 ={ \Pi(m+n-1)\over \Pi(n)\Pi(m-1)}
\tag{6}
\end{equation}
to be correct.

Whole homogeneous functions are also called \gesperrt{forms}. According to the number of variables one distinguishes \gesperrt{unary forms} (simple powers), \gesperrt{binary, ternary, quaternary forms}. It is the binary forms that especially interest us here; their theory is essentially identical with the theory of whole rational functions of one variable. One passes from the binary forms back to these functions by regarding one of the homogeneous variables as constant, for example by giving it the value $1$.


% ===== scan leaf 84, printed 59 =====
% !PRINTED 59
\weberpage{59}
\pdfbookmark[2]{§ 16. The derivatives of functions of several variables}{bm19}

\begin{center}
{\large \S.\ 16.}

\vspace{0.5em}
{\large\gesperrt{The Derivatives of Functions of Several Variables.}}
\end{center}

In \S~13 we defined the derived functions of a whole rational function of one variable. The concept is immediately transferable to functions of several variables, by forming the derivatives with respect to each variable for itself, as though it were the only one. Thus, if one writes for instance as in \S~15
\begin{equation}
 F(x,y,z,\ldots)=\sum A_{\alpha,\beta,\gamma,\ldots}x^\alpha y^\beta z^\gamma\cdots,
\tag{1}
\end{equation}
one obtains the first derivative with respect to $x$:
\begin{equation}
 F'(x)=\sum \alpha A_{\alpha,\beta,\gamma,\ldots}x^{\alpha-1}y^\beta z^\gamma\cdots,
\tag{2}
\end{equation}
or with respect to $y$:
\begin{equation}
 F'(y)=\sum \beta A_{\alpha,\beta,\gamma,\ldots}x^\alpha y^{\beta-1}z^\gamma\cdots
\tag{3}
\end{equation}
and so on. From these functions one can again form, according to the same rules, the derivatives with respect to the different variables, and thus obtains the higher derivatives.

In order to display the results more clearly, let $\Phi(x_1,x_2,\ldots,x_m)$ be a whole rational function of order $n$ in the $m$ variables $x_1,x_2,\ldots,x_m$. We replace these variables by binomials
\[
 x_1+\xi_1,\quad x_2+\xi_2,\ldots,\quad x_m+\xi_m
\]
and expand, in each term of the function
\[
 \Phi(x_1+\xi_1,x_2+\xi_2,\ldots,x_m+\xi_m)=\Phi(x+\xi),
\]
by carrying out the multiplication
\[
 (x_1+\xi_1)^{\mu_1}(x_2+\xi_2)^{\mu_2}\cdots(x_m+\xi_m)^{\mu_m}
\]
according to powers of $\xi_1,\xi_2,\ldots,\xi_m$. If equal powers and products of the variables $\xi$ are collected each into one term, then for $\Phi(x+\xi)$, in the notation (2) of \S~15, one obtains a representation which in differential calculus is called the \gesperrt{Taylor expansion}:
\begin{equation}
 \Phi(x+\xi)=\sum
 {\xi_1^{\alpha_1}\xi_2^{\alpha_2}\cdots\xi_m^{\alpha_m}
 \over \Pi(\alpha_1)\Pi(\alpha_2)\cdots\Pi(\alpha_m)}
 D_{\alpha_1,\alpha_2,\ldots,\alpha_m}\Phi .
\tag{4}
\end{equation}
The coefficients, which we denote by
\[
 D_{\alpha_1,\alpha_2,\ldots,\alpha_m}\Phi,
\]
are functions of the variables $x$ and are called, when
%CONT>>


% ===== scan leaf 85, printed 60 =====
% !PRINTED 60
\weberpage{60}
%<<CONT

\[
 \alpha_1+\alpha_2+\cdots+\alpha_m=\nu,
\]
the \gesperrt{derivatives of order $\nu$ of the function $\Phi$}.

They are also written, according to the notation usual in differential calculus, as
\begin{equation}
 D_{\alpha_1,\alpha_2,\ldots,\alpha_m}\Phi
 =
 {\partial^\nu\Phi\over
 \partial x_1^{\alpha_1}\partial x_2^{\alpha_2}\cdots\partial x_m^{\alpha_m}}.
\tag{5}
\end{equation}

The law of formation of the derivatives can be summarized in the following propositions, where for brevity we omit the indices on the symbol $D$.

\gesperrt{I. If $C$ is a constant, then}
\[
 D(C\Phi)=CD\Phi.
\]
\gesperrt{II. If $\Phi$ and $\Psi$ are any two functions, then}
\[
 D(\Phi+\Psi)=D\Phi+D\Psi.
\]
Both follow immediately from (4).

Thus we can easily form the derived functions in general if we know them for the special case in which $\Phi$ is a product of powers, that is, if we know
\[
 D_{\alpha_1,\alpha_2,\ldots,\alpha_m}(x_1^{\mu_1}x_2^{\mu_2}\cdots x_m^{\mu_m}),
\]
where the $\mu$ are arbitrary non-negative exponents.

But
\[
 (x_1+\xi_1)^{\mu_1}
 =\sum^{\alpha_1}{\Pi(\mu_1)\over\Pi(\alpha_1)\Pi(\mu_1-\alpha_1)}
 \xi_1^{\alpha_1}x_1^{\mu_1-\alpha_1},
\]
and therefore
\begin{equation}
\begin{aligned}
&(x_1+\xi_1)^{\mu_1}\cdots(x_m+\xi_m)^{\mu_m}\\
&=\sum^{\alpha_1\ldots\alpha_m}
 {\Pi(\mu_1)\cdots\Pi(\mu_m)x_1^{\mu_1-\alpha_1}\cdots x_m^{\mu_m-\alpha_m}
 \over
 \Pi(\mu_1-\alpha_1)\cdots\Pi(\mu_m-\alpha_m)\Pi(\alpha_1)\cdots\Pi(\alpha_m)}
 \xi_1^{\alpha_1}\cdots\xi_m^{\alpha_m}.
\end{aligned}
\tag{6}
\end{equation}
Comparison with (4) and (6) therefore gives
\begin{equation}
\begin{aligned}
&D_{\alpha_1,\ldots,\alpha_m}(x_1^{\mu_1}\cdots x_m^{\mu_m})\\
&\quad =
{\Pi(\mu_1)\cdots\Pi(\mu_m)\over
 \Pi(\mu_1-\alpha_1)\cdots\Pi(\mu_m-\alpha_m)}
 x_1^{\mu_1-\alpha_1}\cdots x_m^{\mu_m-\alpha_m},
\end{aligned}
\tag{7}
\end{equation}
as long as $\alpha_1\leq\mu_1,\ldots,\alpha_m\leq\mu_m$. On the other hand,
\begin{equation}
 D_{\alpha_1,\ldots,\alpha_m}(x_1^{\mu_1}\cdots x_m^{\mu_m})=0,
\tag{8}
\end{equation}
as soon as one of the indices $\alpha$ is greater than the corresponding exponent $\mu$.


% ===== scan leaf 86, printed 61 =====
% !PRINTED 61
\weberpage{61}

Now let $\beta_1,\beta_2,\ldots,\beta_m$ denote a second system of indices, and form from the function (7) the derivative $D_{\beta_1,\beta_2,\ldots,\beta_m}$. By a repeated application of the same formulae (7) and (8) one obtains
\begin{equation}
 D_{\beta_1,\ldots,\beta_m}D_{\alpha_1,\ldots,\alpha_m}
 (x_1^{\mu_1}\cdots x_m^{\mu_m})
 =
 D_{\beta_1+\alpha_1,\ldots,\beta_m+\alpha_m}
 (x_1^{\mu_1}\cdots x_m^{\mu_m}),
\tag{9}
\end{equation}
and from this, by II, follows the general validity of the theorem
\[
\text{III.}\qquad
D_{\beta_1,\ldots,\beta_m}D_{\alpha_1,\ldots,\alpha_m}\Phi
 =
D_{\beta_1+\alpha_1,\ldots,\beta_m+\alpha_m}\Phi,
\]
which is a generalization of the theorem (8), \S~13; hence the higher derivatives can be formed by continued differentiation of the lower ones.

For the first derivatives of a function $\Phi$,
\[
D_{1,0,\ldots,0}\Phi,\quad D_{0,1,\ldots,0}\Phi,\ldots,D_{0,0,\ldots,1}\Phi,
\]
we also use the shorter signs
\[
\Phi'(x_1),\quad \Phi'(x_2),\ldots,\Phi'(x_m),
\]
and likewise for the second derivatives
\[
D_{2,0,\ldots,0},\quad D_{1,1,\ldots,0},\ldots
\]
the signs
\[
\Phi''(x_1,x_1),\quad \Phi''(x_1,x_2)=\Phi''(x_2,x_1),\ldots.
\]
This notation too can be generalized and would lead to an expression corresponding to formula (4), \S~15.

Because of their frequent use, we mention the formulae for the quadratic forms especially. Let
\begin{equation}
 \Phi(x)=\sum a_{i,k}x_i x_k,
\tag{10}
\end{equation}
where $i,k$ independently run through the range of values $1,2,\ldots,m$ and $a_{i,k}=a_{k,i}$. Then
\begin{equation}
\begin{aligned}
 {1\over2}\Phi'(x_1)&=a_{1,1}x_1+a_{1,2}x_2+\cdots+a_{1,m}x_m,\\
 {1\over2}\Phi'(x_2)&=a_{2,1}x_1+a_{2,2}x_2+\cdots+a_{2,m}x_m,\\
 &\quad\cdots\\
 {1\over2}\Phi'(x_m)&=a_{m,1}x_1+a_{m,2}x_2+\cdots+a_{m,m}x_m,
\end{aligned}
\tag{11}
\end{equation}
and we further put
\begin{equation}
 \Phi(x,\xi)=\Phi(\xi,x)=\sum \xi_i\Phi'(x_i)=\sum x_i\Phi'(\xi_i).
\tag{12}
\end{equation}
Then
\begin{equation}
 \Phi(x+\xi)=\Phi(x)+\Phi(x,\xi)+\Phi(\xi).
\tag{13}
\end{equation}


% ===== scan leaf 87, printed 62 =====
% !PRINTED 62
\weberpage{62}
\pdfbookmark[2]{§ 17. Euler's theorem on homogeneous functions}{bm20}

The function $\Phi(x,\xi)$ is called the \gesperrt{polar} of $\Phi$. It is linear and homogeneous both with respect to the $x$ and with respect to the $\xi$.

It can be expressed by
\begin{equation}
 \Phi(x,\xi)=2\sum a_{ik}\xi_i x_k
\tag{14}
\end{equation}
and satisfies the condition
\begin{equation}
 \Phi(x,x)=2\Phi(x).
\tag{15}
\end{equation}

\begin{center}
{\large \S.\ 17.}

\vspace{0.5em}
{\large\gesperrt{Euler's Theorem on Homogeneous Functions.}}
\end{center}

From the preceding developments one can easily derive a fundamental theorem on homogeneous functions, discovered by Euler and named after him.

We obtain it most simply from formula (4) of the preceding paragraph if, denoting by $t$ an arbitrary variable, we put
\begin{equation}
 \xi_1=tx_1,\quad \xi_2=tx_2,\ldots,\quad \xi_m=tx_m
\tag{1}
\end{equation}
and then apply the fundamental formula \S~15 (1) for homogeneous functions. We thus obtain first:
\begin{equation}
 (1+t)^n\Phi(x)
 =
 \sum { (tx_1)^{\alpha_1}(tx_2)^{\alpha_2}\cdots(tx_m)^{\alpha_m}
 \over \Pi(\alpha_1)\Pi(\alpha_2)\cdots\Pi(\alpha_m)}
 D_{\alpha_1,\alpha_2,\ldots,\alpha_m}\Phi.
\tag{2}
\end{equation}
If one applies the binomial theorem to the left-hand side of (2) and then equates on both sides the coefficients of equal powers of $t$, one obtains for every $\nu=1,2,\ldots,n$:
\begin{equation}
 {\Pi(n)\over\Pi(n-\nu)}\Phi(x_1,x_2,\ldots,x_m)
 =
 \sum^{\alpha}
 {\Pi(\nu)\over\Pi(\alpha_1)\Pi(\alpha_2)\cdots\Pi(\alpha_m)}
 x_1^{\alpha_1}x_2^{\alpha_2}\cdots x_m^{\alpha_m}
 D_{\alpha_1,\alpha_2,\ldots,\alpha_m}\Phi,
\tag{3}
\end{equation}
where the sum is extended over all systems of values of the $\alpha$ satisfying the condition
\begin{equation}
 \alpha_1+\alpha_2+\cdots+\alpha_m=\nu.
\tag{4}
\end{equation}

In this form the theorem to be proved is contained in its full generality. For the special case $\nu=1$ we obtain the formula
%CONT>>


% ===== scan leaf 88, printed 63 =====
% !PRINTED 63
\weberpage{63}
%<<CONT

\begin{equation}
 n\Phi(x_1,x_2,\ldots,x_m)
 =x_1\Phi'(x_1)+x_2\Phi'(x_2)+\cdots+x_m\Phi'(x_m),
\tag{5}
\end{equation}
of which formula (15) of the preceding paragraph is a special case, and for $\nu=2$:
\begin{equation}
 n(n-1)\Phi(x_1,x_2,\ldots,x_m)
 =
 \sum_{i,k=1}^{m}x_i x_k\Phi''(x_i,x_k),
\tag{6}
\end{equation}
where the sum is to be extended from $i=1$ to $i=m$ and from $k=1$ to $k=m$, so that every term with unequal $i,k$ occurs twice in the sum.

If, when the $\alpha$ are subject to the condition (4), we put
\begin{equation}
 \Phi_\nu(\xi,x)=
 \sum {\xi_1^{\alpha_1}\xi_2^{\alpha_2}\cdots\xi_m^{\alpha_m}
 \over \Pi(\alpha_1)\Pi(\alpha_2)\cdots\Pi(\alpha_m)}
 D_{\alpha_1,\alpha_2,\ldots,\alpha_m}\Phi,
\tag{7}
\end{equation}
then by (4) of \S~16,
\begin{equation}
 \Phi(x+\xi)=\Phi(x)+\Phi_1(x,\xi)+\Phi_2(x,\xi)+\cdots+\Phi_n(x,\xi),
\tag{8}
\end{equation}
and since the left-hand side remains unchanged when $x$ and $\xi$ are interchanged, the relation follows:
\begin{equation}
 \Phi_{n-\nu}(x,\xi)=\Phi_\nu(\xi,x),
\tag{9}
\end{equation}
and therefore in particular
\begin{equation}
 \Phi_n(x,\xi)=\Phi(\xi).
\tag{10}
\end{equation}
The function $\Phi_\nu(x,\xi)$, regarded as a function of $x$, is called the \gesperrt{$\nu$th polar} of the function $\Phi$ for the system of values $\xi$.

We shall now specialize formula (3) to the case of a binary form ($m=2$). We denote the variables by $x,y$ and put, for abbreviation,
\[
 \Phi(x,y)=u,\qquad D_{h,\nu-h}\Phi=u_h.
\]
% [source-original sic] Printed p.63 cites (4), although this specialization is derived from (3); corrected in translation.
From (3) we obtain
\begin{equation}
 {\Pi(n)\over\Pi(n-\nu)}u
 =
 \sum_{h=0}^{\nu}{\Pi(\nu)\over\Pi(h)\Pi(\nu-h)}
 u_h x^h y^{\nu-h},
\tag{9}% [source-original sic] Printed p.63 repeats (9) after (9) and (10); retained verbatim.
\end{equation}
where $\nu$ can assume any value not greater than $n$.


% ===== scan leaf 89, printed 64 =====
% !PRINTED 64
\weberpage{64}
\pdfbookmark[1]{Second Section. Determinants.}{bm21}
\pdfbookmark[2]{§ 18. Permutations of n elements}{bm22}

\begin{center}
{\large\gesperrt{Second Section.}}\\[0.5em]
{\Large\gesperrt{Determinants.}}\\[0.4em]
\end{center}

\begin{center}
{\large \S.\ 18.}

\vspace{0.5em}
{\large\gesperrt{Permutations of \(n\) Elements.}}
\end{center}

We consider a system of \(n\) distinct elements of any kind, for example the \(n\) digits
\[
 1,2,3,\ldots,n,
\]
whose complex, in this definite arrangement, we shall denote by \(\mathfrak A\). The elements of \(\mathfrak A\) can be arranged in different ways, for example
\[
 2,1,3,\ldots,n.
\]
The passage from one arrangement to another is called a \gesperrt{permutation}.

If we denote by \(\Pi(n)\) the number of the different arrangements, which can depend only on the number \(n\) of the elements, then first
\[
 \Pi(1)=1,\qquad \Pi(2)=2.
\]
To determine the number in general, we think of an \(n\)-th element as added to \(n-1\) elements. In every arrangement of the \(n-1\) elements the \(n\)-th element can now be put in \(n\) different places, namely before the first, between the first and the second, between the second and the third, and so on, finally after the \((n-1)\)-st; and all arrangements thus obtained are distinct. Hence
\begin{equation}
 \Pi(n)=n\Pi(n-1),
\tag{1}
\end{equation}
from which, by complete induction,
\begin{equation}
 \Pi(n)=1\cdot2\cdot3\cdots n
\tag{2}
\end{equation}
%CONT>>


% ===== scan leaf 90, printed 65 =====
% !PRINTED 65
\weberpage{65}
\pdfbookmark[2]{§ 19. Permutations of the first and second kind}{bm23}
%<<CONT

is obtained, so that here the sign \(\Pi(n)\) has the same meaning as in the first section (\S~7).

Any arrangement of the system \(\mathfrak A\) will be denoted by \(\mathfrak A'\), or more fully, if \(\alpha_1,\alpha_2,\ldots,\alpha_n\) are the digits \(1,2,\ldots,n\) in some order, by
\begin{equation}
 \mathfrak A'=\alpha_1,\alpha_2,\ldots,\alpha_n.
\tag{3}
\end{equation}

One can derive any arrangement from any other in very different ways, that is, one can carry out a permutation in different ways. Among the several possibilities, those carried out by so-called \gesperrt{transpositions}, that is, by successive interchange of only two elements, are of special interest for us. By several transpositions performed successively one can derive every arrangement, for example \(\mathfrak A'\), from \(\mathfrak A\). For this purpose one may proceed, for instance, by first interchanging in \(\mathfrak A\) the element \(1\) with that which stands in the first place in \(\mathfrak A'\), that is, with \(\alpha_1\) (unless \(\alpha_1=1\)); then, if \(\alpha_2\) is not already \(=2\), interchanging \(2\) with \(\alpha_2\), and so on.

Thus, for example, to pass from \((1,2,3,4)\) to \((4,3,2,1)\), one forms the arrangements
\[
 (1,2,3,4),\qquad (4,2,3,1),\qquad (4,3,2,1).
\]
If we briefly denote a transposition by the two interchanged digits, and therefore the interchange of \(1\) with \(2\) by \((1,2)\), then here we have performed successively the transpositions
\[
 (1,4),\quad (2,3).
\]

It is to be observed that the passage from one arrangement to a definite other one can be reached in infinitely many different ways by successive transpositions. Thus the arrangement \((1,2,3,4)\) also passes into \((4,3,2,1)\) by the transpositions
\[
 (1,2),\quad (1,3),\quad (2,4),\quad (1,2).
\]

\begin{center}
{\large \S.\ 19.}

\vspace{0.5em}
{\large\gesperrt{Permutations of the First and Second Kind.}}
\end{center}

The \(\Pi(n)\) arrangements of \(n\) elements can be separated into two kinds according to the following point of view.

From the \(n\) elements of our system one can form
\[
 {n(n-1)\over 2}
\]
and no more pairs. We now assign to the \(n\) elements
%CONT>>


% ===== scan leaf 91, printed 66 =====
% !PRINTED 66
\weberpage{66}
%<<CONT

\(1,2,\ldots,n\), in a definite way, \(n\) real numerical values \(a_1,a_2,\ldots,a_n\), and from these numerical values form the \(\frac{n(n-1)}2\) differences
\[
 a_1-a_2,\quad a_1-a_3,\ldots,
\]
where the lower index is to belong to the minuend. The product of differences
\begin{equation}
 P=(a_1-a_2)(a_1-a_3)\cdots(a_1-a_n)
    (a_2-a_3)\cdots(a_2-a_n)\cdots(a_{n-1}-a_n)
\tag{1}
\end{equation}
will, when \(a_1,a_2,\ldots,a_n\) are mutually distinct numerical values, have a value different from zero, for example a positive one if \(a_1>a_2>a_3>\cdots>a_n\) was assumed.

If now we interchange the indices \(1,2,3,\ldots,n\) in some way, and so pass, say, from \(\mathfrak A\) to \(\mathfrak A'\), then \(P\) passes into
\begin{equation}
\begin{aligned}
 P'={}&(a_{\alpha_1}-a_{\alpha_2})(a_{\alpha_1}-a_{\alpha_3})\cdots
       (a_{\alpha_1}-a_{\alpha_n})\\
     &\quad\cdot(a_{\alpha_2}-a_{\alpha_3})\cdots(a_{\alpha_2}-a_{\alpha_n})
       \cdots(a_{\alpha_{n-1}}-a_{\alpha_n}),
\end{aligned}
\tag{2}
\end{equation}
and, apart from the sign, this product consists of the same factors as \(P\). Thus \(P'\) is either equal to \(P\) or opposite to \(P\).

\gesperrt{I. We count the arrangement \(\mathfrak A'\), and therefore also the permutation which transforms \(\mathfrak A\) into \(\mathfrak A'\), as of the first or of the second kind according as \(P\) is equal or opposite to \(P'\); hence \(\mathfrak A\) itself belongs to the first kind.}

\gesperrt{II. By a simple transposition \((h,k)\), in which \(h,k\) denote any two of the digits \(1,2,\ldots,n\), both \(P\) and \(P'\) change their sign.}

For the factors which contain neither \(h\) nor \(k\) are not affected by this transposition; furthermore, in \(P\) and \(P'\) there is the factor \(\pm(a_h-a_k)\), and the pairs of factors
\[
 \pm(a_h-a_\nu)(a_k-a_\nu),
\]
where \(\nu\) runs through the series \(1,2,\ldots,n\) with the exception of \(h,k\). The first factor changes its sign, while the pair of factors remains unchanged under the transposition \((h,k)\). Hence follows:
%CONT>>


% ===== scan leaf 92, printed 67 =====
% !PRINTED 67
\weberpage{67}
%<<CONT

\gesperrt{III. The permutations of the first kind are composed of an even number of transpositions, and those of the second kind of an odd number.}

The so-called identical permutation, which leaves \(\mathfrak A\) unchanged, is therefore also to be counted as of the first kind.

From this one obtains further the conclusion: in however many different ways one may derive \(\mathfrak A'\) from \(\mathfrak A\) by transpositions, the number of these transpositions is in all these ways consistently even or odd, according as \(\mathfrak A'\) belongs to the first or to the second kind.

If in all arrangements
\[
 \mathfrak A,\mathfrak A',\mathfrak A'',\ldots
\]
of the \(n\) elements we perform one transposition, say \((1,2)\), then each of these arrangements passes into a definite other one, say \(\mathfrak A\) into \(\mathfrak B\), \(\mathfrak A'\) into \(\mathfrak B'\), \(\mathfrak A''\) into \(\mathfrak B''\), and so on; and if the same transposition is repeated once more, then \(\mathfrak B\) again passes into \(\mathfrak A\), \(\mathfrak B'\) again into \(\mathfrak A'\), and so on. Hence the arrangements \(\mathfrak B,\mathfrak B',\mathfrak B'',\ldots\) are all distinct from one another, and consequently in their totality agree with the totality of the \(\mathfrak A\)'s. Since, as we have seen above, all difference-products
\[
 P,P',P'',\ldots
\]
which are formed from \(P\) with the different arrangements \(\mathfrak A,\mathfrak A',\mathfrak A'',\ldots\) change their sign under a transposition, it follows that to each \(\mathfrak A\) of the first kind there corresponds a \(\mathfrak B\) of the second kind, and to each \(\mathfrak A\) of the second kind a \(\mathfrak B\) of the first kind.

\gesperrt{IV. Accordingly the number of arrangements of the first kind is just as large as the number of arrangements of the second kind, namely} \({1\over2}\Pi(n)\)\footnote{These theorems have here been obtained from the consideration of the product \(P\), hence from a numerical quantity. How one can arrive at the same results without using such a function will be seen in the fourteenth section.}.

For \(n=3\) we have the following six arrangements, of which the first horizontal row forms the first kind:
\begin{equation}
\begin{array}{ccc}
(1,2,3),&(2,3,1),&(3,1,2),\\
(3,2,1),&(2,1,3),&(1,3,2).
\end{array}
\tag{3}
\end{equation}


% ===== scan leaf 93, printed 68 =====
% !PRINTED 68
\weberpage{68}
\pdfbookmark[2]{§ 20. Determinants}{bm24}

\begin{center}
{\large \S.\ 20.}

\vspace{0.5em}
{\large\gesperrt{Determinants.}}
\end{center}

We now consider a system of \(n^2\) arbitrary quantities, with which the rational operations of arithmetic can be performed. For a simple notation of these quantities we choose one letter with a double index \(a_i^{(k)}\), where both \(i\) and \(k\) are to run through the series of digits \(1,2,3,\ldots,n\). For better survey we arrange these quantities in a square, so that all \(a\)'s with the same upper index stand in a horizontal row, and all \(a\)'s with the same lower index in a vertical row, and denote this square by \(\Delta\):
\begin{equation}
\Delta=
\begin{vmatrix}
 a_1^{(1)}&a_2^{(1)}&a_3^{(1)}&\cdots&a_n^{(1)}\\
 a_1^{(2)}&a_2^{(2)}&a_3^{(2)}&\cdots&a_n^{(2)}\\
 a_1^{(3)}&a_2^{(3)}&a_3^{(3)}&\cdots&a_n^{(3)}\\
 \vdots&\vdots&\vdots&&\vdots\\
 a_1^{(n)}&a_2^{(n)}&a_3^{(n)}&\cdots&a_n^{(n)}
\end{vmatrix}.
\tag{1}
\end{equation}
For brevity one calls the horizontal rows \gesperrt{rows}, and the vertical rows \gesperrt{columns}.

But by the square enclosed between vertical lines we shall understand not only the complex of the quantities \(a\), but a definite arithmetical combination of these quantities, which can be calculated as soon as the \(a\)'s are given numerically, and which we shall now describe.

Form the product of the terms standing in the diagonal running from upper left to lower right:
\begin{equation}
 M=a_1^{(1)}a_2^{(2)}a_3^{(3)}\cdots a_n^{(n)}.
\tag{2}
\end{equation}
From this derive \(\Pi(n)\) products \(M,M',M'',\ldots\) by permuting the lower indices, and give to each product so obtained the \gesperrt{positive or negative sign} according as the applied permutation belongs to the first or to the second kind; thus, in the notation of the preceding paragraph,
\begin{equation}
 M'=\pm a_{\alpha_1}^{(1)}a_{\alpha_2}^{(2)}\cdots a_{\alpha_n}^{(n)}.
\tag{3}
\end{equation}
The sum of these products
\[
M+M'+M''+\cdots=\Sigma M
\]
%CONT>>


% ===== scan leaf 94, printed 69 =====
% !PRINTED 69
\weberpage{69}
%<<CONT

is to be \(\Delta\). The quantity \(\Delta\) is called the \gesperrt{determinant} of the \(n^2\) elements \(a_i^{(k)}\), more precisely, when the distinction is necessary, an \gesperrt{\(n\)-rowed determinant} (also a determinant of the \(n\)-th degree or \(n\)-th order). The term \(M\) of this sum, that is, the product of all elements standing in the diagonal of the square, is called the \gesperrt{principal term}.

For example, if \(n=2\), we obtain
\begin{equation}
\Delta=a_1^{(1)}a_2^{(2)}-a_2^{(1)}a_1^{(2)},
\tag{4}
\end{equation}
and for \(n=3\) [according to (3) of the preceding paragraph]:
\begin{equation}
\begin{aligned}
\Delta={}&a_1^{(1)}a_2^{(2)}a_3^{(3)}
+a_2^{(1)}a_3^{(2)}a_1^{(3)}
+a_3^{(1)}a_1^{(2)}a_2^{(3)}\\
&-a_1^{(1)}a_3^{(2)}a_2^{(3)}
-a_2^{(1)}a_1^{(2)}a_3^{(3)}
-a_3^{(1)}a_2^{(2)}a_1^{(3)}.
\end{aligned}
\tag{5}
\end{equation}
Or in another notation:
\begin{equation}
\begin{vmatrix}
a&b\\
c&d
\end{vmatrix}
=ad-bc,
\tag{6}
\end{equation}
\begin{equation}
\begin{vmatrix}
a&b&c\\
a'&b'&c'\\
a''&b''&c''
\end{vmatrix}
=ab'c''+bc'a''+ca'b''-ac'b''-ba'c''-cb'a''.
\tag{7}
\end{equation}

The reader is advised to practise the calculation of such determinants on numerical examples.

The notation (1) is in many cases too cumbersome; other, shorter symbols are therefore also in use. Thus \gesperrt{Jacobi}, writing out only the principal term of the expanded determinant, puts
\begin{equation}
\Delta=\sum \pm a_1^{(1)}a_2^{(2)}\cdots a_n^{(n)},
\tag{8}
\end{equation}
and \gesperrt{Kronecker}, still more briefly,
\begin{equation}
\Delta=\left|a_i^{(k)}\right|.
\tag{9}
\end{equation}
Both notations are entirely clear only when the elements are denoted by two indices in the manner presupposed here, and are wholly inapplicable when, for example, the elements are given numerically.

Determinants sometimes occur for which
\[
a_i^{(k)}=a_k^{(i)},
\]
so that in (1) the elements standing symmetrically with respect to the diagonal of the square are equal to one another. In these cases we shall ordinarily place both indices below (to indicate their equivalence), and therefore put
\[
a_{i,k}=a_{k,i}.
\]
Such determinants are called \gesperrt{symmetric}.


% ===== scan leaf 95, printed 70 =====
% !PRINTED 70
\weberpage{70}
\pdfbookmark[2]{§ 21. Principal theorems on determinants}{bm25}

\begin{center}
{\large \S.\ 21.}

\vspace{0.5em}
{\large\gesperrt{Principal Theorems on Determinants.}}
\end{center}

From the concept of the determinant the first theorems useful for application are easily obtained.

If in the product [\S~20, (3)]
\begin{equation}
 M'=\pm a_{\alpha_1}^{(1)}a_{\alpha_2}^{(2)}\cdots a_{\alpha_n}^{(n)}
\tag{1}
\end{equation}
we rearrange the factors, its value is not changed. We can therefore also order the factors in such a way that the lower indices occur in their natural order \(1,2,\ldots,n\). The upper indices will then occur in a certain permuted order; thus \(M'\) will receive the form
\begin{equation}
 \pm a_1^{(\beta_1)}a_2^{(\beta_2)}\cdots a_n^{(\beta_n)},
\tag{2}
\end{equation}
where
\[
 (\beta_1,\beta_2,\ldots,\beta_n)=\mathfrak B
\]
is, just as \((\alpha_1,\alpha_2,\ldots,\alpha_n)=\mathfrak A\), an arrangement of the digits \(1,2,\ldots,n\). One obtains the arrangement \(\mathfrak B\) by undoing, in the factors of \(M'\), the transpositions which led to \(\mathfrak A\), beginning with the last, in order to recover the original arrangement in the row of the lower indices. The resulting order of the upper indices is then the arrangement \(\mathfrak B\). It follows that \(\mathfrak B\) belongs to the first or the second kind according as \(\mathfrak A\) belongs to the first or the second kind, since both arise through the same number of transpositions. The totality of the \(\mathfrak B\)'s represents, just like the totality of the \(\mathfrak A\)'s, all permutations of the \(n\) elements, since two different \(\mathfrak A\)'s can never lead to the same \(\mathfrak B\). Thus is proved:

\gesperrt{I. The determinant \(\Delta\) can also be formed by permuting the upper indices in the principal term \(a_1^{(1)}a_2^{(2)}\cdots a_n^{(n)}\) in all possible ways, assigning to every product so formed the positive or negative sign according as the applied permutation belongs to the first or second kind, and then taking the sum of all these products.}


% ===== scan leaf 96, printed 71 =====
% !PRINTED 71
\weberpage{71}

In the representation \S~20, (1) of \(\Delta\), the rows are characterized by the upper indices, the columns by the lower indices. Hence this theorem may also be expressed as follows:

\gesperrt{II. A determinant is not changed when the rows are made columns and the columns are made rows.}

If, in all arrangements \(\mathfrak A,\mathfrak A',\mathfrak A'',\ldots\) of the \(n\) elements, we interchange any two elements with one another, the totality of these arrangements remains unchanged, but every arrangement of the first kind passes into an arrangement of the second kind, and conversely. If, therefore, in the terms \(M,M',M'',\ldots\) of which \(\Delta\) is composed, any two lower indices are interchanged, every term with positive sign passes into another which in \(\Delta\) bore the negative sign, and conversely; therefore \(\Delta\) changes its sign. With the help of II follows the theorem:

\gesperrt{III. If in \(\Delta\) two lower or two upper indices are interchanged with one another, the determinant changes only its sign.}

Expressed somewhat differently: \gesperrt{if two rows or two columns are interchanged with one another, the determinant changes only its sign}; and from this more generally:

\gesperrt{IV. If in a determinant the rows or the columns are permuted, the absolute value is not changed, and the sign is not changed or passes into the opposite sign according as the applied permutation belongs to the first or second kind.}

From III one obtains the following fundamental theorem:

\gesperrt{V. If in two rows or in two columns the terms standing in the same place are equal to one another (more briefly, if two rows are equal to one another), then the determinant has the value zero.}

For the interchange of the two rows changes the sign by III, but on the other hand, since both rows are identical,
%CONT>>


% ===== scan leaf 97, printed 72 =====
% !PRINTED 72
\weberpage{72}
\pdfbookmark[2]{§ 22. Minors}{bm26}
%<<CONT

cannot change anything; thus for \(\Delta\) only the value zero remains.

One only expresses theorem V differently when one says:

\gesperrt{VI. One obtains a vanishing determinant if the elements of one row are replaced by the corresponding elements of another row, or, shortly, if one lower or upper index is replaced by another.}

\begin{center}
{\large \S.\ 22.}

\vspace{0.5em}
{\large\gesperrt{Subdeterminants.}}
\end{center}

In every term of the expanded determinant
\[
 \sum \pm a_1^{(1)}a_2^{(2)}\cdots a_n^{(n)},
\]
whose value we now wish to denote by \(A\), every one of the numbers \(1,2,\ldots,n\) occurs once and only once as a lower index. Thus a certain complex of terms will contain the factor \(a_1^{(1)}\), another complex the factor \(a_1^{(2)}\), and so on, finally a complex the factor \(a_1^{(n)}\); every term of the determinant occurs in one and only one of these complexes.

If, therefore, we denote the first of these complexes by \(a_1^{(1)}A_1^{(1)}\), the second by \(a_1^{(2)}A_1^{(2)}\), and the last by \(a_1^{(n)}A_1^{(n)}\), then the determinant can be represented in the following way:
\begin{equation}
 A=a_1^{(1)}A_1^{(1)}+a_1^{(2)}A_1^{(2)}+\cdots+a_1^{(n)}A_1^{(n)}.
\tag{1}
\end{equation}
Instead of the lower index \(1\), we could equally well single out any other index \(\nu\), and hence put
\begin{equation}
 A=a_\nu^{(1)}A_\nu^{(1)}+a_\nu^{(2)}A_\nu^{(2)}+\cdots+a_\nu^{(n)}A_\nu^{(n)}.
\tag{2}
\end{equation}
Here the product \(a_\nu^{(\mu)}A_\nu^{(\mu)}\) denotes the complex of all terms of the determinant which contain the factor \(a_\nu^{(\mu)}\).

Since the same rules hold for upper as for lower indices, the determinant can also be written in the following way:
\begin{equation}
 A=a_1^{(\mu)}A_1^{(\mu)}+a_2^{(\mu)}A_2^{(\mu)}
 +\cdots+a_n^{(\mu)}A_n^{(\mu)},
\tag{3}
\end{equation}
where \(\mu\), too, may denote any one of the indices \(1,2,\ldots,n\).


% ===== scan leaf 98, printed 73 =====
% !PRINTED 73
\weberpage{73}

The quantities \(A_\nu^{(\mu)}\) thereby completely defined are called the \gesperrt{subdeterminants} of the determinant \(A\). To become precisely acquainted with their mode of formation, let us first consider the complex \(a_1^{(1)}A_1^{(1)}\). It is obtained by leaving the lower index \(1\) unchanged in the product
\[
 a_1^{(1)}a_2^{(2)}\cdots a_n^{(n)}
\]
and permuting only the remaining indices \(2,3,\ldots,n\) in all ways, and forming the sum of the terms so obtained with regard to the rule of signs. Hence \(A_1^{(1)}\) is the \((n-1)\)-rowed determinant
\begin{equation}
 A_1^{(1)}=
\begin{vmatrix}
a_2^{(2)}&a_3^{(2)}&\cdots&a_n^{(2)}\\
a_2^{(3)}&a_3^{(3)}&\cdots&a_n^{(3)}\\
\vdots&\vdots&&\vdots\\
a_2^{(n)}&a_3^{(n)}&\cdots&a_n^{(n)}
\end{vmatrix},
\tag{4}
\end{equation}
or the determinant obtained from \(A\) by deleting, in the square representing \(A\) [\S~20, (1)], the first row and the first column.

From this one easily obtains the meaning of \(A_\nu^{(\mu)}\). By making \(\nu-1\) interchanges of rows one can bring the \(\nu\)-th row to the first, and by adding \(\mu-1\) interchanges of columns one can bring the \(\mu\)-th column to the first; apart from this the remaining rows and columns keep their relative order. The determinant itself has taken the factor \((-1)^{\mu+\nu}\), and has remained unchanged in absolute value. In the arrangement thus changed the element \(a_\nu^{(\mu)}\) has taken the place of the element \(a_1^{(1)}\), and from this one infers the following rule of formation:

\gesperrt{One obtains the subdeterminant \(A_\nu^{(\mu)}\) by deleting from the square representing the determinant the row and column which intersect in \(a_\nu^{(\mu)}\), and adjoining the factor \((-1)^{\mu+\nu}\).}

For example, for the determinant of three rows, one obtains the representation
\begin{equation}
\begin{aligned}
\begin{vmatrix}
a&b&c\\
a'&b'&c'\\
a''&b''&c''
\end{vmatrix}
&=a\begin{vmatrix}b'&c'\\b''&c''\end{vmatrix}
-b\begin{vmatrix}a'&c'\\a''&c''\end{vmatrix}
+c\begin{vmatrix}a'&b'\\a''&b''\end{vmatrix}\\
&=a(b'c''-c'b'')+b(c'a''-a'c'')+c(a'b''-b'a'').
\end{aligned}
\tag{5}
\end{equation}


% ===== scan leaf 99, printed 74 =====
% !PRINTED 74
\weberpage{74}

Since the lower index \(\nu\) does not occur at all in \(A_\nu^{(\mu)}\), \(A_\nu^{(\mu)}\) is not changed if the lower index \(\nu\) is replaced by another. But then, by \S~21, VI, the determinant vanishes. We therefore obtain from (2) the important relation, in which \(\mu,\nu\) can be any two distinct digits \(1,2,\ldots,n\):
\begin{equation}
0=a_\mu^{(1)}A_\nu^{(1)}+a_\mu^{(2)}A_\nu^{(2)}
 +\cdots+a_\mu^{(n)}A_\nu^{(n)},
\tag{6}
\end{equation}
and likewise from (3)
\begin{equation}
0=a_1^{(\mu)}A_1^{(\nu)}+a_2^{(\mu)}A_2^{(\nu)}
 +\cdots+a_n^{(\mu)}A_n^{(\nu)}.
\tag{7}
\end{equation}

For instance from (5), if \(a,b,c\) are replaced by \(a',b',c'\), there follows
\begin{equation}
 a'(b'c''-c'b'')+b'(c'a''-a'c'')+c'(a'b''-b'a'')=0,
\tag{8}
\end{equation}
a formula of whose correctness one is convinced by the simplest calculation.

If we multiply the relation (6) by an arbitrary factor \(\lambda\) and add it to (2), we obtain the formula
\begin{equation}
\begin{aligned}
 A={}&(a_\nu^{(1)}+\lambda a_\mu^{(1)})A_\nu^{(1)}
 +(a_\nu^{(2)}+\lambda a_\mu^{(2)})A_\nu^{(2)}\\
&\quad+\cdots+(a_\nu^{(n)}+\lambda a_\mu^{(n)})A_\nu^{(n)},
\end{aligned}
\tag{9}
\end{equation}
which expresses for us the following theorem:

\gesperrt{VII. The determinant does not change its value if to the elements of one row one adds the corresponding elements of another row multiplied by an arbitrary common factor.}

The same theorem holds also for the columns. It is often used with advantage for the simplification and numerical calculation of determinants. We add the following theorems, which can be read off at once from the representations (2), (3).

\gesperrt{VIII. If all elements of a row or of a column have a common factor, then this factor can be omitted from the row or column and put as a factor before the determinant.}

For by (2) it is
\[
 p a_\nu^{(1)}A_\nu^{(1)}+p a_\nu^{(2)}A_\nu^{(2)}+\cdots+p a_\nu^{(n)}A_\nu^{(n)}=pA.
\]

\gesperrt{IX. If in a row or in a column all elements but one vanish, then the determinant reduces}
%CONT>>


% ===== scan leaf 100, printed 75 =====
% !PRINTED 75
\weberpage{75}
%<<CONT

\gesperrt{to the product of this one element and the corresponding subdeterminant.}

For if \(a_\nu^{(1)},a_\nu^{(2)},\ldots,a_\nu^{(n)}\), with the exception of \(a_\nu^{(\mu)}\), vanish, then by (2)
\[
 A=a_\nu^{(\mu)}A_\nu^{(\mu)}.
\]
The value of \(A\) is then wholly independent of the \(a_1^{(\mu)},a_2^{(\mu)},\ldots,a_n^{(\mu)}\) with the exception of \(a_\nu^{(\mu)}\).

To make an application of these theorems, we determine the value of the determinant
\[
 \Delta=
 \begin{vmatrix}
 1&a&a^2\\
 1&b&b^2\\
 1&c&c^2
 \end{vmatrix},
\]
where \(a,b,c\) are arbitrary quantities. If we multiply the second column by \(a\) and subtract it from the third, then multiply the first by \(a\) and subtract it from the second, then by VII
\[
 \Delta=
 \begin{vmatrix}
 1&0&0\\
 1&b-a&b(b-a)\\
 1&c-a&c(c-a)
 \end{vmatrix},
\]
and by IX
\[
 \Delta=
 \begin{vmatrix}
 b-a&b(b-a)\\
 c-a&c(c-a)
 \end{vmatrix},
\]
and finally by VIII
\begin{equation}
 \Delta=(b-a)(c-a)
 \begin{vmatrix}
 1&b\\
 1&c
 \end{vmatrix}
 =(b-a)(c-a)(c-b).
\tag{10}
\end{equation}

In the same way one can also treat the \(n\)-rowed determinant
\[
\begin{vmatrix}
1&a_1&a_1^2&\cdots&a_1^{n-1}\\
1&a_2&a_2^2&\cdots&a_2^{n-1}\\
\vdots&\vdots&\vdots&&\vdots\\
1&a_n&a_n^2&\cdots&a_n^{n-1}
\end{vmatrix}
\]
and finds its value equal to
\begin{equation}
 (a_2-a_1)(a_3-a_1)\cdots(a_n-a_1)
 (a_3-a_2)\cdots(a_n-a_2)\cdots(a_n-a_{n-1}).
\tag{11}
\end{equation}


% ===== scan leaf 101, printed 76 =====
% !PRINTED 76
\weberpage{76}

If one orders the columns in the reverse sequence, then, according as \(n\) is even or odd,
\[
 {n\over2}\quad\hbox{or}\quad {n-1\over2}
\]
interchanges are necessary, so that the determinant thus ordered differs from \(\Delta\) by the factor
\[
 (-1)^{n(n-1)/2}.
\]
It comes to the same thing if one gives the opposite sign to the
\[
 {n(n-1)\over2}
\]
factors of the product (11). Thus, along with (11), the equality also holds:
\begin{equation}
\begin{vmatrix}
a_1^{\,n-1}&a_1^{\,n-2}&\cdots&a_1&1\\
a_2^{\,n-1}&a_2^{\,n-2}&\cdots&a_2&1\\
\vdots&\vdots&&\vdots&\vdots\\
a_n^{\,n-1}&a_n^{\,n-2}&\cdots&a_n&1
\end{vmatrix}
=
\begin{gathered}
(a_1-a_2)(a_1-a_3)\cdots(a_1-a_n)\\
(a_2-a_3)\cdots(a_2-a_n)\\
\cdots\\
(a_{n-1}-a_n).
\end{gathered}
\tag{12}
\end{equation}

We still mention here a notation for the subdeterminants which is borrowed from differential calculus and is often used with advantage, especially when derivatives are to be formed. If the quantities \(a_i^{(k)}\) are regarded as independent variables, then the determinant \(A\) is only of the first degree with respect to each of them. The derivative or differential quotient taken with respect to \(a_i^{(k)}\) is therefore equal to the coefficient of \(a_i^{(k)}\), hence
\[
 {\partial A\over\partial a_i^{(k)}}=A_i^{(k)}.
\]
If therefore, for example, the \(a_i^{(k)}\) are functions of a variable \(t\), then \(A\) is also a function of \(t\), and by the first rules of differential calculus one obtains the derivative of \(A\) with respect to \(t\) in the form
\[
 A'(t)=\sum A_i^{(k)}\,{\partial a_i^{(k)}\over\partial t},
\]
where \({\partial a_i^{(k)}}/{\partial t}\) denotes the derivative of \(a_i^{(k)}\) with respect to \(t\).


% ===== scan leaf 102, printed 77 =====
% !PRINTED 77
\weberpage{77}
\pdfbookmark[2]{§ 23. Minors in the wider sense}{bm27}

\begin{center}
{\large \S.\ 23.}

\vspace{0.5em}
{\large\gesperrt{The Subdeterminants in the Wider Sense.}}
\end{center}

We can now generalize the considerations of the preceding paragraph in the following way.  Just as there we started from the problem of seeking out all terms in the developed determinant \(A\) which contain the factor \(a_1^{(1)}\), so now we seek all terms which contain the factor
\[
        a_1^{(1)}a_2^{(2)}\cdots a_\nu^{(\nu)},
\]
where \(\nu\) may be any number less than \(n\).  These terms are obtained from the principal term
\[
        a_1^{(1)}a_2^{(2)}\cdots a_\nu^{(\nu)}a_{\nu+1}^{(\nu+1)}\cdots a_n^{(n)}
\]
if, in permuting the lower indices, one leaves \(1,2,\ldots,\nu\) unchanged and permutes only \(\nu+1,\ldots,n\) in all possible ways, with observance of the rule of signs.  Thus the aggregate of the required terms is
\begin{equation}
        a_1^{(1)}a_2^{(2)}\cdots a_\nu^{(\nu)}
        \begin{vmatrix}
        a_{\nu+1}^{(\nu+1)}&\cdots&a_n^{(\nu+1)}\\
        \vdots&&\vdots\\
        a_{\nu+1}^{(n)}&\cdots&a_n^{(n)}
        \end{vmatrix}.
\tag{1}
\end{equation}

\gesperrt{I. The determinant of \(n-\nu\) rows which occurs here as a factor, and which we denote by}
\[
        A_{1,2,\ldots,\nu}^{1,2,\ldots,\nu},
\]
\gesperrt{arises from \(A\) by suppressing the first \(\nu\) rows and columns.}

We shall now generalize this result in the following manner.  Choose any \(\nu\) elements
\[
        a_{\alpha_1}^{(\beta_1)},
        a_{\alpha_2}^{(\beta_2)},\ldots,
        a_{\alpha_\nu}^{(\beta_\nu)}
\]
so that no two of them occur in the same row or in the same column, that is, so that neither a lower nor an upper index is repeated.  We denote by
%CONT>>


% ===== scan leaf 103, printed 78 =====
% !PRINTED 78
\weberpage{78}
%<<CONT

\begin{equation}
        a_{\alpha_1}^{(\beta_1)}a_{\alpha_2}^{(\beta_2)}\cdots a_{\alpha_\nu}^{(\beta_\nu)}\,
        A_{\alpha_1,\alpha_2,\ldots,\alpha_\nu}^{\beta_1,\beta_2,\ldots,\beta_\nu}
\tag{2}
\end{equation}
the aggregate of those terms of the determinant which contain the product of these elements as a factor.

By interchanging rows and columns, whereby at most the sign of the determinant is changed, one can always cause the elements
\begin{equation}
        a_{\alpha_1}^{(\beta_1)},
        a_{\alpha_2}^{(\beta_2)},\ldots,
        a_{\alpha_\nu}^{(\beta_\nu)}
\tag{3}
\end{equation}
to take the place of the elements. % [sic] Weber prints a full stop in mid-sentence before the following display; retained source-faithfully
\[
        a_1^{(1)},a_2^{(2)},\ldots,a_\nu^{(\nu)}
\]
Then rule I applies to the determination of \(A_{\alpha_1,\alpha_2,\ldots,\alpha_\nu}^{\beta_1,\beta_2,\ldots,\beta_\nu}\), and one obtains:

\gesperrt{II. Apart from sign, \(A_{\alpha_1,\alpha_2,\ldots,\alpha_\nu}^{\beta_1,\beta_2,\ldots,\beta_\nu}\) is obtained as an \((n-\nu)\)-rowed determinant by deleting from \(A\) all rows and all columns which meet in one of the elements (3), and by leaving the remaining rows and columns in their order.}

For the determination of the sign one uses the following rule.  Arrange the lower and upper indices \(1,2,\ldots,n\) in the orders
\begin{equation}
        \alpha_1,\alpha_2,\ldots,\alpha_\nu,\alpha_{\nu+1},\ldots,\alpha_n,
\tag{4}
\end{equation}
\begin{equation}
        \beta_1,\beta_2,\ldots,\beta_\nu,\beta_{\nu+1},\ldots,\beta_n,
\tag{5}
\end{equation}
where the missing \(\alpha_{\nu+1},\ldots,\alpha_n\), and likewise \(\beta_{\nu+1},\ldots,\beta_n\), are taken \gesperrt{in increasing order}.

\gesperrt{III. The \((n-\nu)\)-rowed determinant described in II receives the positive or negative sign according as the two arrangements (4), (5) of the digits \(1,2,\ldots,n\) both belong to the same kind or to different kinds.}

Indeed, the determinant changes its sign through each interchange of two lower or two upper indices.  In order to reduce the general case (2) to the special case (1), one has to make just so many transpositions of upper and lower indices that the permutations (4) and (5) both pass into the original arrangement \(1,2,3,\ldots,n\); the same number of changes of sign has then taken place.


% ===== scan leaf 104, printed 79 =====
% !PRINTED 79
\weberpage{79}

The quantities so defined
\[
        A_{\alpha_1,\alpha_2,\ldots,\alpha_\nu}^{\beta_1,\beta_2,\ldots,\beta_\nu}
\]
\gesperrt{are called the \(\nu\)-th subdeterminants, or subdeterminants of the \(\nu\)-th order.} They are represented by \((n-\nu)\)-rowed determinants. With respect to these subdeterminants, III implies the theorem:

\gesperrt{IV. The subdeterminant \(A_{\alpha_1,\alpha_2,\ldots,\alpha_\nu}^{\beta_1,\beta_2,\ldots,\beta_\nu}\) changes only its sign if two of its lower or two of its upper indices are interchanged; more generally, it remains unaltered in absolute value if the order of the indices \(\alpha_1,\alpha_2,\ldots,\alpha_\nu\) is replaced by any other order, and changes sign or not according as this permutation is of the second or of the first kind.}

If \(\alpha_1',\alpha_2',\ldots,\alpha_\nu'\) denotes any ordering of \(\alpha_1,\alpha_2,\ldots,\alpha_\nu\), then the determinant \(A\) also contains the complex of terms
\[
        \pm a_{\alpha_1'}^{(\beta_1)}a_{\alpha_2'}^{(\beta_2)}\cdots
        a_{\alpha_\nu'}^{(\beta_\nu)}
        A_{\alpha_1,\alpha_2,\ldots,\alpha_\nu}^{\beta_1,\beta_2,\ldots,\beta_\nu}.
\]
Collecting all these terms gives the complex
\begin{equation}
        A_{\alpha_1,\alpha_2,\ldots,\alpha_\nu}^{\beta_1,\beta_2,\ldots,\beta_\nu}
        \Sigma \pm a_{\alpha_1}^{(\beta_1)}a_{\alpha_2}^{(\beta_2)}\cdots
        a_{\alpha_\nu}^{(\beta_\nu)}.
\tag{6}
\end{equation}
The \(\nu\)-rowed determinant which occurs in (6),
\[
        \Sigma\pm a_{\alpha_1}^{(\beta_1)}a_{\alpha_2}^{(\beta_2)}\cdots a_{\alpha_\nu}^{(\beta_\nu)}=
        \begin{vmatrix}
        a_{\alpha_1}^{(\beta_1)}&a_{\alpha_2}^{(\beta_1)}&\cdots&a_{\alpha_\nu}^{(\beta_1)}\\
        a_{\alpha_1}^{(\beta_2)}&a_{\alpha_2}^{(\beta_2)}&\cdots&a_{\alpha_\nu}^{(\beta_2)}\\
        \vdots&\vdots&&\vdots\\
        a_{\alpha_1}^{(\beta_\nu)}&a_{\alpha_2}^{(\beta_\nu)}&\cdots&a_{\alpha_\nu}^{(\beta_\nu)}
        \end{vmatrix},
\]
is called the subdeterminant complementary to
\(A_{\alpha_1,\alpha_2,\ldots,\alpha_\nu}^{\beta_1,\beta_2,\ldots,\beta_\nu}\), and is denoted by
\(B_{\alpha_1,\alpha_2,\ldots,\alpha_\nu}^{\beta_1,\beta_2,\ldots,\beta_\nu}\).
It contains exactly the rows and columns missing in \(A_{\alpha_1,\alpha_2,\ldots,\alpha_\nu}^{\beta_1,\beta_2,\ldots,\beta_\nu}\), and,
%CONT>>


% ===== scan leaf 105, printed 80 =====
% !PRINTED 80
\weberpage{80}
%<<CONT

apart from sign, agrees with the subdeterminant of the \((n-\nu)\)-th order
\[
        A_{\alpha_{\nu+1}\ldots\alpha_n}^{\beta_{\nu+1}\ldots\beta_n}.
\]
Thus the complex of terms (6) is denoted by
\begin{equation}
        A_{\alpha_1,\alpha_2,\ldots,\alpha_\nu}^{\beta_1,\beta_2,\ldots,\beta_\nu}
        B_{\alpha_1,\alpha_2,\ldots,\alpha_\nu}^{\beta_1,\beta_2,\ldots,\beta_\nu}.
\tag{7}
\end{equation}

If now, for \(\alpha_1,\alpha_2,\ldots,\alpha_\nu\), one chooses each combination of \(\nu\) of the digits \(1,2,\ldots,n\), whose number (according to \S~7) is \(B_\nu^{(n)}\), while \(\beta_1,\beta_2,\ldots,\beta_\nu\) is held fixed, one obtains the same number of complexes of the form (7), and every term of the determinant \(A\) occurs in one and only one of them.

\gesperrt{V. Consequently, by summing all expressions (7), one obtains the determinant \(A\):}
\begin{equation}
        A=\sum^{\alpha}
        A_{\alpha_1,\alpha_2,\ldots,\alpha_\nu}^{\beta_1,\beta_2,\ldots,\beta_\nu}
        B_{\alpha_1,\alpha_2,\ldots,\alpha_\nu}^{\beta_1,\beta_2,\ldots,\beta_\nu}.
\tag{8}
\end{equation}
Of course one may also keep the combination of the \(\alpha\)'s fixed and sum with respect to the \(\beta\)'s.

This is Laplace's theorem.

Another representation of the determinant \(A\) by the first and second subdeterminants is obtained as follows.

Choose in \(A\) a row and a column which intersect in an element, say \(a_\nu^{(\mu)}\). In every term of \(A\) there occurs one element with lower index \(\nu\) and one element with upper index \(\mu\). Thus in \(A\) we have first the complex \(a_\nu^{(\mu)}A_\nu^{(\mu)}\), and furthermore the various complexes \(a_\nu^{(i)}a_k^{(\mu)}A_{\nu,k}^{i,\mu}\), in which \(i\) may denote any index different from \(\mu\), and \(k\) any index different from \(\nu\).

\gesperrt{VI. We can therefore put:}
\begin{equation}
        A=a_\nu^{(\mu)}A_\nu^{(\mu)}+
        \sum^{i,k} a_\nu^{(i)}a_k^{(\mu)}A_{\nu,k}^{i,\mu},
\tag{9}
\end{equation}
or, according to IV,
\begin{equation}
        A=a_\nu^{(\mu)}A_\nu^{(\mu)}-
        \sum^{i,k} a_\nu^{(i)}a_k^{(\mu)}A_{\nu,k}^{\mu,i}.
\tag{10}
\end{equation}
We further observe concerning this theorem that \(A_{\nu,k}^{\mu,i}\) is the \gesperrt{first subdeterminant}, corresponding to the element \(a_k^{(i)}\), of the \((n-1)\)-rowed determinant \(A_\nu^{(\mu)}\); for \(A_{\nu,k}^{\mu,i}\) is the coefficient
%CONT>>


% ===== scan leaf 106, printed 81 =====
% !PRINTED 81
\weberpage{81}
\pdfbookmark[2]{§ 24. Linear homogeneous equations}{bm28}
%<<CONT

of \(a_\nu^{(\mu)}a_k^{(i)}\) in the expansion of \(A\), and \(A_\nu^{(\mu)}\) is the coefficient of \(a_\nu^{(\mu)}\); consequently \(A_{\nu,k}^{\mu,i}\) is the coefficient of \(a_k^{(i)}\) in the determinant \(A_\nu^{(\mu)}\).

According to this theorem, the so-called bordered determinant
\begin{equation}
        U=\begin{vmatrix}
        a_1^{(1)}&a_2^{(1)}&\cdots&a_n^{(1)}&u_1\\
        a_1^{(2)}&a_2^{(2)}&\cdots&a_n^{(2)}&u_2\\
        \vdots&\vdots&&\vdots&\vdots\\
        a_1^{(n)}&a_2^{(n)}&\cdots&a_n^{(n)}&u_n\\
        v_1&v_2&\cdots&v_n&q
        \end{vmatrix}
\tag{11}
\end{equation}
can be expanded with respect to the elements of the last row and column, and one obtains:
\begin{equation}
        U=qA-\sum^{i}\sum^{k}u_i v_k A_k^{(i)}.
\tag{12}
\end{equation}
This equation is obtained from (10) by replacing \(n\) by \(n+1\), and distinguishing the elements of the last row and column by a different notation.

For the higher subdeterminants also, notation by differential quotients is sometimes convenient, so that, for example,
\begin{equation}
        A_{\nu,k}^{i,\mu}=
        {\partial^2 A\over \partial a_\nu^{(i)}\,\partial a_k^{(\mu)}}
\tag{13}
\end{equation}
is put.

\begin{center}
{\large \S.\ 24.}

\vspace{0.5em}
{\large\gesperrt{Linear Homogeneous Equations.}}
\end{center}

The principal application of determinants, to which the whole theory owes its origin, is the solution of linear equations.

We shall attack the problem here at once in its most general form, since the special form affords scarcely any simplification and can afterwards easily be derived from the general result.

We consider a system of \(m\) equations of the first degree in which \(n\) unknowns \(x_1,x_2,\ldots,x_n\) occur homogeneously:
%CONT>>


% ===== scan leaf 107, printed 82 =====
% !PRINTED 82
\weberpage{82}
%<<CONT

\begin{equation}
\begin{aligned}
 a_1^{(1)}x_1+a_2^{(1)}x_2+\cdots+a_n^{(1)}x_n&=0,\\
 a_1^{(2)}x_1+a_2^{(2)}x_2+\cdots+a_n^{(2)}x_n&=0,\\
 &\ \vdots\\
 a_1^{(m)}x_1+a_2^{(m)}x_2+\cdots+a_n^{(m)}x_n&=0.
\end{aligned}
\tag{1}
\end{equation}
where the coefficients \(a_i^{(k)}\) are regarded as given quantities. For the present we shall make no assumption whatever about the numbers \(m,n\), but set ourselves the general task of determining \gesperrt{all systems of values} of \(x_1,x_2,\ldots,x_n\) which satisfy equations (1).

We can give one solution of equations (1) at once: whatever the coefficients \(a_i^{(k)}\) may be, they are satisfied when
\begin{equation}
        x_1=0,\quad x_2=0,\quad\ldots,\quad x_n=0.
\tag{2}
\end{equation}
We may also mention another extreme case: if all the coefficients \(a_i^{(k)}\) have the value zero, then equations (1) are satisfied for arbitrary values of \(x_1,x_2,\ldots,x_n\).

We precede the general answer to the question by the following remarks. We write the system of coefficients of (1) in the form of a rectangle
\begin{equation}
\begin{array}{cccc}
 a_1^{(1)}&a_2^{(1)}&\cdots&a_n^{(1)}\\
 a_1^{(2)}&a_2^{(2)}&\cdots&a_n^{(2)}\\
 \vdots&\vdots&&\vdots\\
 a_1^{(m)}&a_2^{(m)}&\cdots&a_n^{(m)}.
\end{array}
\tag{3}
\end{equation}
Such a scheme, which by itself has no numerical meaning, is called a \gesperrt{matrix}, insofar as it is regarded as the source of a larger number of determinants. \gesperrt{The determinants arising from the matrix are obtained by deleting arbitrary rows and columns, in an arbitrary number limited only so far that the remaining elements form a square, and by regarding this square as a determinant.} Thus one obtains from the matrix one-rowed, two-rowed, and so forth, determinants. The highest determinants are \(n\)- or \(m\)-rowed, according as \(n\) or \(m\) is the smaller number (or \(n\)-rowed if \(n=m\)).


% ===== scan leaf 108, printed 83 =====
% !PRINTED 83
\weberpage{83}

We now make the assumption that among the \(\nu\)-rowed determinants of the matrix \gesperrt{at least one} is different from zero, while the \((\nu+1)\)-rowed and consequently also the higher determinants, if such exist, shall all vanish. \(\nu\) may be any number which is not greater than the smaller of the two numbers \(n\) and \(m\) (or, if \(n=m\), does not exceed this common value).

Such a number \(\nu\) can always be found if we exclude the already disposed of, wholly uninteresting case in which all the coefficients \(a_i^{(k)}\) vanish.

Without limiting generality, we may, to simplify the notation, assume that the non-vanishing \(\nu\)-rowed determinant is
\begin{equation}
        A=\begin{vmatrix}
        a_1^{(1)}&a_2^{(1)}&\cdots&a_\nu^{(1)}\\
        a_1^{(2)}&a_2^{(2)}&\cdots&a_\nu^{(2)}\\
        \vdots&\vdots&&\vdots\\
        a_1^{(\nu)}&a_2^{(\nu)}&\cdots&a_\nu^{(\nu)}
        \end{vmatrix}.
\tag{4}
\end{equation}
For it is plainly open to us to arrange the system of equations (1) in any order, and furthermore we can choose the notation of the unknowns \(x\) so that any \(\nu\) of them are the first \(\nu\). We denote the subdeterminants of \(A\), as before, by \(A_i^{(k)}\), where \(i,k\) range from \(1\) to \(\nu\).

\gesperrt{I. If first \(\nu=n\), which presupposes that \(m\) is not smaller than \(n\), then equations (1) have no solution other than that contained in equations (2).}

For let us take the first \(n\) equations of (1):
\begin{equation}
\begin{aligned}
 a_1^{(1)}x_1+a_2^{(1)}x_2+\cdots+a_n^{(1)}x_n&=0,\\
 a_1^{(2)}x_1+a_2^{(2)}x_2+\cdots+a_n^{(2)}x_n&=0,\\
 &\vdots\\
 a_1^{(n)}x_1+a_2^{(n)}x_2+\cdots+a_n^{(n)}x_n&=0,
\end{aligned}
\tag{5}
\end{equation}
and multiply these in succession by \(A_\mu^{(1)},A_\mu^{(2)},\ldots,A_\mu^{(n)}\), where \(\mu\) may be any one of the indices \(1,2,\ldots,n\), and add them. Then, since by \S.\,22, (2) and (6),
\[
 \sum_{1,\nu}^{i} A_\mu^{(i)}a_\lambda^{(i)}=0\quad\hbox{or}\quad=A
\]
%CONT>>


% ===== scan leaf 109, printed 84 =====
% !PRINTED 84
\weberpage{84}
%<<CONT

holds according as \(\lambda\) is different from \(\mu\) or not,
\[
 A x_\mu=0,
\]
and since by our hypothesis \(A\) is different from zero,
\[
 x_\mu=0.
\]

We single out the most frequently used special case \(m=n\), and give the theorem for this case in the following form:

\gesperrt{II. If a system of \(n\) linear homogeneous equations with \(n\) unknowns has a determinant different from zero, then all the unknowns have the value zero; or:}

\gesperrt{If a system of \(n\) linear equations with as many homogeneously occurring unknowns has a solution in which not all the unknowns vanish, then the determinant of the system vanishes.}

Here, by the determinant of a system of \(n\) linear homogeneous equations with \(n\) unknowns is meant the determinant formed from the \(n^2\) coefficients of these equations.

We next consider the case in which \(\nu\) is smaller than \(n\). Since \(m\) must be equal to or greater than \(\nu\), we choose the first \(\nu\) equations of the system (1), and write them thus:
\begin{equation}
\begin{aligned}
 a_1^{(1)}x_1+a_2^{(1)}x_2+\cdots+a_\nu^{(1)}x_\nu&=-a_{\nu+1}^{(1)}x_{\nu+1}-\cdots-a_n^{(1)}x_n,\\
 a_1^{(2)}x_1+a_2^{(2)}x_2+\cdots+a_\nu^{(2)}x_\nu&=-a_{\nu+1}^{(2)}x_{\nu+1}-\cdots-a_n^{(2)}x_n,\\
 &\vdots\\
 a_1^{(\nu)}x_1+a_2^{(\nu)}x_2+\cdots+a_\nu^{(\nu)}x_\nu&=-a_{\nu+1}^{(\nu)}x_{\nu+1}-\cdots-a_n^{(\nu)}x_n.
\end{aligned}
\tag{6}
\end{equation}
We again let \(\mu\) denote one of the indices \(1,2,\ldots,\nu\), multiply equations (6) in succession by \(A_\mu^{(1)},A_\mu^{(2)},\ldots,A_\mu^{(\nu)}\), and add them. It follows, as before, using \S.\,22, (2), (6):
\begin{equation}
 A x_\mu=-C_{\nu+1,\mu}x_{\nu+1}-\cdots-C_{n,\mu}x_n,
\tag{7}
\end{equation}
where, for abbreviation, we have put
\begin{equation}
 C_{\lambda,\mu}=\sum_{1,\nu}^{i}a_\lambda^{(i)}A_\mu^{(i)},
 \qquad \lambda=\nu+1,\nu+2,\ldots,n.
\tag{8}
\end{equation}
By \S.\,22, (2), \(C_{\lambda,\mu}\) is the determinant obtained from the determinant defined by (4) by replacing the
%CONT>>


% ===== scan leaf 110, printed 85 =====
% !PRINTED 85
\weberpage{85}
%<<CONT

elements of the \(\mu\)th column \(a_\mu^{(1)},a_\mu^{(2)},\ldots,a_\mu^{(\nu)}\) by \(a_\lambda^{(1)},a_\lambda^{(2)},\ldots,a_\lambda^{(\nu)}\). By (7), since \(A\) is different from zero, \(x_1,x_2,\ldots,x_\nu\) are now expressed linearly in terms of \(x_{\nu+1},\ldots,x_n\) and the known quantities.

It remains to show

\gesperrt{III. that the expressions (7) satisfy equations (1), whatever values \(x_{\nu+1},\ldots,x_n\) may have; thus \(n-\nu\) of the unknowns remain arbitrary, and the other \(\nu\) depend on them according to (7).}

To see the truth of this theorem, we need only substitute the expressions (7) in equations (1). The calculation is greatly simplified by use of a summation sign \(\Sigma\), to which, as above, we append the summation letters above and the limits below. We can then first write equations (7) as
\begin{equation}
 A x_\mu=-\sum_{h=\nu+1}^{n}\sum_{i=1}^{\nu}a_h^{(i)}A_\mu^{(i)}x_h,
 \qquad \mu=1,2,\ldots,\nu.
\tag{9}
\end{equation}
We multiply by \(a_\mu^{(k)}\), where \(k\) denotes any one of the digits \(1,2,\ldots,m\), and sum with respect to \(\mu\):
\begin{equation}
 A\sum_{\mu=1}^{\nu}a_\mu^{(k)}x_\mu
 =-\sum_{h=\nu+1}^{n}x_h\sum_{i=1}^{\nu}\sum_{\mu=1}^{\nu}
 a_h^{(i)}a_\mu^{(k)}A_\mu^{(i)},
 \qquad k=1,2,\ldots,m.
\tag{10}
\end{equation}
To both sides we add the sum \(A\sum_{h=\nu+1}^{n}a_h^{(k)}x_h\), and obtain
\begin{equation}
 A\sum_{\mu=1}^{n}a_\mu^{(k)}x_\mu
 =\sum_{h=\nu+1}^{n}x_h\left(Aa_h^{(k)}-
 \sum_{i=1}^{\nu}\sum_{\mu=1}^{\nu}a_h^{(i)}a_\mu^{(k)}A_\mu^{(i)}\right).
\tag{11}
\end{equation}
The factor of \(x_h\) in the sum on the right-hand side is, by \S.\,23, (12), the determinant
\begin{equation}
\begin{vmatrix}
 a_1^{(1)}&a_2^{(1)}&\cdots&a_\nu^{(1)}&a_h^{(1)}\\
 a_1^{(2)}&a_2^{(2)}&\cdots&a_\nu^{(2)}&a_h^{(2)}\\
 \vdots&\vdots&&\vdots&\vdots\\
 a_1^{(\nu)}&a_2^{(\nu)}&\cdots&a_\nu^{(\nu)}&a_h^{(\nu)}\\
 a_1^{(k)}&a_2^{(k)}&\cdots&a_\nu^{(k)}&a_h^{(k)}
\end{vmatrix},
\tag{12}
\end{equation}
and therefore vanishes, if \(k\le\nu\), by \S.\,21, V, because
%CONT>>


% ===== scan leaf 111, printed 86 =====
% !PRINTED 86
\weberpage{86}
%<<CONT

two rows coincide, while if \(k>\nu\), it vanishes by the hypothesis, since (12) is then a \((\nu+1)\)-rowed determinant of the matrix (3). Thus from (11), since \(A\) is different from zero, we obtain
\begin{equation}
        \sum_{\mu=1}^{n}a_\mu^{(k)}x_\mu=0,\qquad k=1,2,\ldots,m,
\tag{13}
\end{equation}
that is, the system of equations (1) is satisfied by (7).

We shall single out from the theorem just proved the special case in which \(m=n-1\) and \(\nu=n-1\). In this case only \gesperrt{one} of the unknowns remains arbitrary, and the ratios of the unknowns are completely determined. We may state this result as follows:

Denote the \((n-1)\)-rowed determinants of the matrix
\begin{equation}
\begin{array}{cccc}
 a_1^{(1)}&a_2^{(1)}&\cdots&a_n^{(1)}\\
 a_1^{(2)}&a_2^{(2)}&\cdots&a_n^{(2)}\\
 \vdots&\vdots&&\vdots\\
 a_1^{(n-1)}&a_2^{(n-1)}&\cdots&a_n^{(n-1)}
\end{array}
\tag{14}
\end{equation}
taken with alternating sign by \(A_1,A_2,\ldots,A_n\), and assume that at least one of these quantities is different from zero. Then the solution of the system
\begin{equation}
\begin{aligned}
 a_1^{(1)}x_1+a_2^{(1)}x_2+\cdots+a_n^{(1)}x_n&=0,\\
 a_1^{(2)}x_1+a_2^{(2)}x_2+\cdots+a_n^{(2)}x_n&=0,\\
 &\vdots\\
 a_1^{(n-1)}x_1+a_2^{(n-1)}x_2+\cdots+a_n^{(n-1)}x_n&=0
\end{aligned}
\tag{15}
\end{equation}
is given by the ratios
\begin{equation}
        x_1:x_2:\cdots:x_n=A_1:A_2:\cdots:A_n.
\tag{16}
\end{equation}
Thus for \(n=3\) we obtain the solution of the system of equations, often occurring in geometry,
\begin{equation}
        ax+by+cz=0,\qquad a'x+b'y+c'z=0
\tag{17}
\end{equation}
in the form
\begin{equation}
        x:y:z=bc'-cb':ca'-ac':ab'-ba'.
\tag{18}
\end{equation}


% ===== scan leaf 112, printed 87 =====
% !PRINTED 87
\weberpage{87}
\pdfbookmark[2]{§ 25. Elimination from linear equations}{bm29}

\begin{center}
{\large \S.\ 25.}

\vspace{0.5em}
{\large\gesperrt{Elimination from Linear Equations.}}
\end{center}

It sometimes happens that, for a given system of linear equations, the question concerns not so much the actual determination of the unknowns as the decision whether their solution is possible, that is, the setting up of the condition equations which must hold among the coefficients if solutions, or solutions of a specified kind, are to exist at all. The setting up of these condition equations is called \gesperrt{elimination}. The solution of this problem is already contained implicitly in what precedes; but we shall return explicitly to some questions belonging here.

As in the preceding paragraph, we consider a system of \(m\) linear equations with \(n\) homogeneously occurring unknowns, and ask: when does this system have a solution in which not all the unknowns vanish? We have already seen that this is always the case if \(n>m\). But if \(n\le m\), the \gesperrt{necessary and sufficient} condition for such a solution is that all the \(n\)-rowed determinants of the matrix vanish. For if one of them does not vanish, then by \S.\,24, II the values of the unknowns are necessarily zero; whereas, if they all vanish, a number \(\nu<n\) can be found such that all \((\nu+1)\)-rowed determinants of the matrix are zero while at least one of the \(\nu\)-rowed determinants does not vanish, so that by \S.\,24, III a solution of the required kind exists.

From the matrix one can form, when \(n\le m\),
\[
        {m(m-1)\cdots(m-n+1)\over 1\cdot2\cdots n}
\]
\(n\)-rowed determinants, and this, therefore, would be the number of conditions. If \(n=m\), this number is \(=1\), and we obtain the case \S.\,24, II and, as was to be expected, one condition. In general, however, although all these conditions must be fulfilled, their number is greater than necessary, because some of them are necessary consequences of the others.


% ===== scan leaf 113, printed 88 =====
% !PRINTED 88
\weberpage{88}

To obtain a system of necessary, sufficient, and mutually independent conditions, we state the question somewhat more precisely and ask for the conditions:

\gesperrt{that from a system of \(m\) linear homogeneous equations with \(n\) unknowns, \(\nu\) of the unknowns can be completely determined by \(n-\nu\) which remain arbitrary.}

This question too is in fact already answered in \S.\,24. Among the \(\nu\)-rowed determinants one must be different from zero, while all the \((\nu+1)\)-rowed determinants vanish. But it is already sufficient if it is known of a smaller number of the \((\nu+1)\)-rowed determinants that they vanish.

Assume that the unknowns \(x_{\nu+1},x_{\nu+2},\ldots,x_n\) are to remain arbitrary, that \(x_1,x_2,\ldots,x_\nu\) are to be determined by them, and take the determinant
\begin{equation}
        A=\begin{vmatrix}
        a_1^{(1)}&a_2^{(1)}&\cdots&a_\nu^{(1)}\\
        a_1^{(2)}&a_2^{(2)}&\cdots&a_\nu^{(2)}\\
        \vdots&\vdots&&\vdots\\
        a_1^{(\nu)}&a_2^{(\nu)}&\cdots&a_\nu^{(\nu)}
        \end{vmatrix}
\tag{1}
\end{equation}
to be different from zero. We calculate the unknowns \(x_1,x_2,\ldots,x_\nu\) according to \S.\,24, (9), and form the sums (11), whose vanishing says that the given system of equations is actually satisfied. The conditions for this are therefore:
\begin{equation}
\begin{vmatrix}
        a_1^{(1)}&a_2^{(1)}&\cdots&a_\nu^{(1)}&a_h^{(1)}\\
        a_1^{(2)}&a_2^{(2)}&\cdots&a_\nu^{(2)}&a_h^{(2)}\\
        \vdots&\vdots&&\vdots&\vdots\\
        a_1^{(\nu)}&a_2^{(\nu)}&\cdots&a_\nu^{(\nu)}&a_h^{(\nu)}\\
        a_1^{(k)}&a_2^{(k)}&\cdots&a_\nu^{(k)}&a_h^{(k)}
\end{vmatrix}=0,
\tag{2}
\end{equation}
or, written otherwise,
\begin{equation}
        A a_h^{(k)}-\sum_{i=1}^{\nu}\sum_{\mu=1}^{\nu}
        a_h^{(i)}a_\mu^{(k)}A_\mu^{(i)}=0,
\tag{3}
\end{equation}
and these conditions are also sufficient. Equation (2) or (3), however, is identically satisfied when \(h=1,2,\ldots,\nu\), or \(k=1,2,\ldots,\nu\), and therefore gives no condition on the coefficients for these values. Such conditions arise only for
%CONT>>


% ===== scan leaf 114, printed 89 =====
% !PRINTED 89
\weberpage{89}
\pdfbookmark[2]{§ 26. Inhomogeneous linear equations}{bm30}
%<<CONT

\begin{equation}
 h=\nu+1,\nu+2,\ldots,n,
 \qquad
 k=\nu+1,\nu+2,\ldots,m,
\tag{4}
\end{equation}
and hence
\begin{equation}
        (n-\nu)(m-\nu)
\tag{5}
\end{equation}
in number.

These conditions are, however, actually mutually independent, that is, none follows from the others; for the left-hand sides of (3), by suitable assumptions concerning the coefficients \(a\), can receive an entirely arbitrary value for every combination of indices in the series (4), as is seen by setting all \(a_h^{(i)},a_\mu^{(k)}\), with the exception of \(a_h^{(k)}\), equal to zero.

\begin{center}
{\large \S.\ 26.}

\vspace{0.5em}
{\large\gesperrt{Non-homogeneous Linear Equations.}}
\end{center}

We have simplified the problem of solving linear equations in the preceding considerations, and reduced it to more general laws, to no inconsiderable degree by assuming the equations homogeneous with respect to the unknowns. In applications, however, the unknowns frequently do not occur homogeneously, and although in principle one case is not essentially different from the other, we shall nevertheless consider the case of non-homogeneous equations separately.

We can derive it from the case of homogeneous equations by adding the requirement that one specified unknown shall have the value \(1\). If this unknown occurs among those which the problem leaves arbitrary, no difficulty arises, because we need only put it \(=1\). But if it belongs to those determined by the others, certain further conditions must be satisfied, which say that the value \(1\) is admissible for this unknown.

We shall treat the question here independently and with a somewhat altered notation, but for simplicity restrict ourselves to the most important case, in which the number of unknowns agrees with the number of equations.

Let the following system of equations be solved with respect to the unknowns \(x_1,x_2,\ldots,x_n\):
%CONT>>


% ===== scan leaf 115, printed 90 =====
% !PRINTED 90
\weberpage{90}
%<<CONT

\begin{equation}
\begin{aligned}
 a_1^{(1)}x_1+a_2^{(1)}x_2+\cdots+a_n^{(1)}x_n&=y_1,\\
 a_1^{(2)}x_1+a_2^{(2)}x_2+\cdots+a_n^{(2)}x_n&=y_2,\\
 &\vdots\\
 a_1^{(n)}x_1+a_2^{(n)}x_2+\cdots+a_n^{(n)}x_n&=y_n.
\end{aligned}
\tag{1}
\end{equation}
where the coefficients \(a_i^{(k)}\) and the independent terms \(y_1,y_2,\ldots,y_n\) are regarded as given. We denote by
\begin{equation}
        A=\sum\pm a_1^{(1)}a_2^{(2)}\cdots a_n^{(n)}
\tag{2}
\end{equation}
the determinant of the system, and, as before, denote by \(A_i^{(k)}\) the first subdeterminants of \(A\). If \(k\) is one of the indices \(1,2,\ldots,n\), and we multiply equations (1) in succession by \(A_k^{(1)},A_k^{(2)},\ldots,A_k^{(n)}\) and then add them, we obtain
\begin{equation}
 x_1\sum^{i}a_1^{(i)}A_k^{(i)}+
 x_2\sum^{i}a_2^{(i)}A_k^{(i)}+
 \cdots+
 x_n\sum^{i}a_n^{(i)}A_k^{(i)}
 =\sum^{i}y_iA_k^{(i)}.
\tag{3}
\end{equation}
By \S.\,22, (2) and (6), the coefficients of \(x_1,x_2,\ldots,x_n\) vanish here with the exception of the single coefficient of \(x_k\), which receives the value \(A\), and consequently:
\begin{equation}
 A x_k=\sum^{i}y_iA_k^{(i)},
\tag{4}
\end{equation}
and in the case in which \(A\) is different from zero, the values of the unknowns can be calculated uniquely from this in such a way that equations (1) are actually satisfied [\S.\,23, (3), (7)]. If we put \(A_k^{(i)}=A\alpha_k^{(i)}\), we can write the solution in the form
\begin{equation}
\begin{aligned}
 x_1&=\alpha_1^{(1)}y_1+\alpha_1^{(2)}y_2+\cdots+\alpha_1^{(n)}y_n,\\
 x_2&=\alpha_2^{(1)}y_1+\alpha_2^{(2)}y_2+\cdots+\alpha_2^{(n)}y_n,\\
 &\vdots\\
 x_n&=\alpha_n^{(1)}y_1+\alpha_n^{(2)}y_2+\cdots+\alpha_n^{(n)}y_n,
\end{aligned}
\tag{5}
\end{equation}
whose analogy with system (1) is evident.

But if \(A=0\), equations (4) teach nothing about the values of the \(x\)'s; they give only a system of conditions which the \(y\)'s must necessarily satisfy if equations (1) are to be soluble at all. To arrive at a generally valid result, we proceed quite similarly to the case of homogeneous equations. We assume that, besides the determinant \(A\), all subdeterminants down to a certain order are also zero. We denote a non-vanishing \(\nu\)-rowed subdeterminant by
%CONT>>


% ===== scan leaf 116, printed 91 =====
% !PRINTED 91
\weberpage{91}
%<<CONT

\begin{equation}
        B=\begin{vmatrix}
        a_1^{(1)}&a_2^{(1)}&\cdots&a_\nu^{(1)}\\
        a_1^{(2)}&a_2^{(2)}&\cdots&a_\nu^{(2)}\\
        \vdots&\vdots&&\vdots\\
        a_1^{(\nu)}&a_2^{(\nu)}&\cdots&a_\nu^{(\nu)}
        \end{vmatrix},
\tag{6}
\end{equation}
denote its subdeterminants by \(B_k^{(h)}\), where \(h,k\) run only through the values from \(1\) to \(\nu\), and now suppose that \gesperrt{all \((\nu+1)\)-rowed subdeterminants of \(A\) vanish}. If we multiply the first \(\nu\) equations of the system (1) by \(B_k^{(1)},B_k^{(2)},\ldots,B_k^{(\nu)}\), and add, it follows [\S.\,22, (2), (6)]:
\begin{equation}
 Bx_k=\sum_{i=1}^{\nu}B_k^{(i)}y_i-
 \sum_{s=\nu+1}^{n}x_s\sum_{i=1}^{\nu}B_k^{(i)}a_s^{(i)}.
\tag{7}
\end{equation}
Hereby \(\nu\) of the unknowns \(x\) are determined by the others, and, to see how far equations (1) are thereby satisfied, we substitute the expressions in (1). For this purpose we multiply (7) by \(a_k^{(\lambda)}\), where \(\lambda\) runs through the values \(1,2,\ldots,n\), and sum with respect to \(k\). Thus follows:
\begin{equation}
 B\sum_{k=1}^{\nu}a_k^{(\lambda)}x_k
 =\sum_{i=1}^{\nu}\sum_{k=1}^{\nu}B_k^{(i)}a_k^{(\lambda)}y_i
 -\sum_{s=\nu+1}^{n}x_s\sum_{i=1}^{\nu}\sum_{k=1}^{\nu}B_k^{(i)}a_s^{(i)}a_k^{(\lambda)}.
\tag{8}
\end{equation}
and if we add to both sides
\[
 B\sum_{k=\nu+1}^{n}a_k^{(\lambda)}x_k
\]
we obtain:
\begin{equation}
 B\sum_{k=1}^{n}a_k^{(\lambda)}x_k
 =\sum_{i=1}^{\nu}\sum_{k=1}^{\nu}B_k^{(i)}a_k^{(\lambda)}y_i
 +\sum_{s=\nu+1}^{n}x_s\left(Ba_s^{(\lambda)}-
 \sum_{i=1}^{\nu}\sum_{k=1}^{\nu}B_k^{(i)}a_s^{(i)}a_k^{(\lambda)}\right).
\tag{9}
\end{equation}
Now, by \S.\,23, (12), the coefficient of \(x_s\) on the right-hand side is a \((\nu+1)\)-rowed subdeterminant of \(A\), and hence is zero by the hypothesis. According to (1),
\[
 \sum_{k=1}^{n}a_k^{(\lambda)}x_k=y_\lambda
\]
is to hold, and consequently the expressions (7) satisfy equations (1) if and only if
\begin{equation}
 B y_\lambda=\sum_{i=1}^{\nu}\sum_{k=1}^{\nu}B_k^{(i)}a_k^{(\lambda)}y_i
\tag{10}
\end{equation}
holds. This equation is always satisfied when \(\lambda\le\nu\), by \S.\,22, (2), (6). But if \(\lambda=\nu+1,\nu+2,\ldots,n\), then (10) contains \(n-\nu\) conditions on the \(y\)'s which, since each
%CONT>>


% ===== scan leaf 117, printed 92 =====
% !PRINTED 92
\weberpage{92}
\pdfbookmark[2]{§ 27. Multiplication of determinants}{bm31}
%<<CONT

condition contains a new one of the quantities \(y\), are independent of one another. \gesperrt{If one of these conditions is not fulfilled, the given system of equations has no solution.}

If in the system of equations (1) we regard the \(x\)'s not as unknowns but as variables, then the \(y\)'s too will be variable quantities. Under this interpretation we call the system (1) a \gesperrt{linear substitution}, insofar as it mediates the passage from one system of variables to another. Only when the determinant \(A\), which we now call the \gesperrt{determinant of the substitution}, is different from zero can the \(y\)'s be regarded as independent variables. If this is the case, the system (5) gives the representation of the variables \(x\) in terms of the \(y\)'s, or the \gesperrt{substitution inverse} to (1).


\begin{center}
{\large \S.\ 27.}

\vspace{0.5em}
{\large\gesperrt{Multiplication of Determinants.}}
\end{center}

The theorem which we still have to prove teaches how one can represent the product of two determinants with the same number of rows by a single determinant with the same number of rows. One is led to it most simply by considering the solution of two systems of linear equations.

Let now the coefficients $a_i^{(k)}, b_i^{(k)}$, when $i,k$ run through the sequence of numbers $1,2,\ldots,n$, be arbitrary \gesperrt{variable quantities}; likewise let the quantities $x_i,y_i,z_i$ be bound only by the relations
\begin{align}
 \sum_i a_i^{(k)}x_i&=y_k,\tag{1}\\
 \sum_k b_h^{(k)}y_k&=z_h.\tag{2}
\end{align}
If the problem is now posed of determining the variables $x$ by the variables $z$, this problem can be solved in two ways. By § 26, (4) one may solve equations (1) with respect to $x$, equations (2) with respect to $y$, and insert the latter expressions into the former. If $A,B$ denote the determinants of the two systems of equations, that is,
%CONT>>


% ===== scan leaf 118, printed 93 =====
% !PRINTED 93
\weberpage{93}
%<<CONT

\[
A=\begin{vmatrix}
 a_1^{(1)}&a_2^{(1)}&\cdots&a_n^{(1)}\\
 a_1^{(2)}&a_2^{(2)}&\cdots&a_n^{(2)}\\
 \cdots&\cdots&&\cdots\\
 a_1^{(n)}&a_2^{(n)}&\cdots&a_n^{(n)}
\end{vmatrix},\qquad
B=\begin{vmatrix}
 b_1^{(1)}&b_2^{(1)}&\cdots&b_n^{(1)}\\
 b_1^{(2)}&b_2^{(2)}&\cdots&b_n^{(2)}\\
 \cdots&\cdots&&\cdots\\
 b_1^{(n)}&b_2^{(n)}&\cdots&b_n^{(n)}
\end{vmatrix},
\]
and if $A_i^{(k)},B_i^{(k)}$ denote the first subdeterminants, one obtains
\begin{align}
 A x_k&=\sum_i A_k^{(i)}y_i,\tag{3}\\
 B y_i&=\sum_h B_h^{(i)}z_h,\tag{4}
\end{align}
and, by substitution of (4) in (3),
\begin{equation}
 ABx_k=\sum_h z_h\sum_i A_k^{(i)}B_h^{(i)}.\tag{5}
\end{equation}
But one may also proceed by inserting from (1) the expressions for $y$ into (2), whereby, if
\begin{equation}
 c_i^{(h)}=\sum_s a_i^{(s)}b_h^{(s)}\tag{6}
\end{equation}
is put, one obtains
\begin{equation}
 \sum_i c_i^{(h)}x_i=z_h.\tag{7}
\end{equation}
Now put
\[
C=\begin{vmatrix}
 c_1^{(1)}&c_2^{(1)}&\cdots&c_n^{(1)}\\
 c_1^{(2)}&c_2^{(2)}&\cdots&c_n^{(2)}\\
 \cdots&\cdots&&\cdots\\
 c_1^{(n)}&c_2^{(n)}&\cdots&c_n^{(n)}
\end{vmatrix}\tag{8}
\]
If $C_i^{(k)}$ denotes the subdeterminants of $C$, the solution of (7) gives
\begin{equation}
 Cx_k=\sum_h z_h C_k^{(h)},\tag{9}
\end{equation}
and comparison of (5) with (9) gives
\begin{equation}
 AB=C,\tag{10}
\end{equation}
and for the subdeterminants
\begin{equation}
 C_k^{(h)}=\sum_i A_k^{(i)}B_h^{(i)}.\tag{11}
\end{equation}
In this conclusion, however, there is still a gap, caused by the doubt whether (5) and (9) might not differ by a factor common to both sides. To remove this doubt we shall prove equation (10) directly; the correctness of (11) then follows.


% ===== scan leaf 119, printed 94 =====
% !PRINTED 94
\weberpage{94}

We shall form the principal term of the determinant $C$ from (6). Let $s_1,s_2,\ldots,s_n$ be mutually independent summation letters which run from $1$ to $n$. Then
\begin{align}
 c_1^{(1)}c_2^{(2)}\cdots c_n^{(n)}
 &=\sum_{s_1}a_1^{(s_1)}b_1^{(s_1)}
   \sum_{s_2}a_2^{(s_2)}b_2^{(s_2)}\cdots
   \sum_{s_n}a_n^{(s_n)}b_n^{(s_n)}\notag\\
 &=\sum_{s_1,s_2,\ldots,s_n}a_1^{(s_1)}a_2^{(s_2)}\cdots a_n^{(s_n)}
 b_1^{(s_1)}b_2^{(s_2)}\cdots b_n^{(s_n)}.\tag{12}
\end{align}
The permutation of the lower indices of the $c$ corresponds to the permutation of the lower indices of the $a$, and if, with regard to the rule of signs, one therefore forms the determinant $C$, one obtains
\begin{align}
 \sum \pm c_1^{(1)}c_2^{(2)}\cdots c_n^{(n)}
 =\sum_{s_1,s_2,\ldots,s_n} b_1^{(s_1)}b_2^{(s_2)}\cdots b_n^{(s_n)}
 \sum\pm a_1^{(s_1)}a_2^{(s_2)}\cdots a_n^{(s_n)}.\tag{13}
\end{align}
Now
\[
\sum\pm a_1^{(s_1)}a_2^{(s_2)}\cdots a_n^{(s_n)}
\]
is, by § 21, VI, always zero if the same digit occurs twice among $s_1,s_2,
\ldots,s_n$; hence only those terms in (13) keep a non-zero value for which $s_1,s_2,
\ldots,s_n$ is an arrangement of the indices $1,2,
\ldots,n$. This value is $+A$ or $-A$ according as this arrangement belongs to the first or the second kind (§ 21, IV).

Therefore the right-hand side of (13) becomes
\[
 A\sum_{s_1,s_2,\ldots,s_n}\pm b_1^{(s_1)}b_2^{(s_2)}\cdots b_n^{(s_n)},
\]
where the sum is extended over all permutations $s_1,s_2,
\ldots,s_n$. But this sum is precisely the determinant $B$, and hence the formula
\[
 C=AB
\]
is proved.

The law of formation of the elements $c_i^{(h)}$ contained in formula (6) can be expressed in words as follows: \gesperrt{To obtain the elements of the determinant $C$, which is the product of the two determinants $A,B$, one multiplies the elements of each column of $A$ by the corresponding elements of a column of $B$ and adds the products.}

By the theorems on determinants the form of $C$ can be varied in many ways. We record the following.


% ===== scan leaf 120, printed 95 =====
% !PRINTED 95
\weberpage{95}
\pdfbookmark[2]{§ 28. Determinants of the minors}{bm32}

If two columns in $A$ or in $B$ are interchanged, the elements $c_i^{(h)}$ do not change, but only interchange among themselves. If, however, two rows in $A$ or in $B$ are interchanged, then the $c_i^{(h)}$ change because the factors in the individual products of the sum are grouped differently; but if \gesperrt{corresponding} interchanges are made simultaneously in the rows of $A$ and of $B$, the $c_i^{(h)}$ remain unchanged, since only the individual terms of the sum are thereby interchanged.

By turning the rows into columns in $A$ or in $B$, or in both at once, one obtains three further different ways of forming the product of two determinants in determinant form. This can also be expressed as follows: \gesperrt{To form the product of two determinants, one may multiply the elements of the separate rows or columns of one factor by the corresponding elements of the rows or columns of the other factor, add the products, and regard these sums of products as the elements of a new determinant.}

The multiplication rule can be applied to a product of two determinants with different numbers of elements by transforming the determinant with the smaller number of rows, by the theorem § 22, IX, into another with more rows.

\begin{center}
{\large \S.\ 28.}

\vspace{0.5em}
{\large\gesperrt{Determinants of Subdeterminants.}}
\end{center}

We immediately make an application here of the multiplication law of determinants. Let, as before,
\[
A=\begin{vmatrix}
 a_1^{(1)}&a_2^{(1)}&\cdots&a_n^{(1)}\\
 a_1^{(2)}&a_2^{(2)}&\cdots&a_n^{(2)}\\
 \cdots&\cdots&&\cdots\\
 a_1^{(n)}&a_2^{(n)}&\cdots&a_n^{(n)}
\end{vmatrix},\tag{1}
\]
and let
%CONT>>


% ===== scan leaf 121, printed 96 =====
% !PRINTED 96
\weberpage{96}
%<<CONT

\[
\begin{matrix}
 A_1^{(1)}&A_2^{(1)}&\cdots&A_n^{(1)}\\
 A_1^{(2)}&A_2^{(2)}&\cdots&A_n^{(2)}\\
 \cdots&\cdots&&\cdots\\
 A_1^{(n)}&A_2^{(n)}&\cdots&A_n^{(n)}
\end{matrix}\tag{2}
\]
be the system of subdeterminants. If we form from (2) the determinant which we shall denote by $\Delta$, we may apply the multiplication rule to the product $A\Delta$. By § 22, (3) and (7), this gives
\[
 A\Delta=\begin{vmatrix}
 A&0&\cdots&0\\
 0&A&\cdots&0\\
 \cdots&\cdots&&\cdots\\
 0&0&\cdots&A
\end{vmatrix}=A^n,
\]
and, by division by $A$,
\begin{equation}
 \Delta=A^{n-1}.\tag{3}
\end{equation}
Thus $\Delta$ is the $(n-1)$st power of $A$. In this derivation it has indeed first been assumed that $A$ is different from zero. But since (3) is an identity with respect to the elements $a_i^{(h)}$, that is, since it also holds when these quantities are independent variables, it follows that the formula (3) still holds in the exceptional case; that is, if $A$ vanishes, then $\Delta$ also vanishes.

This result is a special case of a more general theorem by which any determinant of the matrix (2) may be formed. Consider the $\nu$-rowed subdeterminant
\[
\Delta_\nu=\begin{vmatrix}
 A_1^{(1)}&A_2^{(1)}&\cdots&A_\nu^{(1)}\\
 A_1^{(2)}&A_2^{(2)}&\cdots&A_\nu^{(2)}\\
 \cdots&\cdots&&\cdots\\
 A_1^{(\nu)}&A_2^{(\nu)}&\cdots&A_\nu^{(\nu)}
\end{vmatrix},\tag{4}
\]
from which all other $\nu$-rowed subdeterminants can be derived by permutation of the upper and lower indices. One may apply the multiplication rule by transforming $\Delta_\nu$, according to the concluding remark of the last paragraph, into an $n$-rowed determinant.


% ===== scan leaf 122, printed 97 =====
% !PRINTED 97
\weberpage{97}

Put
\[
\Delta_\nu=\begin{vmatrix}
 A_1^{(1)}&\cdots&A_\nu^{(1)}&A_{\nu+1}^{(1)}&\cdots&A_n^{(1)}\\
 \cdots&&\cdots&\cdots&&\cdots\\
 A_1^{(\nu)}&\cdots&A_\nu^{(\nu)}&A_{\nu+1}^{(\nu)}&\cdots&A_n^{(\nu)}\\
 0&\cdots&0&1&\cdots&0\\
 \cdots&&\cdots&\cdots&&\cdots\\
 0&\cdots&0&0&\cdots&1
\end{vmatrix},\tag{5}
\]
where $n-\nu$ rows and columns have been adjoined, the former containing only zeros except in the diagonal terms. If one now forms the product $A\Delta_\nu$, it follows that
\begin{equation}
A\Delta_\nu=
\begin{vmatrix}
 A&0&\cdots&0&a_{\nu+1}^{(1)}&\cdots&a_n^{(1)}\\
 0&A&\cdots&0&a_{\nu+1}^{(2)}&\cdots&a_n^{(2)}\\
 \cdots&\cdots&&\cdots&\cdots&&\cdots\\
 0&0&\cdots&A&a_{\nu+1}^{(\nu)}&\cdots&a_n^{(\nu)}\\
 0&0&\cdots&0&a_{\nu+1}^{(\nu+1)}&\cdots&a_n^{(\nu+1)}\\
 \cdots&\cdots&&\cdots&\cdots&&\cdots\\
 0&0&\cdots&0&a_{\nu+1}^{(n)}&\cdots&a_n^{(n)}
\end{vmatrix},\tag{6}
\end{equation}
and by Theorem IX, § 22, this is
\[
= A^\nu
 \begin{vmatrix}
 a_{\nu+1}^{(\nu+1)}&\cdots&a_n^{(\nu+1)}\\
 \cdots&&\cdots\\
 a_{\nu+1}^{(n)}&\cdots&a_n^{(n)}
 \end{vmatrix}.
\]
Dividing here by $A$ and applying the notation of § 23, one obtains
\begin{equation}
 \Delta_\nu=A^{\nu-1}A_{1,2,\ldots,\nu}^{1,2,\ldots,\nu}.\tag{7}
\end{equation}
For $\nu=2$ this gives the special result
\begin{equation}
 A_1^{(1)}A_2^{(2)}-A_1^{(2)}A_2^{(1)}=A A_{1,2}^{1,2}.\tag{8}
\end{equation}
We shall use formula (8) often later. In the notation by differential quotients it can be represented as
\begin{equation}
 A\frac{\partial^2A}{\partial a_1^{(1)}\partial a_2^{(2)}}
 =\frac{\partial A}{\partial a_1^{(1)}}\frac{\partial A}{\partial a_2^{(2)}}
 -\frac{\partial A}{\partial a_1^{(2)}}\frac{\partial A}{\partial a_2^{(1)}}.\tag{9}
\end{equation}
It is especially important in the case where $A$ is a symmetric determinant, so that $a_i^{(k)}=a_k^{(i)}$; then also
\[
 \frac{\partial A}{\partial a_1^{(2)}}=\frac{\partial A}{\partial a_2^{(1)}},
\]
%CONT>>


% ===== scan leaf 123, printed 98 =====
% !PRINTED 98
\weberpage{98}
\pdfbookmark[2]{§ 29. Interpolation}{bm33}
%<<CONT

where in the differentiation no account is taken of the dependence $a_i^{(k)}=a_k^{(i)}$. Then formula (9) becomes
\begin{equation}
 A\frac{\partial^2A}{\partial a_1^{(1)}\partial a_2^{(2)}}
 =\frac{\partial A}{\partial a_1^{(1)}}\frac{\partial A}{\partial a_2^{(2)}}
 -\left(\frac{\partial A}{\partial a_1^{(2)}}\right)^2.\tag{10}
\end{equation}

\begin{center}
{\large \S.\ 29.}

\vspace{0.5em}
{\large\gesperrt{Interpolation.}}
\end{center}

In § 9, under a special assumption, we already solved the problem of determining an integral rational function which assumes prescribed values for a sufficient number of prescribed values of the argument. We shall now solve this problem generally by applying determinants. First, however, we shall put forward the following observations.

In § 22, (12), the value of the determinant
\[
\begin{vmatrix}
 \alpha_1^{n-1}&\alpha_1^{n-2}&\cdots&\alpha_1&1\\
 \alpha_2^{n-1}&\alpha_2^{n-2}&\cdots&\alpha_2&1\\
 \cdots&\cdots&&\cdots&\cdots\\
 \alpha_n^{n-1}&\alpha_n^{n-2}&\cdots&\alpha_n&1
\end{vmatrix}
\]
was found to be equal to the product of differences
\begin{equation}
(\alpha_1-\alpha_2)(\alpha_1-\alpha_3)\cdots(\alpha_1-\alpha_n)
(\alpha_2-\alpha_3)\cdots(\alpha_2-\alpha_n)\cdots(\alpha_{n-1}-\alpha_n).\tag{1}
\end{equation}
For brevity we now denote this product of differences by
\begin{equation}
 [\alpha_1,\alpha_2,\alpha_3,\ldots,\alpha_n],\tag{2}
\end{equation}
so that the quantity $[\alpha_1,\alpha_2,\alpha_3,\ldots,\alpha_n]$ changes its sign when two of the elements $\alpha_1,\alpha_2,
\ldots,\alpha_n$ are interchanged.

We further define an integral rational function $f(x)$ of degree $n$ in the variable $x$ by the equation
\begin{equation}
 f(x)=(x-\alpha_1)(x-\alpha_2)\cdots(x-\alpha_n),\tag{3}
\end{equation}
so that, by § 13,
\begin{align}
 f'(\alpha_1)&=(\alpha_1-\alpha_2)(\alpha_1-\alpha_3)\cdots(\alpha_1-\alpha_n),\notag\\
 f'(\alpha_2)&=(\alpha_2-\alpha_1)(\alpha_2-\alpha_3)\cdots(\alpha_2-\alpha_n),\tag{4}\\
 &\vdots\notag\\
 f'(\alpha_n)&=(\alpha_n-\alpha_1)(\alpha_n-\alpha_2)\cdots(\alpha_n-\alpha_{n-1}).\notag
\end{align}
%CONT>>


% ===== scan leaf 124, printed 99 =====
% !PRINTED 99
\weberpage{99}
%<<CONT

In the product of all these expressions each factor of the product (1) occurs twice with opposite sign; since the number of factors of (1) is $\frac{n(n-1)}{2}$, we obtain the equation
\begin{equation}
 f'(\alpha_1)f'(\alpha_2)\cdots f'(\alpha_n)=(-1)^{\frac{n(n-1)}2}[\alpha_1,\alpha_2,\ldots,\alpha_n]^2.\tag{5}
\end{equation}
But with the aid of the relations (4) we may also represent our product (1) as follows:
\begin{align}
 [\alpha_1,\alpha_2,\ldots,\alpha_n]
 &= f'(\alpha_1)[\alpha_2,\alpha_3,\ldots,\alpha_n]\notag\\
 &=-f'(\alpha_2)[\alpha_1,\alpha_3,\ldots,\alpha_n]\tag{6}\\
 &=+f'(\alpha_3)[\alpha_1,\alpha_2,\ldots,\alpha_n]\notag\\% [sic] Weber prints the full alpha_1,alpha_2,...,alpha_n sequence here, although the factor f'(alpha_3) requires the bracket with alpha_3 omitted and the immediately following sentence explicitly says each right-hand bracket contains one element fewer. Source-original mathematical error preserved rather than silently normalized.
 &=\cdots=(-1)^{n-1}f'(\alpha_n)[\alpha_1,\alpha_2,\ldots,\alpha_{n-1}],\notag
\end{align}
where on the right-hand side each bracketed expression contains one element fewer than the bracketed expression on the left.

Now pose the problem of determining an integral rational function of degree $(n-1)$, $\varphi(x)$, which, for the argument values $x=\alpha_1,\alpha_2,
\ldots,\alpha_n$, assumes in order the prescribed values $\varphi(\alpha_1),\varphi(\alpha_2),\ldots,\varphi(\alpha_n)$.

If one puts
\begin{equation}
 \varphi(x)=a_1x^{n-1}+a_2x^{n-2}+\cdots+a_{n-1}x+a_n,\tag{7}
\end{equation}
then the $n$ unknown coefficients $a_1,a_2,
\ldots,a_n$ are to be determined from the linear equations
\begin{align}
 \varphi(\alpha_1)&=a_1\alpha_1^{n-1}+a_2\alpha_1^{n-2}+\cdots+a_{n-1}\alpha_1+a_n,\notag\\
 \varphi(\alpha_2)&=a_1\alpha_2^{n-1}+a_2\alpha_2^{n-2}+\cdots+a_{n-1}\alpha_2+a_n,\tag{8}\\
 &\vdots\notag\\
 \varphi(\alpha_n)&=a_1\alpha_n^{n-1}+a_2\alpha_n^{n-2}+\cdots+a_{n-1}\alpha_n+a_n.\notag
\end{align}
The determinant of this system is precisely the quantity (2)
\[
[\alpha_1,\alpha_2,\ldots,\alpha_n],
\]
and the problem is therefore always soluble by § 26 when this quantity is different from zero, that is, when no two of the $n$ quantities $\alpha_1,\alpha_2,
\ldots,\alpha_n$ are equal.

Instead of solving the equations (8) by § 26 and substituting the expressions so found into (7), it is preferable, by § 24, II, to eliminate the homogeneous quantities
\[
1,a_1,a_2,\ldots,a_n
\]
%CONT>>


% ===== scan leaf 125, printed 100 =====
% !PRINTED 100
\weberpage{100}
%<<CONT

from the $n+1$ equations (7) and (8). This gives the determinant equation
\[
\begin{vmatrix}
 \varphi(x)&x^{n-1}&x^{n-2}&\cdots&x&1\\
 \varphi(\alpha_1)&\alpha_1^{n-1}&\alpha_1^{n-2}&\cdots&\alpha_1&1\\
 \varphi(\alpha_2)&\alpha_2^{n-1}&\alpha_2^{n-2}&\cdots&\alpha_2&1\\
 \cdots&\cdots&\cdots&&\cdots&\cdots\\
 \varphi(\alpha_n)&\alpha_n^{n-1}&\alpha_n^{n-1}&\cdots&\alpha_n&1% [sic] Weber prints the third element of the last row as \alpha_n^{n-1}, not \alpha_n^{n-2}; source-original error preserved to match the scan.
\end{vmatrix}=0.\tag{9}
\]
We now expand this determinant according to the elements of the first column. For the subdeterminants which occur there we can use the notation (2), and we obtain
\begin{align}
&\varphi(x)[\alpha_1,\alpha_2,\ldots,\alpha_n]
-\varphi(\alpha_1)[x,\alpha_2,\alpha_3,\ldots,\alpha_n]
+\varphi(\alpha_2)[x,\alpha_1,\alpha_3,\ldots,\alpha_n]
\notag\\
&\qquad+\cdots(-1)^n\varphi(\alpha_n)[x,\alpha_1,\alpha_2,\ldots,\alpha_{n-1}]=0.\tag{10}
\end{align}
Now by (1)
\[
 [x,\alpha_2,\alpha_3,
 \ldots,\alpha_n]=(x-\alpha_2)(x-\alpha_3)\cdots(x-\alpha_n)[\alpha_2,\alpha_3,
 \ldots,\alpha_n],
\]
for which we may also write
\[
 [x,\alpha_2,\alpha_3,\ldots,\alpha_n]=\frac{f(x)}{x-\alpha_1}[\alpha_2,\alpha_3,\ldots,\alpha_n],
\]
and likewise
\[
 [x,\alpha_1,\alpha_3,\ldots,\alpha_n]=\frac{f(x)}{x-\alpha_2}[\alpha_1,\alpha_3,\ldots,\alpha_n],
\]
\[\cdots\]
\[
 [x,\alpha_1,\alpha_2,\ldots,\alpha_{n-1}]=\frac{f(x)}{x-\alpha_n}[\alpha_1,\alpha_2,\ldots,\alpha_{n-1}],
\]
Dividing (10) by $[\alpha_1,\alpha_2,
\ldots,\alpha_n]$ and taking account of the several expressions (6), one finally obtains
\begin{equation}
 \varphi(x)=f(x)\left(\frac{\varphi(\alpha_1)}{(x-\alpha_1)f'(\alpha_1)}+\frac{\varphi(\alpha_2)}{(x-\alpha_2)f'(\alpha_2)}+\cdots+\frac{\varphi(\alpha_n)}{(x-\alpha_n)f'(\alpha_n)}\right).\tag{11}
\end{equation}
This completely determines the function $\varphi(x)$. This expression for $\varphi(x)$ is called the \gesperrt{interpolation formula of Lagrange}. That it satisfies the stipulated requirements can be verified afterwards very easily by putting $x_1=\alpha_1,\alpha_2,
\ldots,\alpha_n$.% [sic] Weber prints x_1 rather than x in this verification clause; source-original error preserved to match the scan.
\ The denominators $x-\alpha_1,x-\alpha_2,
\ldots,x-\alpha_n$ are present in it only apparently.


% ===== scan leaf 126, printed 101 =====
% !PRINTED 101
\weberpage{101}
\pdfbookmark[1]{Third Section. The roots of algebraic equations.}{bm34}
\pdfbookmark[2]{§ 30. Concept of the roots. Multiple roots}{bm35}

\begin{center}
{\Large\gesperrt{Third Section.}}

\vspace{1em}
{\large\gesperrt{The Roots of Algebraic Equations.}}
\end{center}

\begin{center}
{\large \S.\ 30.}

\vspace{0.5em}
{\large\gesperrt{Concept of Roots. Multiple Roots.}}
\end{center}

After, in the first two sections, the algebraic magnitudes had been considered more from the formal side, where the matter was identical transformations of letter-expressions in which the letters throughout could be conceived as symbols for variable magnitudes, the numerical magnitudes now come more into the foreground.

Here, in accordance with what was fixed in the Introduction, by numbers we understand real or imaginary magnitudes of the form $a+bi$, and for illustration we represent the real magnitudes by points of a straight line and the imaginary magnitudes by points of a plane. By the \gesperrt{absolute value} of an imaginary magnitude $a+bi$ we understand $\sqrt{a^2+b^2}$ and denote it, following Weierstrass, by
\[
 |a+bi|.
\]
Let now
\begin{equation}
 f(x)=a_0x^n+a_1x^{n-1}+a_2x^{n-2}+\cdots+a_n
\tag{1}
\end{equation}
be an integral function of $x$, in which the coefficients are arbitrary real or imaginary numbers, and the first, $a_0$, is assumed different from zero. If $\alpha$ is a number which, when put for $x$, makes the function $f(x)$ equal to zero, that is, which satisfies the condition $f(\alpha)=0$, then $\alpha$ is called a root of the equation
\[
 f(x)=0.
\]


% ===== scan leaf 127, printed 102 =====
% !PRINTED 102
\weberpage{102}

We also say briefly that $\alpha$ is a root of $f(x)$. By § 4, $f(x)$ may then be divided by $x-\alpha$ without remainder, so that one may put
\begin{equation}
 f(x)=(x-\alpha)f_1(x),
\tag{2}
\end{equation}
where $f_1(x)$ is only of the $(n-1)$st degree and has the same first coefficient $a_0$ as $f(x)$, thus
\[
 f_1(x)=a_0x^{n-1}+a'_1x^{n-2}+\cdots .
\]
The coefficients of $f_1(x)$ are given in § 4, (6).

Every root of $f_1(x)$ is therefore at the same time a root of $f(x)$; and conversely every root of $f(x)$ is either $=\alpha$ or is a root of $f_1(x)$. If one root $\alpha$ of $f(x)$ is known, then the task of finding the remaining ones is reduced to the solution of an equation of the $(n-1)$st degree.

The aim of the considerations of this section consists in proving that every function $f(x)$ of the $n$th degree has at least one root. This theorem is called the \gesperrt{fundamental theorem of algebra}. First we draw from (2) the important consequence:

\medskip
\noindent I. \gesperrt{An equation of the $n$th degree cannot have more than $n$ roots.}
\medskip

For if $f(x)$ had more than $n$ roots, then $f_1(x)$, which is only of the $(n-1)$st degree, would have more than $n-1$ roots. But an equation of the first degree has not more than one root, whence the correctness of our theorem follows by complete induction.

Sometimes one also gives it the expression:

\medskip
\noindent II. \gesperrt{If a function $f(x)$ of the $n$th degree has more than $n$ roots, then all its coefficients must be zero.}
\medskip

If $\beta$ is a root of $f_1(x)$, then one may likewise put
\[
 f_1(x)=(x-\beta)f_2(x),\quad f(x)=(x-\alpha)(x-\beta)f_2(x),
\]
where $f_2(x)=a_0x^{n-2}+\cdots$ is only of the $(n-2)$nd degree. Thus, if one assumes that each of the functions obtained by this division, $f_1(x), f_2(x),\ldots$, has at least one root, one finally obtains
\begin{equation}
 f(x)=a_0(x-\alpha)(x-\beta)\cdots(x-\nu),
\tag{3}
\end{equation}
and we can therefore also express the theorem which we earlier described as the aim of our considerations as follows. It is to be proved:

\medskip


% ===== scan leaf 128, printed 103 =====
% !PRINTED 103
\weberpage{103}
\pdfbookmark[2]{§ 31. Continuity of integral functions}{bm36}

\noindent A function of the $n$th degree can be decomposed into $n$ linear factors.
\medskip

The magnitudes $\alpha,\beta,\ldots,\nu$ are then all roots of $f(x)$, and their number is therefore $n$. It is, however, not excluded that the same number occurs more than once among $\alpha,\beta,\ldots,\nu$. Then $f(x)$ would have fewer than $n$ roots, while the theorem would still remain that $f(x)$ is decomposable into $n$ linear factors. To establish agreement, it has been agreed that, if $x-\alpha$ goes into $f(x)$ several times, $\alpha$ is to be counted several times among the roots, and hence that one speaks of simple, double, triple, etc. roots.

By the concept of the derived functions [§ 13, (2)], if $\alpha$ is an arbitrary magnitude,
\begin{equation}
 f(x)=f(\alpha)+(x-\alpha)f'(\alpha)+\frac{(x-\alpha)^2}{1\cdot2}f''(\alpha)
 +\frac{(x-\alpha)^3}{1\cdot2\cdot3}f'''(\alpha)+\cdots .
\tag{4}
\end{equation}
If it is now again assumed that $f(\alpha)$ vanishes, then there results
\[
 f_1(x)=f'(\alpha)+\frac{x-\alpha}{1\cdot2}f''(\alpha)+\frac{(x-\alpha)^2}{1\cdot2\cdot3}f'''(\alpha)+\cdots,
\]
and hence
\[
 f_1(\alpha)=f'(\alpha).
\]
Thus $\alpha$ is a double root of $f(x)$ when $f'(\alpha)$ vanishes simultaneously with $f(\alpha)$. Formula (4) also shows that only under this assumption is $f(x)$ divisible by $(x-\alpha)^2$. This conclusion may be carried further and leads to the theorem:

\medskip
\noindent III. \gesperrt{The necessary and sufficient condition that $\alpha$ be an $m$-fold root of $f(x)$ is that $f(\alpha),f'(\alpha),f''(\alpha),\ldots,f^{(m-1)}(\alpha)$ vanish, while $f^{(m)}(\alpha)$ does not vanish.}

\begin{center}
{\large \S.\ 31.}

\vspace{0.5em}
{\large\gesperrt{Continuity of Integral Functions.}}
\end{center}

If, as before,
\[
 f(x)=a_0x^n+a_1x^{n-1}+\cdots+a_n
\]
is an integral function of $x$, we first prove the following theorem:

\medskip


% ===== scan leaf 129, printed 104 =====
% !PRINTED 104
\weberpage{104}

\noindent IV. \gesperrt{$f(x)$ becomes infinitely large at the same time as $x$.}
\medskip

This means that, if $C$ is an arbitrarily given positive number, the positive number $R$ can be so chosen that
\[
 |f(x)|>C
\]
as soon as the absolute value of $x$ becomes greater than $R$; or, expressed geometrically: one can describe in the $x$-plane a circle around the origin such that outside this circle the absolute value of $f(x)$ no longer falls below $C$, however large $C$ is assumed. The theorem is evident for the case of a simple power $x^n$; for if $r$ is the absolute value of $x$, then $r^n$ is the absolute value of $x^n$, and $r^n$, as soon as $r$ has become greater than $1$, is greater than $r$ and hence grows with $r$ into the infinite.

To prove the theorem in general, let us denote the absolute values of $x,a_0,a_1,a_2,\ldots,a_n$ by $r,c_0,c_1,c_2,\ldots,c_n$. Then
\begin{equation}
 |f(x)|=r^n\left|a_0+\frac{a_1}{x}+\frac{a_2}{x^2}+\cdots+\frac{a_n}{x^n}\right|,
\tag{2}% [sic] Section 31 has no printed equation (1); its opening polynomial is unnumbered and the first printed tag is (2). The source-original numbering gap is preserved.
\end{equation}
and, since by the theorems proved in the Introduction the absolute value of a sum cannot be smaller than the difference and cannot be greater than the sum of the absolute values of the summands,
\[
\left|a_0+\frac{a_1}{x}+\frac{a_2}{x^2}+\cdots+\frac{a_n}{x^n}\right|
\ge c_0-\left|\frac{a_1}{x}+\frac{a_2}{x^2}+\cdots+\frac{a_n}{x^n}\right|
\ge c_0-\frac{c_1}{r}-\frac{c_2}{r^2}-\cdots-\frac{c_n}{r^n}.
\]
Since $c_0,c_1,\ldots,c_n$ are fixed constants, one can choose $R$ so large that, as soon as $r>R$,
\[
 \frac{c_1}{r}+\frac{c_2}{r^2}+\cdots+\frac{c_n}{r^n}
\]
becomes arbitrarily small, for example smaller than $\tfrac12c_0$. Then
\[
 \left|a_0+\frac{a_1}{x}+\frac{a_2}{x}+\cdots+\frac{a_n}{x^n}\right|>\tfrac12c_0,% [sic] Weber prints a_2/x without exponent 2 in this isolated line; source-original error preserved to match the scan.
\]
and hence by (2)
\[
 |f(x)|>\tfrac12c_0R^n,
\]
so that, if one assumes
\[
 R^n>\frac{2C}{c_0}
\]
%CONT>>


% ===== scan leaf 130, printed 105 =====
% !PRINTED 105
\weberpage{105}
%<<CONT

\begin{equation}
 |f(x)|>C.
\tag{3}
\end{equation}
To this is attached the theorem:

\medskip
\noindent V. \gesperrt{$f(x)$ is a continuous function of $x$.}
\medskip

By this the following is meant. If the absolute value of $x$ is smaller than an arbitrary finite magnitude $R$, then, however small the positive magnitude $\omega$ is assumed, one can find a positive magnitude $\varepsilon$ such that
\begin{equation}
 |f(x+h)-f(x)|<\omega
\tag{4}
\end{equation}
as long as the absolute value $\varrho$ of $h$ remains smaller than $\varepsilon$, and indeed in such a way that $\varepsilon$ depends only on the choice of $\omega$, not on $x$.

By § 13, (2),
\begin{equation}
\begin{split}
 f(x+h)-f(x)&=h f'(x)+\frac{h^2}{1\cdot2}f''(x)+\cdots+\frac{h^n}{\Pi(n)}f^{(n)}(x)\\
 &=h\left\{f'(x)+\frac{h}{2}f''(x)+\cdots+a_0h^{n-1}\right\}.
\end{split}
\tag{5}
\end{equation}
If now we assume the absolute value of $h$ smaller than $1$, then the absolute value of the magnitude in parentheses
\begin{equation}
 f'(x)+\frac{h}{2}f''(x)+\cdots+a_0h^{n-1}
\tag{6}
\end{equation}
is smaller than a certain non-zero number $k$ which one obtains if in (6) one replaces $h$ by $1$, the coefficients $a_0,a_1,\ldots,a_{n-1}$ by their absolute values and $x$ by $R$, because thereby every individual term of the sum (6) is replaced by its absolute value or by a larger number. If one chooses $\varepsilon$ so small that
\[
 \varepsilon k<\omega
\]
then by (5) the requirement (4) is fulfilled as long as
\[
 \varrho<\varepsilon
\]
remains.

We can give the theorem of the continuity of the function $f(x)$ a somewhat different expression, which is important for what follows. By the theorem that the absolute value of a sum of two terms lies, in magnitude, between the sum and the difference of the absolute values of the terms (cf. Introduction, p. 19), we obtain from (5), if we put
%CONT>>


% ===== scan leaf 131, printed 106 =====
% !PRINTED 106
\weberpage{106}
%<<CONT

\[
 \varphi(h)=h f'(x)+\frac{h^2}{1\cdot2}f''(x)+\cdots+h^n a_0
\]
\[
 |f(x)|-|\varphi(h)|<|f(x+h)|<|f(x)|+|\varphi(h)|,
\]
or
\[
 -|\varphi(h)|<|f(x+h)|-|f(x)|<|\varphi(h)|.
\]
Now, if $\varrho<\varepsilon$, then $|\varphi(h)|<\omega$, and we can therefore also say:

\medskip
\noindent VI. \gesperrt{The absolute value of the difference}
\[
 |f(x+h)|-|f(x)|
\]
\gesperrt{is smaller than an arbitrarily given magnitude $\omega$ as soon as the absolute value of $h$ is smaller than a sufficiently small $\varepsilon$.}
\medskip

From this we can still conclude, if we put $x=y+iz$ and assume $h$ either real or purely imaginary, that $|f(y+zi)|$ is, for fixed $z$, a continuous function of $y$, and for fixed $y$, a continuous function of $z$.

\medskip
\noindent VII. \gesperrt{If the coefficients $a_0,a_1,\ldots,a_n$ and the variable $x$ are real, then $x$ may be chosen, in absolute value, so large that the sign of $f(x)$ agrees with the sign of the first term $a_0x^n$; thus, for positive $a_0$ and even $n$, it becomes positive, while for odd $n$ it becomes negative for negative $x$ and positive for positive $x$.}

For in
\[
 f(x)=x^n\left(a_0+\frac{a_1}{x}+\frac{a_2}{x^2}+\cdots+\frac{a_n}{x^n}\right)
\]
one can choose $x$ so large in absolute value that the absolute value of
\[
 \frac{a_1}{x}+\frac{a_2}{x^2}+\cdots+\frac{a_n}{x^n}
\]
falls below that of $a_0$, and therefore the expression in parentheses
\[
 a_0+\frac{a_1}{x}+\frac{a_2}{x^2}+\cdots+\frac{a_n}{x^n}
\]
has the sign of $a_0$, which it then retains for every absolutely larger $x$.


% ===== scan leaf 132, printed 107 =====
% !PRINTED 107
\weberpage{107}
\pdfbookmark[2]{§ 32. Change of sign of f(x). Roots of equations of odd degree and of pure equations}{bm37}

\begin{center}
{\large \S.\ 32.}

\vspace{0.5em}
{\large\gesperrt{Change of Sign of $f(x)$. Roots of Equations of Odd Degree and of Pure Equations.}}
\end{center}

On the basis of the theorems of the preceding paragraph we can prove the following important theorem.

\medskip
\noindent VIII. \gesperrt{If the coefficients of $f(x)$ are real and there exist two real values $a,b$ of $x$ such that $f(a)$ and $f(b)$ have opposite signs, then between $a$ and $b$ there is at least one root $\xi$ of the equation $f(x)=0$.}
\medskip

We shall assume that $a<b$ and that $f(a)$ is negative, $f(b)$ positive. We divide the numbers between $a$ and $b$ into two parts $\mathfrak A$ and $\mathfrak B$ as follows: a number $\alpha$ belongs to $\mathfrak A$ if, between $\alpha$ and $a$, including the endpoints, there is no $x$ for which $f(x)$ becomes positive; and a number $\beta$ belongs to $\mathfrak B$ if between $\beta$ and $a$ there is at least one value of $x$ for which $f(x)$ becomes positive. A third possibility is plainly not possible, and every number between $a$ and $b$ is therefore placed in $\mathfrak A$ or in $\mathfrak B$. We also count the numbers smaller than $a$ as belonging to $\mathfrak A$, and those greater than $b$ as belonging to $\mathfrak B$.

If $\alpha$ belongs to $\mathfrak A$, then every $x$ smaller than $\alpha$ also belongs to $\mathfrak A$; and if $\beta$ belongs to $\mathfrak B$, every larger $x$ belongs to $\mathfrak B$. Thus every $\beta$ is greater than every $\alpha$, and the two numerical domains $\mathfrak A,\mathfrak B$ form a \gesperrt{cut} (Introduction, p. 5), generated by a number $\xi$, such that every number $\alpha$ smaller than $\xi$ belongs to $\mathfrak A$, and every number $\beta$ greater than $\xi$ belongs to $\mathfrak B$.

If now, for some number $x_0$, the value of $f(x_0)$ is negative or positive, then by V of the preceding paragraph a positive magnitude $\varepsilon$ can be so determined that for every $x$ between $x_0-\varepsilon$ and $x_0+\varepsilon$ the values of $f(x)$ are always negative or always positive. It follows from this that $f(\xi)$ can be neither negative nor positive and must therefore be zero. For if $f(\xi)$ were negative, then $f(x)$ would also be negative between $\xi$ and $\xi+\varepsilon$; these values would therefore belong to $\mathfrak A$, although they are greater than $\xi$. If $f(\xi)$ were positive, then $f(x)$ would be positive as long as $x$ lay between $\xi-\varepsilon$ and $\xi$; these values, although smaller than $\xi$, would therefore belong to $\mathfrak B$. Both alternatives are impossible by the definition of
%CONT>>


% ===== scan leaf 133, printed 108 =====
% !PRINTED 108
\weberpage{108}
%<<CONT

$\xi$. The same conclusion may be drawn in quite the same way when $f(a)$ is positive and $f(b)$ negative, and so our theorem is proved.

From this we draw two important consequences. If $f(\xi)$ vanishes, then an interval
\[
 \xi-\varepsilon<x<\xi+\varepsilon
\]
can be so determined that $f(x)$ does not vanish in this interval except for $x=\xi$. Then four cases are possible with respect to the signs of $f(x)$ in the two intervals $\xi-\varepsilon<x<\xi$ and $\xi<x<\xi+\varepsilon$, which are put together in the following scheme:
\[
\begin{array}{lcc}
&\multicolumn{2}{c}{\xi-\varepsilon<\xi<\xi+\varepsilon}\\
1.\ f(x)&+&-\\
2.\ f(x)&-&+\\
3.\ f(x)&+&+\\
4.\ f(x)&-&-.
\end{array}
\]
If $f^{(\nu)}(x)$ is the first among the derivatives of $f(x)$ which is different from zero for $x=\xi$, then the formula [§ 13, (2)]
\[
 f(\xi+h)=\frac{h^\nu}{\Pi(\nu)}f^{(\nu)}(\xi)+\cdots,
\]
shows that these four cases are distinguished by the following signs:
\[
\begin{array}{lll}
1.&\nu\text{ odd}& f^{(\nu)}(\xi)<0\\
2.&\nu\text{ odd}& f^{(\nu)}(\xi)>0\\
3.&\nu\text{ even}& f^{(\nu)}(\xi)>0\\
4.&\nu\text{ even}& f^{(\nu)}(\xi)<0
\end{array}
\]
We say that in the first case $f(x)$ passes \gesperrt{decreasing} through zero, and in the second \gesperrt{increasing} through zero. In case 3, $f(\xi)=0$ is a \gesperrt{minimum}, in case 4 a \gesperrt{maximum}.

In particular, if $f'(\xi)$ is different from zero, so that $\xi$ is a simple root, then the positive or negative sign of $f'(\xi)$ is the characteristic mark for the increase or decrease of the function $f(x)$ as $x$ passes through $\xi$.

With these aids we still prove two further important theorems.

\medskip
\noindent IX. \gesperrt{The equation}\par
\[
 x^n-a=0
\]
\gesperrt{has, when $a$ is a positive number, one and only one positive root.}\par
\medskip


% ===== scan leaf 134, printed 109 =====
% !PRINTED 109
\weberpage{109}

For if we put
\begin{equation}
 f(x)=x^n-a,
\tag{1}
\end{equation}
then by § 31, VII, $f(x)$ is positive for sufficiently large positive $x$, and negative for $x=0$; hence there is a positive value of $x$, denoted by $\sqrt[n]{a}$, for which $f(x)$ vanishes. That there can be only one such value follows from the fact that $x^n$, as long as $x$ remains positive, increases continually with $x$.

If $n$ is odd, then also for negative $a$ there exists one and only one negative root of (1), which is proved in the same way.

\medskip
\noindent X. \gesperrt{An equation of odd degree with real coefficients always has at least one real root.}
\medskip

For by theorem VII, $f(x)$, if the degree is odd, can become both positive and negative; hence, by VIII, $f(x)=0$ has a root.

Finally we prove:

\medskip
\noindent XI. \gesperrt{A pure equation always has a root.}
\medskip

By a pure equation we understand an equation of the form
\begin{equation}
 x^n-a=0,
\tag{2}
\end{equation}
where $a$ is an arbitrary complex magnitude. That this equation always has a root for real positive $a$ has already been proved above; and for $a=0$ it is plainly satisfied by $x=0$.

The general proof can be divided into two parts; it is sufficient to assume first $n=2$, and then $n$ odd. For if the magnitudes $\sqrt[n]{a}$ and $\sqrt[2]{a}$, or $\sqrt{a}$, exist for every $a$, then the existence of
\[
 \sqrt[2n]{a}=\sqrt[2]{\sqrt[n]{a}},\qquad \sqrt[4n]{a}=\sqrt[2]{\sqrt[2n]{a}}\ldots,
\]
and hence of $\sqrt[n]{a}$ for every arbitrary $n$, is secured. It would even suffice to prove the existence of $\sqrt[n]{a}$ for the case in which $n$ is a prime number. But from this no particular advantage is at first to be drawn.

We denote the complex magnitude $a$ by $b+ci$ and put, since $x$ will in general also be complex,
\[
 x=y+iz.
\]


% ===== scan leaf 135, printed 110 =====
% !PRINTED 110
\weberpage{110}

Then we first have to show that
\begin{equation}
 (y+iz)^2=b+ic,
\tag{3}
\end{equation}
whatever real values $b,c$ may have, can be satisfied by real $y,z$. Equation (3), however, is fulfilled if
\begin{equation}
 y^2-z^2=b,
\qquad 2yz=c,
\tag{4}
\end{equation}
and hence
\[
 (y^2+z^2)^2=b^2+c^2
\]
therefore
\begin{equation}
 y^2+z^2=\sqrt{b^2+c^2}.
\tag{5}
\end{equation}
By IX there always exists a positive square root $\sqrt{b^2+c^2}$. From (4) and (5) it now follows that
\[
 y^2=\frac{b+\sqrt{b^2+c^2}}{2},\qquad
 z^2=\frac{-b+\sqrt{b^2+c^2}}{2},
\]
and the two magnitudes $\pm b+\sqrt{b^2+c^2}$ are plainly positive, since $b^2+c^2>b^2$, hence also $\sqrt{b^2+c^2}>\sqrt{b^2}$.% [sic] The strict inequalities fail when c=0: then b^2+c^2=b^2 and one of the two magnitudes is 0 rather than positive. Non-negativity suffices for the construction; the source-original error is preserved rather than silently normalized.
\ Thus we obtain
\[
 y=\pm\sqrt{\frac{b+\sqrt{b^2+c^2}}{2}},\qquad
 z=\pm\sqrt{\frac{-b+\sqrt{b^2+c^2}}{2}},
\]
but one sign is determined by the other through the condition
\[
 2yz=c;
\]
one of the two signs remains arbitrary, so that we obtain not merely one, but two opposite solutions of (3).

Further let $n$ be odd; we take the equation to be solved in the form
\begin{equation}
 (y+iz)^n=b+ci.
\tag{6}
\end{equation}
If $c=0$, so that $a$ is real, then we can put $z=0$ and by IX obtain a real value of $x$. But if we assume $c$ different from zero, then from (6) (Introduction, p. 19) there follows
\begin{equation}
 (y-iz)^n=b-ci.
\tag{7}
\end{equation}
Multiplying the two equations (6), (7) by one another gives
\[
 (y^2+z^2)^n=b^2+c^2,
\]
and hence, by IX,
%CONT>>


% ===== scan leaf 136, printed 111 =====
% !PRINTED 111
\weberpage{111}
\pdfbookmark[2]{§ 33. Solution of pure equations by trigonometric functions}{bm38}
%<<CONT

\begin{equation}
 y^2+z^2=\sqrt[n]{b^2+c^2},
\tag{8}
\end{equation}
which determines the absolute value of $x$. Furthermore, from (6) and (7) one obtains
\begin{equation}
 \varphi(y,z)={ (y-iz)^n(b+ic)-(y+iz)^n(b-ic) \over 2i}=0.
\tag{9}
\end{equation}
Now $\varphi(y,z)$ is an integral homogeneous function of the $n$th degree in the two variables $y,z$, and indeed with real coefficients, since replacing $i$ by $-i$ changes nothing in $\varphi(y,z)$. If it is ordered by descending powers of $y$, the first term is $cy^n$, hence, by assumption, different from zero. Putting
\[
 y=\lambda z,
\]
we obtain from (9)
\begin{equation}
 \varphi(\lambda,1)=0,
\tag{10}
\end{equation}
therefore an equation of the $n$th degree for $\lambda$, which by X certainly has a root. Once $\lambda$ is determined, it follows from (8) that
\begin{equation}
 z=\sqrt{\frac{\sqrt[n]{b^2+c^2}}{1+\lambda^2}},\qquad
 y=\lambda\sqrt{\frac{\sqrt[n]{b^2+c^2}}{1+\lambda^2}}.
\tag{11}
\end{equation}
Through this the equations (8), (9) are in fact satisfied; from (9), however, follows
\[
 \frac{(y+iz)^n}{b+ic}=\frac{(y-iz)^n}{b-ic},
\]
and from (8) that the common value of these two fractions is $\pm1$. Therefore we have
\[
 (y+iz)^n=\pm(b+ci),\qquad (y-iz)^n=\pm(b-ci),
\]
and thus the sign of the square root in (11) can be so determined that equations (6) and (7) are fulfilled.

\begin{center}
{\large \S.\ 33.}

\vspace{0.5em}
{\large\gesperrt{Solution of Pure Equations by Trigonometric Functions.}}
\end{center}

The solvability of a pure equation can be shown far more completely and simply if one presupposes the trigonometric functions and their simplest properties as known. These functions are, it is true, foreign to
%CONT>>


% ===== scan leaf 137, printed 112 =====
% !PRINTED 112
\weberpage{112}
%<<CONT

algebra proper, and for that reason it is more satisfactory to answer the \gesperrt{questions of principle} without their aid, as has been done in the preceding and will also be done later. For application and for a more convenient visualization, however, we shall not dispense with this aid.

If, when $b,c$ are real numbers, we put
\begin{equation}
 b=r\cos\varphi,
\qquad c=r\sin\varphi,
\tag{1}
\end{equation}
then, if $r$ is assumed positive, the angle $\varphi$ is thereby determined up to a multiple of $2\pi$. It will be completely determined if we fix an interval of length $2\pi$ in which it is to lie, for example
\[
 -\pi<\varphi\le \pi.
\]
In geometrical interpretation $r,\varphi$ are the polar coordinates of the point whose rectangular coordinates are $b,c$. The absolute value of the complex magnitude $a=b+ci$ is $r$, and we shall call $\varphi$ the \gesperrt{phase} of this complex magnitude.

The product of two complex magnitudes
\begin{equation}
\begin{aligned}
 a&=r(\cos\varphi+i\sin\varphi),\\
 a'&=r'(\cos\varphi'+i\sin\varphi'),
\end{aligned}
\tag{2}
\end{equation}
is obtained in the form
\begin{equation}
 aa'=rr'[\cos(\varphi+\varphi')+i\sin(\varphi+\varphi')],
\tag{3}
\end{equation}
and from this, if $n$ is an arbitrary positive whole number, or even a negative one, one obtains the theorem named after \gesperrt{Moivre},
\begin{equation}
 (\cos\varphi+i\sin\varphi)^n=\cos n\varphi+i\sin n\varphi,
\tag{4}
\end{equation}
which has led to regarding the magnitude $\cos\varphi+i\sin\varphi$ as an exponential function with imaginary exponent, and hence to setting
\[
 \cos\varphi+i\sin\varphi=e^{i\varphi}
\]
Thereby formulas (3) and (4) take the form
\begin{equation}
 e^{i\varphi}e^{i\varphi'}=e^{i(\varphi+\varphi')},\qquad (e^{i\varphi})^n=e^{in\varphi}.
\tag{5}
\end{equation}
If accordingly we put
\begin{equation}
 a=r(\cos\varphi+i\sin\varphi)
\tag{6}
\end{equation}
and by $\sqrt[n]{r}$ understand the positive value existing and completely determined by § 32, IX, whose $n$th power $=r$, then
%CONT>>


% ===== scan leaf 138, printed 113 =====
% !PRINTED 113
\weberpage{113}
%<<CONT

\begin{equation}
 x=\sqrt[n]{r}\left(\cos\frac{\varphi}{n}+i\sin\frac{\varphi}{n}\right)
\tag{7}
\end{equation}
is a magnitude whose $n$th power is $=a$, hence a root of the equation
\begin{equation}
 x^n=a,
\tag{8}
\end{equation}
where $n$ is understood to be an arbitrary positive whole number. But every one of the values
\begin{equation}
 x_k=\sqrt[n]{r}\left[\cos\left(\frac{\varphi+2k\pi}{n}\right)+i\sin\left(\frac{\varphi+2k\pi}{n}\right)\right],
\tag{9}
\end{equation}
satisfies the same requirement, when $k$ is an arbitrary integer. In (9), however, not more than $n$ different values are contained, namely those obtained by putting
\begin{equation}
 k=0,1,2,\ldots,n-1;
\tag{10}
\end{equation}
for if $k$ is increased by a multiple of $n$, the value of (9) is not changed, while the $n$ values (10) give only different values of $x$, because they have different phases.

Equation (8) consequently has not merely one, but $n$ different roots. The geometrical images of the values $x_k$ all lie on a circle with radius $\sqrt[n]{r}$, and are separated from one another by the angle $\frac{2\pi}{n}$. The first of these points is obtained by dividing the given angle $\varphi$ into $n$ equal parts. The radius $\sqrt[n]{r}$ is greater or smaller than $r$ according as $r$ is smaller or greater than $1$. The adjacent figure shows these relations for $n=3$ under the assumption $r>1$.


% Exact 300-dpi crop from the registered Volume-I scan, printed Fig. 1.
\weberfigureimage{weber-b1-fig01.png}{0.50}


One can obtain the different values $x_k$ by multiplying the first of them,
\[
 x_0=\sqrt[n]{r}\left(\cos\frac{\varphi}{n}+i\sin\frac{\varphi}{n}\right),
\]
by the different values of
%CONT>>


% ===== scan leaf 139, printed 114 =====
% !PRINTED 114
\weberpage{114}
\pdfbookmark[2]{§ 34. Removal of the second term of an equation}{bm39}
%<<CONT

\begin{equation}
 \varepsilon_k=\cos\frac{2k\pi}{n}+i\sin\frac{2k\pi}{n}.
\tag{11}
\end{equation}
The quantities $\varepsilon_k$ are roots of the equation
\begin{equation}
 x^n=1,
\tag{12}
\end{equation}
and are called the $n$th \gesperrt{roots of unity}. There are $n$, and only $n$, different values of them; for odd $n$ only one is real, namely for $k=0$, and for even $n$ two are real, namely for $k=0$ and $k=\frac n2$. The geometrical images of the $n$th roots of unity are the vertices of a regular $n$-gon inscribed in the circle of radius $1$. Their algebraic determination is the subject of the \gesperrt{theory of cyclotomy}. If we denote by $\varepsilon$ the value of $\varepsilon_k$ for $k=1$, thus put
\[
 \varepsilon=\cos\frac{2\pi}{n}+i\sin\frac{2\pi}{n},
\]
then, by \gesperrt{Moivre's} theorem,
\[
 \varepsilon_k=\varepsilon^k,
\]
so that all $n$th roots of unity are represented as powers of one among them.


\begin{center}
{\large \S.\ 34.}

\vspace{0.5em}
{\large\gesperrt{Removal of the Second Term of an Equation.}}
\end{center}

The certainty of the existence of the roots of pure equations puts us in a position to prove the existence of roots for equations of the second, third, and fourth degrees, or, as one also says, for quadratic, cubic, and biquadratic equations, by reducing the determination of their roots to the solution of pure equations. We shall return to this question later from various more general points of view, but here we shall briefly discuss the older methods of solution, which, even if they give little insight into the general connection of this question, are nevertheless very simple in application. They arouse the appearance as if one were dealing with a method capable of generalization to higher equations, which is not the case. First the following general remark.

Let us, for the sake of simplicity, take in the function
%CONT>>


% ===== scan leaf 140, printed 115 =====
% !PRINTED 115
\weberpage{115}
%<<CONT

\begin{equation}
f(x)=x^n+a_1x^{n-1}+a_2x^{n-2}+\cdots+a_n
\tag{1}
\end{equation}
the coefficient of $x^n$ equal to $1$, which in general is obtained by division by the coefficient of $x^n$. Then by the simple substitution
\begin{equation}
x=y-\frac{a_1}{n},
\tag{2}
\end{equation}
we can transform $f(x)$ into another function $\varphi(y)$ of the $n$th degree in which the second term $y^{n-1}$ does not occur; for, by the binomial theorem,
\[
 x^n=y^n-a_1y^{n-1}+\frac{n-1}{2n}a_1^2y^{n-2}+\cdots,
\]
\[
 x^{n-1}=y^{n-1}-\frac{n-1}{n}a_1y^{n-2}+\cdots,
 \qquad x^{n-2}=y^{n-2}-\cdots,
\]
and consequently
\[
 f(x)=y^n+\left(a_2-\frac{n-1}{2n}a_1^2\right)y^{n-2}+\cdots .
\]

If we actually carry out the transformation for the cases $n=2,3,4$, we obtain
\begin{equation}
\begin{array}{rcl}
n=2, & \varphi(y)&=y^2+a,\\[0.25em]
n=3, & \varphi(y)&=y^3+ay+b,\\[0.25em]
n=4, & \varphi(y)&=y^4+ay^2+by+c,
\end{array}
\tag{3}
\end{equation}
provided that we put
\begin{equation}
\begin{array}{rcll}
\text{for }n=2:&& a&=-\dfrac{a_1^2}{4}+a_2,\\[0.75em]
\text{for }n=3:&& a&=-\dfrac{a_1^2}{3}+a_2,\\[0.75em]
&& b&=\dfrac{2a_1^3}{27}-\dfrac{a_1a_2}{3}+a_3,\\[0.75em]
\text{for }n=4:&& a&=-\dfrac{3a_1^2}{8}+a_2,\\[0.75em]
&& b&=\dfrac{a_1^3}{8}-\dfrac{a_1a_2}{2}+a_3,\\[0.75em]
&& c&=-\dfrac{3a_1^4}{256}+\dfrac{a_1^2a_2}{16}-\dfrac{a_1a_3}{4}+a_4.
\end{array}
\tag{4}
\end{equation}
Thus the solution of the equation $f(x)=0$ is reduced by (2) to the solution of $\varphi(y)=0$. This first gives, for $n=2$,
%CONT>>


% ===== scan leaf 141, printed 116 =====
% !PRINTED 116
\weberpage{116}
\pdfbookmark[2]{§ 35. Cubic equations. Cardano's formula}{bm40}
%<<CONT

\[
 y=\pm\sqrt{-a},
\]
and consequently
\[
 x=-\frac{a_1}{2}\pm\sqrt{\frac{a_1^2}{4}-a_2}
   =\frac{-a_1\pm\sqrt{a_1^2-4a_2}}{2}.
\]
We therefore have two roots
\begin{equation}
\alpha=\frac{-a_1+\sqrt{a_1^2-4a_2}}{2},\qquad
\beta =\frac{-a_1-\sqrt{a_1^2-4a_2}}{2}.
\tag{5}
\end{equation}

\begin{center}
{\large \S.\ 35.}

\vspace{0.5em}
{\large\gesperrt{Cubic Equations. Cardano's Formula.}}
\end{center}

We take the equation of the third degree in the reduced form
\begin{equation}
y^3+ay+b=0
\tag{1}
\end{equation}
and introduce two new unknowns $u,v$ by putting
\begin{equation}
y=u+v.
\tag{2}
\end{equation}
Substitution in (1) gives
\begin{equation}
u^3+v^3+(3uv+a)(u+v)+b=0.
\tag{3}
\end{equation}

We determine one of the two quantities $u,v$ by the other according to the equation
\begin{equation}
3uv=-a,
\tag{4}
\end{equation}
and obtain from (3)
\begin{equation}
u^3+v^3=-b.
\tag{5}
\end{equation}
But from (4) and (5), $u^3$ and $v^3$ can be determined by a square root. Namely one obtains
\begin{equation}
(u^3-v^3)^2=(u^3+v^3)^2-4u^3v^3
=b^2+\frac{4a^3}{27}.
\tag{6}
\end{equation}
Put, for abbreviation,
\begin{equation}
R=\frac{b^2}{4}+\frac{a^3}{27};
\tag{7}
\end{equation}
then
\[
 u^3-v^3=2\sqrt R.
\]


% ===== scan leaf 142, printed 117 =====
% !PRINTED 117
\weberpage{117}

We give $\sqrt R$ one of the two signs and, according to (5), obtain
\begin{equation}
\begin{array}{rclcrcl}
u^3&=&-\dfrac b2+\sqrt R,&& v^3&=&-\dfrac b2-\sqrt R,\\[0.75em]
u&=&\sqrt[3]{-\dfrac b2+\sqrt R},&& v&=&\sqrt[3]{-\dfrac b2-\sqrt R},
\end{array}
\tag{8}
\end{equation}
and therefore
\begin{equation}
y=\sqrt[3]{-\frac b2+\sqrt R}+\sqrt[3]{-\frac b2-\sqrt R}.
\tag{9}
\end{equation}

Multiplication of the two expressions (8) gives
\[
 u^3v^3=-\frac{a^3}{27}
\]
and shows that, if for $u$ a determinate one among the three values of the cube root is taken,
\begin{equation}
v=-\frac{a}{3u}=-\frac{a}{3\sqrt[3]{-\dfrac b2+\sqrt R}}
\tag{10}
\end{equation}
is in any case one of the three values of $\sqrt[3]{-\dfrac b2-\sqrt R}$. But it follows from (4) that only when we take this value for $v$ is equation (1) satisfied by $y=u+v$. Thus, corresponding to the three values of $u$, we obtain three roots of equation (1). From the existence of one root of the cubic equation, however, the existence of three already follows also from the fact that it has already been proved that the quadratic equation has two roots.

To derive the other two roots from the one root (9), put
\[
 \alpha=u+v
\]
and divide $y^3+ay+b$ by $y-\alpha$. The quotient whose roots are the two other roots $\beta,\gamma$ of the cubic equation (1) is
\[
 y^2+\alpha y+\alpha^2+a=0,
\]
whence, by (4) of the preceding paragraph, one finds
\[
 y=\frac{-\alpha\pm\sqrt{-3\alpha^2-4a}}{2}.
\]
If here one puts $\alpha=u+v$ and, by (4), $a=-3uv$, then it follows that
%CONT>>


% ===== scan leaf 143, printed 118 =====
% !PRINTED 118
\weberpage{118}
\pdfbookmark[2]{§ 36. Cayley's expression for Cardano's formula}{bm41}
%<<CONT

\[
 y=\frac{-(u+v)\pm\sqrt{-3(u+v)^2+12uv}}{2}
  =\frac{-(u+v)\pm\sqrt{-3}(u-v)}{2}.
\]
We therefore put
\begin{equation}
\varepsilon=\frac{-1+\sqrt{-3}}{2},
\tag{11}
\end{equation}
from which
\begin{equation}
\varepsilon^3=1,\qquad \varepsilon^2+\varepsilon+1=0,
\qquad
\varepsilon^2=\frac{-1-\sqrt{-3}}{2}
\tag{12}
\end{equation}
follows. Thus one obtains the three roots of the cubic equation (6)
% [sic] Source-original mathematical cross-reference: the printed text says (6), although the cubic equation itself is (1).
\begin{equation}
\begin{array}{rcl}
\alpha&=&u+v,\\
\beta&=&\varepsilon u+\varepsilon^2v,\\
\gamma&=&\varepsilon^2u+\varepsilon v.
\end{array}
\tag{13}
\end{equation}
Here $\varepsilon$ is a third root of unity different from $1$, expressed algebraically without the trigonometric functions. If one substitutes $\varepsilon u$ or $\varepsilon^2u$ for $u$, and hence $\varepsilon^2v$ or $\varepsilon v$ for $v$, the three quantities $\alpha,\beta,\gamma$ are interchanged with one another.

The expression (9) for the root of a cubic equation is called \gesperrt{Cardano's formula}.

\begin{center}
{\large \S.\ 36.}

\vspace{0.5em}
{\large\gesperrt{Cayley's Expression for Cardano's Formula.}}
\end{center}

Every cube root has, as we have seen, three different values, which one obtains from one of them by multiplication by $1,\varepsilon,\varepsilon^2$. Thus each of the two quantities $u,v$, as they are defined by (8), \S\,35, also has three values, and the sum $u+v$, if one takes account only of the conditions (8), has \gesperrt{nine different} values. Of these, however, only three give solutions of the cubic equation, and only by adjoining the relation (4), \S\,35,
\begin{equation}
3uv=-a
\tag{1}
\end{equation}
are the usable values singled out from among the nine.


% ===== scan leaf 144, printed 119 =====
% !PRINTED 119
\weberpage{119}
\pdfbookmark[2]{§ 37. The biquadratic equations}{bm42}

This is a defect of \gesperrt{Cardano's formula}, which accordingly by itself does not yet suffice to give the roots of the cubic equation unambiguously. \gesperrt{Cayley} remedied this defect by the following representation. Define two new quantities $\xi,\eta$ by the equations
\begin{equation}
u=\xi^2\eta,
\qquad
v=\xi\eta^2,
\tag{2}
\end{equation}
from which, by multiplication and use of (1), one obtains
\begin{equation}
\xi\eta=\sqrt[3]{-\frac a3}.
\tag{3}
\end{equation}
Consequently, if this value is inserted in (2) and the expressions \S\,35, (8) for $u,v$ are substituted, then
\begin{equation}
\xi=\sqrt[3]{\frac{3b}{2a}+\sqrt{\frac{9b^2}{4a^2}+\frac a3}},\qquad
\eta=\sqrt[3]{\frac{3b}{2a}-\sqrt{\frac{9b^2}{4a^2}+\frac a3}}.
\tag{4}
\end{equation}
Thus, if one now puts the root $y$ of the cubic equation in the form
\begin{equation}
y=\xi\eta(\xi+\eta),
\tag{5}
\end{equation}
one obtains an expression which, when $\xi$ is replaced, independently, by $\xi,\varepsilon\xi,\varepsilon^2\xi$ and $\eta$ by $\eta,\varepsilon\eta,\varepsilon^2\eta$, assumes not nine but \gesperrt{only the three values}\footnotemark
\[
\begin{array}{c}
\xi^2\eta+\xi\eta^2,\\
\varepsilon\xi^2\eta+\varepsilon^2\xi\eta^2,\\
\varepsilon^2\xi^2\eta+\varepsilon\xi\eta^2
\end{array}
\]
\footnotetext{\gesperrt{Cayley}, Phil. Mag. vol. XXI, 1861. Collectet mathematical papers vol. V, No. 310.}% [sic] The registered source prints ‘Collectet’; preserved as a source-original bibliographic anomaly. The callout is attached to the translated phrase introducing the three displayed values so it cannot strand after the display.



\begin{center}
{\large \S.\ 37.}

\vspace{0.5em}
{\large\gesperrt{The Biquadratic Equation.}}
\end{center}

A similar path to that used in solving the cubic equation by \gesperrt{Cardano's formula} may be taken for the solution of the equation of the fourth degree
\begin{equation}
 y^4+ay^2+by+c=0.
\tag{1}
\end{equation}
We put
%CONT>>


% ===== scan leaf 145, printed 120 =====
% !PRINTED 120
\weberpage{120}
%<<CONT

\[
 2y=u+v+w
\]
and obtain, if we write, for abbreviation,
\begin{equation}
 s=u^2+v^2+w^2,\qquad
 t=v^2w^2+w^2u^2+u^2v^2,
\tag{2}
\end{equation}
then
\[
 4y^2=s+2(vw+wu+uv),
\]
\[
 16y^4=s^2+4s(vw+wu+uv)+4t+8uvw(u+v+w).
\]
Substitution in (1) gives
\[
 s^2+4t+4as+16c+8(uvw+b)(u+v+w)
  +4(s+2a)(vw+wu+uv)=0.
\]
This equation is satisfied by the assumptions
\[
 s+2a=0,\qquad uvw+b=0,\qquad s^2+4t+4as+16c=0.
\]
With the help of the first of these three equations the third becomes
\[
 t=a^2-4c,
\]
and, putting back the values (2) for $s,t$,
\begin{equation}
\begin{array}{rcl}
 u^2+v^2+w^2&=&-2a,\\
 v^2w^2+w^2u^2+u^2v^2&=&a^2-4c,\\
 uvw&=&-b.
\end{array}
\tag{3}
\end{equation}
By \S\,7 these equations are satisfied if and only if $u^2,v^2,w^2$ are the roots of the cubic equation
\begin{equation}
 z^3+2az^2+(a^2-4c)z-b^2=0,
\tag{4}
\end{equation}
and if the signs of $u,v,w$ are so chosen that the last of equations (3) is satisfied. This last condition still allows four different determinations of signs, so that the four roots of the biquadratic equation are obtained in the following way:
\begin{equation}
\begin{array}{rcl}
2\alpha&=&u+v+w,\\
2\beta&=&u-v-w,\\
2\gamma&=&-u+v-w,\\
2\delta&=&-u-v+w.
\end{array}
\tag{5}
\end{equation}
Thus the solution of the biquadratic equation has been reduced to that of the cubic equation (4).

This equation is called a \gesperrt{cubic resolvent} of the biquadratic equation.


% ===== scan leaf 146, printed 121 =====
% !PRINTED 121
\weberpage{121}
\pdfbookmark[2]{§ 38. Proof of the fundamental theorem of algebra}{bm43}

\begin{center}
{\large \S.\ 38.}

\vspace{0.5em}
{\large\gesperrt{Proof of the Fundamental Theorem.}}
\end{center}

We now proceed to the proof of the \gesperrt{fundamental theorem of algebra}, namely that every equation $f(x)=0$ has at least one root. We first put forward some general propositions.

1. \gesperrt{If $S$ denotes an arbitrary system of real numbers all greater than a certain positive or negative number $C$ (that is, a system which contains no infinitely negative numbers), then there exists a lower bound for the numbers $S$.}

By a \gesperrt{lower bound} $g$ one is to understand a number which is not undercut by any number of the system $S$, but which is of such a nature that, if $\delta$ is an arbitrarily given positive quantity, then between $g$ and $g+\delta$ (including the endpoints) there always lies at least one number of the system $S$, however small $\delta$ may be.

The proof follows directly from the possibility of a cut $(\mathfrak A,\mathfrak B)$, constructed as follows: one throws a number $\alpha$ into $\mathfrak A$ if it is not undercut by any number of the system $S$, and a number $\beta$ into $\mathfrak B$ if it is undercut by at least one number in $S$. The number $g$ determined by this cut is not undercut by any number in $S$, for otherwise there would also be numbers smaller than $g$, and therefore belonging to $\mathfrak A$, which nevertheless are undercut by numbers in $S$; while, on the other hand, every number $\beta$ lying even ever so little above $g$ belongs to $\mathfrak B$ and is therefore undercut by a number in $S$. These are precisely the characteristic marks of the \gesperrt{lower bound}.

2. \gesperrt{If $S'$ is a part of $S$, then the numbers $S'$ also have a lower bound $g'$, and this is either equal to $g$ or greater than $g$.}

For it cannot be smaller than $g$, since otherwise numbers would occur in $S'$, and therefore also in $S$, which lie below $g$.

Now let $f(x)$ be a real function of $x$ which has only finite values in the interval
\begin{equation}
 a\le x\le b
\tag{1}
\end{equation}
%CONT>>


% ===== scan leaf 147, printed 122 =====
% !PRINTED 122
\weberpage{122}
%<<CONT

and which is also \gesperrt{continuous} in this interval. In agreement with \S\,31, V and VI, this means that
\begin{equation}
 f(x\pm h)-f(x)
\tag{2}
\end{equation}
falls, in absolute value, below an arbitrarily prescribed number $\eta$ throughout the interval whenever the positive $h$ is smaller than a certain number $\varepsilon$. For $x=a$ we take in (2) only the upper sign, and for $x=b$ only the lower sign, so as not to pass outside the interval with $x\pm h$.

For the values of such a function in the interval there is, by 1., a \gesperrt{lower bound} $g$, and we now prove the proposition:

3. \gesperrt{The function $f(x)$ assumes the value $g$ for some value $\xi$ of the interval, so that the lower bound becomes a minimum of the function $f(x)$ in the interval.}

The proof is as follows. By 2., the \gesperrt{lower bound} of the function values in a part of the interval (1) is either equal to or greater than $g$.

If now $f(a)=g$, what is to be proved is true. If, however, $f(a)>g$, then by the continuity of $f(x)$ one can give an interval $a\le x\le a+h$ in which $f(x)>g$ remains true, and hence the \gesperrt{lower bound} of $f(x)$ is greater than $g$.

We now construct in the interval (1) a cut $(\mathfrak A,\mathfrak B)$ of the following kind: we count a value $\alpha$ of the interval (1) as belonging to $\mathfrak A$ if the \gesperrt{lower bound} of $f(x)$ in the interval $a\le x\le \alpha$ is greater than $g$, and a value $\beta$ as belonging to $\mathfrak B$ if the \gesperrt{lower bound} in the interval $a\le x\le \beta$ is equal to $g$.

It is possible that $\mathfrak B$ consists only of the single value $b$; but then $f(b)=g$ must hold, and our assertion is fulfilled for $x=b$. For if $f(b)>g$, one could choose a quantity $g'$ between $f(b)$ and $g$, and, by continuity, an interval $b-h\le x\le b$ such that in this interval all function values $f(x)$ are greater than $g'$, so that their \gesperrt{lower bound} is equal to or greater than $g'$, and hence certainly greater than $g$. But since $b-h$ belongs to $\mathfrak A$, the \gesperrt{lower bound} is also greater than $g$ in the interval $a\le x\le b-h$, and consequently in the whole interval (1), contrary to the assumption.


% ===== scan leaf 148, printed 123 =====
% !PRINTED 123
\weberpage{123}

If therefore $f(b)>g$, the cut $(\mathfrak A,\mathfrak B)$ defines a number $\xi$ in the interior of the interval (1), and we can now show that
\begin{equation}
 f(\xi)=g
\tag{3}
\end{equation}
must hold.

If $f(\xi)>g$ and $f(\xi)>g'>g$, then, by the continuity of $f(x)$, we can determine an interval
\[
 \xi-h\le x\le \xi+h
\]
in which all function values $f(x)$ are greater than $g'$, and hence their \gesperrt{lower bound} is greater than $g$. But since $\xi-h$ belongs to $\mathfrak A$, the \gesperrt{lower bound} of $f(x)$ in the whole interval
\[
 a\le x\le \xi+h
\]
is also greater than $g$, while $\xi+h$ belongs to $\mathfrak B$, and this is the contradiction.

Let $f(x)$ now be an integral rational function with arbitrary complex or real coefficients, and let the variable $x$ also be allowed to take complex values.

By \S\,31, IV, if $C$ is an arbitrary positive value, the positive quantity $R$ can be so chosen that
\begin{equation}
 |f(x)|>C
\tag{4}
\end{equation}
whenever
\begin{equation}
 |x|\ge R.
\tag{5}
\end{equation}
For visualization we represent the region $(R)$ bounded by the inequality
\[
 |x|<R
\]
for the variable $x$ by a circular disc of radius $R$ in the plane of the variable $x=y+iz$.

If we take $C$ larger than any value of $|f(x)|$ in the interior of the region $(R)$, for example larger than $|f(0)|$, then $|f(x)|$ will certainly become smaller in the interior of $(R)$ than on its boundary. Since now in the whole interior of $(R)$ the function $|f(x)|$, which for abbreviation we shall now denote by $X$, is not negative, there is a \gesperrt{lower bound} $g$ for the values of $X$, and we have to prove the proposition:

4. \gesperrt{There exists a value $\xi$ of $x$ in the interior of the region $(R)$ such that}
\begin{equation}
 |f(\xi)|=g
\tag{6}
\end{equation}
%CONT>>


% ===== scan leaf 149, printed 124 =====
% !PRINTED 124
\weberpage{124}
%<<CONT

\gesperrt{so that here also the lower bound is a minimum.}

The \gesperrt{lower bound} of $X$ in any part of the region is either equal to $g$ or greater than $g$.

A domain of magnitudes determined by the inequalities
\begin{equation}
 |x|\le R,\qquad -R\le y\le \alpha
\tag{7}
\end{equation}
is represented in our Figure 2 by the segment $(PQ\alpha Q')$, which we shall call the segment $(P,\alpha)$.

% Exact 300-dpi crop from the registered Volume-I scan, printed Fig. 2.
\weberfigureimage{weber-b1-fig02.png}{0.62}

We now first determine a value $\eta$ of $y$ by a cut $(\mathfrak A,\mathfrak B)$ as follows: a number $\alpha$ between $-R$ and $+R$ is admitted into $\mathfrak A$ if the \gesperrt{lower bound} of $X$ in the segment $(P,\alpha)$ is greater than $g$, and a value $\beta$ is admitted into $\mathfrak B$ if the \gesperrt{lower bound} of $X$ in the segment $(P,\beta)$ is equal to $g$. This cut $(\mathfrak A,\mathfrak B)$ defines a number $\eta$ with the property that the \gesperrt{lower bound} of $X$ in the region $(P,y)$ is greater than $g$ if $y<\eta$, and equal to $g$ if $y>\eta$.

Now, by \S\,31, VI,
\[
 |f(\eta+iz)|
\]
is a continuous function of $z$ in the interval
\begin{equation}
 -\sqrt{R^2-\eta^2}\le z\le \sqrt{R^2-\eta^2},
\tag{8}
\end{equation}
which is denoted in Figure 2 by $M,M'$.

Thus by 3. this function receives, for some value $\zeta$ of $z$ in this interval, a minimum value $\gamma$, so that, if
\[
 \xi=\eta+i\zeta
\]
is put,
\begin{equation}
 |f(\xi)|=\gamma.
\tag{9}
\end{equation}


% ===== scan leaf 150, printed 125 =====
% !PRINTED 125
\weberpage{125}

It is now further easy to see that $\gamma=g$ must be the case; for since $g$ is the \gesperrt{lower bound} of all values of $X$ within $(R)$, $\gamma$ cannot first of all be smaller than $g$.

But $\gamma$ also cannot be greater than $g$; for all values which $X$ assumes on the segment $MM'$ are $\ge\gamma$. If, however, $\gamma>g$, then one can choose a quantity $g'$ such that $\gamma>g'>g$, and, by the continuity of $X$ proved in \S\,31, VI, numbers $\alpha,\beta$ can be so determined that
\[
 \alpha<\eta<\beta
\]
and such that in the whole region $(P,\beta)-(P,\alpha)=(\alpha,\beta)$ ($QNN'Q'$ in Figure 2) $X$ remains greater than $g'$. Consequently the \gesperrt{lower bound} of $X$ in $(\alpha,\beta)$ is greater than $g$. Since now the \gesperrt{lower bound} of $X$ in $(P,\alpha)$ is also greater than $g$, it is also greater than $g$ in $(P,\beta)$, which is composed of $(P,\alpha)$ and $(\alpha,\beta)$; but $\beta$ belongs to $\mathfrak B$, and this would give a contradiction with the definition of $\mathfrak B$.

Thus there remains only the possibility that $\gamma=g$, and proposition 4 is thereby proved; namely, there is a value $\xi$ in the interior of $(R)$ for which
\[
 |f(\xi)|=g
\]
holds, so that $|f(x)|$ attains a minimum value. We now prove:

5. \gesperrt{If $\alpha$ is any value of $x$ for which $f(\alpha)$ is different from zero, then $h$ can be chosen so that}
\begin{equation}
 |f(\alpha+h)|<|f(\alpha)|
\tag{10}
\end{equation}
\gesperrt{results.}

It then follows that, if $f(\xi)$ were different from zero, $g$ could not be the minimum of the function $|f(x)|$; but since this has been proved, one must have
\begin{equation}
 f(\xi)=0,
\tag{11}
\end{equation}
and $\xi$ is a root of the equation $f(x)=0$. The fundamental theorem will therefore be proved.

In the proof of 5. we make use of the theorem already proved (\S\,32), that a pure equation always has a root.

Some of the derived functions $f'(\alpha),f''(\alpha),\ldots$ may vanish; the last, $f^{(n)}(\alpha)$, which is equal to $\Pi(n)a_0$, is however
%CONT>>


% ===== scan leaf 151, printed 126 =====
% !PRINTED 126
\weberpage{126}
%<<CONT

different from zero. Thus suppose that $f'(\alpha),f''(\alpha),\ldots,f^{(m-1)}(\alpha)$ vanish, but $f^{(m)}(\alpha)$ does not vanish. We then have, by \S\,13,
\begin{equation}
 f(\alpha+h)=f(\alpha)+\frac{h^m}{\Pi(m)}f^{(m)}(\alpha)
 +\frac{h^{m+1}}{\Pi(m+1)}f^{(m+1)}(\alpha)+\cdots .
\tag{12}
\end{equation}

We choose, as is always possible by \S\,31, V, a positive number $\varepsilon$ so that
\begin{equation}
\left|\frac{h}{m+1}\frac{f^{(m+1)}(\alpha)}{f^{(m)}(\alpha)}
+\frac{h^2}{(m+1)(m+2)}\frac{f^{(m+2)}(\alpha)}{f^{(m)}(\alpha)}+\cdots\right|<1
\tag{13}
\end{equation}
whenever
\begin{equation}
 |h|<\varepsilon,
\tag{14}
\end{equation}
and a positive proper fraction $\delta$ such that
\begin{equation}
 \delta |f(\alpha)|<\left|\frac{\varepsilon^m f^{(m)}(\alpha)}{\Pi(m)}\right|,
\tag{15}
\end{equation}
which, if $f(\alpha)$ does not vanish, is likewise possible. We then determine $h$ (on the basis of \S\,32) from the equation
\begin{equation}
 \frac{h^m f^{(m)}(\alpha)}{\Pi(m)}=-\delta f(\alpha),
\tag{16}
\end{equation}
whence, by (15),
\[
 |h|<\varepsilon
\]
follows. From (12), when the value from (16) is substituted for $h^m$, we now get
\[
 f(\alpha+h)=(1-\delta)f(\alpha)
 -\delta f(\alpha)\left(\frac{h}{m+1}\frac{f^{(m+1)}(\alpha)}{f^{(m)}(\alpha)}
 +\frac{h^2}{(m+1)(m+2)}\frac{f^{(m+2)}(\alpha)}{f^{(m)}(\alpha)}+\cdots\right),
\]
and hence
\begin{align}
 |f(\alpha+h)|\le{}& |f(\alpha)|(1-\delta)\notag\\
 &+\delta |f(\alpha)|\left|\frac{h}{m+1}\frac{f^{(m+1)}(\alpha)}{f^{(m)}(\alpha)}
 +\frac{h^2}{(m+1)(m+2)}\frac{f^{(m+2)}(\alpha)}{f^{(m)}(\alpha)}+\cdots\right|.
\tag{17}
\end{align}
By (13), therefore,
\begin{equation}
 |f(\alpha+h)|<|f(\alpha)|.\tag{18}
\end{equation}
\noindent\hfill q.e.d.

Thus the proof of the fundamental theorem is completed.


% ===== scan leaf 152, printed 127 =====
% !PRINTED 127
\weberpage{127}
\pdfbookmark[2]{§ 39. Algorithm for computing the roots}{bm44}

\begin{center}
{\large \S.\ 39.}

\vspace{0.5em}
{\large\gesperrt{Algorithm for the Calculation of the Roots.}}
\end{center}

The proof given in the preceding paragraph for the existence of a root of an algebraic equation, although it leaves nothing to be desired in conclusiveness, nevertheless still has the defect that it does not make clear the steps by which a root can be calculated by a convergent computational procedure. We therefore add the following considerations, intended to remedy this defect, even if only theoretically. They also give a second proof of the fundamental theorem, due in its essentials to \gesperrt{Lipschitz} (\emph{Lehrbuch der Analysis}, vol. I, \S\,61 ff.; I owe a simplification to a communication of \gesperrt{Dedekind}).

If the two integral rational functions $f(x)$ and $f'(x)$ have a common divisor, then this can be found and removed by rational operations according to \S\,6. We may therefore assume that $f(x)$ and $f'(x)$ have no common divisor, that is, that they do not vanish simultaneously for any value of $x$.

We now assume that the calculation of the roots has already succeeded for a function of degree $(n-1)$, and accordingly suppose $f'(x)$ resolved into linear factors. Thus, if
\begin{equation}
 f(x)=x^n+a_1x^{n-1}+a_2x^{n-2}+\cdots,
\tag{1}
\end{equation}
then
\begin{equation}
 f'(x)=nx^{n-1}+(n-1)a_1x^{n-2}+\cdots
      =n(x-\beta_1)(x-\beta_2)\cdots(x-\beta_{n-1}),
\tag{2}
\end{equation}
where $\beta_1,\beta_2,\ldots,\beta_{n-1}$ are to be regarded as known numbers, some of which may also be identical. By our assumption, the absolute values
\begin{equation}
 |f(\beta_1)|,\ |f(\beta_2)|,\ldots, |f(\beta_{n-1})|,
\tag{3}
\end{equation}
which we denote by
\begin{equation}
 b_1,b_2,\ldots,b_{n-1}
\tag{4}
\end{equation}
are all different from zero; let the notation be so chosen that $b_1$ is the smallest among them,
%CONT>>


% ===== scan leaf 153, printed 128 =====
% !PRINTED 128
\weberpage{128}
%<<CONT

or at least exceeds none of the others. By theorem 5 of \S\,38 a value $\alpha$ of $x$ can then be determined so that the absolute value $a$ of $f(\alpha)$ is smaller than $b_1$, hence
\begin{equation}
 a<b_1\le b_2\le b_3\le\cdots\le b_{n-1}.
\tag{5}
\end{equation}

We now bound a domain $G$ for the variable $x$ in such a way that \gesperrt{outside} this domain the absolute value of $f(x)$ is always \gesperrt{greater} than $a$, so that $G$ contains all points $x$ for which $|f(x)|<a$ (but also other points). By \S\,31, IV this is possible in the following way. About the origin as centre draw a circle $(R)$ of sufficiently large radius $R$ containing the points $\alpha,\beta_1,\beta_2,
\ldots,\beta_{n-1}$, and cut out of this circle the points $\beta_1,\beta_2,
\ldots,\beta_{n-1}$, at which $|f(x)|$ is certainly greater than $a$, by circular envelopes of sufficiently small but non-zero radii $\varrho_1,\varrho_2,
\ldots,\varrho_{n-1}$, so that within all these circles $(\varrho_1),(\varrho_2),\ldots,(\varrho_{n-1})$ the absolute value $|f(x)|$ remains greater than $a$ (\S\,31, V). None of these circles then encloses the point $\alpha$, for at that point $|f(x)|=a$. The connected piece of surface bounded by the circles $(R),(\varrho_1),\ldots,(\varrho_{n-1})$ is the domain $G$.

% Exact 300-dpi crop from the registered Volume-I scan, printed Fig. 3.
\weberfigureimage{weber-b1-fig03.png}{0.58}

We now determine a sequence of numbers
\[
 \alpha,\alpha',\alpha'',\alpha''',\ldots
\]
in such a way that the absolute values of $f(\alpha),f(\alpha'),f(\alpha''),\ldots$, which we denote by
\[
 a,a',a'',a''',\ldots,
\]
always decrease. The points $\alpha,\alpha',\alpha''$ so determined all remain in the interior of the domain $G$, since for them $|f(x)|<a$.

For this we use a procedure quite similar to that used for the proof of theorem 5 in \S\,38, but simplified by the fact that $f'(\alpha)$ is already different from zero. If $\delta$ denotes a positive proper fraction, still to be determined more precisely, we put in the expansion
%CONT>>


% ===== scan leaf 154, printed 129 =====
% !PRINTED 129
\weberpage{129}
%<<CONT

\begin{equation}
 f(\alpha+h)=f(\alpha)+h f'(\alpha)+\frac{h^2}{1\cdot2}f''(\alpha)+\cdots
\tag{6}
\end{equation}
\begin{equation}
 h=-\delta\frac{f(\alpha)}{f'(\alpha)},
\tag{7}
\end{equation}
and obtain
\begin{align}
 f(\alpha+h)={}&(1-\delta)f(\alpha)
 +\delta^2\left\{\frac{f(\alpha)^2 f''(\alpha)}{1\cdot2\,f'(\alpha)^2}
 -\delta\frac{f(\alpha)^3 f'''(\alpha)}{1\cdot2\cdot3\,f'(\alpha)^3}+\cdots\right\}.
\tag{8}
\end{align}

Now by (2)
\[
 f'(\alpha)=n(\alpha-\beta_1)(\alpha-\beta_2)\cdots(\alpha-\beta_{n-1});
\]
therefore, if the absolute values of $(\alpha-\beta_1),(\alpha-\beta_2),
\ldots,(\alpha-\beta_{n-1})$ are denoted by $r_1,r_2,
\ldots,r_{n-1}$,
\[
 |f'(\alpha)|=nr_1r_2\cdots r_{n-1}.
\]
Since $\alpha$ lies outside the circle $\varrho_1$ drawn about $\beta_1$, we have $r_1>\varrho_1,r_2>\varrho_2,\ldots$, and consequently
\[
 |f'(\alpha)|>n\varrho_1\varrho_2\cdots\varrho_{n-1}.
\]
Thus $|f'(\alpha)|$ is greater than a positive number $k$ independent of the position of $\alpha$ in $G$. It follows that one can choose a sufficiently large positive number $Q$, likewise independent of the position of $\alpha$ and of the proper fraction $\delta$, so that
\begin{equation}
 \left|\frac{f(\alpha)^2f''(\alpha)}{1\cdot2\,f'(\alpha)^2}
 -\delta\frac{f(\alpha)^3f'''(\alpha)}{1\cdot2\cdot3\,f'(\alpha)^3}+\cdots\right|<Q .
\tag{9}
\end{equation}

One can obtain $Q$ by replacing in the denominators on the left of (9) $f'(\alpha)$ by $k$, replacing $\delta$ by $1$, replacing all the other terms by their absolute values, replacing finally the absolute value of $\alpha$ by the greater value $R$, and then making the number so obtained arbitrarily larger. From (8) one then obtains
\begin{equation}
 |f(\alpha+h)|<(1-\delta)|f(\alpha)|+\delta^2Q.
\tag{10}
\end{equation}
Assume now $Q$ so chosen that, for every $\alpha$ in $G$,
\[
 2Q>|f(\alpha)|=a
\]
which is plainly compatible with the above determination. We may then put
\[
 \delta=\frac{a}{2Q},\qquad
 h=-\frac{a f(\alpha)}{2Q f'(\alpha)}
\]
%CONT>>


% ===== scan leaf 155, printed 130 =====
% !PRINTED 130
\weberpage{130}
%<<CONT

and, denoting $\alpha+h$ by $\alpha'$,
\begin{equation}
 \alpha'=\alpha-\frac{a f(\alpha)}{2Q f'(\alpha)},
\tag{11}
\end{equation}
and, denoting $|f(\alpha')|$ by $a'$, obtain from (10)
\begin{equation}
 a'<a\left(1-\frac{a}{4Q}\right).
\tag{12}
\end{equation}

By the same procedure we derive from $\alpha',a'$ a second system $\alpha'',a''$, then from $\alpha'',a''$ a third system $\alpha''',a'''$, and so on; for the sequence of decreasing positive numbers $a,a',a'',\ldots$ we obtain the inequalities
\begin{equation}
\begin{aligned}
 a'&<a\left(1-\frac{a}{4Q}\right),\\
 a''&<a'\left(1-\frac{a'}{4Q}\right),\\
 a'''&<a''\left(1-\frac{a''}{4Q}\right),\\
 &\hspace{1.5em}\vdots
\end{aligned}
\tag{13}
\end{equation}
We now prove that this sequence approaches the limit zero. If this were not the case, a positive number $\omega$ could be given such that all $a^{(\nu)}$ remain above $\omega$; then the differences
\[
 \left(1-\frac{a}{4Q}\right),\quad
 \left(1-\frac{a'}{4Q}\right),\quad
 \left(1-\frac{a''}{4Q}\right)\cdots
\]
would all be smaller than the positive proper fraction
\[
 \Theta=1-\frac{\omega}{4Q},
\]
and from (13) it would follow that
\[
 a'<a\Theta,\quad
 a''<a'\Theta,\quad
 a'''<a''\Theta,\quad
 \ldots,\quad
 a^{(\nu)}<a^{(\nu-1)}\Theta.
\]
Whence, by multiplication,
\begin{equation}
 a^{(\nu)}<a\Theta^\nu.
\tag{14}
\end{equation}
This is a contradiction, since $\Theta^\nu$ becomes infinitely small as $\nu$ increases without bound. The result so obtained is expressed by
\begin{equation}
 \lim a^{(\nu)}=0.
\tag{15}
\end{equation}


% ===== scan leaf 156, printed 131 =====
% !PRINTED 131
\weberpage{131}

If the inequalities (13) are written in the form
\[
\begin{aligned}
 a-a'&>\frac{a^2}{4Q},\\
 a'-a''&>\frac{a'^2}{4Q},\\
 &\hspace{1.5em}\vdots\\
 a^{(\nu-1)}-a^{(\nu)}&>\frac{(a^{(\nu-1)})^2}{4Q},
\end{aligned}
\]
and added from the $(\mu+1)$-st to the $(\nu+1)$-st, one obtains
\[
 (a^{(\mu)})^2+(a^{(\mu+1)})^2+\cdots+(a^{(\nu)})^2
 <4Q\bigl(a^{(\mu)}-a^{(\nu+1)}\bigr),
\]
and therefore, when $\mu$ and $\nu$ both tend to infinity,
\begin{equation}
 \lim\bigl((a^{(\mu)})^2+(a^{(\mu+1)})^2+\cdots+(a^{(\nu)})^2\bigr)=0.
\tag{16}
\end{equation}

For the sequence of numbers $\alpha,\alpha',\alpha'',\ldots$ the law of formation resulting from (11) is
\begin{equation}
\begin{aligned}
 \alpha-\alpha'&=\frac{a f(\alpha)}{2Q f'(\alpha)},\\
 \alpha'-\alpha''&=\frac{a' f(\alpha')}{2Q f'(\alpha')},\\
 &\hspace{1.5em}\vdots\\
 \alpha^{(\nu-1)}-\alpha^{(\nu)}&=
 \frac{a^{(\nu-1)} f(\alpha^{(\nu-1)})}{2Q f'(\alpha^{(\nu-1)})}.
\end{aligned}
\tag{17}
\end{equation}
If an arbitrary number of these equations are added, from the $(\mu+1)$-st to the $(\nu+1)$-st, then
\begin{equation}
 \alpha^{(\mu)}-\alpha^{(\nu+1)}=
 \frac{a^{(\mu)}f(\alpha^{(\mu)})}{2Q f'(\alpha^{(\mu)})}+\cdots+
 \frac{a^{(\nu)}f(\alpha^{(\nu)})}{2Q f'(\alpha^{(\nu)})}.
\tag{18}
\end{equation}
Since the inequalities
\[
 |f'(\alpha)|>k
\]
are satisfied, it follows for the absolute value of this difference that
\begin{equation}
 |\alpha^{(\mu)}-\alpha^{(\nu+1)}|<
 \frac{(a^{(\mu)})^2+(a^{(\mu+1)})^2+\cdots+(a^{(\nu)})^2}{2Qk},
\tag{19}
\end{equation}
and hence, by (16),
\begin{equation}
 \lim |\alpha^{(\mu)}-\alpha^{(\nu+1)}|=0.
\tag{20}
\end{equation}

This says (according to the concept of a sequence of numbers explained in the Introduction, p. 15) that the numbers $\alpha,\alpha',\alpha'',\ldots$ approach a definite limit, which we shall denote by $\xi$.
%CONT>>


% ===== scan leaf 157, printed 132 =====
% !PRINTED 132
\weberpage{132}
\pdfbookmark[2]{§ 40. Continuity of the roots}{bm45}
%<<CONT

From the continuity of the function $f(x)$ it then follows easily that $f(\xi)=0$; for if $|f(\xi)|>0$, then, since the $\alpha^{(\nu)}$ come arbitrarily close to $\xi$, the absolute values $a^{(\nu)}$ of $f(\alpha^{(\nu)})$ could not sink below every positive value, which is precisely what we have proved for them.


\begin{center}
{\large \S.\ 40.}

\vspace{0.5em}
{\large\gesperrt{Continuity of the Roots.}}
\end{center}

We close this section with the proof of the theorem:

\begin{center}
\gesperrt{The roots of an algebraic equation are continuous functions of the coefficients.}
\end{center}

We have first to explain the meaning of this theorem.

Let
\begin{equation}
 f(x)=x^n+a_1x^{n-1}+a_2x^{n-2}+\cdots
\tag{1}
\end{equation}
be an integral rational function of $x$ of degree $n$. By what has been proved in the preceding paragraphs, $f(x)$ can be decomposed into $n$ linear factors, some of which may be equal to one another. Combining equal factors and denoting the roots by $\alpha,\beta,\gamma,\ldots$, we put
\begin{equation}
 f(x)=(x-\alpha)^a(x-\beta)^b(x-\gamma)^c\cdots,
\tag{2}
\end{equation}
where $a,b,c,\ldots$ are positive whole numbers whose sum is equal to $n$.

The roots $\alpha,\beta,\gamma,\ldots$ will change with the coefficients $a_1,a_2,a_3,\ldots$; the degree of their multiplicity can also become another one.

We denote the changes of $a_1,a_2,
\ldots$ by $\varepsilon_1,\varepsilon_2,\ldots$ and put
\begin{align}
 \varphi(x)&=\varepsilon_1x^{n-1}+\varepsilon_2x^{n-2}+\cdots,
\tag{3}\\
 f(x)+\varphi(x)&=f_1(x).
\tag{4}
\end{align}
We surround the points $\alpha,\beta,\gamma,
\ldots$ by regions of arbitrary smallness, but in such a way that these regions exclude one another; for instance, we enclose the points $\alpha,\beta,\gamma,
\ldots$ by circular peripheries with radii $\varrho,\varrho',\varrho'',\ldots$, and denote these regions by $(\varrho),(\varrho'),(\varrho''),\ldots$.


% ===== scan leaf 158, printed 133 =====
% !PRINTED 133
\weberpage{133}

If the absolute values of $\varepsilon_1,\varepsilon_2,
\ldots$ lie below sufficiently small values, then we can prove first of the function $f_1(x)$
\begin{center}
\gesperrt{that it has no roots outside the regions $(\varrho),(\varrho'),(\varrho''),\ldots$,}
\end{center}
and secondly,
\begin{center}
\gesperrt{that the number of roots of $f_1(x)$ within $(\varrho)$ is exactly $a$, within $(\varrho')$ exactly $b$, within $(\varrho'')$ exactly $c$, and so on.}
\end{center}
In the last part of the theorem, however, it must be observed that, if $f_1(x)$ has multiple roots, these must be counted according to their multiplicity.

We construct, according to \S 31, IV, in the $x$-plane a circle with radius $R$ and with the origin as centre, which encloses the regions $(\varrho),(\varrho'),(\varrho''),\ldots$, so that outside this circle no more roots of $f_1(x)$ lie, and we denote by $G$ the region lying within this circle, but outside $(\varrho),(\varrho'),(\varrho''),\ldots$. It must also be remarked that $R$ can be assumed independently of $\varepsilon_1,\varepsilon_2,\ldots$, so long as a definite upper bound is fixed for the absolute values of these quantities.

% Exact unretouched crop from the registered Volume-I scan, printed Fig. 4.
\weberfigureimage{weber-b1-fig04.png}{0.42}

If now $x$ is a point of the region $G$, then the absolute value of $x-\alpha$ is greater than $\varrho$, and consequently, by (2),
\begin{equation}
 |f(x)|>\varrho^a\varrho'^b\varrho''^c\cdots .
\tag{5}
\end{equation}
If now $\alpha_1$ is a root of $f_1(x)$, then from (4) it follows that
\begin{equation}
 f(\alpha_1)=-\varphi(\alpha_1),
\tag{6}
\end{equation}
and $\alpha_1$ therefore cannot lie in the region $G$ if one takes care that everywhere within $G$
\begin{equation}
 |\varphi(x)|<\varrho^a\varrho'^b\varrho''^c\cdots,
\tag{7}
\end{equation}
which is always possible by sufficient diminution of the upper bound of $\varepsilon_1,\varepsilon_2,
\ldots$. Thus the first part of our assertion is proved, namely that, when the inequality (7) is satisfied, no roots of $f_1(x)$ lie outside the regions $(\varrho),(\varrho'),(\varrho''),\ldots$.


% ===== scan leaf 159, printed 134 =====
% !PRINTED 134
\weberpage{134}

To prove the second part, we put
\begin{align}
 f(x)&=\psi(x)(x-\alpha)^a,
\tag{8}\\
 \psi(x)&=(x-\beta)^b(x-\gamma)^c\cdots .
\tag{9}
\end{align}
We now choose a positive number $A$, which shall however be independent of $\varrho$, in such a way that, as long as $x$ lies in the interior or on the periphery of the region $(\varrho)$,
\begin{equation}
 |\psi(x)|>A \qquad (\text{for } |x-\alpha|\le \varrho)
\tag{10}
\end{equation}
remains true. Such a number is obtained, for example, if $l$ is taken equal to half the least of the distances $(\alpha,\beta),(\alpha,\gamma),\ldots$ and one puts
\[
 A=l^{b+c+\cdots};
\]
then the requirement expressed in (10) is fulfilled at least as long as $\varrho$ is smaller than $l$. If only one point $\alpha$ is present, then $\psi(x)=1$ is to be put, and for $A$ any proper fraction can be chosen.

Let now $\alpha_1,\alpha_2,
\ldots$ be the roots of $f_1(x)$ within $(\varrho)$, and let $a_1,a_2,
\ldots$ be the degrees of their multiplicity; and let $\beta_1,\beta_2,
\ldots$, $b_1,b_2,
\ldots$, $\gamma_1,\gamma_2,
\ldots$, $c_1,c_2,
\ldots$ have the same meaning for the regions $(\varrho'),(\varrho''),\ldots$. Then
\begin{equation}
 n=a_1+a_2+\cdots+b_1+b_2+\cdots+c_1+c_2+\cdots,
\tag{11}
\end{equation}
since the total number of all roots must be equal to $n$. The function $f_1(x)$ can then be represented as
\begin{equation}
 f_1(x)=\psi_1(x)(x-\alpha_1)^{a_1}(x-\alpha_2)^{a_2}\cdots,
\tag{12}
\end{equation}
where
\begin{equation}
 \psi_1(x)=(x-\beta_1)^{b_1}(x-\beta_2)^{b_2}\cdots
          (x-\gamma_1)^{c_1}(x-\gamma_2)^{c_2}\cdots .
\tag{13}
\end{equation}

Now we can determine a positive number $B$ independent of $\varrho,\varrho',\varrho''$ such that, as long as $x$ remains within or on the boundary of $(\varrho)$,
\begin{equation}
 |\psi_1(x)|<B \qquad (\text{for } |x-\alpha|\le \varrho).
\tag{14}
\end{equation}
For example, one can choose a magnitude $L$ which is greater than twice the distance of the point $\alpha$ from any of the points $\beta,\gamma,\ldots$ and in every case greater than $1$, and then assume
\[
 B>L^n .
\]
Now from (4) it follows that
% [sic] The source prints a strict triangle inequality in (15); equality is possible, so the universally valid relation is <=. The later strict bounds still suffice for the argument.
\begin{equation}
 |f(x)|<|f_1(x)|+|\varphi(x)|,
\tag{15}
\end{equation}
%CONT>>


% ===== scan leaf 160, printed 135 =====
% !PRINTED 135
\weberpage{135}
%<<CONT

hence by (8) and (12)
\begin{equation}
 |\psi(x)|\,|x-\alpha|^a<|\psi_1(x)|\,|x-\alpha_1|^{a_1}\,|x-\alpha_2|^{a_2}\cdots+|\varphi(x)|.
\tag{16}
\end{equation}
If we now assume $x$ on the periphery of $\varrho$, and therefore put
\[
 |x-\alpha|=\varrho,
\]
then
\[
 |x-\alpha_1|<2\varrho,\qquad |x-\alpha_2|<2\varrho,\ldots,
\]
and consequently, by (10) and (16),
\begin{equation}
 \varrho^a A < (2\varrho)^{a_1+a_2+\cdots}B+|\varphi(x)|.
\tag{17}
\end{equation}
We now choose the upper bounds for the coefficients $\varepsilon_1,\varepsilon_2,
\ldots$ of $\varphi$ so small that
\[
 |\varphi(x)|<\varrho^aA'
\]
where $A'$ denotes a number which is smaller than $A$; then from (17) it follows that\begingroup\renewcommand{\thefootnote}{E1}\footnote[1]{[Source note:] Set $s=a_1+a_2+\cdots$. From (17) one obtains $A-A'<2^s\varrho^{s-a}B$. The contradiction concerns $s>a$, not $a>s$: for $0<\varrho<1$, one has $2^s\varrho^{s-a}B\le 2^nB\varrho\to0$. Thus the signs in (19), (20) should be $\ge$; summing still gives equality. The source reference (20), (21) means (19), (20). Printed text and formulas are retained.}\endgroup
\begin{equation}
 A-A'<2^{a}\varrho^{-a+a_1+a_2+\cdots}B.
\tag{18}
\end{equation}
But, for sufficiently small $\varrho$, this would no longer be possible if $a$ were greater than the sum $a_1+a_2+\cdots$. Hence
\begin{equation}
 a\le a_1+a_2+\cdots .
\tag{19}
\end{equation}
In the same way one proves that
\begin{equation}
 b\le b_1+b_2+\cdots,
 \qquad
 c\le c_1+c_2+\cdots,
 \qquad \ldots .
\tag{20}
\end{equation}
But since the sums of the left sides as well as of the right sides in the inequalities (20), (21) must be equal to $n$,
% [sic] Source-original mathematical cross-reference: the printed page says (20), (21), although the displayed inequalities are numbered (19), (20).
only equality signs can hold, and therefore
\[
 a=a_1+a_2+\cdots,
 \quad b=b_1+b_2+\cdots,
 \quad c=c_1+c_2+\cdots,
 \quad \ldots,
\]
whereby the second part of our theorem is also proved.

We shall add to the theorem just proved the following remark concerning the case of multiple roots.

The necessary and sufficient condition that $f(x)=0$ should have multiple roots at all is that $f(x)$ and $f'(x)$ should have a common factor. Therefore, if the algorithm
%CONT>>


% ===== scan leaf 161, printed 136 =====
% !PRINTED 136
\weberpage{136}
%<<CONT

of the greatest common divisor is applied to these two functions according to the principles developed in the first section, the condition for common factors is obtained by setting the last constant remainder equal to zero, in the form of an equation between the coefficients,
\[
 F(a_1,a_2,\ldots,a_n)=0,
\]
where $F$ is an integral rational function of the arguments $a_1,a_2,
\ldots,a_n$, which is called the \gesperrt{discriminant} of $f$, and whose properties and mode of formation will be treated more fully in the next section. In any case, since equations of degree $n$ without multiple roots do exist at all, $F$ cannot vanish identically.

Consider now
\[
 F(a_1+\varepsilon_1,a_2+\varepsilon_2,
\ldots,a_n+\varepsilon_n).
\]
This function too, if $a_1,a_2,
\ldots,a_n$ are definite and $\varepsilon_1,\varepsilon_2,
\ldots,\varepsilon_n$ are indeterminate numbers, will not vanish identically; and one can choose the $\varepsilon$'s so that its value comes out different from zero, even when any arbitrary upper bound has been fixed for the absolute values of the $\varepsilon$'s.\footnote[1]{This will only later be rigorously justified. For the time being we presuppose it in these preliminary remarks, on which no further conclusions will be founded.}

One therefore sees from these considerations how an $a$-fold root $\alpha$ of $f(x)$, under continuous variation of the coefficients, radiates apart into $a$ simple roots, so that from this point of view too one is justified in considering the $a$-fold root as having arisen from the coincidence of $a$ simple roots.

If we do not assume the coefficient of the highest power of $x$ equal to $1$, and therefore examine the equation in the form
\begin{equation}
 a_0x^n+a_1x^{n-1}+a_2x^{n-2}+\cdots+a_{n-1}x+a_n=0,
\tag{21}
\end{equation}
then, if $a_n=0$, one or more of the roots are equal to zero, and equation (21) can be divided by $x$ or by a higher power of $x$, whereby an equation of correspondingly lower degree arises in which the last term is different from zero. We shall therefore assume from the outset that $a_n$ is different from zero, and now put in (21)
\begin{equation}
 x=\frac1y.
\tag{22}
\end{equation}
%CONT>>


% ===== scan leaf 162, printed 137 =====
% !PRINTED 137
\weberpage{137}
%<<CONT

By multiplication by $y^n$ and division by $a_n$, (21) then passes into
\begin{equation}
 y^n+\frac{a_{n-1}}{a_n}y^{n-1}+\cdots+\frac{a_1}{a_n}y+\frac{a_0}{a_n}=0,
\tag{23}
\end{equation}
and we can apply to the roots of this equation the theorem just stated, by assuming $a_0=0$, so that one of the roots of (23) vanishes. We thus obtain the theorem:

One can take $a_0$ in the equation (21) so small that one of its roots becomes large beyond all bounds, while the others differ arbitrarily little from the roots of the equation of degree $(n-1)$
\[
 a_1x^{n-1}+a_2x^{n-2}+\cdots+a_{n-1}x+a_n=0.
\]
This is also expressed by saying that, \gesperrt{as $a_0$ vanishes, one of the roots of (21) becomes infinite.}

If $a_1$ also vanishes, or $a_1$ and $a_2$, and so on, two or three roots become infinite.

This justifies the expression sometimes used, that an equation of degree $n$ passes, by the becoming-infinite of one root, into an equation of degree $(n-1)$.

\begin{center}
\rule{0.16\textwidth}{0.4pt}
\end{center}


% ===== scan leaf 163, printed 138 =====
% !PRINTED 138
\weberpage{138}
\pdfbookmark[1]{Fourth Section. Symmetric functions.}{bm46}
\pdfbookmark[2]{§ 41. Concept of the symmetric functions. Fundamental symmetric functions}{bm47}

\begin{center}
{\Large Fourth Section.}\\[0.8em]
{\Large\gesperrt{Symmetric Functions.}}\\[0.8em]
\rule{0.16\textwidth}{0.4pt}
\end{center}

\begin{center}
{\large \S.\ 41.}

\vspace{0.5em}
{\large\gesperrt{Concept of Symmetric Functions.}}

{\large\gesperrt{Symmetric Fundamental Functions.}}
\end{center}

In this section we consider \gesperrt{whole functions}, of arbitrary degree, in an arbitrary number $n$ of variables
\[
 \alpha_1,\alpha_2,\ldots,\alpha_n.
\]

Such a function is called \gesperrt{symmetric} \gesperrt{if it remains unchanged when the variables $\alpha_1,\alpha_2,\ldots,\alpha_n$ are subjected to an arbitrary permutation;} and it is such symmetric functions whose properties and laws of formation we now have to learn more precisely.

In order that a function $\Phi(\alpha_1,\alpha_2,\ldots,\alpha_n)$ be symmetric, it is sufficient that it not be changed when any two of the arguments $\alpha_1,\alpha_2,\ldots,\alpha_n$ are interchanged, since by \S\,18 all permutations can be formed by a succession of transpositions.

The function $\Phi$ will in general not be homogeneous, but will contain terms of different dimensions. If, however, one collects all terms of the same dimension, every $\Phi$ can be represented as a sum of \gesperrt{homogeneous functions} of different degrees; and, if $\Phi$ is to be symmetric, then each homogeneous constituent of a definite degree must by itself be symmetric, since the dimensions of the terms are not changed by the permutations of the variables.


% ===== scan leaf 164, printed 139 =====
% !PRINTED 139
\weberpage{139}

Accordingly we may confine ourselves to the consideration of \gesperrt{homogeneous symmetric functions}, from which all the others can be put together.

After this it is easy to give the general form of a symmetric function. We obtain it if, in one term of such a function,
\[
 \alpha_1^{\nu_1}\alpha_2^{\nu_2}\cdots\alpha_n^{\nu_n},
\]
as in the formation of determinants, we permute the lower indices in all possible ways and take the sum of all terms so formed. Such a function may be called an \gesperrt{elementary constituent} of a symmetric function. If we take several such elementary constituents, multiply them by arbitrary factors independent of the $\alpha$'s, and add them, then we obtain the most general symmetric function. The number of terms of one of these elementary constituents is $\Pi(n)$ when the exponents $\nu_1,\nu_2,\ldots,\nu_n$ are all different; but if one and the same exponent $\nu$ occurs several times, the permutations which give no different terms are to be omitted. For example, for three variables, if $\nu_1,\nu_2,\nu_3$ are different, an elementary constituent is
\[
\begin{aligned}
 &\alpha_1^{\nu_1}\alpha_2^{\nu_2}\alpha_3^{\nu_3}
+\alpha_1^{\nu_1}\alpha_3^{\nu_2}\alpha_2^{\nu_3}
+\alpha_2^{\nu_1}\alpha_1^{\nu_2}\alpha_3^{\nu_3}
+\alpha_2^{\nu_1}\alpha_3^{\nu_2}\alpha_1^{\nu_3}\\
 &\qquad+\alpha_3^{\nu_1}\alpha_1^{\nu_2}\alpha_2^{\nu_3}
+\alpha_3^{\nu_1}\alpha_2^{\nu_2}\alpha_1^{\nu_3},
\end{aligned}
\]
whereas, if $\nu_3=\nu_2$,
\[
 \alpha_1^{\nu_1}\alpha_2^{\nu_2}\alpha_3^{\nu_2}
+\alpha_2^{\nu_1}\alpha_1^{\nu_2}\alpha_3^{\nu_2}
+\alpha_3^{\nu_1}\alpha_1^{\nu_2}\alpha_2^{\nu_2}.
\]

The simplest example of a symmetric function is the sum of the variables
\[
 \alpha_1+\alpha_2+\alpha_3+\cdots+\alpha_n.
\]
Likewise the product $\alpha_1\alpha_2\cdots\alpha_n$ belongs here.

These two are the extreme cases of a series of symmetric functions which we call the \gesperrt{symmetric fundamental functions} and obtain as follows.

The product
\begin{equation}
 f(x)=(x-\alpha_1)(x-\alpha_2)\cdots(x-\alpha_n)
\tag{1}
\end{equation}
is, whatever $x$ may be, provided only that $x$ is independent of the $\alpha$'s, a symmetric function of $\alpha_1,\alpha_2,\ldots,\alpha_n$. If therefore we carry out the multiplication of the several factors and
%CONT>>


% ===== scan leaf 165, printed 140 =====
% !PRINTED 140
\weberpage{140}
\pdfbookmark[2]{§ 42. The power sums}{bm48}
%<<CONT

arrange according to powers of $x$, then the coefficients of the separate powers of $x$ are likewise symmetric functions. For they change no more under interchange of the $\alpha$'s than does the function $f(x)$ (\S\,1). We put
\begin{equation}
 f(x)=x^n+a_1x^{n-1}+a_2x^{n-2}+\cdots+a_{n-1}x+a_n.
\tag{2}
\end{equation}

The functions $a_1,a_2,\ldots,a_n$ are the coefficients of the algebraic equation whose roots are $\alpha_1,\alpha_2,\ldots,\alpha_n$. They are called the \gesperrt{symmetric fundamental functions}. We already formed these symmetric fundamental functions in \S\,7; as we recall, $a_1$ is the sum of the $\alpha$'s taken with negative sign, $a_2$ the sum of the products two at a time, $a_3$ the sum of the products three at a time taken with negative sign, and so on; finally $a_n$ is the product of all the $\alpha$'s with positive or negative sign according as $x$ is even or odd.
% [sic] The registered source prints x here; the parity determining the sign is that of the degree n.
We write, for short,
\begin{equation}
\begin{aligned}
 a_1&=-\sum \alpha_1,\\
 a_2&=+\sum \alpha_1\alpha_2,\\
 a_3&=-\sum \alpha_1\alpha_2\alpha_3,\\
 &\ \vdots\\
 a_n&=\pm \alpha_1\alpha_2\cdots\alpha_n.
\end{aligned}
\tag{3}
\end{equation}

The goal of our considerations is the proof of the principal theorem \gesperrt{that all symmetric functions of the $\alpha$'s can be expressed rationally by the fundamental functions.}

\begin{center}
{\large \S.\ 42.}

\vspace{0.5em}
{\large\gesperrt{The Power Sums.}}
\end{center}

We occupy ourselves first with another special kind of symmetric functions, the \gesperrt{power sums}. Namely, if $\nu$ denotes any positive integral exponent, then
\begin{equation}
 s_\nu=\alpha_1^\nu+\alpha_2^\nu+\cdots+\alpha_n^\nu
\tag{1}
\end{equation}
obviously belongs to the symmetric functions, and $s_\nu$ is called the $\nu$th power sum. We shall derive formulae according to which the power sums are expressible by the fundamental functions.

In what follows, whenever $\varphi(x)$ is any function of $x$, we always denote by
\[
 S[\varphi(\alpha)]
\]
%CONT>>


% ===== scan leaf 166, printed 141 =====
% !PRINTED 141
\weberpage{141}
%<<CONT

the sum obtained by replacing $x$ in $\varphi(x)$ successively by each of the variables $\alpha_1,\alpha_2,\ldots,\alpha_n$ and adding the functions so formed, thus
\[
 S[\varphi(\alpha)]=\varphi(\alpha_1)+\varphi(\alpha_2)+\cdots+\varphi(\alpha_n).
\]
Accordingly, for example,
\[
 S(\alpha^\nu)=s_\nu
\]
is the $\nu$th power sum.

We now divide $f(x)$, by \S\,4, by an arbitrary linear function $x-\alpha$ and obtain
\begin{equation}
 \frac{f(x)}{x-\alpha}
   =x^{n-1}+f_1(\alpha)x^{n-2}+f_2(\alpha)x^{n-3}+\cdots+f_{n-1}(\alpha)+\frac{f(\alpha)}{x-\alpha},
\tag{2}
\end{equation}
where, by \S\,4, (8),
\begin{equation}
\begin{aligned}
 f_1(\alpha)&=\alpha+a_1,\\
 f_2(\alpha)&=\alpha^2+a_1\alpha+a_2,\\
 f_3(\alpha)&=\alpha^3+a_1\alpha^2+a_2\alpha+a_3,\\
 &\ \vdots\\
 f_{n-1}(\alpha)&=\alpha^{n-1}+a_1\alpha^{n-2}+a_2\alpha^{n-3}+\cdots+a_{n-1}.
\end{aligned}
\tag{3}
\end{equation}

We put in (2), for $\alpha$, each of the quantities $\alpha_1,\alpha_2,\ldots,\alpha_n$, whereby $f(\alpha)=0$, and form the sum $S$. The left hand side then gives, by \S\,14, (11),
\begin{equation}
\begin{aligned}
 S\left(\frac{f(x)}{x-\alpha}\right)&=f'(x)\\
 &=n x^{n-1}+(n-1)a_1x^{n-2}+(n-2)a_2x^{n-3}+\cdots+a_{n-1},
\end{aligned}
\tag{4}
\end{equation}
whereas the right hand side becomes
\begin{equation}
 n x^{n-1}+x^{n-2}S[f_1(\alpha)]+x^{n-3}S[f_2(\alpha)]+\cdots+S[f_{n-1}(\alpha)].
\tag{5}
\end{equation}
The comparison of coefficients of like powers of $x$ in (4) and (5) gives
\begin{equation}
\begin{aligned}
 S[f_1(\alpha)]&=(n-1)a_1,\\
 S[f_2(\alpha)]&=(n-2)a_2,\\
 &\ \vdots\\
 S[f_{n-1}(\alpha)]&=a_{n-1}.
\end{aligned}
\tag{6}
\end{equation}
But by (1) and (3)
\[
\begin{aligned}
 S[f_1(\alpha)]&=s_1+na_1,\\
 S[f_2(\alpha)]&=s_2+a_1s_1+na_2,\\
 &\ \vdots\\
 S[f_{n-1}(\alpha)]&=s_{n-1}+a_1s_{n-2}+a_2s_{n-3}+\cdots+na_{n-1}.
\end{aligned}
\]
%CONT>>


% ===== scan leaf 167, printed 142 =====
% !PRINTED 142
\weberpage{142}
%<<CONT

Consequently from (6) one obtains the following system of formulae due to \gesperrt{Newton}\footnote[1]{\gesperrt{Newton}, Arithmetica universalis, edit. 's Gravesande, p.~592.}:
\begin{equation}
\begin{aligned}
0&=s_1+a_1,\\
0&=s_2+a_1s_1+2a_2,\\
0&=s_3+a_1s_2+a_2s_1+3a_3,\\
&\ \vdots\\
0&=s_{n-1}+a_1s_{n-2}+a_2s_{n-3}+\cdots+(n-1)a_{n-1}.
\end{aligned}
\tag{7}
\end{equation}

This system of formulae can, however, be continued still further; for, since
\[
 S[f(\alpha)]=0,\qquad S[\alpha f(\alpha)]=0,\qquad S[\alpha^2f(\alpha)]=0,\ldots,
\]
it follows that
\begin{equation}
\begin{aligned}
0&=s_n+a_1s_{n-1}+a_2s_{n-2}+\cdots+n a_n,\\
0&=s_{n+1}+a_1s_n+a_2s_{n-1}+\cdots+a_ns_1,\\
0&=s_{n+2}+a_1s_{n+1}+a_2s_n+\cdots+a_ns_2,\quad\ldots
\end{aligned}
\tag{8}
\end{equation}

By the formulae (7), (8) the problem is now solved of representing, one after another, the functions $s_1,s_2,s_3,\ldots$ up to any height as whole rational functions of the symmetric fundamental functions; for example
\begin{equation}
\begin{aligned}
 s_1&=-a_1,\\
 s_2&=+a_1^2-2a_2,\\
 s_3&=-a_1^3+3a_1a_2-3a_3,\\
 s_4&=+a_1^4-4a_1^2a_2+4a_1a_3+2a_2^2-4a_4,\\
 &\hspace{1.5em}\vdots
\end{aligned}
\tag{9}
\end{equation}
The manner of formation of these expressions shows that the coefficients in these representations are \gesperrt{whole numbers}.

Conversely, by means of the formulae (7) and (8), one can express the symmetric fundamental functions $a_1,a_2,a_3,\ldots$ rationally by the power sums $s_1,s_2,s_3,\ldots$. This representation is, however, far less simple in so far as the coefficients are not whole but fractional numbers, for example
\begin{equation}
\begin{aligned}
 a_1&=-s_1,\\
 2a_2&=+s_1^2-s_2,\\
 6a_3&=-s_1^3+3s_1s_2-2s_3,\\
 &\hspace{1.5em}\vdots
\end{aligned}
\tag{10}
\end{equation}


% ===== scan leaf 168, printed 143 =====
% !PRINTED 143
\weberpage{143}
\pdfbookmark[2]{§ 43. Proof of the main theorem for two variables}{bm49}

One can also continue the system of formulae (8) in the opposite direction, if one forms the equations
\[
 S[\alpha^{-1}f(\alpha)]=0,\qquad
 S[\alpha^{-2}f(\alpha)]=0,\ldots .
\]
In this way one obtains a means of expressing the power sums $s_\nu$ also for negative exponents $\nu$ by the fundamental functions. These sums of negative powers likewise belong to the symmetric functions, although no longer to the \gesperrt{whole}, but to the \gesperrt{fractional}, ones. They pass into whole functions only after multiplication by powers of the product $\alpha_1\alpha_2\cdots\alpha_n$. For completeness we put down here the first two of these formulae:
\begin{equation}
\begin{aligned}
0&=s_{n-1}+a_1s_{n-2}+a_2s_{n-3}+\cdots+n a_{n-1}+a_ns_{-1},\\
0&=s_{n-2}+a_1s_{n-3}+a_2s_{n-4}+\cdots+a_{n-1}s_{-1}+a_ns_{-2}.
\end{aligned}
\tag{11}
\end{equation}
which, by (7), can also be written as
\[
 a_{n-1}+a_ns_{-1}=0,\qquad 2a_{n-2}+a_{n-1}s_{-1}+a_ns_{-2}=0.
\]

\begin{center}
{\large \S.\ 43.}

\vspace{0.5em}
{\large\gesperrt{Proof of the Principal Theorem for Two Variables.}}
\end{center}

We now pass to the proof of the fundamental theorem of the theory of symmetric functions, that they are all rationally expressible by the symmetric fundamental functions. Since we use complete induction as the means of proof, we first derive the theorem under the assumption that only two independent variables $\alpha,\beta$ are given, but in a way which at the same time will bring out the guiding thought for the general proof.

We denote the symmetric fundamental functions by
\begin{equation}
 a=-(\alpha+\beta),\qquad b=\alpha\beta
\tag{1}
\end{equation}
and accordingly put
\begin{equation}
 f(x)=(x-\alpha)(x-\beta)=x^2+ax+b.
\tag{2}
\end{equation}

Let $S(\alpha,\beta)$ be any whole rational and symmetric function of $\alpha$ and $\beta$. We can substitute for $\beta$ from (1) the value $-(a+\alpha)$, and obtain, if we arrange according to powers of $\alpha$,
\begin{equation}
 S(\alpha,\beta)=S(\alpha,-a-\alpha)=A_0\alpha^m+A_1\alpha^{m-1}+\cdots+A_{m-1}\alpha+A_m,
\tag{3}
\end{equation}
%CONT>>


% ===== scan leaf 169, printed 144 =====
% !PRINTED 144
\weberpage{144}
\pdfbookmark[2]{§ 44. General proof of the main theorem}{bm50}
%<<CONT

where the coefficients $A_0,A_1,\ldots,A_m$ depend only on $a$ and on the coefficients still occurring in $S$. We note expressly, however, that if in $S(\alpha,\beta)$ no fractional numerical coefficients occur, then no fractions occur in the coefficients $A_0,A_1,\ldots,A_m$ either.

We now put
\begin{equation}
 \Phi(x)=A_0x^m+A_1x^{m-1}+\cdots+A_{m-1}x+A_m
\tag{4}
\end{equation}
and divide $\Phi(x)$ by $f(x)$, by \S\,3. We obtain a quotient $Q$ and a remainder which, with respect to $x$, is at most of the first degree; thus
\begin{equation}
 \Phi(x)=Qf(x)+A+Bx;
\tag{5}
\end{equation}
Here $A$ and $B$ are now whole functions of $a$ and $b$, and they contain no fractional numerical coefficients if none occur in $S$. If now we put $x=\alpha$, it follows from (5) and (3), since $f(\alpha)$ vanishes, that
\begin{equation}
 S(\alpha,\beta)=A+B\alpha.
\tag{6}
\end{equation}
But since $S(\alpha,\beta)$ and likewise $A,B$ are symmetric, it follows by interchanging $\alpha$ and $\beta$ that
\begin{equation}
 S(\alpha,\beta)=A+B\beta.
\tag{7}
\end{equation}
From this one concludes, since $\alpha$ and $\beta$ are independent variables, that $B=0$, and consequently
\begin{equation}
 S(\alpha,\beta)=A,
\tag{8}
\end{equation}
which proves the fundamental theorem in this case.

\begin{center}
{\large \S.\ 44.}

\vspace{0.5em}
{\large\gesperrt{General Proof of the Principal Theorem.}}
\end{center}

We now suppose that the fundamental theorem has been proved for symmetric functions of \(n-1\) variables, and we derive it, by a procedure entirely analogous to that applied in \S\,43, for \(n\) variables.

Let again
\begin{equation}
 S=S(\alpha_1,\alpha_2,\ldots,\alpha_n)
\tag{1}
\end{equation}
be a whole symmetric function of the \(n\) variables \(\alpha_1,\alpha_2,\ldots,\alpha_n\). If we arrange it according to powers of \(\alpha_1\), and accordingly put
\begin{equation}
 S=S_0\alpha_1^\mu+S_1\alpha_1^{\mu-1}+\cdots+S_{\mu-1}\alpha_1+S_\mu,
\tag{2}
\end{equation}
%CONT>>


% ===== scan leaf 170, printed 145 =====
% !PRINTED 145
\weberpage{145}
%<<CONT
then the coefficients $S_0,S_1\ldots S_\mu$ are whole \gesperrt{symmetric} functions of the $n-1$ variables $\alpha_2,\ldots\alpha_n$.

If we denote the fundamental symmetric functions of these latter variables by $a'_1,a'_2\ldots a'_{n-1}$, then, by our supposition, we can express the coefficients $S_0,S_1\ldots S_\mu$ rationally in terms of these. But by \S\,42, (2) and (3),
\[
\begin{array}{l}
 a'_1=f_1(\alpha_1)=\alpha_1+a_1\\
 a'_2=f_2(\alpha_1)=\alpha_1^2+a_1\alpha_1+a_2\\
 a'_3=f_3(\alpha_1)=\alpha_1^3+a_1\alpha_1^2+a_2\alpha_1+a_3\\
 \makebox[23em][l]{\dotfill},
\end{array}
\]
that is, the coefficients $S_0,S_1\ldots S_\mu$ can be expressed wholly and rationally in terms of $\alpha_1,a_1,a_2\ldots a_n$. If, therefore, after introducing these expressions, we arrange (2) anew according to powers of $\alpha_1$, we obtain
\weq{3}{S=A_0\alpha_1^m+A_1\alpha_1^{m-1}+\cdots+A_{m-1}\alpha_1+A_m,}
where in general $m$ will be an exponent different from $\mu$, and the coefficients $A_0,A_1\ldots A_m$ depend wholly and rationally on $a_1,a_2\ldots a_n$. We put again
\weq{4}{\Phi(x)=A_0x^m+A_1x^{m-1}+\cdots+A_{m-1}x+A_m}
and divide $\Phi(x)$ by
\[
\begin{aligned}
 f(x)&=x^n+a_1x^{n-1}+a_2x^{n-2}+\cdots+a_n\\
     &=(x-\alpha_1)(x-\alpha_2)\ldots(x-\alpha_n).
\end{aligned}
\]
This gives a quotient and a remainder which, with respect to $x$, is at most of degree $n-1$. We therefore put
\weq{5}{\Phi(x)=Qf(x)+\psi(x),}
\weq{6}{\psi(x)=C_0x^{n-1}+C_1x^{n-2}+\cdots+C_{n-2}x+C_{n-1},}
and here $C_0,C_1\ldots C_{n-1}$ are whole rational functions of $a_1,a_2\ldots a_n$ in which, if $S$ in its original form contains no fractional coefficients, no fractions occur either.

But now, since $f(\alpha_1)$ vanishes and $S=\Phi(\alpha_1)$, it follows from (5) that
\weq{7}{S=\psi(\alpha_1)}
and here, since $S$ is symmetric, $\alpha_1$ may be replaced by $\alpha_2,\alpha_3\ldots\alpha_n$. This gives the following $n$ equations:
%CONT>>


% ===== scan leaf 171, printed 146 =====
% !PRINTED 146
\weberpage{146}
%<<CONT
\weq{8}{\begin{gathered}
 C_0\alpha_1^{n-1}+C_1\alpha_1^{n-2}+\cdots+C_{n-2}\alpha_1+(C_{n-1}-S)=0\\
 C_0\alpha_2^{n-1}+C_1\alpha_2^{n-2}+\cdots+C_{n-2}\alpha_2+(C_{n-1}-S)=0\\
 \makebox[30em]{\dotfill}\\
 C_0\alpha_n^{n-1}+C_1\alpha_n^{n-2}+\cdots+C_{n-2}\alpha_n+(C_{n-1}-S)=0.
\end{gathered}}
If we regard this system as a system of homogeneous linear equations with the $n$ unknowns
\[
 C_0,C_1\ldots C_{n-2},C_{n-1}-S,
\]
then its determinant,
\[
\begin{vmatrix}
 \alpha_1^{n-1},&\alpha_1^{n-2}&\ldots&\alpha_1,&1\\
 \alpha_2^{n-1},&\alpha_2^{n-2}&\ldots&\alpha_2,&1\\
 \multicolumn{5}{c}{\makebox[12em]{\dotfill}}\\
 \alpha_n^{n-1},&\alpha_n^{n-2}&\ldots&\alpha_n,&1
\end{vmatrix},
\]
which by \S\,22, formula (12), equals the product of all differences $\alpha_1-\alpha_2,\alpha_1-\alpha_3\ldots$, is different from zero, and consequently, by Theorem II in \S\,24,
\[
 C_0=0,\quad C_1=0\ldots\quad C_{n-2}=0
\]
and
\weq{9}{S=C_{n-1},}
which contains the fundamental theorem to be proved.

The same result is also obtained by Theorem II, \S\,30; for since the function of degree $(n-1)$ in $x$
\[
 C_0x^{n-1}+C_1x^{n-2}+\cdots+C_{n-2}x+C_{n-1}-S
\]
vanishes for $x=\alpha_1,\alpha_2\ldots\alpha_n$, that is, for more than $n-1$ values, all its coefficients must be zero.

The proof which we have given here for the fundamental theorem, following \gesperrt{Cauchy},\footnote[1]{\gesperrt{Cauchy}, Exercices de mathématiques, $4^{\text{ème}}$ année.} also provides a means of actually computing, in particular cases, the expression of a symmetric function in terms of the fundamental functions. The other known proofs of the theorem offer the same possibility. We shall give one more proof here, which often leads to a simpler method of computation\footnote[2]{\gesperrt{Waring}, Meditationes algebraicae. \gesperrt{Gauss}, Demonstratio nova altera theorematis omnem functionem algebraicam rationalem integram unius variabilis in factores reales primi vel secundi gradus resolvi posse. Werke, vol.\ III, p.\ 36.}.


% ===== scan leaf 172, printed 147 =====
% !PRINTED 147
\weberpage{147}
\pdfbookmark[2]{§ 45. Second proof of the theorem on symmetric functions}{bm51}
\begin{center}
{\large \S.\ 45.}

\vspace{0.5em}
{\large\gesperrt{Second Proof of the Theorem on Symmetric Functions.}}
\end{center}

Let $S$ be a whole symmetric function of the variables $\alpha_1,\alpha_2\ldots\alpha_n$. The individual terms of this function are all of the form
\[
 M\alpha_1^{\nu_1}\alpha_2^{\nu_2}\ldots\alpha_n^{\nu_n},
\]
where $M$ is a coefficient independent of the $\alpha$. These terms are now to be arranged in a definite manner. After the order of the variables $\alpha_1,\alpha_2\ldots\alpha_n$ has been fixed, of two terms
\[
 A=M\alpha_1^{\nu_1}\alpha_2^{\nu_2}\ldots\alpha_n^{\nu_n},\qquad
 A'=M'\alpha_1^{\nu_1}\alpha_2^{\nu'_2}\ldots\alpha_n^{\nu'_n}
\]
% [sic] Source p147 omits the prime in the first exponent of A'; retained, not silently normalized.
$A$ is to be called the \gesperrt{higher} when the first of the differences\begingroup\renewcommand{\thefootnote}{E1}\footnote[1]{[Source note:] The first exponent of the second term is printed as $\nu_1$, without a prime. The following definition requires $\nu'_1$ there; the printed expression is retained unchanged.}\endgroup
\[
 \nu_1-\nu'_1,\nu_2-\nu'_2\ldots\nu_n-\nu'_n,
\]
that is different from zero has a \gesperrt{positive} value, that is, when either $\nu_1>\nu'_1$ or $\nu_1=\nu'_1,\nu_2>\nu'_2$ or $\nu_1=\nu'_1,\nu_2=\nu'_2,\nu_3>\nu'_3$, and so on.

Since we suppose that all terms in which all the exponents $\nu_1,\nu_2\ldots\nu_n$ agree have been combined into one term, this decides which of any two terms is higher; and if $A$ is higher than $A'$, and $A'$ higher than $A''$, then $A$ is also higher than $A''$.

If now, according to this ordering,
\[
 A=M\alpha_1^{\nu_1}\alpha_2^{\nu_2}\ldots\alpha_n^{\nu_n}
\]
is the highest term of our function $S$, it follows that
\[
 \nu_1\eqslantgtr\nu_2
\]
must hold; for if $\nu_1<\nu_2$, the term
\[
 M\alpha_1^{\nu_2}\alpha_2^{\nu_1}\ldots\alpha_n^{\nu_n},
\]
which must also occur in $S$ because of symmetry, would stand higher in the order than $A$, contrary to the supposition. Likewise, $\nu_2\eqslantgtr\nu_3$ must hold; for if $\nu_2<\nu_3$, then
\[
 M\alpha_1^{\nu_1}\alpha_2^{\nu_3}\alpha_3^{\nu_2}\ldots\alpha_n^{\nu_n}
\]
would stand higher than $A$, and so forth. We therefore conclude that the exponents in $A$
%CONT>>


% ===== scan leaf 173, printed 148 =====
% !PRINTED 148
\weberpage{148}
%<<CONT
\[
 \nu_1,\nu_2,\nu_3\ldots\nu_n
\]
\gesperrt{form a decreasing, or at least never increasing, sequence of numbers}, or that the differences
\[
 \nu_1-\nu_2,\nu_2-\nu_3,\ldots\nu_{n-1}-\nu_n
\]
are all positive or at least not negative.

We conclude secondly that no term of the function $S$ can contain an exponent higher than $\nu_1$; for otherwise a term would also occur in which $\alpha_1$ had this higher exponent, and this would be of higher order than $A$, contrary to the supposition. Thus the terms that can occur at all in a symmetric function are, apart from their coefficients, completely determined by the highest term, and the number of possible terms is finite when the highest term is given.

We can use the exponent sequence of its highest term,
\[
 (\nu_1,\nu_2\ldots\nu_n)
\]
to denote the \gesperrt{order} of the symmetric function.

Apart from the case in which all exponents are zero, so that the function is a constant, the lowest possible order is
\[
 (1,0,0\ldots0),
\]
to which the single symmetric function
\[
 M(\alpha_1+\alpha_2+\cdots+\alpha_n)=-Ma_1
\]
corresponds.\begingroup\renewcommand{\thefootnote}{E2}\footnote[1]{[Source note:] These examples leave an arbitrary constant summand unstated; adding such a summand does not change the stated order. The formulas are retained as printed.}\endgroup{} The next higher order is
\[
 (1,1,0\ldots0),
\]
to which corresponds the symmetric function
\[
 Ma_1+M'a_2
\]
and thus in the first $n$ orders we obtain the fundamental functions themselves and their linear combinations, and no others.

The next order is characterized by
\[
 (2,0,0\ldots0)
\]
and contains the linear combinations of the fundamental functions with the sum of the squares.

In all these cases the representability of the symmetric functions has already been proved, and we shall therefore now give the general proof by deriving the representation
%CONT>>


% ===== scan leaf 174, printed 149 =====
% !PRINTED 149
\weberpage{149}
%<<CONT
of the function $S$ of order $(\nu_1,\ \nu_2 \ldots \nu_n)$ by the
fundamental functions under the assumption that it is already known
for the functions of lower order, that is, by the method of
inference of complete induction.

We remark on this that by the multiplication of two
symmetric functions $S$ and $S'$, whose highest terms are
\[
A = M\,\alpha_1^{\nu_1}\,\alpha_2^{\nu_2} \ldots \alpha_n^{\nu_n},\qquad
A' = M'\,\alpha_1^{\nu'_1}\alpha_2^{\nu'_2} \ldots \alpha_n^{\nu'_n}
\]
a symmetric function arises whose highest term is
the product
\[
A\,A' = M\,M'\,\alpha_1^{\nu_1+\nu'_1}\,\alpha_2^{\nu_2+\nu'_2} \ldots \alpha_n^{\nu_n+\nu'_n}
\]
For suppose there were in $S\,S'$ a higher term
than $A\,A'$, arising from the multiplication of the two terms
\[
B = N\,\alpha_1^{\mu_1}\,\alpha_2^{\mu_2} \ldots \alpha_n^{\mu_n},\qquad
B' = N'\,\alpha_1^{\mu'_1}\alpha_2^{\mu'_2} \ldots \alpha_n^{\mu'_n}
\]
so that
\[
B\,B' = N\,N'\,\alpha_1^{\mu_1+\mu'_1}\,\alpha_2^{\mu_2+\mu'_2} \ldots \alpha_n^{\mu_n+\mu'_n},
\]
and suppose that we do not have simultaneously $B = A$ and $B' = A'$; then the
first of the differences
\[
\mu_1 - \nu_1 + \mu'_1 - \nu'_1,\quad \mu_2 - \nu_2 + \mu'_2 - \nu'_2, \ldots \mu_n - \nu_n + \mu'_n - \nu'_n,
\]
that is different from zero would be positive, which is impossible, since the
first non-vanishing one among the differences
\[
\mu_1 - \nu_1,\quad \mu_2 - \nu_2 \ldots \mu_n - \nu_n
\]
and among the differences
\[
\mu'_1 - \nu'_1,\quad \mu'_2 - \nu'_2 \ldots \mu'_n - \nu'_n
\]
is by hypothesis negative. By repeating this
argument we obtain the more general theorem, that the
highest term of a product of several factors is obtained
by multiplying the highest terms of the individual factors.

The highest terms of the fundamental symmetric functions
\[
a_1,\ a_2,\ a_3 \ldots a_n
\]
are now, apart from sign,
\[
\alpha_1,\quad \alpha_1\,\alpha_2,\quad \alpha_1\,\alpha_2\,\alpha_3, \ldots, \alpha_1\,\alpha_2\,\alpha_3 \ldots \alpha_n,
\]
and if we therefore form the product
\[
P = \pm\,M\,a_1^{\nu_1-\nu_2}\,a_2^{\nu_2-\nu_3} \ldots a_{n-1}^{\nu_{n-1}-\nu_n}\,a_n^{\nu_n},
\]
we obtain a symmetric function whose highest
term is likewise $A$, as in $S$.


% ===== scan leaf 175, printed 150 =====
% !PRINTED 150
\weberpage{150}
\pdfbookmark[2]{§ 46. Discriminants}{bm52}
The difference
\[
S - P = S'
\]
is therefore again a symmetric function whose highest term is
lower than that of $S$, and which by hypothesis we can represent
rationally by the fundamental functions.
Since $P$ is represented in the same way, the goal is attained. It is
also evident here once again that no fractions are introduced by the
representation by means of the fundamental functions, if originally
none are contained in $S$.

This method of computing symmetric functions at the same
time gives us information about the degree of the integral rational
function of $a_1,\ a_2 \ldots a_n$ by which a given
symmetric function of the $\alpha_1,\ \alpha_2 \ldots \alpha_n$ is expressed.
Namely, if $(\nu_1,\ \nu_2 \ldots \nu_n)$ is the order of the symmetric
function (measured by the sequence of exponents of the first
term), then $\nu_1$ is the degree of $P$ with respect to all of
$a_1, a_2 \ldots a_n$, and in the developed expression for $S$ no
term of higher degree can occur; nor can the term $P$ be
destroyed by any of the following terms, since $P$ is completely
determined by the order of $S$, and the order of $S'$
and of all the following symmetric functions is lower than
the order of $S$.

If, then, in place of $a_1, a_2 \ldots a_n$ we put:
\[
\frac{a_1}{a_0},\quad \frac{a_2}{a_0} \ldots \frac{a_n}{a_0},
\]
then $a_0^{\nu_1}\,S$ is an \gesperrt{integral homogeneous function
of degree} $\nu_1$ \gesperrt{in} $a_0,\ a_1 \ldots a_n$, and one cannot turn $S$ into an integral
function of the $a$ by multiplication with a lower power of $a_0$.

\begin{center}
\S.\,46.

\gesperrt{Discriminants.}
\end{center}

The theorems obtained in the foregoing teach us to express
symmetric functions of the roots of an equation rationally by
the coefficients of the equation; the results are
at first derived under the assumption that the roots,
or, what by the third section amounts to the same thing,
%CONT>>


% ===== scan leaf 176, printed 151 =====
% !PRINTED 151
\weberpage{151}
%<<CONT
the coefficients were independent variables; but they cannot
become false when definite numerical values are put in their
place, since, at least for the integral symmetric functions, no
denominators occur that could vanish.

We first make some applications to particular symmetric
functions which are important for what follows. Already in
\S.\,19 it was remarked that the difference product

\weq{1}{\begin{aligned}
P = {}&(\alpha_1 - \alpha_2)\,(\alpha_1 - \alpha_3)\ldots(\alpha_1 - \alpha_n)\\
&\phantom{(\alpha_1 - \alpha_2)\,}(\alpha_2 - \alpha_3)\ldots(\alpha_2 - \alpha_n)\\
&\cdot\;\;\cdot\;\;\cdot\;\;\cdot\;\;\cdot\;\;\cdot\;\;\cdot\;\;\cdot\;\;\cdot\;\;\cdot\;\;\cdot\;\;\cdot\\
&\phantom{(\alpha_1 - \alpha_2)\,(\alpha_1 - \alpha_3)\ldots}(\alpha_{n-1} - \alpha_n)
\end{aligned}}

merely changes its sign under an interchange of two elements.
The square of this product will therefore remain unchanged
under all permutations of two of the variables $\alpha$, and hence
also under all permutations; that is, it will be a symmetric function
of the variables. The order of this symmetric function
is, as one sees from (1), equal to

\[
(2\,n - 2,\ \ 2\,n - 4 \ldots 2,\, 0).
\]

If, then, we assume the function $f(x)$, whose roots are the $\alpha_1, \alpha_2 \ldots \alpha_n$,
in the form

\[
f(x) = a_0\,x^n + a_1\,x^{n-1} + a_2\,x^{n-2} + \cdots + a_{n-1}\,x + a_n,
\]

so that the highest power of $x$ does not necessarily have the coefficient
1, then $P^2$ is an integral rational function of

\weq{2}{\frac{a_1}{a_0},\ \ \frac{a_2}{a_0} \ldots \frac{a_n}{a_0},}

which, on multiplication by $a_0^{2n-2}$, goes over into a \gesperrt{homogeneous} function
of $a_0, a_1, \ldots a_n$. We call this homogeneous function,
which is thus defined by

\weq{3}{a_0^{2\,n-2}\,P^2 = D}

the \gesperrt{discriminant of the function} $f(x)$. If we set
$a_0 = 1$, we obtain the \gesperrt{inhomogeneous form of the
discriminant}\footnote{The discriminant is, apart from the factor $(-\,1)^{\frac{n\,(n-1)}{2}}\ a_0^{2\,n-2}$,
the same as what \gesperrt{Gauss} calls the determinant of the function (l. c., art. 6).}.


% ===== scan leaf 177, printed 152 =====
% !PRINTED 152
\weberpage{152}

\gesperrt{The vanishing of the discriminant is the necessary
and sufficient condition that two of the values}
$\alpha_1,\ \alpha_2 \ldots \alpha_n$ \gesperrt{become equal to one another.}

We can set up yet a second expression for the
discriminant.

By \S.\,29, (4) we have

% [sic: this block of formulas is unnumbered in print; the numbering jumps from (3) to (5)]
\[
\begin{aligned}
f'(\alpha_1) &= a_0\,(\alpha_1 - \alpha_2)\,(\alpha_1 - \alpha_3)\ldots(\alpha_1 - \alpha_n)\\
f'(\alpha_2) &= a_0\,(\alpha_2 - \alpha_1)\,(\alpha_2 - \alpha_3)\ldots(\alpha_2 - \alpha_n)\\
\cdot\;\;\cdot\;\;\cdot\;\;\cdot\;\;\cdot\;\;\cdot &\;\;\cdot\;\;\cdot\;\;\cdot\;\;\cdot\;\;\cdot\;\;\cdot\;\;\cdot\;\;\cdot\;\;\cdot\\
f'(\alpha_n) &= a_0\,(\alpha_n - \alpha_1)\,(\alpha_n - \alpha_2)\ldots(\alpha_n - \alpha_{n-1}),
\end{aligned}
\]

and consequently, since here every factor $(\alpha_h - \alpha_k)$ occurs twice with
opposite signs,

\weq{5}{D = (-\,1)^{\frac{n\,(n-1)}{2}}\ a_0^{n-2}\,f'(\alpha_1)\,f'(\alpha_2) \ldots f'(\alpha_n).}

In order actually to represent the discriminant, we have, first of all a general means. % [sic: comma as printed]
By \S.\,22 we had

\weq{6}{(-\,1)^{\frac{n\,(n-1)}{2}}\ P = \left|\,
\begin{array}{l}
1,\ \alpha_1,\ \alpha_1^2 \ldots \alpha_1^{n-1}\\
1,\ \alpha_2,\ \alpha_2^2 \ldots \alpha_2^{n-1}\\
\cdot\;\;\cdot\;\;\cdot\;\;\cdot\;\;\cdot\;\;\cdot\;\;\cdot\\
1,\ \alpha_n,\ \alpha_n^2 \ldots \alpha_n^{n-1}
\end{array}
\,\right|,}

and if, following \S.\,27, we represent the square of this determinant,
by multiplication along vertical rows, as a new determinant,
and set

\[
s_m = \alpha_1^m + \alpha_2^m + \cdots \alpha_n^m
\]

then

\weq{7}{D = a_0^{2\,n-2}\left|\,
\begin{array}{l}
s_0,\ \ \ s_1,\ s_2\ \ \ldots s_{n-1}\\
s_1,\ \ \ s_2,\ s_3\ \ \ldots s_n\\
\cdot\;\;\cdot\;\;\cdot\;\;\cdot\;\;\cdot\;\;\cdot\;\;\cdot\;\;\cdot\;\;\cdot\\
s_{n-1},\ s_n,\ s_{n+1} \ldots s_{2\,n-2}
\end{array}
\,\right|.}

The quantities $s_m$, however, are the power sums which we have already
expressed in \S.\,42 by the coefficients; only, in place
of the $a_1,\ a_2 \ldots a_n$, the quotients (2) are to be put.

Thus we find, e.g., for $n = 2$

\weq{8}{D = a_0^2\,(s_0\,s_2 - s_1^2) = a_1^2 - 4\,a_0\,a_2}

and for $n = 3$

\weq{9}{D = a_0^4\,(s_0\,s_2\,s_4 - s_0\,s_3^2 + 2\,s_1\,s_2\,s_3 - s_1^2\,s_4 - s_2^3).}


% ===== scan leaf 178, printed 153 =====
% !PRINTED 153
\weberpage{153}
\pdfbookmark[2]{§ 47. Discriminants of the forms of the third and fourth order}{bm53}

% [restored 2026-08-06] omission raised by this page's own translation pass and confirmed
% against the scan; the German fragment now carries the line too.
Herein one is to set, according to \S.\,42,
\[
\begin{aligned}
s_0 &= 3, \quad a_0 s_1 = -a_1\\
a_0^2 s_2 &= a_1^2 - 2 a_0 a_2\\
a_0^3 s_3 &= -a_1^3 + 3 a_0 a_1 a_2 - 3 a_0^2 a_3\\
a_0^4 s_4 &= a_1^4 - 4 a_0 a_1^2 a_2 + 4 a_0^2 a_1 a_3 + 2 a_0^2 a_2^2,
\end{aligned}
\]
whereby one obtains

\weq{10}{D = a_1^2 a_2^2 + 18 a_0 a_1 a_2 a_3 - 4 a_0 a_2^3 - 4 a_1^3 a_3 - 27 a_0^2 a_3^2.}

\begin{center}
\S.\,47.

\gesperrt{Discriminants of the Forms of Third and Fourth Order.}
\end{center}

For the functions of third and fourth order we have, in the solutions of the equations of these two degrees (\S.\,35, 37), a simple means of computing the discriminants.

As regards first the equation of third order, the three roots of
\weq{1}{x^3 + a x + b = 0}
were represented in the form:
\weq{2}{\begin{aligned}
\alpha &= u + v\\
\beta &= \varepsilon u + \varepsilon^2 v\\
\gamma &= \varepsilon^2 u + \varepsilon v,
\end{aligned}}
in which
\weq{3}{\varepsilon = \frac{-1 + \sqrt{-3}}{2}, \quad \varepsilon^2 = \frac{-1 - \sqrt{-3}}{2}}
are the complex third roots of unity, and $u$, $v$ are determined by \S.\,35, (8).

From (3) it follows that
\weq{4}{\varepsilon - \varepsilon^2 = \sqrt{-3}, \quad 1 - \varepsilon = \varepsilon^2 \sqrt{-3}, \quad 1 - \varepsilon^2 = -\varepsilon \sqrt{-3}}
and from (2)
\[
\begin{aligned}
\alpha - \beta &= \sqrt{-3}\,(\varepsilon^2 u - \varepsilon v)\\
\alpha - \gamma &= -\sqrt{-3}\,(\varepsilon u - \varepsilon^2 v)\\
\beta - \gamma &= \sqrt{-3}\,(u - v).
\end{aligned}
\]
From this, by multiplication,
\weq{5}{(\alpha - \beta)(\alpha - \gamma)(\beta - \gamma) = 3 \sqrt{-3}\,(u^3 - v^3).}


% ===== scan leaf 179, printed 154 =====
% !PRINTED 154
\weberpage{154}

Squaring, one obtains for the discriminant of the cubic equation (1), by \S.\,35, (6)
\weq{6}{D = -(27 b^2 + 4 a^3).}
From this one obtains the discriminant of the general function
\weq{7}{a_0 x^3 + a_1 x^2 + a_2 x + a_3,}
if for $a$ and $b$ one substitutes the expressions \S.\,34, (4), in which one replaces $a_1$, $a_2$, $a_3$ by
\[
\frac{a_1}{a_0}, \quad \frac{a_2}{a_0}, \quad \frac{a_3}{a_0}
\]
and multiplies [\S.\,46, (3)] by $a_0^4$. One thus obtains
\[
\begin{aligned}
a &= \frac{3 a_0 a_2 - a_1^2}{3 a_0^2},\\
b &= \frac{2 a_1^3 - 9 a_0 a_1 a_2 + 27 a_0^2 a_3}{27 a_0^3},
\end{aligned}
\]
and consequently
\weq{8}{D = -\frac{1}{27 a_0^2}\left\{(2 a_1^3 - 9 a_0 a_1 a_2 + 27 a_0^2 a_3)^2 + 4(3 a_0 a_2 - a_1^2)^3\right\}.}

If one works out the square and the cube, one comes back to formula (10) of the preceding paragraph; it is remarkable, however, that formula (8) apparently contains the denominator $27 a_0^2$, which cancels out in the computation.

We can proceed in the same way with the biquadratic equations. If $\alpha$, $\beta$, $\gamma$, $\delta$ are the roots of the biquadratic equation
\weq{9}{x^4 + a x^2 + b x + c = 0}
we obtain, by \S.\,37, (5)
\weq{10}{\begin{aligned}
\alpha - \beta &= v + w, & \gamma - \delta &= v - w,\\
\alpha - \gamma &= w + u, & \delta - \beta &= w - u,\\
\alpha - \delta &= u + v, & \beta - \gamma &= u - v.
\end{aligned}}
By this the discriminant of equation (9) becomes
\[
D = (v^2 - w^2)^2 (w^2 - u^2)^2 (u^2 - v^2)^2,
\]
that is, $D$ is equal to the discriminant of the cubic resolvent \S.\,37, (4)
\weq{11}{z^3 + 2 a z^2 + (a^2 - 4c) z - b^2 = 0,}
whose roots are $u^2$, $v^2$, $w^2$.

But if one forms this discriminant according to formula (10), \S.\,46, by
%CONT>>


% ===== scan leaf 180, printed 155 =====
% !PRINTED 155
\weberpage{155}

%<<CONT
$a_0 = 1, \quad a_1 = 2\,a, \quad a_2 = a^2 - 4\,c, \quad a_3 = -\,b^2$

then it follows that
\weq{12}{
\begin{aligned}
D ={}& -\,4\,a^3 b^2 + 144\,a c b^2 + 16\,a^4 c \\
& - 128\,a^2 c^2 + 256\,c^3 - 27\,b^4.
\end{aligned}
}

If, on the other hand, one applies formula (8) to form the discriminant of the cubic equation (11), one obtains

\weq{13}{
27\,D = 4\,(a^2 + 12\,c)^3 - (2\,a^3 - 72\,a c + 27\,b^2)^2.
}

From this one derives the discriminant of the general function of the fourth degree

\weq{14}{
f(x) = a_0\,x^4 + a_1\,x^3 + a_2\,x^2 + a_3\,x + a_4
}

by expressing $a$, $b$, $c$, according to \S.\,34, (4), by

\[
\frac{-\,3\,a_1^2}{8} + a_2
\]

\[
\frac{a_1^3}{8} - \frac{a_1 a_2}{2} + a_3
\]

\[
\frac{-\,3\,a_1^4}{256} + \frac{a_1^2 a_2}{16} - \frac{a_1 a_3}{4} + a_4
\]

then replacing $a_1$, $a_2$, $a_3$, $a_4$ by

\[
\frac{a_1}{a_0}, \quad \frac{a_2}{a_0}, \quad \frac{a_3}{a_0}, \quad \frac{a_4}{a_0}
\]

and multiplying by $a_0^6$.

If, after carrying out this simple calculation, one sets for abbreviation

\weq{15}{
A = a_2^2 - 3\,a_1 a_3 + 12\,a_0 a_4,
}

\weq{16}{
B = 27\,a_1^2 a_4 + 27\,a_0 a_3^2 + 2\,a_2^3 - 72\,a_0 a_2 a_4 - 9\,a_1 a_2 a_3,
}

then one obtains from (13)

\weq{17}{
27\,D = 4\,A^3 - B^2.
}

The quantities $A$ and $B$, which will occupy us further in the next section, are called the first and second invariant of the biquadratic function $f$. By (17), then, the discriminant can be expressed rationally through these quantities, though not in such a way that the coefficients become integral; if one carries out the calculation prescribed in (17), in order to represent $D$ through the coefficients $a_0, a_1, a_2, a_3$ themselves, % [sic] printed coefficient list stops at a_3 although (14) contains a_4
the factor 27 must still cancel out. We shall not set down here the long expression, whose computation offers no difficulty whatever,
%CONT>>


% ===== scan leaf 181, printed 156 =====
% !PRINTED 156
\weberpage{156}
\pdfbookmark[2]{§ 48. Resultants}{bm54}

\begin{center}
\S.\,48.

\gesperrt{Resultants.}
\end{center}

Another application of the theory of symmetric functions, which moreover contains the formation of the discriminant as a special case, is the formation of resultants.

Let

\weq{1}{
f(x) = a_0\,x^n + a_1\,x^{n-1} + a_2\,x^{n-2} + \cdots + a_n,
}

\weq{2}{
\varphi(x) = b_0\,x^m + b_1\,x^{m-1} + b_2\,x^{m-2} + \cdots + b_m
}

be two integral rational functions of degree $n$ and $m$, and let $\alpha_1, \alpha_2 \ldots \alpha_n$ be the roots of the first, so that the following factorization holds:

\weq{3}{
f(x) = a_0\,(x - \alpha_1)\,(x - \alpha_2) \ldots (x - \alpha_n)
}

Now the product

\[
\varphi(\alpha_1)\,\varphi(\alpha_2) \ldots \varphi(\alpha_n)
\]

is a symmetric function of the roots of $f(x)$, and can therefore be expressed integrally and rationally in terms of the following quotients:

\[
\frac{a_1}{a_0}, \quad \frac{a_2}{a_0} \ldots \frac{a_n}{a_0}
\]

The order of this symmetric function is $(m, m \ldots m)$ (\S.\,45), and hence

\weq{4}{
R = a_0^m\,\varphi(\alpha_1)\,\varphi(\alpha_2) \ldots \varphi(\alpha_n)
}

is an integral rational and homogeneous function of order $m$ in $a_0, a_1 \ldots a_n$. It is also, as its expression shows, an integral homogeneous function of order $n$ in $b_0, b_1 \ldots b_m$. It is called the \gesperrt{resultant} of the two functions $f(x)$ and $\varphi(x)$. This designation is justified by a second representation.

If we denote the roots of $\varphi(x)$ by $\beta_1, \beta_2 \ldots \beta_m$, and thus put

\weq{5}{
\varphi(x) = b_0\,(x - \beta_1)\,(x - \beta_2) \ldots (x - \beta_m),
}

then by (4) we obtain

\weq{6}{
R = a_0^m\,b_0^n \overset{i,\,k}{\Pi}(\alpha_i - \beta_k),
}

in which the product sign $\Pi$ extends over all the indices $i = 1, 2 \ldots n$, $k = 1, 2 \ldots m$,
%CONT>>


% ===== scan leaf 182, printed 157 =====
% !PRINTED 157
\weberpage{157}

This expression now shows that $R$, apart from the factor
$(-1)^{mn}$, depends on $\varphi(x)$ in the same way as
on $f(x)$, and that we may therefore also set

\weq{7}{(-1)^{mn}\, R = b_0^n\, f(\beta_1)\, f(\beta_2)\, \ldots\, f(\beta_m).}

The discriminants are accordingly a special case of the
resultants, which one obtains when one takes for $\varphi(x)$ the derivative
$f'(x)$.

\gesperrt{The vanishing of the resultant of $f(x)$ and
$\varphi(x)$ is the necessary and sufficient condition
for the two equations}

\[ f(x) = 0 \quad \text{and} \quad \varphi(x) = 0 \]

\gesperrt{to have a common root.}

The equation that arises by setting the resultant equal to zero
is also called the \gesperrt{final equation} from $f(x) = 0$ and $\varphi(x) = 0$,
and the setting up of the \gesperrt{final equation} the \gesperrt{elimination}
of $x$.

For the computation of resultants one may apply the algorithm
of the greatest common divisor (\S.\,6). For if
$m \leqq n$, one can determine an integral function $Q$ and a like
 function $\psi(x)$, whose degree is smaller than $m$, in such a way
 that

\[ f(x) = Q\varphi(x) + \psi(x), \]

and hence, if one puts for $x$ one of the roots $\beta_i$,

\[ f(\beta_i) = \psi(\beta_i), \]

hence

\weq{8}{b_0^n \overset{i}{\Pi} f(\beta_i) = b_0^n \overset{i}{\Pi} \psi(\beta_i);}

if $m'$ is the degree of $\psi(x)$, then

\weq{9}{b_0^{m'} \overset{i}{\Pi} \psi(\beta_i) = R'}

is the resultant of $\varphi(x)$ and $\psi(x)$, and therefore

\weq{10}{R = b_0^{n-m'}\, R'.}

The formation of the resultant $R$ is hereby reduced to the formation
of $R'$. One can now divide $\varphi(x)$ again by $\psi(x)$,
and thus reduces the formation of $R'$ to the formation of a
resultant of functions of ever lower degree, until
one finally arrives at a constant, i.\,e. a function of the
coefficients $a$, $b$, which must be identical with the resultant
 sought.


% ===== scan leaf 183, printed 158 =====
% !PRINTED 158
\weberpage{158}

The computation of resultants is usually very laborious. We shall cite only the simplest example, the resultant of two quadratic functions:
\[
f(x) = a_0 x^2 + a_1 x + a_2
\]
\[
\varphi(x) = b_0 x^2 + b_1 x + b_2.
\]
Here one can compute the product
\[
R = a_0^2 (b_0 \alpha_1^2 + b_1 \alpha_1 + b_2) (b_0 \alpha_2^2 + b_1 \alpha_2 + b_2)
\]
directly, and obtains, on setting
\[
- a_0 (\alpha_1 + \alpha_2) = a_1, \qquad a_0 \alpha_1 \alpha_2 = a_2
\]
the result
\weq{11}{
R = a_0^2 b_2^2 + a_2^2 b_0^2 - b_0 b_1 a_1 a_2 - a_0 a_1 b_1 b_2
+ b_0 b_2 a_1^2 + a_0 a_2 b_1^2 - 2 a_0 b_0 a_2 b_2,
}
or, what is identical with it,
\weq{12}{
R = (a_0 b_2 - a_2 b_0)^2 - (a_0 b_1 - a_1 b_0) (a_1 b_2 - a_2 b_1).
}
The individual terms of the resultant of the two functions $f(x)$ and $\varphi(x)$ of degree $n$ and $m$ have, apart from a numerical (integer) factor, the form
\weq{13}{
a_0^{\nu_0} a_1^{\nu_1} \ldots a_n^{\nu_n} b_0^{\mu_0} b_1^{\mu_1} \ldots b_m^{\mu_m},
}
in which, as follows from the degrees determined above,
% [source-original mathematical slip, confirmed 2026-08-25 against scan p. 183 / printed p. 158]
% The printed right-hand sides are reversed: R has degree m in the a-coefficients and degree n
% in the b-coefficients, so the consistent relations would be sum(nu_i)=m and sum(mu_i)=n.
% The facing translation preserves the printed formulas verbatim.
\weq{14}{
\begin{aligned}
\nu_0 + \nu_1 + \cdots + \nu_n &= n \\
\mu_0 + \mu_1 + \cdots + \mu_m &= m.
\end{aligned}
}
But we can give yet another relation between these exponents.

For since [by (6)] $R$ is a homogeneous function of degree $n\,m$ of the $n+m$ variables $\alpha, \beta$, and since further $a_0$ is of the zeroth, $a_1$ of the first, $a_2$ of the second, etc., degree in the $\alpha$, and in general $a_k$ and $b_k$ are of degree $k$ in these variables, it follows that
\weq{15}{
\begin{aligned}
\nu_1 + 2 \nu_2 + 3 \nu_3 + \cdots + n \nu_n \\
{}+ \mu_1 + 2 \mu_2 + 3 \mu_3 + \cdots + m \mu_m &= n m,
\end{aligned}
}
a relation from which an important conclusion can be drawn, one that we shall discuss somewhat more fully in the following section.


% ===== scan leaf 184, printed 159 =====
% !PRINTED 159
\weberpage{159}
\pdfbookmark[2]{§ 49. Elimination. Bezout's theorem}{bm55}

\begin{center}
\S.\,49.

\gesperrt{Elimination. Theorem of Bezout.}
\end{center}

If in the two functions

\weq{1}{
\begin{aligned}
f(x) &= a_0 x^n + a_1 x^{n-1} + \cdots + a_n \\
\varphi(x) &= b_0 x^m + b_1 x^{m-1} + \cdots + b_m
\end{aligned}
}

the coefficients $a$, $b$ are themselves again integral rational functions
of a quantity $y$, then the resultant will also be an integral
rational function of $y$, which we shall denote by $F(y)$.
The formation of this equation is the elimination of
$x$ from the two equations (1).

The equation

\weq{2}{F(y) = 0}

has as roots all those values of $y$ for which the two
equations

\weq{3}{f(x) = 0, \qquad \varphi(x) = 0}

have a common root. If the condition (2) is satisfied,
one finds the values of $x$ that satisfy the two equations
(3) by seeking the greatest common divisor
of $f(x)$ and $\varphi(x)$. This common divisor
is of the first degree, and thus determines $x$ rationally, when
only one such common root is present; in the other
case it will be of correspondingly higher degree, and the
determination of $x$ still requires the solution of an equation
of higher degree.

If to each root of the resultant there belongs only one common
root $x$, then the number of pairs of values $x$, $y$ that
satisfy the equations (3) is equal to the degree of the resultant
with respect to $y$. Thus if the coefficients $a$ are of degree $\nu$, $b$
of degree $\mu$ with respect to $y$, then by the determination of degree
of the preceding paragraph the degree of the resultant is

% [source-original mathematical slip, confirmed 2026-08-25 against scan p. 184 / printed p. 159]
% The printed expression n*nu + m*mu reverses the bidegrees; the consistent degree is
% m*nu + n*mu. The facing translation preserves the printed formula verbatim.
\weq{5}{n\nu + m\mu;} % [sic] equation numbering jumps from (3) to (5) in print; no (4) occurs here

so great, then, is the number of common pairs of roots.
This number can, however, still be modified by the fact that
the coefficients may possibly be so constituted that
%CONT>>


% ===== scan leaf 185, printed 160 =====
% !PRINTED 160
\weberpage{160}
%<<CONT
the highest power of $y$ drops out of the final equation. One
could then save the agreement of the number with (5) only by
the arbitrary addition of infinite roots. Multiple roots, too, give
cause for misgivings. These drawbacks are met in part by taking
as a basis the \gesperrt{homogeneous} form of the equations, a form in
which the elimination problem presents itself especially in
geometry.

For if the function $f(x)$ has arisen from a homogeneous
function of order $n$ of the three variables $x, y, z$ (a ternary
form) by arranging according to powers of $x$, then the
coefficients $a_0, a_1, a_2 \ldots a_n$ are homogeneous forms of the
two variables $y, z$ of the degree indicated by the index, so
that $a_0$ is a constant, $a_1$ a linear, $a_2$ a quadratic form,
and so forth, and the corresponding holds of the coefficients
$b_0, b_1 \ldots b_m$ of the function $\varphi(x)$.

If $x, y, z$ denote triangular coordinates in the plane, then
$f(x) = 0$, $\varphi(x) = 0$ are the equations of a curve of order $n$ and
one of order $m$, and the common roots of these two
equations (the ratios $x:y:z$) denote the points of intersection
of the two curves. The final equation, which one obtains by
eliminating $x$, is a homogeneous equation whose degree
comes out, by (15), \S.\,48, equal to $nm$.

A lowering of the degree cannot take place here; the
final equation is always a homogeneous equation of degree $nm$
in the two unknowns $y, z$. The only exception to which
attention must be paid here is that the resultant vanishes
identically (for all $y, z$), so that infinitely many common
roots are present. In geometry this case has the meaning
that the two curves have a portion of a curve in common, and
not merely isolated points.

If the equations $f = 0$, $\varphi = 0$ have only a finite
number of common roots, then in the expressions

\[
y_1 = y + \beta x, \qquad z_1 = z + \gamma x
\]

we can arrange the constants $\beta, \gamma$ so that, whenever to one value
of the ratio $y : z$ there correspond several distinct values of $x$
satisfying the equations (1), the associated values
of $y_1 : z_1$ turn out different from one another.


% ===== scan leaf 186, printed 161 =====
% !PRINTED 161
\weberpage{161}
\pdfbookmark[2]{§ 50. Elimination from several equations}{bm56}
% [source-original mathematical slip, confirmed 2026-08-25 against scan p. 186 / printed p. 161]
% The print assigns degrees m and n, respectively, to f and phi, although printed p. 156
% defined f with degree n and phi with degree m. The facing translation preserves that order.
We can then also regard $f$ and $\varphi$ as homogeneous functions
of degree $m$ and $n$ in the variables $x, y_1, z_1$, and
the resultant then becomes a homogeneous function $R(y_1, z_1)$
of degree $mn$ in $y_1$ and $z_1$. If we set in it

\[
y_1 = b' + b\lambda, \qquad z_1 = c' + c\lambda,
\]

we obtain an equation of degree $mn$ in $\lambda$, and,
provided $R(y_1, z_1)$ does not vanish identically, we can choose the constants
$b, b', c, c'$ so that the coefficient of the highest
power of $\lambda$, namely $R(b, c)$, is different from zero. Then
to each root of the resultant there corresponds only one system of values
of the ratio $x:y:z$, and so the number of common
roots of $f$ and $\varphi$ is at most equal to $mn$. In order to be able to
state the theorem that the number of common roots
is always equal to $mn$, one must further adopt a convention
according to which certain of these pairs of values are to be counted multiply.

The theorem which, in geometric dress, says that
two curves of the $m$th and $n$th order intersect in $mn$ points
is called \gesperrt{Bezout's theorem}.

\begin{center}
§.\,50.

\gesperrt{Elimination from several equations.}
\end{center}

The principles of elimination, which we have applied in the
foregoing to two equations, can be extended without
difficulty to several equations with a corresponding
number of unknowns.

Let us first take the equations

\weq{1}{f(x) = 0, \quad \varphi(x) = 0, \quad \psi(x) = 0,}

of degrees $m, n, p$, whose coefficients $a_i, b_i, c_i$ are at first to be
regarded as independent variables.

We apply to two of them, say to $f(x), \varphi(x)$, the
algorithm of the greatest common divisor and
obtain a last remainder, which is a rational function of the $a$
and $b$ and which, freed of denominators, agrees with the resultant $R$
of $f(x)$ and $\varphi(x)$, and a last remainder but one, which is a
\gesperrt{linear function} of $x$. The condition for the existence
of a common root of $f = 0$ and $\varphi = 0$
is the vanishing of the resultant, and the last remainder but one
%CONT>>


% ===== scan leaf 187, printed 162 =====
% !PRINTED 162
\weberpage{162}
%<<CONT
gives, under this assumption, a rational expression for the common
root in terms of the coefficients.

If one substitutes this expression into $\psi(x)$, then, on clearing
away all denominators, one obtains an integral rational function
of the coefficients $a$, $b$, $c$, which we shall denote by $S$, and
\gesperrt{whose vanishing, in conjunction with $R = 0$, contains the necessary and
sufficient condition for the three equations (1) to have a common
root.}

Let us now replace the coefficients $a$, $b$, $c$ of the functions
(1) by integral rational functions of a second unknown
$y$, whose coefficients shall be denoted by $\alpha$, $\beta$, $\gamma$ and for the time being
regarded as independent variables.

$R$ and $S$ then become integral rational functions of
$\alpha$, $\beta$, $\gamma$ and $y$, and if we therefore regard the equations

\weq{2}{R = 0, \qquad S = 0}

as equations for $y$, then by the rules of the preceding
paragraph we can form their resultant $T$, which,
if it does not vanish identically, will be an integral rational
function of the coefficients $\alpha$, $\beta$, $\gamma$, and the equation

\weq{3}{T = 0}

is the necessary and sufficient condition for the
system of equations (1) to be satisfiable by one or more pairs of values $x$, $y$.
$T$ is called the \gesperrt{resultant of the system
(1).} % [source-original logical slip, confirmed 2026-08-25 against scan p. 187 / printed p. 162]
% The print says that eliminating x and y forms a final equation "for x", although x has just
% been eliminated. The facing translation preserves the printed variable rather than guessing.
One can therefore, from the three equations (1), by \gesperrt{elimination of the unknowns $x$ and $y$, form a final equation for
$x$.}

If one now replaces the $\alpha$, $\beta$, $\gamma$ in turn by integral rational
functions of a third unknown $z$, then $T$ passes over into an
equation for $z$. The functions $f$, $\varphi$, $\psi$ become functions
of the three unknowns $x$, $y$, $z$. The degree of equation (3) gives
the number of systems of values $x$, $y$, $z$ for which the three equations
(1) are simultaneously satisfied.

But if the coefficients are so constituted that $T$ vanishes
identically (for all values of $z$), then there are infinitely many
systems of values $x$, $y$, $z$ which satisfy the equations (1) simultaneously.

If one regards $x$, $y$, $z$ as Cartesian coordinates, then
the equations (1), each by itself, represent a surface.
The degree of the final equation $T = 0$ with respect to $z$ gives the
%CONT>>


% ===== scan leaf 188, printed 163 =====
% !PRINTED 163
\weberpage{163}

%% [running head, not body text -- mirrored 2026-08-07] Page furniture (running
%% head, folio, signature) is excluded from this edition by standing policy.
%% \noindent\S.\,50.\hfill Elimination from several equations.\hfill 163

%<<CONT
the number of points of intersection of these three surfaces. But if
$T$ is identically zero, then the three surfaces have either a
portion of surface or a curve in common.

We can determine the degree of the final equation in the following
way, which requires no lengthy considerations at all about the
degrees of the functions $R$, $S$, $T$. In order to make ourselves
independent of exceptional cases, we introduce homogeneous variables
by adding a fourth, $u$, and thus consider three homogeneous
quaternary equations
\weq{4}{F(x, y, z, u) = 0, \quad \Phi(x, y, z, u) = 0, \quad \Psi(x, y, z, u) = 0}
of the orders $m$, $n$, $p$.

If we assume the special case in which each of these three functions
splits into purely linear factors, then it is clear that, if the
number of common solutions is not infinite. % [sic] printed period followed by lowercase "their"
their number must be
equal to the product $mnp$; for we can set one linear factor of each
of $F$, $\Phi$, $\Psi$ equal to zero simultaneously, and so obtain
$mnp$ systems of three linear equations each, which, if they are
independent of one another, yield just as many systems of values for
the ratios $x : y : z : u$.

The conclusion that in the general case the degree of the final
equation of (4) is just as great appears questionable at first
sight. But following a remark of \gesperrt{F. Klein}, the doubt can
be completely removed by the following simple consideration.

Let us first imagine the final equation formed from the equations
(4) under the assumption that the coefficients in the equations (4)
are variables independent of one another. Then the resultant which
results after elimination of $x$, $y$ is certainly not identically
zero, since it does not vanish identically even for special systems
of coefficients, e.g.\ when $F$, $\Phi$, $\Psi$ split into linear
factors; and we obtain a homogeneous final equation for the two
unknowns $z$, $u$ of a certain degree, which is to be determined.

If one specialises the coefficients in any way whatever, the final
equation may indeed possibly be satisfied identically; but if this
case does not occur, then its order cannot be changed (since all
terms are of the same order). Now since, by the specialisation
assumed above, which expresses itself in the splitting of the
functions $F$, $\Phi$, $\Psi$, the
%CONT>>


% ===== scan leaf 189, printed 164 =====
% !PRINTED 164
\weberpage{164}
\pdfbookmark[2]{§ 51. Reducible and irreducible functions}{bm57}
%<<CONT
identical vanishing does not take place, but rather a final equation of order $mnp$ appears, then this must also be the
degree of the final equation in the general case.

With this the \gesperrt{Bezout} theorem is proved for the case of
three homogeneous quaternary equations.

That we have restricted ourselves here to three equations
was done merely in the interest of simplicity. The extension required for the case of $n$ equations in $n + 1$ homogeneous unknowns follows of itself, and
may be left to the reader.

With regard to our general problem we wish only to remark further that from these considerations it follows that the
problem of algebra, in so far as it relates to the solution of any
algebraic equations whatever, is thereby reduced
% German Sperrung starts at "ist" and excludes "damit zurückgeführt"; English keeps
% "is thereby reduced" outside since the verb cluster straddles the boundary
\gesperrt{to the treatment of a chain of equations each in
a single unknown.}

\begin{center}
\S.\,51.

\gesperrt{Decomposable and indecomposable functions.}
\end{center}

Among the integral rational functions of several variables, \gesperrt{decomposable and indecomposable} ones must be
distinguished.

An integral rational function of any variables whatever is called \gesperrt{decomposable} if it can be represented as a product of at least
two integral rational functions of the same variables; if it cannot be represented as such a product,
it is called \gesperrt{indecomposable}.

If a function is decomposable, then the degree of each
of the factors is lower than the degree of the function itself. A
\gesperrt{linear} function is therefore always indecomposable, and \gesperrt{every integral
rational function can be decomposed into a finite number
of indecomposable functions.}

We call an integral function $W$ divisible by another $w$
if a third integral function $w'$ exists such that

\[
W = w\,w'
\]

holds. From this it then follows that, if $U$ is an integral function divisible by $w$
and $V$ is an arbitrary integral function, the product $UV$
%CONT>>


% ===== scan leaf 190, printed 165 =====
% !PRINTED 165
\weberpage{165}
%<<CONT
is divisible by $w$, and that, if $U$, $U'$, $U''\ldots$ are integral rational
functions divisible by $w$, and $V$, $V'$, $V''\ldots$ are arbitrary integral
rational functions, then $UV + U'V' + U''V''\ldots$ is also divisible by $w$.

If $W$ is divisible by $w$, one also says that $w$ goes into $W$.

Two integral rational functions $U$, $V$ that are not divisible by one and the
same integral rational function are called \gesperrt{relatively prime}
or \gesperrt{coprime}.

We prove the theorem:

I. \gesperrt{If $U$, $V$, $v$ are integral rational functions of any
variables whatever, if $U$ and $v$ are relatively prime and $UV$ is divisible
by $v$, then $V$ is divisible by $v$.}

This theorem corresponds exactly to a well-known fundamental theorem from the
theory of the integers, namely that if a product of two integers is divisible
by a third integer which is coprime to one of the factors, then the other
factor must be divisible by this number.

We prove it by complete induction. If $U$, $V$, $v$ depend on only
\gesperrt{one} variable $x$, the theorem is true; for by \S.\,6 (Theorem I)
one can in this case, when $U$ and $v$ are relatively prime, determine two
other integral functions $P$ and $p$ of $x$ so that
\[
PU + pv = 1,
\]
from which, by multiplication by $V$,
\[
PUV + pVv = V
\]
follows, and from this one sees that if $UV$ is divisible by $v$, then $V$
must also be divisible by $v$.

We therefore assume that the theorem we wish to prove has been proved for
functions of $n$ and fewer variables, and from this we derive its validity
for functions of $n + 1$ variables.

For this it is necessary that we draw from Theorem I, assumed to be true,
some consequences, whose conclusion then yields the validity of the theorem
for the next higher number of variables.

If $v$ is an indecomposable function, then another function $U$ of the same
variables is either relatively prime to $v$ or divisible by $v$. From this it
follows by I. that a product of
%CONT>>


% ===== scan leaf 191, printed 166 =====
% !PRINTED 166
\weberpage{166}

%<<CONT
two functions $UV$ can be divisible by $v$ only if one of the two
factors is divisible by $v$. The same holds for a product of
several functions, and so from I. there follows the theorem:

\gesperrt{II. A product of several integral rational functions is divisible by an indecomposable function $v$ only if at least one of the factors of the product is divisible by $v$.}

If an integral rational function $U$ is decomposed in two
ways into indecomposable factors,
\[
U = v\,v'\,v''\dots = w\,w'\,w''\dots,
\]
then by II. at least one of the factors $v, v', v''\dots$ must be
divisible by $w$, say $v$. But then, since $v$ too is indecomposable,
$v$ can differ from $w$ only by a constant factor.

Accordingly, if $c$ is this constant factor,
\[
c\,v'\,v''\cdots = w'\,w''\dots,
\]
from which it follows that one of the functions $v', v''\dots$, say $v'$,
is divisible by $w'$, and thus differs from $w'$ only by a constant
factor.

We therefore obtain, as a second consequence of Theorem I:

\gesperrt{III. An integral rational function can, apart from constant factors, be decomposed into indecomposable factors in only one way.}

From this arises the concept of the \gesperrt{greatest common
divisor} of two or more integral rational functions
$U, V\dots$ By this one understands the product of all
indecomposable factors that occur in the decompositions of each of the
functions $U, V\dots$, or the function of highest possible
degree that divides all the functions $U, V\dots$ By III. this
function is, apart from a constant factor, completely determined
for every system of functions $U, V\dots$

Several functions $U, V, W\dots$ are called relatively prime if
no two of them have a common divisor.

All these definitions and theorems are exactly analogous to
very well-known theorems of elementary number theory, which there
appear as consequences of the algorithm of the greatest common
%CONT>>


% ===== scan leaf 192, printed 167 =====
% !PRINTED 167
\weberpage{167}

%<<CONT
divisor occur, only that here the indecomposable functions
take over the role of the primes.

We add one more theorem of this kind:

\gesperrt{IV. If $u$, $v$ are relatively prime and $Uu$, $Uv$ are
divisible by $w$, then $U$ is divisible by $w$.}

For if one decomposes the integral rational functions $U, u, v, w$
into their indecomposable factors, then any factor of
$w$, since it cannot occur in $u$ and $v$ at the same time, must
divide $U$. If one removes it from $U$ and $w$, one can conclude
in the same way for a next factor of $w$, and so on.

We now consider integral rational functions of one vari\-
able $t$
\[
f(t) = u_0 t^m + u_1 t^{m-1} + \cdots + u_{m-1} t + u_m,
\]
whose coefficients are integral rational functions of $n$ vari\-
ables $x$, of which $t$ is independent.

If the coefficients $u_0, u_1 \ldots u_m$ have no common
divisor, then $f(t)$ is called \gesperrt{primitive}; in the other case the
greatest common divisor of the coefficients $u_0, u_1 \ldots u_m$
is called \gesperrt{the divisor of the function} $f(t)$ (cf. §.\,2).

We conclude, always under the assumption of the validity
of I:

\gesperrt{V. The product of two primitive functions $f(t)$
and $\varphi(t)$ is again a primitive function, and
the divisor of a product of two imprimitive
functions is equal to the product of the divisors
of the two factors.}

The proof is exactly as for the corresponding theorem
in §.\,2.

Let
\weq{1}{
\begin{aligned}
f(t) &= u_0 t^m + u_1 t^{m-1} + \cdots + u_m \\
\varphi(t) &= v_0 t^\mu + v_1 t^{\mu-1} + \cdots + v_\mu
\end{aligned}
}
be two primitive functions and
\weq{2}{
F(t) = U_0 t^{m+\mu} + U_1 t^{m+\mu-1} + \cdots + U_{m+\mu}
}
their product. Then, if $r$ and $s$ denote any two numbers from
the series $0, 1 \ldots m$, and $0, 1, 2 \ldots \mu$ (§.\,2), we have
\weq{3}{
\begin{aligned}
U_{r+s} &= u_r v_s + u_{r-1} v_{s+1} + \cdots \\
&\qquad + u_{r+1} v_{s-1} + \cdots
\end{aligned}
}
%CONT>>


% ===== scan leaf 193, printed 168 =====
% !PRINTED 168
\weberpage{168}

Now if $w$ is any indecomposable function which divides neither
all the $u_r$ nor all the $v_s$, then we choose $r$ and $s$ in (3)
so that $u_r$ is the first $u$ not divisible by $w$, so that
$u_{r-1}, u_{r-2}\dots$ are divisible by $w$, and likewise so that $v_s$ is the
first $v$ not divisible by $w$. Then $U_{r+s}$ also cannot
be divisible by $w$, because all its terms, with the exception
of the first, are divisible by $w$; that is, $F(t)$ is primitive.

From this it follows immediately, if $p$ and $q$ are any integral
rational functions of the $x$, that $pq$ is the divisor of the product
of the two imprimitive functions $pf(t)$, $q\varphi(t)$, and hence
the second part of Theorem V. From this it follows further:

\gesperrt{VI. If an integral rational function. $F(t)$ of the $n+1$ variables $x$ and $t$ is decomposable into two factors which are integral with respect to $t$ and at least rational with respect to the $x$, then it is also decomposable into two functions which are integral and rational in $x$ and $t$.} % [sic] printed period after "function" before F(t)

For by the hypothesis there is an integral rational
function $w$ of the $x$ alone and two integral rational functions
of the $x$ and $t$, $f_1(t)$, $\varphi_1(t)$, such that
\[
w\,F(t) = f_1(t)\,\varphi_1(t),
\]
and by (V) $w$ must divide the product of the divisors of $f_1(t)$ and
$\varphi_1(t)$, so that we also obtain
\[
F(t) = f(t)\,\varphi(t)
\]
where $f(t)$, $\varphi(t)$ are likewise integral rational functions
of $x$ and $t$, and indeed with respect to $t$ of the same degree as
$f_1(t)$ and $\varphi_1(t)$, q. e. d. --- From this we conclude further
that two functions $F(t)$, $f(t)$ which, regarded as integral rational functions
of the $n+1$ variables $x$ and $t$, are relatively prime in the sense
defined above, must also prove to be relatively prime when regarded
as functions of $t$ alone and treated by the algorithm of the greatest
common divisor (\S.\,6). For if they had a common divisor which
was integral with respect to $t$ but fractional with respect to the $x$,
then an integral function $T$ of $x$ and $t$ and an integral function
$P$ of the $x$ alone could be determined so that
\[
P\,F(t) = T\,F_1(t), \qquad P\,f(t) = T\,f_1(t)
\]
would hold, where $F_1, f_1$ are integral functions without a common divisor
%CONT>>


% ===== scan leaf 194, printed 169 =====
% !PRINTED 169
\weberpage{169}

%<<CONT
are. By IV, $P$ would then have to divide the divisor of $T$, and
$F(t)$ and $f(t)$ could not be relatively prime. By \S.\,6 we can
therefore determine two integral rational functions $Q$, $q$ of $x$
and $t$ and an integral rational function $X$ of the $x$ alone in such a way
that the identity

\weq{4}{QF(t) + qf(t) = X}
holds.

Now if $\Phi(t)$ is a further integral rational function of $x$
and $t$, then from (4), by multiplication with $\Phi(t)$, it follows that
\[
Q\Phi(t)F(t) + q\Phi(t)f(t) = X\Phi(t),
\]
and hence, if $\Phi(t)F(t)$ is divisible by $f(t)$, then by this
$X\Phi(t)$ is also divisible by $f(t)$.

Accordingly, if $\varphi(t)$ and $\psi(t)$ again denote two integral
functions of $x$ and $t$, we may set

\weq{5}{\begin{aligned} X\Phi(t) &= \varphi(t)f(t) \\ F(t)\Phi(t) &= \psi(t)f(t). \end{aligned}}

If one multiplies the second of these equations by $X$ and
substitutes from the first the expression $\varphi(t)f(t)$ for $X\Phi(t)$, then
$f(t)$ can be cancelled and it follows that
\[
\varphi(t)F(t) = X\psi(t).
\]
By IV., $X$ must therefore divide both the divisor of $\varphi(t)f(t)$ and
the divisor of $\varphi(t)F(t)$, and since the functions $f(t)$
and $F(t)$, and consequently also their divisors, are relatively prime, $X$ must
divide the divisor of $\varphi(t)$. If we accordingly set
\[
\varphi(t) = X\varphi_1(t),
\]
then from (5) it follows that
\[
\Phi(t) = \varphi_1(t)f(t),
\]
that is, $\Phi(t)$ is divisible by $f(t)$; hence:

\gesperrt{VII. If $F(t)$ and $f(t)$ are relatively prime and $\Phi(t)F(t)$
is divisible by $f(t)$, then $\Phi(t)$ is divisible by $f(t)$.}

But this is nothing other than Theorem I for functions
of $n + 1$ variables, and I. is thus proved in general.


% ===== scan leaf 195, printed 170 =====
% !PRINTED 170
\weberpage{170}
\pdfbookmark[2]{§ 52. Tschirnhausen transformation}{bm58}

\begin{center}
§. 52.

\gesperrt{The Tschirnhausen Transformation.}
\end{center}

We now make one more important application of the theory of symmetric functions and of the formation of resultants, to the transformation of algebraic equations that goes back to \gesperrt{Tschirnhausen}\textsuperscript{1)}.

Let
\weq{1}{f(x) = x^n + a_1 x^{n-1} + a_2 x^{n-2} + \cdots + a_{n-1} x + a_n}
be any function of degree $n$, and let
\weq{2}{y = \frac{\chi(x)}{\psi(x)}}
be any (integral or fractional) rational function of $x$, subject to the single assumption that $f(x)$ and $\psi(x)$ \gesperrt{are to have no common divisor}. By the concluding theorem of \S.\,6 we can then determine the functions $F(x)$ and $\varphi(x)$, the latter of degree at most $n-1$, so that
\weq{3}{\chi(x) = F(x) f(x) + \varphi(x) \psi(x)}
holds; consequently, as soon as a root of the equation
\weq{4}{f(x) = 0}
is put for $x$, we have
\[
\frac{\chi(x)}{\psi(x)} = \varphi(x). % [sic] period as printed before the following continuation
\]
Since we shall consider the values of $y$ only for those values of $x$ that are roots of (4), we lose nothing in generality if from the outset we set
\[
y = \varphi(x),
\]
that is, equal to an integral function of degree $(n-1)$. Let us therefore put
\weq{5}{y = \varphi(x) = \alpha_0 + \alpha_1 x + \alpha_2 x^2 + \cdots + \alpha_{n-1} x^{n-1},}
in which the coefficients $\alpha$ remain for the present entirely undetermined.

If we denote by
\weq{6}{x_1, x_2 \ldots x_n}
%CONT>>

\par\smallskip
{\small 1) Acta eruditorum, Leipzig 1683.}\par


% ===== scan leaf 196, printed 171 =====
% !PRINTED 171
\weberpage{171}

%<<CONT
the roots of the equation (4), then the corresponding values of $y$ result from (5):
\weq{7}{y_1 = \varphi(x_1), \quad y_2 = \varphi(x_2) \ldots y_n = \varphi(x_n),}
and the product
\weq{8}{\Phi(y) = (y - y_1)(y - y_2) \ldots (y - y_n)}
is a \gesperrt{symmetric function of the quantities (6), which can therefore be expressed rationally in terms of the coefficients of $f(x)$.}

Represented in this way, let $\Phi(y)$ have the expression
\weq{9}{\Phi(y) = y^n + b_1 y^{n-1} + b_2 y^{n-2} + \cdots + b_{n-1} y + b_n.}

The coefficients $b_\nu$ are, by \S.\,7, determined through the $y_i$ by means of the formulas
\[
b_1 = -\Sigma y_1, \quad b_2 = \Sigma y_1 y_2, \quad b_3 = -\Sigma y_1 y_2 y_3 \ldots,
\]
\gesperrt{from which it follows that $b_\nu$ is an integral rational and homogeneous function of degree $\nu$ of the variables $\alpha_0, \alpha_1 \ldots \alpha_{n-1}$.}

$\Phi(y)$ is nothing other than the resultant of the two equations
\weq{10}{f(x) = 0, \quad \varphi(x) - y = 0,}
and by \S.\,49 one can also, conversely, express the common root of these two equations (10), when $y$ becomes equal to one of the values $y_1, y_2 \ldots y_n$, rationally in terms of the coefficients of (10), that is, rationally in terms of $a_i$, $\alpha_i$, $y$, provided that \gesperrt{only one such common root exists.}

This will occur when the $n$ values $y_1, y_2 \ldots y_n$ are different from one another. One then obtains an expression
\[
x = \beta_0 + \beta_1 y + \beta_2 y^2 + \cdots + \beta_{n-1} y^{n-1},
\]
which is valid when one of the values (6) is put for $x$ and the corresponding value (7) for $y$. The quantities $\beta_i$ are composed rationally of the $a_i$ and $\alpha_i$.

But if among the values $y_1, y_2 \ldots y_n$ several equal ones occur, then the determination of the corresponding values of $x$ still requires the solution of an equation of the corresponding degree. \gesperrt{This degree, however, is always smaller than $n$}, since the function $\varphi(x)$ can take the same value $y$ for at most $n-1$ values of $x$. If, then, the equation
% [source-original mathematical caveat, confirmed 2026-08-25 against scan p. 197 / printed p. 171]
% The assertion omits the constant-phi case, in which all n roots x_i can have the same y-value.
% The facing translation preserves the printed claim rather than silently qualifying it.
\weq{11}{\Phi(y) = 0}
%CONT>>


% ===== scan leaf 197, printed 172 =====
% !PRINTED 172
\weberpage{172}

%<<CONT
is completely solved, then the equation (1) is thereby also completely solved [if (11) has $n$ distinct roots], or at least reduced to the solution of an equation of lower degree [if (11) has equal roots].

Accordingly we regard (11) as a transformation of the equation (1), and the formation of (11) from (1) is called the \gesperrt{Tschirnhausen transformation} of the equation (1).

The purpose of such a transformation is to determine the coefficients $\alpha_i$ in such a way that $\Phi(y)$ takes a simpler form than $f(x)$.

One can form the coefficients $b_\nu$ if one knows the power sums of the $y_i$ (\S.\,42).

If we denote the sums of the $m$th powers of the $x_i$ by $s_m$, and those of the $y_i$ by $\sigma_m$, that is
\[
s_m = S x_i^m, \qquad \sigma_m = S y_i^m,
\]
then we can form $\sigma_m$ from the $s_m$ in the following way.

We first expand $y^m$ by the polynomial theorem (\S.\,12), that is, we set
% [sic] The print omits the plus sign before \cdots in the following exponent.
\weq{12}{
y^m = \sum^{\nu} \frac{\Pi(m)}{\Pi(\nu_0)\,\Pi(\nu_1)\,\dots\,\Pi(\nu_{n-1})}\, \alpha_0^{\nu_0}\alpha_1^{\nu_1}\dots\alpha_{n-1}^{\nu_{n-1}}\, x^{\nu_1+2\nu_2\cdots+(n-1)\nu_{n-1}},
}
where the sum extends over all non-negative integral values of $\nu_0, \nu_1 \dots \nu_{n-1}$ that satisfy the condition
\weq{13}{
\nu_0 + \nu_1 + \cdots + \nu_{n-1} = m
}
If we then substitute $x_1, x_2 \dots x_n$ for $x$ in (12) and form the sum of the values of $y_1^m, y_2^m \dots y_n^m$ so obtained, it follows that
\weq{14}{
\sigma_m = \sum^{\nu} \frac{\Pi(m)}{\Pi(\nu_0)\,\Pi(\nu_1)\,\dots\,\Pi(\nu_{n-1})}\, s_{\nu_1+2\nu_2+\cdots+(n+1)\nu_{n-1}}\, \alpha_0^{\nu_0}\alpha_1^{\nu_1}\dots\alpha_{n-1}^{\nu_{n-1}}. % [sic] exponent reads (n+1), cf. (n-1) in (12)
}
In this way $\sigma_m$ is represented as an integral homogeneous function of order $m$ of the $\alpha_0, \alpha_1 \dots \alpha_{n-1}$. A somewhat simplified notation is obtained by the second representation of forms (\S.\,15). If we let $\mu_1, \mu_2 \dots \mu_m$ run independently of one another through the sequence of numbers $0, 1 \dots n-1$, then, if in a combination the index 0 occurs $\nu_0$ times, the index 1 occurs $\nu_1$ times, the index 2 occurs $\nu_2$ times, \dots the index $n-1$ occurs $\nu_{n-1}$ times, we have
\[
\nu_0 + \nu_1 + \nu_2 + \cdots + \nu_{n-1} = m
\]
and
%CONT>>


% ===== scan leaf 198, printed 173 =====
% !PRINTED 173
\weberpage{173}
\pdfbookmark[2]{§ 53. Application to the cubic and biquadratic equations}{bm59}

\[
\mu_1 + \mu_2 + \cdots + \mu_m = \nu_1 + 2\nu_2 + \cdots + (n-1)\nu_{n-1},
\]
and accordingly the expression for $\sigma_m$ can be represented thus:
\weq{15}{
\sigma_m = \overset{\mu}{\textstyle\sum} s_{\mu_1+\mu_2+\ldots+\mu_m}\, \alpha_{\mu_1}\alpha_{\mu_2}\cdots\alpha_{\mu_m}.
}
For example, we obtain
\[
\sigma_1 = \alpha_0 s_0 + \alpha_1 s_1 + \cdots + \alpha_{n-1} s_{n-1}
\]
\weq{16}{
\begin{aligned}
\sigma_2 &= \overset{i,k}{\textstyle\sum} s_{i+k}\, \alpha_i \alpha_k \\
\sigma_3 &= \overset{h,i,k}{\textstyle\sum} s_{h+i+k}\, \alpha_h \alpha_i \alpha_k,
\end{aligned}
}
in which $h$, $i$, $k$ run through the sequence of numbers $0, 1, 2 \ldots n-1$.

\begin{center}
\S.\,53.

\gesperrt{Application to the cubic and biquadratic equations.}
\end{center}

The goal that \gesperrt{Tschirnhausen} already had in view with this transformation consisted in making so many terms of the transformed equation $\Phi(y) = 0$ vanish, by determining the substitution coefficients $\alpha_0, \alpha_1 \ldots \alpha_{n-1}$, that the equation becomes solvable. This succeeds easily for the equations of the third and fourth degree, though it is not simple, and does not seem possible without lengthy calculations, to derive the otherwise well-known formulas in this way.

Let us first take the cubic equation
\weq{1}{
x^3 + a x^2 + b x + c = 0,
}
which we transform by means of the following substitution:
\weq{2}{
y = \alpha + \beta x + \gamma x^2
}
In order to reduce the equation for $y$
\weq{3}{
y^3 + A y^2 + B y + C = 0
}
to a pure cubic equation, we have to determine the ratios $\alpha:\beta:\gamma$ from the two equations
\weq{4}{
A = 0, \quad B = 0
}
of which the first is linear, the second quadratic.

The equations (4) are equivalent to
\weq{5}{
\sigma_1 = 0, \quad \sigma_2 = 0
}
%CONT>>


% ===== scan leaf 199, printed 174 =====
% !PRINTED 174
\weberpage{174}

%<<CONT
and thus, by (16) of the preceding section, they yield
\[
\alpha s_0 + \beta s_1 + \gamma s_2 = 0,
\]
\[
\alpha^2 s_0 + 2\alpha\beta s_1 + 2\alpha\gamma s_2 + \beta^2 s_2 + 2\beta\gamma s_3 + \gamma^2 s_4 = 0,
\]
in which $s_0 = 3$. If one multiplies the latter equation by $s_0$ and subtracts from it the square of the first, it follows that
\weq{6}{
\beta^2(s_0 s_2 - s_1^2) + 2\beta\gamma(s_0 s_3 - s_1 s_2) + \gamma^2(s_0 s_4 - s_2^2) = 0.
}
The discriminant of this quadratic equation
\[
(s_0 s_3 - s_1 s_2)^2 - (s_0 s_2 - s_1^2)(s_0 s_4 - s_2^2)
\]
gives, when expanded,
\[
-s_0(s_0 s_2 s_4 - s_0 s_3^2 - s_1^2 s_4 - s_2^3 + 2 s_1 s_2 s_3),
\]
and is therefore, if $D$ denotes the discriminant of the cubic equation (1) [\S.\,46, (9)], equal to $-3D$; and the solution of (6) accordingly gives
\weq{7}{
\frac{\beta}{\gamma} = \frac{-(s_0 s_3 - s_1 s_2) + \sqrt{-3D}}{s_0 s_2 - s_1^2}.
}
With this, equation (3) has been reduced to a pure equation, which for its solution requires only the extraction of a cube root. Which sign we wish to give the square root occurring in (7) is left to our discretion.

To make use of the \gesperrt{Tschirnhausen} transformation for the biquadratic equation, one may proceed, for instance, by determining the $\alpha_0, \alpha_1, \alpha_2, \alpha_3$ in the transformed equation
\weq{8}{
y^4 + b_1 y^3 + b_2 y^2 + b_3 y + b_4 = 0
}
so that
\weq{9}{
b_1 = 0, \quad b_3 = 0
}
whereby (8) passes into a quadratic equation for $y^2$, from whose roots the square roots must still be extracted, so that, after the equations (9) have been solved, a square root must still be extracted twice more.

If, by means of the first equation (9), one eliminates $\alpha_0$ from the second, there arises a homogeneous cubic equation among $\alpha_1, \alpha_2, \alpha_3$. One may therefore assume one of the two ratios $\alpha_1:\alpha_2:\alpha_3$ arbitrarily, e.g. set $\alpha_3 = 0$, and obtains for the determination of the other a \gesperrt{cubic equation}. By this, then, the solution of the biquadratic equation is reduced to that of a cubic and to square roots.


% ===== scan leaf 200, printed 175 =====
% !PRINTED 175
\weberpage{175}
\pdfbookmark[2]{§ 54. The Tschirnhausen transformation of the equation of the 5th degree}{bm60}

\begin{center}
§. 54.

\gesperrt{The Tschirnhausen Transformation of the Equation of the Fifth Degree.}
\end{center}

If one wishes to apply the Tschirnhausen transformation to the equation of the fifth degree
\weq{1}{x^5 + a_1 x^4 + a_2 x^3 + a_3 x^2 + a_4 x + a_5 = 0}
one has to set
\weq{2}{y = \alpha_0 + \alpha_1 x + \alpha_2 x^2 + \alpha_3 x^3 + \alpha_4 x^4}
and obtains the transformed equation
\weq{3}{y^5 + b_1 y^4 + b_2 y^3 + b_3 y^2 + b_4 y + b_5 = 0.}

The most obvious thought would now be to determine the ratios of the
five unknowns $\alpha$ in such a way that
\weq{4}{b_1 = 0, \quad b_2 = 0, \quad b_3 = 0, \quad b_4 = 0}
would hold, which would reduce (3) to a pure equation; and
this was no doubt also the goal that \gesperrt{Tschirnhausen} had in
view. But of the four equations (4) the first is linear,
the second quadratic, the third of the third and the fourth of the
fourth degree. The final equation that one can derive from them
will therefore (by \S.\,50) be of degree $2.3.4 = 24$, and hence
no use can be drawn from it for the solution of the equation of the
fifth degree. One therefore seeks at first to satisfy only the
two equations
\weq{5}{b_1 = 0, \quad b_2 = 0,}
whereby the equation (3) passes into a form which, following
\gesperrt{F. Klein} (Lectures on the Icosahedron), is called a \gesperrt{principal equation} of the fifth degree.

To satisfy the two equations (5) we have at our
disposal the four ratios $\alpha_0 : \alpha_1 : \alpha_2 : \alpha_3 : \alpha_4$, and we
can still use these quantities in manifold ways to simplify the equation of the fifth degree.

We shall return to this subject in more detail in a later
section, and confine ourselves here to indicating the
aims in general terms. For the sake of vividness
we make use of a geometrical dress, which, incidentally,
is by no means essential to the understanding of the algebraic theory
as it is to be given later.


% ===== scan leaf 201, printed 176 =====
% !PRINTED 176
\weberpage{176}

If, with the help of the linear equation $b_1 = 0$, we
express one of the five quantities $\alpha$, say $\alpha_0$, in terms of the others, then we can regard the four remaining ones $\alpha_1, \alpha_2, \alpha_3, \alpha_4$ as homogeneous coordinates of a point in space. The equation
$b_2 = 0$ then represents a surface of the second order, $b_3 = 0$
a surface of the third order, $b_4 = 0$ a surface of the fourth order, $b_5 = 0$ a surface of the fifth order.

If one wished, then, to make $b_2$, $b_3$, $b_4$ vanish simultaneously and
thereby reduce the given equation to a pure one, one
would have to determine one of the 24 points of intersection of three surfaces
of the second, third, and fourth order, which depends on an equation of degree $24$.

Instead of proceeding in this way, one seeks to determine a straight line on the surface
$b_2 = 0$. Of such straight lines
there are on every surface of the second order one or two
families (real or imaginary), and one of these lines can be determined
\gesperrt{by square roots}, as will be carried out further below. The determination of a point of intersection of one
of these straight lines with the surface $b_3 = 0$ then leads to
an equation of the third degree, and if, by determining a common factor of the $\alpha$, one further makes the coefficient
$b_4 = 1$, the equation for $y$ takes the form
\weq{6}{y^5 + y + b = 0,}
whereby $y$ is determined as a function of \gesperrt{one} variable $b$.

This form of the equation is usually named after the English
mathematician \gesperrt{Jerrard}, who made it known in the year 1834.
But it had already been published much earlier (1786) by \gesperrt{E. J. Bring}
in Lund\footnote{Cf. F. Klein, Vorlesungen über das Ikosaëder, p. 143.}. We therefore call it the \gesperrt{Bring-Jerrard} form.

Just as well one could make $b_4 = 0$ by an equation of the fourth degree,
and would obtain a second normal form of the equation
of the fifth degree:
\weq{7}{y^5 + y^2 + b = 0,}
which has the same justification as the \gesperrt{Bring-Jerrard}
form.

For the solution of the equation of the fifth degree, to be sure, nothing
is gained by these considerations. Their usefulness consists
%CONT>>


% ===== scan leaf 202, printed 177 =====
% !PRINTED 177
\weberpage{177}

in this, that the equation of the fifth degree is reduced to a normal form which depends on only one variable coefficient. This can be achieved in infinitely many different ways, and is accomplished most simply, namely without the solution of an equation of the third or fourth degree, if one regards the equation of the fifth degree $b_5 = 0$ itself as a transformation of the given equation of the fifth degree. For if one has solved the equation $b_5 = 0$, then one of the five values of $y$ becomes equal to 0, and the given equation of the fifth degree is thereby solved at the same time.

% Printer's furniture excluded by policy: the footer reads "Weber, Algebra. I." with gathering number 12; the remaining faint show-through is not legible source text.


% ===== scan leaf 203, printed 178 =====
\weberpage{178}% page number not printed on this page
\pdfbookmark[1]{Fifth Section. Linear transformation. Invariants.}{bm61}
\pdfbookmark[2]{§ 55. Linear transformation}{bm62}
% !PRINTED none

\begin{center}
\gesperrt{Fifth Section.}

\gesperrt{Linear Transformation.}

\rule{0.3\linewidth}{0.4pt}

\S.\,55.

\gesperrt{Introduction of the linear transformation.}
\end{center}

If, in any functions that depend on the $m$ variables $x_1, x_2 \ldots x_m$, we substitute for these variables linear expressions that depend on $m$ new variables $x_1', x_2' \ldots x_m'$
\weq{1}{\begin{aligned}
x_1 &= \alpha_1^{(1)} x_1' + \alpha_1^{(2)} x_2' + \cdots + \alpha_1^{(m)} x_m', \\
x_2 &= \alpha_2^{(1)} x_1' + \alpha_2^{(2)} x_2' + \cdots + \alpha_2^{(m)} x_m', \\
& \cdots\cdots\cdots\cdots\cdots\cdots\cdots\cdots \\
x_m &= \alpha_m^{(1)} x_1' + \alpha_m^{(2)} x_2' + \cdots + \alpha_m^{(m)} x_m',
\end{aligned}}
then we call this a \gesperrt{linear substitution}, and the result a \gesperrt{linear transformation}.

The quantity
\weq{2}{r = \begin{vmatrix}
\alpha_1^{(1)}, & \alpha_1^{(2)} & \ldots & \alpha_1^{(m)} \\
\alpha_2^{(1)}, & \alpha_2^{(2)} & \ldots & \alpha_2^{(m)} \\
\multicolumn{4}{c}{\cdot\ \cdot\ \cdot\ \cdot\ \cdot} \\
\alpha_m^{(1)}, & \alpha_m^{(2)} & \ldots & \alpha_m^{(m)}
\end{vmatrix}}
is called the \gesperrt{substitution determinant}. It must always be different from zero, since otherwise a dependence among the variables $x_1, x_2 \ldots x_m$ would be expressed by (1).

By a linear substitution every homogeneous function of the variables $x$ passes over into a homogeneous function of the same degree in the variables $x'$.


% ===== scan leaf 204, printed 179 =====
% !PRINTED 179
\weberpage{179}
\pdfbookmark[2]{§ 56. Quadratic forms}{bm63}

Thus, for example, every linear function
\[
y = \overset{i}{\Sigma} a_i x_i
\]
passes into a function of the same kind
\[
y = \overset{i}{\Sigma} a_i' x_i'
\]
in which
\weq{3}{a_k' = \overset{i}{\Sigma} a_i \alpha_i^{(k)}.}
If we consider $m$ such linear functions
\[
y_\nu = \overset{i}{\Sigma} a_{i,\nu} x_i \qquad \nu = 1, 2 \ldots m,
\]
we obtain $m$ transformed functions
\[
y_\nu = \overset{i}{\Sigma} a_{i,\nu}' x_i' \qquad \nu = 1, 2 \ldots m.
\]
The determinants of the systems $y_\nu, y_\nu'$
\[
\Delta = \Sigma \pm a_{1,1} a_{2,2} \ldots a_{m,m}, \qquad \Delta' = \Sigma \pm a_{1,1}' a_{2,2}' \ldots a_{m,m}'
\]
stand, by the multiplication law of determinants, in the relation
\weq{4}{\Delta' = r \Delta,}
so that the one does not vanish without the other.

\begin{center}
\S.\,56.

\gesperrt{Quadratic Forms.}
\end{center}

Of particular importance is the transformation of the homogeneous functions of the second degree, or the quadratic forms. We denote them, as already in the first section, thus:
\weq{1}{\varphi(x_1, x_2 \ldots x_m) = \overset{i,k}{\Sigma} a_{i,k} x_i x_k, \qquad a_{i,k} = a_{k,i}.}
By the substitution \S.\,55, (1), $\varphi$ passes into
\weq{2}{\Phi(x_1', x_2' \ldots x_m') = \overset{i,k}{\Sigma} a_{i,k}' x_i' x_k'.}
If we put $x + \xi$ for $x$ and likewise $x' + \xi'$ for $x'$, so that the $\xi'$ stand in the same relation to the $\xi$ as the $x'$ to the $x$, then, if for abbreviation we set
\[
\tfrac{1}{2}\varphi'(x_i) = u_i, \qquad \tfrac{1}{2}\Phi'(x_i') = u_i'
\]
we obtain, by comparison of the terms of the first dimension (\S.\,16):
%CONT>>


% ===== scan leaf 205, printed 180 =====
% !PRINTED 180
\weberpage{180}
\weq{3}{\Sigma \xi_\nu u_\nu = \Sigma \xi'_\nu u'_\nu,}
and therefore, as in \S.\,55, (3)
\weq{4}{u'_k = \overset{i}{\Sigma} \alpha_i^{(k)} u_i.}
Now
\weq{5}{u_k = \overset{i}{\Sigma} a_{k,i} x_i, \qquad u'_k = \overset{i}{\Sigma} a'_{k,i} x'_i,}
and accordingly (4) gives us the transformation of $m$ linear
functions. The coefficients of the given system are $a'_{k,h}$,
and those of the transformed system, which one obtains by substituting
the first system (5) into (4), are
\weq{6}{b_{k,h} = \overset{i}{\Sigma} \alpha_i^{(k)} a_{i,h}.}
The determinant of the former system
\[
H' = \Sigma \pm a'_{1,1} a'_{2,2} \ldots a'_{m,m}
\]
is, by \S.\,55, (4),
\[
= r \Sigma \pm b_{1,1} b_{2,2} \ldots b_{m,m},
\]
and by (6) the determinant of the $b_{k,h}$ is
\[
= r \Sigma \pm a_{1,1} a_{2,2} \ldots a_{m,m}.
\]
If, then, we set
\weq{7}{H = \Sigma \pm a_{1,1} a_{2,2} \ldots a_{m,m}}
we have the important theorem
\weq{8}{H' = r^2 H.}
The determinant $H$ is called the \gesperrt{determinant of the form}
$\varphi(x)$.

Thus under linear transformation of the function $\varphi$ it
changes only by the factor $r^2$, and it cannot vanish for the
transformed form if it does not do so for the original one.

\gesperrt{The vanishing of the determinant $H$ has the meaning
that the function $\varphi(x)$ can be regarded as
a function of fewer than $m$ variables, which
are linear functions of the $x$.}

For if, in the first place, $\varphi(x)$ passes by a linear transformation
into the function $\Phi(x')$, which is independent of one of the variables,
say of $x'_1$, then $H'$ vanishes, since the elements
in one of the rows (and columns) all vanish, and hence
by (8) also $H = 0$.

In the second place, however, $H$ is the determinant of the system of linear
equations
%CONT>>


% ===== scan leaf 206, printed 181 =====
% !PRINTED 181
\weberpage{181}
\pdfbookmark[2]{§ 57. Transformation of quadratic forms into a sum of squares}{bm64}

\weq{9}{\varphi'(x_i) = 0 \qquad i = 1, 2, \ldots m,}

and if this determinant vanishes, then there is (\S.\,24, III) a system of quantities

\weq{10}{x_1^{(0)}, x_2^{(0)} \ldots x_m^{(0)},}

whose elements do not all vanish, and which, substituted for the $x$, satisfy the equations (9).

But now, if $x_i$, $x_i'$ are any two systems of quantities, and $\lambda$ likewise denotes an arbitrary quantity, one has the identical equations

\weq{11}{\begin{aligned}
\varphi(x + \lambda x') &= \varphi(x) + \lambda \overset{i}{\Sigma} x_i \varphi'(x_i') + \lambda^2 \varphi(x')\\
\varphi(x') &= \tfrac{1}{2}\Sigma x_i' \varphi'(x_i'),
\end{aligned}}

[\S.\,16, (12) to (15)], so that, if for $x_i'$ the values $x_i^{(0)}$ are substituted,

\weq{12}{\varphi(x_i + \lambda x_i^{(0)}) = \varphi(x),}

becomes entirely independent of $\lambda$. If we now assume that $x_1^{(0)}$ is different from zero, and we set

\[
\lambda = -\frac{x_1}{x_1^{(0)}},
\]

and further:

\weq{13}{\begin{aligned}
y_2 &= x_2 x_1^{(0)} - x_1 x_2^{(0)}\\
y_3 &= x_3 x_1^{(0)} - x_1 x_3^{(0)}\\
&\cdots\cdots\cdots\cdots\\
y_m &= x_m x_1^{(0)} - x_1 x_m^{(0)},
\end{aligned}}

then it follows that

\weq{14}{x_1^{(0)\,2} \varphi(x) = \varphi(0, y_2, y_3 \ldots y_m),}

whereby $\varphi(x)$ is expressed in terms of the $m - 1$ variables $y_2, y_3 \ldots y_m$.

\begin{center}
\S.\,57.

\gesperrt{Transformation of Quadratic Forms into a\\
Sum of Squares.}
\end{center}

We at once make an important application of the theorem just proved to the discriminant of the function (11) (regarded as a function of the second degree in $\lambda$), in which
%CONT>>


% ===== scan leaf 207, printed 182 =====
% !PRINTED 182
\weberpage{182}
%<<CONT
it is no longer assumed that the determinant of the function $\varphi$
vanishes. This discriminant is
\weq{1}{\psi(x) = [\Sigma x_i \varphi'(x'_i)]^2 - 4\varphi(x)\varphi(x').}
If we put $x_i + \lambda x'_i$ in place of $x_i$ in it, then, substituting
for $\varphi(x + \lambda x')$ the expression \S.\,56, (11)
and setting $\Sigma x'_i \varphi'(x'_i) = 2\varphi(x')$, we obtain
\[
\psi(x + \lambda x') = \psi(x),
\]
from which it follows that $\psi(x)$ depends only on $m - 1$ linear combinations
of the $m$ variables. If one puts for $x'_1, x'_2 \ldots x'_m$
arbitrary special values for which $\varphi(x')$ does not vanish, and
takes $x'_1$ different from zero, then
% [sic] The print has x_n; the variables otherwise run through x_m.
\[
x_1'^2 \psi(x_1, x_2, \ldots x_n) = \psi(0, x_2 x'_1 - x_1 x'_2, \ldots x_m x'_1 - x_1 x'_m),
\]
hence a function of the $m - 1$ variables
\[
y_2 = x_2 x'_1 - x_1 x'_2, \ldots y_m = x_m x'_1 - x_1 x'_m.
\]
One can also confirm this by observing that
\[
\psi'(x_1), \ \psi'(x_2) \ldots \psi'(x_m)
\]
all vanish simultaneously for $x_1, x_2 \ldots x_m = x'_1, x'_2 \ldots x'_m$,
so that the determinant of the function $\psi$ must vanish.

Formula (1) now contains a very important result, which
can best be surveyed if we give this formula the following
form.

We set
\weq{2}{\begin{aligned}
\psi(x) &= -4\varphi(x')\chi(y_2, \ldots y_m) \\
\Sigma x_i \varphi'(x'_i) &= 2\sqrt{\varphi(x')}\, Y_1,
\end{aligned}}
and obtain
\weq{3}{\varphi(x) = Y_1^2 + \chi(y_2, y_3 \ldots y_m),}
whereby $\varphi(x)$ is represented as the sum of a square of a
% [sic] The print has lowercase y_1 here, while (2) and (3) use Y_1.
linear function $y_1$ and a function $\chi$ of $m - 1$ variables.

But from this follows the important theorem

\gesperrt{A quadratic form of $m$ variables can always be
represented as a sum of squares of $m$ or fewer
linear functions.}

This theorem is evidently true for a function of one
variable, which is itself a square. But from (3) the truth
of the theorem follows for a function of $m$ variables,
provided its truth is assumed for a function of $m - 1$
variables.


% ===== scan leaf 208, printed 183 =====
% !PRINTED 183
\weberpage{183}
\pdfbookmark[2]{§ 58. Law of inertia of quadratic forms}{bm65}

The representation is, as (2) shows, possible in infinitely many different ways, since the entirely arbitrary quantities $x'$ still occur in this formula.

A representation by fewer than $m$ squares is possible only when the determinant of the function $\varphi$ vanishes.

Since a square root occurs in the expression (2) for $Y_1$, $Y_1$ can be real or imaginary, even when the coefficients of $\varphi$ and the $x$ and $x'$ are assumed to be real.

But if, in the case where $Y_1$ is imaginary, we put $iY_1$ in place of $Y_1$, we can also state our theorem as follows:

\gesperrt{A real quadratic form in $m$ variables can be represented as a sum of at most $m$ positive or negative squares of real linear functions of the $x$}\footnote{For the general theory of the linear transformation of forms of the second degree, besides the investigations of \gesperrt{Jacobi} and \gesperrt{Hesse} (\gesperrt{Hesse}, ``Vorlesungen über analytische Geometrie des Raumes'', third edition, edited by \gesperrt{Gundelfinger}, Leipzig 1876), the works of \gesperrt{Weierstrass} and \gesperrt{Kronecker} (Monatsberichte der Berliner Akademie 1868, 1874) are of particular importance. A thorough historical account of the development of the question is given in the ``Bericht über den gegenwärtigen Stand der Invariantentheorie'' by \gesperrt{Franz Meyer} (first Jahresbericht der deutschen Mathematiker-Vereinigung, Berlin 1892).}.

\begin{center}
\S. 58.

\gesperrt{Law of inertia of quadratic forms.}
\end{center}

For the decomposition of a quadratic form into squares the following theorem now holds, which received from \gesperrt{Sylvester} the name of the law of inertia of quadratic forms (Philos. Mag. 1852).

\gesperrt{However one may decompose a real quadratic form $\varphi(x)$ into a sum of positive and negative squares of linear functions, the number of the positive and of the negative among these squares, and hence also their total number, is always the same, provided that no linear dependence exists among these linear functions.}

The proof of this important theorem is very simple.

Let
%CONT>>


% ===== scan leaf 209, printed 184 =====
% !PRINTED 184
\weberpage{184}
%<<CONT
\weq{1}{\begin{aligned}
\varphi(x) &= Y_1^2 + Y_2^2 + \cdots Y_\nu^2 - Y_1'^2 - Y_2'^2 - \cdots - Y_{\nu'}'^2\\
&= Z_1^2 + Z_2^2 + \cdots Z_\mu^2 - Z_1'^2 - Z_2'^2 - \cdots - Z_{\mu'}'^2
\end{aligned}}
be two decompositions of $\varphi(x)$ into squares. Thus let $Y_1, Y_2 \ldots Y_\nu, Y_1', Y_2' \ldots Y_{\nu'}'$ be homogeneous linear functions of $x$ between which no linear relation with constant coefficients holds, so that $\nu + \nu' \leqq m$; the same assumptions are made about the functions $Z_1, Z_2 \ldots Z_\mu, Z_1', Z_2' \ldots Z_{\mu'}'$, so that $\mu + \mu' \leqq m$ as well.

Suppose that
\[
\nu < \mu,
\]
then we can subject the variables $x$ to the linear equations
\weq{2}{\begin{aligned}
Y_1 &= 0, \quad Y_2 = 0 \ldots Y_\nu = 0\\
Z_1' &= 0, \quad Z_2' = 0 \ldots Z_{\mu'}' = 0
\end{aligned}}
whose number $\nu + \mu'$ is smaller than $m$.

Now if all the functions $Z_1, Z_2 \ldots Z_\mu$ were linearly dependent on the functions $Y_1, Y_2 \ldots Y_\nu, Z_1', Z_2' \ldots Z_{\mu'}'$, then from these $\mu$ equations one could, by eliminating the $\nu$ variables $Y_1, Y_2 \ldots Y_\nu$, derive a linear equation between the $Z_1, Z_2 \ldots Z_\mu, Z_1' \ldots Z_{\mu'}'$, which by hypothesis is \gesperrt{not} to hold. Hence the vanishing of all the $Z_1, Z_2 \ldots Z_\mu$ is not a necessary consequence of the equations (2), and we can add to the equations (2) one further equation that is not homogeneous, say
\[
Z_1 = 1,
\]
then for the $m$ unknowns $x$ we have a system of $m$ or fewer equations which are independent of one another and therefore do not contradict one another.

But then the second representation (1) gives a positive value for $\varphi(x)$, while the first gives a negative or vanishing value, and in this lies a contradiction. From this it follows that
\[
\nu \geqq \mu,
\]
and since one concludes in the same way
\[
\mu \geqq \nu,
\]
only $\nu = \mu$ remains. In the same way one can prove that $\nu' = \mu'$, whereby our theorem is completely proved.


% ===== scan leaf 210, printed 185 =====
% !PRINTED 185
\weberpage{185}
\pdfbookmark[2]{§ 59. Transformation of forms of the nth degree}{bm66}

\begin{center}
§. 59.

\gesperrt{Transformation of forms of degree $n$.}
\end{center}

When a form of degree $n$ in $m$ variables, $F(x)$, is transformed by a linear substitution [\S.\,55, (1)] into $\Phi(x')$, then at the same time the linear function of the variables $\xi_1, \xi_2 \ldots \xi_m$

\weq{1}{\overset{i}{\varSigma}\, \xi_i F'(x_i)}

is transformed, by the same substitution applied to the $\xi_i$, into

% [sic] In (2) the print omits the derivative prime on \Phi, although (1), the coefficient argument, and the following identity require \Phi'. The print is preserved.
\weq{2}{\overset{i}{\varSigma}\, \xi_i' \Phi(x_i')}

because both expressions represent the coefficient of the first power of $t$ in the expansion of
\[
F(x_i + t\xi_i) = \Phi(x_i' + t\xi_i')
\]
in powers of $t$ (\S.\,16). Likewise the quadratic function

\weq{3}{\varSigma\, \xi_i \xi_k F''(x_i x_k)}

is transformed into

\weq{4}{\varSigma\, \xi_i' \xi_k' \Phi''(x_i', x_k')}

If we apply the first result to a system of $m$ functions $F_1(x), F_2(x) \ldots F_m(x)$ of arbitrary degrees, and denote by $\varDelta$ the determinant
\[
\begin{vmatrix}
F_1'(x_1),\ F_1'(x_2) \ldots F_1'(x_m) \\
F_2'(x_1),\ F_2'(x_2) \ldots F_2'(x_m) \\
\cdot\ \cdot\ \cdot\ \cdot\ \cdot\ \cdot\ \cdot\ \cdot\ \cdot\ \cdot \\
F_m'(x_1),\ F_m'(x_2) \ldots F_m'(x_m)
\end{vmatrix}\! ,
\]
while we denote by $\varDelta'$ the corresponding determinant formed from the functions $\Phi_1, \Phi_2 \ldots \Phi_m$, then by \S.\,55, (4), if we again denote by $r$ the substitution determinant, we obtain

\weq{5}{\varDelta' = r\varDelta.}

The function $\varDelta$ is called the \gesperrt{functional determinant} of the $m$ functions $F_1, F_2 \ldots F_m$.

Likewise we obtain from the transformation (3), (4), if we denote by $H$ the determinant
%CONT>>


% ===== scan leaf 211, printed 186 =====
% !PRINTED 186
\weberpage{186}
\pdfbookmark[2]{§ 60. Invariants and Covariants}{bm67}
%<<CONT
\[
\left|
\begin{array}{l}
F''(x_1, x_1),\ F''(x_1, x_2)\ldots F''(x_1, x_m) \\
F''(x_2, x_1),\ F''(x_2, x_2)\ldots F''(x_2, x_m) \\
\cdots\cdots\cdots\cdots\cdots\cdots \\
F''(x_m, x_1),\ F'''(x_m, x_2)\ldots F''(x_m, x_m) % [sic] printed F''' in the second entry of the last row, inconsistent with F'' elsewhere
\end{array}
\right|
\]
and by $H'$ the corresponding determinant for the transformed
function $\Phi$, then by \S.\,56, (8)

\weq{6}{H' = r^2 H.}

$H$ is called the \gesperrt{Hessian determinant}\footnote{Named after \gesperrt{Hesse}, who made frequent use of this determinant in geometrical investigations, in particular in the theory of curves of the third order and in elimination problems, and who recognised its significance.} of the function
$F$. In both formulas (5) and (6), $r$ denotes the substitution
determinant.

\begin{center}
\S.\,60.\\
\gesperrt{Invariants and Covariants.}
\end{center}

In the linear transformation of homogeneous functions
there always occur certain functions of the coefficients which,
when one replaces the coefficients of the given form by the coefficients
of the transformed form, change in no other way than
that they take on a factor which is a power of the substitution
determinant. Such functions also occur in the simultaneous
transformation of a system of several forms, and they
depend on the coefficients of all the forms of the
system. Such functions are, for example, in a
system of $n$ linear forms the determinant of the system,
and in a quadratic form the determinant of this form.

These functions are called \gesperrt{invariants of the form}, or,
when they depend on several forms, \gesperrt{simultaneous invariants of the system}.

If, then, $I(a)$ is a function of the coefficients $a$ of a
form $F(x)$, and if $a'$ are the coefficients of the transformed
form $\Phi(x')$, then $I(a)$ is called an invariant if the identical relation
%CONT>>


% ===== scan leaf 212, printed 187 =====
% !PRINTED 187
\weberpage{187}
%<<CONT
\weq{1}{I(a') \doteq r^{\lambda} I(a)}
holds. The exponent $\lambda$ we shall call the \gesperrt{weight} of the
invariant. If this exponent is zero, then $I$ is called an
\gesperrt{absolute invariant}.

One mostly considers only \gesperrt{integral rational invariants},
that is, functions $I$ which depend rationally on the coefficients $a$
and are moreover integral homogeneous functions of
them. The simultaneous invariants are required to be homogeneous
with respect to the coefficients of each of the forms on which
they depend.

Besides the invariants there are also functions which, in addition to
the coefficients, still contain the variables themselves, but in other
respects have the same properties as the invariants (1), ``the invariant
property''; thus, if $C(x, a)$ is such a
function,

\weq{2}{C(x', a') \doteq r^{\lambda} C(x, a).}

Such functions are called \gesperrt{covariants of the form $F(x)$}
% [sic] The print omits the sentence-ending period after F(x).
The functional determinant is a simultaneous covariant of a
system of $m$ functions, and the Hesse determinant
of a function of higher than the second degree is a
covariant of a single form. Among the covariants too
one considers preferably integral rational and homogeneous
functions of $x$.

If $m$ is the number of the variables $x$, $n$ the degree of the function
$F(x)$, and $\nu$ and $\mu$ the degrees of $C(x, a)$ with respect to the $x$ and
$a$, then among the numbers $\lambda, \mu, \nu, m, n$ there holds the relation

\weq{3}{n\mu = m\lambda + \nu,}

which one proves by comparing the degree of the right and
left sides of (2) with respect to the substitution coefficients.
The $a'$ are namely of degree $n$ in the substitution coefficients,
$r$ of degree $m$, $x$ of the first degree, from which
(3) follows. The corresponding formula for the invariants
is obtained by setting $\nu = 0$; indeed, the
invariants may altogether be regarded as a special case of the
covariants.

A general law for the formation of invariants and
covariants, which we have already used in the particular cases of \S.\,59,
we shall now briefly discuss.


% ===== scan leaf 213, printed 188 =====
% !PRINTED 188
\weberpage{188}

If, following \S.\,16, (4), we arrange the function
\[
F(x_1 + t\xi_1, x_2 + t\xi_2 \ldots x_m + t\xi_m)
\]
in powers of $t$, we obtain for the coefficient of $t^\nu$
\weq{4}{F_\nu(x,\xi) = \frac{1}{\Pi(\nu)}\sum^{\alpha}\frac{\Pi(\nu)}{\Pi(\alpha_1) \ldots \Pi(\alpha_m)}\xi_1^{\alpha_1}\xi_2^{\alpha_2}\ldots\xi_m^{\alpha_m}D_{\alpha_1,\alpha_2\ldots\alpha_m}(F),}
\[
\alpha_1 + \alpha_2 + \cdots + \alpha_m = \nu.
\]
If we apply the linear transformation (1), \S.\,55 to the variables
$\xi$ and $x$ simultaneously, then $F_\nu(x,\xi)$ passes over into $\Phi_\nu(x',\xi')$,
which is obtained from $F_\nu$ by simultaneously replacing all the letters
$x, \xi, a$ by $x', \xi', a'$ and $F$ by $\Phi$.

From this it follows:

I. \gesperrt{If one regards $F_\nu(x,\xi)$ as a form of degree $\nu$
in $\xi$ and forms an invariant of this form, whose coefficients therefore
still depend on the $x$, one obtains a covariant of the form $F$.}

And since the variables $x$ and $\xi$ are subject to the same
transformation, it likewise follows:

II. \gesperrt{If one forms a covariant of the function $F_\nu(x,\xi)$
regarded as a function of $\xi$, and in it replaces the $\xi$ by the $x$,
one obtains a covariant of the original function $F(x)$.}

The same theorems hold also for simultaneous invariants
and covariants.

For the functions $F_\nu(x,\xi)$ we have the theorem
\weq{5}{F_\nu(x,\xi) = F_{n-\nu}(\xi,x),}
which follows by comparing corresponding powers of $t$
on the two sides of the following identity:
\[
F(x_1+t\xi_1\ldots x_m+t\xi_m) = t^n F(t^{-1}x_1+\xi_1\ldots t^{-1}x_m+\xi_m)
\]

The functions
\[
P_\nu(x,\xi) = \frac{\Pi(\nu)\;\Pi(n-\nu)}{\Pi(n)}F_\nu(x,\xi)
\]
are called the \gesperrt{polars of the function $F(x)$}, and specifically
$P_\nu(x,\xi)$ is called the $\nu$th polar. For them too the following theorem holds:
\[
P_\nu(x,\xi) = P_{n-\nu}(\xi,x).
\]


% ===== scan leaf 214, printed 189 =====
% !PRINTED 189
\weberpage{189}
\pdfbookmark[2]{§ 61. Linear transformation of binary forms}{bm68}

The constant factor $\Pi(\nu)\Pi(n-\nu):\Pi(n)$ is added so that,
when the form $F(x)$ is written with the polynomial coefficients,
the functions $P_\nu$ receive the simplest possible coefficients.

For the binary form $f(x,y)$, to which we shall chiefly apply
these definitions in what follows, the following expressions
result for the polars:
\[
P_1(x,\xi) = \frac{1}{n}\left[\xi f'(x) + \eta f'(y)\right]
\]
\[
P_2(x,\xi) = \frac{1}{n(n-1)}\left\{\xi^2 f''(x,x) + 2\,\xi\,\eta\, f''(x,y) + \eta^2 f''(y,y)\right\},
\]
and in general
\[
n(n-1)\ldots(n-\nu+1)\,P_\nu(x,\xi) =
\]
\[
\xi^\nu\frac{\partial^\nu f}{\partial x^\nu} + \nu\,\xi^{\nu-1}\eta\,\frac{\partial^\nu f}{\partial x^{\nu-1}\partial y} + \cdots + \eta^\nu\frac{\partial^\nu f}{\partial y^\nu}.
\]

\begin{center}
\S.\,61.

\gesperrt{Linear Transformation of Binary Forms.}
\end{center}

An integral rational function of degree $n$ in $x$
\weq{1}{f(x) = a_0x^n + a_1x^{n-1} + a_2x^{n-2} + \cdots + a_{n-1}x + a_n}
is turned into a binary form by replacing $x$ with $x:y$
and multiplying by $y^n$:
\weq{2}{f(x,y) = a_0x^n + a_1x^{n-1}y + a_2x^{n-2}y^2 + \cdots + a_ny^n.}

If one knows the roots of $f(x)$, one can decompose $f(x)$
into $n$ linear factors:
\weq{3}{f(x) = a_0(x-\alpha_1)(x-\alpha_2)\ldots(x-\alpha_n),}
and from this it follows that
\weq{4}{f(x,y) = a_0(x-\alpha_1y)(x-\alpha_2y)\ldots(x-\alpha_ny).}

But if one also replaces $\alpha_1, \alpha_2\ldots\alpha_n$ by
ratios $\alpha_1:\beta_1, \alpha_2:\beta_2, \ldots \alpha_n:\beta_n$, then
one obtains
\weq{5}{f(x,y) = A(x\beta_1 - y\alpha_1)(x\beta_2 - y\alpha_2)\ldots(x\beta_n - y\alpha_n),}
if
\[
a_0 = A\beta_1\beta_2\ldots\beta_n
\]
is set.


% ===== scan leaf 215, printed 190 =====
% !PRINTED 190
\weberpage{190}

For the factors of $f(x, y)$, which we shall now encounter more often, we introduce as an abbreviation the sign

\weq{6}{x\beta_1 - y\alpha_1 = (x\beta_1).}

Likewise we shall set

\weq{7}{x_1 y_2 - x_2 y_1 = (x_1 y_2),}

in which $x_1, y_1$ and $x_2, y_2$ are two arbitrary pairs of variables. When we employ the inhomogeneous representation, we set still more briefly

\weq{8}{\begin{aligned}
x - \alpha_1 &= (0, 1),\ x - \alpha_2 = (0, 2) \ldots x - \alpha_n = (0, n) \\
\alpha_1 - \alpha_2 &= (1, 2) \ldots \alpha_i - \alpha_k = (i, k).
\end{aligned}}

If we make in the form (2) a linear substitution

\weq{9}{\begin{aligned}
x &= \alpha x' + \beta y' \\
y &= \gamma x' + \delta y'
\end{aligned}}

with the determinant, different from zero,

\weq{10}{r = \alpha\delta - \beta\gamma,}

then the form $f(x, y)$ passes over into a form $F(x', y')$ of degree $n$, and we have

\weq{11}{F(x', y') = a'_0 x'^n + a'_1 x'^{n-1} y' + a'_2 x'^{n-2} y'^2 + \cdots a'_n y'^n.}

The coefficients $a'_0, a'_1 \ldots a'_n$ of the transformed form are obtained as linear and homogeneous expressions in the coefficients $a_0, a_1 \ldots a_n$, and as homogeneous expressions of degree $n$ in the substitution coefficients $\alpha, \beta, \gamma, \delta$. One thus obtains, on setting $x' = 1, y' = 0$ or $x' = 0, y' = 1$,

\weq{12}{a'_0 = f(\alpha, \gamma), \qquad a'_n = f(\beta, \delta).}

The remaining coefficients $a'$ can be expressed by means of the derivatives of the function $f(x, y)$, e.g.

% [sic] The print has the factor (n - 1); the coefficient form in (2) does not support it algebraically.
\[
(n - 1) a'_1 = \beta f'(\alpha) + \delta f'(\gamma),
\]

which we shall not carry out further.

If we apply the transformation (9) to the two pairs of variables $x, y$ and $\alpha_k, \beta_k$ simultaneously, we obtain the linear factors of $f(x, y)$ transformed, namely

\weq{13}{(x'\beta'_k) = r(x\beta_k),}

as one finds from the multiplication rule for determinants, or also by direct computation. These factors $(x\beta_k)$ are therefore also \gesperrt{covariants} of $f(x, y)$, though to be sure irrational ones, since they do not depend rationally on the coefficients of $f(x)$. In the same way one obtains
%CONT>>


% ===== scan leaf 216, printed 191 =====
% !PRINTED 191
\weberpage{191}
%<<CONT

\weq{14}{(\alpha'_h\beta'_k) = r(\alpha_h\beta_k),}

according to which the determinants $(\alpha_h\beta_k)$ are to be regarded as irrational \gesperrt{invariants}. We now wish to set out how rational invariants and covariants can be formed from them.

For this purpose it is best to consider the linear transformation in the inhomogeneous form

\weq{15}{x = \frac{\alpha x' + \beta}{\gamma x' + \delta}}

and to apply it to the function $f(x)$ in (1). It then follows, if

\[
F(x') = a'_0 x'^n + a'_1 x'^{n-1} + a'_2 x'^{n-2} + \cdots + a'_n
\]

is set, that

\weq{16}{F(x') = (\gamma x' + \delta)^n f(x).}

We also write this equation thus

% [sic] The print omits the otherwise required plus sign between the a_2/x^2 term and the ellipsis.
\[
F(x') = (\alpha x' + \beta)^n \left(a_0 + a_1 \frac{1}{x} + a_2 \frac{1}{x^2} \cdots\right),
\]

from which, if one sets $x' = -\delta : \gamma$, and hence $x = \infty$,

\[
a_0 r^n = (-\gamma)^n F\left(-\frac{\delta}{\gamma}\right),
\]

or, if $\alpha'_1, \alpha'_2, \ldots \alpha'_n$ depend on $\alpha_1, \alpha_2, \ldots \alpha_n$ in the same way as $x'$ does on $x$, so that

\[
F(x') = a'_0 (x' - \alpha'_1)(x' - \alpha'_2) \ldots (x' - \alpha'_n)
\]

is set,

\weq{17}{r^n a_0 = a'_0(\gamma\alpha'_1 + \delta)(\gamma\alpha'_2 + \delta) \ldots (\gamma\alpha'_n + \delta)}

follows.

Now by the transformation (15) we have

\weq{18}{\begin{aligned}
x - \alpha_i &= \frac{r(x' - \alpha'_i)}{(\gamma x' + \delta)(\gamma\alpha'_i + \delta)} \\
\alpha_i - \alpha_k &= \frac{r(\alpha'_i - \alpha'_k)}{(\gamma\alpha'_i + \delta)(\gamma\alpha'_k + \delta)}.
\end{aligned}}

We now form a product $P$ out of any of the factors

\[
\begin{aligned}
&(0, 1), (0, 2) \ldots (0, n) \\
&(1, 2), (1, 3) \ldots (n - 1, n),
\end{aligned}
\]

but in such a way that each of the digits $1, 2 \ldots n$ occurs in it equally often in all, say $\mu$ times. The index $0$, that is, the variable $x$, may occur $\nu$ times, and the total number of the
%CONT>>


% ===== scan leaf 217, printed 192 =====
% !PRINTED 192
\weberpage{192}
\pdfbookmark[2]{§ 62. Binary cubic forms}{bm69}
%<<CONT
factors shall be $q$. Then, since each factor contains two
indices,
\weq{19}{\nu + n\mu = 2q.}
Then from (17) and (18), if $P'$ arises from $P$ by the
permutation of $x, \alpha_i$ with $x', \alpha_i'$, it follows that
\weq{20}{{a_0'}^{\mu} P' = a_0^{\mu} r^{n\mu-q} (\gamma x' + \delta)^{\nu} P.}
We now form
\weq{21}{C(x, y, a) = a_0^{\mu} y^{\nu} \Sigma P\left(\frac{x}{y}\right),}
the sum being taken over all the values that one obtains from $P$
when the indices $1, 2, 3 \ldots n$ are permuted with one another
in all possible ways. This sum is then a symmetric function of
the roots of (1), and can therefore be expressed as an integral
rational function of the ratios $a_1 : a_0, a_2 : a_0 \ldots a_n : a_0$,
in such a way that no term exceeds the order $\mu$. The
function $C(x, y, a)$ is therefore an integral homogeneous function of
$a_0, a_1 \ldots a_n$ of order $\mu$, and an integral homogeneous
function of $x, y$ of order $\nu$. Since by (19), (20) it satisfies
the condition
% [sic] The print omits the comma between y' and a' in C(x', y' a').
\weq{22}{C(x', y' a') = r^{\lambda} C(x, y, a),}
\weq{23}{\lambda = n\mu - q, \quad 2\lambda = n\mu - \nu}
it is a \gesperrt{covariant}, and in the particular case
where $\nu = 0$, an \gesperrt{invariant}.

The weight $\lambda$ is, as the second expression shows, always
positive or zero, since $n\mu$ is at least equal to $\nu$. The extreme
value $n\mu = \nu$, or $\lambda = 0$, occurs only when in $P$
each of the indices $1, 2, 3 \ldots n$ occurs joined only with 0,
so that $P$ is a power of $f(x)$ itself.

\begin{center}
\S. 62.

\gesperrt{Binary cubic forms.}
\end{center}

A binary quadratic form gives no invariant
formations other than the discriminant. With these we therefore
do not concern ourselves, and pass at once to the consideration of the
cubic forms
\weq{1}{f(x, y) = a_0 x^3 + a_1 x^2 y + a_2 x y^2 + a_3 y^3}
%CONT>>


% ===== scan leaf 218, printed 193 =====
% !PRINTED 193
\weberpage{193}
%<<CONT
over. Here we form, as the first \gesperrt{quadratic} covariant, the \gesperrt{Hessian} determinant. We shall here always follow the rule of writing the forms with integer numerical coefficients having no common factor. Then, in the determinant formed from the second derivatives of (1), we must discard the factor 4, and we obtain as the first covariant
\[
H = \left|\begin{array}{cc} 3\,a_0\,x + a_1\,y, & a_1\,x + a_2\,y \\ a_1\,x + a_2\,y, & a_2\,x + 3\,a_3\,y \end{array}\right|.
\]
or
\weq{2}{H = A_0\,x^2 + A_1\,x\,y + A_2\,y^2,}
if for abbreviation
\weq{3}{
\begin{aligned}
A_0 &= 3\,a_0\,a_2 - a_1^2 \\
A_1 &= 9\,a_0\,a_3 - a_1\,a_2 \\
A_2 &= 3\,a_1\,a_3 - a_2^2
\end{aligned}
}
is set.

The discriminant of this quadratic form is an invariant of the given cubic form. It contains, however, the factor 3 in all its numerical coefficients, and accordingly we set
\weq{4}{3\,D = 4\,A_0\,A_2 - A_1^2,}
and then $D$ agrees with the discriminant of the cubic form [\S.\,46, (10)]:
\weq{5}{D = a_1^2\,a_2^2 + 18\,a_0\,a_1\,a_2\,a_3 - 4\,a_0\,a_2^3 - 4\,a_1^3\,a_3 - 27\,a_0^2\,a_3^2.}

We obtain a further cubic covariant if we form the functional determinant of $f(x,y)$ and $H(x,y)$:
\weq{6}{Q(x,y) = f'(x)\,H'(y) - f'(y)\,H'(x)}
or
\weq{7}{
Q = \left|\begin{array}{cc} 3\,a_0\,x^2 + 2\,a_1\,x\,y + a_2\,y^2, & 2\,A_0\,x + A_1\,y \\ a_1\,x^2 + 2\,a_2\,x\,y + 3\,a_3\,y^2, & A_1\,x + 2\,A_2\,y \end{array}\right|,
}
in which no numerical factor can be cancelled. We set down here, in order, the computed coefficients of $x^3$, $x^2y$, $xy^2$, $y^3$, whose formation presents no difficulty:
\[
\begin{aligned}
&27\,a_0^2\,a_3 - 9\,a_0\,a_1\,a_2 + 2\,a_1^3,\\
&27\,a_0\,a_1\,a_3 - 18\,a_0\,a_2^2 + 3\,a_1^2\,a_2,\\
-\,&27\,a_0\,a_2\,a_3 + 18\,a_1^2\,a_3 - 3\,a_1\,a_2^2,\\
-\,&27\,a_0\,a_3^2 + 9\,a_1\,a_2\,a_3 - 2\,a_2^3.
\end{aligned}
\]


% ===== scan leaf 219, printed 194 =====
% !PRINTED 194
\weberpage{194}
The behaviour of $H$, $D$, $Q$ under linear transformation follows from \S.\,60, (2), (3); if $H'$, $D'$, $Q'$ are the corresponding constructions for a transformed form, one has
\weq{8}{H' = r^2\,H,\quad D' = r^6\,D,\quad Q' = r^3\,Q.}

We can use the covariants to transform the cubic form to a normal form which at the same time gives the solution of the cubic equation.

We choose the normal form
\weq{9}{f(x,y) = F(\xi,\eta) = \xi^3 + \eta^3,}
in which $\xi$, $\eta$ are linear functions of $x$, $y$; a form which, as will appear presently, can always be produced when $D$ does not vanish, that is, when the equation $f = 0$ does not have two equal roots. The solutions of $f = 0$ are then obtained from the linear equations
\[
\xi + \eta = 0,\quad \xi + \varepsilon\,\eta = 0,\quad \xi + \varepsilon^2\,\eta = 0,
\]
in which $\varepsilon$ is a third root of unity.

To form the expressions $H'$, $D'$, $Q'$ for the normal form (9), we have to set $a_0' = a_3' = 1$, $a_1' = a_2' = 0$, and obtain
\[
H' = 9\,\xi\,\eta,\quad D' = -\,27,\quad Q' = 27\,(\xi^3 - \eta^3).
\]

If the normal form (9) is derived from the general form $f(x,y)$ by linear transformation, then by (8) we obtain
\weq{10}{
\begin{aligned}
9\,\xi\,\eta &= r^2\,H \\
-\,27 &= r^6\,D \\
27\,(\xi^3 - \eta^3) &= r^3\,Q \\
\xi^3 + \eta^3 &= f.
\end{aligned}
}
From this it follows that
\[
r^6 = \frac{-\,27}{D},\quad 2\,\xi^3 = \frac{r^3\,Q}{27} + f,\quad 2\,\eta^3 = -\,\frac{r^3\,Q}{27} + f,
\]
hence, if for $r^3$ the value $-\,9 : \sqrt{-3\,D}$ from the first equation is substituted,
\weq{11}{
\begin{aligned}
6\,\xi^3\,\sqrt{-3\,D} &= +\,Q + 3\,\sqrt{-3\,D}\,f \\
6\,\eta^3\,\sqrt{-3\,D} &= -\,Q + 3\,\sqrt{-3\,D}\,f,
\end{aligned}
}
from which $\xi$ and $\eta$ are determined by two cube roots, of which, by the first equation (10), the one is determined by the other.


% ===== scan leaf 220, printed 195 =====
% !PRINTED 195
\weberpage{195}
In this also lies the proof that, if $D$ is different from zero, the normal form can be produced by a linear transformation. For under this assumption $H$ decomposes into two linear factors distinct from one another, and if one chooses these, apart from constant factors, as the new variables $\xi$, $\eta$ of a linear transformation, then $A_0' = 0$, $A_2' = 0$; but it follows from the identical relations
\[
3\,a_3\,A_0 + a_1\,A_2 = a_2\,A_1,\quad a_2\,A_0 + 3\,a_0\,A_2 = a_1\,A_1,
\]
unless at the same time $A_1' = 0$, which is excluded by the non-vanishing $D$, that $a_1' = 0$, $a_2' = 0$, i.\,e. the transformed form of $f$ contains only the cubes of $\xi$ and $\eta$.

If one raises the first equation (10) to the cube, and eliminates $\xi^3$, $\eta^3$, $r^6$ by means of (11) and the second equation (10), one obtains among the covariants the following identical relation
\weq{12}{4\,H^3 + Q^2 + 27\,D\,f^2 = 0.}

To carry out the transformation into the normal form, i.\,e. actually to find the functions $\xi$, $\eta$, one decomposes the function $H$ into its linear factors
% [sic] The print has A_2 y in the second factor. Algebraically it must be A_1 y, as in the first factor, because A_1^2-4A_0A_2=-3D. The print is preserved.
\[
4\,A_0\,H = \bigl(2\,A_0\,x + A_1\,y + \sqrt{-3\,D}\,y\bigr)\,\bigl(2\,A_0\,x + A_2\,y - \sqrt{-3\,D}\,y\bigr).
\]
Then $\xi$, $\eta$ differ from these two factors of $H$ only by a constant factor each. We therefore set, if we denote two constant factors by $h$, $k$,
\weq{13}{
\begin{aligned}
2\,\xi &= h\,\bigl(2\,A_0\,x + A_1\,y - \sqrt{-3\,D}\,y\bigr) \\
2\,\eta &= k\,\bigl(2\,A_0\,x + A_1\,y + \sqrt{-3\,D}\,y\bigr),
\end{aligned}
}
from which, by multiplication, with regard to the first two equations (10), there results
\weq{14}{3\,h\,k\,A_0\,\sqrt[3]{-D} = 1,}
and the equations (11) give, by comparison of the coefficients of $x^3$,
\weq{15}{
\begin{aligned}
6\,h^3\,A_0^3\,\sqrt{-3\,D} &= +\,q_0 + 3\,\sqrt{-3\,D}\,a_0 \\
6\,k^3\,A_0^3\,\sqrt{-3\,D} &= -\,q_0 + 3\,\sqrt{-3\,D}\,a_0,
\end{aligned}
}
where $q_0$ is the coefficient of $x^3$ in $Q$, that is
\weq{16}{q_0 = 27\,a_0^2\,a_3 - 9\,a_0\,a_1\,a_2 + 2\,a_1^3}
.


% ===== scan leaf 221, printed 196 =====
% !PRINTED 196
\weberpage{196}
\pdfbookmark[2]{§ 63. The complete system of forms of the binary cubic form}{bm70}
That the signs of $\sqrt{-3\,D}$ are to be taken, with regard to the signs in (11), as has been done in the formulas (13), follows most simply if one carries out the computation for the special case $a_0 = 0$, $a_3 = 0$.

\begin{center}
\S.\,63.
\end{center}

\begin{center}
\gesperrt{The full system of forms of the binary cubic form.}
\end{center}

We wish to show further that the invariants and covariants of the cubic form are exhausted by the functions $D, f, H, Q$, i.\,e. that all invariants and covariants of a cubic form can be expressed rationally in terms of these four forms.

So let
\[
C = C(x,y,a)
\]
be a covariant of degree $\nu$ in the variables and of degree $\mu$ in the coefficients, of which we shall assume that it has only integer numerical coefficients.

Then, for any transformed form,
\weq{1}{C' = C(\xi,\eta,a') = r^{\lambda}\,C(x,y,a),}
\weq{2}{\lambda = \frac{3\,\mu - \nu}{2}.}

We now understand by $\xi$, $\eta$ the variables of the normal form considered above. If $\varepsilon$ denotes a third root of unity, then
\[
\xi = \varepsilon\,\xi',\quad \eta = \varepsilon^2\,\eta'
\]
is a linear substitution with determinant 1, by which the normal form \S.\,62, (9) passes into itself. It follows, therefore, that $C'$ cannot change when $\xi$, $\eta$ are replaced by $\varepsilon\,\xi$, $\varepsilon^2\,\eta$; hence $C'$ can depend rationally only on $\xi\,\eta$, $\xi^3$, $\eta^3$.

The permutation of $\xi$ with $\eta$ corresponds to a substitution with determinant $-1$; under it, therefore, $C'$ remains unchanged or changes its sign, according as $\lambda$ is even or odd. Hence $C'$ consists of a sum of terms of the form
\[
M\,(\xi\,\eta)^h\,(\xi^{3k} \pm \eta^{3k}),
\]
in which $M$ is an integer factor (since we assume the numerical coefficients in $C$, and hence also in $C'$, to be
%CONT>>


% ===== scan leaf 222, printed 197 =====
% !PRINTED 197
\weberpage{197}
%<<CONT
have set) and the upper or lower sign holds according as
$\lambda$ is even or odd. In the latter case this expression is
divisible by $\xi^3 - \eta^3$, and the quotient can be expressed rationally
and in integers through $\xi^3 + \eta^3$ and $\xi\eta$ (because a
symmetric function of two quantities can be represented in this way
through the sum and the product). If we therefore set $\tau = 0$
or $= 1$ according as $\lambda$ is even or odd, there results
for $C'$ an expression of the form
\[
C'(\xi, \eta) = (\xi^3 - \eta^3)^\tau \; \Sigma \, M (\xi\eta)^\alpha (\xi^3 + \eta^3)^\beta,
\]
in which the $M$ are integer coefficients and $\alpha$, $\beta$ run through a series
% [sic] The print says positive integers; as exponents alpha and beta may also be 0, so nonnegative integers are meant.
of positive integers which satisfy the condition
\weq{3}{3\tau + 3\beta + 2\alpha = \nu}
From (1) there results, then, if we multiply by a sufficiently high power of 3 and express $\xi\eta$, $\xi^3+\eta^3$, $\xi^3-\eta^3$
by $H$, $f$, $Q$ according to \S.\,62, (10), and $r$ by $D^{-\frac{1}{6}}$,
\weq{4}{3^\sigma C = Q^\tau \, \Sigma \, M D^\gamma H^\alpha f^\beta,}
in which
\weq{5}{6\gamma = \lambda - 3\tau - 2\alpha,}
and in which the $M$ are likewise integer factors.

To each value of $\gamma$ there belongs, by (5), a definite
value of $\alpha$, and by (3) a definite value of $\beta$; likewise
the two remaining numbers are in each case determined by a value
of $\alpha$ or of $\beta$, so that in (4) every exponent of
$D$, $H$ or $f$ occurs only once.

That $\gamma$ is an integer follows easily from (2), (3), (5);
for since by the definition $\lambda - 3\tau$ is an even number,
$3\gamma$ is first of all an integer, and since further, by (3),
\[
\lambda - 2\alpha = \frac{3}{2}(\mu - \nu + 2\beta + 2\tau),
\]
$6\gamma$, and hence $3\gamma$, is divisible by 3. From this it follows that
% [sic] The print names sigma here. The immediately preceding divisibility argument actually proves that gamma is integral; sigma had already been introduced as a sufficiently high exponent of 3.
$\sigma$ too is an integer.

That $\gamma$ must, however, also be positive, we see as follows:

Suppose that on the right-hand side of (4) negative powers
of $D$ also occur; then we multiply this equation
on both sides by so high a power of $D$ that on the right no
negative powers of $D$ appear any longer, but that the
right-hand side also does not acquire the factor $D$. If in the equation
so obtained we set
%CONT>>


% ===== scan leaf 223, printed 198 =====
% !PRINTED 198
\weberpage{198}
%<<CONT
\weq{6}{a_0 = 0, \quad a_1 = 1, \quad a_2 = 0, \quad a_3 = 0,}
then
\weq{7}{D = 0, \quad H = -x^2, \quad f = x^2 y, \quad Q = 2 x^3.}
Thus the left side would vanish while the right side does not
vanish, and in this lies a contradiction.

But finally we can also prove that on the left side of (4)
no power of 3 can occur which is not contained in all the
factors $M$, and which therefore cancels out.

In \S.\,2 we proved the theorem that the product of two
primitive integral rational functions of arbitrary variables is
again a primitive function. Here a primitive function is understood
to be one whose coefficients are integers without a common
divisor. The point now is to show that a sum of the form
\[
Q^\tau \, \Sigma \, M D^\gamma H^\alpha f^\beta,
\]
in which the $M$ are integers without a common divisor, cannot
become imprimitive, and in particular cannot acquire the factor 3,
when $Q$, $D$, $H$, $f$ are replaced by their expressions in the
$a$, $x$, $y$. By the theorem just mentioned it suffices, since $Q$
is a primitive function, to show this for the form
\weq{8}{\Sigma \, M D^\gamma H^\alpha f^\beta.}
In this we may moreover omit all terms whose $M$ already has
the factor 3, and finally, again by the theorem mentioned, we
may assume that the expression is not divisible by $D$, and
hence that it contains a term with $\gamma = 0$.
Suppose then that under these assumptions the expanded
expression (8) has the divisor 3, and now substitute the
particular values (6), (7), i.\,e.\ set $D = 0$, $H = -x^2$,
$f = x^2 y$; then the expression reduces to the single term
in which $\gamma = 0$, $2\alpha = \lambda$,
\[
(-1)^\alpha M x^{2(\alpha + \beta)} y^\beta
\]
and it would therefore follow that this $M$ must have the
divisor 3, which is contrary to the assumption. Thus we have
proved:

\gesperrt{Every integer covariant of the cubic form
can be expressed integrally and rationally and with integer coeffi-}
%CONT>>


% ===== scan leaf 224, printed 199 =====
% !PRINTED 199
\weberpage{199}
\pdfbookmark[2]{§ 64. Biquadratic forms}{bm71}
%<<CONT
\gesperrt{cients $M$ be represented by one of the two
expressions}
\[
\Sigma \, M D^\gamma H^\alpha f^\beta, \quad Q \, \Sigma \, M D^\gamma H^\alpha f^\beta.
\]
The invariants are contained among the covariants as a
special case, and are all powers of $D$.

\begin{center}
\S.\,64.

\gesperrt{Biquadratic Forms.}
\end{center}

We now pass on to the consideration of the following biquadratic form:
\weq{1}{f(x, y) = a_0 x^4 + a_1 x^3 y + a_2 x^2 y^2 + a_3 x y^3 + a_4 y^4}

We have first the \gesperrt{Hessian} determinant as a covariant of the fourth order, which, in order to free it of a numerical
factor, we divide by 3:
\weq{2}{H = \frac{1}{3}
\begin{vmatrix}
12 a_0 x^2 + 6 a_1 x y + 2 a_2 y^2, & 3 a_1 x^2 + 4 a_2 x y + 3 a_3 y^2 \\
3 a_1 x^2 + 4 a_2 x y + 3 a_3 y^2, & 2 a_2 x^2 + 6 a_3 x y + 12 a_4 y^2
\end{vmatrix},}
or
\[
H = A_0 x^4 + A_1 x^3 y + A_2 x^2 y^2 + A_3 x y^3 + A_4 y^4,
\]
in which
\[
\begin{aligned}
A_0 &= 8 a_0 a_2 - 3 a_1^2, & A_4 &= 8 a_2 a_4 - 3 a_3^2, \\
A_1 &= 24 a_0 a_3 - 4 a_1 a_2, & A_3 &= 24 a_1 a_4 - 4 a_2 a_3, \\
A_2 &= 48 a_0 a_4 + 6 a_1 a_3 - 4 a_2^2, & &
\end{aligned}
\]
and a covariant of the sixth degree
\weq{3}{T = \frac{1}{12}
\begin{vmatrix}
f'(x), & f'(y) \\
H'(x), & H'(y)
\end{vmatrix}.}

The invariants of the biquadratic form are most simply
formed, following \S.\,61, from the differences of the roots.

We thus obtain an invariant of the second order and one
of the third order in the coefficients:
\weq{4}{A = \tfrac{1}{2} a_0^2 [(12)^2(34)^2 + (13)^2(24)^2 + (14)^2(23)^2],}
\weq{5}{B = a_0^3 \, \Sigma \, (12)^2(34)^2(13)(42),}
expressions which can be written more perspicuously if one
sets
\weq{6}{U = (12)(34), \quad V = (13)(42), \quad W = (14)(23)}
namely
%CONT>>


% ===== scan leaf 225, printed 200 =====
% !PRINTED 200
\weberpage{200}
%<<CONT
\[
\begin{aligned}
& A = \tfrac{1}{2} a_0^2 (U^2 + V^2 + W^2) \\
(7)\quad & B = a_0^3 \bigl[U^2(V - W) + V^2(W - U) + W^2(U - V)\bigr] \\
& \phantom{B} = a_0^3 (W - V)(U - W)(V - U).
\end{aligned}
\]
For the calculation of these quantities the necessary formulas have
already been developed in \S.\,47. By the formulas (10) given there we obtain
\[
U = v^2 - w^2, \quad V = w^2 - u^2, \quad W = u^2 - v^2,
\]
if $u^2$, $v^2$, $w^2$ are the roots of the cubic resolvent
\[
z^3 + 2 a z^2 + (a^2 - 4 c) z - b^2 = 0
\]
and from this it follows that
\[
\begin{aligned}
A &= a_0^2 \bigl[(u^2 + v^2 + w^2)^2 - 3 (v^2 w^2 + w^2 u^2 + u^2 v^2)\bigr] \\
&= a_0^2 (a^2 + 12 c) \\
B &= a_0^3 (2 u^2 - v^2 - w^2)(2 v^2 - w^2 - u^2)(2 w^2 - u^2 - v^2) \\
&= a_0^3 (3 u^2 + 2 a)(3 v^2 + 2 a)(3 w^2 + 2 a) \\
&= a_0^3 (2 a^3 - 72 a c + 27 b^2),
\end{aligned}
\]
so that $A$, $B$ have the same meaning as in \S.\,47, and, expressed
in terms of the coefficients $a_0, a_1, a_2, a_3, a_4$, are represented thus:
\weq{8}{A = a_2^2 - 3 a_1 a_3 + 12 a_0 a_4}
\weq{9}{B = 27 a_1^2 a_4 + 27 a_0 a_3^2 + 2 a_2^3 - 72 a_0 a_2 a_4 - 9 a_1 a_2 a_3.}
The discriminant of the biquadratic equation as well, as
a third invariant, though one dependent on $A$ and $B$,
\weq{10}{D = a_0^6 \, U^2 V^2 W^2}
we have already formed at the place mentioned, and found
\weq{11}{27 D = 4 A^3 - B^2.}
% [sic] The print refers to section 67; the formula quoted here verbatim is the binary specialization of section 60, (3), while section 67 begins only later.
By the formula of \S.\,67, $2\lambda = n\mu - \nu$, we obtain for the
invariant constructions found so far, which we denote for a
transformed form with accents, the following relations
\weq{12}{
\begin{aligned}
& H' = r^2 H, \quad T' = r^3 T, \\
& A' = r^4 A, \quad B' = r^6 B, \quad D' = r^{12} D.
\end{aligned}
}


% ===== scan leaf 226, printed 201 =====
% !PRINTED 201
\weberpage{201}
\pdfbookmark[2]{§ 65. Solution of the biquadratic equation}{bm72}

\begin{center}
§. 65.

\gesperrt{Solution of the biquadratic equation.}
\end{center}

We wish to produce a normal form by a linear substitution
with determinant 1; for its derivation we assume the
given biquadratic form decomposed into linear factors
:
\weq{1}{f(x,\,y) = a_0\,(x - \alpha y)\,(x - \beta y)\,(x - \gamma y)\,(x - \delta y).}
We set
\weq{2}{\begin{gathered}
x\,\sqrt{\alpha - \beta} = \beta\,\xi - \alpha\,\eta, \quad r = 1,\\
y\,\sqrt{\alpha - \beta} = \xi - \eta,
\end{gathered}}
hence
\weq{3}{\begin{gathered}
x - \alpha y = -\sqrt{\alpha - \beta}\,\xi\\
x - \beta y = -\sqrt{\alpha - \beta}\,\eta\\
(x - \gamma y)\,\sqrt{\alpha - \beta} = (\beta - \gamma)\,\xi - (\alpha - \gamma)\,\eta\\
(x - \delta y)\,\sqrt{\alpha - \beta} = (\beta - \delta)\,\xi - (\alpha - \delta)\,\eta,
\end{gathered}}
and thereby $f(x,\,y)$ passes over into
\weq{4}{F(\xi,\,\eta) = a_1'\,\xi^3\eta + a_2'\,\xi^2\eta^2 + a_3'\,\xi\eta^3.}

The coefficients $a_1', a_2', a_3'$ are easily expressed by (3) in terms of the
$\alpha, \beta, \gamma, \delta$:
\weq{5}{\begin{gathered}
a_1' = \phantom{-{}} a_0\,(\beta - \gamma)\,(\beta - \delta)\\
a_3' = \phantom{-{}} a_0\,(\alpha - \gamma)\,(\alpha - \delta)\\
a_2' = -a_0\,[(\alpha - \gamma)\,(\beta - \delta) + (\alpha - \delta)\,(\beta - \gamma)]
\end{gathered}}

By permutation of the roots we can produce the normal form
 (4) in six different ways, which however yield only
three different values of $a_2'$. These three values
are, if as above
\weq{6}{\begin{gathered}
U = (\alpha - \beta)\,(\gamma - \delta),\\
V = (\alpha - \gamma)\,(\delta - \beta),\\
W = (\alpha - \delta)\,(\beta - \gamma),\\
U + V + W = 0
\end{gathered}}
is set,
%CONT>>


% ===== scan leaf 227, printed 202 =====
% !PRINTED 202
\weberpage{202}
%<<CONT
\weq{7}{\begin{gathered}
a_2' = a_0\,(V - W),\\
a_2'' = a_0\,(W - U),\\
a_2''' = a_0\,(U - V).
\end{gathered}}

If the $a_2', a_2'', a_2'''$ are known, then the $U, V, W$
can also be determined by
\weq{8}{\begin{gathered}
3\,a_0\,U = a_2''' - a_2''\\
3\,a_0\,V = a_2' - a_2'''\\
3\,a_0\,W = a_2'' - a_2',
\end{gathered}}
from which it appears that if two of the three quantities $a_2', a_2'', a_2'''$
are equal to each other, then one of the quantities $U, V, W$ vanishes
, and hence two of the roots of the biquadratic equation
are equal to each other, and consequently its discriminant vanishes.

If we further set
\weq{9}{\begin{gathered}
\alpha\beta + \gamma\delta = u,\\
\alpha\gamma + \delta\beta = v,\\
\alpha\delta + \beta\gamma = w,
\end{gathered}}
then
\[
a_0\,(u + v + w) = a_2,
\]
[the sum of the products of the roots of $f(x,\,y)$ taken two at a time],
\[
\begin{gathered}
U = v - w,\\
V = w - u,\\
W = u - v,
\end{gathered}
\]
and hence
\weq{10}{\begin{gathered}
3\,a_0\,u = a_2 - a_2',\\
3\,a_0\,v = a_2 - a_2'',\\
3\,a_0\,w = a_2 - a_2'''.
\end{gathered}}

Thus with the quantities $a_2', a_2'', a_2'''$ the $u, v, w$ are known
at the same time. But if one knows the quantities $u, v, w$, then
the solution of the biquadratic equation is reduced to square roots.
For we have
\[
\alpha\beta + \gamma\delta = u, \quad a_0\,\alpha\beta\,.\,\gamma\delta = a_4
\]
and from this
\[
2\,\alpha\beta = v_1 + \sqrt{\frac{a_0\,u^2 - 4\,a_4}{a_0}}, \quad
2\,\gamma\delta = v_1 - \sqrt{\frac{a_0\,u^2 - 4\,a_4}{a_0}}.
\]
% [resolved 2026-08-06, 5x zoom strip 5/5] v_1 is what is printed, in BOTH halves of this
% display, and the subscript is unambiguous at zoom. Weber introduces v without a subscript
% earlier, so the inconsistency is the author's, not the transcription's -- reproduced as
% printed, per the standing rule on his own oddities.

Since one can now in the same way determine the products
\[
\alpha\gamma, \quad \alpha\delta, \quad \beta\gamma, \quad \beta\delta
\]
%CONT>>


% ===== scan leaf 228, printed 203 =====
% !PRINTED 203
\weberpage{203}
\pdfbookmark[2]{§ 66. The covariants}{bm73}
%<<CONT
can be determined, then $\alpha^2$, for example, can be computed from the following quotient:
\[
\frac{\alpha\beta\,.\,\alpha\gamma}{\beta\gamma}
\]

It now remains, therefore, to form the cubic equation whose
roots are $a_2', a_2'', a_2'''$. But this is obtained at once from
the invariants.

If one forms the invariants $A', B'$ from the transformed
form (4), then by §.\,64, (8), (9), (12), since $r = 1$, it follows that
\weq{11}{\begin{gathered}
A = a_2'^{\,2} - 3\,a_1'\,a_3',\\
B = 2\,a_2'^{\,3} - 9\,a_1'\,a_2'\,a_3',\\
D = a_1'^{\,2}\,a_2'^{\,2}\,a_3'^{\,2} - 4\,a_1'^{\,3}\,a_3'^{\,3},
\end{gathered}}
and from these, by elimination of $a_1'\,a_3'$,
\[
a_2'^{\,3} - 3\,a_2'\,A + B = 0.
\]
Hence $a_2', a_2'', a_2'''$ are the roots of the equation, cubic in $z$,
\weq{12}{z^3 - 3\,A\,z + B = 0.}

This is surely the simplest form that can be given to the cubic
resolvent of the biquadratic equation.

\begin{center}
§. 66.

\gesperrt{The Covariants.}
\end{center}

If we form the covariant $H$, according to §.\,64, (2), for the transformed
 form
\weq{1}{f(x,\,y) = F(\xi,\,\eta) = a_1'\,\xi^3\eta + a_2'\,\xi^2\eta^2 + a_3'\,\xi\eta^3}
we obtain
\weq{2}{\begin{gathered}
-H = 3\,a_1'^{\,2}\,\xi^4 + 4\,a_1'\,a_2'\,\xi^3\eta - (6\,a_1'\,a_3' - 4\,a_2'^{\,2})\,\xi^2\eta^2\\
{} + 4\,a_2'\,a_3'\,\xi\eta^3 + 3\,a_3'^{\,2}\,\eta^4,
\end{gathered}}
and from this it follows
\weq{3}{H + 4\,a_2'\,f = -3\,(a_1'\,\xi^2 - a_3'\,\eta^2)^2.}

This combination is thus, apart from the factor $-3$, the
square of a quadratic form, and the same holds if we replace
$a_2'$ by $a_2'', a_2'''$. We therefore set, as an
abbreviation,
%CONT>>


% ===== scan leaf 229, printed 204 =====
% !PRINTED 204
\weberpage{204}
%<<CONT
\weq{4}{\begin{gathered}
H + 4\,a_2'\,f = -3\,\psi_1^2\\
H + 4\,a_2''\,f = -3\,\psi_2^2\\
H + 4\,a_2'''\,f = -3\,\psi_3^2.
\end{gathered}}

Of the three functions $\psi_1, \psi_2, \psi_3$ \gesperrt{no two can have
a common divisor}, provided the discriminant
$D$ is assumed different from zero. For
then the three quantities $a_2', a_2'', a_2'''$ are also different from one another.
If then $\psi_1$ and $\psi_2$ vanish, $H$ and
$f$ must vanish, i.\,e.\ a common divisor of $\psi_1$ and $\psi_2$
would have to be a common divisor of $H$ and $f$. But now $\xi$
was an arbitrary linear divisor of $f$. By (2) this can be
a divisor of $H$ only if $a_3' = 0$, that is, again, only if
the discriminant $D$ vanishes [§.\,65, (11)].

Now the functional determinant $T$ of $f$ and $H$ is identical
with the functional determinant of $f$ and $H + \lambda f$, whatever
$\lambda$ may be, and from this it follows, when $\lambda = 4\,a_2'$ is
set,
\[
T = -\tfrac{1}{2}\,\psi_1\,[F'(\xi)\,\psi_1'(\eta) - F'(\eta)\,\psi_1'(\xi)].
\]

Thus $T$ is divisible by $\psi_1$, and for the same reasons
by $\psi_2$ and $\psi_3$, and consequently $T$ is, up to a factor independent
of $\xi, \eta$, identical with the product $\psi_1\,\psi_2\,\psi_3$.
If we accordingly form the product of the three equations (4),
it follows, having regard to the cubic equation §.\,65, (12),
which the quantities $a_2', a_2'', a_2'''$ satisfy:
\[
H^3 - 48\,A\,H\,f^2 - 64\,B\,f^3 = c\,T^2,
\]
where $c$ is a constant still to be determined. By §.\,64, (12)
$c$ remains unchanged under a linear transformation, and one
therefore finds its value by setting the terms with the highest
power of $\xi$ on the two sides of the normal form (1), (2)
equal to each other, $c = -27$.

We thus have between the invariants and covariants
the following identical relation
\weq{5}{H^3 - 48\,A\,H\,f^2 - 64\,B\,f^3 = -27\,T^2.}

The functions $\psi_1, \psi_2, \psi_3$ are likewise covariants, though
with coefficients which are not rational in the coefficients
of $f$. We can easily express them in terms of the roots $\alpha, \beta, \gamma, \delta$
of $f = 0$. If for simplification one sets $y = 1$,
then by §.\,65, (3) and (5) it follows that
%CONT>>


% ===== scan leaf 230, printed 205 =====
% !PRINTED 205
\weberpage{205}
%<<CONT
\[
\psi_1 = a_1'\xi^2 - a_3'\eta^2
\]
\[
= a_0\,\frac{(\beta-\gamma)(\beta-\delta)(x-\alpha)^2 - (\alpha-\gamma)(\alpha-\delta)(x-\beta)^2}{\alpha-\beta}.
\]
In this expression, however, $\alpha - \beta$ can be cancelled from numerator and denominator, and
$\psi_1$ can be represented in two forms; to obtain $\psi_2$ and $\psi_3$, one then
need only permute $\beta$, $\gamma$, $\delta$ cyclically.

One thus finds
\[
\begin{aligned}
\frac{\psi_1}{a_0} &= (\gamma-\beta)(x-\alpha)(x-\delta) - (\alpha-\delta)(x-\beta)(x-\gamma) \\
&= (\delta-\beta)(x-\alpha)(x-\gamma) - (\alpha-\gamma)(x-\beta)(x-\delta)
\end{aligned}
\]
\weq{6}{
\begin{aligned}
\frac{\psi_2}{a_0} &= (\delta-\gamma)(x-\alpha)(x-\beta) - (\alpha-\beta)(x-\gamma)(x-\delta) \\
&= (\beta-\gamma)(x-\alpha)(x-\delta) - (\alpha-\delta)(x-\gamma)(x-\beta) \\
\frac{\psi_3}{a_0} &= (\beta-\delta)(x-\alpha)(x-\gamma) - (\alpha-\gamma)(x-\delta)(x-\beta) \\
&= (\gamma-\delta)(x-\alpha)(x-\beta) - (\alpha-\beta)(x-\delta)(x-\gamma).
\end{aligned}
}
The quantities $\psi_1^2$, $\psi_2^2$, $\psi_3^2$ may be regarded as the roots
of a cubic equation whose coefficients are covariants. One obtains
this equation easily from (4) in the form
\[
\left(z + \frac{H+4a_2'f}{3}\right)\left(z + \frac{H+4a_2''f}{3}\right)\left(z + \frac{H+4a_2'''f}{3}\right) = 0.
\]
By \S.\,65, (12), and making use of the expression (5) for $T^2$, this equation yields
\weq{7}{
z^3 + Hz^2 + \tfrac{1}{3}(H^2 - 16Af^2)z - T^2 = 0.
}
It is also of interest to form the discriminant of this equation. One
obtains it most simply from the expression
\[
\varDelta = (\psi_1^2 - \psi_2^2)^2(\psi_1^2 - \psi_3^2)^2(\psi_2^2 - \psi_3^2)^2,
\]
if one substitutes the expressions (4) and then, for $(a_2'-a_2'')^2(a_2'-a_2''')^2(a_2''-a_2''')^2$,
puts the discriminant of the equation \S.\,65, (12)
\[
27(4A^3 - B^2) = 3^6 D
\]
where then $D$ is the discriminant of $f$. One thus obtains
\weq{8}{
\varDelta = 2^{12} D f^6.
}


% ===== scan leaf 231, printed 206 =====
% !PRINTED 206
\weberpage{206}
\pdfbookmark[2]{§ 67. The complete system of invariants of the binary biquadratic form}{bm74}
\begin{center}
\S.\,67.

\gesperrt{The full system of invariants of the binary\\
biquadratic form.}
\end{center}

We wish further to prove that with the constructs $A$, $B$ the system of
independent invariants of the binary biquadratic form is exhausted.
But whereas one usually contents oneself with the proof that every
invariant is an integral rational function of $A$, $B$, we wish, as we
have already done in \S.\,63 for the covariants of the cubic form, also
to touch on the question of the numerical coefficients, and here a new
phenomenon shows itself.

The invariant $D$, for example, is indeed by \S.\,64, (11) expressed rationally
in terms of $A$, $B$. But the numerical coefficients are not integers; rather
they have the denominator 27, although the coefficients in the expanded function
$D$ are all integers.

We therefore now consider as integral invariants integral rational homogeneous
functions of degree $\mu$ of the five variables $a_0, a_1, a_2, a_3, a_4$
\weq{1}{
I(a_0, a_1, a_2, a_3, a_4) = I(a),
}
whose coefficients are integers, and to which the invariant property
\weq{2}{
I' = r^{2\mu} I
}
belongs, when $I'$ or $I(a')$ is the same function of the coefficients of a
transformed function.

Such invariants are $A$, $B$, $D$, between which there holds the relation
\weq{3}{
27 D = 4A^3 - B^2,
}
and our goal is to prove that all integral invariants are integral rational
functions, with integer coefficients, of these three.

If we form the function $I'$ in the normal form (4) of \S.\,65, then
we first obtain from (2), since $r = 1$,
\[
I(a) = I'(0, a_1', a_2', a_3', 0),
\]
and this function $I$ can depend only on $a_2'$ and the product $a_1'a_3'$;
for the substitution
%CONT>>


% ===== scan leaf 232, printed 207 =====
% !PRINTED 207
\weberpage{207}
%<<CONT
\[
\xi = \lambda\xi', \quad \eta = \frac{1}{\lambda}\eta',
\]
whose determinant is $= 1$, leaves the normal form \S.\,65, (4), and in it
$a_2'$, unchanged, while $a_1'$, $a_3'$ pass into $\lambda^2 a_1'$, $a_3' : \lambda^2$,
and here $\lambda$ is an arbitrary quantity.

If, then, we set
\[
a_2' = z
\]
we have
\weq{4}{
I = \varphi(a_1'a_3', z),
}
where $\varphi$ is an integral, rational function with integer coefficients of the two
arguments $a_1'a_3'$ and $z$.

But now, by \S.\,65, (11),
\[
3a_1'a_3' = z^2 - A,
\]
and if we therefore multiply (4) by a suitable power of 3, say $3^\nu$,
we obtain
\weq{5}{
3^\nu I = \psi(A, z),
}
where $\psi$ is again a function of $A$ and $z$ with integer coefficients. But by
\S.\,65, (12),
\[
z^3 = 3Az - B,
\]
and by means of this we can express all higher powers of $z$ through the first and
second, and thus obtain
\weq{6}{
3^\nu I = \chi(A, B, z),
}
where $\chi$ is at most of the second degree with respect to $z$.

But the left side of (6) remains unchanged when $z$ is replaced by $a_2'$,
$a_2''$, $a_2'''$, that is, by three different values, and
consequently the function $\chi$ cannot contain the variable $z$ at all any longer
(\S.\,30, II).

Therefore
\weq{7}{
3^\nu I = \chi(A, B),
}
where $\chi$ is an integral, rational function with integer coefficients.

If in (7), by (3), we set
\[
B^2 = 4A^3 - 27 D
\]
it follows that (7) has one of the two following forms
\weq{8}{
3^\nu I = \Phi(A, D), \quad B\Phi(A, D),
}
according as the degree of $I$ is even or odd.


% ===== scan leaf 233, printed 208 =====
% !PRINTED 208
\weberpage{208}
If $I$ is a \gesperrt{primitive} function of the $a$, that is, one whose
numerical coefficients have no common divisor, then the coefficients of the
functions $\Phi$ can in any case have no common divisor other than a power
of 3.

What we have to prove is that they are all divisible by $3^\nu$, or, what is
the same thing, that if we assume $\Phi(A, D)$ to be a primitive function,
then $\nu = 0$ must hold.

The question is therefore this: Can functions with the divisor 3 arise from
the primitive function $\Phi(A, D)$ or $B\Phi(A, D)$ when $A$, $B$, $D$ are
replaced by their expressions in the $a$? Since $B$ is primitive, by
\S.\,2 the function $B\Phi(A, D)$, expressed in the $a$, can have the
divisor 3 only if $\Phi(A, D)$ has it.

The question is thus reduced to this: Can the primitive function
$\Phi(A, D)$ acquire the divisor 3 when $A$, $D$ are replaced by their
expressions?

That this question must be answered in the negative we can easily see as
follows. We first imagine all those terms removed from $\Phi(A, D)$ whose
coefficients are divisible by 3; for evidently the remaining terms must
still retain the same property of acquiring the divisor 3 upon substitution
of the expressions for $A$, $D$.

Secondly, we may assume that $\Phi(A, D)$ does not have the factor $D$, for
by omitting this factor the property in question would, by the theorem
already mentioned (\S.\,2), not be destroyed.

But then a number divisible by 3 would have to arise from $\Phi(A, D)$ when
all the $a$, with the exception of $a_2$, are set equal to zero and
$a_2 = 1$. This gives $D = 0$, $A = 1$, and so $\Phi(1, 0)$ would have to be
a number divisible by 3; but this is the coefficient of the highest power of
$A$ in $\Phi(A, D)$, which by hypothesis is not divisible by 3.
With this we have therefore proved:

\gesperrt{Every integral invariant of the biquadratic form can be
represented rationally and with integer coefficients by $A$, $B$, $D$, and
indeed in one of the two forms}
\[
\Phi(A, D), \quad B\Phi(A, D).
\]


% ===== scan leaf 234, printed 209 =====
% !PRINTED 209
\weberpage{209}

That conversely every such expression, if it is homogeneous in
the coefficients $a$, represents an invariant with integer coefficients,
is clear of itself\footnote{The theory of invariants is presented in detail in
the works: \gesperrt{Clebsch}, ``Theorie der bin\"aren algebraischen Formen''.
Leipzig 1872. \gesperrt{Fa\`a di Bruno}, ``Einleitung in die Theorie der bin\"aren
Formen''. German by \gesperrt{Walter}. Leipzig 1881. \gesperrt{P. Gordans}, ``Vorlesungen
\"uber Invariantentheorie'', edited by \gesperrt{Kerschensteiner}. Leipzig
1885. Cf. also \gesperrt{Franz Meyer}, ``Bericht \"uber den gegenw\"artigen Stand
der Invariantentheorie'' in the Jahresbericht der deutschen Mathematiker-
Vereinigung 1890/91 (Berlin 1892).}.

% Printer's furniture excluded by policy: footer "Weber, Algebra. I." and gathering number 14.


% ===== scan leaf 235, printed 210 =====
\weberpage{210}% page number not printed on this page
\pdfbookmark[1]{Sixth Section. Tschirnhausen Transformation.}{bm75}
\pdfbookmark[2]{§ 68. Hermite's form of the Tschirnhausen transformation}{bm76}
% !PRINTED none

\begin{center}
\gesperrt{Sixth Section.}

\gesperrt{Tschirnhausen Transformation.}

\rule{4cm}{0.4pt}

\S.\,68.

\gesperrt{The Hermite Form of the Tschirnhausen Transformation.}
\end{center}

We have already become acquainted with the basic idea of the
Tschirnhausen transformation in the fourth section.

The task was this: to transform an algebraic equation of degree $n$
\weq{1}{f(x) = a_0 x^n + a_1 x^{n-1} + \cdots + a_{n-1} x + a_n = 0}
by a substitution of degree $(n-1)$
\weq{2}{y = \alpha_0 + \alpha_1 x + \alpha_2 x^2 + \cdots + \alpha_{n-1} x^{n-1}}
so as to gain, in the arbitrary coefficients $\alpha$, means of
simplifying the equation.

Hermite, by giving the substitution (2) a particular form, greatly
simplified this task and brought it into connection with the theory of
invariants\footnote{\gesperrt{Hermite}, Sur quelques th\'eor\`emes d'alg\`ebre et la r\'esolution de
l'\'equation du quatri\`eme degr\'e, published separately from the Comptes rendus of the Paris Academy. Paris 1859.}.

In \S.\,4 we became acquainted with a series of functions,
$f_0, f_1, f_2 \ldots f_{n-1}$, by which the powers of $x$ up to the
$(n-1)$st can be expressed rationally, and it is these functions
that Hermite uses for the representation of the substitution (2).

Here we denote by $x$ a root of the equation (1),
while $t$ is to be understood as an indeterminate variable.
Then by \S.\,4
%CONT>>


% ===== scan leaf 236, printed 211 =====
% !PRINTED 211
\weberpage{211}
%<<CONT
\weq{3}{\frac{f(t)}{t-x} = t^{n-1}f_0(x) + t^{n-2}f_1(x) + \cdots + t f_{n-2}(x) + f_{n-1}(x)}
and
\begin{align*}
f_0(x) &= a_0\\
f_1(x) &= a_0 x + a_1\\
f_2(x) &= a_0 x^2 + a_1 x + a_2 \tag{4}\\
&\ \cdots\cdots\cdots\cdots\cdots\cdots\cdots\cdots\cdots\cdots\cdots\cdots\\
f_{n-1}(x) &= a_0 x^{n-1} + a_1 x^{n-2} + a_2 x^{n-3} + \cdots + a_{n-1}.
\end{align*}
We now assume the substitution (2) in the form
\weq{5}{y = t_{n-1} f_0(x) + t_{n-2} f_1(x) + \cdots + t_1 f_{n-2}(x) + t_0 f_{n-1}(x),}
in which the $t_{n-1}, t_{n-2} \ldots t_1, t_0$ are the indeterminate quantities
that have taken the place of the $\alpha$, and are not to be confused
with the powers of $t$. But $y$ arises \gesperrt{from the expression (3)
when $t^k$ is replaced by $t_k$.}

If, as in \S.\,42, we denote by the sign $S$ placed before
a function that the sum is to be taken over all the roots
$x$ of equation (1), then we have
\begin{align*}
S[f_0(x)] &= n a_0\\
S[f_1(x)] &= (n-1) a_1\\
S[f_2(x)] &= (n-2) a_2 \tag{6}\\
&\ \cdots\cdots\cdots\cdots\cdots\cdots\cdots\\
S[f_{n-1}(x)] &= a_{n-1},
\end{align*}
hence
\weq{7}{S(y) = n a_0 t_{n-1} + (n-1) a_1 t_{n-2} + \cdots + 2 a_{n-2} t_1 + a_{n-1} t_0,}
an expression which results from $f'(t)$ when $t^k$ is replaced
by $t_k$.

If with the help of (7) we eliminate the variable $t_{n-1}$ from (5),
it follows that
\weq{8}{y - \frac{1}{n} S(y) = t_{n-2} F_0(x) + t_{n-3} F_1(x) + \cdots + t_0 F_{n-2}(x),}
if
\begin{align*}
F_0 &= f_1 - \frac{n-1}{n} a_1 = a_0 x + \frac{a_1}{n},\\
F_1 &= f_2 - \frac{n-2}{n} a_2 = a_0 x^2 + a_1 x + \frac{2 a_2}{n}, \tag{9}\\
&\ \cdots\cdots\cdots\cdots\cdots\cdots\cdots\cdots\cdots\cdots\cdots\cdots\\
F_{n-2} &= f_{n-1} - \frac{1}{n} a_{n-1} = a_0 x^{n-1} + a_1 x^{n-2} + \cdots + \frac{n-1}{n} a_{n-1},
\end{align*}
so that
%CONT>>


% ===== scan leaf 237, printed 212 =====
% !PRINTED 212
\weberpage{212}
\pdfbookmark[2]{§ 69. Invariant property of the Tschirnhausen transformation}{bm77}
%<<CONT
\weq{10}{S[F_0(x)] = 0, \quad S[F_1(x)] = 0, \ \ldots \ S[F_{n-2}(x)] = 0.}
If, in accordance with (3), we set
\begin{align*}
F(t, x) &= \frac{f(t)}{t-x} - \frac{1}{n} f'(t) \tag{11}\\
&= t^{n-2} F_0(x) + t^{n-3} F_1(x) + \cdots + t F_{n-3}(x) + F_{n-2}(x),
\end{align*}
then the right-hand side of (8) arises from $F(t, x)$ by the
replacement of $t^k$ by $t_k$.

If we assume $y$ from the outset in the form
\weq{12}{y = t_{n-2} F_0(x) + t_{n-3} F_1(x) + \cdots + t_0 F_{n-2}(x)}
then the equation
\weq{13}{S(y) = 0}
is identically satisfied. What advantage this form of the substitution
affords, the next considerations will show.

\begin{center}
\S.\,69.

\gesperrt{Invariant Property of the Tschirnhausen Transformation.}
\end{center}

We have now to investigate the influence which a linear
transformation, to which we subject the function $f(x)$, has on the
Tschirnhausen transformation.

We make in $f(x)$ the linear substitution
\weq{1}{x = \frac{\alpha \xi + \beta}{\gamma \xi + \delta}, \quad \alpha \delta - \beta \gamma = r,}
whereby we obtain
\weq{2}{\varphi(\xi) = (\gamma \xi + \delta)^n f\left(\frac{\alpha \xi + \beta}{\gamma \xi + \delta}\right),}
so that $\varphi(\xi)$ is an integral rational function of degree $n$, and
$\varphi(\xi) = 0$ is the equation $f(x) = 0$ transformed by (1).

We now derive a function $\Phi(\tau, \xi)$ from $\varphi(\xi)$ in exactly the same
way as we derived $F(t, x)$ from
$f(x)$ in the preceding paragraph, namely
\begin{align*}
\Phi(\tau, \xi) &= \frac{\varphi(\tau)}{\tau - \xi} - \frac{1}{n} \varphi'(\tau) \tag{3}\\
&= \tau^{n-2} \Phi_0(\xi) + \tau^{n-3} \Phi_1(\xi) + \cdots + \tau \Phi_{n-3}(\xi) + \Phi_{n-2}(\xi),
\end{align*}
in which $\tau$, just like $t$, is a variable, and the functions
$\Phi_0(\xi), \Phi_1(\xi) \ldots \Phi_{n-2}(\xi)$, likewise from $\xi$ and the coefficients
%CONT>>


% ===== scan leaf 238, printed 213 =====
% !PRINTED 213
\weberpage{213}
%<<CONT
of $\varphi(\xi)$, just as the corresponding functions
$F_0(x), F_1(x) \cdots F_{n-2}(x)$ are formed from $x$ and the coefficients of $f(x)$.

If now, on the other hand, we carry out in $F(t, x)$, simultaneously with the
substitution (1), the substitution

\weq{4}{t = \frac{\alpha\tau+\beta}{\gamma\tau+\delta}}

then we obtain

\weq{5}{(\gamma\tau+\delta)^n f(t) = \varphi(\tau),}
\weq{6}{(\gamma\xi+\delta)(\gamma\tau+\delta)(t-x) = r(\tau-\xi)}

and, by forming from (5) the derivative, most simply by
differentiation by means of the formula

\[
\frac{dt}{d\tau} = \frac{r}{(\gamma\tau+\delta)^2}
\]

we obtain

\weq{7}{r(\gamma\tau+\delta)^{n-1}\frac{1}{n}f'(t) = (\gamma\tau+\delta)\frac{1}{n}\varphi'(\tau) - \gamma\varphi(\tau).}

From (5), (6) and (7) it follows that

\[
r(\gamma\tau+\delta)^{n-2}\frac{f(t)}{t-x} = \frac{\varphi(\tau)}{\tau-\xi}\frac{\gamma\xi+\delta}{\gamma\tau+\delta}
\]

\[
r(\gamma\tau+\delta)^{n-2}\frac{1}{n}f'(t) = -\frac{\gamma(\tau-\xi)\varphi(\tau)}{(\gamma\tau+\delta)(\tau-\xi)} + \frac{1}{n}\varphi'(\tau),
\]

and by subtraction of the two formulas

\[
r(\gamma\tau+\delta)^{n-2}\left(\frac{f(t)}{t-x} - \frac{1}{n}f'(t)\right) = \frac{\varphi(\tau)}{\tau-\xi} - \frac{1}{n}\varphi'(\tau),
\]

hence

\weq{8}{r(\gamma\tau+\delta)^{n-2}F(t, x) = \Phi(\tau, \xi)}

or, written more explicitly,

\weq{9}{\tau^{n-2}\Phi_0 + \tau^{n-3}\Phi_1 + \cdots + \Phi_{n-2} = r[(\alpha\tau+\beta)^{n-2}F_0 +(\alpha\tau+\beta)^{n-3}(\gamma\tau+\delta)F_1 +\cdots+(\gamma\tau+\delta)^{n-2}F_{n-2}].}

We now set, replacing in $F(t, x)$ the powers $t^k$ by
arbitrary variables $t_k$, and likewise in $\Phi(\tau, k)$ % [sic] printed \textquotedbl{}$\Phi(\tau,k)$\textquotedbl{}, not $\Phi(\tau,\xi)$
$\tau^k$ by $\tau_k$,

\weq{10}{\begin{aligned}
Y(t, x) &= t_{n-2}F_0 + t_{n-3}F_1 + \cdots + t_0 F_{n-2},\\
H(\tau, \xi) &= \tau_{n-2}\Phi_0 + \tau_{n-3}\Phi_1 + \cdots + \tau_0\Phi_{n-2},
\end{aligned}}

so that the substitutions of the Tschirnhausen transformation

\weq{11}{y = Y(t, x), \qquad \eta = H(\tau, \xi)}

read [\S.\,68, (12)].
%CONT>>


% ===== scan leaf 239, printed 214 =====
% !PRINTED 214
\weberpage{214}
%<<CONT
Equation (9), which is an identity with respect to $\tau$, remains
valid, however, when after carrying out the indicated raising to
powers $\tau^k$ is replaced on the right by $\tau_k$, and it therefore teaches
that between the functions $Y(t, x)$, $H(\tau, \xi)$ the relation holds

\weq{12}{H(\tau, \xi) = r\, Y(t, x),}

if we let the $t$ depend on the $\tau$ in the following way:
Set

\weq{13}{\begin{aligned}
t_{n-2} &= (\alpha\tau+\beta)^{n-2}\\
t_{n-3} &= (\alpha\tau+\beta)^{n-3}(\gamma\tau+\delta)\\
&\ \cdots\cdots\cdots\cdots\cdots\cdots\cdots\cdots\cdots\\
t_0 &= (\gamma\tau+\delta)^{n-2},
\end{aligned}}

and \gesperrt{after carrying out the powers} replace $\tau^k$ by $\tau_k$.

Now let $z$ be an arbitrary variable, and before we carry out
the raising to powers in (13), let us multiply these equations
in order by

\[
1, \quad -(n-2)z, \quad \frac{(n-2)(n-3)}{1\,.\,2}z^2, \cdots \pm z^{n-2};
\]

then on the left we obtain an integral rational function of $z$
of degree $(n-2)$

\weq{14}{\begin{aligned}
T(z) &= t_{n-2} - (n-2)t_{n-3}z\\
&\quad + \frac{(n-2)(n-3)}{1\,.\,2}t_{n-4}z^2 - \cdots \pm t_0 z^{n-2},
\end{aligned}}

and the right side yields, by the binomial theorem,

\[
[(\alpha\tau+\beta) - z(\gamma\tau+\delta)]^{n-2} = [\tau(\alpha-\gamma z) + (\beta-\delta z)]^{n-2},
\]

which, if we set

\weq{15}{z = \frac{\alpha\xi+\beta}{\gamma\xi+\delta}, \qquad \xi = \frac{\delta z-\beta}{-\gamma z+\alpha}}

goes over into

\[
(\alpha-\gamma z)^{n-2}(\tau-\xi)^{n-2} = \frac{r^{n-2}(\tau-\xi)^{n-2}}{(\gamma\xi+\delta)^{n-2}}
\]

If we carry out the power $(\tau-\xi)^{n-2}$ here by the binomial
theorem, and then replace $\tau^k$ by $\tau_k$, we obtain a transformation
of the function $T(z)$. We therefore set

% [sic] The print has a plus sign before \cdots; the binomial expansion and (14) require a minus sign there.

\[
\begin{aligned}
\Theta(\xi) &= \tau_{n-2} - (n-2)\tau_{n-3}\xi\\
&\quad + \frac{(n-2)(n-3)}{1\,.\,2}\tau_{n-4}\xi^2 + \cdots \pm \tau_0\xi^{n-2}
\end{aligned}
\]

and obtain the identical transformation
%CONT>>


% ===== scan leaf 240, printed 215 =====
% !PRINTED 215
\weberpage{215}
\pdfbookmark[2]{§ 70. Elaborations on Hermite's theorem}{bm78}
%<<CONT

\weq{16}{(\gamma\xi+\delta)^{n-2}\, T\!\left(\frac{\alpha\xi+\beta}{\gamma\xi+\delta}\right) = r^{n-2}\Theta(\xi),}

in which the connection between the coefficients $t$ and $\tau$
is expressed by the symbolic equations (13), and is therefore the
same as in the functions $Y$ and $H$.

\gesperrt{Accordingly, $r^{n-2}\Theta(\xi)$ is the transformation of the
function $(\gamma\xi+\delta)^{n-2}T(z)$ by the same linear substitution
by which $(\gamma\xi+\delta)^n f(z)$ passes over into $\varphi(\xi)$.}

If we denote the coefficients of this transformed function,
that is, the quantities $r^{n-2}\tau_k$, by $t'$, then equation (12),
which is linear with respect to $\tau$, yields

\weq{17}{H(t', \xi) = r^{n-1}Y(t, x).}

Now from $Y(t, x)$ there arise $n$ functions when one puts
for $x$ the $n$ roots of $f(x)$. Every homogeneous, rational,
symmetric function of these $n$ functions is rationally expressible
in terms of the coefficients of $f(x)$; and if we therefore denote such
a function by $K(t, a)$, then equation (17) yields, if we denote
the coefficients of $\varphi(\xi)$ by $a'$,

\weq{18}{K(t', a') = r^{\nu(n-1)}K(t, a),}

where $\nu$ denotes the degree of the symmetric function. With
this the beautiful theorem of \gesperrt{Hermite} is proved:

\gesperrt{The coefficients in the equation $f(x) = 0$ transformed
by the Tschirnhausen transformation}

\[
y = Y(t, x)
\]

\gesperrt{and all symmetric functions of the $n$ values $y$ are
simultaneous invariants of the two functions}

\[
f(z), \quad T(z).
\]

\begin{center}
\S.\,70.\\[2pt]
\gesperrt{Developments concerning the Hermite theorem.}
\end{center}

From the considerations of the last paragraph it has
emerged that if we set

\weq{1}{y = t_{n-2}F_0(x) + t_{n-3}F_1(x) \cdots + t_0 F_{n-2}(x)}

then the symmetric functions of the $n$ values $y_1, y_2 \cdots y_n$,
which one obtains from $y$ when one puts for $x$ the $n$ roots of
%CONT>>


% ===== scan leaf 241, printed 216 =====
% !PRINTED 216
\weberpage{216}
%<<CONT
$f(x) = 0$, are simultaneous invariants of two forms of degrees $n$ and
$(n-2)$, $f(z)$, $T(z)$.

Such an integral rational and homogeneous symmetric
function of order $\nu$

\[
P(y_1, y_2 \cdots y_n)
\]

is, by (1), evidently a homogeneous function of order $\nu$ of the
variables $t$. It is likewise a homogeneous function of order $\nu$
in the $a$, since the expression (1), so far as it contains the $a$ explicitly,
is linear in them, and since the symmetric functions of the $x$ depend only
on the ratios $a_1 : a_0, a_2 : a_0, \ldots a_n : a_0$. In the
expression of $P$ by the $a$, however, a power of $a_0$ might possibly
remain in the denominator, and we still have to show
that this does not occur.

If we set

\[
a_0^{\lambda} P(y_1, \ldots y_n) = K(t, a)
\]

and determine the power $a_0^{\lambda}$ so that $K(t, a)$ is an integral function
of the $a$, but no longer vanishes for $a_0 = 0$, then it
will follow that $\lambda = 0$ must hold, if we show that for
$a_0 = 0$ all the $y$ retain finite values.

If we let $a_0 = 0$, while the remaining $a$ stay
unchanged, then one of the roots $x$ will, as we have seen in
\S.\,40, grow infinite. But that the corresponding
$y$ nevertheless remains finite follows from the equation
\S.\,68, (11), which is an identity with respect to $t$,

\weq{2}{\frac{f(t)}{t-x} - \frac{1}{n}f'(t) = t^{n-2}F_0(x) + t^{n-3}F_1(x) + \cdots + F_{n-2}(x),}

from which it can be seen that, when $a_0 = 0$ and $x$ becomes infinite,
the functions

\[
F_0(x), \quad F_1(x) \cdots F_{n-2}(x)
\]

become equal to the coefficients of $-\dfrac{1}{n}f'(t)$, that is, equal to

\[
0, \quad -\frac{n-1}{n}a_1 \ldots, \quad -\frac{1}{n}a_{n-1}
\]

Thus the $y$ also remain finite, and $P(y_1, y_2 \cdots y_n)$
is equal to an \gesperrt{integral rational invariant} of degree $\nu$
both in the $t$ and in the $a$.

If we therefore assume the equation for $y$ in the form

\weq{3}{y^n + P_2 y^{n-2} + P_3 y^{n-3} + \cdots + P_n = 0,}

then $P_{\nu}$ is such an invariant of degree $\nu$.


% ===== scan leaf 242, printed 217 =====
% !PRINTED 217
\weberpage{217}

An invariant of degree $n(n-1)$ is also the discriminant $\Delta$ of equation (3), or the square of the difference product

\weq{4}{\begin{gathered} \Pi = (y_1 - y_2)(y_1 - y_3)\cdots \\ (y_2-y_3)\cdots \end{gathered}}

To obtain closer insight into the formation of this quantity, we recall that we obtain $y$ when, in

\[ \frac{f(t)}{t-x} - \frac{1}{n}f'(t) \]

we replace $t^k$ by $t_k$. We therefore obtain $y_1-y_2$ from the formula

\weq{5}{y_1-y_2 = (x_1-x_2)\frac{f(t)}{(t-x_1)(t-x_2)},}

when we make the same replacement.

The quotient

\[ \frac{f(t)}{(t-x_i)(t-x_k)} \]

is an integral rational function of $t$ of degree $n-2$. If we replace $t^k$ in it by $t_k$, let it pass over into $t_{i,k}$; we then have

\weq{6}{y_i - y_k = (x_i-x_k)t_{i,k}.}

If we therefore set

\weq{7}{\begin{gathered} a_0^{n-1}\Theta = t_{1,2}t_{1,3}\ldots t_{1,n} \\ t_{2,3}\ldots t_{2,n} \\ t_{n-1,n} \end{gathered}}

then $\Theta$ is a homogeneous integral function of degree $\frac{1}{2}n(n-1)$ with respect to the $t$, and moreover, as a symmetric function of the $x$, it can be expressed rationally in terms of the coefficients $a$.

But from (6) it follows, if $D$, as in \S.\,46, (3), denotes the discriminant of $f(x)$, that

\weq{8}{\Delta = D\Theta^2,}

from which it is to be concluded that $\Theta$ is an invariant \gesperrt{which no longer contains the denominator $a_0$}. $\Delta$ is of degree $n(n-1)$ in the coefficients $a$, while $D$ by \S.\,46 is of degree $2n-2$; from this it follows that $\Theta$ is of degree $\frac{1}{2}(n-1)(n-2)$ in the coefficients $a$.


% ===== scan leaf 243, printed 218 =====
% !PRINTED 218
\weberpage{218}

$\Theta$ is a so-called decomposable form of degree $\frac{1}{2}n(n-1)$ in the variables $t$; for by (7) it can be decomposed entirely into linear factors, which, to be sure, are not rational in the $a$.

We have already pointed out in \S.\,52 that in the Tschirnhausen transformation $x$ is also expressible rationally in terms of $y$; and the same therefore holds for every rational function of $x$ as well.

If we consider any such function $\varphi(x)$, which may also contain the coefficients $a$ and $t$, but is always assumed to be integral and homogeneous in both kinds of variables, then we may set

\weq{9}{\varphi(x) = C_0 + C_1y + C_2y^2 + \cdots + C_{n-1}y^{n-1},}

in which the $C_0, C_1, \ldots C_{n-1}$ are expressible rationally in $a$ and in $t$.

If we write down equation (9) for $x=x_1,x_2\ldots x_n$, $y=y_1,y_2,\ldots y_n$, we obtain for the determination of the $C$ a system of linear equations whose determinant

\weq{10}{\begin{vmatrix} 1, & y_1, & y_1^2 & \ldots & y_1^{n-1} \\ 1, & y_2, & y_2^2 & \ldots & y_2^{n-1} \\ \cdot & \cdot & \cdot & \cdot & \cdot \\ 1, & y_n, & y_n^2 & \ldots & y_n^{n-1} \end{vmatrix},}

is therefore equal to $\Theta\sqrt{D}$. The numerator of the expression for $C_\nu$ arises from the determinant (10) by replacing the elements of the $(\nu+1)$th column by $\varphi(x_1), \varphi(x_2)\ldots\varphi(x_n)$, and consequently the numerator, when everything is expressed in terms of the $x_i$, has the factor $\sqrt{D}$ (because it vanishes when two of the $x_i$ become equal to each other).

We may therefore set

\[ \Theta C_\nu = Q_\nu, \]

in which $Q_\nu$ is an integral rational function of the $t$ and the $a$, which can at most still contain a power of $a_0$ in the denominator. We then obtain

\weq{11}{\Theta\varphi(x) = Q_0 + Q_1y + Q_2y^2 + \cdots + Q_{n-1}y^{n-1}.}

The coefficients $Q_\nu$, however, are not invariants, and their computation is in most cases difficult.


% ===== scan leaf 244, printed 219 =====
% !PRINTED 219
\weberpage{219}
\pdfbookmark[2]{§ 71. Transformation of the cubic equation}{bm79}

\begin{center}
\S.\,71.

\gesperrt{Transformation of the cubic equation.}
\end{center}

The \gesperrt{theorem of Hermite} (\S.\,69) suffices, for the cubic equation, to carry out the transformation without further computation.

For this purpose, if we employ homogeneous variables $z_1, z_2$, we have to form simultaneous invariants of the cubic form

\weq{1}{f(z_1,z_2) = a_0z_1^3 + a_1z_1^2z_2 + a_2z_1z_2^2 + a_3z_2^3}

and of the linear form

\weq{2}{T = -t_0z_1 + t_1z_2}

But if we apply the linear substitution

\weq{3}{\begin{gathered} z_1 = \alpha z_1' + \beta z_2' \\ z_2 = \gamma z_1' + \delta z_2' \end{gathered}}

to $T$, there results

\[ -t_0z_1 + t_1z_2 = -t_0'z_1' + t_1'z_2', \]

provided that

\[ \begin{gathered} t_0' = +t_0\alpha - t_1\gamma, \\ t_1' = -t_0\beta + t_1\delta, \end{gathered} \]

or

\weq{4}{\begin{gathered} rt_1 = \alpha t_1' + \beta t_0' \\ rt_0 = \gamma t_1' + \delta t_0'. \end{gathered}}

But this substitution arises from (3) when one replaces $z_1, z_2, z_1', z_2'$ by $rt_1, rt_0, t_1', t_0'$. Now if $I(a,t_1,t_0)$ is a simultaneous invariant of $f$ and $T$, homogeneous and of degree $\nu$ in $t_0, t_1$ and of weight $\lambda$, so that

\[ I(a',t_1',t_0') = r^\lambda I(a,t_1,t_0), \]

then by this permutation it follows that

\[ I(a',z_1',z_2') = r^{\lambda-\nu}I(a,z_1,z_2), \]

i.\,e.\ $I(a,z_1,z_2)$ is a covariant of $f$.

\gesperrt{One therefore obtains all simultaneous invariants of $f$ and $T$ from the covariants of $f$, by replacing $z_1, z_2$ in them by $t_1, t_0$.}

But the covariants of $f$ we have come to know completely in \S.\,62, 63.

After this it is easy to form the cubic equation which, from the substitution \S.\,68, (12),
%CONT>>


% ===== scan leaf 245, printed 220 =====
% !PRINTED 220
\weberpage{220}
%<<CONT

\weq{5}{y = t_1\Bigl(a_0x + \frac{a_1}{3}\Bigr) + t_0\Bigl(a_0x^2 + a_1x + \frac{2a_2}{3}\Bigr)}

results for $y$. If we write the equation in the form

\[ y^3 + P_2y + P_3 = 0, \]

then $P_2, P_3$ are covariants of $f$, and indeed $P_2$ of the second, $P_3$ of the third order, both in $t$ and in $a$.

Now since $P_2$ and $P_3$ are covariants of the cubic form, by \S.\,62 and 63 they can differ from the two functions defined there,

\[ H(t_1,t_0), \qquad Q(t_1,t_0) \]

only by constant, that is, numerical factors. These constant factors can be determined by some special assumption.

We may, for example, assume

\[ t_1 = 1, \quad t_0 = 0, \quad a_0 = 1, \quad a_1 = 0, \]

and then $y=x$, hence

\[ P_2 = a_2, \qquad P_3 = a_3. \]

On the other hand, by the formulas (2), (3), (7) given in \S.\,62, this particular assumption yields

\[ H(1,0) = 3a_2, \qquad Q(1,0) = 27a_3, \]

from which one concludes in general

\[ P_2 = \frac{1}{3}H(t_1,t_0), \qquad P_3 = \frac{1}{27}Q(t_1,t_0), \]

so that we obtain for $y$ the cubic equation

\weq{6}{y^3 + \frac{1}{3}H(t_1,t_0)y + \frac{1}{27}Q(t_1,t_0) = 0.}

The discriminant $\Delta$ of this cubic equation is

\[ \Delta = -\frac{1}{27}(4H^3 + Q^2), \]

hence, applying the relation \S.\,62, (12),

\weq{7}{4H^3 + Q^2 + 27Df^2 = 0}

\weq{8}{\Delta = Df(t_1,t_0)^2,}

if $D$ is the discriminant of the given cubic equation, in agreement with the general results of the preceding section.

If we wish to base the solution of the cubic equation on this, we must first, following the prescriptions of the preceding section, express $x$ rationally in terms of $y$. We set

\weq{9}{a_0f(t_1,t_0)x = Q_0 + Q_1y + Q_2y^2.}


% ===== scan leaf 246, printed 221 =====
% !PRINTED 221
\weberpage{221}

The computation of the coefficients $Q$ is easy to carry out if one sets up this equation for $x_1$, $x_2$, $x_3$ and correspondingly $y_1$, $y_2$, $y_3$.

We shall give a few indications of the course of the computation, which requires only the representation of symmetric functions of the roots of a cubic equation by the coefficients, according to the prescriptions of the fourth section (\S.\,42, 45). If one first takes the sum of the three equations (9), one obtains by (6)
\[
3 Q_0 = \frac{2}{3} H(t_1, t_0) Q_2 - a_1 f(t_1, t_0),
\]
so that only $Q_1$ and $Q_2$ still need to be computed; these one finds by multiplying the equation (9) formed for $x = x_1$ by
\[
y_2 - y_3 = a_0 (x_2 - x_3) (t_1 - t_0 x_1)
\]
and by
\[
y_2^2 - y_3^2 = -a_0 (x_2 - x_3) y_1 (t_1 - t_0 x_1)
\]
and adding the three analogous equations formed by cyclic permutation of the indices 1, 2, 3.

One then has only to make use of the two formulae
\[
a_0^2 \Sigma x_1^2 (x_2 - x_3) = -\sqrt{D}
\]
\[
a_0^3 \Sigma x_1^3 (x_2 - x_3) = a_1 \sqrt{D},
\]
and thus finds
\[
Q_0 = \frac{-a_1}{3} f(t_1, t_0) - \frac{2}{3} t_0 H(t_1, t_0)
\]
\[
Q_1 = a_0 t_1^2 + \frac{2}{3} a_1 t_0 t_1 + \frac{1}{3} a_2 t_0^2
\]
\[
Q_2 = -t_0.
\]

If for simplification we set $t_1 = t$, $t_0 = 1$, thus passing to the inhomogeneous expressions, we obtain
\weq{10}{
(3 a_0 x + a_1) f(t) = -\frac{2}{3} H(t) + y f'(t) - 3 y^2,
}
while the equation for $y$ reads
\weq{11}{
y^3 + \frac{1}{3} H(t) y + \frac{1}{27} Q(t) = 0,
}
if we write $H(t)$, $Q(t)$ for $H(t, 1)$, $Q(t, 1)$.

Now, in order to solve the cubic equation, we determine $t$ from the quadratic equation
\[
H(t) = 0.
\]


% ===== scan leaf 247, printed 222 =====
% !PRINTED 222
\weberpage{222}

For the moment, as an abbreviation, we shall set
\weq{12}{
\varrho = \sqrt{-3 D}
}
then from (7) it follows that
\weq{13}{
Q = 3 \varrho f(t),
}
and from (11)
\weq{14}{
y = -\frac{1}{3} \sqrt[3]{3 \varrho f(t)},
}
and the expression (10) yields
\weq{15}{
(3 a_0 x + a_1) f(t) = y f'(t) - 3 y^2.
}
In order to derive the agreement of this result with the \gesperrt{Cardano formula}, we go back to \S.\,62 and set $x = t$, $y = 1$ in the formulas there.

On account of (13), by \S.\,62, (11), (13), we get
\[
\eta = 0, \quad \xi = -h \varrho,
\]
and by \S.\,62, (10)
\[
f(t) = \xi^3, \quad f'(t) = 3 \xi^2 \xi',
\]
in which $\xi' = h A_0$ is, by \S.\,62, (13), the derivative of $\xi$ with respect to $t$,
\[
f(t) = -h^3 \varrho^3, \quad f'(t) = 3 h^3 A_0 \varrho^2.
\]
Accordingly it follows from (14) that
\[
y = \frac{h \varrho \sqrt[3]{3 \varrho}}{3}
\]
and from (15)
\[
3 a_0 x + a_1 = -h A_0 \sqrt[3]{3 \varrho} - \frac{1}{h \sqrt[3]{3 \varrho}},
\]
or, by \S.\,62, (14), (15),
\weq{16}{
3 a_0 x + a_1 = A_0 (k - h) \sqrt[3]{3 \varrho},
}
in which, by \S.\,62, (15),
\[
h A_0 \sqrt[3]{3 \varrho} = \sqrt[3]{\frac{-q_0 + 3 \varrho a_0}{2}}
\]
\[
k A_0 \sqrt[3]{3 \varrho} = \sqrt[3]{\frac{+q_0 + 3 \varrho a_0}{2}}.
\]

If we assume $a_0 = 1$, $a_1 = 0$, then from (16) we obtain the \gesperrt{Cardano formula} (\S.\,35)
\weq{17}{
x = \sqrt[3]{\frac{-a_3}{2} + \sqrt{\frac{a_2^3}{27} + \frac{a_3^2}{4}}} + \sqrt[3]{\frac{-a_3}{2} - \sqrt{\frac{a_2^3}{27} + \frac{a_3^2}{4}}}.
}


% ===== scan leaf 248, printed 223 =====
% !PRINTED 223
\weberpage{223}
\pdfbookmark[2]{§ 72. General execution of the transformation}{bm80}

\begin{center}
\S.\,72.
\end{center}

\begin{center}
\gesperrt{General execution of the transformation.}
\end{center}

In the general carrying-out of the \gesperrt{Tschirnhausen} transformation in the \gesperrt{Hermite} form we can go a significant step further.

We first consider the general substitution \S.\,68, (5)
\weq{1}{
y = t_{n-1} f_0 + t_{n-2} f_1 + t_{n-3} f_2 + \cdots + t_0 f_{n-1},
}
from which we obtain the more special form \S.\,68 (12) if we bind the variables $t$ to the one linear condition
\weq{2}{
S(y) = n a_0 t_{n-1} + (n-1) a_1 t_{n-2} + \cdots + a_{n-1} t_0 = 0
}
— which, however, is not yet to be done for the present.

By means of the functions $f_0, f_1, f_2 \ldots f_{n-1}$, according to \S.\,68, (4), all powers of $x$, and hence also all rational functions of $x$, can be represented \gesperrt{linearly} and homogeneously.

If we can find this representation for the various powers of $y$, then by elimination of the $f_i$ the equation of degree $n$ for $y$ can be formed.

But we attain the same end still more simply if we represent the $n$ products $y f_s$ linearly in terms of the $f$. For if we have the equations
\weq{3}{
y f_s = E_{0,s} f_0 + E_{1,s} f_1 + E_{2,s} f_2 + \cdots + E_{n-1,s} f_{n-1}
}
\[
s = 0, 1, 2 \ldots n-1,
\]
then by elimination of $f_0, f_1 \ldots f_{n-1}$ from the system (3) we obtain
\weq{4}{
E =
\begin{vmatrix}
E_{0,0} - y, & E_{1,0}, & E_{2,0} & \cdots & E_{n-1,0} \\
E_{0,1}, & E_{1,1} - y, & E_{2,1} & \cdots & E_{n-1,1} \\
\cdots & \cdots & \cdots & \cdots & \cdots \\
E_{0,n-1}, & E_{1,n-1}, & E_{2,n-1} & \cdots & E_{n-1,n-1} - y
\end{vmatrix}
= 0,
}
that is, an equation of degree $n$ for $y$, which is the transformed equation sought. Everything is therefore reduced to the determination of the coefficients $E_{r,s}$, which, as appears from (1), are in any case linear functions of the variables $t$. In particular,
\weq{5}{
E_{0,0} = a_0 t_{n-1}, \quad E_{1,0} = a_0 t_{n-2}, \quad \ldots \quad E_{n-1,0} = a_0 t_0.
}

To compute the remaining $E$, we observe that by the definition \S.\,68, (4) the following relations hold between the functions $f$:
%CONT>>


% ===== scan leaf 249, printed 224 =====
% !PRINTED 224
\weberpage{224}
%<<CONT
\weq{6}{
\begin{aligned}
x f_0 &= f_1 - a_1 \\
x f_1 &= f_2 - a_2 \\
&\cdots\cdots\cdots\cdots \\
x f_{n-2} &= f_{n-1} - a_{n-1} \\
x f_{n-1} &= -a_n;
\end{aligned}
}
of which the last is a consequence of the equation $f(x) = 0$.

If we now continue the series of variables $t_0, t_1 \ldots t_{n-1}$ by defining the new variables $t_n, t_{n+1}, t_{n+2} \ldots$ through the following equations
\weq{7}{
\begin{aligned}
a_0 t_n &+ a_1 t_{n-1} + a_2 t_{n-2} + \cdots + a_n t_0 = 0, \\
a_0 t_{n+1} &+ a_1 t_n + a_2 t_{n-1} + \cdots + a_n t_1 = 0, \\
a_0 t_{n+2} &+ a_1 t_{n+1} + a_2 t_n + \cdots + a_n t_2 = 0, \\
&\cdots\cdots\cdots\cdots\cdots\cdots\cdots
\end{aligned}
}
then, multiplying the equations (6) in order by $t_{n-1}, t_{n-2} \ldots t_0$ and adding, we obtain
\[
x y = t_n f_0 + t_{n-1} f_1 + \cdots + t_1 f_{n-1},
\]
and if we now multiply (6) in the same way by $t_n, t_{n-1} \ldots t_1$,
\[
x^2 y = t_{n+1} f_0 + t_n f_1 + \cdots + t_2 f_{n-1}.
\]
Thus the following system results
\weq{8}{
\begin{aligned}
y &= t_{n-1} f_0 + t_{n-2} f_1 + \cdots + t_0 f_{n-1} \\
x y &= t_n f_0 + t_{n-1} f_1 + \cdots + t_1 f_{n-1} \\
x^2 y &= t_{n+1} f_0 + t_n f_1 + \cdots + t_2 f_{n-1} \\
&\cdots\cdots\cdots\cdots\cdots\cdots\cdots\cdots \\
x^s y &= t_{n+s-1} f_0 + t_{n+s-2} f_1 + \cdots + t_s f_{n-1},
\end{aligned}
}
and if we multiply these equations in turn by $a_s, a_{s-1} \ldots a_0$ and add, there follows the expression sought
\[
y f_s = E_{0,s} f_0 + E_{1,s} f_1 + \cdots + E_{n-1,s} f_{n-1},
\]
in which, by (7),
\weq{9}{
\begin{aligned}
E_{0,s} &= a_0 t_{n+s-1} + a_1 t_{n+s-2} + \cdots + a_s t_{n-1} \\
&= -a_{s+1} t_{n-2} \cdots - a_n t_{s-1} \\
E_{1,s} &= a_0 t_{n+s-2} + a_1 t_{n+s-3} + \cdots + a_s t_{n-2} \\
&= -a_{s+1} t_{n-3} \cdots - a_n t_{s-2} \\
&\cdots\cdots\cdots\cdots\cdots\cdots\cdots\cdots \\
E_{s-1,s} &= a_0 t_n + a_1 t_{n-1} + \cdots + a_s t_{n-s} \\
&= -a_{s+1} t_{n-s-1} \cdots - a_n t_0 \\
E_{s,s} &= a_0 t_{n-1} + a_1 t_{n-2} + \cdots + a_s t_{n-s-1} \\
E_{s+1,s} &= a_0 t_{n-2} + a_1 t_{n-3} + \cdots + a_s t_{n-s-2} \\
&\cdots\cdots\cdots\cdots\cdots\cdots\cdots\cdots \\
E_{n-1,s} &= a_0 t_s + a_1 t_{s-1} + \cdots + a_s t_0.
\end{aligned}
}
%CONT>>


% ===== scan leaf 250, printed 225 =====
% !PRINTED 225
\weberpage{225}
\pdfbookmark[2]{§ 73. The Bezoutiant}{bm81}
To this it should further be remarked that the superfluously introduced
variables $t_n, t_{n+1} \ldots$ have already been eliminated again in the
second form of $E_{0,s}, E_{1,s} \ldots E_{s-1,s}$, and that the variable
$t_{n-1}$ occurs only in $E_{s,s}$.

\begin{center}
\S.\,73.

\gesperrt{The Bezoutiant.}
\end{center}

Equation (4) of the preceding section, when expanded, will have
the form
\[
y^n + P_1 y^{n-1} + P_2 y^{n-2} + \cdots + P_n = 0,
\]
or, if we determine $t_{n-1}$ so that $S(y) = 0$,
\weq{1}{y^n + P_2 y^{n-2} + \cdots + P_n = 0;}
in it, then,
\[
- 2 P_2 = S(y^2)
\]
is a function of the second degree with respect to the $t$ and with
respect to the $a$. This function is especially important for many
applications and has received from \gesperrt{Sylvester} the name Bezoutiant
of the function $f(x)$, in honour of the French mathematician
\gesperrt{Bezout}, who already in the last century gave the first
correct treatments of elimination.

By the formulas of the last section this function can be
computed comparatively simply. We first introduce, alongside
the variables $t_{n-1}, t_{n-2} \ldots t_0$, a second
system of variables independent of them, $\tau_{n-1}, \tau_{n-2} \ldots \tau_0$, and
set
\weq{2}{
\begin{aligned}
y &= t_{n-1} f_0 + t_{n-2} f_1 + \cdots + t_0 f_{n-1} \\
z &= \tau_{n-1} f_0 + \tau_{n-2} f_1 + \cdots + \tau_0 f_{n-1}.
\end{aligned}
}
Now instead of forming $S(y^2)$, we first compute $S(yz)$,
from which $S(y^2)$ then results when one sets the $\tau$ equal to the $t$.

If one multiplies formula (3) of the preceding section
\weq{3}{y f_s = E_{0,s} f_0 + E_{1,s} f_1 + \cdots + E_{n-1,s} f_{n-1}}
by $\tau_{n-s-1}$ and then sums with respect to $s$ from $0$ to $n-1$,
one finds
\weq{4}{
\begin{aligned}
y z = {} & f_0 \sum_{0,\,n-1}^{s} E_{0,s} \tau_{n-s-1} \\
& + f_1 \sum_{0,\,n-1}^{s} E_{1,s} \tau_{n-s-1} + \cdots + f_{n-1} \sum_{0,\,n-1}^{s} E_{n-1,s} \tau_{n-s-1}.
\end{aligned}
}


% ===== scan leaf 251, printed 226 =====
% !PRINTED 226
\weberpage{226}
Now by \S.\,68, (6),
\[
S f_0 = n a_0, \quad S f_1 = (n-1) a_1, \ldots S f_{n-1} = a_{n-1},
\]
so that one obtains
\[
\begin{aligned}
S(y z) = {} & n a_0 \sum_{0,\,n-1}^{s} E_{0,s} \tau_{n-s-1} \\
& + (n-1) a_1 \sum_{0,\,n-1}^{s} E_{1,s} \tau_{n-s-1} \\
& \cdots\cdots\cdots\cdots \\
& + a_{n-1} \sum_{0,\,n-1}^{s} E_{n-1,s} \tau_{n-s-1}.
\end{aligned}
\]
In this, now, the coefficient of $\tau_{n-1}$ is
\[
n a_0 E_{0,0} + (n-1) a_1 E_{1,0} + \cdots + a_{n-1} E_{n-1,0},
\]
hence by (3)
\[
= a_0 S(y).
\]
Now since $S(z)$ begins with the term $n a_0 \tau_{n-1}$, in the
difference
\weq{5}{S(y z) - \frac{1}{n} S(y) S(z)}
no term will occur that is multiplied by $\tau_{n-1}$. But since this
expression is moreover symmetric with respect to $t$ and $\tau$,
it also does not contain $t_{n-1}$, and in forming it we may
% [sic] The print omits “=0” after the first term; the immediately following paragraph confirms that both t_{n-1} and tau_{n-1} are set equal to zero.
simply assume $t_{n-1}$ and $\tau_{n-1} = 0$.

We now set
\[
S(y) = n a_0 t_{n-1} + \varphi(t),
\]
in which
\weq{6}{\varphi(t) = (n-1) a_1 t_{n-2} + (n-2) a_2 t_{n-3} + \cdots + a_{n-1} t_0.}
Then we have
\weq{7}{
\begin{aligned}
S(y z) - \frac{1}{n} S(y) S(z) = {} & n a_0 \sum_{1,\,n-1}^{s} E_{0,s} \tau_{n-s-1} \\
& + (n-1) a_1 \sum_{1,\,n-1}^{s} E_{1,s} \tau_{n-s-1} \\
& \cdots\cdots\cdots\cdots \\
& + a_{n-1} \sum_{1,\,n-1}^{s} E_{n-1,s} \tau_{n-s-1} - \frac{1}{n} \varphi(t) \varphi(\tau),
\end{aligned}
}
where, however, it is to be noted that $t_{n-1} = 0$ is to be assumed, and hence
\[
E_{s,s} = a_1 t_{n-2} + \cdots + a_s t_{n-s-1}
\]
is to be set, while the remaining $E$ are determined by the formulas (9) of the
preceding section.


% ===== scan leaf 252, printed 227 =====
% !PRINTED 227
\weberpage{227}
The expression on the right-hand side of (7) is a \gesperrt{bilinear}
function of the $t$ and $\tau$; when developed, let it have the form
\[
\sum_{0,\,n-2}^{h} \sum_{0,\,n-2}^{k} B_{h,k} t_h \tau_k,
\]
with the condition $B_{i,k} = B_{k,i}$. The \gesperrt{Bezoutiant} is then
obtained by setting $\tau = t$, thus
\weq{8}{B = \sum_{0,\,n-2}^{h,k} B_{h,k} t_h t_k,}
in which the coefficients $B_{i,k}$ are quadratic functions of the $a$;
and from (7) it follows
\weq{9}{S(y z) - \frac{1}{n} S(y) S(z) = \sum_{0,\,n-2}^{h,k} B_{h,k} t_h \tau_k,}
\weq{10}{S(y^2) = \frac{1}{n} [S(y)]^2 + B.}
By formula (8), \S.\,68, we have
\[
\begin{aligned}
y - \frac{1}{n} S(y) &= t_{n-2} F_0 + t_{n-3} F_1 + \cdots + t_0 F_{n-2} \\
z - \frac{1}{n} S(z) &= \tau_{n-2} F_0 + \tau_{n-3} F_1 + \cdots + \tau_0 F_{n-2},
\end{aligned}
\]
and if we multiply the two together,
\[
y z - \frac{1}{n} [y S(z) + z S(y)] + \frac{1}{n^2} S(y) S(z)
= \sum_{0,\,n-2}^{h,k} \tau_h t_k F_{n-h-2} F_{n-k-2}.
\]
If we sum this formula over all the roots $x$, that is, take
the sum denoted by $S$, it follows that
\[
S(y z) - \frac{1}{n} S(y) S(z) = \sum^{h,k} \tau_h t_k S(F_{n-h-2} F_{n-k-2}),
\]
and comparison with (9) shows
\weq{11}{B_{h,k} = S(F_{n-h-2} F_{n-k-2})}
We denote by $\varDelta$ the determinant of the function $B$,
and thus set
\[
\varDelta = \sum \pm B_{0,0} B_{1,1} \ldots B_{n-2,\,n-2}.
\]
By the multiplication theorem for determinants, and having
regard to the formulas
\[
S F_0 = 0, \quad S F_1 = 0, \ldots S F_{n-2} = 0
\]
%CONT>>


% ===== scan leaf 253, printed 228 =====
%<<CONT
% !PRINTED 228
\weberpage{228}
is, however, if we set
\[
\Phi =
\begin{vmatrix}
1, & F_0(x_1), & F_1(x_1) & \ldots & F_{n-2}(x_1) \\
1, & F_0(x_2), & F_1(x_2) & \ldots & F_{n-2}(x_2) \\
\multicolumn{5}{c}{\cdots\cdots\cdots\cdots\cdots\cdots\cdots\cdots} \\
1, & F_0(x_n), & F_1(x_n) & \ldots & F_{n-2}(x_n)
\end{vmatrix}
\]
by (11)
\weq{12}{\Phi^2 = n \varDelta.}
If one now takes account of the expressions (9) in \S.\,68 for the functions
$F_0, F_1 \ldots F_{n-2}$, then by the theorem that in a determinant
a column multiplied by an arbitrary factor may be added to
another column (\S.\,22, VII), one easily obtains for $\Phi$ the expression
\[
a_0^{n-1}
\begin{vmatrix}
1, & x_1, & x_1^2 & \ldots & x_1^{n-1} \\
1, & x_2, & x_2^2 & \ldots & x_2^{n-1} \\
\multicolumn{5}{c}{\cdots\cdots\cdots\cdots\cdots\cdots} \\
1, & x_n, & x_n^2 & \ldots & x_n^{n-1}
\end{vmatrix},
\]
whose square, by \S.\,46, is the discriminant of $f(x)$. Accordingly,
by (12), we thus have the important theorem:

\gesperrt{The determinant of the Bezoutiant of $f(x)$ is the
$n$th part of the discriminant of $f(x)$.}

The computation of the Bezoutiant presents, by our formulas,
no further difficulty at all, and turns out to be fairly simple.
For $n = 3$ one obtains the result, already to be inferred from
the formulas of \S.\,71,
\weq{13}{B = - \frac{2}{3} H(t_1, t_0).}
For the computation of the Bezoutiant of the biquadratic
form we shall, to illustrate the foregoing, set out the
system of formulas in full, leaving the carrying out of the
slightly longer computation for the form of the fifth degree
to the reader.

For the form of the fourth degree we have, by (7) and \S.\,72, (9)
\[
f(x) = a_0 x^4 + a_1 x^3 + a_2 x^2 + a_3 x + a_4
\]
\[
\begin{array}{rcl|r}
E_{0,1} &=& -a_2 t_2 - a_3 t_1 - a_4 t_0 & 4 a_0 \tau_2 \\
E_{1,1} &=& a_1 t_2 & 3 a_1 \tau_2 \\
E_{2,1} &=& a_0 t_2 + a_1 t_1 & 2 a_2 \tau_2 \\
E_{3,1} &=& a_0 t_1 + a_1 t_0 & a_3 \tau_2 \\
\hline
\end{array}
\]
%CONT>>


% ===== scan leaf 254, printed 229 =====
% !PRINTED 229
\weberpage{229}

%% [running head, not body text -- mirrored 2026-08-07] Page furniture is
%% excluded from this edition by standing policy.
%% \S.\,73. \hfill The Bezoutiant.

\[
\begin{array}{rcl|r}
E_{0,2} &=& - a_3\,t_2 - a_4\,t_1 & 4\,a_0\,\tau_1 \\
E_{1,2} &=& - a_3\,t_1 - a_4\,t_0 & 3\,a_1\,\tau_1 \\
E_{2,2} &=& a_1\,t_2 + a_2\,t_1 & 2\,a_2\,\tau_1 \\
E_{3,2} &=& a_0\,t_2 + a_1\,t_1 + a_2\,t_0 & a_3\,\tau_1 \\
\hline
E_{0,3} &=& - a_4\,t_2 & 4\,a_0\,\tau_0 \\
E_{1,3} &=& - a_4\,t_1 & 3\,a_1\,\tau_0 \\
E_{2,3} &=& - a_4\,t_0 & 2\,a_2\,\tau_0 \\
E_{3,3} &=& a_1\,t_2 + a_2\,t_1 + a_3\,t_0 & a_3\,\tau_0.
\end{array}
\]

One has to multiply these expressions by the factors placed
to the right of them and to add, then further append
\[
\begin{array}{r}
-\dfrac{1}{4}\,\varphi(t)\,\varphi(\tau) \\[1ex]
= -\dfrac{1}{4}\,.(3\,a_1\,t_2 + 2\,a_2\,t_1 + a_3\,t_0)\,(3\,a_1\,\tau_2 + 2\,a_2\,\tau_1 + a_3\,\tau_0)
\end{array}
\]
and to seek out the coefficients of $\tau_h t_k$. In this way
one obtains
\weq{14}{
\begin{array}{rl@{\qquad}rl}
B_{2,2} &= -\dfrac{1}{4}(8\,a_0\,a_2 - 3\,a_1^2), & B_{0,1} &= -\dfrac{1}{2}\;(6\,a_1\,a_4 - a_2\,a_3) \\[1.5ex]
B_{1,1} &= -4\,a_0\,a_4 - 2\,a_1\,a_3 + a_2^2, & B_{1,2} &= -\dfrac{1}{2}\;(6\,a_0\,a_3 - a_1\,a_2) \\[1.5ex]
B_{0,0} &= -\dfrac{1}{4}(8\,a_2\,a_4 - 3\,a_3^2), & B_{0,2} &= -\dfrac{1}{4}(16\,a_0\,a_4 - a_1\,a_3).
\end{array}
}
In the same way one can proceed with a form of the fifth
degree, and one finds for the coefficients of the Bezoutiant the following expressions
\[
\begin{array}{rcl@{\qquad}rcl}
5\,B_{0,0} &=& 4\,a_4^2 - 10\,a_3\,a_5 \\
5\,B_{1,1} &=& 6\,a_3^2 - 10\,a_2\,a_4 - 20\,a_1\,a_5 \\
5\,B_{2,2} &=& 6\,a_2^2 - 10\,a_1\,a_3 - 20\,a_0\,a_4 \\
5\,B_{3,3} &=& 4\,a_1^2 - 10\,a_0\,a_2 \\
\end{array}
\]
\[
\begin{array}{rcl}
5\,B_{0,1} &=& 3\,a_3\,a_4 - 15\,a_2\,a_5 \\
5\,B_{0,2} &=& 2\,a_2\,a_4 - 20\,a_1\,a_5 \\
5\,B_{0,3} &=& a_1\,a_4 - 25\,a_0\,a_5 \\
5\,B_{1,2} &=& 4\,a_2\,a_3 - 15\,a_1\,a_4 - 25\,a_0\,a_5 \\
5\,B_{1,3} &=& 2\,a_1\,a_3 - 20\,a_0\,a_4 \\
5\,B_{2,3} &=& 3\,a_1\,a_2 - 15\,a_0\,a_3.
\end{array}
\]
%CONT>>


% ===== scan leaf 255, printed 230 =====
% !PRINTED 230
\weberpage{230}
\pdfbookmark[2]{§ 74. Transformation of the equation of the fifth degree}{bm82}

%% [running head, not body text -- commented out 2026-08-07] Page furniture (running
%% head, folio, signature) is excluded from this edition by standing policy.
%% Sixth Section. \hfill \S.\,74.

\begin{center}
\S.\,74.

\gesperrt{Transformation of the equation of the fifth degree.}
\end{center}

\indent These developments shall now be applied in order to carry
out the transformation of the equation of the fifth degree, as we
sketched it in \S.\,54.

\indent A system of values $t_0, t_1, t_2, t_3$ can be determined in
infinitely many ways so that the Bezoutiant $B$ vanishes; in general
the solution of a quadratic equation is required for this. We can,
for example, set
\[
t_2 = 0, \quad t_3 = 0
\]
and determine the ratio $t_0 : t_1$ from the quadratic equation
\[
B_{0,0}\,t_0^2 + 2\,B_{0,1}\,t_0\,t_1 + B_{1,1}\,t_1^2 = 0\,.
\]
These values of $t$ then contain the square root
\[
\sqrt{B_{0,1}^2 - B_{0,0}\,B_{1,1}}\,.
\]
Instead of this, however, one can also make other choices, and one
obtains ever different square roots. In the selection of this square
root there lies something arbitrary and undetermined.

\indent Among these various modes of determination we pick out one
and with it transform the given equation of the fifth degree into a
\gesperrt{principal equation}, that is, into one in which the third
and fourth powers of the unknown do not occur; or we assume, following
F. Klein's precedent, that the equation to be transformed is \gesperrt{from the
outset a principal equation}
\[
% The source's principal-form specialization intentionally has no x^4 or x^3 terms.
\text{(1)} \qquad f(x) = a_0\,x^5 + a_3\,x^2 + a_4\,x + a_5 = 0\,.
\]
But we wish to emphasise expressly that, by \S.\,70, (8), through this
preliminary Tschirnhausen transformation the discriminant of the given
equation is changed only by a factor $\Theta^2$, in which $\Theta$
depends \gesperrt{rationally} on the $t$ employed, and hence also on the square
root used for their determination.

\indent Under the assumption (1) the coefficients of the Bezoutiant,
by the formulas of the preceding paragraph, now become the following:
%CONT>>


% ===== scan leaf 256, printed 231 =====
% !PRINTED 231
\weberpage{231}

%% [running head, not body text -- commented out 2026-08-07] Page furniture (running
%% head, folio, signature) is excluded from this edition by standing policy.
%% \S.\,74. \hfill Principal equation of the fifth degree \hfill 231

%<<CONT
\weq{2}{
\begin{array}{rcl@{\qquad}rcl}
5\,B_{0,0} &=& 4\,a_4^2 - 10\,a_3\,a_5, & 5\,B_{0,1} &=& 3\,a_3\,a_4 \\[1ex]
5\,B_{1,1} &=& 6\,a_3^2, & 5\,B_{0,2} &=& 0 \\[1ex]
5\,B_{2,2} &=& -20\,a_0\,a_4, & 5\,B_{0,3} &=& -25\,a_0\,a_5 \\[1ex]
5\,B_{3,3} &=& 0 & 5\,B_{1,2} &=& -25\,a_0\,a_5 \\[1ex]
5\,B_{1,3} &=& -20\,a_0\,a_4, & 5\,B_{2,3} &=& -15\,a_0\,a_3.
\end{array}
}
If one computes the determinant from these, one obtains
without difficulty, by \S.\,73, (12), for the discriminant
of the principal equation of the fifth degree
\[
\text{(3)} \qquad D = a_0^2\bigl(5^5\,a_0^2\,a_5^4 + 2^8\,a_0\,a_4^5 + 2250\,a_0\,a_3^2\,a_4\,a_5^2
- 1600\,a_0\,a_3\,a_4^3\,a_5 + 108\,a_3^5\,a_5 - 27\,a_3^4\,a_4^2\bigr)\,.
\]
Now since $B_{3,3} = 0$ here, we have a system of values
of the $t$ for which the Bezoutiant vanishes, namely
\[
t_0 = 0, \quad t_1 = 0, \quad t_2 = 0\,.
\]
We now take our Tschirnhausen substitution in the following form [\S.\,68, (12)]
\[
\text{(4)} \qquad y = t_3\,F_0 + t_2\,(\alpha_2\,F_1 + \alpha_1\,F_2 + \alpha_0\,F_3)
\]
and we wish to dispose of the $\alpha_0, \alpha_1, \alpha_2$ in such a way that the equation
of the fifth degree for $y$ becomes a principal equation independently of $t_3, t_2$ --- that is, that $S(y^2)$ vanishes identically for all $t_2, t_3$.

\indent Since $S(F_0^2) = B_{3,3}$ vanishes here, this
requirement is satisfied, by \S.\,73, (11), if we determine $\alpha_0, \alpha_1, \alpha_2$ so that
\[
\text{(5)} \qquad \alpha_0^2\,B_{0,0} + \alpha_1^2\,B_{1,1} + \alpha_2^2\,B_{2,2}
+ 2\,\alpha_0\,\alpha_1\,B_{0,1} + 2\,\alpha_0\,\alpha_2\,B_{0,2} + 2\,\alpha_1\,\alpha_2\,B_{1,2} = 0\,,
\]
\[
\text{(6)} \qquad \alpha_0\,B_{0,3} + \alpha_1\,B_{1,3} + \alpha_2\,B_{2,3} = 0\,.
\]
These equations lead, by elimination of one of the
three quantities $\alpha$, to a quadratic equation for the ratio of the other two. In itself it is not necessary
to exclude any particular case, not even the one in which
$B_{0,3}, B_{1,3}, B_{2,3}$ all three vanish; one would then merely have
to take for the $\alpha$ some solution of equation (5).

\indent But in order not to be long-winded, we shall assume
that $B_{0,3}$ is different from zero, because the cases thereby excluded, in which a root of the given equation becomes zero or infinite, really offer no interest
here, since they come back to an equation of lower degree. The treatment, moreover, remains entirely the same if
%CONT>>


% ===== scan leaf 257, printed 232 =====
% !PRINTED 232
\weberpage{232}

%% [running head, not body text -- commented out 2026-08-07] Page furniture (running
%% head, folio, signature) is excluded from this edition by standing policy.
%% 232 \hfill Sixth Section. \hfill \S.\,74.

%<<CONT
we assume that another of the three quantities $B_{0,3}, B_{1,3}, B_{2,3}$
is different from zero.

\indent In order now to treat the equations (5), (6) in a symmetric
way, and in particular to recognise which square root is required
for their solution, we proceed as follows. For brevity we set
\weq{7}{
\begin{array}{rcl}
B_0 &=& \alpha_0\,B_{0,0} + \alpha_1\,B_{0,1} + \alpha_2\,B_{0,2} \\
B_1 &=& \alpha_0\,B_{1,0} + \alpha_1\,B_{1,1} + \alpha_2\,B_{1,2} \\
B_2 &=& \alpha_0\,B_{2,0} + \alpha_1\,B_{2,1} + \alpha_2\,B_{2,2}\,.
\end{array}
}
Then we can represent the equations (5) and (6) as follows:
\weq{8}{
\begin{array}{rcccccl}
\alpha_0\,B_0 & + & \alpha_1\,B_1 & + & \alpha_2\,B_2 & = & 0 \\
\alpha_0\,B_{0,3} & + & \alpha_1\,B_{1,3} & + & \alpha_2\,B_{2,3} & = & 0\,.
\end{array}
}
We multiply the second by an indeterminate factor
$\alpha_3$ and add it to the first, whereby we obtain
\[
\text{(9)} \qquad \alpha_0\,(B_0 + \alpha_3\,B_{0,3}) + \alpha_1\,(B_1 + \alpha_3\,B_{1,3}) + \alpha_2\,(B_2 + \alpha_3\,B_{2,3}) = 0\,,
\]
an equation which is satisfied if we set
\weq{10}{
\begin{array}{rcl}
B_0 + \alpha_3\,B_{0,3} &=& 0 \\
B_1 + \alpha_3\,B_{1,3} &=& \sigma\,\alpha_2 \\
B_2 + \alpha_3\,B_{2,3} &=& -\,\sigma\,\alpha_1\,.
\end{array}
}
From these three equations, in combination with the equation
(6), we have to determine the unknowns $\alpha_0, \alpha_1, \alpha_2, \alpha_3, \sigma$.
We first write the equations out here in full
\weq{11}{
\begin{array}{rcrcrcrcl}
\alpha_0\,B_{0,0} & + & \alpha_1\,B_{0,1} & + & \alpha_2\,B_{0,2} & + & \alpha_3\,B_{0,3} & = & 0 \\
\alpha_0\,B_{1,0} & + & \alpha_1\,B_{1,1} & + & \alpha_2\,(B_{1,2} - \sigma) & + & \alpha_3\,B_{1,3} & = & 0 \\
\alpha_0\,B_{2,0} & + & \alpha_1\,(B_{2,1} + \sigma) & + & \alpha_2\,B_{2,2} & + & \alpha_3\,B_{2,3} & = & 0 \\
\alpha_0\,B_{3,0} & + & \alpha_1\,B_{3,1} & + & \alpha_2\,B_{3,2} & & & = & 0\,,
\end{array}
}
and if we eliminate $\alpha_0, \alpha_1, \alpha_2, \alpha_3$ from these, there results
a quadratic equation for $\sigma$, after whose solution the
ratios $\alpha_0 : \alpha_1 : \alpha_2 : \alpha_3$ can be determined from linear equations.

\indent The quadratic equation for $\sigma$, however, reads in determinant form
\[
\begin{vmatrix}
B_{0,0}, & B_{0,1}, & B_{0,2}, & B_{0,3} \\
B_{1,0}, & B_{1,1}, & B_{1,2} - \sigma, & B_{1,3} \\
B_{2,0}, & B_{2,1} + \sigma, & B_{2,2}, & B_{2,3} \\
B_{3,0}, & B_{3,1}, & B_{3,2}, & 0
\end{vmatrix} = 0\,.
\]


% ===== scan leaf 258, printed 233 =====
% !PRINTED 233
\weberpage{233}
\pdfbookmark[2]{§ 75. Normal form of the equation of the fifth degree}{bm83}

Since the left side is not changed by the permutation of $\sigma$ with $-\sigma$, it is a pure quadratic equation, and, having regard to \S.\,73, (12), it gives
\[
B_{0,3}^2\,\sigma^2 - 5\,D = 0,
\]
where $D$ is the discriminant of the given equation. By (2) we therefore obtain
\weq{12}{5\,a_0\,a_5\,\sigma = \sqrt{5\,D},}
in which the sign is arbitrary.

\begin{center}
\S.\ 75.

\gesperrt{Normal form of the equation of the fifth degree.}
\end{center}

As already remarked earlier, the principal aim of these considerations is to establish a normal form of the equation of the fifth degree which depends on only \gesperrt{one} undetermined coefficient, \gesperrt{one parameter}. Such a normal form is the \gesperrt{Bring-Jerrard} form. To obtain it, by (4) of the preceding paragraph we have to set
\weq{1}{y = t_3\,F_0 + t_2\,(\alpha_2\,F_1 + \alpha_1\,F_2 + \alpha_0\,F_3)}
so that identically
\[
S(y) = 0, \quad S(y^2) = 0
\]
holds, and then the ratio $t_3 : t_2$ is to be formed from the cubic equation
\[
S(y^3) = 0
\]
This cubic equation can actually be formed, even though its expression becomes long. The formulae needed for this have been computed by \gesperrt{Cayley}\textsuperscript{1)}.

It is the result of a remarkable investigation of \gesperrt{Gordan} that one can obtain another normal form, the \gesperrt{Brioschi} form, without any new irrationality, and to conclude these considerations on the Tschirnhausen transformation we shall now derive this result\textsuperscript{2)}.

We hold fast to the assumptions of the preceding paragraph and, for the moment, set as an abbreviation
%CONT>>

\bigskip
\noindent\rule{4cm}{0.4pt}

\noindent\textsuperscript{1)} \gesperrt{Cayley}, on Tschirnhausen's Transformation, Phil. Trans. 1861, Mathematical Papers, Tom. IV, Nr. 275.

\noindent\textsuperscript{2)} \gesperrt{Gordan}, Mathematische Annalen, Bd. 28, 1886.


% ===== scan leaf 259, printed 234 =====
% !PRINTED 234
\weberpage{234}
%<<CONT
\weq{2}{u = F_0, \quad v = \alpha_2\,F_1 + \alpha_1\,F_2 + \alpha_0\,F_3,}
in which $\alpha_0, \alpha_1, \alpha_2$ are to have the values determined in the preceding paragraph.

By the formula \S.\,73, (4) the three functions
\[
u^2, \quad u\,v, \quad v^2
\]
can be represented linearly and homogeneously in terms of $f_0, f_1, f_2, f_3, f_4$, and since
\weq{3}{S(u^2) = 0, \quad S(u\,v) = 0, \quad S(v^2) = 0}
holds, these expressions also become linear and homogeneous in $F_0, F_1, F_2, F_3$.

The calculation is easy to carry out by the formulas of \S\S.\,72, 73, but it will not be pursued further here, since we are concerned only with setting out the fundamental idea. We wish only to remark that the coefficients in the expressions for $u^2$, $u\,v$, $v^2$ are linear in the coefficients of the original equation of the fifth degree and quadratic in the $\alpha_0, \alpha_1, \alpha_2$. If now, with the aid of (2), we eliminate the functions $F_0, F_1, F_2, F_3$ from these expressions, we obtain a relation of the form
\weq{4}{p\,u^2 + 2\,q\,u\,v + r\,v^2 = a\,u + b\,v,}
% [sic] The source prints a colon between a and b in this coefficient list; preserved diplomatically.
in which $p, q, r, a: b$ depend rationally on the coefficients $a_i$ and $\alpha_i$, and in any case do not all vanish at once.

If for the moment we set
\[
p\,u^2 + 2\,q\,u\,v + r\,v^2 = \varphi(u,v),
\]
then by \S.\,57
\[
4\,\varphi(u',v')\,\varphi(u,v) - [u'\,\varphi'(u) + v'\,\varphi'(v)]^2
\]
is the square of a linear function of $u, v$, so that $\varphi(u,v)$ is thereby decomposed into the sum of two squares. If we choose $u' = b$, $v' = -a$, it follows that
\[
\varphi(b,-a)\,\varphi(u,v) = [b\,(p\,u + q\,v) - a\,(q\,u + r\,v)]^2 + (p\,r - q^2)\,(a\,u + b\,v)^2,
\]
or, if for abbreviation one sets
\weq{5}{\begin{aligned}
p\,b^2 - 2\,q\,a\,b + r\,a^2 &= m,\\
b\,p - a\,q &= a',\\
b\,q - a\,r &= b',\\
q^2 - p\,r &= c
\end{aligned}}
then
\weq{6}{m\,(p\,u^2 + 2\,q\,u\,v + r\,v^2) = (a'\,u + b'\,v)^2 - c\,(a\,u + b\,v)^2,}
%CONT>>


% ===== scan leaf 260, printed 235 =====
% !PRINTED 235
\weberpage{235}
%<<CONT
and by (4) this equation can also be represented thus:
\weq{7}{m\,(a\,u + b\,v) = (a'\,u + b'\,v)^2 - c\,(a\,u + b\,v)^2.}
Now set
\weq{8}{y = \frac{a'\,u + b'\,v}{a\,u + b\,v},}
and form the equation of the fifth degree whose roots are the five values of $y$:
\weq{9}{y^5 + c_1\,y^4 + c_2\,y^3 + c_3\,y^2 + c_4\,y + c_5 = 0;}
in it the coefficients $c$ can be determined more closely in the following way.

From (8) we derive
\[
y + \sqrt{c} = \frac{a'\,u + b'\,v + \sqrt{c}\,(a\,u + b\,v)}{a\,u + b\,v}
\]
and from this, by (7),
\[
y + \sqrt{c} = \frac{m}{a'\,u + b'\,v - \sqrt{c}\,(a\,u + b\,v)},
\]
\[
\frac{m}{y + \sqrt{c}} = a'\,u + b'\,v - \sqrt{c}\,(a\,u + b\,v).
\]
But from this it follows, by (3), that
\[
S\left[\frac{1}{y + \sqrt{c}}\right] = 0, \quad S\left[\frac{1}{(y + \sqrt{c})^2}\right] = 0.
\]
If, then, we derive from (9) the equation whose roots are the values
\[
\xi = \frac{1}{y + \sqrt{c}}
\]
that is, if we set
\[
y = \frac{1 - \xi\,\sqrt{c}}{\xi}
\]
an equation for $\xi$ must arise in which the coefficients of $\xi^4$ and $\xi^3$ vanish, and indeed whichever sign we give to the square root $\sqrt{c}$. In this way one obtains four equations between the coefficients $c_\nu$ of equation (9).

The equation for $\xi$ becomes, namely,
\[
(1 - \xi\,\sqrt{c})^5 + c_1\,\xi\,(1 - \xi\,\sqrt{c})^4 + c_2\,\xi^2\,(1 - \xi\,\sqrt{c})^3 + c_3\,\xi^3\,(1 - \xi\,\sqrt{c})^2 + c_4\,\xi^4\,(1 - \xi\,\sqrt{c}) + c_5\,\xi^5 = 0.
\]
If one here sets the coefficients of $\xi^4$ and $\xi^3$ equal to 0, it follows that
%CONT>>


% ===== scan leaf 261, printed 236 =====
% !PRINTED 236
\weberpage{236}
%<<CONT
\[
\begin{aligned}
5\,\sqrt{c}^{\,4} - 4\,c_1\,\sqrt{c}^{\,3} + 3\,c_2\,\sqrt{c}^{\,2} - 2\,c_3\,\sqrt{c} + c_4 &= 0,\\
-10\,\sqrt{c}^{\,3} + 6\,c_1\,\sqrt{c}^{\,2} - 3\,c_2\,\sqrt{c} + c_3 &= 0,
\end{aligned}
\]
and on account of the double sign of $\sqrt{c}$ these equations split into the four
\[
\begin{aligned}
5\,c^2 + 3\,c_2\,c + c_4 &= 0,\\
4\,c_1\,c + 2\,c_3 &= 0,\\
10\,c + 3\,c_2 &= 0,\\
6\,c_1\,c + c_3 &= 0.
\end{aligned}
\]
From the second and fourth of these equations it follows that
\[
c_1 = 0, \quad c_3 = 0,
\]
and then from the third and first
\[
c_2 = -\frac{10}{3}\,c, \quad c_4 = -5\,c^2,
\]
so that the equation for $y$ thus takes the form
\weq{10}{y^5 - \frac{10}{3}\,c\,y^3 + 5\,c^2\,y + c_5 = 0.}
This equation still depends on the two parameters $c_5$ and $c$; one can reduce it to an equation with one parameter by the substitution
\weq{11}{y = \sqrt{\frac{c}{3}}\,z, \quad \gamma = c_5\,\sqrt{\frac{3^5}{c^5}},}
whereby it takes the simple and elegant form
\weq{12}{z^5 - 10\,z^3 + 45\,z + \gamma = 0.}
This is the \gesperrt{Brioschi} normal form. The substitution (11), however, suffers from the drawback that it still contains a square root; whereas in the formulas (8) and (10) only the two square roots discussed in the previous section occur. But, following a remark of \gesperrt{Klein}, one can also pass from (10) by rational means to a normal form that contains only one parameter.

If, for example, one sets
\weq{13}{y = \frac{3\,c_5}{c^2}\,\frac{1}{z}, \quad \gamma = \frac{9\,c_5^2}{c^5},}
then from (10) one obtains the equation
\weq{14}{z^5 + 15\,z^4 - 10\,\gamma\,z^2 + 3\,\gamma^2 = 0.}


% ===== scan leaf 262, printed 237 =====
% !PRINTED 237
\weberpage{237}

In this transformation the assumption has been tacitly made that the quantities denoted by $m$ and $c$ in (5) are not zero.

If $m = 0$, then in formula (4) $p\,u^2 + 2\,q\,u\,v + r\,v^2$ is divisible by $a\,u + b\,v$; hence, at least for some of the roots, the equation
\[
a\,u + b\,v = 0,
\]
that is, an equation of degree no higher than the fourth, must hold.

If $c = 0$, so that $p\,u^2 + 2\,q\,u\,v + r\,v^2$ is the square of a linear function, then it follows from (7) that, if we set
\[
y = a'\,u + b'\,v
\]
then
\[
S(y) = 0, \quad S(y^2) = 0, \quad S(y^3) = 0
\]
and hence that the \gesperrt{Bring-Jerrard} form can be produced in a rational way.

Finally, we may summarise the results of these considerations as follows:

\gesperrt{The principal equation of the fifth degree can be transformed, by a transformation which contains as its only irrationality the square root of the discriminant, into a normal form with one parameter.}

\begin{center}
\rule{0.25\linewidth}{0.4pt}
\end{center}


% ===== scan leaf 263, printed 238 =====
\weberpage{238}% page number not printed on this page
% !PRINTED none
% Blank verso after printed 237: no direct-printed text, running head, folio, or rule.
% Pale show-through of the recto's line and formula field is visible across much of
% the leaf; it is not page content.


% ===== scan leaf 264, printed 239 =====
\weberpage{239}% page number not printed on this page
% !PRINTED none
% Part-title recto opening Book Two: no running page number. The capitals are
% letterspaced. There is no leading rule: a short rule separates the two title
% lines, and a longer rule follows the second line.
% [leaf swap mirrored 2026-08-07] The direct-printed title belongs on i0264;
% i0265 is its blank, ghosted verso.
\begin{center}
\gesperrt{BOOK TWO.}

\rule{0.15\linewidth}{0.4pt}

\bigskip
\gesperrt{THE ROOTS.}
\bigskip

\rule{0.25\linewidth}{0.4pt}
\end{center}


% ===== scan leaf 265, printed 240 =====
\weberpage{240}% page number not printed on this page
% !PRINTED none
% Blank verso of the Book Two part-title leaf. No direct-printed text, rule,
% running head, or folio. The two title lines are visible only as faint,
% horizontally reversed show-through from the recto at i0264.
% [leaf swap mirrored 2026-08-07] The English title block was moved to i0264.


% ===== scan leaf 266, printed 241 =====
\weberpage{241}% page number not printed on this page
\pdfbookmark[0]{Second Book. The Roots.}{bm84}
\pdfbookmark[1]{Seventh Section. Reality of the Roots.}{bm85}
\pdfbookmark[2]{§ 76. Generalities on the reality of roots of equations and on discriminants}{bm86}
% !PRINTED none
% section-opening page: "Siebenter Abschnitt. Realität der Wurzeln." heading block with §.76;
% no printed folio number was visible on this page in the scan (section-opening layout
% replaces the usual running head with the centred section title). Footer furniture
% "Weber, Algebra. I." and signature "16" are excluded under standing policy.

\begin{center}
\gesperrt{Seventh Section.}

\gesperrt{Reality of the Roots.}

\rule{3cm}{0.4pt}

\S.\,76.

\gesperrt{Generalities on the reality of roots of equations and on the discriminant.}
\end{center}

In this section we shall occupy ourselves with the question of how many roots of an algebraic equation are real. The coefficients of the equation are thereby assumed to be \gesperrt{real numbers}. We shall call such equations, for short, \gesperrt{real equations}, and begin with some general considerations.

The real equation of degree $n$
\[
f(x) = 0
\]
has, as we have seen in the third section, $n$ roots, which are either real or imaginary, i.\,e. of the form $\xi + i\eta$, with $\xi$, $\eta$ real. But this count of $n$ roots is obtained in general only if, in certain circumstances, we count a root multiply, namely $(m+1)$-fold when, together with $f(x)$, the first $m$ derivatives of $f(x)$ vanish. Now since a complex expression of the form $X + iY$ is equal to zero only when the two real constituents $X$ and $Y$ vanish separately, that is, when together with $X+iY$ also $X-iY$ vanishes, and since further, with real coefficients, if $f(\xi+i\eta) = X+iY$, one obtains $f(\xi-i\eta) = X-iY$, it follows from $f(\xi+i\eta)=0$ that at the same time $f(\xi-i\eta)=0$ must hold, and hence that to every imaginary root $\xi+i\eta$ there belongs a second imaginary root distinct from it, $\xi-i\eta$, i.\,e. the conjugate imaginary root
%CONT>>


% ===== scan leaf 267, printed 242 =====
% !PRINTED 242
\weberpage{242}
%<<CONT
belongs. Since, together with the derivatives $f'(\xi+i\eta), f''(\xi+i\eta)\ldots$, the conjugate quantities $f'(\xi-i\eta), f''(\xi-i\eta)\ldots$ vanish at the same time, it follows that \gesperrt{conjugate imaginary roots have the same degree of multiplicity.}

From this it follows that the imaginary roots always occur in even number, and that an equation of odd degree must always have at least one real root.

Our next aim will be to determine the number of real roots, without solving the equation, directly from the values of \gesperrt{certain rational functions of the coefficients of the equation}.

A very important part is played here by the discriminant, whose significance for our question we must first consider.

In \S.\,46 we defined the discriminant as the product of the squares of all the root differences
\[
(x_1 - x_2)^2,
\]
further multiplied by $a_0^{2n-2}$, where $a_0$ is the coefficient of the highest power of the unknown in the given equation, and we showed there how the discriminant can be computed as a rational function of the coefficients. The factor $a_0^{2n-2}$ is always positive, and could here also be assumed equal to 1 without essential restriction of generality. As before, we shall use the sign $D$ for the discriminant.

\gesperrt{The discriminant vanishes when, and only when, two equal roots occur among the roots.}

If we assume that $f(x)$ has been freed of multiple factors by separating off the greatest common divisor of $f(x)$ and $f'(x)$, which is done by rational computation, then $D$ will not vanish.

If $x_1$ and $x_2$ are real, then $(x_1 - x_2)^2$ is positive. If $x_1$ is real and $x_2$ imaginary, then there is a root $x_2'$ conjugate to $x_2$, and the product
\[
(x_1 - x_2)(x_1 - x_2')
\]
is positive, being the product of two conjugate imaginary quantities, and hence so is its square.


% ===== scan leaf 268, printed 243 =====
% !PRINTED 243
\weberpage{243}
\pdfbookmark[2]{§ 77. Discussion of the quadratic and cubic equations}{bm87}
If $x_1$ and $x_2$ are both imaginary, but not conjugate, then the product
\[
(x_1 - x_2)(x_1' - x_2')
\]
and its square are positive.

But if, finally, $x_1$ and $x_2$ are conjugate imaginary, then their difference is \gesperrt{purely imaginary} and its square negative.

Thus in the product by which we have defined the discriminant there occur exactly as many negative factors as there are pairs of conjugate imaginary roots, and from this we conclude the important fundamental theorem:

\gesperrt{If $f(x) = 0$ is a real equation without multiple roots, then the discriminant is positive or negative according as the number of pairs of conjugate imaginary roots is even or odd.}

The case in which only real roots are present belongs to those in which the discriminant is positive.

For the cases of the quadratic and cubic equations, in which only one pair of imaginary roots can occur, this principal theorem already completes the determination among the various cases possible with regard to the reality of the roots.

\begin{center}
\S.\,77.

\gesperrt{Discussion of the quadratic and cubic equation.}
\end{center}

For the quadratic equation
\[
a_0 x^2 + a_1 x + a_2 = 0
\]
we have (\S.\,46)
\[
\begin{array}{l}
D = a_1^2 - 4a_0a_2 > 0 \quad \text{real roots.}\\
\phantom{D = a_1^2 - 4a_0a_2 } < 0 \quad \text{imaginary roots.}
\end{array}
\]
For the cubic equation
\[
a_0 x^3 + a_1 x^2 + a_2 x + a_3 = 0
\]
we have
\[
D = a_1^2a_2^2 + 18a_0a_1a_2a_3 - 4a_0a_2^3 - 4a_1^3a_3 - 27a_0^2a_3^2,
\]
\[
\begin{array}{l}
D > 0 \quad \text{three real roots,}\\
D < 0 \quad \text{one real, two imaginary roots.}
\end{array}
\]


% ===== scan leaf 269, printed 244 =====
% !PRINTED 244
\weberpage{244}

The case $D=0$ likewise gives no occasion here for further distinctions; for in this case the two roots of the quadratic equation are equal to one another and real, and of the three roots of the cubic equation two are equal to one another and all three are real; for an imaginary root cannot occur doubly here, since otherwise its conjugate would also occur doubly, and just as little can the single root be imaginary, since otherwise a second one would have to be present.

For the cubic equation it only remains to seek the condition under which all three roots are equal to one another.

In this case the left-hand side of the cubic equation must be a complete cube, so that
\[
a_0 x^3 + a_1 x^2 + a_2 x + a_3 = a_0 (x - \alpha)^3.
\]
Comparison of the two sides of this identical equation yields
\[
a_1 = -3a_0\alpha, \quad a_2 = 3a_0\alpha^2, \quad a_3 = -a_0\alpha^3,
\]
from which, by elimination of $\alpha$, one obtains
\weq{1}{\begin{aligned} a_1^2 - 3a_0a_2 &= 0,\\ a_1a_2 - 9a_0a_3 &= 0. \end{aligned}}
From this there results as a consequence (by multiplying the first of these equations by $a_2$, the second by $a_1$, and subtracting)
\[
a_2^2 - 3a_1a_3 = 0;
\]
and when these conditions are satisfied, there follows from them conversely
\[
f(x) = a_0\left(x + \frac{a_1}{3a_0}\right)^3,
\]
i.\,e. the equality of all three roots.

For the cubic equation
\[
x^3 + a_1 x^2 + a_2 x + a_3 = 0
\]
we can decide completely not only about the reality but also about the signs of the roots.

Namely, if we denote by $\alpha$, $\beta$, $\gamma$ the three roots, then
\weq{2}{a_1 = -(\alpha+\beta+\gamma), \quad a_2 = \alpha\beta+\alpha\gamma+\beta\gamma, \quad a_3 = -\alpha\beta\gamma.}

Now if $D<0$, so that only one root, say $\alpha$, is real, then $\beta\gamma$ is positive and $\alpha$ will be negative or positive according as $a_3$ is positive or negative. But if $D$ is positive, so that all three roots are real, then, when $\alpha$, $\beta$, $\gamma$ are positive, by (2)
%CONT>>


% ===== scan leaf 270, printed 245 =====
% !PRINTED 245
\weberpage{245}

\weq{3}{a_1 < 0, \quad a_2 > 0, \quad a_3 < 0;}
these conditions are necessary in order that all three roots be
positive. But they are also sufficient; for if one or three roots
are negative, then $a_3 > 0$. If, however, two roots, say
$\beta, \gamma$, are negative, then either
\[
a_1 \geqq 0 \quad \text{or} \quad \alpha > -(\beta + \gamma).
\]
But in the latter case it follows that
\[
a_2 < \beta\gamma - (\beta + \gamma)^2 = -\beta^2 - \beta\gamma - \gamma^2,
\]
so $a_2$ is negative.

Finally, if one root is equal to 0, then $a_3$ must necessarily
vanish.

By replacing $x$ with $-x$, one concludes that for three
negative roots the necessary and sufficient condition is
\weq{4}{a_1 > 0, \quad a_2 > 0, \quad a_3 > 0.}
We therefore obtain, under the assumption of a positive
discriminant, the following table:
\[
\begin{array}{ccc|c}
a_1 & a_2 & a_3 & \\ \hline
- & + & - & 3\\[2pt]
- & - & - & \\
+ & + & - & \raisebox{1.2em}{$\big\}1$}\\
+ & - & - & \\ \hline
+ & + & + & 0\\[2pt]
- & - & + & \\
- & + & + & \raisebox{1.2em}{$\big\}2$}\\
+ & - & + &
\end{array}
\]
where in the last column the number of positive roots is given.
We can express the result of this consideration as follows:

\gesperrt{For a positive discriminant, the number of
positive roots of the cubic equation is equal to the number of
sign changes in the sequence}
\[
1, \; a_1, \; a_2, \; a_3,
\]
where by a sign change we understand the succession of a
positive and a negative or of a negative and a positive
quantity.


% ===== scan leaf 271, printed 246 =====
% !PRINTED 246
\weberpage{246}
\pdfbookmark[2]{§ 78. Discussion of the biquadratic equation}{bm88}

If one of the two quantities $a_1$, $a_2$ vanishes, then, since in that case not all the roots can be of the same sign, we have one or two positive roots.

If $a_3 = 0$, the discriminant reduces to
\[
a_2^2(a_1^2 - 4a_2),
\]
and if this is positive, the cubic equation has one vanishing root and two further real roots. These are positive if
\[
a_1 < 0, \quad a_2 > 0,
\]
negative if
\[
a_1 > 0, \quad a_2 > 0,
\]
and one of them is positive and one negative if
\[
a_1 \gtreqless 0, \quad a_2 < 0
\]
holds.

\begin{center}
\S.\,78.

\gesperrt{Discussion of the biquadratic equation.}
\end{center}

For equations of the fourth and fifth degree there exist either no imaginary roots, or two, or four. If the discriminant is negative, there are two imaginary roots. With a positive discriminant, either four imaginary roots or none may be present. To distinguish these two cases will be our task in what follows. But first we turn to an elementary consideration of the biquadratic equation, which for the sake of simplicity we shall assume in the abbreviated form
\weq{1}{x^4 + ax^2 + bx + c = 0}
from which we can easily return to the general form.

If we denote the roots by $\alpha$, $\beta$, $\gamma$, $\delta$ and set, as in \S.\,47,
\weq{2}{\begin{aligned} \alpha - \beta &= v+w, & \gamma - \delta &= v-w,\\ \alpha - \gamma &= w+u, & \delta - \beta &= w-u,\\ \alpha - \delta &= u+v, & \beta - \gamma &= u-v, \end{aligned}}
then $u^2$, $v^2$, $w^2$ are the roots of the cubic resolvent
\weq{3}{y^3 + 2ay^2 + (a^2-4c)y - b^2 = 0,}
and the discriminant $D$ of this cubic equation, which by
%CONT>>


% ===== scan leaf 272, printed 247 =====
% !PRINTED 247
\weberpage{247}
%<<CONT
\weq{4}{27D = 4(a^2+12c)^3 - (2a^3-72ac+27b^2)^2}
is determined, is at the same time the discriminant of the biquadratic equation (1), and the signs of the quantities $u$, $v$, $w$ are restricted by the condition
\weq{5}{uvw = -b}
[\S.\,37, (3)].

If the discriminant is negative, then the equation (3), just like (1), has two conjugate imaginary roots.

If $D$ is positive and all four roots $\alpha$, $\beta$, $\gamma$, $\delta$ are real, then $u$, $v$, $w$ also become real, and hence their squares, i.\,e. the roots of (3), are \gesperrt{positive}.

If all four roots are imaginary, say $\alpha$ conjugate to $\beta$ and $\gamma$ to $\delta$, then it follows from (2) that $v$ and $w$ are purely imaginary and $u$ is real; hence in this case the equation (3) has one positive and two negative roots.

In both cases, however, if $b$ vanishes, one of the quantities $u$, $v$, $w$ may also be equal to zero.

Having regard to the results of the last section on the cubic equation, we thus arrive at the following result:

\gesperrt{The necessary and sufficient condition for the existence of four distinct real roots is}
\weq{6}{D > 0, \quad a < 0, \quad a^2 - 4c > 0.}

In all other cases in which $D$ is positive, the equation has four imaginary roots.

We can also carry the discussion through completely for $D=0$, that is, in the case of the equality of two roots.

If we assume $\gamma = \delta$, then, since we have made the assumption $\alpha+\beta+\gamma+\delta = 0$, it follows from (2) that
\[
2v = 2w = \alpha - \beta
\]
\[
2u = \alpha+\beta-2\gamma = 2(\alpha+\beta)
\]
and from this, by (3),
\weq{7}{\begin{aligned} -4a &= 2(\alpha+\beta)^2 + (\alpha-\beta)^2\\ 16(a^2-4c) &= (\alpha-\beta)^2[8(\alpha+\beta)^2+(\alpha-\beta)^2]. \end{aligned}}
To dispose first of the case in which also $\alpha=\beta$, we obtain from (7) for this the condition $a^2-4c=0$, and the first equation (7) shows that $a$ is negative if $\alpha$ and $\gamma$
%CONT>>


% ===== scan leaf 273, printed 248 =====
% !PRINTED 248
\weberpage{248}
%<<CONT
real, but positive when $\alpha$ and $\gamma$ are conjugate (and then, because $\alpha+\gamma=0$, purely imaginary). From this it follows:

The biquadratic equation has two pairs of equal real roots when
\weq{8}{D = 0, \quad a < 0, \quad a^2-4c = 0}
and two pairs of equal imaginary roots when
\weq{9}{D = 0, \quad a > 0, \quad a^2-4c = 0.}

But if $\alpha$ is different from $\beta$, then $\gamma$ must be real, and (7) shows that when $\alpha$ and $\beta$ are real, $a$ is negative and $a^2-4c$ is positive; that, on the other hand, when $\alpha$ and $\beta$ are conjugate imaginary, either $a$ must be positive or $a^2-4c$ negative, since in this case $(\alpha-\beta)^2$ is negative, and if $2(\alpha+\beta)^2+(\alpha-\beta)^2$ is positive, then certainly $8(\alpha+\beta)^2+(\alpha-\beta)^2$ must also be positive.

We can therefore supplement (6) by saying:

\gesperrt{The necessary and sufficient condition for the existence of four real roots, of which two (but not more) may also be equal to one another, is}
\weq{10}{D \geqq 0, \quad a < 0, \quad a^2-4c > 0.}

\gesperrt{The necessary and sufficient condition for three equal roots,} which are necessarily all real, is by (2)
\[
u^2 = v^2 = w^2,
\]
hence the cubic resolvent (3) must have three equal roots, and for this, by \S.\,77 (1), the necessary and sufficient conditions are
\weq{11}{\begin{aligned} a^2+12c &= 0\\ 2a^3-8ac+9b^2 &= 0. \end{aligned}}
By \S.\,47, (15), (16) the two invariants $A$ and $B$ of the biquadratic equation are
\[
A = a^2+12c, \quad B = 2a^3-72ac+27b^2.
\]
Hence we can also write the equations (11) thus
\[
A = 0, \quad B+4aA = 0,
\]
so that (11) is equivalent to
\[
A = 0, \quad B = 0.
\]
By elimination of $c$ one can also derive from (11) the equation:
\weq{12}{8a^3+27b^2 = 0,}


% ===== scan leaf 274, printed 249 =====
% !PRINTED 249
\weberpage{249}

Finally, there remains the condition for the equality of all four roots
\weq{13}{a = 0, \quad b = 0, \quad c = 0.}

These relations can be made very vivid by a geometric interpretation, and although geometric considerations do not properly lie within our plan, we shall nevertheless not neglect to point them out to the reader from time to time.

If one interprets $a$, $b$, $c$ as rectangular coordinates in space, then every point of space is to be regarded as the carrier of a certain biquadratic equation of the form (1), $f(x) = 0$; all equations that have a given number $x$ as a root are represented by the points of a plane $[f(x) = 0]$. The equation $D = 0$ is the equation of a curved surface (of the fifth degree), the discriminant surface, which is generated by the lines of intersection of the planes $f(x) = 0$, $f'(x) = 0$, and is therefore a developable surface. It is the envelope of all the planes $f(x) = 0$.

The surface has an edge of regression consisting of two branches, which is represented by the equations (11), (12), and a double line, which is determined by the equations $b = 0$, $a^2 - 4c = 0$ and therefore has the shape of a parabola. On the side of negative $a$ the surface crosses itself in this parabola. On the side of positive $a$ the parabola continues as an isolated line.

The discriminant surface divides the whole of space into three compartments, of which two are connected only along the double parabola, and these compartments contain the points to which no, two, and four real roots correspond. If for the moment we call these three compartments $[0]$, $[2]$, $[4]$, then $[0]$ borders $[4]$ only along the double parabola, while $[0]$ as well as $[4]$ border $[2]$ along the parts of the surface.

% [sic] The source defines [4,2] but here prints the second boundary pair as [0],[4]; preserved diplomatically.
We shall denote by $[0, 2]$ and $[4, 2]$ the boundary surfaces of $[0]$, $[2]$ and of $[0]$, $[4]$. The isolated part of the parabola continues into the interior of $[0]$. The edges of regression lie on the part of the surface that separates $[4]$ from $[2]$.

At the origin of coordinates the three parts of space and their boundary surfaces and boundary lines meet. The points of the surface parts $[0, 2]$ represent equations with two equal and
%CONT>>


% ===== scan leaf 275, printed 250 =====
% !PRINTED 250
\weberpage{250}
%<<CONT
two imaginary roots, and the points of $[2, 4]$ equations with two equal roots and two real roots distinct from them.

On the double parabola two equal roots occur twice, and indeed on the part where $[0]$ borders on $[4]$ they are real, on the isolated part imaginary.

% Source figure: Weber_Algebra_Bd1_IA.pdf, PDF page 276, printed p.250, Fig.5; exact caption-inclusive crop.
% Former placeholder description retained: perspective line drawing of the discriminant surface, showing two cuspidal-edge branches meeting at a shaded curved vault-like shell; dotted lines indicate hidden edges and the coordinate framework of the model.
\weberfigureimage{b1_i0275_fig5_source_20260831.png}{0.56}

The points of the cuspidal edges represent equations with three equal roots, and the origin of coordinates the equation with four equal roots.

It was \gesperrt{Kronecker} who first drew attention to the discriminant surface and its significance for the discussion of the biquadratic equation. A model of the surface has been constructed by \gesperrt{Kerschensteiner}.\footnote{\gesperrt{Kronecker}, Monatsbericht der Berliner Akademie of 14 February 1878. A description of the model is to be found in the catalogue of mathematical models edited by Dyck, Munich 1892, from which Fig. 5 above is taken.}

The analytic expression for the conditions of reality of the roots can be brought into a form in which only the covariants of the biquadratic form occur.\footnote{\gesperrt{Clebsch}, ``Theorie der binären Formen'' \S.\,47. \gesperrt{Faà di Bruno}, ``Theorie der binären Formen'', German translation by Walter (Leipzig 1881), \S.\,20, where an essential addition by \gesperrt{Nöther} is to be found.}

The starting point for this is furnished by the expressions for the functions $\psi_1$, $\psi_2$, $\psi_3$ given in \S.\,66, (6). If we, in
%CONT>>


% ===== scan leaf 276, printed 251 =====
% !PRINTED 251
\weberpage{251}
%<<CONT
those expressions we interchange $\alpha$ with $\beta$, then $\psi_1$, $\psi_2$, $\psi_3$ pass into $\psi_1$, $-\psi_3$, $-\psi_2$, and if we simultaneously interchange $\alpha$ with $\beta$ and $\gamma$ with $\delta$, then $\psi_1$, $\psi_2$, $\psi_3$ pass into $\psi_1$, $-\psi_2$, $-\psi_3$. From this the following results:

If one sets an arbitrary real value for the variable $x$, and if the roots $\alpha$, $\beta$, $\gamma$, $\delta$ of the biquadratic form are all real, then $\psi_1$, $\psi_2$, $\psi_3$ are also real.

If $\alpha$, $\beta$ are conjugate imaginary and $\gamma$, $\delta$ real, then the interchange of $\alpha$ and $\beta$ is equivalent to the interchange of $i$ with $-i$. In this case, therefore, $\psi_1$ is real, and $\psi_2$, $\psi_3$ are conjugate imaginary.

If, finally, $\alpha$, $\beta$ and $\gamma$, $\delta$ are two pairs of conjugate imaginary roots, then $\alpha$ is interchanged with $\beta$ and $\gamma$ with $\delta$ when $i$ passes into $-i$; hence in this case $\psi_1$ is real, and $\psi_2$, $\psi_3$ are \gesperrt{purely imaginary}.

If four roots are real, then $\psi_1^2$, $\psi_2^2$, $\psi_3^2$ become real and positive.

If two roots are real, then $\psi_1^2$ is real and positive, and $\psi_2^2$, $\psi_3^2$ are conjugate imaginary.

If four roots are imaginary, then $\psi_1^2$ is positive, and $\psi_2^2$, $\psi_3^2$ are real and negative.

Now in \S.\,66, (7) we have set up the cubic equation whose roots are $\psi_1^2$, $\psi_2^2$, $\psi_3^2$, and have also formed its discriminant. In the previous section we became acquainted with a criterion for the number of positive roots of a cubic equation, from which the following result is obtained:

If $D < 0$, the biquadratic equation has two real and two imaginary roots.

If $D > 0$, then, if four real roots are to exist, in the sequence
\[
1, \quad H, \quad H^2 - 16 A f^2, \quad -T^2
\]
three changes of sign must occur, i.e.\ we must have
\[
H < 0, \quad H^2 - 16 A f^2 > 0
\]
hold. In the three other still possible combinations of signs all four roots are imaginary. Here, however, an arbitrary real value may be set for $x$.


% ===== scan leaf 277, printed 252 =====
% !PRINTED 252
\weberpage{252}
\pdfbookmark[2]{§ 79. The Bezoutiant and its significance for the reality of the roots}{bm89}

\begin{center}
\S.\,79.

\gesperrt{The Bezoutiant and its significance for the reality of the roots.}
\end{center}

For the investigation of the reality of the roots of an arbitrary real equation in more general cases, the function which we designated in \S.\,73 as the Bezoutiant of the equation can be applied with advantage.

If
\weq{1}{f(x) = x^n + a_1 x^{n-1} + a_2 x^{n-2} + \cdots + a_n = 0}
is any real equation of degree $n$, then the Bezoutiant was defined as follows:

One sets
\weq{2}{y = t_{n-1} f_0(x) + t_{n-2} f_1(x) + \cdots + t_0 f_{n-1}(x),}
in which the $t_0, t_1 \ldots t_{n-1}$ denote indeterminate variables and
\[
f_0(x) = 1, \quad f_1(x) = x + a_1, \quad f_2(x) = x^2 + a_1 x + a_2 \ldots
\]
Then [\S.\,73, (10)]
\weq{3}{S(y^2) = \frac{1}{n}\left[S(y)\right]^2 + B,}
where $B$ is a quadratic form of the $n - 1$ variables $t_0, t_1 \ldots t_{n-2}$ with coefficients that are composed rationally from the coefficients $a_i$. This function $B$ we have defined as the Bezoutiant, and we have actually constructed it for the cases $n = 4$ and $n = 5$.

Now in the sum on the left-hand side of (3)
\weq{4}{y_1^2 + y_2^2 + y_3^2 + \cdots + y_n^2,}
the $y_1, y_2 \ldots y_n$ are linear functions of the variables $t_0, t_1 \ldots t_{n-1}$, and thus in the formula (3) we have a transformation of the quadratic form
\weq{5}{\frac{1}{n}\left[S(y)\right]^2 + B = \Phi(t_0, t_1 \ldots t_{n-1})}
\gesperrt{into a sum of $n$ squares.}

For the present we make no restrictive assumptions whatever about equality or distinctness of the roots of $f(x)$, and must now first investigate to what extent the linear functions $y_1, y_2 \ldots y_n$ are independent of one another. In this regard the following theorem holds:
%CONT>>


% ===== scan leaf 278, printed 253 =====
% !PRINTED 253
\weberpage{253}
%<<CONT

% [corrected 2026-08-07] Removed a bare \S.\,79. copied from the running head;
% the actual section heading and title are on printed 252 (i0277).
\gesperrt{If $x_1, x_2 \ldots x_\nu$ are roots of equation (1) that are distinct
from one another, then $y_1, y_2 \ldots y_\nu$ are linearly independent.}

For suppose there existed between $y_1, y_2 \ldots y_\nu$ a linear
relation, identical with respect to $t$,
\[
\alpha_1 y_1 + \alpha_2 y_2 + \cdots + \alpha_\nu y_\nu = 0,
\]
whose coefficients $\alpha$ are independent of the $t$ and not all equal
to zero; then this would break up into the following equations
\[
\begin{aligned}
\alpha_1 f_0(x_1) + \alpha_2 f_0(x_2) + \cdots + \alpha_\nu f_0(x_\nu) &= 0, \\
\alpha_1 f_1(x_1) + \alpha_2 f_1(x_2) + \cdots + \alpha_\nu f_1(x_\nu) &= 0, \\
\cdot\ \cdot\ \cdot\ \cdot\ \cdot\ \cdot\ \cdot\ \cdot\ \cdot\ \cdot\ \cdot\ \cdot\ \cdot\ \cdot & \\
\alpha_1 f_{n-1}(x_1) + \alpha_2 f_{n-1}(x_2) + \cdots + \alpha_\nu f_{n-1}(x_\nu) &= 0.
\end{aligned}
\]
Since now $\nu \leqq n$,
the determinant of the first $\nu$
of these equations must vanish, hence
\[
\begin{vmatrix}
f_0(x_1), & f_1(x_1) & \ldots & f_{\nu-1}(x_1) \\
f_0(x_2), & f_1(x_2) & \ldots & f_{\nu-1}(x_2) \\
\cdot & \cdot & \cdot\ \cdot\ \cdot\ \cdot\ \cdot & \cdot \\
f_0(x_\nu), & f_1(x_\nu) & \ldots & f_{\nu-1}(x_\nu)
\end{vmatrix}
= 0,
\]
or
\[
\begin{vmatrix}
1, & x_1 + a_1, & x_1^2 + a_1 x_1 + a_2 \ldots \\
1, & x_2 + a_1, & x_2^2 + a_1 x_2 + a_2 \ldots \\
\cdot & \cdot & \cdot\ \cdot\ \cdot\ \cdot\ \cdot\ \cdot\ \cdot\ \cdot \\
1, & x_\nu + a_1, & x_\nu^2 + a_1 x_\nu + a_2 \ldots
\end{vmatrix}
= 0.
\]
But this determinant can be brought, by repeated application
of the theorem on the addition of columns \S.\,22, (VII),
to the form
\[
\pm
\begin{vmatrix}
1, & x_1, & x_1^2 & \ldots & x_1^{\nu-1} \\
1, & x_2, & x_2^2 & \ldots & x_2^{\nu-1} \\
\cdot & \cdot & \cdot & \cdot\ \cdot\ \cdot\ \cdot & \cdot \\
1, & x_\nu, & x_\nu^2 & \ldots & x_\nu^{\nu-1}
\end{vmatrix}
=
\begin{aligned}
&(x_1 - x_2)(x_1 - x_3) \ldots (x_1 - x_\nu) \\
&\hphantom{(x_1 - x_2)}(x_2 - x_3) \ldots (x_2 - x_\nu) \\
&\hphantom{(x_1 - x_2)(x_2 - x_3)}\cdot\ \cdot\ \cdot\ \cdot\ \cdot \\
&\hphantom{(x_1 - x_2)(x_2 - x_3)\cdot\cdot}(x_{\nu-1} - x_\nu),
\end{aligned}
\]
and it therefore cannot vanish if, as was assumed,
the $x_1, x_2 \ldots x_\nu$ are distinct.

If, therefore, in the sum (4) we collect all terms that are equal
to one another, \gesperrt{there remain as many squares
of linearly independent functions as equation (1)
has distinct roots.}


% ===== scan leaf 279, printed 254 =====
% !PRINTED 254
\weberpage{254}

Now if $x_1$ is a real root of (1), then $y_1^2$ is a positive square in the sum (4). But if $x_1$ and $x_2$ are conjugate imaginary, then $y_1$ and $y_2$ likewise split into two conjugate imaginary constituents
\[
y_1 = u + vi, \quad y_2 = u - vi,
\]
where $u$ and $v$ are real linear functions of the $t$. Accordingly we have
\[
y_1^2 + y_2^2 = 2u^2 - 2v^2.
\]

% The source visibly letterspaces only this occurrence of "verschiedenen Wurzeln".
Thus by (4) the function $\Phi(t)$ defined by (5) is decomposed into a sum of squares, whose number is equal to the number of \gesperrt{distinct roots} of $f(x) = 0$, and among which there are as many negative ones as there are pairs of imaginary roots among these mutually distinct roots.

But by the law of inertia of quadratic forms (\S.\,58) these numbers remain the same under every other real linear transformation of the quadratic form into a sum of squares. Now $S(y)$ is a real linear function of the $t$, and $B$ no longer contains the variable $t_{n-1}$; from this, formula (5) yields the theorem:

\gesperrt{I. If the Bezoutiant $B$ is decomposed by a real linear transformation into a sum of $\pi$ positive and $\nu$ negative squares, which cannot be reduced to a smaller number, then $\pi + \nu + 1$ is the number of distinct roots of $f(x) = 0$; and $\nu$ is the number of pairs of conjugate imaginary roots contained among them, $\pi - \nu + 1$ that of the real ones.}

If the determinant of $B$, and hence also the discriminant of $f(x)$ (\S.\,73), is different from zero, then $\pi + \nu = n - 1$ and $\nu$ is less than or at most equal to $\frac{1}{2}n$.

The Bezoutiant is thus a quadratic form of a \gesperrt{special nature}, which expresses itself precisely in the fact that among the squares into which it can be decomposed, at most $\frac{1}{2}n$ negative ones can occur.


% ===== scan leaf 280, printed 255 =====
% !PRINTED 255
\weberpage{255}
\pdfbookmark[2]{§ 80. The inertia of forms of the second degree}{bm90}

\begin{center}
\S.\,80.

\gesperrt{The Inertia of Forms of the Second Degree.}
\end{center}

By the theorem of the preceding section we are directed to the investigation of the homogeneous functions of the second degree with real coefficients. Concerning these functions we became acquainted in \S.\,58, under the name of the law of inertia, with a theorem which forms the foundation for what follows.

The theorem consisted in this: however one may transform a quadratic form of $m$ variables into a sum of positive and negative squares of linearly independent linear functions, which is possible in infinitely many different ways, the number of the positive squares, and likewise that of the negative ones, is always the same.

If we denote by $\pi$ the number of the positive squares and by $\nu$ the number of the negative ones, then $\pi + \nu$ is at most equal to $m$.

If we denote the difference $m - \pi - \nu$ by $\varrho$, then, when we represent the general quadratic form of $m$ variables as a sum of $m$ squares, $\varrho$ can be designated as the number of squares with vanishing coefficients.

The three numbers $\pi$, $\nu$, $\varrho$, of which none can be negative and whose sum is equal to the number of the variables,
\weq{1}{\pi + \nu + \varrho = m,}
are therefore unalterable for a given quadratic form.

We have to seek means of determining the numbers $\pi$, $\nu$, $\varrho$ from the coefficients of the quadratic form. If we apply these means to the Bezoutiant, we obtain criteria for inferring, from the coefficients of an equation, the number of its real roots.

Now, then, let, as in \S.\,56
\weq{2}{\varphi(x_1, x_2 \ldots x_m) = \sum_{1,m}^{i,k} a_{i,k} x_i x_k, \qquad a_{i,k} = a_{k,i}}
be a quadratic form of $m$ variables with \gesperrt{real} coefficients $a_{i,k}$.


% ===== scan leaf 281, printed 256 =====
% !PRINTED 256
\weberpage{256}

The determinant

\weq{3}{R = \Sigma \pm a_{1,1},\, a_{2,2} \ldots a_{m,m}}

shows by its vanishing that $\varphi$ can be transformed by a linear transformation into a function of fewer than $m$ variables, and hence that $\varrho$ is greater than zero.

The determinant $R$ is a symmetric determinant. Among its subdeterminants of various orders there occur certain ones which are again symmetric determinants, namely those which one obtains from $R$ by striking out rows and columns which intersect in diagonal terms; those, therefore, which, if $\alpha, \beta, \gamma \ldots$ denote any of the indices $1, 2 \ldots m$, are to be denoted by

\weq{4}{\Sigma \pm a_{\alpha,\alpha}\, a_{\beta,\beta}\, a_{\gamma,\gamma} \ldots}

These we shall call the principal subdeterminants of $R$.

We shall now determine the number of the principal subdeterminants.

The number of $\mu$-rowed principal subdeterminants is equal to the number of ways in which groups of $\mu$ different indices can be selected from the sequence of indices $1, 2, 3 \ldots m$, hence equal to the number of combinations of $m$ elements taken $\mu$ at a time, and this number is equal to the binomial coefficient $B_\mu^m$. The number of all the principal subdeterminants is therefore, if we count the determinant $R$ itself and in addition unity as a zero-rowed determinant,
\[
1 + m + \frac{m(m-1)}{1.2} + \cdots + m + 1 = 2^m.
\]

But usually only a small part of these actually needs to be taken into account.

The indices $1, 2, 3 \ldots m$ can be arranged in $\Pi(m) = 1.2.3 \ldots m$ different ways; if we form the determinant with any one of these arrangements

% [the printed determinant carries a dotted-line annotation boxing the entries that
%  intersect the diagonal (a_{1,1}; a_{2,1},a_{2,2}; a_{3,1},a_{3,2},a_{3,3}) -- a graphical
%  overlay in the print pointing out the rows/columns that get struck out; reproduced here
%  as the plain determinant, the box itself is not a mathematical symbol]
\[
\begin{vmatrix}
a_{1,1}, & a_{1,2}, & a_{1,3} & \cdots & a_{1,m} \\
a_{2,1}, & a_{2,2}, & a_{2,3} & \cdots & a_{2,m} \\
a_{3,1}, & a_{3,2}, & a_{3,3} & \cdots & a_{3,m} \\
\cdot & \cdot & \cdot & \cdots & \cdot \\
a_{m,1}, & a_{m,2}, & a_{m,3} & \cdots & a_{m,m}
\end{vmatrix},
\]

then from it a series of principal subdeterminants can be
%CONT>>


% ===== scan leaf 282, printed 257 =====
% !PRINTED 257
\weberpage{257}
\pdfbookmark[2]{§ 81. Quadratic forms with vanishing determinant}{bm91}
%<<CONT
derive, as is indicated by the strokes, that is, in such a way that
one obtains each preceding one from the one that follows it by
omitting the last row and column; thus
\[
R_0 = 1, \quad R_1 = a_{1,1}, \quad R_2 = a_{1,1}a_{2,2} - a_{1,2}^2,
\]
\[
R_3 = \begin{vmatrix} a_{1,1}, & a_{1,2}, & a_{1,3} \\ a_{2,1}, & a_{2,2}, & a_{2,3} \\ a_{3,1}, & a_{3,2}, & a_{3,3} \end{vmatrix}, \ldots R_m = R.
\]
Such a system, when it is ordered in the descending sequence
\[
R_m,\ R_{m-1} \ldots R_1,\ R_0
\]
shall be called a chain of principal subdeterminants.
Of such chains, $\Pi(m)$ different ones can be formed, in all of
which the first term is $R$ and the last term is 1.

\begin{center}
\S.\,81.

\gesperrt{Quadratic Forms with Vanishing Determinant.}
\end{center}

If the determinant $R$ of the quadratic form
$\varphi(x_1, x_2 \ldots x_m)$, defined in (3), \S.\,80, vanishes, then, as we
have already seen in \S.\,57, $\varphi$ can be expressed by fewer than $m$
mutually independent linear functions of $x$. If we denote by
$m - \varrho$ the smallest number of variables by which $\varphi$ can be
expressed, then $\varrho$ has the same meaning as in the preceding
section.

Now if among the principal subdeterminants of $R$ there is
found a $k$-rowed one which is different from zero, say
\[
R_k = \Sigma \pm a_{1,1}a_{2,2} \ldots a_{k,k},
\]
then, if we set
\[
x_{k+1} = 0,\quad x_{k+2} = 0, \ldots x_m = 0
\]
the expression
\[
\varphi(x_1, x_2 \ldots x_k, 0, 0, 0)
\]
is a quadratic form of $k$ variables $x_1, x_2 \ldots x_k$ whose
determinant $R_k$ is different from zero, and which accordingly cannot
be expressed by fewer than $k$ variables. Still less, then, can the
function $\varphi$, when $x_1, x_2 \ldots x_m$ vary without restriction,
depend on fewer than $k$ variables, and it follows that
% The source prints the rank/nullity bound with a less-than-or-equal sign.
\weq{1}{\varrho \leqq m - k.}


% ===== scan leaf 283, printed 258 =====
% !PRINTED 258
\weberpage{258}
We now assume that the form $\varphi(x_1, x_2 \ldots x_m)$ can be
expressed by means of $k$ mutually independent linear forms in $x$,
which we shall denote as follows:
\weq{2}{
\begin{aligned}
z_1 &= \beta_1^{(1)}x_1 + \beta_2^{(1)}x_2 + \cdots + \beta_k^{(1)}x_k + \beta_{k+1}^{(1)}x_{k+1} + \cdots + \beta_m^{(1)}x_m \\
&\cdot\ \cdot\ \cdot\ \cdot\ \cdot\ \cdot\ \cdot\ \cdot\ \cdot\ \cdot\ \cdot\ \cdot\ \cdot\ \cdot\ \cdot\ \cdot\ \cdot\ \cdot \\
z_k &= \beta_1^{(k)}x_1 + \beta_2^{(k)}x_2 + \cdots + \beta_k^{(k)}x_k + \beta_{k+1}^{(k)}x_{k+1} + \cdots + \beta_m^{(k)}x_m
\end{aligned}
}

Since $z_1, z_2 \ldots z_k$ are to be linearly independent of one
another, at least one of the $k$-rowed determinants that can be formed
from the matrix
\[
\begin{matrix}
\beta_1^{(1)}, & \beta_2^{(1)} & \ldots & \beta_m^{(1)} \\
\cdot & \cdot & \cdot & \cdot \\
\beta_1^{(k)}, & \beta_2^{(k)} & \ldots & \beta_m^{(k)}
\end{matrix}
\]
must be different from zero. For if they were all equal to zero,
then the unknowns $h_1, h_2 \ldots h_k$, which are not all to vanish,
could be determined from the equations
\[
h_1 \beta_s^{(1)} + h_2 \beta_s^{(2)} + \cdots + h_k \beta_s^{(k)} = 0, \quad s = 1, 2, 3 \ldots m
\]
(\S.\,24), and $z_1, z_2 \ldots z_k$ would satisfy an equation
\[
h_1 z_1 + h_2 z_2 + \cdots + h_k z_k = 0
\]
and so would not be independent.

We may, however, without restricting the generality, assume
that the determinant
\[
\Sigma \pm \beta_1^{(1)}\beta_2^{(2)} \ldots \beta_k^{(k)}
\]
is different from zero, and then we can solve the equations
(2) for the variables $x_1, x_2 \ldots x_k$.

To these solutions we give the form
\weq{3}{
\begin{aligned}
x_1 &= y_1 + \alpha_{k+1}^{(1)}x_{k+1} + \cdots + \alpha_m^{(1)}x_m \\
x_2 &= y_2 + \alpha_{k+1}^{(2)}x_{k+1} + \cdots + \alpha_m^{(2)}x_m \\
&\cdot\ \cdot\ \cdot\ \cdot\ \cdot\ \cdot\ \cdot\ \cdot\ \cdot\ \cdot\ \cdot\ \cdot \\
x_k &= y_k + \alpha_{k+1}^{(k)}x_{k+1} + \cdots + \alpha_m^{(k)}x_m,
\end{aligned}
}
in which the $y_s$ are linear homogeneous combinations of the $z_s$,
which are determined from the equations
\[
\begin{aligned}
\beta_1^{(1)}y_1 + \beta_2^{(1)}y_2 + \cdots + \beta_k^{(1)}y_k &= z_1, \\
\beta_1^{(2)}y_1 + \beta_2^{(2)}y_2 + \cdots + \beta_k^{(2)}y_k &= z_2, \\
&\cdot\ \cdot\ \cdot\ \cdot\ \cdot\ \cdot\ \cdot\ \cdot\ \cdot\ \cdot\ \cdot\ \cdot \\
\beta_1^{(k)}y_1 + \beta_2^{(k)}y_2 + \cdots + \beta_k^{(k)}y_k &= z_k
\end{aligned}
\]
%CONT>>


% ===== scan leaf 284, printed 259 =====
% !PRINTED 259
\weberpage{259}

%<<CONT
be determined, while the $\alpha$ are new coefficients, which are derived
from the $\beta$ in an easily surveyed manner, and with whose
formation we are not further concerned here.

Now by hypothesis the function
$\varphi(x_1, x_2 \ldots x_m)$ can be expressed by $z_1, z_2 \ldots z_k$ alone, and hence also by $y_1, y_2 \ldots y_k$
alone. If we denote this expression by
$\Phi(y_1, y_2 \ldots y_k)$, then [by the mediation of (3)] we have the
identity
\weq{4}{\varphi(x_1, x_2 \ldots x_k, x_{k+1} \ldots x_m) = \Phi(y_1, y_2 \ldots y_k),}
and if we set $x_{k+1}, x_{k+2}, \ldots x_m = 0$ in it,
\weq{5}{\Phi(x_1, x_2 \ldots x_k) = \varphi(x_1, x_2 \ldots x_k, 0, 0 \ldots 0),}
whereby, since the designation of the variables does not matter,
$\Phi$ is completely determined; we may therefore set
\begin{align}
\Phi(y_1, y_2 \ldots y_k) &= \varphi(y_1, y_2 \ldots y_k, 0, 0 \ldots 0) \tag{6}\\
&= \varphi(x_1, x_2 \ldots x_m);
\end{align}
$\Phi$ thus arises from $\varphi$ by setting the $m-k$ last
variables (in the arrangement chosen here) equal to zero:
\weq{7}{\Phi(x_1, \ldots x_k) = \sum_{1,k}^{r,s} a_{r,s} x_r x_s.}

The determinant of $\Phi$ is one of the principal subdeterminants of $R$, and indeed one obtains it by striking out the
$m-k$ last rows and columns in $R$.

If we assume that $\varphi$ cannot be expressed by still fewer than
$k$ variables, so that $k = m - \varrho$, then
the determinant of this function $\Phi$ cannot vanish.

From this and from what was proved above there now results
the theorem:

\gesperrt{II.} \gesperrt{If all principal subdeterminants of $R$ of
more than $k$ rows vanish, while among
the principal subdeterminants of $k$ rows
at least one is different from zero, then}
\[
\varrho = m - k.
\]
For it follows from (7) that, if $\varrho = m - k$, at least
one $k$-rowed principal subdeterminant must be different from
zero, and from (1) that, if in addition a principal subdeterminant of more than $k$ rows is different from zero, then
\[
\varrho < m - k
\]
holds.


% ===== scan leaf 285, printed 260 =====
% !PRINTED 260
\weberpage{260}
\pdfbookmark[2]{§ 82. Quadratic forms with non-vanishing determinant}{bm92}

% Running head "Seventh Section." excluded as page furniture.

%<<CONT
If $\pi, \nu, \varrho$ are the number of positive, negative, vanishing squares into which $\varphi$ can be decomposed, then $\pi, \nu$ are the
numbers of positive and negative squares into which the
form $\Phi$ determined by (7) can be decomposed, and the determination of the numbers $\pi, \nu$ therefore need only be carried out for the latter function, whose determinant is different from
zero.

\begin{center}
\S.\,82.

\gesperrt{Quadratic forms with non-vanishing\\
determinant.}
\end{center}

In the investigation of the forms $\varphi(x_1, x_2 \ldots x_m)$ with non-vanishing determinant we make use of a theorem on determinants which we proved at the end of \S.\,28, and which
we wish to recall here.

If $A$ is a determinant of $m$ rows, and if $A_i^{(k)}$ are its
first, $A_{i,i'}^{k,k'}$ its second subdeterminants, then
\weq{1}{A_1^{(1)} A_i^{(i)} - A_1^{(i)} A_i^{(1)} = A A_{1,i}^{1,i}.}
If $A$ is a \gesperrt{symmetric determinant} different from zero, then $A_1^{(i)} = A_i^{(1)}$ and we can write formula (1)
thus:
\weq{2}{A_1^{(1)} A_i^{(i)} - A_1^{(i)2} = A A_{1,i}^{1,i}.}
In this, $i$ may denote any of the indices $2, 3 \ldots m$. If
$A_1^{(1)} = 0$ and $A_{1,i}^{1,i} = 0$, it follows from this that $A_1^{(i)} = 0$
must also hold, and from this we conclude that
\[
A_1^{(1)}, \; A_{1,2}^{1,2}, \; A_{1,3}^{1,3} \ldots A_{1,m}^{1,m}
\]
cannot vanish simultaneously, since otherwise
\[
A_1^{(2)}, \; A_1^{(3)} \ldots A_1^{(m)},
\]
and hence also, contrary to the hypothesis, $A$ would vanish (\S.\,22).

If we apply this to the determinant $R = R_m$ of our function $\varphi$, now assumed to be different from zero, it
follows that though the first principal subdeterminant $R_{m-1}$ may vanish,
not all the second principal subdeterminants $R_{m-2}$ can
vanish. The same mode of reasoning can be applied
when we replace $R$ by a non-vanishing $R_{m-1}, R_{m-2} \ldots$,
and we thus arrive at the following important theorem:
%CONT>>


% ===== scan leaf 286, printed 261 =====
% !PRINTED 261
\weberpage{261}
\pdfbookmark[2]{§ 83. Number of the positive and negative squares}{bm93}

III.\ \gesperrt{If $R$ is different from zero, one can arrange the indices $1, 2 \ldots m$ so that in the chain of principal subdeterminants}
\weq{3}{R_m,\ R_{m-1} \ldots R_1,\ R_0}
\gesperrt{no two consecutive terms vanish.}

The determinant relation (2) yields, if one sets
\[
A = R_{k+1}, \quad (k = 1, 2 \ldots m-1)
\]
an equation between three consecutive terms of the chain (3), which we can write thus:
\weq{4}{R_k S_k - T_k^2 = R_{k-1} R_{k+1},}
in which $S_k$ and $T_k$ are certain subdeterminants of $R$, and hence integral rational functions of the coefficients $a_{i,k}$.

More precisely described, $R_k$, $S_k$, $T_k$ are first subdeterminants of $R_{k+1}$, and indeed
\[
R_k = \frac{\partial R_{k+1}}{\partial a_{k+1,k+1}}, \quad S_k = \frac{\partial R_{k+1}}{\partial a_{k,k}}, \quad T_k = \frac{\partial R_{k+1}}{\partial a_{k,k+1}}.
\]

An arrangement of the indices corresponding to theorem (3) we shall suppose chosen for what follows. Then, if $R_k$ vanishes, $R_{k-1}$ and $R_{k+1}$ are different from zero, and (4) shows that they have opposite signs; hence:

IV.\ \gesperrt{If an inner term of a chain of principal subdeterminants vanishes, then the two adjoining terms have opposite signs.}

\begin{center}
\S.\,83.\\[4pt]
\gesperrt{Number of the positive and negative squares.}
\end{center}

In order to determine the number of positive and negative squares of a form with non-vanishing determinant, we assume first that in the chain of principal subdeterminants
\[
R_m,\ R_{m-1},\ R_{m-2} \ldots R_1,\ R_0 = 1
\]
no term vanishes. We can then determine the coefficient $\lambda$ so that the determinant of the form
%CONT>>


% ===== scan leaf 287, printed 262 =====
% !PRINTED 262
\weberpage{262}

%<<CONT
\weq{1}{\psi = \varphi(x_1, x_2, x_3 \ldots x_m) - \lambda x_m^2}
vanishes. For the determinant is
\[
\begin{vmatrix}
a_{1,1},\, a_{1,2} \ldots a_{1,m} \\
a_{2,1},\, a_{2,2} \ldots a_{2,m} \\
\cdots\cdots\cdots\cdots \\
a_{m,1},\, a_{m,2} \ldots a_{m,m} - \lambda
\end{vmatrix}
= R_m - \lambda R_{m-1},
\]
and it vanishes when we set
\[
\lambda = \frac{R_m}{R_{m-1}}
\]

The function $\psi$ can then be expressed in terms of $m-1$ variables $y$, and by formula (6), \S.\,81, one obtains
\[
\psi(x_1, x_2 \ldots x_m) = \psi(y_1, y_2 \ldots y_{m-1}, 0),
\]
or, by (1),
\weq{2}{\varphi(x_1, x_2 \ldots x_m) = \frac{R_m}{R_{m-1}}\, x_m^2 + \varphi(y_1, y_2 \ldots y_{m-1}, 0),}
and the investigation of the function $\varphi$ of $m$ variables is thereby reduced to the investigation of $\varphi_1(y) = \varphi(y_1, y_2 \ldots y_{m-1}, 0)$ of $m-1$ variables, whose determinant is equal to $R_{m-1}$ and hence different from zero. According as $R_m : R_{m-1}$ is positive or negative, the function $\varphi(x)$ will have one more positive or one more negative square than $\varphi_1(y)$.

By applying the same procedure to $\varphi_1(y)$ and the following functions, we obtain the theorem:

The number of positive and negative terms of the sequence
\weq{3}{\frac{R_m}{R_{m-1}}, \quad \frac{R_{m-1}}{R_{m-2}}, \ldots \frac{R_1}{R_0}}
agrees with the number of positive and negative squares into which the function $\varphi(x)$ can be decomposed.

We wish to give this theorem a somewhat different expression, but first premise the following explanation.

If a sequence of real numbers different from zero is given in a definite arrangement, the signs of these quantities can change in manifold ways; if two quantities of the same sign follow one another, there is a \gesperrt{succession of signs} (permanence), but if a quantity is followed by another of opposite sign, we have a \gesperrt{change of sign} (variation).


% ===== scan leaf 288, printed 263 =====
% !PRINTED 263
\weberpage{263}

If we now consider the sequence of quantities from this point of view,
\[
R_m,\ R_{m-1},\ R_{m-2} \ldots R_1,\ R_0,
\]
then in passing from $R_k$ to $R_{k-1}$ there occurs a permanence of sign or a change of sign, according as the quotient $R_k : R_{k-1}$ is positive or negative.

We can therefore also state the following theorem:

V.\ \gesperrt{If $\pi$ is the number of positive, $\nu$ that of the negative squares of $\varphi$, then $\pi$ is equal to the number of permanences of sign, $\nu$ equal to the number of changes of sign in the chain}
\weq{4}{R_m,\ R_{m-1},\ R_{m-2} \ldots R_1,\ R_0.}

This formulation of the theorem has the advantage that it can be carried over to the case in which individual interior terms of the sequence (4) vanish, provided only that no two consecutive terms are zero (which by \S.\,82, III. may always be assumed). For if $R_k = 0$, while $R_{k-1}$, $R_{k+1}$ are different from zero, then by \S.\,82, Theorem IV, $R_{k-1}$ and $R_{k+1}$ have different signs. In the sequence
\[
R_{k+1},\ R_k,\ R_{k-1}
\]
there thus occurs one change of sign and one permanence of sign, no matter whether we replace the vanishing $R_k$ by a positive or by a negative quantity.

If now we replace the chain (4)
\[
R_m,\ R_{m-1} \ldots R_1,\ R_0,
\]
in which individual terms vanish, by another,
\weq{5}{R'_m,\ R'_{m-1} \ldots R'_1,\ R'_0,}
in which no term vanishes, and in which to the non-vanishing terms of the sequence (4) there correspond terms of the same sign, then the sequences (4) and (5) have equally many changes of sign, whatever signs the $R'$ corresponding to the vanishing $R$ may have.

Now let $\psi(x)$ be a second arbitrary quadratic form of the variables $x$, with which we form the form
\weq{6}{\varphi' = \varphi + \varepsilon \psi}
in which $\varepsilon$ is a still undetermined coefficient. We shall now show at once that we can choose $\varepsilon$ in such a way that $\varphi'$ and $\varphi$ have the same number of positive and negative squares
%CONT>>


% ===== scan leaf 289, printed 264 =====
% !PRINTED 264
\weberpage{264}
%<<CONT
have; that, however, in the chain of the principal subdeterminants
$R'_k$ of the form $\varphi'$ no vanishing terms occur; and that,
finally, to a non-vanishing $R_k$ there corresponds an $R'_k$ of the
same sign. Then the numbers $\pi$, $\nu$ for
$\varphi$ and for $\varphi'$ can be determined both from the series (4) and from the series
(5), and the number of \gesperrt{changes of sign}, which
is the same in both, gives the number $\nu$ of the negative squares.

Now, in order to prove that the coefficient $\varepsilon$ can
be determined in the manner indicated, we assume
that $\varphi$ has somehow been transformed into a sum of squares
\[
\varphi = \lambda_1 y_1^2 + \lambda_2 y_2^2 + \cdots + \lambda_m y_m^2,
\]
and that $\psi$, expressed in the variables $y_1, y_2 \ldots y_m$, has the
expression
\[
\psi = \Sigma \beta_{i,k} y_i y_k.
\]
The numbers $\pi$, $\nu$ for the function $\varphi'$ are then determined from the
chain of the principal subdeterminants of
\[
\begin{vmatrix}
\lambda_1 + \varepsilon\beta_{1,1}, & \varepsilon\beta_{1,2} \ldots & \varepsilon\beta_{1,m} \\
\varepsilon\beta_{2,1}, & \lambda_2 + \varepsilon\beta_{2,2} \ldots & \varepsilon\beta_{2,m} \\
\cdot & \cdot & \cdot \\
\varepsilon\beta_{m,1}, & \varepsilon\beta_{m,2} \ldots & \lambda_m + \varepsilon\beta_{m,m}
\end{vmatrix}
\]
according to Theorem V.

Now, however, one may take $\varepsilon$ so small that these principal
subdeterminants agree \gesperrt{in sign} with
\[
\lambda_1\lambda_2\lambda_3 \ldots \lambda_m, \quad \lambda_1\lambda_2\lambda_3 \ldots \lambda_{m-1}, \quad \ldots, \quad \lambda_1\lambda_2, \quad \lambda_1, \quad 1,
\]
and then the number of positive and negative squares
of $\varphi'$ is equal to the number of positive and negative ones among the
coefficients $\lambda$, of which none vanishes, i.\,e. the numbers
$\pi$ and $\nu$ are the same for $\varphi$ and $\varphi'$.

If now the coefficients of $\psi$, expressed in the original
variables, are $b_{i,k}$, so that
\[
\psi = \Sigma b_{i,k} x_i x_k,
\]
then one of the principal subdeterminants of $\varphi'$ is
\[
R'_k = \begin{vmatrix}
a_{1,1} + \varepsilon b_{1,1} & \ldots & a_{1,k} + \varepsilon b_{1,k} \\
\cdot & \cdot & \cdot \\
a_{k,1} + \varepsilon b_{k,1} & \ldots & a_{k,k} + \varepsilon b_{k,k}
\end{vmatrix},
\]
and is thus an integral rational function of degree $k$ in $\varepsilon$,
\[
R'_k = R_k + \varepsilon M_1 + \varepsilon^2 M_2 + \ldots \varepsilon^k M_k,
\]
%CONT>>


% ===== scan leaf 290, printed 265 =====
% !PRINTED 265
\weberpage{265}
\pdfbookmark[2]{§ 84. Application to the Bezoutiant}{bm94}
%<<CONT
in which the $M_1, M_2 \ldots M_k$ depend rationally on the $a_{i,k}, b_{i,k}$. In particular,
\[
M_k = \Sigma \pm b_{1,1} b_{2,2} \ldots b_{k,k},
\]
and one can always choose the $b_{i,k}$ so that $M_k$ is different from zero. By \S.\,31 one can therefore choose $\varepsilon$ so that, if $R_k$ is different from zero, $R_k'$ has the same sign as $R_k$, and if $R_k$ vanishes, $R_k'$ does not vanish.

We can combine what has thus been proved with the results of \S.\,81 into a general rule for determining the number of negative squares, valid also in the case of vanishing determinant.

% The source continues to use the same variant rho (varrho) introduced in section 80.
VI. \gesperrt{To determine the numbers $\pi$, $\nu$, $\varrho$ of the function $\varphi(x)$, arrange the variables $x_1, x_2 \ldots x_m$ so that in the chain of principal subdeterminants}
\weq{7}{R_m, R_{m-1} \ldots R_1, R_0}
\gesperrt{as small a number of initial terms as possible vanishes, and so that of the following terms no two standing next to one another vanish; $\varrho$ is then the number of vanishing initial terms, $\nu$ the number of changes of sign, and $\pi = m - \nu - \varrho$.}

Finally let it be remarked that the number of changes of sign in the series (7) can also be counted from right to left, that is, one finds the same number of changes of sign if the series is written in the reverse order.

\begin{center}
\S.\,84.

\gesperrt{Application to the Bezoutiant.}
\end{center}

The exact discussion of the inertia of quadratic forms had for us only the one purpose of determining the number of real and imaginary roots of an algebraic equation by counting the positive and negative squares of the Bezoutiant. If $n$ is the degree of the equation, its Bezoutiant is a quadratic form in $n - 1$ variables, and in \S.\,73 we formed it completely for $n = 3, 4, 5$
%CONT>>


% ===== scan leaf 291, printed 266 =====
% !PRINTED 266
\weberpage{266}

%<<CONT
and have learned the ways by which one can arrive at their computation in other cases as well.

We have already remarked earlier that the Bezoutiant is not a general quadratic form, but that it has the special peculiarity that the number $\nu$ of its negative squares can never become greater than $\frac{1}{2}n$. This peculiarity must find its expression in certain inequality conditions between the coefficients, which, however, are not yet known at present. Only in the case $n = 3$ can we state the algebraic character of this restriction completely.

For the cubic equation
\[
f(x) = a_0 x^3 + a_1 x^2 + a_2 x + a_3 = 0
\]
the Bezoutiant is, by \S.\,73, (13), nothing other than the \gesperrt{Hessian} covariant
\[
-\tfrac{2}{3} H(t_1, t_0) = -\tfrac{2}{3}(A_0 t_1^2 + A_1 t_1 t_0 + A_2 t_0^2).
\]
But if $D$ is the discriminant, so that
\[
3D = 4A_0 A_2 - A_1^2
\]
there holds the relation [\S.\,62, (4), (12)]
\[
4H^3 + Q^2 + 27Df^2 = 0,
\]
and hence, if we set $t_0 = 0$ in it,

\weq{1}{-4A_0^3 + q_0^2 + 27Da_0^2 = 0,}

where the coefficient appearing above is given by
\[
q_0 = 27a_0^2 a_3 - 9a_0 a_1 a_2 + 2a_1^3\text{.}
\]
This relation shows that $A_0$ and $D$ cannot both be positive at the same time; and hence that the form $-H$ is not decomposable into two negative squares.

We wish to apply the general theorem of the preceding section to the biquadratic equation as well, which for the sake of simplicity we assume in the form

\weq{2}{x^4 + ax^2 + bx + c = 0}

To form the coefficients of the Bezoutiant, we have, in the formulas \S.\,73, (14), to replace
\[
a_0,\ a_1,\ a_2,\ a_3,\ a_4
\]
by
\[
1,\ 0,\ a,\ b,\ c
\]
We thus obtain
\[
\begin{aligned}
&B_{2,2} = -2a, \quad B_{1,1} = a^2 - 4c, \quad B_{0,0} = -\tfrac{1}{4}(8ac - 3b^2), \\
&B_{0,1} = \tfrac{1}{2}ab, \quad B_{2,0} = -4c, \quad B_{1,2} = -3b.
\end{aligned}
\]
% [resolved 2026-08-06, 8x zoom strip 8/8] The smudged sign inside B_{0,0} is a MINUS.
% Compared against the clean minus signs on these same two lines (-2a, -4c, -3b, -1/4) and
% the plus signs three lines above (-4A_0^3 + q_0^2 + 27 D a_0^2): the disputed stroke has
% no vertical bar, only a horizontal one with a speck of ink. B_{0,0} = -1/4(8ac - 3b^2)
% stands as transcribed. Same strip also confirms B_{2,2} = -2a, B_{1,1} = a^2 - 4c,
% B_{0,1} = 1/2 ab, B_{2,0} = -4c, B_{1,2} = -3b.


% ===== scan leaf 292, printed 267 =====
% !PRINTED 267
\weberpage{267}
We arrange the determinant thus:
\[
\begin{vmatrix}
-2a, & -3b, & -4c \\
-3b, & a^2 - 4c, & \tfrac{1}{2} a b \\
-4c, & \tfrac{1}{2} a b, & -\tfrac{1}{4} (8ac - 3b^2)
\end{vmatrix}
\]
and obtain the chain of principal subdeterminants
\weq{3}{D, \; -2a^3 + 8ac - 9b^2, \; -2a, \; 1,}
where $D$ denotes the discriminant, which by \S.\,64, (11) is determined by
\weq{4}{27 D = 4(a^2 + 12c)^3 - (2a^3 - 72ac + 27b^2)^2.}
The criterion for all four roots to be real is accordingly
\weq{5}{D > 0, \quad -2a^3 + 8ac - 9b^2 > 0, \quad -2a > 0.}
This criterion is apparently different from, and in any case less simple than, the one set up in \S.\,78
\weq{6}{D > 0, \quad a < 0, \quad a^2 - 4c > 0.}

But we now wish to show further that the two say exactly the same thing.

First of all it is immediately seen that (6) must be satisfied when (5) is satisfied, for
\[
- 2a(a^2 - 4c) - 9b^2
\]
cannot be positive when $a$ and $a^2 - 4c$ are negative; but to see conversely that the conditions (5) follow from (6), we must take the value of $D$ into consideration.

For brevity we set
\weq{7}{\alpha = -2a, \quad \beta = -6a^3 + 24ac - 27b^2, \quad \gamma = 3a^2 - 12c,}
so that the conditions (5) read
\[
\alpha > 0, \quad \beta > 0, \quad D > 0,
\]
while the conditions (6) read
\[
\alpha > 0, \quad \gamma > 0, \quad D > 0\text{.}
\]
It is therefore to be shown that if $\alpha$, $\gamma$ and $D$ are positive, $\beta$ must also be positive. To this end we express $D$ in (4), by means of (7), in terms of $\alpha, \beta, \gamma$ and obtain
\[
27 D = 4(\alpha^2 - \gamma)^3 - [2\alpha(\alpha^2 - \gamma) - \beta]^2.
\]
If $D$ is positive, then by this $\alpha^2 - \gamma$ must certainly be positive. If we now set
%CONT>>


% ===== scan leaf 293, printed 268 =====
% !PRINTED 268
\weberpage{268}
%<<CONT
\[
\alpha^2 - \gamma = \delta^2,
\]
then $\alpha^2 > \delta^2$, and hence, if $\delta$ is taken positive,
\[
\alpha > \delta
\]
and
\[
27 D = 4 \delta^6 - (2\alpha\delta^2 - \beta)^2.
\]
By this, $D$ cannot be positive when $\beta$ is negative; for if $\beta$ is negative, then
\[
2\alpha\delta^2 - \beta > 2\alpha\delta^2 > 2\delta^3,
\]
and so $4\delta^6 - (2\alpha\delta^2 - \beta)^2$ is negative.

It is hereby directly proved that the criteria (5) and (6) say exactly the same thing.

According to the geometric interpretation given in \S.\,78, $\beta = 0$ means a surface of the third order which passes through the parabola $b = 0$, $\gamma = 0$, and which, so far as negative values of $a$ come into consideration, runs entirely within the region of space in which $D$ is negative.

In order to have before our eyes an example for the equation of the fifth degree, we consider the equation
\[
x^5 + x^2 + a = 0.
\]
We therefore have to set, in our general expressions of \S.\,73 or also of \S.\,74, (2), (3),
\[
a_0, a_1, a_2, a_3, a_4, a_5 =
\]
\[
1, 0, 0, 1, 0, a,
\]
and we obtain for the determinant of the Bezoutiant
\[
\begin{vmatrix}
-2a, & 0, & 0, & -5a \\
0, & \tfrac{6}{5}, & -5a, & 0 \\
0, & -5a, & 0, & -3 \\
-5a, & 0, & -3, & 0
\end{vmatrix}
\]
and for the discriminant
\[
D = 108 a + 3125 a^4.
\]
A chain of principal subdeterminants is
\weq{8}{D, \; 50a^3, \; \frac{-12a}{5}, \; -2a, \; 1.}
The discriminant is negative when
\weq{9}{0 < -a < \sqrt[3]{\frac{108}{3125}},}
otherwise positive, and so in particular positive for all positive $a$. One sees that when the discriminant is positive there are always two
%CONT>>


% ===== scan leaf 294, printed 269 =====
% !PRINTED 269
\weberpage{269}
%<<CONT
sign changes, whereas for negative discriminant, where $a$ is also negative, \gesperrt{one} sign change takes place.

So long as $a$ lies in the interval (9), our equation therefore has two imaginary roots; if it lies outside this interval, four imaginary roots, never four real roots.

A purely numerical example is offered by the equation, drawn from the complex multiplication of the elliptic functions,
\[
x^5 - x^3 - 2x^2 - 2x - 1 = 0.
\]
For the determinant of the Bezoutiant we obtain from \S.\,73
\[
\begin{vmatrix}
\dfrac{-4}{5}, & \dfrac{-3}{5}, & \dfrac{4}{5}, & 5 \\[6pt]
\dfrac{-3}{5}, & \dfrac{4}{5}, & \dfrac{33}{5}, & 8 \\[6pt]
\dfrac{4}{5}, & \dfrac{33}{5}, & \dfrac{46}{5}, & 6 \\[6pt]
5, & 8, & 6, & 2
\end{vmatrix}
\]
the chain of principal subdeterminants
\[
\frac{47^2}{5}, \; \frac{94}{5}, \; -1, \; \frac{-4}{5}, \; 1,
\]
so that the discriminant of our equation is $47^2$. The equation therefore has two pairs of imaginary roots and one real root. The second term of this chain, $94 : 5$, need not be computed if one knows the discriminant; indeed, the mere knowledge of the negative sign in the penultimate term already suffices to reach the foregoing decision on the reality of the roots.


% ===== scan leaf 295, printed 270 =====
% !PRINTED 270
\weberpage{270}
\pdfbookmark[1]{Eighth Section. Sturm's Theorem.}{bm95}
\pdfbookmark[2]{§ 85. Sturm's problem}{bm96}
\begin{center}
\gesperrt{Eighth Section.}

\gesperrt{Sturm's Theorem.}

\rule{0.3\linewidth}{0.4pt}

\S.\,85.

\gesperrt{Sturm's Problem.}
\end{center}

The question treated in the previous section, of the number of real roots of an equation, is a special case of a more general problem, which is to form the subject of this section, and which is the foundation of all methods for the approximate numerical computation of the roots of equations.

The question is this: \gesperrt{how many real roots of a real numerical equation lie between two given real values} $a$ and $b$?

Once this is decided, the further task is to narrow the bounds $a$, $b$ until only \gesperrt{one root} of the given equation lies between them, and finally to bring them so close to one another that each of them can be regarded as an \gesperrt{approximate value} of this root.

If we take $a = -\infty$, $b = +\infty$, or at any rate $a$ negative and $b$ positive and so large that no roots can lie beyond these bounds, then this problem coincides with the one treated in the previous section, of determining the number of all the real roots.

We shall first show how the general problem can be reduced to the special one, and thus how the question we have now raised is, in principle at least, answered by the considerations of the previous section. Let us first consider
%CONT>>


% ===== scan leaf 296, printed 271 =====
% !PRINTED 271
\weberpage{271}
\pdfbookmark[2]{§ 86. Sturm's chains}{bm97}
%<<CONT
the problem of determining the number of \gesperrt{positive roots} of an equation $f(x) = 0$. If we set $x = y^2$, then to every real value of $y$ there corresponds a positive value of $x$, and conversely to every positive value of $x$ there correspond two real, opposite values of $y$. % Source-original edge case: y=0 gives x=0, not a positive value, so the halving statement needs f(0) != 0 (or separate treatment of the zero root); the print is preserved.
\gesperrt{The number of positive roots of $f(x) = 0$ is therefore half as great as the number of real roots of $f(y^2) = 0$.} If we further set
\[
y = \frac{x - a}{b - x},
\]
% F(z^2) is correct: after f(x) has been transformed into F(y), the new substitution y=z^2 converts positive y-roots into pairs of real z-roots.
then, while $x$ goes from $a$ to $b$, $y$ will run through positive values from $0$ to infinity; and if by this substitution we transform $f(x)$ into $F(y)$, then $f(x) = 0$ will have just as many roots between $a$ and $b$ as $F(y) = 0$ has positive roots, and hence half as many as $F(z^2) = 0$ has real roots. Our task is thus in fact reduced to determining the number of real roots of a certain other equation, whose degree is twice as great as the degree of the given equation. But this reduction of the problem does not give its solution in the simplest form, and we must look for a simpler answer to the question.

\begin{center}
\S.\,86.

\gesperrt{The Sturm chains.}
\end{center}

The following consideration leads us to a solution of the problem that we posed in the preceding section.

Let $\alpha$ and $\beta$ be any two real numbers with $\alpha < \beta$. Further, let $f(x)$ be a given integral rational function of $x$, for which it is to be determined how many of its roots lie in the interval $\Delta$
\[
\alpha < x < \beta\text{.}
\]
% the German verb "liegen" follows the display; English absorbs "lie" before it
We assume that $\alpha$, $\beta$ do not themselves belong to the roots of $f(x) = 0$, so that $f(\alpha)$ and $f(\beta)$ are different from zero. Moreover, we shall assume for the present that $f(x)$ has no multiple root, at least in the interval $\Delta$, so that for no value of the interval are $f(x)$ and $f'(x)$ to vanish simultaneously.


% ===== scan leaf 297, printed 272 =====
% !PRINTED 272
\weberpage{272}
We assume that we can in some way construct a series of $m + 1$ continuous functions
\weq{1}{f(x), f_1(x), f_2(x) \ldots f_m(x), \quad (\mathfrak{R})}
possessing the following properties:

1. Of the functions $f_\nu$, no two consecutive ones shall vanish simultaneously in the interval $\Delta$.

2. The last of them, $f_m(x)$, shall not vanish at all in the interval $\Delta$, and hence shall keep an unchanging sign.

3. If a middle term, say $f_\nu(x)$, vanishes for some $x$ in the interval $\Delta$, then for this $x$ the two adjacent functions $f_{\nu-1}(x)$ and $f_{\nu+1}(x)$ shall have opposite signs.

4. If $f(x)$ vanishes in the interval $\Delta$, then $f_1(x)$ shall have the same sign for this value of $x$ as $f'(x)$.

Such a series of functions ($\mathfrak{R}$) we shall call a \gesperrt{Sturm chain}. How one can form it we shall see later; first, consequences are to be drawn from the definition.

For every value of $x$ for which none of the functions of ($\mathfrak{R}$) vanishes, each of these functions has a definite sign. We count a sign change whenever, in running through the chain from left to right, a positive term is followed by a negative one, or a negative term by a positive one. The number of sign changes (variations) counted in this way for a given value of $x$ we shall denote by $V(x)$. If a middle term $f_\nu(x)$ vanishes, then by 3. we are to count one sign change in passing from $f_{\nu-1}(x)$ to $f_{\nu+1}(x)$.

If $\alpha$, $\beta$ are two values for which none of the functions $f_\nu$ vanishes, and if we now let $x$ increase continuously from $\alpha$ to $\beta$, then a change in the number of sign changes can occur only when one of the functions of ($\mathfrak{R}$) changes its sign, and hence passes through the value zero (because of the assumed continuity).

But if this is a middle function, the number of sign changes does not alter (by 3). For in
\[
f_{\nu-1}, \; f_\nu, \; f_{\nu+1}
\]
%CONT>>


% ===== scan leaf 298, printed 273 =====
% !PRINTED 273
\weberpage{273}
\pdfbookmark[2]{§ 87. First example: spherical harmonics}{bm98}

%<<CONT
a change of sign always takes place before and after the passage
of $f_\nu$ through zero, because $f_{\nu+1}$ and $f_{\nu-1}$
have opposite signs.

But when $f(x)$ passes through zero, then by 4. a change of sign
is lost in the passage through $x$. For if $f(x)$ goes from
negative to positive values, then (cf. \S.\,32) $f'(x)$,
and hence $f_1(x)$, is positive, and the signs change thus:
\[
\begin{array}{c}
f(x), \; f_1(x) \\
- \quad + \\
+ \quad +,
\end{array}
\]
and when $f(x)$ passes from positive to negative values, then
$f'(x)$ and $f_1(x)$ are negative, so the signs are thus:
\[
\begin{array}{c}
f(x), \; f_1(x) \\
+ \quad - \\
- \quad -,
\end{array}
\]
hence in both cases a change of sign has been lost.

Now since $f_m(x)$ by 2. does not change its sign, it follows
that $V(\alpha) - V(\beta)$ is equal to the number of roots of
$f(x) = 0$ lying between $\alpha$ and $\beta$, or:

\gesperrt{The number of roots of $f(x) = 0$ between $\alpha$ and
$\beta$ is equal to the excess of the number of changes of
sign of the chain for $x = \alpha$ over the number of
changes of sign for $x = \beta$.}

If for $x = \alpha$ or $x = \beta$ one or several of the
intermediate functions of $(\mathfrak{R})$ should vanish,
then by 3. it makes no difference whether we replace this vanishing
value by a positive or by a negative one.

\begin{center}
\S.\,87.

\gesperrt{First example: spherical functions.}
\end{center}

Before we now pass on to the general methods by which
Sturm chains are to be formed, we shall consider a few examples
in which such chains present themselves through special
circumstances.


% ===== scan leaf 299, printed 274 =====
% !PRINTED 274
\weberpage{274}

We consider first the so-called \gesperrt{spherical functions}, which are used
in many ways in mathematical physics and mechanics.

By spherical functions one understands a system of integral rational
functions of increasing degree, which is defined as follows:
\[
\begin{array}{l}
P_0(x) = 1 \\
P_1(x) = x \\
P_2(x) = \frac{3}{2}(x^2 - \frac{1}{3}) \\
P_3(x) = \frac{5}{2}(x^3 - \frac{3}{5}x) \\
\hbox{\quad\quad . \quad . \quad . \quad . \quad . \quad . \quad . \quad . \quad .}
\end{array}
\]
\[
P_n(x) = \frac{1.3 \ldots (2n-1)}{1.2 \ldots n}
\left(x^n - \frac{n(n-1)}{2(2n-1)}x^{n-2}
+ \frac{n(n-1)(n-2)(n-3)}{2.4(2n-1)(2n-3)}x^{n-4} - \cdots\right),
\]
so that $P_n(x)$ is an integral rational function of degree $n$,
and, according as $n$ is even or odd, contains only the even or
only the odd powers. By use of the sign $\Pi(n)$ one can also
represent the general expression for the spherical functions thus:
\weq{1}{P_n(x) = \sum^{\mu} \frac{(-1)^\mu}{2^n}
\frac{\Pi(2n-2\mu)x^{n-2\mu}}{\Pi(\mu)\,\Pi(n-2\mu)\,\Pi(n-\mu)},}
in which the sum with respect to $\mu$ is to be extended from
$\mu = 0$ so far as $n - 2\mu$ does not become negative.

Between these functions there hold the following relations, which,
after substituting the expression (1), can be verified by simple
comparison of the coefficients of equal powers of $x$,
\weq{2}{\begin{aligned}
n P_n(x) - (2n-1) x P_{n-1}(x) + (n-1) P_{n-2}(x) &= 0 \\
P_1(x) - x P_0(x) &= 0
\end{aligned}}
\weq{3}{(1-x^2) P_n'(x) + n x P_n(x) - n P_{n-1}(x) = 0.}
From the last equation it follows for $x = \pm 1$
\[
\pm P_n(\pm 1) = P_{n-1}(\pm 1),
\]
hence
\weq{4}{P_n(1) = 1, \quad P_n(-1) = (-1)^n.}
From (2) and (3) it now follows that the functions $P_n$, ordered
downwards from an arbitrary $m$,
%CONT>>


% ===== scan leaf 300, printed 275 =====
% !PRINTED 275
\weberpage{275}
\weq{5}{P_m, \; P_{m-1}, \; P_{m-2} \ldots P_0}

have, in the interval from $-1$ to $+1$, the properties of a \gesperrt{Sturm} chain. For if $P_n$ and $P_{n-1}$ were to vanish simultaneously for any $x$, then by (2) $P_{n-2}$, and consequently $P_{n-3}$ etc. down to $P_0$, would also have to vanish, which is impossible, since $P_0=1$. Accordingly the condition \S.\,86, 2. is also satisfied for every arbitrary interval.

If $P_{n-1}=0$, then by (2) $P_n$ and $P_{n-2}$ are of opposite sign, as the condition \S.\,86, 3. requires.

From (3) it follows that $P_n(x)$ and $P_n'(x)$ never vanish simultaneously, since otherwise $P_{n-1}(x)$ would also have to vanish; and finally it follows from (3) that, if $P_n(x)$ vanishes and $-1<x<1$, then $P_n'(x)$ and $P_{n-1}(x)$ have the same sign.

If we now substitute the values $-1, +1$ for $x$ in (5), it follows from (4) that for $x=-1$ there are nothing but sign changes in (5), and for $x=+1$ no sign change at all; and from this it follows:

\gesperrt{The equation $P_m(x)=0$ has $m$ real roots between $x=-1$ and $x=+1$, and since $P_m$ is of degree $m$, these are all the roots.}

But we can draw a conclusion that goes somewhat further still. If we take an interval $\varDelta=(\alpha,\beta)$ in which two roots $\xi, \eta$ of $P_m=0$, and no more, are contained, then in the passage from $\alpha$ to $\beta$ two sign changes are lost in the chain (5). We take $\alpha$ so close to $\xi$, and $\beta$ so close to $\eta$, that $P_{m-1}(\alpha)$ has the same sign as $P_{m-1}(\xi)$, and $P_{m-1}(\beta)$ the same sign as $P_{m-1}(\eta)$. If $P_m$ passes from negative to positive as it goes through $\xi$, then it passes from positive to negative as it goes through $\eta$; hence $P_m'(\xi)$, and consequently $P_{m-1}(\xi)$ and $P_{m-1}(\alpha)$, are positive, while $P_m'(\eta)$, and consequently $P_{m-1}(\eta)$ and $P_{m-1}(\beta)$, are negative; and the reverse takes place if $P_m$ passes from positive to negative as it goes through $\xi$. One sees from this, then, that in the passage from $\alpha$ to $\beta$, in

\[
P_m, \quad P_{m-1}
\]

one sign change is lost; hence in the same passage, and consequently also in the passage from $\xi$ to $\eta$, there must, in the chain

\[
P_{m-1}, \quad P_{m-2} \ldots P_0
\]
%CONT>>


% ===== scan leaf 301, printed 276 =====
% !PRINTED 276
\weberpage{276}
\pdfbookmark[2]{§ 88. Second example}{bm99}
%<<CONT
a change of sign is likewise lost; from this we conclude the theorem:

\gesperrt{Between two successive roots of $P_m=0$ there lies one and only one root of $P_{m-1}=0$.}

\begin{center}
\S.\,88.

\gesperrt{Second example.}
\end{center}

As a second example we consider an equation which, for brevity, we shall call the \gesperrt{secular equation}, because the investigation of the secular perturbations of the planets first led to it\footnote{\gesperrt{Laplace}, Histoire de l'Acad\'emie des Sciences 1772. Among more recent works one may consult on this \gesperrt{Dziobek}, Theorie der Planetenbewegung, Leipzig 1888.}). It is indeed otherwise well known under this name, but it also occurs in many other investigations, e.g. in the determination of the principal axes of a surface of the second degree, and in the theory of small oscillations; its general analytic significance lies in the fact that it furnishes a particular transformation, the so-called orthogonal transformation, of a quadratic form into a sum of squares. We shall here take the equation as it stands, without reference to any application, and shall discuss it according to our general theorems.

Let $a_{i,k}$, where $i$ and $k$ run through the series of indices $1, 2, 3 \ldots n$, be any system of real quantities, and let

\weq{1}{a_{i,k} = a_{k,i}}

be assumed. We consider the symmetric determinant

\weq{2}{L_n(x) = \begin{vmatrix} a_{1,1}-x, & a_{1,2} & \ldots & a_{1,n} \\ a_{2,1}, & a_{2,2}-x & \ldots & a_{2,n} \\ \cdot & \cdot & \cdot & \cdot \\ a_{n,1}, & a_{n,2} & \ldots & a_{n,n}-x \end{vmatrix}}

which is an integral rational \gesperrt{function} of degree $n$ in $x$, in which $x^n$ has the coefficient $(-1)^n$. The equation $L_n(x)=0$ shall be the object of our consideration.

We form the chain of principal subdeterminants taken with alternating signs

\weq{3}{L_n(x), \; -L_{n-1}(x), \; L_{n-2}(x) \ldots, \; (-1)^{n-1}L_1(x), \; (-1)^n,}
%CONT>>


% ===== scan leaf 302, printed 277 =====
% !PRINTED 277
\weberpage{277}
%<<CONT
and we shall now show that if we add the assumption \gesperrt{that of the quantities (3) no two consecutive ones vanish at the same time}, we have before us a \gesperrt{Sturm} chain.

This is a simple consequence of the formula which we already used in \S.\,82 for a similar purpose, and which we can represent thus:

\weq{4}{L_k S_k - T_k^2 = L_{k-1} L_{k+1},}

in which $S_k$ and $T_k$ are integral rational functions of $x$.

We have only to show that the requirements \S.\,86, 1. to 4. are satisfied. Of these, however, 1. is included in the assumption, and 2. is fulfilled, since the last element is equal to $\pm 1$. From the formula (4) it follows that when $L_k=0$, $L_{k-1}$ and $L_{k+1}$ have opposite signs, so that condition 3. is fulfilled.

Only condition 4. now remains. But if we apply the formula (4) to $k=n-1$, it reads

\[
\frac{\partial L_n}{\partial a_{n,n}} \cdot \frac{\partial L_n}{\partial a_{n-1,n-1}} - \left(\frac{\partial L_n}{\partial a_{n,n-1}}\right)^2 = L_n L_{n-2},
\]

and from this it follows that, when $L_n=0$,

\[
\frac{\partial L_n}{\partial a_{n,n}}, \quad \frac{\partial L_n}{\partial a_{n-1,n-1}}
\]

have equal signs; and in the same way one can conclude that all the principal subdeterminants of $L_n$

\[
\frac{\partial L_n}{\partial a_{i,i}}
\]

have the same sign.

Now, however, the derivative of $L_n$ is (cf. \S.\,22)

\[
L_n'(x) = -\sum_{1,n}^{i} \frac{\partial L_n}{\partial a_{i,i}},
\]

and hence, for a value $x$ for which $L_n(x)$ vanishes, it has the opposite sign to $L_{n-1}$, which is precisely what requirement 4. demands. From this follows the theorem:

I.\ \gesperrt{If $\alpha, \beta$ are two real values, $\alpha<\beta$, then the number of roots of $L_n(x)$ lying between $\alpha$ and $\beta$ is equal to the excess of the number of sign changes of the chain (3) for $x=\alpha$ over the number of sign changes for $x=\beta$.}

The highest power of $x$ that occurs in $L_k(x)$ is, as already remarked above,
%CONT>>


% ===== scan leaf 303, printed 278 =====
% !PRINTED 278
\weberpage{278}
%<<CONT
\[
(-1)^k x^k,
\]

and when the absolute value of $x$ is sufficiently large, the sign of this term will decide the sign of $L_k(x)$. If, then, we take for $\alpha$ a sufficiently large negative value, and for $\beta$ a sufficiently large positive value, then in (3) there occur only changes of sign for $x=\alpha$, and only permanences of sign for $x=\beta$. Hence $n$ roots will lie between $\alpha$ and $\beta$, and from this follows the theorem:

II.\ \gesperrt{The equation $L_n(x)=0$ has only real roots.}

But these two theorems are of only limited applicability so long as we cannot free ourselves from the assumption that in the chain of the $L_k$ no two consecutive terms are to vanish at the same time. From this restriction, however, we can free the theorem by the following simple consideration.

Suppose that for some value of $x$, $L_k$ vanishes but $L_{k-1}$ does not; we can then determine quantities $a_{k+1,i}$, which do not occur in $L_k$, in such a way that $L_{k+1}$ does not vanish for this value of $x$; for $L_{k+1}$ cannot vanish identically for all $a_{k+1,i}$, because it contains the term

\[
-a_{k+1,k}^2 L_{k-1}
\]

[\S.\,23, (12)].

Accordingly, if in the chain

\weq{5}{L_n, \; -L_{n-1}, \; L_{n-2} \cdots \pm L_1, \; \mp 1}

some consecutive terms vanish for some value of $x$, we can, by altering the $a_{i,k}$, derive another sequence

\weq{6}{L_n', \; -L_{n-1}', \; L_{n-2}' \cdots \pm L_1', \; \mp 1}

in which no two consecutive terms vanish for any value $x$ (of the interval $\alpha \ldots \beta$).

At the same time the $a_{i,k}'$, on which the $L'$ depend, can be so chosen that they \gesperrt{differ from the $a_{i,k}$ by less than an arbitrarily given quantity $\omega$}.

If now the numbers $\alpha, \beta$ are so chosen that in the sequence (5) no term vanishes for $x=\alpha$ or $x=\beta$, then we can take $\omega$ so small that corresponding terms of (5) and (6) have the same sign. But by our theorem the number of roots of $L_n'=0$ between $\alpha$ and $\beta$ is determined by the signs of the sequence (6).


% ===== scan leaf 304, printed 279 =====
% !PRINTED 279
\weberpage{279}
\pdfbookmark[2]{§ 89. Sturm's functions}{bm100}
Now on the other hand $\omega$ can again be taken so small that the roots of $L_n' = 0$ differ as little as we please from those of $L_n = 0$ (\S.\,40).

A double root of $L_n=0$ may indeed pass over into two simple roots of $L_n'=0$; but since $L_n'=0$ has no imaginary roots, these are real; and the same holds when $L_n=0$ has multiple roots.

\gesperrt{Accordingly the theorems I, II retain their validity when the assumption is given up that in the chain of the $L_k$ no consecutive terms vanish; only the multiple roots must then be counted according to their multiplicity.}

\begin{center}
\S.\,89.

\gesperrt{The Sturm functions.}
\end{center}

After these particular examples we turn to the consideration of the procedure by which one obtains a Sturm chain in all cases\footnote{\gesperrt{Sturm}, M\'em. sur la r\'esolution des \'equations num\'eriques. M\'em. de l'acad\'emie de Paris. Sav. \'etrang. VI, 1835. Extract in Bull. de Ferussac XI, 1829.}. The manner of forming these functions is in principle extraordinarily simple, even though in practical execution it is usually not feasible.

We restrict ourselves here to the consideration of equations \gesperrt{without multiple roots}, or we assume that before the application $f(x)$ has been freed of every common factor with its derivative $f'(x)$.

If we then set
\weq{1}{f_1(x) = f'(x)}
the condition \S.\,86, 4. is certainly satisfied. We now proceed as if it were a matter of finding the greatest common divisor of $f(x)$ and $f_1(x)$, reversing each time the sign of the remainder; we thus form by division the equations
\weq{2}{\begin{gathered}
f = q_1 f_1 - f_2\\
f_1 = q_2 f_2 - f_3\\
\cdots\cdots\\
f_{m-2} = q_{m-1} f_{m-1} - f_m,
\end{gathered}}
%CONT>>


% ===== scan leaf 305, printed 280 =====
% !PRINTED 280
\weberpage{280}
\pdfbookmark[2]{§ 90. Hermite's solution of Sturm's problem}{bm101}
%<<CONT
in which the $q_1, q_2 \ldots q_{m-1}$, and likewise the $f_1, f_2 \ldots f_m$, are integral rational functions of $x$; the degrees of the functions $f, f_1, f_2, \ldots f_m$ \gesperrt{decrease}, and one can therefore continue the operation until $f_m$ is constant, or at least no longer vanishes in the interval under consideration. That $f_m$ cannot become zero is a consequence of the assumption that $f$ and $f_1$ are without common divisor.

That in the sequence
\weq{3}{f,\, f_1,\, f_2 \ldots f_m}
one then really has a Sturm chain follows immediately on going through the criteria of \S.\,86, 1. to 4. For if two consecutive ones among the functions (3) vanish simultaneously, then by (2) all the following ones vanish as well; but this is impossible, because $f_m$ is different from zero.

If however $f_\nu(x) = 0$, then from (2) it follows that
\[ f_{\nu-1}(x) = -f_{\nu+1}(x), \]
with which all the requirements of \S.\,86 are satisfied.

It is, as goes without saying, permitted to multiply the functions of the sequence (3) by positive, e.g. constant, factors, without their ceasing to form a \gesperrt{Sturm} chain.

For $n=2$ we can accordingly take as the Sturm functions the following:
\[
\begin{gathered}
f(x) = x^2 + a x + b\\
f_1(x) = 2x + a\\
f_2(x) = a^2 - 4b.
\end{gathered}
\]

\begin{center}
\S.\,90.

\gesperrt{Hermite's Solution of Sturm's Problem.}
\end{center}

Another way of solving Sturm's problem, which leads to simpler results, at least as regards the carrying out of the calculation in detail, has been taken by \gesperrt{Hermite}, who connects it with the law of inertia of quadratic forms and the Tschirnhausen transformation\footnote{\gesperrt{Hermite}, Remarques sur le th\'eor\`eme de M. Sturm. Comptes rendus of the Paris Academy, T.\,36 (1853).}.


% ===== scan leaf 306, printed 281 =====
% !PRINTED 281
\weberpage{281}
Let
\weq{1}{f(x) = a_0 x^n + a_1 x^{n-1} + \cdots + a_{n-1} x + a_n = 0}
be the given equation and
\weq{2}{x_1,\, x_2 \ldots x_n}
its roots, which we \gesperrt{assume to be distinct from one another}.

We use the functions defined in \S.\,68
\weq{3}{f_0(x),\, f_1(x),\, f_2(x) \ldots f_{n-1}(x)}
and, understanding by $t_0, t_1 \ldots t_{n-1}$ independent variables, we set, as there,
\weq{4}{y = t_{n-1} f_0(x) + t_{n-2} f_1(x) + \cdots + t_1 f_{n-2}(x) + t_0 f_{n-1}(x).}
Let $y_1, y_2 \ldots y_n$ be the values that $y$ takes for $x = x_1, x_2 \ldots x_n$.

Now if $\alpha$ is an arbitrary real value, we set
\weq{5}{H_\alpha = (x_1-\alpha)y_1^2 + (x_2-\alpha)y_2^2 + \cdots + (x_n-\alpha)y_n^2.}
This is a quadratic form in the $n$ variables $t$, and we wish first to determine the number of its negative and positive terms. If $x_1$ is a real root, so that $y_1$ is also real, then the square $(x_1-\alpha)y_1^2$ is positive or negative according as $x_1$ is greater or smaller than $\alpha$. But if $x_1, x_2$ form an imaginary pair, then $y_1, y_2$ are also conjugate imaginary, and
\[ (x_1-\alpha)y_1^2 + (x_2-\alpha)y_2^2 \]
decomposes into a \gesperrt{positive} and a \gesperrt{negative} square. This follows as in \S.\,79, if one sets $y_1\sqrt{x_1-\alpha} = u+iv$, $y_2\sqrt{x_2-\alpha} = u-iv$.

From this it follows that the number $N_\alpha$ of negative squares in $H_\alpha$ is equal to the number of imaginary pairs, increased by the number of real roots that are smaller than $\alpha$.

If we now take a second real number $\beta > \alpha$, form the function $H_\beta$, and denote by $N_\beta$ the number of its negative squares, then the \gesperrt{difference}
\[ N_\beta - N_\alpha \]
\gesperrt{is equal to the number of real roots between $\alpha$ and $\beta$.}

One thus obtains a means of determining this number, that is, of solving Sturm's problem, if one expresses the coefficients of the function $H_\alpha$ as functions of the coefficients of $f(x)$ and of $\alpha$, and then investigates the number $N_\alpha$,
% [sic] The printed page ends the complete sentence before § 91 with a comma.


% ===== scan leaf 307, printed 282 =====
% !PRINTED 282
\weberpage{282}
\pdfbookmark[2]{§ 91. Determination of Hermite's form H}{bm102}
\begin{center}
\S.\,91.

\gesperrt{Determination of the Hermite form $H$.}
\end{center}

To determine the Hermite form $H$ we can apply the methods that we became acquainted with in Section VI. In \S.\,72 we derived the formula:
\weq{1}{y f_s = E_{0,s} f_0 + E_{1,s} f_1 + \cdots + E_{n-1,s} f_{n-1}.}
Making use of the relations
\[ x f_0 = f_1 - a_1,\quad x f_1 = f_2 - a_2, \ldots x f_{n-1} = -a_n, \]
it follows from this that
% [sic] Print has superscript $s$ here; equations (1) and (4) have $f_s$.
\weq{2}{\begin{aligned}
x y f^s &= E_{0,s} f_1 + E_{1,s} f_2 + \cdots + E_{n-2,s} f_{n-1}\\
&\quad - a_1 E_{0,s} - a_2 E_{1,s} - \cdots - a_n E_{n-1,s}.
\end{aligned}}
We therefore determine a further series of functions $E_{-1,s}$ by the equation
\weq{3}{a_0 E_{-1,s} + a_1 E_{0,s} + a_2 E_{1,s} + \cdots + a_n E_{n-1,s} = 0}
and then obtain from (1) and (2)
\weq{4}{\begin{aligned}
(x-\alpha) y f_s &= (E_{-1,s} - \alpha E_{0,s}) f_0 + (E_{0,s} - \alpha E_{1,s}) f_1 + \cdots\\
&\quad + (E_{n-2,s} - \alpha E_{n-1,s}) f_{n-1}.
\end{aligned}}
The function $E_{-1,s}$ can, however, by means of the relation [\S.\,72. (7)]
\[ a_0 t_{n+s} + a_1 t_{n+s-1} + \cdots + a_n t_s = 0 \]
be expressed in entirely the same way as the remaining $E$, so that we obtain the following system of formulas [\S.\,72, (9)]:
\weq{5}{\begin{aligned}
E_{-1,s} &= a_0 t_{n+s} + a_1 t_{n+s-1} + \cdots + a_s t_n\\
&= -a_{s+1} t_{n-1} \cdots - a_n t_s\\
E_{0,s} &= a_0 t_{n+s-1} + a_1 t_{n+s-2} + \cdots + a_s t_{n-1}\\
&= -a_{s+1} t_{n-2} \cdots - a_n t_{s-1}\\
E_{1,s} &= a_0 t_{n+s-2} + a_1 t_{n+s-3} + \cdots + a_s t_{n-2}\\
&= -a_{s+1} t_{n-3} \cdots - a_n t_{s-2}\\
&\;\vdots\\
E_{s-1,s} &= a_0 t_n + a_1 t_{n-1} + \cdots + a_s t_{n-s}\\
&= -a_{s+1} t_{n-s-1} \cdots - a_n t_0\\
E_{s,s} &= a_0 t_{n-1} + a_1 t_{n-2} + \cdots + a_s t_{n-s-1}\\
E_{s+1,s} &= a_0 t_{n-2} + a_1 t_{n-3} + \cdots + a_s t_{n-s-2}\\
&\;\vdots\\
% [sic] Print has a multiplication dot after $a_1t_{s-1}$ before the following plus.
E_{n-1,s} &= a_0 t_s + a_1 t_{s-1} \cdot + \cdots + a_s t_0.
\end{aligned}}


% ===== scan leaf 308, printed 283 =====
% !PRINTED 283
\weberpage{283}
\pdfbookmark[2]{§ 92. The determinant of Hermite's form}{bm103}

Let us introduce a second system of variables $\tau$ and set
\[
z = \tau_{n-1} f_0 + \tau_{n-2} f_1 + \cdots + \tau_0 f_{n-1},
\]
then from (4), making use of the formulae [\S.\,68, (6)]
\[
S(f_s) = (n-s)\,a_s
\]
we obtain
\weq{6}{\begin{aligned}
S[(x-\alpha)\,y\,z] &= n\,a_0 \sum_{0,\,n-1}^{s} (E_{-1,s} - \alpha\,E_{0,s})\,\tau_{n-s-1} \\
&\quad + (n-1)\,a_1 \sum_{0,\,n-1}^{s} (E_{0,s} - \alpha\,E_{1,s})\,\tau_{n-s-1} \\
&\quad \cdots\cdots\cdots\cdots\cdots\cdots\cdots \\
&\quad + a_{n-1} \sum_{0,\,n-1}^{s} (E_{n-2,s} - \alpha\,E_{n-1,s})\,\tau_{n-s-1},
\end{aligned}}
and this passes directly into the function $H$ when we set $\tau = t$.

If we take as an example the case $n = 3$, we obtain
\[
\begin{aligned}
E_{-1,0} &= -a_1 t_2 - a_2 t_1 - a_3 t_0, & E_{0,0} &= a_0 t_2, \\
E_{-1,1} &= -a_2 t_2 - a_3 t_1, & E_{0,1} &= -a_2 t_1 - a_3 t_0, \\
E_{-1,2} &= -a_3 t_2, & E_{0,2} &= -a_3 t_1, \\
E_{1,0} &= a_0 t_1, & E_{2,0} &= a_0 t_0, \\
E_{1,1} &= a_0 t_2 + a_1 t_1, & E_{2,1} &= a_0 t_1 + a_1 t_0, \\
E_{1,2} &= -a_3 t_0, & E_{2,2} &= a_0 t_2 + a_1 t_1 + a_2 t_0,
\end{aligned}
\]
and from these the coefficients of the form $H$ can easily be computed according to (6):
\[
\begin{aligned}
H_{2,2} &= -a_0 (a_1 + 3 a_0 \alpha) \\
H_{1,1} &= -3 a_0 a_3 - a_1 a_2 - 2(a_1^2 - a_0 a_2)\alpha \\
H_{0,0} &= -a_2 a_3 + (2 a_1 a_3 - a_2^2)\alpha \\
H_{1,0} &= -2 a_1 a_3 + (3 a_0 a_3 - a_1 a_2)\alpha \\
H_{2,1} &= -2 a_0 (a_2 + a_1 \alpha) \\
H_{0,2} &= -a_0 (3 a_3 + a_2 \alpha).
\end{aligned}
\]

\begin{center}
\S.\,92.

\gesperrt{The Determinant of the Hermite Form.}
\end{center}

For the question of the number of negative squares of the form $H$, a knowledge of its determinant is of importance. This determinant can be computed in general in the following way.


% ===== scan leaf 309, printed 284 =====
% !PRINTED 284
\weberpage{284}

If we set
\[
H = \sum_{0,\,n-1}^{i} \sum^{k} H_{i,k}\, t_i\, t_k,
\]
then the coefficients $H_{i,k}$, whose determinant is to be formed, are given by the formula \S.\,91, (6). Here $H_{i,k}$ is the coefficient of $t_i \tau_k$ in that formula, hence
\[
H_{i,k} = S(x-\alpha)\,f_{n-i-1}\,f_{n-k-1}.
\]
By the multiplication theorem, however, the determinant of these $n^2$ quantities can be put into the form:
\[
\varDelta =
\begin{vmatrix}
(x_1-\alpha)f_0(x_1), & (x_1-\alpha)f_1(x_1), & \dots & (x_1-\alpha)f_{n-1}(x_1) \\
\cdots & \cdots & \cdots & \cdots \\
(x_n-\alpha)f_0(x_n), & (x_n-\alpha)f_1(x_n), & \dots & (x_n-\alpha)f_{n-1}(x_n)
\end{vmatrix}
\times
\begin{vmatrix}
f_0(x_1), & f_1(x_1), & \dots & f_{n-1}(x_1) \\
\cdots & \cdots & \cdots & \cdots \\
f_0(x_n), & f_1(x_n), & \dots & f_{n-1}(x_n)
\end{vmatrix},
\]
or also, since
\[
a_0 (x_1-\alpha)(x_2-\alpha)\cdots(x_n-\alpha) = (-1)^n f(\alpha)
\]
into the form
\[
\frac{(-1)^n f(\alpha)}{a_0}
\begin{vmatrix}
f_0(x_1), & f_1(x_1) & \dots & f_{n-1}(x_1) \\
\cdots & \cdots & \cdots & \cdots \\
f_0(x_n), & f_1(x_n) & \dots & f_{n-1}(x_n)
\end{vmatrix}^2.
\]
The determinant of the $f_k(x_i)$ still occurring here is the product of the two following:
\[
\begin{vmatrix}
a_0, & 0 & \dots & 0 \\
a_1, & a_0 & \dots & 0 \\
\cdots & \cdots & \cdots & \cdots \\
a_{n-1}, & a_{n-2} & \dots & a_0
\end{vmatrix}
\begin{vmatrix}
1, & x_1, & x_1^2 & \dots & x_1^{n-1} \\
1, & x_2, & x_2^2 & \dots & x_2^{n-1} \\
\cdots & \cdots & \cdots & \cdots & \cdots \\
1, & x_n, & x_n^2 & \dots & x_n^{n-1}
\end{vmatrix}.
\]
But the square of this product, if $D$ denotes the discriminant of the function $f(x)$ (\S.\,46), is equal to $a_0^2 D$, so that for $\varDelta$ we obtain
\weq{1}{\varDelta = (-1)^n a_0 f(\alpha)\, D.} % source type in $(-1)^n$ is damaged; reading resolved by the immediately preceding identity
Since by hypothesis $D$ is not to vanish, $\varDelta$ will therefore vanish only when one of the roots of $f(x) = 0$ is put for $\alpha$, and this too is to be excluded.

If we denote a chain of principal subdeterminants of $\varDelta$, beginning with the lowest, by
\[
\varDelta_1(\alpha),\ \varDelta_2(\alpha), \dots, \varDelta_n(\alpha) = \varDelta,
\]
%CONT>>


% ===== scan leaf 310, printed 285 =====
% !PRINTED 285
\weberpage{285}
\pdfbookmark[2]{§ 93. Outline of the theory of characteristics}{bm104}

%<<CONT
then by \S.\,83, V. the number of sign changes in
\weq{2}{1,\ \varDelta_1(\alpha), \varDelta_2(\alpha) \dots \varDelta_n(\alpha)}
is equal to the number $N_\alpha$ of negative squares in $H$, and if we therefore count the corresponding number $N_\beta$ in
\[
1,\ \varDelta_1(\beta), \varDelta_2(\beta) \dots \varDelta_n(\beta)
\]
then $N_\beta - N_\alpha$ is the number of roots of $f(x)$ lying between $\alpha$ and $\beta$.

For the cubic equation the expressions follow from the end of the preceding section. We form the chain
\[
a_0,\quad \frac{1}{a_0} H_{2,2},\quad \frac{1}{a_0}\left(H_{2,2} H_{1,1} - H_{1,2}^2\right),\quad \frac{1}{a_0} \varDelta,
\]
which evidently has the same number of sign changes as (2), and thus obtain the four functions
\[
\begin{aligned}
&a_0,\qquad -3 a_0 \alpha - a_1,\\
&2\alpha^2 a_0(a_1^2 - 3a_0 a_2) + \alpha(2a_1^3 + 9a_0^2a_3 - 7a_0a_1a_2)\\
&\quad + 3a_0a_1a_3 + a_1^2a_2 - 4a_0a_2^2,\qquad -f(\alpha)\,D.
\end{aligned}
\]
If as a test we take $\alpha = -\infty$, $\beta = +\infty$, then, when moreover $a_0 = 1$ is set, we obtain the determination of signs
\[
\begin{aligned}
1, &\quad +1, \quad (a_1^2 - 3a_2), \quad D,\\
1, &\quad -1, \quad (a_1^2 - 3a_2), \quad -D.
\end{aligned}
\]
In the enumeration one must observe that when $D$ is positive, $a_1^2 - 3a_2$ cannot be negative [\S.\,47, (8)].

\begin{center}
\S.\,93.

\gesperrt{Outlines of the theory of characteristics}\footnote{Monatsberichte der Berliner Akademie, March and August 1869, February 1873, February 1878.}).
\end{center}

The theorems of \gesperrt{Sturm} are generalised in an extraordinary manner by the theory of characteristics of \gesperrt{Kronecker}, which pursues the same aim as \gesperrt{Sturm}'s theorem for systems of equations in arbitrarily many variables. We confine ourselves here to the consideration of the simplest case, which we intend to apply to the enclosure of the complex roots of an equation; moreover, we here avail ourselves un-
%CONT>>


% ===== scan leaf 311, printed 286 =====
% !PRINTED 286
\weberpage{286}

%<<CONT
restricted use of geometric intuition and of the notation of the differential calculus.

We first consider two real functions $\varphi(x,y)$, $\psi(x,y)$, and interpret the real variables $x, y$ as rectangular coordinates in a plane. The equations
\weq{1}{\varphi(x,y) = 0,\qquad \psi(x,y) = 0}
then represent two curves, which we shall call briefly the curve $\varphi$ and the curve $\psi$. We assume first that these curves are closed and do not extend to infinity. We assume moreover that $\varphi$ is \gesperrt{negative} in the \gesperrt{interior} of the curve $\varphi$ and \gesperrt{positive} in the \gesperrt{exterior}; likewise that $\psi$ is negative in the interior of $\psi$ and positive in the exterior.

If, as Fig.\ 6 shows, we denote by $\vartheta$ the angle which the normal $n$ at any point of $\varphi$, drawn in the direction in which $\varphi$ increases, that is, outward, makes with the direction of the positive $x$-axis, then, if $\varphi'(x)$ and $\varphi'(y)$ are the partial derivatives of $\varphi$, and the square root is taken positive,
\[
\cos\vartheta = \frac{\varphi'(x)}{\sqrt{\varphi'(x)^2 + \varphi'(y)^2}},\qquad \sin\vartheta = \frac{\varphi'(y)}{\sqrt{\varphi'(x)^2 + \varphi'(y)^2}}.
\]
If we draw the tangent $t$ so that it lies to the normal $n$ just described as the positive $y$-axis lies to the positive $x$-axis, and denote by $\xi$ the angle which it makes with the positive $x$-axis, then
\[
\cos\xi = -\frac{\varphi'(y)}{\sqrt{\varphi'(x)^2 + \varphi'(y)^2}},\qquad \sin\xi = \frac{\varphi'(x)}{\sqrt{\varphi'(x)^2 + \varphi'(y)^2}}.
\]
If we designate the advance along $\varphi$ in the direction $t$ as the \gesperrt{positive} one, and if $dx, dy$ are the projections of this advance, then $dx$ and $dy$ are proportional to, and agree in sign with, $-\varphi'(y)$, $\varphi'(x)$; and if $\Phi$ is any other arbitrary function, then
\[
d\Phi = \Phi'(x)\,dx + \Phi'(y)\,dy
\]
agrees in sign with the \gesperrt{functional determinant}
\weq{2}{[\varphi, \Phi] = \varphi'(x)\Phi'(y) - \varphi'(y)\Phi'(x).}
Under the usual assumption about the direction of the coordinates, the positive direction of advance is the one for which the interior of the surface lies to the left.


% ===== scan leaf 312, printed 287 =====
% !PRINTED 287
\weberpage{287}
\pdfbookmark[2]{§ 94. Characteristic of a system of three functions}{bm105}
Now if, at one of the points of intersection of the curves $\varphi$, $\psi$,
the positive direction of advance on $\varphi$ leads into the interior
of $\psi$, then $d\psi$, and hence also the functional determinant
$[\varphi, \psi]$, is negative, and in the opposite case positive (Fig.~7).
We call a point of intersection of $\varphi$ and $\psi$ a
\gesperrt{point of exit} $A(\varphi, \psi)$ or a point of entry
$E(\varphi,\psi)$, according as

\weberfigureimage{weber-b1-fig06-07.png}{0.94}

the positive direction of advance of $\varphi$ leads from the interior of $\psi$
to the exterior or from the exterior to the interior, and thus we have the theorem:
\weq{4}{
\begin{array}{lll}
\text{at an } A(\varphi, \psi) & \text{is} & [\varphi, \psi] > 0 \\
\text{,,}\ \ \text{,,} & E(\varphi, \psi) & \text{,,}\ \ [\varphi, \psi] < 0.
\end{array}
}
The number of the $A(\varphi, \psi)$ is just as great as that of the $E(\varphi, \psi)$;
and if we interchange $\varphi$ with $\psi$, the $A(\varphi, \psi)$, $E(\varphi, \psi)$
pass into $E(\psi, \varphi)$ and $A(\psi, \varphi)$.

The two functions $\varphi$, $\psi$ determine in a unique manner
a region of the plane, which is characterised by the fact that in it
the product $\varphi\psi$ is negative. We shall call it the \gesperrt{inner region}
$(\varphi, \psi)$ or $(\psi, \varphi)$. In our Fig.~7 it is the shaded
area.

\begin{center}
\S.\,94.

\gesperrt{Characteristic of a system of three functions.}
\end{center}

We now adjoin to the two functions $\varphi$, $\psi$ a third,
$f(x, y)$. The curve represented by the equation $f = 0$,
or the curve $f$, shall likewise be closed in the finite part of the plane
%CONT>>


% ===== scan leaf 313, printed 288 =====
% !PRINTED 288
\weberpage{288}
%<<CONT
and we shall further assume that it passes through none of the
points of intersection of $\varphi$ and $\psi$.

We give these three functions a definite cyclic order
\[
f,\ \varphi,\ \psi,
\]
so that $f$ is again to follow $\psi$, and we associate with this
order a definite direction of advance on each of the three curves,
say the positive one. We shall stipulate that with the reversed
order
\[
f,\ \psi,\ \varphi
\]
the opposite, that is the negative, direction of advance shall be
associated on each of the three curves.

We now traverse, in the first order, the curve $f$ in the
positive sense, and attend to the points of intersection of $f$
and $\varphi$. Such a point shall be called an exit point $A(f; \varphi, \psi)$

\weberfigureimage{weber-b1-fig08.png}{0.48}

if the positive direction of $f$ leads from the interior of the
inner region $(\varphi, \psi)$ into the exterior, and an entry point
$E(f; \varphi, \psi)$ if it leads from the exterior into the interior.

In Fig.~8 the points $a$ are the exit points, the point
$e$ an entry point.

The total number of entry and exit points agrees with the
number of points of intersection of $f$ and $\varphi$, and is therefore, since
both curves are closed, an \gesperrt{even number}. If $a$ is the number
of points $A(f; \varphi, \psi)$, $e$ the number of points $E(f; \varphi, \psi)$, then
$e - a$ is likewise an even number, and
\[
\tfrac{1}{2}(e - a) = k
\]
is an integer, which after \gesperrt{Kronecker} is called the \gesperrt{characteristic
of the system of functions} $f$, $\varphi$, $\psi$. (In the case of Fig.~8
it is equal to $-1$).

If we permute the functions $f$, $\varphi$, $\psi$ cyclically, we
obtain three determinations of the characteristic, and if
we reverse the order and thereby, in accordance with the convention
adopted, also replace the positive directions by the negative
ones, three more. These determinations are, if we now use
the symbols $A(f; \varphi, \psi)$, $E(f; \varphi, \psi)$ at the same time as a designation
for the numbers of the points concerned,
%CONT>>


% ===== scan leaf 314, printed 289 =====
% !PRINTED 289
\weberpage{289}
%<<CONT
\[
\begin{array}{ll}
1. & \tfrac{1}{2}\big[E(f; \varphi, \psi) - A(f; \varphi, \psi)\big], \\
2. & \tfrac{1}{2}\big[E(\varphi; \psi, f) - A(\varphi; \psi, f)\big], \\
3. & \tfrac{1}{2}\big[E(\psi; f, \varphi) - A(\psi; f, \varphi)\big], \\
4. & \tfrac{1}{2}\big[E(f; \psi, \varphi) - A(f; \psi, \varphi)\big], \\
5. & \tfrac{1}{2}\big[E(\psi; \varphi, f) - A(\psi; \varphi, f)\big], \\
6. & \tfrac{1}{2}\big[E(\varphi; f, \psi) - A(\varphi; f, \psi)\big],
\end{array}
\]
\gesperrt{and there now holds the fundamental theorem that these six determinations
yield the same number.}

To prove it, we traverse the curve $f$ in the positive
sense, and attend to all the intersection points with $\varphi$
and $\psi$. The path along the curve $f$ must in doing so enter
the inner region $(\varphi, \psi)$ just as often as it steps out of it,
because it must return again to its starting point.
But the curve $f$ enters at the points $E(f; \varphi, \psi)$
and $A(f; \psi, \varphi)$ and steps out at the points $A(f; \varphi, \psi)$ and
$E(f; \psi, \varphi)$. Hence
\[
E(f; \varphi, \psi) + A(f; \psi, \varphi) = A(f; \varphi, \psi) + E(f; \psi, \varphi),
\]
by which the agreement of 1. and 4. is established.

\weberfigureimage{weber-b1-fig09.png}{0.42}

If now, secondly, we traverse the curve $\varphi$ in the negative sense,
and attend to its intersection points with $f$, it turns out
that every point $E(f; \varphi, \psi)$ is at the same time a point $E(\varphi; f, \psi)$,
and every point $A(f; \varphi, \psi)$ a point $A(\varphi; f, \psi)$. For in
Fig.~9, in which the point $\alpha$ represents any one of the intersection points
of $f$ and $\varphi$, this point is an $E(f; \varphi, \psi)$ and
an $E(\varphi; f, \psi)$ when $\psi$ is positive at $\alpha$, and an $A(f; \varphi, \psi)$ and
an $A(\varphi; f, \psi)$ when $\psi$ is negative at $\alpha$.

From this follows the agreement of 1. with 6. There follows
further, by a renewed application of the first argument, the
agreement of 6. with 2., and then by application of the
second argument that of 2. with 5., and finally by the first that
of 5. with 3., \gesperrt{hence the agreement of all six expressions.}

One is therefore justified in designating the number so determined
simply as the \gesperrt{characteristic of the system of functions} $(f, \varphi, \psi)$.
%CONT>>


% ===== scan leaf 315, printed 290 =====
% !PRINTED 290
\weberpage{290}
\pdfbookmark[2]{§ 95. Relation of the characteristics to the points of intersection}{bm106}
%<<CONT
\begin{center}
\S.\,95.

\gesperrt{Relation of the characteristic to the
points of intersection.}
\end{center}

One must now observe that the designation of the points of intersection
of two curves as points of entry or of exit in \S.\,93
is essentially different from that in \S.\,94. There only
the relation of the points to the two intersecting
curves came into consideration, while in \S.\,94 the relation to
a third curve was also in question; this is completely expressed in the chosen
designations $E(\varphi, \psi)$, $A(\varphi, \psi)$ and

\weberfigureimage{weber-b1-fig10.png}{0.48}

$E(f; \varphi, \psi)$, $A(f; \varphi, \psi)$. But now we must
investigate more precisely how the two modes of designation
are related to one another.

For the sake of clarity we choose a somewhat simpler
figure than above, which at the same time illustrates all possible
circumstances (Fig.~10).

At a point $E(\varphi, \psi)$ the $\varphi$-curve enters the region
of the curve $\psi$, i.\,e.\ $\psi$ passes from positive to negative
values. If at the same time $f$ is positive, then this point is an
$E(\varphi; \psi, f)$, like the point 1 in our figure; but if $f$ is negative,
then it is an $A(\varphi; \psi, f)$, like the point 3. Thus we have:

\[
\begin{array}{ll}
\text{A point } E(\varphi, \psi) \text{ is an} & \\
\quad E(\varphi; \psi, f), \text{ if } f > 0, & \text{(point 1 of Fig. 10)} \\
\quad A(\varphi; \psi, f), \text{ if } f < 0, & \text{(point 3 of Fig. 10).} \\
\text{A point } A(\varphi, \psi) \text{ is an} & \\
\quad E(\varphi; \psi, f), \text{ if } f < 0, & \text{(point 2 of Fig. 10)} \\
\quad A(\varphi; \psi, f), \text{ if } f > 0, & \text{(point 4 of Fig. 10).}
\end{array}
\]

Now at the points $E(\varphi, \psi)$ the functional determinant
$[\varphi, \psi]$ is negative, and at $A(\varphi, \psi)$ positive (\S.\,93), and we can
therefore summarise what has been proved as follows:

$E(\varphi; \psi, f)$ is the number of points of intersection of $\varphi$, $\psi$ at
which $[\varphi, \psi]\, f < 0$,

$A(\varphi; \psi, f)$ the number of points of intersection of $\varphi$, $\psi$ at
which $[\varphi, \psi]\, f > 0$;
%CONT>>


% ===== scan leaf 316, printed 291 =====
% !PRINTED 291
\weberpage{291}
%<<CONT
thus the characteristic is equal to half the excess of the first number over the second.

To this theorem we can give the following expression, in which it appears entirely independent of relations of magnitude and depends only on the relations of position of the curves and of their points of intersection.

If we distinguish the points of intersection of $\varphi$, $\psi$, according as they are exit points $A(\varphi,\psi)$ or entry points $E(\varphi,\psi)$, by the signs $\alpha$ and $\varepsilon$, and likewise, according as they are exterior or interior points with respect to the curve $f$, by $\alpha$ and $\varepsilon$, and then give the points $\alpha,\varepsilon$ and $\varepsilon,\alpha$ the \gesperrt{character $+1$, and the points $\alpha,\alpha$ and $\varepsilon,\varepsilon$ the character $-1$}\textsuperscript{1)}, \gesperrt{then the characteristic of the system of functions $(f,\varphi,\psi)$ is equal to half the sum of the characters of all the points of intersection of $\varphi,\psi$.}

The theorems and definitions that we have given here remain valid unchanged even when the curves considered consist of several closed circuits; double points may also be present on the curves; it is only necessary that, for every direction of advance on one of the curves, it be completely determined whether it is to be regarded as positive or as negative, i.\,e. every piece of a curve $f$ or $\varphi$ or $\psi$ must separate portions of surface from one another in which the functions $f$ or $\varphi$ or $\psi$ have opposite signs.

Only those cases are to be excluded in which one of the curves passes through a double point of the other, or in which the curves touch one another, or in which the three curves $f$, $\varphi$, $\psi$ pass through \gesperrt{one} point; for in these cases the determination of a point as an entry point or exit point would become doubtful.

If the curves extend to infinity, then, in order to make our theorems applicable, one must close them off by arbitrarily added pieces of curve, as we shall see in the next section.

\noindent\rule{4cm}{0.4pt}

\noindent\textsuperscript{1)} \gesperrt{Kronecker} understands by the character of a point something different.


% ===== scan leaf 317, printed 292 =====
% !PRINTED 292
\weberpage{292}
\pdfbookmark[2]{§ 96. Application of the characteristics to the localization of the complex roots of an equation}{bm107}
\begin{center}
§. 96.

\gesperrt{Application of the characteristics to the bounding\\
of the complex roots of an equation.}
\end{center}

Now let
\[
z = x + y i
\]
be a complex variable and
\weq{1}{F(z) = \varphi(x,y) + i\,\psi(x,y)}
an integral rational function of $z$ with real or complex coefficients, and in it let $\varphi(x,y)$, $\psi(x,y)$ be real functions of the real variables $x, y$. We assume that $F(z)$ and $F'(z)$ do not vanish simultaneously. $F(z)$ vanishes only at the intersection points of the two curves $\varphi$ and $\psi$. Differentiation of (1) yields
\[
F'(z) = \varphi'(x) + i\,\psi'(x) = -\,i\,\varphi'(y) + \psi'(y),
\]
hence
\weq{2}{\varphi'(x) = \psi'(y), \quad \varphi'(y) = -\,\psi'(x),}
and consequently
\weq{3}{
\begin{aligned}
[\varphi,\psi] &= \varphi'(x)\,\psi'(y) - \varphi'(y)\,\psi'(x)\\
&= \varphi'(x)^2 + \varphi'(y)^2 = \psi'(x)^2 + \psi'(y)^2.
\end{aligned}
}

From this it follows that the functional determinant $[\varphi,\psi]$ never becomes negative, and vanishes only where $\varphi'(x)$ and $\varphi'(y)$, and hence also $\psi'(x)$ and $\psi'(y)$, vanish simultaneously.

But this never occurs at an intersection point of $\varphi$ and $\psi$, since otherwise $F(z)$ and $F'(z)$ would vanish there simultaneously. The positive direction of advance is completely determined, in every part of the curve $\varphi$ or $\psi$, by the sign of the function in the adjoining portions of the surface; even any double points, at which $\varphi$, $\varphi'(x)$ and $\varphi'(y)$ vanish simultaneously, form no exception to this.

Since the functional determinant $[\varphi,\psi]$ is positive at all intersection points of the curves $\varphi, \psi$, all these points are exit points $A(\varphi,\psi)$, and from this it follows that the curves $\varphi$ and $\psi$ cannot be closed, but must extend to infinity, since otherwise an exit point would necessarily have to be followed by an entry point. For the rest, here too the curves $\varphi, \psi$ determine an inner region in which the product $\varphi\psi$ is negative.


% ===== scan leaf 318, printed 293 =====
% !PRINTED 293
\weberpage{293}
We now add to the two functions $\varphi, \psi$ a third function $f$, which, set equal to zero, represents a \gesperrt{closed} curve, and we determine the characteristic of the system of functions entirely in the earlier manner, by proceeding along the curve $f$
\weq{1}{k = \tfrac{1}{2}\big[E(f;\varphi,\psi) - A(f;\varphi,\psi)\big] = \tfrac{1}{2}\big[E(f;\psi,\varphi) - A(f;\psi,\varphi)\big].}
Thus, for example, in Fig. 11 the characteristic is $k = 2$.

\weberfigureimage{weber-b1-fig11-12.png}{0.96}

In order to be able to apply our theorems, we must make the curves $\varphi, \psi$ closed in some way by arbitrarily constructed connections, all of which, however, are to run outside $f$. The characteristic of this system is then likewise determined by (1). As an example we complete Fig. 11 by Fig. 12, in which these closing connections are drawn. The points of intersection added in this way may be partly exit points $A(\varphi,\psi)$, partly entry points $E(\varphi,\psi)$. We shall denote their numbers by
\[
A'(\varphi,\psi), \quad E'(\varphi,\psi)
\]
while $A(\varphi,\psi)$ denotes the number of the points of intersection originally present. We must now always have
\weq{2}{A(\varphi,\psi) + A'(\varphi,\psi) = E'(\varphi,\psi)}
We divide the points $A(\varphi,\psi)$ into two groups $A_\alpha(\varphi,\psi)$, $A_\varepsilon(\varphi,\psi)$, of which the first are to lie outside, the others inside $f$. Then we have the following characters of the points of intersection
%CONT>>


% ===== scan leaf 319, printed 294 =====
% !PRINTED 294
\weberpage{294}
\pdfbookmark[2]{§ 97. Determination of the characteristic}{bm108}
%<<CONT
\[
\begin{aligned}
A_\alpha(\varphi,\psi) &\quad -1,\\
A_\varepsilon(\varphi,\psi) &\quad +1,\\
A'(\varphi,\psi) &\quad -1,\\
E'(\varphi,\psi) &\quad +1.
\end{aligned}
\]

Twice the value of the characteristic is then, if for the sake of simplicity we omit the designation $(\varphi,\psi)$ (§.\,95),
\weq{3}{-A_\alpha + A_\varepsilon - A' + E' = 2k;}
adding to this the equation (2) yields
\weq{4}{k = A_\varepsilon,}
i.\,e.\ \gesperrt{the characteristic is equal to the number of the intersection points of $\varphi, \psi$ lying in the interior of $f$.}

With this we have thus found a means of determining, from the characteristic, the number of roots of the equation
\[
F(z) = 0,
\]
whose geometric images lie in the interior of an arbitrarily given curve. The outer connections added at will play no further part here at all, and can be suppressed entirely.

\begin{center}
§. 97.

Determination of the Characteristic.
\end{center}

To make the theorem proved in the foregoing applicable, we must still show how we can actually determine the characteristics. For this purpose, let us imagine the curve $f$ divided into segments in such a way that in any one of them there lies only one intersection point with the curve $\varphi$ or one intersection point with the curve $\psi$; and let us denote by $a$ (Fig. 13) the beginning, by $b$ the end of such a segment, and by $\xi$ the intersection point of $f$ with $\varphi$ lying in it. Then the point $\xi$ will be a point $E(f;\varphi,\psi)$ if the product $\varphi\psi$ is positive at $a$, and a point $A(f;\varphi,\psi)$ if $\varphi\psi$ is negative at $a$; and with this the

\weberfigureimage{weber-b1-fig13.png}{0.42}
%CONT>>


% ===== scan leaf 320, printed 295 =====
% !PRINTED 295
\weberpage{295}
\pdfbookmark[2]{§ 98. Gauss's first proof of the fundamental theorem of algebra}{bm109}
%<<CONT
the possibility of determining the characteristic completely from
the signs of the functions $\varphi, \psi$ at the points of division of the
segments. But how we are to divide the curve $f$ into such subsegments
is taught us by the \gesperrt{Sturm theorem}. For this purpose we imagine
$x$ and $y$ represented as functions of a variable $t$, which runs along the curve $f$ through all
real values from $-\infty$ to $+\infty$
(or also only through the values of a
finite interval). Thereby $\varphi$ and $\psi$ also pass into functions of the variable $t$, and one
has only to investigate the position of their roots according to the \gesperrt{Sturm theorem}.

\weberfigureimage{weber-b1-fig14.png}{0.44}

If, for example, as in Fig. 14, the curve $f$ is a circle with the
equation
\[
(x - \alpha)^2 + (y - \beta)^2 = \varrho^2,
\]
then we can set:
\[
x - \alpha = \varrho\, \frac{1 - t^2}{1 + t^2},
\]
\[
y - \beta = \varrho\, \frac{2t}{1 + t^2},
\]
and $t$ simply runs through all values from $-\infty$ to $+\infty$, while
the point $x, y$ moves over the whole circle in the positive sense.
$\varphi$ and $\psi$, multiplied by a suitable power of
$1 + t^2$, pass into integral rational functions of $t$,
which permit the application of the \gesperrt{Sturm theorem}.

\begin{center}
\S.\,98.

\gesperrt{Gauss' first proof of the fundamental theorem of algebra.}
\end{center}

\gesperrt{Gauss} has given three different proofs of the principal theorem
of algebra, that an equation of degree $n$ has $n$ roots.
The first of these, which is found in his doctoral dissertation, and which he later carried out still further and
presented differently, rests, though in a different garb,
%CONT>>


% ===== scan leaf 321, printed 296 =====
% !PRINTED 296
\weberpage{296}
%<<CONT
on the idea that underlies the concept of the characteristic\,\footnote{\gesperrt{Gauss}, Demonstratio nova theorematis omnem functionem algebraicam rationalem integram unius variabilis in factores reales primi vel
secundi gradus resolvi posse (1799). Beiträge zur Theorie der algebraischen
Gleichungen (1849). \gesperrt{Kronecker}, Monatsberichte der Berliner Akademie,
21. Februar 1878.}. We therefore fittingly close this section with an
exposition of this proof.

Let, then, $z = x + yi$ and
\weq{1}{F(z) = \varphi(x, y) + i\,\psi(x, y)}
be an integral rational function of $z$ of degree $n$ with real
or complex coefficients, in which $z^n$ has the coefficient 1,
thus
\weq{2}{F(z) = z^n + a_1 z^{n-1} + a_2 z^{n-2} + \cdots + a_n.}

We assume again that $F(z)$ and $F'(z)$ have no common
divisor, and the proof that $F(z)$ vanishes for $n$
values of $z$ now rests on the fact that, if we
choose for $f$ a circle of sufficiently large radius, the
characteristic of the system of functions $f, \varphi, \psi$ can be determined
directly and is found to be equal to $n$. It then follows by \S.\,96
that in the circle so chosen there lie $n$ points $z$ at which
$F(z)$ vanishes.

We make use of polar coordinates, and set, since
$a_1, a_2, a_3 \dots a_n$ may also be complex,
\[
a_1 = p_1 e^{iq_1},\quad a_2 = p_2 e^{iq_2}, \dots a_n = p_n e^{iq_n}
\]
\[
z = R e^{i\vartheta},
\]
in which $p_1, p_2, \dots p_n, R$ are positive and the angles $q_1, q_2, \dots q_n, \vartheta$
can be chosen in any interval of extent $2\pi$.
Hereby we obtain
\weq{3}{
\begin{aligned}
\varphi(x, y) &= R^n \cos n\vartheta + p_1 R^{n-1} \cos[(n-1)\vartheta + q_1] \\
&\quad + \cdots + p_n \cos q_n \\[4pt]
\psi(x, y) &= R^n \sin n\vartheta + p_1 R^{n-1} \sin[(n-1)\vartheta + q_1] \\
&\quad + \cdots + p_n \sin q_n,
\end{aligned}
}
and we also form the derivatives taken with respect to $\vartheta$
%CONT>>


% ===== scan leaf 322, printed 297 =====
% !PRINTED 297
\weberpage{297}
%<<CONT
\weq{4}{
\begin{aligned}
\frac{d\varphi}{d\vartheta} &= -n R^n \sin n\vartheta - (n-1) p_1 R^{n-1} \sin[(n-1)\vartheta + q_1] \\
&\quad - \cdots - p_{n-1} \sin(\vartheta + q_{n-1}) \\[4pt]
\frac{d\psi}{d\vartheta} &= n R^n \cos n\vartheta + (n-1) p_1 R^{n-1} \cos[(n-1)\vartheta + q_1)] \\
&\quad + \cdots + p_{n-1} \cos(\vartheta + q_{n-1}).
\end{aligned}
} % [sic] stray closing parenthesis after q_1 in the cos-bracket of the second line, as printed

We now divide the periphery of the circle of radius $R$ into
$4n$ parts of the angular magnitude
\[
2\omega = 2\,\frac{\pi}{4n},
\]
beginning at $\vartheta = -\omega$. The points of division are therefore
\[
-\omega,\ \omega,\ 3\omega,\ 5\omega,\ \dots (8n - 3)\omega,
\]
and let the intervals be denoted in order by $1, 2, 3 \dots 4n$.
Fig. 15 shows this division for $n = 3$.

\weberfigureimage{weber-b1-fig15.png}{0.54}

In the intervals
\[
1, 3, 5 \dots (4n - 1)
\]
$\cos n\vartheta$, taken absolutely, is greater than $1 : \sqrt{2}$, and
is alternately of positive and negative sign. In the
intervals
\[
2, 4, 6 \dots 4n
\]
$\sin n\vartheta$ is absolutely greater than $1 : \sqrt{2}$ and likewise
alternately of positive and negative sign.

But $R$ may be taken so large that in the intervals
$1, 3, 5 \dots (4n - 1)$ the sign of $\varphi$ and of $\dfrac{d\psi}{d\vartheta}$ agrees with that
of $\cos n\vartheta$, and in the intervals $2, 4, 6 \dots 4n$ the sign
of $\psi$ and of $\dfrac{d\varphi}{d\vartheta}$ agrees with that of $\sin n\vartheta$.

Since, then, $\varphi$ changes its sign in the intervals $2, 4, 6 \dots 4n$,
it must pass through zero in each of these intervals;
and since the derivative $\dfrac{d\varphi}{d\vartheta}$ does not change its sign, it can
pass through zero only \gesperrt{once} in each.


% ===== scan leaf 323, printed 298 =====
% !PRINTED 298
\weberpage{298}
In the same way it is to be concluded that $\psi$ passes through zero
once and only once in each of the intervals
$1, 3, 5 \dots 4n - 1$.
We thus have a division of the circumference of the circle such as was
required in the preceding section for the determination of the characteristic,
and since at the points of division $\omega, 5\omega, 9\omega \dots (4n - 3)\omega$
the product $\varphi\psi$ is positive, only points $E(f; \varphi, \psi)$
occur on the circumference of the circle, and their number is $2n$; the number
of the roots of $F(z)$ is therefore equal to $n$, which was to be proved.


% ===== scan leaf 324, printed 299 =====
% !PRINTED 299
% ?? uncertain: no folio/running head visible on any of the three bands of this page (section-opening page); number inferred from the printed "300" on the immediately following page.
\weberpage{299}
\pdfbookmark[1]{Ninth Section. Estimation of the Roots.}{bm110}
\pdfbookmark[2]{§ 99. The Budan-Fourier theorem}{bm111}

\begin{center}
\gesperrt{Ninth Section.}

\gesperrt{Estimation of the Roots.}

\rule{0.3\linewidth}{0.4pt}

\S.\,99.

\gesperrt{The Budan-Fourier Theorem.}
\end{center}

The functions of the \gesperrt{Sturm} chains do indeed solve completely and without exception the problem of determining the number of roots of an equation between given limits; but their actual computation is usually difficult and often impracticable. A number of theorems are known for the \gesperrt{estimation} of the number of roots between given limits, which furnish much simpler rules, but which, to be sure, also do not answer the question completely, giving only a maximal number beyond which the number of roots cannot go. These rules are often very useful and sufficient for applications, and they must therefore not be omitted here.

We consider first a procedure which originates with \gesperrt{Budan}, but which was then used and extended by \gesperrt{Fourier}\,\textsuperscript{1}).

Let us consider, in place of a \gesperrt{Sturm} chain, the sequence of derivatives of a real function $f(x)$ of degree $n$

\weq{1}{f(x),\ f'(x),\ f''(x)\ .\ .\ .\ f^{(n)}(x),}

in which, then, $f^{(n)}(x)$ is a constant, which we assume to be different from zero and positive. It is, apart from a positive numerical factor, the coefficient of $x^n$ in $f(x)$.

\noindent\rule{0.25\linewidth}{0.4pt}

\footnotesize
\textsuperscript{1}) \gesperrt{Budan}, Nouvelle methode pour la r\'esolution des \'equations num\'eriques, (presented to the Paris Academy in 1803). \gesperrt{Fourier}, Analyse des \'equations d\'etermin\'ees. Paris 1831. On the history of this question compare \gesperrt{Lagrange}, Trait\'e de la r\'esolution des \'equations, Note VIII (Works, Vol. 8).
\normalsize


% ===== scan leaf 325, printed 300 =====
% !PRINTED 300
\weberpage{300}

Let $\alpha$ and $\beta > \alpha$ again be two real numbers, which determine the interval $(\alpha, \beta)$ in which the number of real roots of $f(x)$ is to be counted.

We shall first assume that in the interval $(\alpha, \beta)$ no two terms of the sequence (1) vanish simultaneously, and that for $x = \alpha$, $x = \beta$ no term of the sequence vanishes.

If we now let $x$ increase continuously from $\alpha$ to $\beta$, a change can occur in the succession of signs of the sequence (1) only when one of the functions passes through zero. When $f^{(\nu)}(x)$ passes through zero for $x = \xi$, then $f^{(\nu)}(x)$ passes from negative to positive or from positive to negative values, according as $f^{(\nu+1)}(x)$ is positive or negative, and thus a change of sign between $f^{(\nu)}$ and $f^{(\nu+1)}$ is lost in the passage through $\xi$. This holds also when $f^{(\nu)}$ is the function $f(x)$ itself. But when $f^{(\nu)}$ is one of the derivatives, it is preceded by a function $f^{(\nu-1)}$, and since $f^{(\nu-1)}$ is by hypothesis not zero for $x = \xi$, and $f^{(\nu)}$ changes its sign in the passage through $\xi$, either a loss or a gain of one change of sign takes place between $f^{(\nu-1)}$ and $f^{(\nu)}$ in the passage through $\xi$. Thus when an inner term of the sequence (1) passes through zero, the number of changes of sign remains unaltered, or two changes of sign are lost.

But if $f(x)$ itself passes through zero, one change of sign is lost.

From this follows the theorem:

I. \gesperrt{The number of roots of $f(x)$ lying between $\alpha$ and $\beta$ is at most as great as the number of changes of sign lost between $\alpha$ and $\beta$, and if it is smaller, the difference is an even number.}

Making use of a formula, we can also say:

If $V(x)$ is the number of changes of sign (variations) that the sequence (1) presents for any value $x$, then the number of roots of $f(x)$ lying between $\alpha$ and $\beta$ is

\weq{2}{V(\alpha) - V(\beta) - 2h,}

where $h$ is a non-negative integer.

The proof of theorem I. still requires a supplement for the case in which several successive terms of the sequence (1) vanish simultaneously.


% ===== scan leaf 326, printed 301 =====
% !PRINTED 301
\weberpage{301}

Suppose then that for a value $\xi$ of $x$ between $\alpha$ and $\beta$, in the sequence

\weq{3}{f^{(\nu)}(x),\ f^{(\nu+1)}(x),\ .\ .\ .\ f^{(\nu+\mu-1)}(x),\ f^{(\nu+\mu)}(x)}

all the terms, with the exception of the last, vanish, and suppose the last, $f^{(\nu+\mu)}(x)$, has, say, a positive value.

We mark off about $\xi$ two intervals $\delta_1$, $\delta_2$, such that all values of $x$ in the interval $\delta_1$ are smaller than $\xi$, and in the interval $\delta_2$ greater, and we take these intervals so small that the functions (3) do not vanish in them except at $\xi$, so that in particular $f^{(\nu+\mu)}(x)$ remains positive.

Now since, when $f^{(\nu+\mu)}(x)$ is positive, $f^{(\nu+\mu-1)}(x)$ increases together with $x$, we have

\[
\begin{array}{l}
f^{(\nu+\mu-1)}(x) \text{ in } \delta_1 \text{ negative,} \\
\text{in } \delta_2 \text{ positive.}
\end{array}
\]

From this it follows that $f^{(\nu+\mu-2)}(x)$ decreases in $\delta_1$ and increases in $\delta_2$, hence is positive in both intervals; and continuing in this way we infer the succession of signs in the sequence (3):

\[
\begin{array}{l}
(4)\quad f^{(\nu)}(x),\ f^{(\nu+1)}(x),\ f^{(\nu+2)}(x),\ .\ .\ .\ f^{(\nu+\mu-1)}(x),\ f^{\nu+\mu}(x) \\[4pt]
\delta_1,\quad (-1)^\mu \quad -(-1)^\mu \quad (-1)^\mu \quad .\ .\ . \quad - \quad + \\[2pt]
\delta_2,\quad + \quad + \quad + \quad .\ .\ . \quad + \quad +
\end{array}
\] % [sic] last term of row (4) prints as $f^{\nu+\mu}(x)$, without parentheses around the exponent, unlike the other terms in the same row

that is, in the sequence (3), on passing through $\xi$, nothing but changes of sign become successions of signs, and $\mu$ changes of sign are lost in the sequence (3).

If $f^{(\nu+\mu)}(\xi)$ is negative, then all the signs in (4) are the opposite ones, and the conclusion otherwise remains the same.

Now if $\nu = 0$, that is, if $f^{(\nu)}(x)$ is the original function $f(x)$ itself, which then has a $(\mu+1)$-fold root at $\xi$, then on passing through $\xi$ a loss of $\mu$ changes of sign also takes place in the sequence (1).

But if $\nu > 0$ and the function $f^{(\nu-1)}(x)$ preceding $f^{(\nu)}(x)$ is different from zero, then we have the following signs:

1. $\mu$ even:

\[
\begin{array}{ll}
\text{a)}\ f^{(\nu-1)}(x),\ f^{(\nu)}(x), & \text{b)}\ f^{(\nu-1)}(x),\ f^{(\nu)}(x) \\[2pt]
\delta_1,\quad + \quad + & \quad - \quad + \\[2pt]
\delta_2.\quad + \quad + & \quad - \quad +
\end{array}
\]


% ===== scan leaf 327, printed 302 =====
% !PRINTED 302
\weberpage{302}

2. $\mu$ odd:

\[
\begin{array}{ll}
\text{a)}\ f^{(\nu-1)}(x),\ f^{(\nu)}(x), & \text{b)}\ f^{(\nu-1)}(x),\ f^{(\nu)}(x) \\[2pt]
\delta_1,\quad + \quad - & \quad - \quad - \\[2pt]
\delta_2,\quad + \quad + & \quad - \quad +
\end{array}
\]

Thus for even $\mu$ no change of sign is lost between $f^{(\nu-1)}$ and $f^{(\nu)}$; for odd $\mu$ a change of sign is lost or one is gained. Hence the loss of changes of sign on passing through $\xi$ in the sequence

\weq{5}{f^{(\nu-1)},\ f^{(\nu)},\ f^{(\nu+1)}\ .\ .\ .\ f^{(\nu+\mu)}}

is equal to $\mu$ for even $\mu$, and to $\mu \pm 1$ for odd $\mu$; it is therefore always an \gesperrt{even number} and never negative.

We can give these results a clear analytic expression, which lets their meaning be recognised better.

As before, we denote by $V(x)$ the number of changes of sign in the sequence (1) for a definite value of $x$, but with the further stipulation that, if for a value of $x$ some of the terms of the sequences vanish, these are simply passed over in counting the changes of sign. We can then regard $V(x)$ as a function of $x$, which however can change only by integers, and is therefore discontinuous for those values of $x$ for which some terms of the sequence (1) vanish. We shall then denote by $V(x-0)$ and $V(x+0)$ the values of the function $V(x)$ immediately before and immediately after such a place. Then the consideration of (4) shows that in all cases

\weq{6}{V(\xi + 0) = V(\xi)}

and that, if $\xi$ is a $\mu$-fold root of $f(x)$,

\weq{7}{V(\xi - 0) = V(\xi) + \mu + 2h,}

where $h$ is a non-negative integer.

From this it follows, then, if for the present we hold to the assumption that for $x = \alpha$ and $x = \beta$ none of the functions (1) vanishes, that $V(\alpha) - V(\beta)$ differs from the sum of all the numbers $\mu$, that is, from the number of the roots lying in the interval $(\alpha, \beta)$, each counted according to the degree of its multiplicity, by an even, non-negative number; and with this Theorem I is thus freed from its restriction.


% ===== scan leaf 328, printed 303 =====
% !PRINTED 303
\weberpage{303}

But we can also give up the assumption that none of the functions (1)
vanishes for $\alpha$ and $\beta$, provided that $\alpha$ and $\beta$ do
not occur among the roots of $f(x)$.

For then our theorem is applicable to two values, of which the first is a
little greater than $\alpha$, the other a little smaller than $\beta$, and
which enclose the same roots as $\alpha$ and $\beta$; that is, if we
denote by $\mathfrak{z}$ the number of these roots, we have
\[
V(\alpha + 0) - V(\beta - 0) = \mathfrak{z} + 2h,
\]
and since
\[
V(\alpha + 0) = V(\alpha), \quad V(\beta - 0) = V(\beta) + 2h',
\]
it follows that
\[
V(\alpha) - V(\beta) = \mathfrak{z} + 2(h + h'),
\]
which is again the statement of our theorem.

But if $\alpha$ or $\beta$ itself belongs to the roots of $f(x)$, then our
theorem is correct only \gesperrt{when these limiting values are not counted as
belonging to the interval}.

The \gesperrt{Budan-Fourier} theorem does not, it is true, give a
completely certain decision about the number of roots in an interval, as
the \gesperrt{Sturm} theorem does; but it can nevertheless in many cases
replace the \gesperrt{Sturm} theorem; for if one has restricted the
interval $(\alpha, \beta)$ so far that no change of sign, or only one, is
lost from $\alpha$ to $\beta$, then it follows with certainty that in the
first case no root, and in the second case one and only one root, lies in
the interval.

The separation of the roots can therefore be completely given by the Budan
sequence only when, in the passage through a value $\xi$, several changes
of sign are not lost simultaneously. This behaviour cannot be recognised
with certainty, however small the interval $(\alpha, \beta)$ may be
chosen, and so, if after a suitable narrowing of the interval the
separation has not succeeded, one will after all have to resort to another
procedure, in the last instance to the \gesperrt{Sturm} theorem.


% ===== scan leaf 329, printed 304 =====
% !PRINTED 304
\weberpage{304}
\pdfbookmark[2]{§ 100. Newton's rule}{bm112}

\begin{center}
\S.\,100.

\gesperrt{Newton's Rule.}
\end{center}

A supplement to the preceding rule, which permits a further
approximation to the exact determination of the limits of the roots, is
given by the prescription now to be discussed, which originates with
\gesperrt{Newton} but was first proved by
\gesperrt{Sylvester}\footnote{\gesperrt{Newton},
Arithmetica universalis. \gesperrt{Sylvester}, Transactions of the
R. Irish Academy, t.\,24 (1871). Phil. Mag., 4.\,ser., t.\,31. Cf. also
\gesperrt{Petersen}, Theorie der algebr. Gleichungen. Kopenhagen 1878.}.

Again let
\weq{1}{f(x) = a_0 x^n + a_1 x^{n-1} + \cdots + a_{n-1} x + a_n = 0}
be the equation to be investigated. Denoting by $\Pi(m)$, as always,
the product $1.\,2.\,3 \ldots m$, and by $f^{(\nu)}(x)$ the $\nu$th
derivative, we set
\weq{2}{f_\nu(x) = \frac{\Pi(n - \nu)}{\Pi(n)} f^{(\nu)}(x),}
\weq{3}{F_\nu(x) = f_\nu^2(x) - f_{\nu+1}(x) f_{\nu-1}(x),}
with the addition that
\weq{4}{F_0 = F = 1, \quad F_n = f_n^2,}
and hence, which is the essential point, are to be positive.

For the derivatives of the functions $f_\nu$, $F_\nu$ we obtain from
(2) and (3)
\weq{5}{f_\nu'(x) = (n - \nu) f_{\nu+1},}
\weq{6}{f_\nu F_\nu'(x) = (n - \nu - 1)(F_\nu f_{\nu+1} + F_{\nu+1} f_{\nu-1}).}

Since the positive numerical factors do not matter, we can, in place of
the Budan sequence, take the sequence of functions
\[
f, f_1, f_2 \ldots f_n
\]
instead.

\gesperrt{Newton's} rule now makes use of the double sequence:
\weq{7}{
\begin{array}{l}
f,\ f_1,\ f_2 \ldots f_n \\
F,\ F_1,\ F_2 \ldots F_n.
\end{array}
}


% ===== scan leaf 330, printed 305 =====
% !PRINTED 305
\weberpage{305}

If we consider two successive terms of this double sequence, such as
\[
\begin{array}{l}
f_\nu,\ f_{\nu+1} \\
F_\nu,\ F_{\nu+1}
\end{array}
\]
for a value of $x$ at which none of the four functions vanishes, then with
regard to the succession of signs the following cases are possible:
\[
\text{a)}\quad
\begin{array}{c}+\ +\\+\ +\end{array},\
\begin{array}{c}+\ +\\-\ -\end{array},\
\begin{array}{c}-\ -\\+\ +\end{array},\
\begin{array}{c}-\ -\\-\ -\end{array},\quad PP,
\]
\[
\text{b)}\quad
\begin{array}{c}+\ -\\+\ -\end{array},\
\begin{array}{c}+\ -\\-\ +\end{array},\
\begin{array}{c}-\ +\\+\ -\end{array},\
\begin{array}{c}-\ +\\-\ +\end{array},\quad VV,
\]
\[
\text{c)}\quad
\begin{array}{c}+\ +\\+\ -\end{array},\
\begin{array}{c}+\ +\\-\ +\end{array},\
\begin{array}{c}-\ -\\+\ -\end{array},\
\begin{array}{c}-\ -\\-\ +\end{array},\quad PV,
\]
\[
\text{d)}\quad
\begin{array}{c}+\ -\\+\ +\end{array},\
\begin{array}{c}+\ -\\-\ -\end{array},\
\begin{array}{c}-\ +\\+\ +\end{array},\
\begin{array}{c}-\ +\\-\ -\end{array},\quad VP.
\]
We denote this behaviour by the following expressions and symbols:

a) succession-succession (permanence-permanence) $PP$,

b) change-change (variation-variation) $VV$,

c) succession-change (permanence-variation) $PV$,

d) change-succession (variation-permanence) $VP$.

The two theorems that we wish to prove read:

\gesperrt{II. The number of roots of $f(x)$ lying between $\alpha$ and
$\beta$ is either exactly as great as, or smaller by an even number than,
the number of change-successions lost in the double sequence (7) in the
passage from $\alpha$ to $\beta$.}

\gesperrt{III. The number of roots of $f(x)$ lying between $\alpha$ and
$\beta$ is either exactly as great as, or smaller by an even number than,
the number of succession-successions gained in the double sequence (7) in
the passage from $\alpha$ to $\beta$.}

Of these two theorems, which do not always give the same upper bound for
the number of roots, one will apply the one that gives the lowest bound.


% ===== scan leaf 331, printed 306 =====
% !PRINTED 306
\weberpage{306}

In the proof of this theorem we make the following assumptions, from
some of which we shall later free ourselves again.

1. $f(x)$ has no double or multiple roots.

2. Of the functions $f_\nu(x)$, no two adjacent ones vanish for
the same value of $x$.

3. Likewise, of the functions $F_\nu(x)$, no two adjacent ones vanish
for the same value of $x$.

4. A consequence of 2. is [by formula (3)] that $f_\nu(x)$ and
$F_\nu(x)$ do not vanish for the same value of $x$.

5. For $x = \alpha$ and $x = \beta$ themselves, none of the functions
$f_\nu$, $F_\nu$ is to vanish.

The proof of Theorems II and III is now as follows:

If we let $x$ go, continuously increasing, from $\alpha$ to $\beta$,
then a change in the succession of signs can occur only when one of the
functions $f_\nu$ or $F_\nu$ passes through zero, and it is only a
matter of investigating the effect that such a passage through zero
produces in the various cases. Let us first assume that one of the
middle functions $f$, say $f_\nu$, passes through zero, where
therefore $0 < \nu < n$. According as $f_{\nu+1}$ is positive or negative,
$f_\nu$ will increase or decrease. Thus in the first case
$f_\nu(x - 0)$ will be negative and $f_\nu(x + 0)$ positive, and in the
second case the reverse; moreover it follows from (3) that, when $f_\nu = 0$,
$F_{\nu-1}$ and $F_{\nu+1}$ are positive, and that $F_\nu$ has the
sign opposite to that of $f_{\nu-1} f_{\nu+1}$. We collect
the possibilities that remain after this in tabular form, but omit
from the enumeration those cases which arise from the ones listed by a
simultaneous change of sign of the three functions $f_{\nu-1}$, $f_\nu$,
$f_{\nu+1}$, since these evidently yield no
different results. There remain to us after this only two
cases:

% [mirrored 2026-08-07] The German fragment's lower third was restored after this page's
% translation pass raised the omission (two sign-table displays and the connecting line);
% the displays below are copied character for character from the restored German.
\[
\begin{array}{|c|ccc|}
 & \nu-1, & \nu, & \nu+1 \\ \hline
f(x-0) & + & - & + \\
f(x+0) & + & + & + \\ \hline
F & + & - & +
\end{array}
\qquad
\begin{array}{|ccc|}
\nu-1, & \nu, & \nu+1 \\ \hline
- & - & + \\
- & + & + \\ \hline
+ & + & +
\end{array}.
\]
In these cases we have accordingly
\[
\begin{array}{|c|cc|}
x-0 & VV, & VV \\
x+0 & PV, & PV
\end{array}
\qquad
\begin{array}{|cc|}
PP, & VP \\
VP, & PP
\end{array}.
\]


% ===== scan leaf 332, printed 307 =====
% !PRINTED 307
\weberpage{307}

\gesperrt{Thus neither a $VP$ nor a $PP$ is lost or gained.}

But if $f(x)$ itself passes through zero, then, assuming $f_1(x)$ positive, which suffices in view of the remark made above, we have
\[
\begin{array}{c|cc}
 & 0 & 1 \\ \hline
f(x-0) & - & + \\
f(x+0) & + & + \\ \hline
F & + & +
\end{array}
\text{.}
\]
Here, then, a $VP$ passes into a $PP$; that is, \gesperrt{a $VP$ is lost and a $PP$ gained.}

If, finally, $F_\nu$ passes through zero, where of course we must have $0 < \nu < n$, the behaviour is as follows:

If $F_\nu = 0$, then $f_{\nu-1}$ and $f_{\nu+1}$ must necessarily have the same sign, since otherwise $F_\nu$ would consist of two positive terms and could not vanish. Further, by (6), $F_\nu'(x)$ has the same sign as the product $f_{\nu-1} f_\nu F_{\nu+1}$, and hence $F_\nu(x-0)$ has the opposite sign to the product, and $F_\nu(x+0)$ the same sign as it.

Here too we may omit from the enumeration those cases which arise from the ones listed by a simultaneous change of sign in all three functions $F_{\nu-1}, F_\nu, F_{\nu+1}$; there remain, then, the following cases:
% [sic] The first right-hand header is printed with a prime after $\nu+1$; the corresponding later header is unprimed.
\[
\begin{array}{c|ccc||ccc}
f & \nu-1, & \nu, & \nu+1 & \nu-1, & \nu, & \nu+1' \\
 & + & + & + & + & + & + \\ \hline
F(x-0) & + & - & + & - & - & + \\ \hline
F(x+0) & + & + & + & - & + & +
\end{array}
\]
\[
\begin{array}{c|ccc||ccc}
f & \nu-1, & \nu, & \nu+1 & \nu-1, & \nu, & \nu+1 \\
 & + & - & + & + & - & + \\ \hline
F(x-0) & + & + & + & - & + & + \\ \hline
F(x+0) & + & - & + & - & - & +
\end{array}
\text{.}
\]
The changes of sign in these four cases are:
%CONT>>


% ===== scan leaf 333, printed 308 =====
% !PRINTED 308
\weberpage{308}
\pdfbookmark[2]{§ 101. Descartes's theorem}{bm113}
%<<CONT
\[
\left|\begin{array}{c|cc|cc|cc|cc}
x-0 & PV, & PV & PP, & PV & VP, & VP & VV, & VP \\ \hline
x+0 & PP, & PP & PV, & PP & VV, & VV & VP, & VV
\end{array}\right|
\text{.}
\]
In the first case two $PP$ are \gesperrt{gained}, in the second and fourth cases no change occurs, in the third case two $VP$ are \gesperrt{lost}.

But with this the theorems II, III are proved.

The assumptions 2., 3., that no two adjacent ones of the functions $f_\nu$ and $F_\nu$ shall vanish for the same $x$, we may now give up. For if we hold fast to the first assumption, that no multiple roots are present, we can impart to the coefficients $a_\nu$ changes so small that, first, the assumptions 2., 3. are satisfied, and secondly, the roots lying between $\alpha$ and $\beta$ change so little that none of them leaves the interval. Since the roots occur only singly, real roots cannot in this process pass into imaginary ones either.

Whether the theorem, with correct counting of the multiple roots, remains valid also in the case of multiple roots may be left undecided.

\begin{center}
\S.\,101.

The Cartesian Theorem.
\end{center}

The \gesperrt{Cartesian} theorem, which is also named after \gesperrt{Harriot}\footnote{Harriot, Artis analyticae praxis ad aequationes algebraicas resolvendas 1631. Gauss' Werke, Bd. III, S. 67. Lagrange, l. c.}, gives a rule for estimating the number of positive roots. We can regard it simply as a special case of \gesperrt{Budan}'s theorem, by taking $\alpha = 0$ and $\beta = \infty$, or at least so large that the functions
\weq{1}{f(x), f'(x), f''(x) \ldots f^{(n)}(x)}
all have the same sign for $x > \beta$. This sign agrees with the sign of the coefficient of the highest power of $x$ in $f(x)$, and may be assumed positive.

If
\[
f(x) = a_0 x^n + a_1 x^{n-1} + a_2 x^{n-2} + \cdots + a_{n-1} x + a_n,
\]
%CONT>>


% ===== scan leaf 334, printed 309 =====
% !PRINTED 309
\weberpage{309}
%<<CONT
so the functions (1), apart from positive numerical factors, take for $x = 0$ the values
\weq{2}{a_n, a_{n-1}, a_{n-2} \ldots a_0,}
and from this, accordingly, we obtain the theorem:

\gesperrt{IV. The number of positive roots of the equation $f(x) = 0$ is equal to, or less by an even number than, the number of changes of sign in the sequence of the coefficients of $f(x)$, where any vanishing coefficients are simply to be passed over, and the root $x = 0$, if it is present, is not counted.}

To estimate the number of negative roots as well, one can either apply the \gesperrt{Budan-Fourier} theorem to the interval $(-\infty, 0)$, or one replaces $x$ by $-x$, that is, one changes the signs of $a_1, a_3, a_5 \ldots$ and then applies Theorem IV.

In a similar way we shall specialise the \gesperrt{Newton} criterion of \S.\,100.

We remark for this purpose that
\[
f_\nu(0) = \frac{\Pi(n-\nu)\,\Pi(\nu)}{\Pi(n)}\, a_{n-\nu},
\]
and consequently
\[
F_\nu(0) = \frac{\Pi(\nu)\,\Pi(\nu-1)\,\Pi(n-\nu)\,\Pi(n-\nu-1)}{\Pi(n)\,\Pi(n)}\left[\nu(n-\nu)\,a_{n-\nu}^2 - (\nu+1)(n-\nu+1)\,a_{n-\nu-1}\,a_{n-\nu+1}\right]
\]
and further that the sequence of the $f_\nu(x)$ presents, for $x = +\infty$, only successions of signs, and for $x = -\infty$ only changes of sign.

If, then, we set
\weq{3}{A_\nu = \nu(n-\nu)\,a_\nu^2 - (\nu+1)(n-\nu+1)\,a_{\nu-1}\,a_{\nu+1},}
\[
A_0 = 1, \quad A_n = 1,
\]
so that $A_\nu$, apart from a positive factor, agrees with $F_{n-\nu}(0)$, it follows that for $x = 0$ the double sequence (7), \S.\,100, written in reverse order, agrees in sign with
\weq{4}{
\begin{array}{c}
a_0,\ a_1,\ a_2 \ldots a_n \\
A_0,\ A_1,\ A_2 \ldots A_n.
\end{array}
}
On passing to $x = +\infty$, in the sequence $f(x), f_1(x) \ldots f_n(x)$, all the changes of sign are, and on passing to
%CONT>>


% ===== scan leaf 335, printed 310 =====
% !PRINTED 310
\weberpage{310}
%<<CONT
$x = -\infty$ all successions of signs have been lost; hence in the double sequence of the $f_\nu, F_\nu$, for $x = +\infty$ all $VP$, and for $x = -\infty$ all $PP$.

Accordingly, by \S.\,100, II, III, we can state the following two theorems:

\gesperrt{V. The number of positive roots of $f(x)$ is as great as, or less by an even number than, the number of change-successions in the double sequence (4).}

\gesperrt{VI. The number of negative roots of $f(x)$ is as great as, or less by an even number than, the number of succession-successions in the double sequence (4).}

But the sum of the numbers of the $VP$ and of the $PP$ is equal to the number of successions of signs in the simple sequence
\weq{5}{A_0, A_1, A_2 \ldots A_n,}
and we can therefore combine V. and VI. into the single theorem:

\gesperrt{VII. The number of all real roots of $f(x)$ is as great as, or less by an even number than, the number of successions of signs in the sequence (5).}

The total number of changes of sign and successions of signs in the sequence (5) is as great as the number of all the real and imaginary roots taken together, namely $n$. For the case in which some terms of the sequence (5) vanish is here excluded. Since, moreover, the first and last terms of the sequence (5) have the same sign, the number of changes of sign in this sequence is an even number. We can therefore, without saying anything new, add to VII. the further theorem:

\gesperrt{VIII. The number of imaginary roots of $f(x)$ is at least equal to the number of changes of sign in the sequence (5).}

In the formulation of the last two theorems it is immaterial whether some of the coefficients $a_1, a_2 \ldots a_n$ vanish, provided only that none of the $A_\nu$ is zero.

For the cubic equation
\[
a_0 x^3 + a_1 x^2 + a_2 x + a_3 = 0
\]
we have
\[
A_1 = 2(a_1^2 - 3a_0a_2)
\]
\[
A_2 = 2(a_2^2 - 3a_1a_3),
\]
%CONT>>


% ===== scan leaf 336, printed 311 =====
% !PRINTED 311
\weberpage{311}
\pdfbookmark[2]{§ 102. Jacobi's criterion}{bm114}
%<<CONT
and in the sequence
\[
1,\ A_1,\ A_2,\ 1
\]
there thus occur two changes of sign whenever, of the two quantities $A_1$, $A_2$, one or both are negative, and in these cases the equation accordingly has two imaginary roots. But if $A_1$ and $A_2$ are positive, the decision cannot be made from this theorem alone; it is then necessary to take into consideration also the quantity
\[
B = 9\,a_0\,a_3 - a_1\,a_2
\]
and to examine the sign of
\[
3\,D = A_1\,A_2 - B^2
\]
(\S.\,62).

In the case of the biquadratic equations one has to consider the three expressions:
\[
\begin{aligned}
A_1 &= 3\,a_1^2 - 8\,a_0\,a_2,\\
A_2 &= 4\,a_2^2 - 9\,a_1\,a_3,\\
A_3 &= 3\,a_3^2 - 8\,a_2\,a_4.
\end{aligned}
\]
If the signs are, in order, $-+-$, one has four imaginary roots. If all three signs are positive, then four or two or no imaginary roots may be present; in all other cases one has two or four imaginary roots.

\begin{center}
\S.\,102.

The Jacobi Criterion.
\end{center}

In order to estimate the number of roots that lie between two limits $\alpha$, $\beta$, \gesperrt{Jacobi} derived a criterion from the \gesperrt{Cartesian} theorem. If, namely,
\weq{1}{f(x) = 0}
is the proposed equation of degree $n$, of which it is to be decided how many roots it has between the two limits $\alpha$, $\beta$, then \gesperrt{Jacobi}\footnote{Observatiunculae ad theoriam aequationum algebraicarum pertinentes. Crelle's Journal, Bd. 13. Gesammelte Werke, Bd. III.} sets
\weq{2}{y = \frac{x-\alpha}{\beta-x},\qquad x = \frac{\alpha+\beta y}{1+y}}
and now forms the equation
%CONT>>


% ===== scan leaf 337, printed 312 =====
% !PRINTED 312
\weberpage{312}
\pdfbookmark[2]{§ 103. Klein's geometric comparison of the various criteria}{bm115}
%<<CONT
\weq{3}{(1+y)^n f\left(\frac{\alpha+\beta y}{1+y}\right) = \varphi(y) = b_0\,y^n + b_1\,y^{n-1} + \cdots + b_n,}
in which $b_0$, $b_1 \ldots b_n$ are integral homogeneous functions of degree $n$ in $\alpha$ and $\beta$, and linear functions of the coefficients of $f(x)$. Now when $x$ increases from $\alpha$ to $\beta$, $y$, likewise increasing, goes from $0$ to $\infty$. Hence $f(x)$ has exactly as many roots between $\alpha$ and $\beta$ as $\varphi(y)$ has positive roots.

From the \gesperrt{Cartesian} theorem it therefore follows:

\gesperrt{IX. The number of roots of $f(x)$ between $\alpha$ and $\beta$ is as great as, or less by an even number than, the number of changes of sign in the sequence}
\[
b_0,\ b_1,\ b_2 \ldots b_n.
\]

While, then, in the \gesperrt{Budan} theorem the changes of sign have to be counted and compared in two sequences of values, the \gesperrt{Jacobi} criterion requires only the counting of the changes of sign in one sequence of values.

In the same way the \gesperrt{Newton} rule can also be transformed, if one takes into consideration the quantities
\[
B_\nu = \nu(n-\nu)\,b_\nu^2 - (\nu+1)(n-\nu+1)\,b_{\nu-1}\,b_{\nu+1}
\]
and one thus obtains the theorem:

\gesperrt{X. The number of roots of $f(x)$ between $\alpha$ and $\beta$ is as great as, or less by an even number than, the number of changes of succession in the double sequence}
\[
\begin{aligned}
&b_0,\ b_1,\ b_2 \ldots b_n,\\
&B_0,\ B_1,\ B_2 \ldots B_n.
\end{aligned}
\]

\begin{center}
\S.\,103.

Klein's geometric comparison of the various criteria.
\end{center}

In order to compare the scope of the various criteria with one another, \gesperrt{Klein} has employed a geometric interpretation, which we shall now discuss. If one does not wish to operate with spaces of more dimensions, and thereby to lose geometric intuition, one must here
%CONT>>


% ===== scan leaf 338, printed 313 =====
% !PRINTED 313
\weberpage{313}
%<<CONT
confine ourselves to the first cases, and so we begin with the quadratic equations
\weq{1}{x^2 + a\,x + b = 0.}
If we regard $a$ and $b$ as rectangular coordinates in a plane, then every point of this plane is the image of one and only one quadratic equation, and conversely every quadratic equation with real coefficients has a point of the plane as its image.

The equation
\weq{2}{a^2 - 4\,b = 0}
represents a parabola, on which lie the images of the equations with equal roots. The inner points, that is, those lying in the hollow of the parabola, represent the equations with imaginary roots, the outer points the equations with real roots.

We now form, according to \S.\,89, the \gesperrt{Sturm} chain
\weq{3}{x^2 + a\,x + b,\quad 2\,x + a,\quad a^2 - 4\,b}
and examine the changes of sign. If $x$ is constant and $a$, $b$ variable, then
\[
x^2 + a\,x + b = 0
\]
is the equation of the tangent to the parabola at the point
\[
a = -2\,x,\quad b = x^2,
\]
and indeed $x^2 + a\,x + b$ is positive in the part of the plane that also contains the parabola.

By this tangent and by the parabola the plane is divided into four regions, in which, as is indicated in Fig. 16 (on the following page), 0, 1, 2 changes of sign occur in the sequence (3).

If we now set $x = \alpha$ and $x = \beta$, we obtain two such tangents and a division of the plane into four regions, in which lie the representative points of those equations that have no root, one root, or two roots between $\alpha$ and $\beta$. In Fig. 17 these four regions are marked by the corresponding numerals.

This is the difference between those points from which no tangent, one tangent, or two tangents can be drawn to the parabola having their point of contact between $\alpha$ and $\beta$. Fig. 17 thus gives the exact and complete distinction of the various cases.


% ===== scan leaf 339, printed 314 =====
% !PRINTED 314
\weberpage{314}
If we now compare with this the theorem of \gesperrt{Budan}, we have the three functions $f(x)$, $f'(x)$, $f''(x)$, and thus have

% SOURCE FIGURE CROP: Bd. I registered scan 340 (printed p. 314), 600 dpi,
% crop x=300,y=700,w=2700,h=1440; source-faithful raster, Fig. 16--17.
\weberfigureimage{weber-b1-fig16-17.png}{0.94}

\weq{4}{x^2 + a\,x + b,\quad 2\,x + a,\quad 1}
to examine with regard to the changes of sign.

The two straight lines
\[
x^2 + a\,x + b = 0,\quad 2\,x + a = 0
\]
divide the plane into three regions, in which the sequence (4) presents the changes of sign indicated in Fig. 18.

% SOURCE FIGURE CROP: Bd. I registered scan 340 (printed p. 314), 600 dpi,
% crop x=350,y=2820,w=2650,h=1450; source-faithful raster, Fig. 18--19.
\weberfigureimage{weber-b1-fig18-19.png}{0.94}

If we again take two values $x = \alpha$ and $x = \beta$, six regions result, in which the loss of changes of sign is given in Fig. 19. In the regions marked 0 and 1
%CONT>>


% ===== scan leaf 340, printed 315 =====
% !PRINTED 315
\weberpage{315}
%<<CONT
regions, then, the loss of changes of sign gives the correct number of
roots. In the region 2, on the other hand, the decision between two
roots and none is not made.

Let us now take the \gesperrt{Jacobi} criterion (\S.\,102). According
to it, in order to estimate the number of roots lying between $\alpha$
and $\beta$, we must first transform the equation (1) into
\[
(\alpha+\beta y)^2+a(\alpha+\beta y)(1+y)+b(1+y)^2=0
\]
or
\[
\begin{aligned}
&(\beta^2+a\beta+b)y^2+(2\alpha\beta+a(\alpha+\beta)+2b)y\\
&\quad{}+(\alpha^2+a\alpha+b)=0,
\end{aligned}
\]
and now count the number of changes of sign in the three functions
\[
\begin{aligned}
&\beta^2+a\beta+b\\
&2\alpha\beta+a(\alpha+\beta)+2b\\
&\alpha^2+a\alpha+b
\end{aligned}
\]
The first and the third of these expressions, set equal to zero,
represent the two parabola tangents considered above. The middle one,
$2\alpha\beta+a(\alpha+\beta)+2b$, vanishes for $a=-2\alpha$,
$b=\alpha^2$ and for $a=-2\beta$, $b=\beta^2$ and thus represents the
line joining the two points $\alpha$, $\beta$, and indeed in such a way
that it becomes negative on the side that contains the point of
intersection of the two tangents [at this point of intersection itself
it is equal to $-(\beta-\alpha)^2$].

% SOURCE FIGURE CROP: Bd. I registered scan 341 (printed p. 315), 600 dpi,
% crop x=100,y=3000,w=700,h=620; source-faithful raster, Fig. 20.
\weberfigureimage{weber-b1-fig20.png}{0.34}
We obtain here five regions with no, one or two changes of sign, as
Fig. 20 indicates.

One sees from this that this figure confines the indecision to a much
smaller region than the previous one. From this point of view, then, it
does not appear justified that \gesperrt{Jacobi} so decidedly gives the
\gesperrt{Budan-Fourier} criterion preference over his
own\footnote{Cf. F.\ \gesperrt{Klein}, Geometrisches zur
Abz\"ahlung der reellen Wurzeln einer algebraischen Gleichung in
\gesperrt{Dyck}'s Catalog mathematischer Modelle (M\"unchen 1892).}.


% ===== scan leaf 341, printed 316 =====
% !PRINTED 316
\weberpage{316}
\pdfbookmark[2]{§ 104. Determination of an upper bound for the roots}{bm116}

\begin{center}
\S.\ 104.

\gesperrt{Determination of an upper limit for the roots.}
\end{center}

We have already, in the third section, made use of the theorem that
one can always find a positive number which is greater than the
absolute values of all the roots of a given equation. But now the
point is to give a determination of such an upper limit that is as
simple and at the same time as accurate as possible.

We consider at first only the real roots, and thus seek a positive
number which is greater than the greatest positive root of a given
equation $f(x)=0$.

Such a determination is given to us by the theorem of \gesperrt{Budan}.

If we take the coefficient of the highest power of $x$ in $f(x)$
equal to 1 (or at least positive), then the positive number
$\alpha$ can always be taken so large that
\weq{1}{f(\alpha),\ f'(\alpha),\ f''(\alpha)\ldots f^{(n)}(\alpha)}
all become positive. Then in the sequence of functions
\[
f(x),\ f'(x),\ f''(x)\ldots f^{(n)}(x)
\]
no change of sign is lost any more between $x=\alpha$ and $x=\infty$, and
hence, by Theorem I, no root of $f(x)$ can lie between
$\alpha$ and $\infty$.

If one wishes to obtain from the same theorem a lower limit for the
negative roots, one takes the positive number $\beta$ so that
\weq{2}{(-1)^n f(-\beta),\ (-1)^{n-1} f'(-\beta),\ \ldots f^{(n)}(-\beta)}
all become positive; then the equation $f(x)=0$ certainly has no
negative root below $-\beta$.

This determination of the limits goes back already to \gesperrt{Newton}.

Another determination of an upper limit, which is simpler for
computation, has been given by \gesperrt{Laguerre}\footnote{\gesperrt{Laguerre},
Nouvelles Annales d.\ math., 2.\ ser., t.\ XIX, 1880. Journal de math.,
3.\ ser., t.\ IX, 1883.}.

Instead of the derivatives $f(x), f'(x), f''(x)\ldots$ he uses the
functions $f_0(x), f_1(x), f_2(x), \ldots f_{n-1}(x)$, which have already
%CONT>>


% ===== scan leaf 342, printed 317 =====
% !PRINTED 317
\weberpage{317}
%<<CONT
many ways, and in particular have rendered good service in the
Tschirnhausen transformation.

We defined these functions in \S.\,4 by
\[
\begin{aligned}
f_0(x) &= 1\\
f_1(x) &= x+a_1\\
f_2(x) &= x^2+a_1x+a_2\\
&\ \cdots\cdots\cdots\cdots\cdots\cdots\cdots\cdots\\
f_{n-1}(x) &= x^{n-1}+a_1x^{n-2}+a_2x^{n-3}+\cdots+a_{n-1}\\
f(x)=f_n(x) &= x^n+a_1x^{n-1}+a_2x^{n-2}+\cdots+a_{n-1}x+a_n.
\end{aligned}
\]
and for them set up the recursion formula
\[
f_v(x)-xf_{v-1}(x)=a_v.
\]
One now has, as has already been shown in the place cited,
\weq{3}{f(x)=(x-\alpha)\left[x^{n-1}f_0(\alpha)+x^{n-2}f_1(\alpha)+\cdots+f_{n-1}(\alpha)\right]+f(\alpha).}
A glance at this formula shows that if $\alpha$ is a number which makes
the functions
\weq{4}{f_0(\alpha),\ f_1(\alpha),\ f_2(\alpha)\ldots f_{n-1}(\alpha),\ f(\alpha)}
positive, then no positive $x$ greater than $\alpha$ can exist which
makes $f(x)$ zero, since in that case nothing but positive terms stand
on the right-hand side of (3).

If, then, a positive $\alpha$ is so determined that the functions
(4) are all positive, this is an upper limit for the positive
roots of $f(x)$.

In the same way one obtains a lower limit $-\beta$ for the negative roots,
if one determines the positive quantity $\beta$ so that
\weq{5}{f_0(-\beta),\ -f_1(-\beta),\ f_2(-\beta),\ \ldots (-1)^n f(-\beta)}
all become positive.

This determination of the limits is indeed simpler to handle in
calculation, but under some circumstances it gives less accurate limits than
the \gesperrt{Newton} determination.

For if one differentiates (3) repeatedly with respect to $x$ and then
sets $x=\alpha$, one obtains, under the assumption that the
quantities (4) are all positive, for the derivatives $f'(x)f''(x)\ldots$
expressions which for $x=\alpha$ consist of nothing but positive terms.
If, then, the expressions (4) are positive, the derivatives (1) are
also necessarily all positive, while the converse is not
necessarily the case.


% ===== scan leaf 343, printed 318 =====
% !PRINTED 318
\weberpage{318}
\pdfbookmark[2]{§ 105. Estimation of the imaginary roots}{bm117}

In order to find an upper limit for the absolute value of the imaginary
roots which remains valid also in the case of complex coefficients, we
consider the function
\[
f(x)=x^n+a_1x^{n-1}+a_2x^{n-2}+\cdots+a_n,
\]
in which the coefficients $a_1, a_2\ldots a_n$ are real or imaginary.
We denote the absolute value of a complex quantity $a$, as before, by
$|a|$, and then, if we assume that the absolute value of $x$ is greater
than 1, we have
\[
\left| a_1+\frac{a_2}{x}+\cdots+\frac{a_n}{x^{n-1}} \right| <
|a_1|+|a_2|+\cdots+|a_n|=\alpha.
\]
If we therefore take $x$ so that its absolute value is greater than
each of the two numbers 1 and $\alpha$, the expression
\[
x^{1-n}f(x)=x+a_1+\frac{a_2}{x}+\cdots+\frac{a_n}{x^{n-1}}
\]
cannot vanish, since the absolute value of $x$ is greater than the
absolute value of the sum of all the remaining terms of the right-hand
side.

If, then, according as $\alpha$ is greater or less than 1, we draw a
circle with the radius $\alpha$ or 1 about the origin as centre,
certainly no roots of $f(x)$ lie outside this circle any longer.

\begin{center}
\S.\ 105.

\gesperrt{Estimation of the imaginary roots.}
\end{center}

We have seen in \S.\,96 how, from the characteristic of a certain
system of curves, we can determine the exact number of the complex
roots of an equation that lie in the interior of a closed curve. This
criterion too can be simplified if one no longer wishes to find the
exact number of the roots, but only an upper limit for this number,
which in many cases also suffices for the exact determination.

Let, then, as in \S.\,96
\weq{1}{F(x)=\varphi(x,y)+i\psi(x,y)}
be a real or complex function of the complex argument
%CONT>>


% ===== scan leaf 344, printed 319 =====
% !PRINTED 319
\weberpage{319}
\pdfbookmark[2]{§ 106. Rolle's theorem}{bm118}
%<<CONT
$z = x + yi$; therein let $\varphi$ and $\psi$ be real functions of $x$ and
$y$, and now let a closed curve $f$ be drawn. We then have the theorem:

\gesperrt{XI. If the curve $f$ is divided by the curve $\varphi$ into $2m$
segments, in which $\varphi$ is alternately positive and negative, then the
number of roots of $F(z)$ lying within $f$ is at most equal to $m$;
and the same also holds when $f$ is divided by $\psi$ into $2m$ such
segments.}

The theorem is an immediate consequence of the theorem on characteristics II,
\S.\,96; for by that theorem
\[
k = \tfrac{1}{2}[E(f;\varphi,\psi) - A(f;\varphi,\psi)];
\]
and if now the curve $f$ is cut by the $\varphi$-curve in $2m$ points,
then
\[
m = \tfrac{1}{2}[E(f;\varphi,\psi) + A(f;\varphi,\psi)],
\]
hence
\[
k = m - A(f;\varphi,\psi).
\]

But here, by the theorem just mentioned, $k$ is equal to the number of roots
of $F$ lying in the interior of $f$, and $A(f;\varphi,\psi)$ is in any case a
non-negative number. In exactly the same way one can proceed with the formula
\[
k = \tfrac{1}{2}[E(f;\psi,\varphi) - A(f;\psi,\varphi)]
\]

\begin{center}
\S.\,106.

The Theorem of Rolle.
\end{center}

If the equation $f(x) = 0$ has two real roots $\alpha$ and $\beta$, of which
the first is an $a$-fold root and the second a $b$-fold root, then we can
set
\weq{1}{f(x) = (x - \alpha)^a (x - \beta)^b f_1(x)}

Now if $\alpha < \beta$, and no root of $f(x) = 0$ lies between $\alpha$ and
$\beta$, then $f(x)$ and $f_1(x)$ will have an unchanged sign throughout
this whole interval. If we form the derivative of (1), we obtain
%CONT>>


% ===== scan leaf 345, printed 320 =====
% !PRINTED 320
\weberpage{320}
%<<CONT
\weq{2}{\dfrac{(x-\alpha)(x-\beta)f'(x)}{f(x)} = a(x-\beta) + b(x-\alpha) + (x-\alpha)(x-\beta)\dfrac{f_1'(x)}{f_1(x)}.}
The right-hand side of this expression becomes negative for $x = \alpha$, namely
$a(\alpha - \beta)$, and positive for $x = \beta$, namely $b(\beta - \alpha)$,
and must therefore, while $x$ goes from $\alpha$ to $\beta$, pass through zero
once, or more generally an odd number of times. On the left-hand
side of (2), $(x-\alpha)(x-\beta)$ and $f(x)$ have an unchanged
sign, and hence $f'(x)$ must pass through zero an odd number of
times. From this follows Rolle's theorem:

\gesperrt{XII. If $\alpha$ and $\beta$ are two successive roots
of $f(x) = 0$, then between $\alpha$ and $\beta$ (not counting the
limits) there lie roots of $f'(x) = 0$ in odd number, and hence at
least one.}

Important consequences follow from this for the case in which the equation
$f(x) = 0$ has \gesperrt{only real roots}. If these, ordered by magnitude, are
$\alpha, \beta, \gamma, \dots$, and if $\alpha$ is an $a$-fold, $\beta$ a
$b$-fold, $\gamma$ a $c$-fold root, etc., then $f'(x) = 0$ has at
$\alpha$ an $(a-1)$-fold, at $\beta$ a $(b-1)$-fold, at $\gamma$ a
$(c-1)$-fold root, and between $\alpha$ and $\beta$, between $\beta$ and
$\gamma$, and so on, at least one real root each. But the total number of these
roots is
\[
a + b + c + \cdots - 1 = n - 1,
\]
and since $f'(x)$ is of degree $(n-1)$, this minimal
number is also not exceeded. From this follows the theorem:

% [page alignment mirrored 2026-08-07] The German fragment boundary was corrected: theorem
% XIII's tail, the connecting sentence, and theorem XIV are printed on page 320 and were moved
% back here from i0346. The English follows the corrected German, including the source
% emphasis boundary: the theorem numeral is plain and the statement is letterspaced.
XIII. \gesperrt{If $f(x) = 0$ has only real roots, then
$f'(x) = 0$ also has only real roots, and of these, those that do not
coincide with multiple roots of $f(x)$ are simple roots, which are
separated by the roots of $f(x)$.}

If we apply this theorem repeatedly, we can generalise it as
follows:

XIV. \gesperrt{If $f(x)$ has only real roots, then $f'(x), f''(x),
f'''(x) \dots$ also have only real roots. If one of these functions $f^{(\nu)}(x)$
has an $a$-fold root $\alpha$, and if $a > 1$, then $\alpha$ is an
$(a+\nu)$-fold root of $f(x)$.}


% ===== scan leaf 346, printed 321 =====
% !PRINTED 321
\weberpage{321}
% [page alignment mirrored 2026-08-07] The three blocks that opened this fragment (tail of
% XIII, connecting sentence, theorem XIV) are printed on page 320 and were moved to i0345,
% following the corrected German. Page 321 begins directly with "Diese Sätze werden…".
These theorems are substantially generalised if one applies a linear
transformation to them\footnote{\gesperrt{Laguerre}, Nouvelles
Annales de Math\'ematiques, 2.\,S\'er., Vol.\,19, 1880.}.

This is more easily carried out if one uses the homogeneous notation,
that is, understands by $f(x,y)$ an integral homogeneous function
of degree $n$. If, for an arbitrary value $\xi:\eta$, one forms the polars
of this function (\S.\,60):
\[
P_1(x,\xi) = \frac{1}{n}\{\xi f'(x) + \eta f'(y)\}
\]
\weq{3}{P_2(x,\xi) = \frac{1}{n(n-1)}\{\xi^2 f''(x,x) + 2\xi\eta f''(x,y) + \eta^2 f''(y,y)\}}
\[
P_3(x,\xi) = \frac{1}{n(n-1)(n-2)}\{\xi^3 f'''(x,x,x) + 3\xi^2\eta f'''(x,x,y) + 3\xi\eta^2 f'''(x,y,y) + \eta^3 f'''(y,y,y)\}
\]
\[
\cdots\cdots\cdots\cdots\cdots\cdots\cdots\cdots\cdots\cdots\cdots\cdots,
\]
then these functions remain unchanged if one transforms $\xi, \eta$ and $x, y$
simultaneously by a linear transformation and puts for $f$ and
its derivatives the corresponding transformed functions. With
respect to the reality of the roots, nothing is changed by a real linear
transformation. But now, if $\xi, \eta$ are real,
the real linear substitution can be so determined that the transformed
values become $\xi' = 1$, $\eta' = 0$, and thereby, if one also sets
$y = 1$, the polars pass into the derivatives of the function $f(x)$ with respect to $x$
(apart from numerical factors).

\gesperrt{From this it follows that in the theorems XIII, XIV the derivatives $f'(x), f''(x)
\dots$ can be replaced by the polars $P_1(x,\xi), P_2(x,\xi) \dots$ for arbitrary
real $\xi, \eta$.}

We make an application of this to the equation which one obtains by
setting the $(n-2)$th polar equal to zero:
\weq{4}{x^2 f''(\xi,\xi) + 2xy f''(\xi,\eta) + y^2 f''(\eta,\eta) = 0.}
If $f(x,y) = 0$ has, as we assume, only real roots, then the
quadratic equation (4) must also have real roots, i.\,e. the
Hesse determinant
\weq{5}{f''(\xi,\xi) f''(\eta,\eta) - f''(\xi,\eta)^2}
must be negative for all real $\xi, \eta$. It can only van-
%CONT>>


% ===== scan leaf 347, printed 322 =====
% !PRINTED 322
\weberpage{322}
\pdfbookmark[2]{§ 107. Laguerre's theorems for equations with only real roots}{bm119}
%<<CONT
ish if (4) has two equal roots, and by XIV this is possible only
when $f(x,y)$ is the $n$th power of a linear function.
In this case the determinant (5) vanishes identically: if we
therefore set this case aside, it follows \gesperrt{that for an equation with only
real roots, the equation obtained by setting the Hesse
determinant equal to zero has only imaginary roots.}

If we wish to return to the inhomogeneous representation, we can, making
use of the Euler theorem
\weq{6}{\begin{aligned}
xf'(x) + yf'(y) &= nf(x)\\
xf''(x,x) + yf''(x,y) &= (n-1)f'(x)\\
xf''(y,x) + yf''(y,y) &= (n-1)f'(y)
\end{aligned}}
eliminate $f'(y), f''(x,y), f''(y,y)$ by means of it. One obtains in this way,
on setting $y = 1$,
\weq{7}{\begin{aligned}
f'(y) &= nf(x) - xf'(x)\\
f''(x,y) &= (n-1)f'(x) - xf''(x)\\
f''(y,y) &= n(n-1)f(x) - 2(n-1)xf'(x) + x^2 f''(x),
\end{aligned}}
and hence, if
\weq{8}{(n-1)H(x,y) = f''(x,x)f''(y,y) - f''(x,y)^2}
is set:
\weq{9}{H(x,y) = H(x) = nf(x)f''(x) - (n-1)f'(x)^2.}
$H(x)$ is, with respect to $x$, at most of degree $2n-4$.

\begin{center}
\S.\,107.

The theorems of Laguerre for equations with only real roots.
\end{center}

Of the theorems last derived, \gesperrt{Laguerre} has made an application
to equations with only real roots, which we shall now become acquainted
with.

Let $f(x) = 0$ be an equation of degree $n$ with $n$ \gesperrt{real}
roots $\alpha, \alpha_1, \alpha_2 \dots$, \gesperrt{which we assume to be distinct
from one another}.

If $x$ is a variable quantity, we have the formula, already applied
several times [\S.\,14, (11)]:
%CONT>>


% ===== scan leaf 348, printed 323 =====
% !PRINTED 323
\weberpage{323}

\weq{1}{S\left(\frac{1}{x-\alpha}\right) = \frac{f'(x)}{f(x)},}

in which the summation sign $S$ refers to the different roots
$\alpha, \alpha_1, \alpha_2 \dots$ If we take the derivative of this once more
with respect to $x$, it follows that

\weq{2}{S\frac{1}{(x-\alpha)^2} = \frac{f'(x)^2 - f(x)f''(x)}{f(x)^2},}

and if on the right-hand side we introduce the function $H$ [\S.\,106, (9)],

\weq{3}{S\frac{1}{(x-\alpha)^2} = \frac{f'(x)^2 - H(x)}{nf(x)^2}.}

If one applies to this formula a linear transformation under which
the value $x = \infty$ corresponds to an arbitrary other real value of the new
variable, one obtains a more general formula,
which we shall first derive. We introduce homogeneous variables
by replacing $x$ by $x:y$ and $\alpha$ by $\alpha:\beta$.
Thereby (3) becomes

\weq{4}{S\frac{\beta^2}{(x\beta - y\alpha)^2} = \frac{f'(x)^2 - y^2H(x,y)}{nf(x,y)^2};}

now we apply an arbitrary linear transformation:

\weq{5}{\begin{aligned}
x &= ax' + by', & \alpha &= a\alpha' + b\beta', \\
y &= cx' + dy', & \beta &= c\alpha' + d\beta',
\end{aligned}}

or, solved:

\weq{6}{\begin{aligned}
rx' &= dx - by, \\
ry' &= -cx + ay, & r &= ad - bc.
\end{aligned}}

If by this transformation
\[
f(x,y) = \varphi(x', y')
\]
then
\[
rf'(x) = d\varphi'(x') - c\varphi'(y'),
\]
and, on account of the covariant property of the function $H$ (\S.\,59),
\[
r^2 H(x,y) = H'(x', y'),
\]
where $H'$ is derived from $\varphi$ in the same way as $H$ is from $f$. From this
it follows from (4):

\weq{7}{S\frac{(c\alpha' + d\beta')^2}{(x'\beta' - y'\alpha')^2} = \frac{[d\varphi'(x') - c\varphi'(y')]^2 - (cx' + dy')^2 H'(x', y')}{n\varphi(x', y')^2}.}

In this formula, since this is only a matter of notation, the accents
may be omitted at $x'$, $y'$, $\alpha'$, $\beta'$, $H'$, and $f$ may be
%CONT>>


% ===== scan leaf 349, printed 324 =====
% !PRINTED 324
%<<CONT
be put in place of $\varphi$. The two arbitrary quantities $c$ and $d$
remain in it.

If, for the sake of a more elegant notation, we further set $c = -\eta$,
$d = \xi$, we thus obtain:

\weq{8}{S\left(\frac{\xi\beta - \eta\alpha}{x\beta - y\alpha}\right)^2 = \frac{[\xi f'(x) + \eta f'(y)]^2 - (x\eta - \xi y)^2 H(x,y)}{nf(x,y)^2},}

and in this form both the right-hand and the left-hand side remain
completely unchanged if we apply an arbitrary linear transformation
simultaneously to the pairs of variables $x, y$; $\xi, \eta$; $\alpha, \beta$.

The formula (8) is an identity; it holds for every function
$f(x, y)$ and for all values of the variables $x, y$; $\xi, \eta$. But now
we make the assumption that $f(x, y)$ has only real roots, that is,
that the $\alpha, \beta$ are real, and we also suppose $\xi, \eta$ and
$x, y$ to be real. Both sides of the equation (8) are then
essentially positive. We denote their common value by
$P$, and now determine a quantity $X:Y$ by the quadratic
equation
\[
P = \left(\frac{X\eta - Y\xi}{Xy - Yx}\right)^2,
\]
or

\weq{9}{\Phi = (Xy - Yx)^2 P - (X\eta - Y\xi)^2 = 0.}

This equation has two real roots, which are obtained from
\[
X\eta - Y\xi = \pm \sqrt{P}\,(Xy - Yx)
\]
About the position of the roots of this equation we can, however,
assert the following:

If we set $X = x$, $Y = y$, then $\Phi$ becomes negative.

But if we set $X = \alpha$, $Y = \beta$, where $\alpha:\beta$ is any one
of the roots of $f = 0$, then
\[
\left(\frac{X\eta - Y\xi}{Xy - Yx}\right)^2 = \left(\frac{\xi\beta - \eta\alpha}{x\beta - y\alpha}\right)^2,
\]
that is, equal to \gesperrt{one term} of the sum
\[
P = S\left(\frac{\xi\beta - \eta\alpha}{x\beta - y\alpha}\right)^2
\]
and therefore smaller than $P$. From this it follows that $\Phi$ \gesperrt{is
positive} for $X:Y = \alpha:\beta$.

If we now arrange the values of all the $\alpha:\beta$ and $x:y$,
ordered cyclically by magnitude, say on a circle,
%CONT>>


% ===== scan leaf 350, printed 325 =====
% !PRINTED 325
%<<CONT
and represent the value $X:Y$ by a variable point of this circle,
and denote the points in question by $\alpha, x, X$, then
$\Phi$ passes through zero when $X$ goes from $x$ forwards or
backwards to the nearest points, which we shall denote by $\alpha_1$ and
$\alpha_2$. If then $X_1$, $X_2$ are the roots of
equation (9), we have the following arrangement on the circle

\weq{10}{\alpha_1,\ X_1,\ x,\ X_2,\ \alpha_2,}

\gesperrt{and $(X_1, X_2)$ represents an interval in which the arbitrarily
chosen point $x$ lies, but no root point of the equation $f = 0$.
$X_1$ and $X_2$ are therefore nearer to the two nearest roots
$\alpha_1, \alpha_2$ than the initial value $x$.}

This holds whatever position the second arbitrary point $\xi$
may have. The points $X_1, X_2$ will indeed depend
on it; but they will always lie in the interval $(\alpha_1, x)$ and $(x, \alpha_2)$.

There now arises the question: how is the point $\xi$ to be chosen so that
$X_1$ may lie as near as possible to $\alpha_1$, or $X_2$ as near as possible to $\alpha_2$;
or rather, since the knowledge of $\xi$ itself does not matter,
which is the point $X_1$ lying as near as possible to $\alpha_1$, and the
point $X_2$ lying as near as possible to $\alpha_2$?

If we think of the point $X$ as given, then (9) is a
quadratic equation for the point $\xi$; to every position of $X$
there thus belong two positions of $\xi$. But only such values of
$X_1$ or $X_2$ are admissible in (10) for which this quadratic
equation yields real roots. Hence only such values of
$X_1$ and $X_2$ can occur for which the discriminant of
(9) with respect to $\xi$ is positive, and we obtain the limiting positions
of $X_1$ and $X_2$ when we set the discriminant of (9) equal to
zero. But the coefficients of $\xi^2$, $2\xi\eta$, $\eta^2$ in (9) are
\[
\frac{f'(x)^2 - y^2 H}{nf^2}(Xy - Yx)^2 - Y^2
\]
\[
\frac{f'(x)f'(y) + xyH}{nf^2}(Xy - Yx)^2 + XY
\]
\[
\frac{f'(y)^2 - x^2 H}{nf^2}(Xy - Yx)^2 - X^2,
\]
and the discriminant is obtained when from the square
%CONT>>


% ===== scan leaf 351, printed 326 =====
% !PRINTED 326
%<<CONT
the product of the two outer ones from that of the middle one, namely,
using the relation $xf'(x) + yf'(y) = nf$:
\[
\frac{(Xy - Yx)^2}{nf^2}\left\{(n-1)H(Xy - Yx)^2 + [Xf'(x) + Yf'(y)]^2\right\},
\]
and one therefore obtains the limiting values of $X:Y$ from the equation

\weq{11}{[Xf'(x) + Yf'(y)]^2 + (n-1)(Xy - Yx)^2 H = 0,}

which is quadratic in $X:Y$.

Since $H$ is negative, this equation has two real roots,
which one finds from

\weq{12}{Xf'(x) + Yf'(y) = \pm(Xy - Yx)\sqrt{(1-n)H}}

.

This formula comprises several particular results, which one
obtains on abandoning the homogeneous form.

Let us first take $y = 0$, so that $x:y$ is infinite, and
set $Y = 1$; then there result two values $X_1$, $X_2$, between
which the totality of all the roots of $f(x)$ is contained.
If we set
\[
f(x, y) = a_0 x^n + a_1 x^{n-1} y + a_2 x^{n-2} y^2 + \cdots,
\]
then

\weq{13}{H = [2na_0a_2 - (n-1)a_1^2]x^{2n-4},}

and one obtains these two values from

\weq{14}{na_0 X = -a_1 \pm \sqrt{(n-1)^2a_1^2 - 2n(n-1)a_0a_2}.}

If one leaves $x:y$ undetermined, but returns to the inhomogeneous
form by setting $y = 1$, $Y = 1$ and applying formulas
(7) and (9) of \S.\,106, one obtains from (12)

\weq{15}{x - X = \frac{nf(x)}{f'(x) \pm \sqrt{(n-1)^2f'(x)^2 - n(n-1)f(x)f''(x)}},}

and these values have the following meaning:

If $x$ lies between two roots $\alpha_1, \alpha_2$ of $f(x)$, then of
the two values of $X$ represented by (15), one lies
nearer to $\alpha_1$, the other nearer to $\alpha_2$, but both still within
the interval $(\alpha_1, \alpha_2)$.

But if $x$ is greater than the greatest, or less than the
least root of $f(x)$, then all the roots lie
between the two values of $X$, while $x$ lies outside.


% ===== scan leaf 352, printed 327 =====
% !PRINTED 327
\weberpage{327}

%<<CONT
If $x$ lies near a root, then one of the two values
of $X$ gives a still more accurate value of this root.

\begin{center}
\rule{2cm}{0.4pt}
\end{center}

Just as an interval is determined by $X_1, X_2$ in which
no root of the equation lies at all, so one can also determine an
interval in which roots certainly lie. Here too
one best starts from the homogeneous form and makes use of
the linear transformation.

Let $f(x, y)$ then be, as before, a form of degree $n$ with none but
real roots; and let further $x:y$, $\xi:\eta$, $\xi':\eta'$ be three real
values, of which two are arbitrary, and of which the
third, through the relation

\weq{16}{\begin{split}
{}[\xi f'(x) + \eta f'(y)]\,[\xi' f'(x) + \eta' f'(y)] \\
{}- H(\xi y - \eta x)\,(\xi' y - \eta' x) = 0
\end{split}}

depends on the two others. The equation (16) remains
unchanged if we subject the three pairs of quantities $x, y$; $\xi, \eta$; $\xi', \eta'$
simultaneously to the same linear transformation. Hence
all conclusions that we draw from a special form of this
equation can be generalised. The linear substitution
can be so determined that for the new variables $y = 0$,
$\xi = 0$; then (16) yields for $\xi'$, $\eta'$ the equation

\[
a_1(n\,a_0\,\xi' + a_1\,\eta') = [2\,n\,a_0\,a_2 - (n-1)\,a_1^2]\eta',
\]

or, if one assumes $a_0 = 1$, $\eta' = 1$,

\[
\xi' = -\frac{a_1^2 - 2\,a_2}{a_1}.
\]

But if, as above, we denote by $\alpha$ the roots of $f$,
we have

\[
a_1^2 - 2\,a_2 = S(\alpha^2), \quad a_1 = -S(\alpha);
\]

and thus

\[
\xi' = \frac{S(\alpha^2)}{S(\alpha)}.
\]

For this we may also write:

\[
S[\alpha(\alpha - \xi')] = 0,
\]

and from this it follows, since all the $\alpha$ are to be real, that the
products $\alpha(\alpha - \xi')$ can be neither all positive nor all negative,
or in other words, one part of the roots must
%CONT>>


% ===== scan leaf 353, printed 328 =====
% !PRINTED 328
\weberpage{328}
%<<CONT
between $0$ and $\xi'$, and another part outside this interval.

From this it now follows generally, by applying the linear transformation, \gesperrt{that on the circle carrying the points $\alpha, x, \xi, \xi'$, if $\xi, \xi'$ are connected by (16), at least one of the roots $\alpha$ lies in each of the two parts into which the circle is divided by $\xi$ and $\xi'$.} The point $x$ and one of the points $\xi, \xi'$ are arbitrary; the third point is determined by (16).

We shall in particular make the assumption that the point $\xi'$ coincides with one of the two points $X$ determined by the equation (12), and thus set, say,
\weq{17}{
X f'(x) + Y f'(y) + (Xy - Yx)\sqrt{(1-n)H} = 0
}

If we introduce this into (16), by setting $\xi', \eta' = X, Y$ and then removing the factor $Xy - Yx$, we obtain
\weq{18}{
(n-1)\,[\xi f'(x) + \eta f'(y)] + (\xi y - \eta x)\sqrt{(1-n)H} = 0,
}
by which the point $\xi$ is uniquely determined.

If we pass to the inhomogeneous form, we obtain [\S.\,106, (7) and (9)] from (17) and (18)
\weq{19}{
x - X = \dfrac{n f(x)}{f'(x) + (n-1)\sqrt{\displaystyle f'(x)^2 - \frac{n}{n-1} f(x) f''(x)}},
}
\weq{20}{
x - \xi = \dfrac{n f(x)}{f'(x) + \sqrt{\displaystyle f'(x)^2 - \frac{n}{n-1} f(x) f''(x)}},
}
in which the square root can have either sign.

The position of the points can then be characterised as follows. There is at least one root $\alpha$ such that
\weq{21}{
\xi,\ \alpha,\ X,\ x,
}
follow one another in order of magnitude, and such that no root lies between $x$ and $X$. We can also choose $\alpha$ so that no further root lies between $\alpha$ and $X$. Between $\alpha$ and $\xi$, however, other roots may possibly lie.


% ===== scan leaf 354, printed 329 =====
% !PRINTED 329
\weberpage{329}

But if we now let $x$ approach the root $\alpha$, then $X$ will approach the same root, and (20) shows that $\xi$ also tends to the same limit. It follows, therefore, that if $x$ lies sufficiently near a root, this nearest root lies in the interval $(\xi, X)$, in which no second root is contained, and that this interval closes more and more about $\alpha$.

As an example we take (with \gesperrt{Laguerre}) the spherical functions, of which we have already shown earlier (\S.\,87) that they have only real roots, which all lie between $-1$ and $+1$. The function
\[
P_n(x) = \frac{1\,.\,3 \ldots (2\,n-1)}{1\,.\,2 \ldots n} \left( x^n - \frac{n(n-1)}{2(2\,n-1)}\,x^{n-2} + \cdots \right)
\]
has the property that
\[
P_n(1) = 1, \qquad P_n'(1) = \frac{n(n+1)}{2},
\]
\[
P_n''(1) = \frac{n(n-1)(n+1)(n+2)}{8}
\]
holds. One proves this easily by complete induction, by verifying the formulas directly in the first cases $n = 2, 3 \ldots$, and then, from the recursion formula
\[
(1 - x^2)\,P_n'(x) + n\,x\,P_n(x) - n\,P_{n-1}(x) = 0
\]
by differentiating twice, establishing their correctness for the index $n$, provided it is assumed for the index $n-1$.

If we then set $x = 1$ in (19) and (20), we obtain two limits between which the root coming nearest to 1 lies, namely
\[
X = 1 - \frac{2}{n + 1 + (n-1)\sqrt{\dfrac{n(n+1)}{2}}},
\]
\[
\xi = 1 - \frac{2}{n + 1 + \sqrt{\dfrac{n(n+1)}{2}}}.
\]
For $n = 7$, for example, this gives
\[
X = 0{,}94967 \ldots
\]
\[
\xi = 0{,}84952 \ldots
\]


% ===== scan leaf 355, printed 330 =====
% !PRINTED 330
\weberpage{330}

The more exact values of the two roots lying nearest to 1, which we cite for the sake of comparison, are, according to the computation of \gesperrt{Gauss} (Gauss Werke, Bd. III, S. 195):
\[
0,94911 \ldots; \quad 0,74153 \ldots
\]
Thus only one root lies here between $X$ and $\xi$, and it is only in the fourth decimal place that $X$ shows the excess over the value of this root.


% ===== scan leaf 356, printed 331 =====
\weberpage{331}% page number not printed on this page
\pdfbookmark[1]{Tenth Section. Approximate Calculation of the Roots.}{bm120}
\pdfbookmark[2]{§ 108. Interpolation. Regula falsi}{bm121}
% !PRINTED none
% [sic] no printed folio number visible on this section-opening page (Tenth Section); page contains the section title, the \S.\,108 heading, and running body text, but no running head/number is legible in any of the three bands.

\begin{center}
\gesperrt{Tenth Section.}

\textbf{Approximate Calculation of the Roots.}

\rule{3cm}{0.4pt}

\S.\ 108.

\gesperrt{Interpolation. Regula falsi.}
\end{center}

With the delimitation of the intervals in each of which only one root of an algebraic equation lies, the possibility is given of calculating the real roots numerically with arbitrary accuracy, by narrowing the interval more and more, e.g.\ by repeated halving. If one divides an interval of size 1 repeatedly into ten parts, then each new step yields a further decimal place.

Once a root has first been separated from all the others, then in the further narrowing of the interval one has only to take into account the sign of the function $f(x)$ itself. The calculations needed for this can be greatly facilitated by the application of the interpolation formula which we came to know in the first section.

Namely, in \S.\,10, (7) we derived a formula by which one can determine a function of degree $n$, which we shall now denote by $\varphi(\xi)$, when $n+1$ successive values $\varphi(0), \varphi(1) \ldots \varphi(n)$ are given. This formula is, if
\weq{1}{\begin{aligned}
B_\nu^{(\xi)} &= \frac{\xi(\xi-1)\ldots(\xi-\nu+1)}{1.2.3\ldots\nu}\\
\Delta_\xi &= \varphi(\xi+1) - \varphi(\xi),
\end{aligned}}
%CONT>>


% ===== scan leaf 357, printed 332 =====
% !PRINTED 332
\weberpage{332}

%<<CONT
\weq{2}{\begin{aligned}
\Delta'_\xi &= \Delta_{\xi+1} - \Delta_\xi. \\
&\cdot\ \cdot\ \cdot\ \cdot\ \cdot\ \cdot\ \cdot\ \cdot\ \cdot\ \cdot \\
\Delta_\xi^{(n-1)} &= \Delta_{\xi+1}^{(n-2)} - \Delta_\xi^{(n-2)}
\end{aligned}}

is set,

\weq{3}{\varphi(\xi) = \varphi(0) + \Delta_0 B_1^{(\xi)} + \Delta'_0 B_2^{(\xi)} + \cdots + \Delta_0^{(n-1)} B_n^{(\xi)}.}

The value of the formula lies in this, that the sequence of differences $\Delta_0, \Delta'_0, \Delta''_0 \ldots$ in many cases presents numerical values decreasing so rapidly that one can content oneself with a few of the first terms of the series (3).

Now if $f(x)$ is the function to be investigated and $(\alpha, \beta)$ the interval in which one of the sought roots lies, then, understanding by $m$ any suitable integer, we set
\[
\delta = \frac{\beta - \alpha}{m}; \qquad x = \alpha + \delta \xi,
\]
and if we accordingly set
\[
\begin{aligned}
\Delta_x &= f(x + \delta) - f(x) \\
\Delta'_x &= \Delta_{x+\delta} - \Delta_x \\
&\cdot\ \cdot\ \cdot\ \cdot\ \cdot\ \cdot\ \cdot\ \cdot
\end{aligned}
\]
then

\weq{4}{f(x) = f(\alpha) + \frac{\Delta_\alpha}{\delta}(x - \alpha) + \frac{\Delta'_\alpha}{\delta^2}\frac{(x-\alpha)(x-\alpha-\delta)}{2} + \cdots}

One could just as well start from the other end point $\beta$ of the interval, and would then have to set in (3)
\[
x = \beta - \delta \xi
\]
One would then obtain
\[
\begin{aligned}
\Delta_x &= f(x - \delta) - f(x) \\
\Delta'_x &= \Delta_{x-\delta} - \Delta_x \\
&\cdot\ \cdot\ \cdot\ \cdot\ \cdot\ \cdot\ \cdot\ \cdot
\end{aligned}
\]

\weq{5}{f(x) = f(\beta) + \frac{\Delta_\beta}{\delta}(x - \beta) + \frac{\Delta'_\beta}{\delta^2}\frac{(x-\beta)(x-\beta+\delta)}{2} + \cdots}

If one interprets the equation
\[
y = f(x)
\]
as the equation of a curve in a rectangular coordinate system $x, y$, then the relations can be illustrated geometrically.

If, for example, one applies equation (4) to the interval from $\alpha$ to $(\alpha + \delta)$ and contents oneself with taking account of the first difference, that means, in the language of geometry,
%CONT>>


% ===== scan leaf 358, printed 333 =====
% !PRINTED 333
\weberpage{333}
%<<CONT
that the arc of the curve running between the two curve points $\alpha, f(\alpha)$ and
$\alpha + \delta, f(\alpha + \delta)$ is replaced by the chord. If the second difference is
also taken into account, then the arc of the curve passing through the three successive
points with the abscissae $\alpha, \alpha + \delta,
\alpha + 2\delta$ is replaced by the arc of a parabola
which passes through the same points and whose axis is parallel to the axis of ordinates;
and this parabola will fit the curve still more closely than the chord.
The smaller the interval $\delta$ is, the less influence the higher
differences will have.

How far one has to go in the approximation, then, depends not only on the
accuracy with which one wishes to know $f(x)$, but essentially also
on the density of the values $\alpha, \alpha + \delta, \alpha + 2\delta
\ldots$ for which the function is known. On these theorems rests the
arrangement of our works of tables, especially of the tables of logarithms. There, to be
sure, we are not dealing with integral rational functions; but for
continuous functions in general the same laws hold here. In the tables
one therefore also finds, beside the values $f(\alpha), f(\alpha + \delta),
f(\alpha + 2\delta) \ldots$, the values of the first difference, or of the first two
differences, given. For the customary seven-place tables the
first difference suffices. In the ten-place table ``Thesaurus
logarithmorum'' of \gesperrt{Vega} the second differences are also
given, and must be taken into account in quite precise computations.

Our interpolation formulae can be applied with advantage to compute the roots
of equations, or, more exactly expressed, to improve the approximate values
found by other means.

We can formulate the problem thus: for a given value of $f(x)$ lying between
$f(\alpha)$ and $f(\alpha + \delta)$, the
corresponding value of $x$ is to be found.

We regard $x = \alpha$ as a first approximate value. If we now set
\[
f(x) - f(\alpha) = \varDelta,
\]
then $\varDelta$ is a given value of the same sign as
$\varDelta_\alpha = f(\alpha + \delta) - f(\alpha)$ and absolutely smaller than
$\varDelta_\alpha$. If we further set $x - \alpha = u$, formula (4) gives
\weq{6}{\varDelta = \frac{\varDelta_\alpha u}{\delta} + \frac{\varDelta'_\alpha}{\delta^2}\,\frac{u(u-\delta)}{2} + \cdots}


% ===== scan leaf 359, printed 334 =====
% !PRINTED 334
\weberpage{334}
If we stop, to begin with, at the first difference, we obtain as
\gesperrt{first correction}
\weq{7}{ u = \delta\frac{\varDelta}{\varDelta_\alpha}, }
and if we now set
\[ u = \delta\frac{\varDelta}{\varDelta_\alpha} + u' \]
then it follows from (6), if one omits $u'$ in the second term, that
\[ 0 = \frac{\varDelta_\alpha}{\delta}u' +
  \frac{\varDelta'_\alpha \varDelta(\varDelta-\varDelta_\alpha)}{2\varDelta_\alpha^2}, \]
and hence as second correction:
\weq{8}{ u' = \delta\frac{\varDelta'_\alpha
  \varDelta(\varDelta_\alpha-\varDelta)}{2\varDelta_\alpha^3}. }

The first correction is obtained by replacing the arc of the curve
running between $\alpha$ and $\alpha+\delta$ by the chord, as above;
the second by replacing the arc between $\alpha,
\alpha+\delta, \alpha+2\delta$ by a parabola which passes through
the same points, but which now has its \gesperrt{axis parallel to the $x$-axis}.

Let us take as an example the equation
\[ f(x) = x^3 - 2x - 2 = 0, \]
which has a real root between 1,7 and 1,8.

One computes
\[
\begin{array}{llll}
x = 1,7, & f(x) = -0,487, & \varDelta_\alpha = 0,719, & \varDelta'_\alpha = 0,108,\\
\phantom{x = {}}1,8, & \phantom{f(x) {}}{=} +0,232, & \phantom{\varDelta_\alpha {}}{=} 0,827, & \\
\phantom{x = {}}1,9, & & & \phantom{\varDelta'_\alpha {}}{=} 1,059.
\end{array}
\]
Since $f(x) = 0$ is to hold, we must set
\[ \varDelta = 0,487 \]
and the first correction to $\alpha = 1,7$ is, by (7),
\[ u = 0,06773, \]
the second correction gives, by (8),
\[ u' = 0,0164, \]
hence
\[ \alpha + u + u' = 1,76937. \]

The more accurate value, computed in another way, is 1,76929\ldots\ We
thus have a result accurate in the first three decimals. But we have
here no other means of judging the accu-
%CONT>>


% ===== scan leaf 360, printed 335 =====
% !PRINTED 335
\weberpage{335}
\pdfbookmark[2]{§ 109. Newton's method of approximation}{bm122}
%<<CONT
racy in advance than the decrease of the differences
$\varDelta, \varDelta', \varDelta''\ldots$ If $\varDelta'$ is so small
that it falls outside the limits of the intended accuracy, then taking
the first difference into account gives an exact result. The rule for
the computation of roots set forth here is called the \gesperrt{Regula falsi}.

The simplest form of this rule, suited to the first approximation, is
this:

If a root $x$ of the equation $f(x)=0$ lies between $\alpha$ and
$\beta$, and if $f(\alpha) = -a, f(\beta) = b$, then
\weq{9}{ x = \frac{b\alpha + a\beta}{a+b} }
is an approximate value of $x$. This is only another way of writing
formula (7).

\begin{center}
§. 109.\\
\gesperrt{The Newton method of approximation.}
\end{center}

A method which, for the approximate computation of the roots of an
equation, is usually better than interpolation is due to
\gesperrt{Newton} and was developed and more closely investigated by
\gesperrt{Fourier}.

The method consists simply in the following. Let
\weq{1}{ f(x) = 0 }
be the equation to be solved, and let a value $x = \alpha$ have been
found which may be regarded as a certain approximation to a root. We
set in (1)
\[ x = \alpha + h, \]
and obtain
\weq{2}{ f(\alpha+h) = f(\alpha) + hf'(\alpha) +
  \frac{h^2}{1.2}f''(\alpha) + \cdots }
If one now determines $h$ from the equation
\weq{3}{ f(\alpha) + hf'(\alpha) = 0, }
which presupposes that $f'(\alpha)$ is different from zero, then
\[ f(\alpha+h) = \frac{h^2}{1.2}f''(\alpha) + \cdots, \]
and thus becomes, if $h$ is a small number, since only the square
%CONT>>


% ===== scan leaf 361, printed 336 =====
% !PRINTED 336
\weberpage{336}
%<<CONT
and higher powers of $h$ occur, has a small value. Under suitable
hypotheses, then,
\weq{4}{ \alpha' = \alpha - \frac{f(\alpha)}{f'(\alpha)} }
is to be regarded as a better approximation to the true value of the
root.

If one then replaces $\alpha$ by $\alpha'$, one will obtain in
\[ \alpha'' = \alpha' - \frac{f(\alpha')}{f'(\alpha')} \]
a still better approximation, and so on.

But there still remain here the following two questions to answer:
\begin{enumerate}
\item[1.] \gesperrt{Under what hypotheses is $\alpha'$
  really a better value than $\alpha$?}
\item[2.] \gesperrt{How can one estimate the degree of accuracy that
  one attains in this way?}
\end{enumerate}
These questions were answered by \gesperrt{Fourier}; but in doing so he
makes the following assumption:

\gesperrt{A root of $f(x)$ is enclosed in an interval $(\alpha,\beta)$
that contains no second root.}

\gesperrt{In the interval $(\alpha,\beta)$, $f'(x)$ is different from zero and
$f''(x)$ is also different from zero.}

As regards the latter assumption, that $f''(x)$ is different from zero
in the interval, it is made only in order to obtain conditions for the
application of the method that can be expressed more simply. In itself
the usefulness of the method, given sufficient narrowing of the
interval, does not depend on it. But if $f(x)$ and $f''(x)$ have no
common divisor, and so do not vanish simultaneously, then by narrowing
the interval one can always satisfy the \gesperrt{Fourier} assumption;
and if $f(x)$ and $f''(x)$ do have a common divisor, one can first
separate it off and then apply the approximation methods to the
individual factors of $f(x)$.

One surveys the situation most simply in geometry, by interpreting
$y=f(x)$ as the equation of a plane curve in a rectangular coordinate
system.

The equation
\weq{3}{ y = f(\alpha) + (x-\alpha)f'(\alpha) } % [sic] equation number (3) repeats the earlier (3) used on the previous printed page for f(alpha)+hf'(alpha)=0; reproduced as printed, Weber restarts the numbering for this geometric sub-passage
is the equation of the tangent to the curve at the point $\alpha, f(\alpha)$
%CONT>>


% ===== scan leaf 362, printed 337 =====
% !PRINTED 337
\weberpage{337}

and the \gesperrt{Newton approximation method} therefore comes down to this,
that in a first approximation one replaces the curve in the interval
$(\alpha, \beta)$ by the tangent at one of the endpoints, instead of
by its chord, as in the interpolation method.

% SOURCE FIGURE CROP: Bd. I registered scan 363 (printed p. 337), 600 dpi,
% crop x=180,y=900,w=2650,h=950; source-faithful raster, Fig. 21--22.
\weberfigureimage{weber-b1-fig21-22.png}{0.94}

If the second differential quotient vanishes in the interval $(\alpha, \beta)$,
so that the curve has a point of inflection, then the case can
occur that both end tangents lead out of the interval;
in that case the \gesperrt{Newton} method is accordingly not ap-

% SOURCE FIGURE CROP: Bd. I registered scan 363 (printed p. 337), 600 dpi,
% crop x=180,y=2280,w=2650,h=1080; source-faithful raster, Fig. 23--24.
\weberfigureimage{weber-b1-fig23-24.png}{0.94}

plicable (Fig. 21). But it can also happen, when the interval $(\alpha, \beta)$
has already been sufficiently narrowed, that the other case occurs,
in which both tangents lead to interior points of the interval
(Fig. 22). In these cases, however, the regula
falsi will always give a better approximation.

But we shall now assume that $f'(x)$ and $f''(x)$ do not
vanish in the interval $(\alpha, \beta)$. Then the curve does not change
the sense of its curvature, and then certainly \gesperrt{one of the
two end tangents} always has its point of intersection with the
$x$-axis in the interior of the interval. This certainly holds when
the tangent is taken at that endpoint of the interval,

% Weber, Algebra. I.        22
%CONT>>


% ===== scan leaf 363, printed 338 =====
% !PRINTED 338
\weberpage{338}
%<<CONT
in which $f(x)$ and $f''(x)$ have \gesperrt{the same sign}. The
two Figs. 23 and 24 illustrate the situation for positive
$f'(x)$ and positive and negative $f''(x)$.

Thus in the first case $\beta'$, and in the second $\alpha'$, is a better
approximate value than $\beta$ and $\alpha$.

If one wishes at the same time to shift the other limit, so as
to obtain a new narrower interval, one may, following
Fourier, draw from the other end point — that is, in the case of Fig. 23,
at $a$ — the straight line parallel to $b\beta'$, and one obtains the
point $\alpha'$ as the lower limit point of the new interval. Then
in the first case
\[
\alpha' = \alpha - \frac{f(\alpha)}{f'(\beta)}, \quad \beta' = \beta - \frac{f(\beta)}{f'(\beta)},
\]
in the second case
\[
\alpha' = \alpha - \frac{f(\alpha)}{f'(\alpha)}, \quad \beta' = \beta - \frac{f(\beta)}{f'(\alpha)},
\]
and $(\alpha', \beta')$ is a narrower interval in which the sought root lies.

% SOURCE FIGURE CROP: Bd. I registered scan 364 (printed p. 338), 600 dpi,
% crop x=250,y=2250,w=1200,h=1150; source-faithful raster, Fig. 25.
\weberfigureimage{weber-b1-fig25.png}{0.42}

It is hardly necessary to consider
separately the two other cases, in
which $f'(x)$ is negative in the
interval.

But one may also, in order to obtain
two new limits, combine the Newton
method with the interpolation method,
with still better success, as Fig. 25 shows.

One then obtains as the two
new limits:
\[
\begin{gathered}
\alpha' = \alpha - \frac{f(\alpha)(\beta - \alpha)}{f(\beta) - f(\alpha)} = \frac{\alpha f(\beta) - \beta f(\alpha)}{f(\beta) - f(\alpha)}, \\
\beta' = \beta - \frac{f(\beta)}{f'(\beta)}.
\end{gathered}
\]

We summarise the results in the following theorem:

\gesperrt{If an interval $(\alpha, \beta)$ has been marked off, in which
$f(x)$ changes its sign once, and $f'(x)$ and $f''(x)$ not
at all, and if $\beta$ is the end point of the interval
in which $f(\beta)$ and $f''(\beta)$ have the same sign,
no matter whether $\beta$ is smaller or greater than}
%CONT>>


% ===== scan leaf 364, printed 339 =====
% !PRINTED 339
\weberpage{339}
%<<CONT
\gesperrt{$\alpha$, then one obtains a narrower interval $(\alpha', \beta')$
with the same properties, if one}
\weq{4}{
\alpha' = \alpha - \frac{f(\alpha)(\beta - \alpha)}{f(\beta) - f(\alpha)}, \quad \beta' = \beta - \frac{f(\beta)}{f'(\beta)}
}
\gesperrt{sets.}

To prove these results of geometric intuition analytically,
we assume that in the interval $f'(x)$ and
$f''(x)$ are positive, hence $\alpha < \beta$, and
\[
f(\alpha) < 0, \quad f(\beta) > 0.
\]

We restrict ourselves to the consideration of this one
case, since the three others are to be treated in exactly the same
way.

We form the function
\[
\varphi(x) = \frac{f(\beta) - f(x)}{\beta - x}
\]
and its first derivative
\[
\varphi'(x) = \frac{-(\beta - x)f'(x) + f(\beta) - f(x)}{(\beta - x)^2}.
\]
The numerator of this expression
\[
\psi(x) = -(\beta - x)f'(x) + f(\beta) - f(x)
\]
vanishes for $x = \beta$, and its derivative is
\[
\psi'(x) = -(\beta - x)f''(x),
\]
hence negative in the interval. From this it follows that $\psi(x)$ decreases
in the interval as $x$ increases, and therefore, because it vanishes for the
greatest value $x = \beta$, remains positive. Consequently
$\varphi'(x)$ is also positive, and $\varphi(x)$ increases steadily in the interval
as $x$ increases. We accordingly have
\weq{5}{
\begin{gathered}
\frac{f(\beta) - f(\alpha)}{\beta - \alpha} < \frac{f(\beta) - f(x)}{\beta - x} < f'(\beta) \\
\alpha < x < \beta.
\end{gathered}
}
If we now set
\[
\alpha' = \alpha - \frac{f(\alpha)(\beta - \alpha)}{f(\beta) - f(\alpha)}, \quad \beta' = \beta - \frac{f(\beta)}{f'(\beta)},
\]
then
\weq{6}{
\beta' - \alpha' = f(\beta)\left\{\frac{\beta - \alpha}{f(\beta) - f(\alpha)} - \frac{1}{f'(\beta)}\right\},
}
and from (5) it follows that this difference is positive, hence
\[
\alpha < \alpha' < \beta' < \beta.
\]
% 22*


% ===== scan leaf 365, printed 340 =====
% !PRINTED 340
\weberpage{340}

If one then sets in (5)
\[
x = \alpha' \quad \text{and} \quad x = \beta'
\]
and observes that
\[
\beta - \alpha' = \frac{(\beta - \alpha)f(\beta)}{f(\beta) - f(\alpha)}, \quad \beta - \beta' = \frac{f(\beta)}{f'(\beta)},
\]
it follows that
\[
f(\alpha') < 0, \quad f(\beta') > 0.
\]

The root sought is therefore contained between $\alpha'$ and $\beta'$,
and the interval $(\alpha', \beta')$ is smaller than $(\alpha, \beta)$.

If we put $\alpha', \beta'$ in place of $\alpha, \beta$ in these expressions,
we obtain a new, still narrower interval, and so on.

As for the convergence of this procedure, we can assure
ourselves of it in the following way. Following (6), we set:
\weq{7}{
\frac{\beta' - \alpha'}{\beta - \alpha} = \Phi(\alpha, \beta) = \frac{f(\beta)\{f(\alpha) - f(\beta) + (\beta - \alpha)f'(\beta)\}}{(\beta - \alpha)f'(\beta)\{f(\beta) - f(\alpha)\}}.
}
Numerator and denominator are here divisible by $(\beta - \alpha)^2$, and
if we cancel the common factor $(\beta - \alpha)^2$ and put $\xi, \eta$
for $\alpha, \beta$, we obtain an expression of the form
\[
\Phi(\xi, \eta) = \frac{\psi_1(\xi, \eta)}{\psi_2(\xi, \eta)},
\]
in which $\psi_1$ and $\psi_2$ are integral rational functions of the two variables
$\xi, \eta$.

If $x$ is the root lying in the interval $(\alpha, \beta)$, then by our
assumptions about $f(x)$ the two functions $\psi_1$ and $\psi_2$
do not become zero when
\weq{8}{
\alpha \leqq \xi < x, \quad x < \eta \leqq \beta.
}
The function $\Phi(\xi, \eta)$ is less than 1 for $\xi = \alpha, \eta = \beta$, but also for
every other pair of values in the interval $(\alpha, \beta)$ for which $f(\xi)$
is negative and $f(\eta)$ positive, since every such interval
$(\xi, \eta)$ could be taken in place of $(\alpha, \beta)$.

Now since $\Phi(\xi, \eta)$ is a continuous function of the two variables
$\xi, \eta$ so long as these remain in the region (8), the
function $\Phi(\xi, \eta)$ must have a maximum value in this region,
and this must be less than 1, because $\Phi(\xi, \eta)$
remains less than 1 even in the limiting cases $\xi = x, \eta = x$.

A positive proper fraction $\Theta$ can therefore be assigned
such that
\[
\Phi(\xi, \eta) < \Theta
\]
%CONT>>


% ===== scan leaf 366, printed 341 =====
% !PRINTED 341
\weberpage{341}
\pdfbookmark[2]{§ 110. Daniel Bernoulli's method of approximation and related methods}{bm123}
%<<CONT
holds, so long as $\xi, \eta$ belong to the domain (8). Then from (7) it follows that
\[
\beta' - \alpha' < (\beta - \alpha)\Theta.
\]
In the same way it follows, when we proceed in the same manner from the interval $(\alpha', \beta')$ to a narrower interval $(\alpha'', \beta'')$, that
\[
\beta'' - \alpha'' < (\beta' - \alpha')\Theta',
\]
in which $\Theta'$ has the same meaning for $\alpha', \beta'$ as $\Theta$ has for $\alpha, \beta$. But since $\alpha', \beta'$ belong to the domain (8), $\Theta'$ is not greater than $\Theta$, and $\Theta$ may also be put in place of $\Theta'$. It follows then:
\[
\beta'' - \alpha'' < (\beta - \alpha)\Theta^2,
\]
and so we conclude further
\[
\beta^{(\nu)} - \alpha^{(\nu)} < (\beta - \alpha)\Theta^\nu.
\]

The intervals therefore decrease at least as strongly as the terms of a decreasing geometric progression.

As an example may serve the equation:
\[
x^3 - 2x^2 - 2 = 0,
\]
which has a root between
\[
\alpha = 2{,}35 \quad \text{and} \quad \beta = 2{,}36
\]

From the formulae (4) one obtains for
\[
\beta' = 2{,}35931\ldots, \quad \alpha' = 2{,}359298\ldots,
\]
so that a value of the root accurate in the fourth decimal is
\[
2{,}3593
\]
For this value itself, as a more exact calculation shows, $f(x)$ is still negative, so that it can be taken for $\alpha'$. The next step of the approximation gives.
\[
2{,}359304.
\]

\begin{center}
\S.\,110.

\gesperrt{The approximation method of Daniel Bernoulli and related methods.}
\end{center}

The method for the approximate solution of an equation which is due to \gesperrt{Daniel Bernoulli} rests on the fact that, if one has a sequence of real quantities, the powers of the greatest among them will outweigh the equally high powers of the rest the more, the higher the powers are.


% ===== scan leaf 367, printed 342 =====
% !PRINTED 342
\weberpage{342}
If $\alpha, \beta, \gamma \ldots$ are arbitrary real or complex quantities, subject however to the condition that the absolute value of $\alpha$ is \gesperrt{greater} than the absolute value of all the others, so that the absolute values of the fractions $\beta:\alpha, \gamma:\alpha \ldots$ are proper fractions, then
\weq{1}{
\frac{\alpha^m + \beta^m + \gamma^m + \cdots}{\alpha^{m-1}+\beta^{m-1}+\gamma^{m-1}+\cdots}
= \alpha\,\frac{1+\left(\dfrac{\beta}{\alpha}\right)^m+\left(\dfrac{\gamma}{\alpha}\right)^m+\cdots}{1+\left(\dfrac{\beta}{\alpha}\right)^{m-1}+\left(\dfrac{\gamma}{\alpha}\right)^{m-1}+\cdots},
}
and the greater $m$ becomes, the more this expression will approach the limit $\alpha$, as the second representation shows.

If $\alpha, \beta, \gamma \ldots$ are the roots of an algebraic equation, then the left side of (1) is the quotient of the $m$th and the $(m-1)$th power sum, and we thus obtain the theorem:

\begin{center}
\gesperrt{The quotient of the $m$th and the $(m-1)$th power sum approaches, as $m$ increases, the absolutely greatest among the roots.}
\end{center}

If one takes $m$ negative, and supposes $\alpha$ absolutely \gesperrt{smaller} than $\beta, \gamma \ldots$, it follows in the same way:

\begin{center}
\gesperrt{The quotient of the $-m$th and the $-(m+1)$th power sum approaches, as $m$ increases, the absolutely smallest among the roots.}
\end{center}

Since the power sums can be computed as symmetric functions from the coefficients, one need only take a sufficiently large $m$ in order to obtain an approximate value of the absolutely greatest and the absolutely smallest root.

But if among the roots of the equation there occur some which have the same absolute value, and this is the greatest or the smallest, then this method is not applicable — thus, for example, not to real equations in which a pair of imaginary roots occurs that surpasses the real ones in absolute value.

But if the greatest root, though it occurs singly, does not much surpass the next following one, then the procedure will yield some accuracy only slowly.

\gesperrt{Jacobi} has supplemented the \gesperrt{Bernoulli} method in one direction\footnote{Observatiunculae etc. Werke, Bd. 3, S. 280.}.


% ===== scan leaf 368, printed 343 =====
% !PRINTED 343
\weberpage{343}
Let us assume that among the roots $x_1, x_2 \ldots x_n$ there are
\[
x_1, x_2 \ldots x_k
\]
whose absolute values are greater than the absolute values of the following $x_{k+1}, x_{k+2} \ldots x_n$, and let us set
\weq{2}{
\begin{array}{rcl}
p_m &=& x_1^m + x_2^m + \cdots + x_k^m\\
p_{m+1} &=& x_1^{m+1}+x_2^{m+1}+\cdots+x_k^{m+1}\\
&\cdots\cdots\cdots\cdots\cdots\cdots\cdots&\\
p_{m+k} &=& x_1^{m+k}+x_2^{m+k}+\cdots+x_k^{m+k},
\end{array}
}
then, if we set
\weq{3}{
(x-x_1)(x-x_2)\ldots(x-x_k) = x^k + A_1 x^{k-1} + \cdots + A_k
}
the coefficients $A_1, A_2 \ldots A_k$ will satisfy the equations
\weq{4}{
\begin{array}{rcl}
p_{m+k} + p_{m+k-1}A_1 + p_{m+k-2}A_2 + \cdots + p_m A_k &=& 0\\
p_{m+k+1} + p_{m+k}A_1 + p_{m+k-1}A_2 + \cdots + p_{m+1}A_k &=& 0\\
&\cdots\cdots\cdots\cdots\cdots\cdots\cdots&\\
p_{m+2k-1}+p_{m+2k-2}A_1+p_{m+2k-3}A_2+\cdots+p_{m+k-1}A_k &=& 0.
\end{array}
}
If one adjoins to the equations (4) the further equation
\[
x^k + A_1 x^{k-1} + A_2 x^{k-2} + \cdots + A_k = 0
\]
and eliminates the $A_1, A_2 \ldots A_k$, there results the equation of degree $k$
\[
\left|
\begin{array}{llll}
x^k, & x^{k-1}, & \cdots & 1\\
p_{m+k}, & p_{m+k-1}, & \cdots & p_m\\
p_{m+k+1}, & p_{m+k}, & \cdots & p_{m+1}\\
\multicolumn{4}{c}{\cdots\cdots\cdots\cdots\cdots\cdots\cdots\cdots\cdots\cdots}\\
p_{m+2k-1}, & p_{m+2k-2}, & \cdots & p_{m+k-1}
\end{array}
\right| = 0,
\]
whose roots are $x_1, x_2 \ldots x_k$. The quantities $p_m, p_{m+1} \ldots$ can, however, be replaced by the corresponding power sums of all the roots, and with all the greater accuracy the greater $m$ is.

Thus one obtains, for example, for $k = 2$:
\[
x^2 + \frac{p_m p_{m+3} - p_{m+1}p_{m+2}}{p_{m+1}^2 - p_m p_{m+2}}\, x + \frac{p_{m+2}^2 - p_{m+1}p_{m+3}}{p_{m+1}^2 - p_m p_{m+2}} = 0,
\]
an equation which one can apply when, in a real equation, a pair of conjugate roots have the absolutely greatest value.

If one applies the \gesperrt{Bernoulli} method to the sum of the negative powers, one obtains, as has already been
%CONT>>


% ===== scan leaf 369, printed 344 =====
% !PRINTED 344
\weberpage{344}
\pdfbookmark[2]{§ 111. Gräffe's method of approximation}{bm124}
%<<CONT
marked, an approximation to the absolutely smallest root. But by shifting the origin one can make any root the absolutely smallest, though for this, admittedly, the knowledge of a value approximated to a certain degree is necessary.

A formula that suffices for all cases has been given by Fr. \gesperrt{Meyer}\footnote{Mathematische Annalen, Bd. 33.}. Let $\alpha_1, \alpha_2, \alpha_3 \ldots \alpha_n$ be the roots of the equation $f(x) = 0$, which may be real or imaginary. One seeks a point $p$ in the $x$-plane such that a circle can be described about $p$ as centre which contains \gesperrt{only one} of the points $\alpha_1, \alpha_2 \ldots \alpha_n$, say the point $\alpha_1$, so that the absolute values of
\weq{5}{
\frac{p-\alpha_1}{p-\alpha_2}, \quad \frac{p-\alpha_1}{p-\alpha_3}, \quad \cdots \quad \frac{p-\alpha_1}{p-\alpha_n}
}
are proper fractions. Such points always exist; if need be, they must be found by the \gesperrt{Sturm} method. If, in addition, one chooses an arbitrary function $\varphi(x)$ that has no common divisor with $f(x)$, e.g. $\varphi(x) = 1$, and forms the symmetric functions
\weq{6}{
y_m = \dfrac{\dfrac{\alpha_1\varphi(\alpha_1)}{(p-\alpha_1)^m} + \dfrac{\alpha_2\varphi(\alpha_2)}{(p-\alpha_2)^m} + \cdots + \dfrac{\alpha_n\varphi(\alpha_n)}{(p-\alpha_n)^m}}{\dfrac{\varphi(\alpha_1)}{(p-\alpha_1)^m} + \dfrac{\varphi(\alpha_2)}{(p-\alpha_2)^m} + \cdots + \dfrac{\varphi(\alpha_n)}{(p-\alpha_n)^m}},
}
then with increasing $m$ these approach the limit $\alpha_1$, and the more rapidly, the smaller the absolute values of the fractions (5) already are.

The correctness of this shows itself at once if one writes the formula (6) in the following way:
\[
y_m = \frac{\alpha_1\varphi(\alpha_1) + \alpha_2\varphi(\alpha_2)\left(\dfrac{p-\alpha_1}{p-\alpha_2}\right)^m + \cdots + \alpha_m\varphi(\alpha_m)\left(\dfrac{p-\alpha_1}{p-\alpha_n}\right)^m}{\varphi(\alpha_1) + \varphi(\alpha_2)\left(\dfrac{p-\alpha_1}{p-\alpha_2}\right)^m + \cdots + \varphi(\alpha_m)\left(\dfrac{p-\alpha_1}{p-\alpha_n}\right)^m}.
\]
% [sic] source prints the final outer subscripts as $m$, although the sequence and the inner $p-\alpha_n$ indicate $n$

\begin{center}
\S.\,111.

\gesperrt{The approximation method of Gräffe.}
\end{center}

On an idea similar to that of the \gesperrt{Bernoulli} approximation there also rests a procedure that \gesperrt{Gräffe}
%CONT>>


% ===== scan leaf 370, printed 345 =====
% !PRINTED 345
\weberpage{345}
%<<CONT
given, and further developed by \gesperrt{Encke}\footnote{Crelle's Journal, Bd.~22 (1841).}, and which lends itself particularly well to a practical execution of the numerical computations.

If, as above, $\alpha, \beta, \gamma \ldots$ are the roots of an equation, and $\alpha$ is the absolutely greatest among them, so that none of the others equals the root $\alpha$ in absolute value, then
\weq{1}{\sqrt[m]{\alpha^m + \beta^m + \gamma^m + \cdots}}
has, as $m$ grows without bound, the value $\alpha$ as its limit; indeed the convergence toward $\alpha$ is even better than for the quotient considered in the preceding section,
\[
\frac{\alpha^m + \beta^m + \gamma^m + \cdots}{\alpha^{m-1} + \beta^{m-1} + \gamma^{m-1} + \cdots},
\]
as one recognises when, according to the generalised binomial theorem, one sets:
\[
\begin{aligned}
\sqrt[m]{\alpha^m + \beta^m + \gamma^m + \cdots} &= \alpha\sqrt[m]{1 + \left(\frac{\beta}{\alpha}\right)^m + \left(\frac{\gamma}{\alpha}\right)^m + \cdots} \\
&= \alpha\left[1 + \frac{1}{m}\left(\frac{\beta}{\alpha}\right)^m + \frac{1}{m}\left(\frac{\gamma}{\alpha}\right)^m + \cdots\right]
\end{aligned}
\]

The absolutely smallest among the roots is obtained by the same principle as the limit of
\weq{2}{\dfrac{1}{\sqrt[m]{\dfrac{1}{\alpha^m} + \dfrac{1}{\beta^m} + \dfrac{1}{\gamma^m} + \cdots}},}
and if $p$ is an arbitrary value, one obtains the root lying nearest to the value $p$ as the limit of
\weq{3}{p + \dfrac{1}{\sqrt[m]{\dfrac{1}{(\alpha-p)^m} + \dfrac{1}{(\beta-p)^m} + \cdots}}.}
If one transforms the equation $f(x)$ by the substitution
\[
y = \frac{1}{x}, \qquad y = \frac{1}{x - p}
\]
then the expressions (2), (3) arise from the expression (1), to which we shall therefore now confine ourselves.


% ===== scan leaf 371, printed 346 =====
% !PRINTED 346
\weberpage{346}
If one applies the formula (1) to a real equation, then $\alpha^m + \beta^m + \gamma^m + \cdots$ is a real number; and if the absolutely greatest root $\alpha$ is real, then the $m$th root is also to be taken real. But if $m$ is an even number, then, in order to decide about the sign, one must further know whether $\alpha$ is positive or negative. If necessary, this is decided by substituting the approximate value found into the given equation.

Likewise, if one applies the formula (1) to an imaginary equation, one will have to decide which of the various $m$th roots gives the correct approximation.

The expressions
\[
s_m = \alpha^m + \beta^m + \gamma^m + \cdots
\]
can be computed especially simply when one takes for $m$ the successive powers of 2, that is $m = 2, 4, 8, 16 \ldots$, and they then yield very good results.

Now $s_2$ is the negative coefficient of the $(n-1)$th power of the unknown in the equation whose roots are the squares of the roots of $f(x)$. This equation, however, is easily formed in the following way.

In $f(x)$ collect the terms with even and with odd powers of $x$, and accordingly set
\[
f(x) = \varphi(x^2) + x\,\psi(x^2).
\]
Then
\[
f_1(x) = \varphi(x)^2 - x\,\psi(x)^2 = 0
\]
is the equation whose roots are the squares of the roots of $f(x)$. If one treats $f_1(x)$ in the same way, one obtains an equation whose roots are the fourth powers of the roots of $f(x)$, and so on.

The carrying out of this computation is very simple, and for the most part leads after a few steps to a good approximation.

We consider some examples.

First we take the equation
\weq{4}{x^3 - 2x - 2 = 0,}
which has one real positive and two imaginary roots. The real root is greater than 1,7, and since the product of all three roots is equal to 2 and $(1,7)^3 > 2$, the absolute value of the two imaginary roots is smaller than the real root. Hence the real root is the absolutely greatest, and the \gesperrt{Gräffe}-
%CONT>>


% ===== scan leaf 372, printed 347 =====
% !PRINTED 347
\weberpage{347}
%<<CONT
method of approximation must lead to it. One now obtains the equations whose roots are the second, fourth, eighth \ldots\ power of equation (4):
\[
\begin{aligned}
x^3 - 4x^2 + 4x - 4 &= 0 \\
x^3 - 8x^2 - 16x - 16 &= 0 \\
x^3 - 96x^2 - 16^2 &= 0
\end{aligned}
\]
and
\[
x = \sqrt[8]{96} = 1{,}7692\ldots
\]
is already a value accurate in the first four decimals. The next step would yield no different result (because the first power of the unknown is missing in the last equation); but the step after that yields the value, accurate in the sixth decimal,
\[
x = \sqrt[32]{85032960} = 1{,}769293\ldots
\]

We shall consider yet a second example of an equation of the fifth degree, one which offers the opportunity to apply several of the theorems of the preceding section. The example is taken, like the earlier ones, from the theory of the complex multiplication of the elliptic functions. It is the equation of the fifth degree
\weq{5}{x^5 - x^3 - 2x^2 - 2x - 1 = 0.}
From this we derive the equation whose roots are the squares of the roots of (5), by substituting, in
\[
(x^5 - x^3 - 2x)^2 - (2x^2 - 1)^2 = 0
\]
$x$ for $x^2$. This gives
\weq{6}{x^5 - 2x^4 - 3x^3 - 1 = 0,}
and if we now apply the same procedure twice more, we obtain the equations whose roots are the fourth and eighth powers of the roots of (5):
\weq{7}{x^5 - 10x^4 + 9x^3 - 4x^2 - 1 = 0,}
\weq{8}{x^5 - 82x^4 + x^3 - 36x^2 - 8x - 1 = 0.}
The last of these equations lends itself to discussion. By the \gesperrt{Cartesian} theorem (\S.\,101) it has at least one positive root.

If, for equation (8), we form the functions occurring in the \gesperrt{Newton} criterion (\S.\,101)
%CONT>>


% ===== scan leaf 373, printed 348 =====
% !PRINTED 348
\weberpage{348}
%<<CONT
\[
1, \quad 4a_1^2 - 10a_0a_2, \quad 6a_2^2 - 12a_1a_3
\]
\[
6a_3^2 - 12a_2a_4, \quad 4a_4^2 - 10a_3a_5, \quad 1
\]
we obtain the signs
\[
+ \quad + \quad - \quad + \quad - \quad +
\]
and (8) therefore has, by \S.\,101, Theorem VIII, four imaginary roots. The positive root is in any case greater than 1, since the left-hand side of (8) is negative for $x = 0$ and $x = 1$.

In order now to decide which of the roots has the greatest absolute value, we shall set $x = 1 : z$ in (8) and then investigate, according to \S.\,96, 97, how many roots lie in the interior of the unit circle in the $z$-plane. If, then, in
\[
z^5 + 8z^4 + 36z^3 - z^2 + 82z - 1
\]
we set $z = \cos\vartheta + i\sin\vartheta$, we obtain the real and imaginary parts:
\[
\varphi = \cos 5\vartheta + 8\cos 4\vartheta + 36\cos 3\vartheta - \cos 2\vartheta + 82\cos\vartheta - 1
\]
\[
\psi = \sin 5\vartheta + 8\sin 4\vartheta + 36\sin 3\vartheta - \sin 2\vartheta + 82\sin\vartheta.
\]
It can now be shown that $\psi$ is positive when $0 < \vartheta < \pi$, and consequently negative when $0 > \vartheta > -\pi$, so that the unit circle is divided into two segments in which $\psi$ has opposite signs. From this it follows, by \S.\,96, that only one root can lie in the interior of the circle, and this is the real root, while the imaginary roots lie outside. It then follows from this, for the equation (8), that the real root has the greatest absolute value, and that it is therefore found by the \gesperrt{Gräffe method}.

In order now to establish this property of the function $\psi$, we set $2\cos\vartheta = \xi$, and apply the goniometric formulae:
\[
\begin{aligned}
\sin 5\vartheta &= \sin\vartheta(\xi^4 - 3\xi^2 + 1) \\
\sin 4\vartheta &= \sin\vartheta(\xi^3 - 2\xi) \\
\sin 3\vartheta &= \sin\vartheta(\xi^2 - 1) \\
\sin 2\vartheta &= \sin\vartheta\,\xi.
\end{aligned}
\]
By this means we obtain
\[
\psi = \sin\vartheta[\xi^4 + \xi^2(8\xi + 33) + (47 - 17\xi)],
\]
from which one sees at once that the factor of $\sin\vartheta$ remains positive so long as $\xi$ lies between $-2$ and $+2$.


% ===== scan leaf 374, printed 349 =====
% !PRINTED 349
\weberpage{349}
\pdfbookmark[2]{§ 112. Trigonometric solution of cubic equations}{bm125}

Equation (8) itself now already gives a fairly good approximate value for the real root

\[
x = \sqrt[8]{82} = 1{,}73471.
\]

A more accurate value, which one obtains by going two steps further, is 1,73469, which is still correct in the fifth decimal place.

This method does not give an enclosure of the sought root between two limits. The criterion for whether a sufficient degree of accuracy has been reached consists in this: that if $s_\mu$, $s_{\mu'}$ are two successive power sums of the roots used in the computation, these are related, with sufficient accuracy, as the $\mu$th to the $\mu'$th power of a quantity, or that approximately

\[
\mu \log s_{\mu'} - \mu' \log s_\mu = 0
\]

holds. One must, however, carry out this test on several successive terms, since in special cases this relation may appear to be satisfied and then cease to hold again in the later computation.

\begin{center}
\S.\ 112.

\gesperrt{Trigonometric Solution of Cubic}

\gesperrt{Equations.}
\end{center}

We now discuss a few further methods of numerical solution applicable to equations of special forms. The aim of these methods is to make the widely available tables of the trigonometric functions and of logarithms useful for the solution of equations. We turn first to the consideration of the cubic equations, which we always assume in the reduced form

\weq{1}{x^3 + a x + b = 0}

in which $a$ and $b$ are real numbers. Since the interchange of $x$ with $-x$ is equivalent to the interchange of $b$ with $-b$, we may restrict ourselves to the assumption that $b$ is negative, and there then remain three different cases to consider. We set $b = -g$ and, according as $a$ is positive or negative, $a = \pm e$. Then we have
%CONT>>


% ===== scan leaf 375, printed 350 =====
% !PRINTED 350
\weberpage{350}
%<<CONT

\[
\begin{aligned}
&\text{a)}\quad x^3 + e x - g = 0\\
&\text{b)}\quad x^3 - e x - g = 0, &&\frac{e^3}{27} < \frac{g^2}{4}\\
&\text{c)}\quad x^3 - e x - g = 0, &&\frac{e^3}{27} > \frac{g^2}{4}.
\end{aligned}
\]

We exclude the limiting cases in which one of the quantities $e$, $g$, $4e^3 - 27g^2$ vanishes.

To apply the \gesperrt{Cardan} formula (\S.\,35), we set

\weq{2}{R = \frac{g^2}{4} \pm \frac{e^3}{27},}

where the upper sign holds in case a), the lower in cases b) and c). In cases a) and b) only one root is real; in case c) (the Casus irreducibilis) all three roots are real. The \gesperrt{Cardan} formula gives

\weq{3}{x = \sqrt[3]{\frac{g}{2} + \sqrt{R}} + \sqrt[3]{\frac{g}{2} - \sqrt{R}}.}

We now proceed in the three cases as follows:

a) We introduce an angle $\vartheta$, which we can choose between $0^0$ and $90^0$, by the equation:

\weq{4}{\frac{g}{2} = \sqrt{\frac{e^3}{27}}\,\mathrm{cotg}\,\vartheta,}

hence

\[
\sqrt{R} = \sqrt{\frac{e^3}{27}}\,\frac{1}{\sin\vartheta},
\]

\[
\begin{aligned}
x &= \sqrt{\frac{e}{3}}\left(\sqrt[3]{\mathrm{cotg}\,\vartheta + \frac{1}{\sin\vartheta}} + \sqrt[3]{\mathrm{cotg}\,\vartheta - \frac{1}{\sin\vartheta}}\right)\\
&= \sqrt{\frac{e}{3}}\left(\sqrt[3]{\mathrm{cotg}\,\frac{\vartheta}{2}} - \sqrt[3]{\mathrm{tg}\,\frac{\vartheta}{2}}\right),
\end{aligned}
\]

and if we further determine an angle $\varphi$ from

\weq{5}{\mathrm{tg}\,\frac{\vartheta}{2} = \mathrm{tg}\,\varphi^3}

then we obtain

\weq{6}{x = 2\sqrt{\frac{e}{3}}\,\mathrm{cotg}\,2\varphi.}

The angles $\vartheta$, $\varphi$ and finally $x$ are found from the logarithmic-trigonometric tables by the formulas (4), (5), (6).


% ===== scan leaf 376, printed 351 =====
% !PRINTED 351
\weberpage{351}

To obtain the imaginary roots, one replaces $\mathrm{tg}\,\varphi$ by $\varrho\,\mathrm{tg}\,\varphi$, where $\varrho$ denotes an imaginary third root of unity.

b) In the second case we determine the angle $\vartheta$, likewise in the first quadrant, from

\weq{7}{\frac{g}{2} = \sqrt{\frac{e^3}{27}}\,\frac{1}{\sin\vartheta}}

\[
\sqrt{R} = \sqrt{\frac{e^3}{27}}\,\mathrm{cotg}\,\vartheta
\]

\[
x = \sqrt{\frac{e}{3}}\left(\sqrt[3]{\mathrm{cotg}\,\frac{\vartheta}{2}} + \sqrt[3]{\mathrm{tg}\,\frac{\vartheta}{2}}\right),
\]

\weq{8}{\mathrm{tg}\,\frac{\vartheta}{2} = \mathrm{tg}\,\varphi^3,}

\weq{9}{x = 2\sqrt{\frac{e}{3}}\,\frac{1}{\sin 2\varphi}.}

c) In the last case, finally, where three real roots are present, we set

\weq{10}{\frac{g}{2} = \sqrt{\frac{e^3}{27}}\cos\vartheta}

\[
\sqrt{R} = i\sqrt{\frac{e^3}{27}}\sin\vartheta
\]

\[
x = \sqrt{\frac{e}{3}}\left(\sqrt[3]{\cos\vartheta + i\sin\vartheta} + \sqrt[3]{\cos\vartheta - i\sin\vartheta}\right),
\]

or, by the theorem of \gesperrt{Moivre}:

\weq{11}{x = 2\sqrt{\frac{e}{3}}\cos\frac{1}{3}\vartheta.}

If one takes $\vartheta$ again in the first quadrant, one obtains for the two other, likewise real, roots

\[
-2\sqrt{\frac{e}{3}}\cos\frac{\pi + \vartheta}{3}, \qquad -2\sqrt{\frac{e}{3}}\cos\frac{\pi - \vartheta}{3}.
\]

All these formulae are arranged for logarithmic computation.


% ===== scan leaf 377, printed 352 =====
% !PRINTED 352
\weberpage{352}
\pdfbookmark[2]{§ 113. Gauss's method for the solution of trinomial equations}{bm126}

\begin{center}
\S.\ 113.

\gesperrt{Gauss's method for the solution of trinomial}

\gesperrt{equations.}
\end{center}

\gesperrt{Gauss} has given a method for solving numerically, in a simple way, the roots of an equation that contains only three terms. Such equations occur frequently and include as special cases all quadratic equations and the reduced cubic equations. Of such an equation, whose general form is

\weq{1}{x^{m+n} + a x^m + b = 0}

only the positive root, if it exists, is sought at first. Any negative root is obtained by replacing $x$ by $-x$.

According to the signs of $a$, $b$ one has three cases to distinguish, since, when both signs are positive, no positive root is present. Denoting by $e$ and $g$ positive numbers, we therefore consider the three cases:

\[
\begin{aligned}
&\text{a)}\quad x^{m+n} + e x^m - g = 0,\\
&\text{b)}\quad x^{m+n} - e x^m - g = 0,\\
&\text{c)}\quad x^{m+n} - e x^m + g = 0.
\end{aligned}
\]

In the first two cases we have, by the \gesperrt{Cartesian} theorem, one positive root each; in the third, two positive roots or none may be present.

Following \gesperrt{Gauss}, we now set

\[
\lambda = \frac{g^n}{e^{m+n}},
\]

and seek to reduce the three equations a), b), c) by suitable substitutions to the form:

\[
\sin^2\Theta + \cos^2\Theta = 1
\]

We set:

\weq{2}{
\begin{aligned}
&\text{a)}\quad \frac{x^{m+n}}{g} = \sin\Theta^2, \quad \frac{e x^m}{g} = \cos\Theta^2, \quad \lambda = \frac{\sin\Theta^{2m}}{\cos\Theta^{2m+2n}},\\
&\text{b)}\quad g x^{-m-n} = \sin\Theta^2, \quad e x^{-n} = \cos\Theta^2, \quad \lambda = \frac{\sin\Theta^{2n}}{\cos\Theta^{2m+2n}},\\
&\text{c)}\quad \frac{x^n}{e} = \sin\Theta^2, \quad \frac{g x^{-m}}{e} = \cos\Theta^2, \quad \lambda = \sin\Theta^{2m}\cos\Theta^{2n}.
\end{aligned}
}


% ===== scan leaf 378, printed 353 =====
% !PRINTED 353
\weberpage{353}

The last equation yields the distinction between the two cases in which the equation c) has two positive roots or none.

For by the rules of the differential calculus one obtains for the maximum of $\sin\Theta^{2m}\cos\Theta^{2n}$
\[
\frac{m^m n^n}{(m+n)^{m+n}},
\]
which is attained for $\mathrm{tg}\,\Theta^2 = m:n$. If therefore $\lambda$ lies below this bound, we have two real roots in c), otherwise none. In the formulae (2), $\lambda$ is a given quantity, and one now has to seek from the tables the angle $\Theta$ that satisfies these equations. If one as yet has no knowledge at all of the approximate position of this angle, it is expedient, for the first approximation, to use a table proceeding say from degree to degree and abridged to two decimal places, in order then to improve the value so found with the help of more exact tables.

\gesperrt{Gauss} uses not the trigonometric tables, but the addition and subtraction logarithms first introduced by him. We shall show this briefly in one of the cases. The details for the practical application of the method are to be sought in \gesperrt{Gauss}'s memoir $^1$).

The tables of addition and subtraction logarithms, as they were first introduced and computed by \gesperrt{Gauss}, and as they are now also to be found in the tables in common use, give, for three numbers which are greater than 1:
\[
a, \quad b = 1+\frac{1}{a}, \quad c = 1+a,
\]
the \gesperrt{Briggian} logarithms:
\[
A, \quad B, \quad C.
\]
If $\Theta$ is an angle in the \gesperrt{first} octant, one may set:
\[
(3)\quad a = \mathrm{cotg}\,\Theta^2, \quad b = \frac{1}{\cos\Theta^2}, \quad c = \frac{1}{\sin\Theta^2}, \quad 0 < \Theta < 45^0,
\]
and if $\Theta$ lies in the \gesperrt{second} octant:
\[
(4)\quad a = \mathrm{tg}\,\Theta^2, \quad b = \frac{1}{\sin\Theta^2}, \quad c = \frac{1}{\cos\Theta^2}, \quad 45^0 < \Theta < 90^0.
\]

\noindent\rule{4cm}{0.4pt}

$^1$) Beiträge zur Theorie der algebraischen Gleichungen, second part (1849). \gesperrt{Gauss} Werke, Vol. III, p. 85.


% ===== scan leaf 379, printed 354 =====
% !PRINTED 354
\weberpage{354}

We wish to apply this to the case b); here we must distinguish whether $\lambda$ is smaller or greater than $2^m$, because on this depends whether $\Theta$ lies in the first or in the second octant. So there are again two cases to distinguish:
\[
\alpha)\quad \lambda < 2^m,
\]
\[
\lambda = \frac{b^m}{a^n} = \frac{c^m}{a^{m+n}} = \frac{b^{m+n}}{c^n},
\]
\[
(5)\quad x^{m+n} = gc, \quad x^n = eb, \quad x^m = \frac{ga}{e}.
\]
\[
\beta)\quad \lambda > 2^m,
\]
\[
\lambda = a^{m+n}b^m = a^n c^m = \frac{c^{m+n}}{b^n},
\]
\[
x^{m+n} = gb, \quad x^n = ec, \quad x^m = \frac{g}{ea}.
\]

In the case $\alpha)$ one would therefore set
\[
(6)\quad \log\lambda = mB - nA
\]
and then look up in the table the corresponding values of $A$ and $B$. Once these have been found, we obtain
\[
(7)\quad m\log x = A + \log g - \log e.
\]
To illustrate the use of these formulas by an example, we shall consider the equation
\[
x^3 - 2x - 2 = 0,
\]
and thus set $e = g = 2$, $m = 1$, $n = 2$, $\lambda = \frac{1}{2}$. The formula (6) therefore gives
\[
(8)\quad 2A - B = 0{,}3010300
\]
and (7)
\[
(9)\quad \log x = A.
\]

To find the first approximate value of $A$, one makes use of a small extract from the table:
\[
\left|
\begin{array}{ll}
A = 0, & B = 0{,}301 \\
A = 0{,}1, & B = 0{,}254 \\
A = 0{,}2, & B = 0{,}212 \\
A = 0{,}3, & B = 0{,}176
\end{array}
\right|.
\]
From this we obtain:
\[
\begin{array}{c|c|c}
A & 2A - B & \text{Error} \\
\hline
0{,}2 & 0{,}188 & -0{,}113 \\
0{,}3 & 0{,}424 & +0{,}123
\end{array}
\]


% ===== scan leaf 380, printed 355 =====
% !PRINTED 355
\weberpage{355}
\pdfbookmark[2]{§ 114. Calculation of the imaginary roots of a trinomial equation}{bm127}

One can now apply the method of interpolation to this in order to obtain a more accurate value of $A$, by distributing the total error of $0{,}236$ over the two values of $A$ in proportion to the partial errors, that is, by putting
\[
A = 0{,}2 + \frac{0{,}1}{0{,}236}\,0{,}113 = 0{,}247\ldots
\]

With the help of a seven-place table of \gesperrt{Zach} one now obtains:
\[
\begin{array}{c|c|c|c}
A & B & 2A - B & \text{Error} \\
\hline
0{,}247 & 0{,}1948581 & 0{,}2991419 & -0{,}0018881 \\
0{,}248 & 0{,}1944969 & 0{,}3015031 & +0{,}0004731
\end{array}
\]
and by a further application of the interpolation
\[
A = 0{,}2477996
\]
or % [sic: the page prints "odre" for "oder"]
\[
x = 1{,}769292.
\]

\begin{center}
\S.\,114.

Computation of the imaginary roots of a trinomial equation.
\end{center}

\gesperrt{Gauss} has also given a procedure for the computation of the imaginary roots of a trinomial equation, which we shall briefly discuss here. We make the assumption of real coefficients, although the method is applicable to the general case as well. We shall put the equation in the form:
\[
(1)\quad x^{m+n} - ex^m - g = 0,
\]
but here we need not restrict ourselves to positive values of $e$ and $g$.

We set
\[
x = re^{i\vartheta}
\]
and, by setting the imaginary part in (1) equal to zero by itself, we obtain
\[
(2)\quad r^n = \frac{e\sin m\vartheta}{\sin(n+m)\vartheta};
\]
%CONT>>


% ===== scan leaf 381, printed 356 =====
% !PRINTED 356
\weberpage{356}
%<<CONT
of such equations, however, we obtain two more by multiplying (1) by $x^{-m}$ and by $x^{-m-n}$. This gives
\[
(3)\quad r^{n+m} = \frac{-g\sin m\vartheta}{\sin n\vartheta},
\]
\[
(4)\quad r^m = \frac{-g\sin(n+m)\vartheta}{e\sin n\vartheta},
\]
and of these three equations each follows from the other two. If one eliminates $r$ from two of them, there results, if, as above,
\[
\lambda = \frac{g^n}{e^{n+m}}
\]
is set,
\[
(5)\quad \lambda = (-1)^n \frac{\sin n\vartheta^n \sin m\vartheta^m}{\sin(n+m)\vartheta^{n+m}}.
\]

Since one need compute only one of the two conjugate roots, one may restrict $\vartheta$ to the first quadrant, but must then in some circumstances also take $r$ negative.

How one computes from formula (5), by means of the trigonometric tables, the angle $\vartheta$, and then from one of the three formulas (2), (3), (4) the corresponding value of $r$, we shall now demonstrate on the example considered above
\[
x^3 - 2x - 2 = 0
\]
as well.

Here equation (5) becomes
\[
(6)\quad \frac{1}{2} = \frac{\sin 2\vartheta^2 \sin\vartheta}{\sin 3\vartheta^3},
\]
or
\[
(7)\quad 3\log\sin 3\vartheta - 2\log\sin 2\vartheta - \log\sin\vartheta = 0{,}3010300,
\]
from which it is first of all evident that the angle $\vartheta$ lies in the interval from 0 to $60^0$. After a few trials one finds first that $\vartheta$ lies between $33^0$ and $34^0$, and if one then computes for these two values the difference
\[
(8)\quad 3\log\sin 3\vartheta - 2\log\sin 2\vartheta - \log\sin\vartheta - 0{,}3010300
\]
one obtains first, to three decimals,
\[
0{,}026, \quad -0{,}013,
\]
and from this, by interpolation, there results the more exact value
\[
\vartheta = 33^0\,40'.
\]


% ===== scan leaf 382, printed 357 =====
% !PRINTED 357
\weberpage{357}

% [sic] source prints $30^0\,40'$; the surrounding calculation implies $33^0\,40'$
One next computes the difference (8) with somewhat greater accuracy, say to five places, for a few angles in the neighbourhood of the values $30^0\,40'$, proceeding from minute to minute, and finds from the change of sign that $\vartheta$ lies between
\[
33^0\,41' \text{ and } 33^0\,42'
\]
If one now carries out the computation to seven places for these two values, there results again, by interpolation, the more exact value
\[
\vartheta = 33^0 \quad 41' \quad 20{,}6''.
\]
From the formulae (3) or (4) one sees that $r$ is negative here, and one obtains $r$ very simply from one of these equations. One finds the \gesperrt{Briggian} logarithms:
\[
\begin{aligned}
\log(-r) &= 0{,}0266148\\
\log\cos\vartheta &= 9{,}9201547 - 10\\
\log\sin\vartheta &= 9{,}7440458 - 10.
\end{aligned}
\]
The two imaginary roots are therefore:
\[
-0{,}884646 \pm 0{,}589740\,i.
\]


% ===== scan leaf 383, printed 358 =====
\weberpage{358}% page number not printed on this page
\pdfbookmark[1]{Eleventh Section. Continued Fractions.}{bm128}
\pdfbookmark[2]{§ 115. Conversion of rational fractions into continued fractions}{bm129}
% !PRINTED none

\begin{center}
\gesperrt{Eleventh Section.}
\end{center}

\begin{center}
\gesperrt{Continued Fractions.}
\end{center}

\begin{center}
\rule{3cm}{0.4pt}
\end{center}

\begin{center}
§. 115.
\end{center}

\begin{center}
Transformation of rational fractions into continued fractions.
\end{center}

If one divides a natural number $m$ by another $n$ according to the ordinary rules, one obtains a quotient and a remainder, which is positive and smaller than the divisor $n$. If $a$ denotes the quotient and $r$ the remainder, then
\weq{1}{m = an + r.}
All numbers $m, m', m''\ldots$ which give the same remainder are called \gesperrt{of equal residue} or \gesperrt{congruent} with respect to the \gesperrt{modulus} $n$, and following \gesperrt{Gauss} one expresses this by the sign:
\weq{2}{m \equiv m' \ (\mathrm{mod}\ n).}
Such a formula is called a \gesperrt{congruence}.

If, then, we consider the system of numbers
\[
r = 0, 1, 2 \ldots n - 1
\]
every arbitrary (positive or negative) integer $m$ is congruent to \gesperrt{one} and only \gesperrt{one} of these numbers with respect to the modulus $n$. But the same property is also possessed by every system of numbers
\[
s = m_0, m_1, m_2 \ldots m_{n-1},
\]
chosen in such a way that among their remainders all the numbers $r$ occur (and each only once). Then every arbitrary number $m$ is, with respect to the modulus $n$, congruent to one and only one of the numbers $s$
%CONT>>


% ===== scan leaf 384, printed 359 =====
% !PRINTED 359
\weberpage{359}

%<<CONT
congruent. Such a system of numbers we shall call a \gesperrt{complete system of residues} for the modulus $n$. Such a complete system of residues is formed, for example, when $n$ is odd, by the numbers
\[
\varrho = \frac{-n+1}{2}, \ \frac{-n+3}{2} \cdots -1, 0, 1, 2 \cdots \frac{n-1}{2}.
\]

The numbers $r$ are called the \gesperrt{least (or least positive) residues}, the numbers $\varrho$ the \gesperrt{absolutely least residues}.

Now, in order to continue the division indicated in (1), we shall denote $m, n, r$ by $m, m_1, m_2$, and then in turn the quotient and the remainder of the division of $m_1$ by $m_2$ by $a_1, m_3$, and so forth. If, moreover, we also write $a_0$ for $a$, there results a series of equations:
\weq{3}{
\begin{aligned}
m &= a_0\, m_1 + m_2\\
m_1 &= a_1\, m_2 + m_3\\
m_2 &= a_2\, m_3 + m_4\\
&\ \cdots\cdots\cdots\cdots\\
m_{\nu-1} &= a_{\nu-1}\, m_\nu + m_{\nu+1},
\end{aligned}
}
which can be continued so long as none of the remainders vanishes, that is, so long as none of the divisions is exact. But because
\[
m_1 > m_2 > m_3 > \cdots
\]
it follows, since there is only a finite number of positive numbers smaller than $m_1$, that after a finite number of divisions the division must be exact. If $m_\nu$ is divisible by $m_{\nu+1}$, then $m_{\nu+1}$ is \gesperrt{the greatest common divisor of $m$ and $m_1$}.

For $m_{\nu+1}$ is then, by the equations (3), a divisor of all the preceding $m_\mu$, and conversely every common divisor of two adjacent $m_\mu$ is also a divisor of all the following ones, and hence also of $m_{\nu+1}$.

If therefore $m$ and $m_1$ have no common divisor other than 1 (that is, if they are \gesperrt{coprime} or \gesperrt{relatively prime}), then we must have $m_{\nu+1} = 1$, and we can represent the equations (3) as follows:
\weq{4}{
\frac{m}{m_1} = a_0 + \frac{m_2}{m_1}, \quad
\frac{m_1}{m_2} = a_1 + \frac{m_3}{m_2}, \cdots
\frac{m_{\nu-1}}{m_\nu} = a_{\nu-1} + \frac{1}{m_\nu}.
}


% ===== scan leaf 385, printed 360 =====
% !PRINTED 360
\weberpage{360}

This gives the development of the fraction $m:m_1$ or $m:n$ into a \gesperrt{continued fraction}. We obtain successively
\weq{5}{
\frac{m}{n} = a_0 + \cfrac{1}{\dfrac{m_1}{m_2}}, \quad
\frac{m}{n} = a_0 + \cfrac{1}{a_1 + \cfrac{1}{\dfrac{m_2}{m_3}}} \cdots,
}
and finally
\weq{6}{
\frac{m}{n} = a_0 + \cfrac{1}{a_1 + \cfrac{1}{a_2 + \cdots + \cfrac{1}{a_{\nu-1} + \cfrac{1}{m_\nu}}}},
}
for which, for the sake of more convenient printing, we also
\weq{7}{
\begin{aligned}
\frac{m}{n} &= \left(a_0, \frac{m_1}{m_2}\right), \quad \frac{m}{n} = \left(a_0, a_1, \frac{m_2}{m_3}\right)\\
\frac{m}{n} &= (a_0, a_1, a_2 \ldots a_{\nu-1}, m_\nu)
\end{aligned}
}
write.

To this the following remark must still be added: every positive rational fraction can be put into the form $m:n$ in such a way that $m$ and $n$ are relatively prime, and hence we can develop every such fraction into a continued fraction of the form (6) or (7). This also holds for negative rational fractions, if we admit negative values for the number $a_0$ as well. The $a_1, a_2 \ldots a_{\nu-1}, m_\nu$ are, however, positive integers even then. These numbers are completely and unambiguously determined. Only in one respect is an arbitrariness still open to us.

Namely, according to the determination made up to this point, the last of the partial denominators $m_\nu$ is greater than 1, since we have assumed that $m_{\nu+1}$ was to be the first among the residues that is equal to 1. Accordingly, if $a_\nu$ is a positive integer, we can set either
\[
m_\nu = a_\nu,
\]
or
\[
m_\nu = a_\nu + \frac{1}{1}
\]
and therefore obtain the two representations by continued fractions:
%CONT>>


% ===== scan leaf 386, printed 361 =====
% !PRINTED 361
\weberpage{361}
\pdfbookmark[2]{§ 116. Continued-fraction expansion of irrational numbers}{bm130}

\weq{8}{
\frac{m}{n} = (a_0, a_1, a_2 \ldots a_\nu)
= (a_0, a_1, a_2 \ldots a_\nu, 1).
}

I. \gesperrt{Thus every rational fraction can be developed into a continued fraction which has, at pleasure, an even or an odd number of partial denominators.}

The numbers occurring in the formulae (5),
\[
\frac{m_1}{m_2}, \ \frac{m_2}{m_3}, \ \frac{m_3}{m_4} \ldots,
\]
we call the \gesperrt{closing numbers} of the continued fraction; the numbers $a_0, a_1, a_2 \ldots$ shall be called the \gesperrt{partial denominators} (although $a_0$ does not properly occur as a denominator).

If we replace $m$ by a number $m'$ congruent to it modulo $n$, the two fractions $m:n$ and $m':n$ differ only by an integer; the continued fraction developments therefore differ only in their first terms $a_0$.

\begin{center}
\S.\,116.

\gesperrt{Continued fraction development of irrational numbers.}
\end{center}

If $x$ is a real irrational number, there will always be one and only one integer $a_0$ such that $x$ lies between $a_0$ and $a_0 + 1$. $a_0$ is positive when $x$ is greater than 1, zero when $x$ is a positive proper fraction, and negative when $x$ is negative. If we therefore set
\[
x = a_0 + \frac{1}{x_1},
\]
then $x_1$ is a positive improper fraction; hence, if $a_1$ is the integer lying next \gesperrt{below} $x_1$, $a_1$ is positive. We set
\[
\begin{aligned}
x_1 &= a_1 + \frac{1}{x_2}, \\
x_2 &= a_2 + \frac{1}{x_3}, \\
&\ \cdots\cdots\cdots \\
x_{n-1} &= a_{n-1} + \frac{1}{x_n},
\end{aligned}
\]
%CONT>>


% ===== scan leaf 387, printed 362 =====
% !PRINTED 362
\weberpage{362}
\pdfbookmark[2]{§ 117. The convergents}{bm131}
%<<CONT
so that $x_1, x_2 \ldots x_n$ are positive improper fractions, and $a_1, a_2 \ldots a_{n-1}$ are positive integers.

By the notation of the preceding section, $x$ is thus equal to the continued fraction
\[
(a_0, a_1, a_2 \ldots a_{n-1}, x_n).
\]

But if $x$ is irrational, this continued-fraction development can be \gesperrt{continued without limit}; that is, we may take $n$ as large as we please. Each such development arises from the preceding one by setting
\[
x_n = a_n + \frac{1}{x_{n+1}}
\]
Here too we call the integers $a_0, a_1, a_2 \ldots a_{n-1}$ the \gesperrt{partial denominators}, and $x_n$ the \gesperrt{closing number}.

The continued-fraction development of the rational numbers is contained in this as a special case. The continued fraction can (by Theorem I.) be continued by at most one term when the closing number is an integer.

\begin{center}
\S.\ 117.

The Convergents.
\end{center}

If we convert a continued fraction
\weq{1}{x = (a_0, a_1, a_2 \ldots a_{n-1}, x_n),}
in which we may regard $x_n$ as a variable quantity, into an ordinary fraction, we obtain in the numerator and in the denominator a linear expression in $x_n$. Thus we have
\weq{2}{x = \frac{P_n x_n + P_{n-1}}{Q_n x_n + Q_{n-1}},}
in which $P_n$, $Q_n$, $P_{n-1}$, $Q_{n-1}$ are integral rational functions of $a_0, a_1, a_2 \ldots a_{n-1}$. That this is correct follows for the first values $n = 1, 2 \ldots$ by direct calculation:
\[
\begin{aligned}
x &= a_0 + \frac{1}{x_1} = \frac{a_0 x_1 + 1}{x_1}\\
&= a_0 + \cfrac{1}{a_1 + \cfrac{1}{x_2}} = \frac{(a_0 a_1 + 1)x_2 + a_0}{a_1 x_2 + 1}.
\end{aligned}
\]


% ===== scan leaf 388, printed 363 =====
% !PRINTED 363
\weberpage{363}

If we assume the formula (2) as already proved for $n$, then, setting
\[
x_n = a_n + \frac{1}{x_{n+1}}
\]
we obtain:
\[
x = \frac{(P_n a_n + P_{n-1}) x_{n+1} + P_n}{(Q_n a_n + Q_{n-1}) x_{n+1} + Q_n}.
\]
\gesperrt{But this is the formula that results from (2) by permutation of $n$ with $n+1$, if we adopt the following law of formation for the $P_n$, $Q_n$:}
\weq{3}{
\begin{aligned}
P_{n+1} &= a_n P_n + P_{n-1} \\
Q_{n+1} &= a_n Q_n + Q_{n-1}.
\end{aligned}
}

These formulas completely determine the $P_n$, $Q_n$ for $n = 2, 3, 4 \ldots$, if we further add the determination:
\weq{4}{
\begin{aligned}
P_0 &= 1, & P_1 &= a_0, \\
Q_0 &= 0, & Q_1 &= 1.
\end{aligned}
}
The fractions
\[
\frac{P_0}{Q_0}, \ \frac{P_1}{Q_1}, \ \frac{P_2}{Q_2} \ldots,
\]
of which the first, with the denominator zero, is introduced only formally, for the sake of uniformity, are called the \gesperrt{convergents} of the continued fraction (1), $P_n$ the numerator and $Q_n$ the denominator of the $n$th convergent.

The law of formation (3) shows that, if the $a$ are integers, then $P_n$ and $Q_n$ are also integers. If the partial denominators from $a_1$ on are positive, then the denominators $Q_1, Q_2, Q_3 \ldots$ are all \gesperrt{positive} and increase with $n$, and indeed, since they are integers, to infinity. Only in the special case where $a_1 = 1$ is $Q_1 = Q_2 = 1$, and the increase begins only from $Q_3$ on, and we can therefore formulate the theorem:

II. \gesperrt{If} $Q_{n-1}, Q_n$ \gesperrt{are the denominators of two successive convergents, then}
\[
0 \leqq Q_{n-1} \leqq Q_n, \quad Q_n = \infty \ \text{for} \ n = \infty,
\]
\gesperrt{where equality with the lower limit 0 can occur only for $n = 1$, and with the upper limit $Q_n$ only for $n = 2$.}


% ===== scan leaf 389, printed 364 =====
% !PRINTED 364
\weberpage{364}

The $P_1, P_2, P_3 \ldots$ are, if $a_0$ is positive, likewise all positive; if $a_0$ is negative, then they are all negative, with the possible exception of $P_2 = a_0 a_1 + 1$, which can be equal to zero. The $Q_n$ are entirely independent of $a_0$. In what manner the $P_n$ depend on $a_0$ we can recognise in the following way.

Let us take the sequence of numbers
\[
a_1, \ a_2, \ a_3, \ a_4 \ldots
\]
as given and consider the recursion formula
\weq{5}{
T_{n+1} = a_n T_n + T_{n-1},
}
then $T_n$ is thereby completely determined for all $n$, once $T_0$ and $T_1$ are given.

If we take two special cases $Q_n$ and $R_n$, which are determined by the conditions
\weq{6}{
\begin{aligned}
R_0 &= 1, & R_1 &= 0 \\
Q_0 &= 0, & Q_1 &= 1
\end{aligned}
}
then one obtains the general solution of equation (5) in the form
\weq{7}{
T_n = T_0 R_n + T_1 Q_n
}
and in this is therefore also contained
\weq{8}{
P_n = R_n + a_0 Q_n
}
If we consider two solutions of (5), $T_n$ and $S_n$, thus:
\[
\begin{aligned}
T_{n+1} &= a_n T_n + T_{n-1} \\
S_{n+1} &= a_n S_n + S_{n-1},
\end{aligned}
\]
then by elimination of $a_n$ it follows that
\[
T_{n+1} S_n - T_n S_{n+1} = -(T_n S_{n-1} - T_{n-1} S_n),
\]
or, if we set
\[
T_n S_{n+1} - S_n T_{n+1} = \varDelta_n
\]
,. % [sic] comma followed by period as printed
\weq{9}{
\varDelta_n = -\varDelta_{n-1} = \varDelta_{n-2} = \cdots = (-1)^n \varDelta_0.
}
But for $T_n = R_n$ and $S_n = Q_n$ we have $\varDelta_0 = 1$, [by (6)], hence by (9)
\[
R_n Q_{n+1} - Q_n R_{n+1} = (-1)^n,
\]
or, by changing $n$ into $n-1$ and reversing the sign:
\weq{10}{
R_n Q_{n-1} - Q_n R_{n-1} = (-1)^n.
}
If one applies this to two functions $T_n$, $S_n$, which are defined by
\[
T_n = T_0 R_n + T_1 Q_n, \quad S_n = S_0 R_n + S_1 Q_n
\]
%CONT>>


% ===== scan leaf 390, printed 365 =====
% !PRINTED 365
\weberpage{365}
\pdfbookmark[2]{§ 118. Solution of indeterminate equations in two unknowns}{bm132}
%<<CONT
then it follows that
\[
T_n S_{n-1} - S_n T_{n-1} = (-1)^n (T_0 S_1 - T_1 S_0),
\]
and so in particular

\weq{11}{P_n Q_{n-1} - Q_n P_{n-1} = (-1)^n.}

This formula, of which we shall yet make manifold application, shows that the numbers $P_n$, $Q_n$ have no common divisor, and hence that the convergents $P_n : Q_n$ cannot be reduced by cancellation.

If the continued fraction is a finite one, that is, if the closing number $x_{n-1} = a_{n-1}$ is an integer, then we carry the formation of the convergents no further than to
\[
\frac{P_n}{Q_n} = \frac{P_{n-1} a_{n-1} + P_{n-2}}{Q_{n-1} a_{n-1} + Q_{n-2}},
\]
and formula (2) shows that then the \gesperrt{last convergent agrees with the value of the continued fraction.}

\begin{center}
§. 118.

Solution of indeterminate equations with two unknowns.
\end{center}

The formulas last found lead to the solution of a problem that occurs in many applications:

\gesperrt{Let $\alpha$, $\beta$ be two given integers with no common divisor; two other integers are to be found which satisfy the condition:}

\weq{1}{\alpha y - \beta x = 1.}

This problem, if it is solvable at all, has infinitely many solutions, all of which can be derived from any one of them.

For if $x_0$, $y_0$ is a solution, so that

\weq{2}{\alpha y_0 - \beta x_0 = 1,}
then by subtraction of (1) and (2) it follows that
\[
\alpha (y - y_0) = \beta (x - x_0),
\]
from which it is to be concluded, since $\alpha$ has with $\beta$ no common divisor
%CONT>>


% ===== scan leaf 391, printed 366 =====
% !PRINTED 366
\weberpage{366}
%<<CONT
that $x - x_0$ is divisible by $\alpha$; if accordingly we set, denoting by $h$ an arbitrary integer,
\[
x = x_0 + h \alpha,
\]
it follows that:
\[
y = y_0 + h \beta,
\]
and all solutions of (1) are contained in this form.

We therefore need only to seek \gesperrt{one} solution of (1), which we can always obtain in the following way. We assume $\beta$ to be positive, which does not essentially restrict the generality, since if need be we can change $x$ into $-x$, and, following \S.\,115, we develop the rational number $\alpha : \beta$ into a continued fraction:
\[
\frac{\alpha}{\beta} = (a_0, a_1, a_2 \ldots a_{n-1}).
\]

From this we form the convergents
\[
\frac{P_0}{Q_0}, \quad \frac{P_1}{Q_1}, \quad \frac{P_2}{Q_2} \cdots \frac{P_{n-1}}{Q_{n-1}}, \quad \frac{P_n}{Q_n},
\]
so that
\[
\frac{\alpha}{\beta} = \frac{P_n}{Q_n}.
\]

But since $\alpha$, $\beta$ as well as $P_n$, $Q_n$ are without common divisor, and since $\beta$ and $Q_n$ are positive, it follows that
\[
\alpha = P_n, \quad \beta = Q_n.
\]
Since further, by \S.\,117,
\[
P_n Q_{n-1} - Q_n P_{n-1} = (-1)^n
\]
it follows that

\weq{3}{x = (-1)^n P_{n-1}, \quad y = (-1)^n Q_{n-1}}
are integers which satisfy the equation (1), and with this the problem is completely solved.

For moderately large numbers the calculation is quite simple.

If, for example, we set $\alpha = 24335$, $\beta = 3588$, we first obtain by \S.\,115 the continued fraction:
\[
\frac{24335}{3588} = (6,\ 1,\ 3,\ 1,\ 1,\ 2,\ 6,\ 2,\ 1,\ 1,\ 4),
\]
and the convergents
\[
\frac{1}{0}, \quad \frac{6}{1}, \quad \frac{7}{1}, \quad \frac{27}{4}, \quad \frac{34}{5}, \quad \frac{61}{9}, \quad \frac{156}{23}, \quad \frac{997}{147}, \quad \frac{2150}{317}, \quad \frac{3147}{464},
\]
\[
\frac{5297}{781}, \quad \frac{24335}{3588},
\]
%CONT>>


% ===== scan leaf 392, printed 367 =====
% !PRINTED 367
\weberpage{367}
%<<CONT
and since here $n = 11$, and thus odd, one must set:
\[
x = -5297, \quad y = -781.
\]
These are the absolutely smallest values of $x$, $y$ that satisfy equation (1).

The smallest positive solution is obtained from this by adding $\alpha$ and $\beta$ to them, thus
\[
x = 19038, \quad y = 2807.
\]
To this is also reduced the solution of the more general equation

\weq{4}{\alpha y - \beta x = \gamma}
in integers $x$, $y$.

First of all it is clear that, if (4) is to be solvable at all, every common divisor of $\alpha$ and $\beta$ must also be a divisor of $\gamma$. Such a common divisor is then removed by division. We therefore assume here too that $\alpha$ and $\beta$ have no common divisor. Then it follows, just as above, that all solutions of (4) are obtained from one of them, $x_0$, $y_0$, in the form
\[
x = x_0 + h \alpha, \quad y = y_0 + h \beta,
\]
in which $h$ is an undetermined integer. If one then sets in (4)
\[
x = \gamma \xi, \quad y = \gamma \eta,
\]
and divides by $\gamma$, then (4) passes into
\[
\alpha \eta - \beta \xi = 1,
\]
that is, into an equation of the form (1).

The solution becomes a completely determinate one when a further condition is given from which $h$ can be determined, e.g. the condition that $y$ shall lie between $0$ and $\beta$, with the inclusion of one of the two limits. We shall summarise this in the following theorem:

\gesperrt{III. The indeterminate equation}
\[
\alpha y - \beta x = \gamma
\]
\gesperrt{has, if $\alpha$, $\beta$, $\gamma$ are integers and $\alpha$ and $\beta$ are without common divisor, always an integral solution $x$, $y$, and only one if the condition}
\[
0 \leqq y < \beta,
\]
\gesperrt{or}
\[
0 < y \leqq \beta
\]
\gesperrt{is added.}


% ===== scan leaf 393, printed 368 =====
% !PRINTED 368
\weberpage{368}
By what was remarked above, however, there are also always solutions of the equation (4) when $\gamma$ is divisible by the greatest common divisor of $\alpha$ and $\beta$. In particular we can state the theorem: \gesperrt{The greatest common divisor $\delta$ of two numbers $\alpha$, $\beta$ can be represented in the form}
\[
\alpha x - \beta y = \delta
\]
\gesperrt{in which $x$, $y$ are integers.}

Theorem III can be generalised in the following way:

\gesperrt{If $\alpha_1$, $\alpha_2$, $\alpha_3 \ldots$ are given integers in arbitrary number, without a divisor common to them all, and $m$ is an arbitrarily given number, then the integers $x_1$, $x_2$, $x_3 \ldots$ can always be determined so that}

\weq{5}{m = \alpha_1 x_1 + \alpha_2 x_2 + \alpha_3 x_3 + \cdots}
\gesperrt{holds.}

If the right-hand side of (5) consists of only two terms, this theorem coincides with Theorem III. If, then, we assume the possibility of satisfying a formula (5) as proved whenever the right-hand side consists of fewer terms, we can, if $\delta$ is the greatest common divisor of $\alpha_2$, $\alpha_3 \ldots$, since $\delta$ must be relatively prime to $\alpha_1$, satisfy the two equations:
\[
\delta = \alpha_2 y_2 + \alpha_3 y_3 + \cdots
\]
\[
m = \alpha_1 x_1 + \delta x
\]
by integral $x$, $x_1$, $y_2$, $y_3 \ldots$ But then the equation (5) is satisfied by
\[
x_1, \quad x_2 = x y_2, \quad x_3 = x y_3 \ldots
\]

The problem solved by Theorem III is also expressed in number theory as follows: If $\alpha$, $\beta$, $\gamma$ are given numbers, a number $y$ is to be found which satisfies the congruence
\[
\alpha y \equiv \gamma \ (\mathrm{mod}\ \beta)
\]
This problem therefore always has a solution when $\gamma$ is divisible by the greatest common divisor of $\alpha$ and $\beta$.


% ===== scan leaf 394, printed 369 =====
% !PRINTED 369
\weberpage{369}
\pdfbookmark[2]{§ 119. Convergence of the approximating fractions}{bm133}

\begin{center}
§. 119.

\gesperrt{Convergence of the Convergents.}
\end{center}

We now assume that the sequence of numbers $a_1, a_2, a_3 \ldots$
is an unlimited one, and form the difference of two successive
convergents:
\weq{1}{\frac{P_n}{Q_n} - \frac{P_{n-1}}{Q_{n-1}} = \frac{(-1)^n}{Q_n Q_{n-1}}.}

This difference therefore has the positive sign for even $n$ and
the negative sign for odd $n$, and decreases without limit in
absolute value as $n$ increases to infinity. If we further form
from (1):
\weq{2}{\begin{aligned}
\frac{P_n}{Q_n} - \frac{P_{n-2}}{Q_{n-2}} &= \frac{(-1)^n}{Q_n Q_{n-1}} + \frac{(-1)^{n+1}}{Q_{n-1} Q_{n-2}} \\
&= \frac{(-1)^n}{Q_{n-1}}\left(\frac{1}{Q_n} - \frac{1}{Q_{n-2}}\right),
\end{aligned}}

then, since $Q_{n-2} < Q_n$, it follows that this difference is
negative for even $n$ and positive for odd $n$.

From this it follows that the sequence of convergents with even
index
\weq{3}{\frac{P_2}{Q_2}, \quad \frac{P_4}{Q_4}, \quad \frac{P_6}{Q_6} \ldots}

is \gesperrt{a decreasing one}, and the sequence of convergents with odd
index
\weq{4}{\frac{P_1}{Q_1}, \quad \frac{P_3}{Q_3}, \quad \frac{P_5}{Q_5} \ldots}

\gesperrt{an increasing one}.

Now by \S.\,117
\[
x = \frac{P_n x_n + P_{n-1}}{Q_n x_n + Q_{n-1}},
\]
hence
\weq{5}{\frac{P_n}{Q_n} - x = \frac{(-1)^n}{Q_n(Q_n x_n + Q_{n-1})},}

and since $Q_n, Q_{n-1}, x_n$ are positive, all the numbers of the
sequence (3) are greater than $x$, all the numbers of the sequence
(4) less than $x$.


% ===== scan leaf 395, printed 370 =====
% !PRINTED 370
\weberpage{370}

The difference (5) sinks below every limit as $n$ increases
without bound, and since $x_n > a_n$, and hence $Q_n x_n + Q_{n-1} > Q_{n+1}$,
it is in absolute value smaller than
\weq{6}{\frac{1}{Q_n Q_{n+1}}}

and all the more smaller than $1 : Q_n^2$.

The numbers of the sequence (3) thus approach the limit $x$
decreasingly, those of the sequence (4) increasingly. The expression (6)
gives a measure of the error one commits when $x$ is replaced by
the convergent $P_n : Q_n$.

The convergents are thus approximate expressions
of irrational numbers by rational
fractions.

That these convergents, for an equal degree of approximation
to the irrational number $x$, are the simplest possible
is expressed by the following theorem:

\gesperrt{IV. Between the two rational
fractions}
\[
\frac{P_n}{Q_n}, \quad \frac{P_{n-1}}{Q_{n-1}}
\]
\gesperrt{no other rational fraction can be inserted
whose denominator is smaller than $Q_n$ or even only
equal to $Q_n$.}

Suppose that the rational fraction $M : N$ lies between
the convergents $P_n : Q_n$ and $P_{n-1} : Q_{n-1}$. Then the
difference
\[
\frac{P_n}{Q_n} - \frac{P_{n-1}}{Q_{n-1}}
\]
is absolutely greater than, and of the same sign as,
\[
\frac{M}{N} - \frac{P_{n-1}}{Q_{n-1}},
\]
hence
\[
(-1)^n \left(\frac{P_n}{Q_n} - \frac{P_{n-1}}{Q_{n-1}}\right) > (-1)^n \left(\frac{M}{N} - \frac{P_{n-1}}{Q_{n-1}}\right).
\]

If one multiplies both sides by $N Q_{n-1}$, then by (1) it follows that
\[
\frac{N}{Q_n} > (-1)^n (M Q_{n-1} - N P_{n-1}),
\]
and since on the right stands a positive integer, which is therefore at least
equal to 1, it follows that
\[
N > Q_n.
\]


% ===== scan leaf 396, printed 371 =====
% !PRINTED 371
\weberpage{371}
\pdfbookmark[2]{§ 120. Equivalent numbers}{bm134}

\begin{center}
§. 120.

\gesperrt{Equivalent Numbers.}
\end{center}

The continued fractions lead us to the consideration of a
special kind of linear substitutions, which are of great importance
in algebra and number theory generally, and which we shall
consider somewhat more closely.

If $x$ and $y$ are two numbers, or also variable quantities,
which stand in the mutual dependence
\weq{1}{y = \frac{\alpha x + \beta}{\gamma x + \delta},}

in which $\alpha, \beta, \gamma, \delta$ are integers satisfying the condition
\weq{2}{\alpha \delta - \beta \gamma = \varepsilon = \pm 1}

then we call $x$ and $y$ \gesperrt{equivalent} to one another$^{1)}$.
We call them \gesperrt{properly} or \gesperrt{improperly equivalent},
according as $\varepsilon = +1$ or $\varepsilon = -1$. This relation is
a reciprocal one, for from (1) it follows that
\weq{3}{x = \frac{\delta y - \beta}{-\gamma y + \alpha}.}

\gesperrt{If two quantities are equivalent to a third,
then they are also equivalent to one another.}

For if
\weq{4}{x = \frac{\alpha' z + \beta'}{\gamma' z + \delta'},}

then from (1) it follows that
\weq{5}{y = \frac{\alpha'' z + \beta''}{\gamma'' z + \delta''},}

where
\weq{6}{\begin{aligned}
\alpha'' &= \alpha \alpha' + \beta \gamma', \quad \beta'' = \alpha \beta' + \beta \delta', \\
\gamma'' &= \gamma \alpha' + \delta \gamma', \quad \delta'' = \gamma \beta' + \delta \delta',
\end{aligned}}
\weq{7}{\alpha'' \delta'' - \beta'' \gamma'' = (\alpha \delta - \beta \gamma)(\alpha' \delta' - \beta' \gamma'),}

or
\weq{8}{\varepsilon'' = \varepsilon \varepsilon'.}

\begingroup
\renewcommand{\thefootnote}{1)}
\footnotetext[1]{Cf. \gesperrt{Dedekind}, letter to Herr \gesperrt{Borchardt} on the theory
of elliptic modular functions. Crelle's Journal, Vol. 83 (1877).}
\endgroup


% ===== scan leaf 397, printed 372 =====
% !PRINTED 372
\weberpage{372}

The substitution (5) is said to be \gesperrt{composed} of (1) and (4).

The need for an abbreviated notation for these substitutions
makes itself felt. Since it is often not the variables that matter,
but only the substitution numbers $\alpha, \beta, \gamma, \delta$,
one denotes the whole substitution (1) by the substitution numbers,
or also merely by a single letter
\weq{9}{S = \begin{pmatrix} \alpha, & \beta \\ \gamma, & \delta \end{pmatrix},}

and then, when it is necessary, writes the equation (1) thus:
\[
y = S(x) = \begin{pmatrix} \alpha, & \beta \\ \gamma, & \delta \end{pmatrix}(x).
\]

The composition of two substitutions is denoted by setting
the symbols side by side, in which, however, attention must be
paid to the \gesperrt{order}; thus
\weq{10}{\begin{pmatrix} \alpha, & \beta \\ \gamma, & \delta \end{pmatrix} \begin{pmatrix} \alpha', & \beta' \\ \gamma', & \delta' \end{pmatrix} = \begin{pmatrix} \alpha'', & \beta'' \\ \gamma'', & \delta'' \end{pmatrix},}

or
\weq{11}{S S' = S''.}

The formulas (6) contain the rule according to which a
composite substitution is to be formed. One can derive it from
the multiplication rule for determinants, but must observe that
in the first factor the summation is to be taken by rows, in the
second by columns, just as the formulas (6) show.

Just as one composes two substitutions, one can also form
the composition of several:
\[
S S_1 S_2 S_3 \ldots
\]

In forming the composite substitution the components may
not, in general at least, be interchanged; but one may at will
combine two adjacent ones into one, then again two, and so
on; that is, the \gesperrt{commutative} law does not hold, but the
\gesperrt{associative} law does.

This follows immediately from the fact that, if $x$ is expressed
in terms of $x_1$, $x_1$ in terms of $x_2$, and $x_2$ in terms of $x_3$,
then the expression of $x$ in terms of $x_3$ can be found either in
this way, that one first expresses $x_1$ in $x$ by means of $x_2$ and
then $x_2$ by means of $x_3$, or in this way, that
%CONT>>


% ===== scan leaf 398, printed 373 =====
% !PRINTED 373
\weberpage{373}
%<<CONT
one first expresses $x_1$ in terms of $x_3$ and substitutes this into the expression of $x$ in terms of $x_1$.

By (8), a composite equivalence is a \gesperrt{proper} or an \gesperrt{improper} one according as an even or an odd number of improper ones is found among the components.

If one composes the two substitutions
% [mirrored 2026-08-07] the German restored the equation number (12) the page prints here
\weq{12}{\begin{pmatrix}\alpha, & \beta\\ \gamma, & \delta\end{pmatrix}, \quad
\begin{pmatrix}\varepsilon\delta, & -\varepsilon\beta\\ -\varepsilon\gamma, & \varepsilon\alpha\end{pmatrix}}
then one obtains, no matter which of the two one puts in the first place,
\weq{13}{\begin{pmatrix}1, & 0\\ 0, & 1\end{pmatrix}.}

This substitution has, by (1), the same meaning as $y = x$; it changes nothing, and is called the \gesperrt{identical substitution}, and indeed is also denoted briefly by (1). Accordingly one calls the two substitutions (12) \gesperrt{reciprocal} to one another, and denotes them by
\[
S, \; S^{-1},
\]
or one sets
\weq{14}{\begin{pmatrix}\varepsilon\delta, & -\varepsilon\beta\\ -\varepsilon\gamma, & \varepsilon\alpha\end{pmatrix} = \begin{pmatrix}\alpha, & \beta\\ \gamma, & \delta\end{pmatrix}^{-1}.}

Under composition with the identical substitution (13) every other substitution remains unchanged.

From this we very simply derive the proof of the theorem:

\gesperrt{V. All rational numbers are equivalent to one another, and indeed both properly and improperly.}

For if $m:n$ and $m':n'$ are two rational fractions, and $m$ and $n$, as well as $m'$ and $n'$, are without common divisor, then, if $\varepsilon, \varepsilon'$ are $\pm 1$ at pleasure, we can determine the numbers $\mu, \nu$ and $\mu', \nu'$ according to \S.\,118 so that
\[
m\nu - n\mu = \varepsilon, \qquad m'\nu' - n'\mu' = \varepsilon'
\]
holds, and so that therefore
\[
M = \begin{pmatrix}m, & \mu\\ n, & \nu\end{pmatrix}, \quad M' = \begin{pmatrix}m', & \mu'\\ n', & \nu'\end{pmatrix}
\]
are two linear substitutions. It is then also
% [mirrored 2026-08-07] the German restored the dropped final display of page 373
\[
M M'^{-1} = \begin{pmatrix}\alpha, & \beta\\ \gamma, & \delta\end{pmatrix}
\]
%CONT>>


% ===== scan leaf 399, printed 374 =====
% !PRINTED 374
\weberpage{374}
\pdfbookmark[2]{§ 121. Expansion of equivalent numbers in continued fractions}{bm135}
%<<CONT
a linear substitution, and
\[
\begin{pmatrix}m, & \mu\\ n, & \nu\end{pmatrix} = \begin{pmatrix}\alpha, & \beta\\ \gamma, & \delta\end{pmatrix}\begin{pmatrix}m', & \mu'\\ n', & \nu'\end{pmatrix},
\]
from which it follows that
\[
m = \alpha m' + \beta n', \qquad n = \gamma m' + \delta n',
\]
hence
\[
\frac{m}{n} = \dfrac{\alpha\dfrac{m'}{n'}+\beta}{\gamma\dfrac{m'}{n'}+\delta},
\]
which was to be proved. The equivalence is a proper or an improper one, according as $\varepsilon = \varepsilon'$ or $\varepsilon = -\varepsilon'$.

Every number is improperly equivalent to its opposite and to its reciprocal, for we have
\[
-x = \begin{pmatrix}-1, & 0\\ 0, & 1\end{pmatrix}(x), \qquad \frac{1}{x} = \begin{pmatrix}0, & 1\\ 1, & 0\end{pmatrix}(x).
\]

\begin{center}
\S.\,121.

\gesperrt{Development of equivalent numbers in continued fractions.}
\end{center}

If of two equivalent numbers one is irrational, then so is the other; and if one is real, then so is the other. We make both of these assumptions, and thus in
\weq{1}{y = \dfrac{\alpha x + \beta}{\gamma x + \delta}, \qquad \alpha\delta - \beta\gamma = \varepsilon}
we suppose $x$ to be irrational and real, and we now wish to investigate how the continued-fraction development of $y$ can be derived from the continued-fraction development of $x$.

Let $x$, then, be developed in a continued fraction with the closing number $x_n$:
\weq{2}{x = (a_0, a_1, a_2 \ldots a_{n-1}, x_n),}
in which we leave $n$ undetermined for the time being.

Expressed by means of the convergents, this becomes
\weq{3}{x = \dfrac{P_n x_n + P_{n-1}}{Q_n x_n + Q_{n-1}}, \qquad P_n Q_{n-1} - Q_n P_{n-1} = (-1)^n.}


% ===== scan leaf 400, printed 375 =====
% !PRINTED 375
\weberpage{375}

If we substitute this in (1), it follows that

\weq{4}{y = \begin{pmatrix} \alpha, & \beta \\ \gamma, & \delta \end{pmatrix}\begin{pmatrix} P_n, & P_{n-1} \\ Q_n, & Q_{n-1} \end{pmatrix}(x_n) = \begin{pmatrix} R_n, & R_{n-1} \\ S_n, & S_{n-1} \end{pmatrix}(x_n),}

in which, by \S.\,120, (10) and (6), (7),

\weq{5}{\begin{aligned}
R_n &= \alpha P_n + \beta Q_n, & R_{n-1} &= \alpha P_{n-1} + \beta Q_{n-1}, \\
S_n &= \gamma P_n + \delta Q_n, & S_{n-1} &= \gamma P_{n-1} + \delta Q_{n-1},
\end{aligned}}

\weq{6}{R_n S_{n-1} - S_n R_{n-1} = (-1)^n \varepsilon.}

If one puts $S_n$ into the form
\[
S_n = Q_n\left(\gamma \frac{P_n}{Q_n} + \delta\right),
\]
then, since $P_n : Q_n$ approaches the value $x$ to any arbitrary degree and $Q_n$ is positive, it follows that, if $n$ is chosen large enough, $S_n$ agrees in sign with
\[
(\gamma x + \delta)
\]
Since $x$ is irrational, this quantity is different from zero, and we shall assume it to be positive. Were it negative, we should only have to change the signs of the four numbers $\alpha, \beta, \gamma, \delta$ simultaneously, whereby (1) remains unaltered, in order to make it positive.

Hence $S_n$ will be positive for sufficiently large $n$. But now from (5), having regard to \S.\,117, it follows that
\[
S_{n+1} = a_n S_n + S_{n-1},
\]
from which one concludes that, if $n$ is so large that $S_{n-2}$ and the following $S$ are positive, $S_n$ increases together with $n$, and hence:

\weq{7}{S_n > S_{n-1} > 0.}

We now develop the rational fraction $R_n : S_n$ into a finite continued fraction

\weq{8}{\frac{R_n}{S_n} = (b_0, b_1, b_2 \ldots b_{m-1}),}

and denote the next-to-last convergent by $R' : S'$, so that we have the relation:

\weq{9}{R_n S' - S_n R' = (-1)^m.}

By Theorem I in \S.\,115, however, we may assume $m$ to be even or odd at pleasure, and we shall dispose of it in such a way that
\[
(-1)^m = (-1)^n \varepsilon
\]
holds. Moreover, as we know from \S.\,117, II,
%CONT>>


% ===== scan leaf 401, printed 376 =====
% !PRINTED 376
\weberpage{376}
%<<CONT
\weq{10}{S_n \geqq S' \geqq 0,}
in which, however, the equality sign can occur in the lower limit only for $m = 1$, and in the upper limit only for $m = 2$ and $b_1 = 1$, hence in general only when $S_n = 1$; and this can be avoided, by (7), if $n$ is taken large enough.

But since, by \S.\,118, III, the numbers $S', R'$ are completely determined by the conditions (9), (10), and $S_{n-1}, R_{n-1}$ satisfy the same conditions by (6), (7), it follows that
\[
S' = S_{n-1}, \qquad R' = R_{n-1}.
\]

From this it now results further that the continued fraction with the closing number $x_n$
\weq{11}{(b_0, b_1, b_2 \ldots b_{m-1}, x_n)}
has the value
\[
\dfrac{R_n x_n + R_{n-1}}{S_n x_n + S_{n-1}},
\]
and therefore agrees with $y$.

If one now develops $x_n$ further into a continued fraction
\[
x_n = (a_n, a_{n+1}, a_{n+2} \ldots),
\]
one obtains from (2) and (11):
\[
x = (a_0, a_1, a_2 \ldots a_{n-1}, a_n, a_{n+1}, a_{n+2} \ldots)
\]
\[
y = (b_0, b_1, b_2 \ldots b_{m-1}, a_n, a_{n+1}, a_{n+2} \ldots),
\]
or, expressed in words, the theorem:

\gesperrt{VI. The continued fraction developments of two equivalent numbers agree with one another from a certain partial denominator on.}

We may add that, if the equivalence is a proper one ($\varepsilon = +1$), the number of partial denominators preceding the agreeing ones is even in both, or odd in both; and that, if the equivalence is improper, this number is even in the one and odd in the other.

That the theorem also holds conversely is easy to see. For if two continued fractions agree from a certain partial denominator on, they can be written in such a way that they have the same closing number. But the value of a continued fraction is always equivalent to each of its closing numbers [\S.\,117, (2)]; \gesperrt{hence two continued fractions with the same closing number are also equivalent to one another.}


% ===== scan leaf 402, printed 377 =====
% !PRINTED 377
\weberpage{377}
\pdfbookmark[2]{§ 122. Quadratic irrational numbers}{bm136}
\begin{center}
§. 122.

Quadratic Irrational Numbers.
\end{center}

Simple and beautiful laws emerge when one applies the continued fraction expansion to the determination of the roots of a quadratic equation. The root of a quadratic equation with integral coefficients has the form
\weq{1}{\omega = x + y\sqrt{d},}
in which $x$, $y$, $d$ are integral or fractional rational numbers. We may, however, without restricting the generality, \gesperrt{assume $d$ to be an integer and without square divisor}; for if $d$ has a denominator, we can multiply through by this denominator, and may reckon the root of the square denominator, and of any square divisor of the numerator, to $y$; we further assume that $d$ is not equal to 1, since otherwise $\omega$ would be rational; $d$ may be positive or negative, and on this it depends whether $\omega$ is real or imaginary.

If we change the sign of the root, we obtain
\weq{2}{\omega' = x - y\sqrt{d},}
which is called the number \gesperrt{conjugate} to $\omega$. The product of the two numbers $\omega$, $\omega'$
\[
\omega\omega' = x^2 - y^2 d
\]
is called the \gesperrt{norm of} $\omega$ (or also of $\omega'$). \gesperrt{The norm of a product of two quadratic irrational numbers is equal to the product of the norms.}

The numbers $\omega$ and $\omega'$ are the roots of an equation with rational coefficients:
\weq{3}{\omega^2 - 2x\omega + (x^2 - y^2 d) = 0.}
To transform it into an equation with integral coefficients, we set
\weq{4}{x^2 - y^2 d = -\frac{a}{c}, \quad 2x = \frac{b}{c},}
in which $a$, $b$, $c$ are integers without common divisor, and obtain from (3)
\weq{5}{c\omega^2 = a + b\omega,}
that is, a \gesperrt{primitive} equation with integral coefficients\par % [sic] source omits terminal punctuation; the next leaf begins a new sentence


% ===== scan leaf 403, printed 378 =====
% !PRINTED 378
\weberpage{378}
The numbers $a$, $b$, $c$ are determined by (4) only up to a common sign.

The discriminant of (3) is
\weq{6}{D = b^2 + 4ac = 4c^2 y^2 d,}
and shall also be called the \gesperrt{discriminant of the irrational number} $\omega$. $D$ is thus an integer which agrees in sign with $d$ and which is divisible by $d$; for if $p$ is a prime divisor of $d$, which therefore by hypothesis divides $d$ only once, then $p$ cannot divide the denominator of $2cy$, and hence by (6) must divide $D$. Therefore $2cy$ is also an integer, and the quotient $D:d$ is a square number.

By (4) and (6), $\omega$ and $\omega'$ can be represented thus:
\weq{7}{\omega = \frac{\sqrt{D}+b}{2c}, \quad \omega' = \frac{-\sqrt{D}+b}{2c},}
or, what is the same thing:
\weq{8}{\omega = \frac{2a}{\sqrt{D}-b}, \quad \omega' = \frac{-2a}{\sqrt{D}+b},}
if the sign of the root is determined from
\weq{9}{\sqrt{D} = 2cy\sqrt{d}}
Since an even square number is divisible by 4, while an odd one, divided by 4, leaves the residue 1, it follows from (6) that
\weq{10}{D \equiv 0 \quad \text{or} \quad \equiv 1 \ (\mathrm{mod}\ 4)}
must hold.

We wish further to investigate how the numbers $a$, $b$, $c$, $D$ change when we pass from $\omega$ to an equivalent number $\omega_1$. Let, then,
\weq{11}{\omega_1 = \frac{\alpha\omega+\beta}{\gamma\omega+\delta}, \quad \omega = \frac{\delta\omega_1-\beta}{-\gamma\omega_1+\alpha_1} \quad \cdot \quad \alpha\delta - \beta\gamma = \varepsilon.} % [sic] denominator prints $\alpha_1$, not $\alpha$; source uses a centred-dot separator, not a comma, before the third identity
From (5) one obtains for $\omega_1$ the quadratic equation
\weq{12}{c_1\omega_1^2 = a_1 + b_1\omega_1,}
if we put
\weq{13}{\begin{aligned}
-a_1 &= -a\alpha^2 + b\alpha\beta + c\beta^2\\
b_1 &= -2a\alpha\gamma + b(\alpha\delta+\beta\gamma) + 2c\beta\delta,\\
c_1 &= -a\gamma^2 + b\gamma\delta + c\delta^2
\end{aligned}}
From (13) one obtains, by solving for $a$, $b$, $c$:
%CONT>>


% ===== scan leaf 404, printed 379 =====
% !PRINTED 379
\weberpage{379}
\pdfbookmark[2]{§ 123. Reduced numbers with negative discriminant}{bm137}
%<<CONT
\weq{14}{\begin{aligned}
a &= a_1\delta^2 + b_1\beta\delta - c_1\beta^2\\
b &= 2a_1\gamma\delta + b_1(\alpha\delta + \beta\gamma) - 2c_1\alpha\beta,\\
-c &= a_1\gamma^2 + b_1\alpha\gamma - c_1\alpha^2,
\end{aligned}}
and further
\weq{15}{b^2 + 4ac = b_1^2 + 4a_1c_1 = D.}
Since $a$, $b$, $c$ are without common divisor, $a_1$, $b_1$, $c_1$ also have no common divisor, as one sees from (14), and $a_1$, $b_1$, $c_1$ have the same meaning for $\omega_1$ as $a$, $b$, $c$ have for $\omega$.

\gesperrt{VII. One sees further from (15) that for equivalent numbers not only the irrationality $\sqrt{d}$, but also the discriminant is the same.}

If one sets
\weq{16}{\omega_1 = x_1 + y_1\sqrt{d} = \frac{\alpha(x+y\sqrt{d})+\beta}{\gamma(x+y\sqrt{d})+\delta},}
and, in order to determine $x_1$, $y_1$, multiplies the numerator and denominator of the last fraction by $\gamma(x-y\sqrt{d})+\delta$, then by (4) and (13) there results
\weq{17}{2x_1 = \frac{b_1}{c_1}, \qquad y_1 = \frac{\varepsilon y c}{c_1},}
hence, if the $\sqrt{D}$ is understood in the same way as in (7) and (8), by (9):
\weq{18}{\omega_1 = \frac{b_1 + \varepsilon\sqrt{D}}{2c_1}.}

\begin{center}
\S.\ 123.

\gesperrt{Reduced Numbers with Negative Discriminant.} % [sic] section number \S.\,123 repeats the number already used earlier on this same page for ``Reducirte Zahlen'' as printed
\end{center}

The considerations up to this point are applicable uniformly to the two cases of real and imaginary quadratic irrational numbers, that is, to positive and negative discriminants. Now, however, the two cases must be separated from one another, and we begin with the negative discriminants. We are concerned, then, with imaginary numbers:
\weq{1}{\omega = x + y\sqrt{d} = \xi + i\eta,}
in which $\xi$ and $\eta$ are real; it suffices to consider one of the two conjugate numbers $\xi+i\eta$, $\xi-i\eta$. We therefore assume $\eta$ to be positive and lay down the following definition:
%CONT>>


% ===== scan leaf 405, printed 380 =====
% !PRINTED 380
\weberpage{380}
%<<CONT

A complex number $\omega = \xi + i\eta$ with positive $\eta$ is called \gesperrt{reduced} if
\weq{2}{-\tfrac12 \leqq \xi \leqq \tfrac12, \quad \xi^2+\eta^2 \geqq 1.}
If one regards $\xi$, $\eta$ as rectangular coordinates in a plane, then every number $\omega$ is represented by a point of this plane, and the position of the points that correspond to the reduced numbers is represented by the region shaded in Fig. 26 (with inclusion of the boundaries), which we call the fundamental region.

% SOURCE FIGURE CROP: Bd. I registered scan 406 (printed p. 380), 600 dpi,
% crop x=350,y=1005,w=1180,h=900; source-faithful raster, Fig. 26.
\weberfigureimage{weber-b1-fig26.png}{0.48}

From the conditions (2) it then further follows, as is also easily confirmed in the figure, that
\weq{3}{\eta^2 \geqq \tfrac34.}
If one now passes from a number $\omega$ to an equivalent number $\omega_1$, formula (16) of the preceding section gives
\weq{4}{\omega_1 = \frac{(\alpha\xi+\beta)+i\alpha\eta}{(\gamma\xi+\delta)+i\gamma\eta} = \frac{[(\alpha\xi+\beta)(\gamma\xi+\delta)+\alpha\gamma\eta^2]+i\varepsilon\eta}{(\gamma\xi+\delta)^2+\gamma^2\eta^2},}
and this formula shows that $\eta_1$ has the same sign as $\eta$ or the opposite one, according as $\varepsilon = +1$ or $=-1$, that is, according as the equivalence is a proper or an improper one.

If we restrict ourselves, then, to numbers with positive imaginary part, only proper equivalence comes into consideration, and we now prove the fundamental theorem:

1. \gesperrt{Every imaginary quadratic irrational number $x+y\sqrt{d}$ with positive imaginary part is equivalent to a reduced number.}

If we understand by $\alpha$ the integer nearest to $\xi$, then $\omega-\alpha$ will satisfy the first of the conditions (2), namely that $-\tfrac12 \leqq \xi-\alpha \leqq \tfrac12$.

If now
\weq{5}{r^2 = (\xi-\alpha)^2+\eta^2 = (\omega-\alpha)(\omega'-\alpha)}
is greater than or equal to 1, then $\omega-\alpha$, which is equivalent to $\omega$, is already reduced; otherwise we set, applying a continued-fraction expansion,
%CONT>>


% ===== scan leaf 406, printed 381 =====
% !PRINTED 381
\weberpage{381}

%% [running head, not body text -- commented out 2026-08-07] Page furniture (running
%% head, folio, signature) is excluded from this edition by standing policy.
%% \noindent\S.\,123. \hfill \textsc{Reduced Numbers.} \hfill 381

%<<CONT
\weq{6}{\omega = \alpha - \dfrac{1}{\omega_1},}

so that $\omega_1 = \xi_1 + i\eta_1$ is also equivalent to $\omega$. We now treat the
number $\omega_1$ again in the same way as $\omega$, by setting, if $\omega_1 - \alpha_1$
is not yet reduced,

\weq{7}{\omega_1 = \alpha_1 - \dfrac{1}{\omega_2}}

and so on. If we consider the sequence of the quantities formed on the analogy of (5),
\[
r_1^2 = (\xi_1-\alpha_1)^2 + \eta_1^2, \quad r_2^2 = (\xi_2-\alpha_2)^2 + \eta_2^2 \ldots,
\]
then all that now matters is to show that after a finite number of steps
of this kind we come to an $r_\nu^2$ that is equal to or greater
than 1.

But by (6) we have
\[
r^2 = (\omega-\alpha)(\omega'-\alpha) = \dfrac{1}{\xi_1^2+\eta_1^2},
\]
and from this the comparison of the imaginary parts on the two
sides of (6) gives:
\[
\eta = r^2 \eta_1.
\]
In the same way there follows the sequence of equations

\weq{8}{\eta = r^2\eta_1, \quad \eta_1 = r_1^2\eta_2, \quad \eta_2 = r_2^2\eta_3 \ldots}

Now, however, $i\eta = y\sqrt{d}$, and hence by the formula (6), \S.\,122
\[
i\eta = \dfrac{\sqrt{D}}{2c},
\]
in which $c$ is an integer. But because equivalent numbers have the same
discriminant, it follows in the same way:
\[
i\eta_1 = \dfrac{\sqrt{D}}{2c_1}, \quad i\eta_2 = \dfrac{\sqrt{D}}{2c_2}, \ldots,
\]
in which $c, c_1, c_2 \ldots$ is a sequence of positive integers. From this
one obtains from (8):

\weq{9}{c_1 = r^2 c, \quad c_2 = r_1^2 c_1, \quad c_3 = r_2^2 c_2 \ldots}

But so long as the numbers $r^2, r_1^2 \ldots r_\nu^2$ are less than 1, it follows from this

\weq{10}{c > c_1 > c_2 \ldots > c_{\nu+1},}

and because there can be only a finite number of positive numbers that are
less than $c$, this sequence must terminate, and our Theorem 1.
is thereby proved.


% ===== scan leaf 407, printed 382 =====
% !PRINTED 382
\weberpage{382}

%% [running head, not body text -- commented out 2026-08-07] Page furniture (running
%% head, folio, signature) is excluded from this edition by standing policy.
%% \noindent 382 \hfill \textsc{Eleventh Section.} \hfill \S.\,123.

The reduced numbers whose images lie on the boundary of the fundamental
region are equivalent in pairs, namely

\noindent a)\quad $\omega = -\tfrac{1}{2} + i\eta$ \quad and \quad $\omega_1 = \omega + 1 = \tfrac{1}{2} + i\eta$

\noindent ($a$ and $a'$ in the figure),

\noindent b)\quad $\omega = -\xi + i\eta$ \quad and \quad $\omega_1 = -\dfrac{1}{\omega} = \xi + i\eta$

\noindent when $\xi^2+\eta^2 = 1$ ($b$ and $b'$ in the figure). But there holds further the
theorem:

\noindent 2.\quad \gesperrt{Apart from the cases a), b), no two reduced numbers are equivalent.}

For let us assume two reduced numbers that are not identical,

\weq{11}{\omega = \xi + \eta i, \quad \omega_1 = \xi_1 + \eta_1 i,}

and suppose, which does not impair the generality,

\weq{12}{\eta_1 \geqq \eta,}

then it follows from formula (4):

\weq{13}{\eta_1 = \dfrac{\eta}{(\gamma\xi+\delta)^2 + \gamma^2\eta^2},}

hence, on account of (12):

\weq{14}{(\gamma\xi+\delta)^2 + \gamma^2\eta^2 \leqq 1,}

hence also, by (3):

\weq{15}{1 \geqq \gamma^2\eta^2 \geqq \tfrac{3}{4}\gamma^2;}

thus there are two possibilities:
\[
\alpha)\ \gamma = 0, \qquad \beta)\ \gamma = \pm 1.
\]

In the case $\alpha$) we have $\alpha\delta = 1$, hence $\alpha = \delta = \pm 1$; from (12) and (13)
it results that $\eta_1 = \eta$, and further from (4):
\[
\xi_1 = \xi \pm \beta.
\]
But this is compatible with the conditions (2) only when
either $\xi_1 = \xi$, $\beta = 0$ and hence $\omega_1$ is identical with $\omega$, or when
$\beta = \pm 1$ and $\xi_1 = -\xi = \pm \tfrac{1}{2}$, that is, in the case a).

In the case $\beta$) we have by (14) $(\xi \pm \delta)^2 \leqq 1$, and hence \gesperrt{either}:
$\delta = 0$, $\beta\gamma = -1$ and $\xi^2+\eta^2 = 1$, $\eta_1 = \eta$, further by (4)
$\xi_1 = \pm\alpha - \xi$; hence, since $\xi_1 = \xi$ is excluded, $\xi_1 = -\xi$,
$\alpha = 0$, which is the exceptional case b),

\gesperrt{or}: $\delta = \pm 1$, hence $(\xi\pm 1)^2 + \eta^2 \leqq 1$, and consequently by (3)
$(\xi\pm 1)^2 \leqq \tfrac{1}{4}$. But since $(\xi\pm 1)$ must by (2) be at least
equal to $\tfrac{1}{2}$ in absolute value, it follows from this that $\xi = \mp\tfrac{1}{2}$,
$\eta^2 = \tfrac{3}{4}$; hence by (13) $\eta_1^2 = \eta^2 = \tfrac{3}{4}$, and by
%CONT>>


% ===== scan leaf 408, printed 383 =====
% !PRINTED 383
\weberpage{383}
\pdfbookmark[2]{§ 124. Reduced numbers with positive discriminant}{bm138}

%% [running head, not body text -- commented out 2026-08-07] Page furniture (running
%% head, folio, signature) is excluded from this edition by standing policy.
%% \noindent\S.\,124. \hfill \textsc{Reducirte Zahlen.} \hfill 383

%<<CONT
(2) $\xi_1^2 = \tfrac{1}{4}$, $\xi_1 = \mp\tfrac{1}{2}$, which is contained both under case a) and under
case b);

and with this, then, Theorem 2. is proved. To this we now add, as a
third theorem:

\noindent 3.\quad \gesperrt{For a given negative discriminant $D$ there is only a finite number of reduced irrational numbers.}

From the representation of the irrational numbers \S.\,122, (7):
\[
\omega = \dfrac{\sqrt{D}+b}{2c}, \qquad \xi = \dfrac{b}{2c}, \qquad \eta = \dfrac{\sqrt{-D}}{2c}
\]
and from the inequality (3) $\eta^2 \geqq \tfrac{3}{4}$ it follows that
\[
\dfrac{-D}{4c^2} \geqq \tfrac{3}{4},
\]
or
\[
c \leqq \sqrt{\dfrac{-D}{3}};
\]
thus, when $D$ is given, there is only a finite number of values
of the positive integer $c$. Further, from $-\tfrac{1}{2} \leqq \xi \leqq \tfrac{1}{2}$ it follows that:
\[
-c \leqq b \leqq c.
\]
Thus for each value of $c$ there is only a finite number of values
for $b$, q.\,e.\,d.

If one unites all numbers equivalent to one another into a
number class, then by 1. we have accordingly proved:

\noindent 4.\quad \gesperrt{The number of classes of imaginary quadratic irrational numbers of a given discriminant is finite.}

\begin{center}
\S.\,124.

\gesperrt{Reduced numbers with positive discriminant.}
\end{center}

We now have to carry out a similar investigation for a
positive $d$, that is, for \gesperrt{real quadratic irrational numbers}. Such a
number $\omega = x + y\sqrt{d}$ shall be called \gesperrt{reduced} if it satisfies the following
conditions:
%CONT>>


% ===== scan leaf 409, printed 384 =====
% !PRINTED 384
\weberpage{384}

%% [running head, not body text -- commented out 2026-08-07] Page furniture (running
%% head, folio, signature) is excluded from this edition by standing policy.
%% \noindent 384 \hfill \textsc{Eleventh Section.} \hfill \S.\,124.

%<<CONT
% [source-reviewed 2026-08-26] The faint mark before the first word at the top edge is a scan speck, not body text; no text is transcribed.
\gesperrt{$\omega$ is positive and greater than 1, and the number $\omega'$
conjugate to $\omega$ is negative and smaller than 1 in absolute value,}

or in symbols:

\weq{1}{0 < y\sqrt{d} - x < 1 < y\sqrt{d} + x.}

If one interprets $x, y$ in a plane as rectangular coordinates, so
that the points of the plane become the images of the numbers $\omega$, then
the images of the reduced numbers lie in that part of the plane which is
shaded in Fig. 27, which corresponds to the assumption $d = 2$.

% SOURCE FIGURE CROP: Bd. I registered scan 410 (printed p. 384), 600 dpi,
% crop x=440,y=1230,w=1160,h=870; source-faithful raster, Fig. 27.
\weberfigureimage{weber-b1-fig27.png}{0.48}

We here take proper and improper equivalence together
and first prove the theorems:

\noindent 1.\quad \gesperrt{Every number $\omega$ is equivalent to a reduced number.}

One sees this easily if one develops $\omega$ into a continued fraction
and goes as far as the closing number $\omega_n$:

\weq{2}{\omega = (\alpha_0, \alpha_1, \alpha_2 \ldots \alpha_{n-1}, \omega_n).}

Then $\omega$ is equivalent to $\omega_n$, and $\omega_n$, as soon as $n$ is greater than 0,
is positive and greater than 1. Since $\omega_n$ is equivalent to $\omega$, it only
remains to prove that, when $n$ is sufficiently large,
the number $\omega'_n$ conjugate to $\omega_n$ is negative and absolutely smaller than 1.
But this one sees from the following formulae: From (2) it follows, if
$P_n : Q_n$ are the convergents,
\[
\omega = \dfrac{P_n \omega_n + P_{n-1}}{Q_n \omega_n + Q_{n-1}},
\]
or by solving

\weq{3}{\omega_n = -\dfrac{Q_{n-1}\omega - P_{n-1}}{Q_n\omega - P_n} = -\dfrac{Q_{n-1}}{Q_n} - \dfrac{(-1)^n}{Q_n(Q_n\omega - P_n)}.}

If in this equation we change $\sqrt{d}$ into $-\sqrt{d}$, then at the same time
$\omega$ passes into $\omega'$, $\omega_n$ into $\omega'_n$, and from (3) we obtain

\weq{4}{\omega'_n = -\dfrac{Q_{n-1}\omega' - P_{n-1}}{Q_n\omega' - P_n} = -\dfrac{Q_{n-1}}{Q_n} \dfrac{\omega' - \dfrac{P_{n-1}}{Q_{n-1}}}{\omega' - \dfrac{P_n}{Q_n}},}
%CONT>>


% ===== scan leaf 410, printed 385 =====
% !PRINTED 385
\weberpage{385}

\weq{5}{\omega_n' + 1 = \frac{1}{Q_n}\left\{ Q_n - Q_{n-1} - \frac{(-1)^n}{Q_n\left(\omega' - \dfrac{P_n}{Q_n}\right)} \right\}.}

If now $n$ is taken sufficiently large, then
\[
\omega' - \frac{P_n}{Q_n}, \quad \omega' - \frac{P_{n-1}}{Q_{n-1}}
\]
% [sic] The source prints $\omega'-\omega=2y\sqrt d$; from (1)--(2) of §122
% the left side is actually $-2y\sqrt d$.  The proof here uses only non-vanishing.
differ arbitrarily little from the non-vanishing difference $\omega' - \omega = 2y\sqrt{d}$, whence (4) shows that $\omega_n'$ is negative. But from (5) it follows that $\omega_n' + 1$ is positive, since $Q_n - Q_{n-1}$, being a positive integer, is at least equal to 1, and
\[
\frac{1}{Q_n\left(\omega' - \dfrac{P_n}{Q_n}\right)}
\]
becomes arbitrarily small for a sufficiently large $n$; and with this the proof of Theorem 1. is complete.

Since the conjugate number of the conjugate $\omega'$ is the original number $\omega$, it follows from the definition of the reduced numbers \gesperrt{that $\omega$ and $-1:\omega'$ are reduced simultaneously.}

\gesperrt{2. If $\omega$ is reduced and, for an integer $\alpha$,}

\weq{6}{\omega = \alpha + \frac{1}{\omega_1},}

\gesperrt{then $\omega_1$ is reduced if and only if $\alpha$ is the greatest integer contained in $\omega$, that is, if}

\weq{7}{\alpha < \omega < \alpha + 1.}

If the condition (7) is fulfilled, then $\omega_1$ is positive and greater than 1, and

\weq{8}{\omega_1' = \frac{1}{\omega' - \alpha}}

is, since $\omega'$ is negative, a negative proper fraction; hence $\omega_1$ is reduced. On the other hand, $\omega_1$ cannot be reduced if $\alpha > \omega$, since then $\omega_1$ would be negative, or if $\alpha + 1 < \omega$, since otherwise $\omega_1$ would be less than 1.

If one puts (8) in the form

\weq{9}{-\frac{1}{\omega_1'} = \alpha + \frac{1}{-\dfrac{1}{\omega'}},}
%CONT>>


% ===== scan leaf 411, printed 386 =====
% !PRINTED 386
\weberpage{386}
\pdfbookmark[2]{§ 125. Expansion of real quadratic irrational numbers in continued fractions}{bm139}

%<<CONT
then, since $-1:\omega'$ and $-1:\omega_1'$ are also reduced, $\alpha$ is the greatest integer contained in $-1:\omega_1'$, and $\omega'$ is completely determined by $\omega_1'$ according to (9), and hence also $\omega$ by $\omega_1$.

From this it follows:

\gesperrt{3. If a reduced number $\omega$ is developed in a continued fraction, then all the closing numbers $\omega_n$ are again reduced numbers, and by any one of these closing numbers $\omega_n$ both the following $\omega_{n+1}$ and the preceding $\omega_{n-1}$ are completely determined.}

\begin{center}
\S.\,125.

\gesperrt{Development of real quadratic irrational numbers}\\
\gesperrt{in continued fractions.}
\end{center}

\gesperrt{4. For a given discriminant there is only a finite number of reduced numbers.}

The correctness of this theorem is easily seen if, according to \S.\,122, one puts the irrational number $\omega$ in the form:

\weq{1}{\omega = \frac{b+\sqrt{D}}{2c} = \frac{-2a}{b-\sqrt{D}}, \quad D = b^2 + 4ac.}

The conjugate number is

\weq{2}{\omega' = \frac{b-\sqrt{D}}{2c} = \frac{-2a}{b+\sqrt{D}}.}

If $\omega$ is to be reduced, then

\weq{3}{0 < \frac{\sqrt{D}-b}{2c} < 1 < \frac{\sqrt{D}+b}{2c}.}

Hence $\sqrt{D}$ must have the same sign as $c$; if we take it to be positive, then

\weq{4}{0 < \sqrt{D} - b < 2c < \sqrt{D} + b.}

A second form of these conditions we obtain from the second representation

\weq{5}{0 < \sqrt{D} - b < 2a < \sqrt{D} + b.}

Hence $a$, $b$, $c$ are positive and $b$ is less than $\sqrt{D}$.

If $\lambda$ denotes the greatest integer contained in $\sqrt{D}$, so that
\[
\lambda < \sqrt{D} < \lambda + 1,
\]
%CONT>>


% ===== scan leaf 412, printed 387 =====
% !PRINTED 387
\weberpage{387}

%<<CONT
so $b$ has one of the values 1, 2, 3 \ldots $\lambda$, but with the restriction that for even $D$ only the even, and for odd $D$ only the odd among these numbers are to be taken for $b$. Then $a$ and $c$ are to be determined so that

\weq{6}{\frac{D-b^2}{4} = ac,}

and so that

\weq{7}{\frac{\lambda-b+1}{2} \leqq c \leqq \frac{\lambda+b}{2},}

and $a$ must lie between the same limits. Only one of these two limiting values is an integer; the other may be replaced, for the lower limit, by the next greater integer, and for the upper limit, by the next smaller. Thus for $b$ only a finite number of values is admissible; then $(D-b^2):4$ is decomposable into two factors in only a finite number of ways, and of these decompositions only those are to be retained in which both factors satisfy the condition (7). Moreover, such combinations are also to be discarded in which $a$, $b$, $c$ acquire a common divisor. Herein lies the means of actually determining all reduced numbers for a given $D$.

For what follows we shall, for the numbers $\omega$ so determined, use the symbol

\weq{8}{\omega = \frac{\sqrt{D}+b}{2c} = \frac{2a}{\sqrt{D}-b} = \{a, b, c\}}

to denote them.

From 3. \S.\,124 and 4. \S.\,125 there now results the following important theorem:

\gesperrt{5. The continued-fraction expansion of a reduced number $\omega$ is periodic,}

that is, if $\omega$ is expanded in a continued fraction

\weq{9}{\omega = (\alpha_0, \alpha_1, \alpha_2 \ldots)}

then after a certain number of partial denominators the same partial denominators always recur in the same order; thus, if the period consists of $\nu$ terms, then
\[
\begin{array}{ccccc}
\alpha_0 &=& \alpha_\nu &=& \alpha_{2\nu} \ldots\\
\alpha_1 &=& \alpha_{\nu+1} &=& \alpha_{2\nu+1} \ldots\\
\multicolumn{5}{c}{\cdots\cdots\cdots\cdots\cdots\cdots\cdots\cdots\cdots}\\
\alpha_{\nu-1} &=& \alpha_{2\nu-1} &=& \alpha_{3\nu-1} \ldots,
\end{array}
\]
%CONT>>


% ===== scan leaf 413, printed 388 =====
% !PRINTED 388
\weberpage{388}

%<<CONT
and it therefore suffices to denote the whole continued fraction by the period, say thus:

\weq{10}{\omega = [\alpha_0, \alpha_1 \ldots \alpha_{\nu-1}].}

This theorem is evidently equivalent to saying that the closing number $\omega_\nu$ of the continued fraction (9) is identical with $\omega$ itself, and in this also lies the proof of the assertion. For the closing numbers $\omega_\nu$, being equivalent numbers, all belong to the same discriminant, and hence by 4. the number of all possible closing numbers is only a finite one. In the sequence of closing numbers $\omega, \omega_1, \omega_2 \ldots$ one that has already occurred must therefore at some point appear for the second time. If $\omega_\nu$ is the first that appears for the second time, then $\omega_\nu = \omega$ must hold, since otherwise by 3. \S.\,124 the number $\omega_{\nu-1}$ would also appear for the second time.

Thus all the reduced numbers belonging to one discriminant can be arranged in periods:

\weq{11}{\omega, \omega_1, \omega_2 \ldots \omega_{\nu-1},}

so that, if
\[
\alpha_0, \alpha_1, \alpha_2 \ldots \alpha_{\nu-1}
\]
are the greatest integers contained in them, we have
\[
\begin{array}{l}
\omega = [\alpha_0, \alpha_1 \ldots \alpha_{\nu-1}]\\
\omega_1 = [\alpha_1, \alpha_2 \ldots \alpha_0]\\
\cdots\cdots\cdots\cdots\cdots\cdots\cdots\cdots
\end{array}
\]
by which the continued fraction expansion is given for each of these numbers.

If the whole system of reduced numbers is not yet exhausted by \gesperrt{one} period, one forms a second period, and so on. Since by a number $\omega$ both the number preceding it in the period and the one following it are completely determined, two different periods never contain a common number.

Of these periods there now holds the theorem:

\gesperrt{6. Reduced numbers from the same period are equivalent; from different periods they are not equivalent.}

The first part of the assertion is clear from the outset, since reduced numbers of the same period are closing numbers of one another; the second part is proved by the theorem \S.\,121, VI, that equivalent numbers, when the continued fraction expansion is carried sufficiently
%CONT>>


% ===== scan leaf 414, printed 389 =====
% !PRINTED 389
\weberpage{389}
%<<CONT
fraction development finally obtain the same closing numbers.
If, then, both continued fractions are periodic, they must belong
to the same period.

If for a given discriminant $D$ one has developed the complete
system of reduced numbers according to the rules given at the
beginning of the section, it is easy to arrange these numbers
in periods and to find the partial denominators $\alpha$ of the continued
fraction. Namely, let
\weq{12}{\omega = \frac{\sqrt{D}+b}{2c} = \frac{2a}{\sqrt{D}-b}}
be one of these numbers and
\weq{13}{\omega_1 = \frac{\sqrt{D}+b_1}{2c_1} = \frac{2a_1}{\sqrt{D}-b_1}}
the number immediately following it in the period, so that
\weq{14}{\omega = \alpha + \frac{1}{\omega_1}, \quad \omega_1 = \frac{1}{\omega-\alpha}}
where $\alpha$ denotes the greatest integer contained in $\omega$.
By (12) and (14) we then have
\[
\omega_1 = \frac{2c}{\sqrt{D}+b-2c\alpha},
\]
and comparison with (13) gives
\weq{15}{a_1 = c, \quad b_1 = 2c\alpha - b.}
If, conversely, two of the reduced numbers $\{a, b, c\}$, $\{a_1, b_1, c_1\}$
belonging to $D$ stand in the relation represented by (15),
in which $\alpha$ is any integer, then the second immediately
follows the first in the period. For from (15)
follows (14), and hence by 2. $\alpha$ must be the greatest number
contained in $\omega$.

The number $\omega_1$ is therefore completely and uniquely determined
from $\omega$ by the conditions (15).

7. \gesperrt{One therefore arranges the numbers $\{a, b, c\}$ from left
to right in such a way that the last number $c$
of the preceding one becomes at the same time the first number $a_1$ of the
following, and that the sum of the two
middle numbers $b + b_1$ is divisible by $2c$.
This arrangement is possible in only one way,
and the quotient $b + b_1 : 2c$ is the number $\alpha$ that
appears as partial denominator in the continued fraction.}


% ===== scan leaf 415, printed 390 =====
% !PRINTED 390
\weberpage{390}
Let us now consider any primitive quadratic equation with
integer coefficients
\weq{16}{A + B\Omega + C\Omega^2 = 0,}
in which $A$, $B$, $C$ are integers without common divisor,
satisfying the condition
\[
B^2 - 4AC = D
\]
The two roots $\Omega$, $\Omega'$ of this equation are quadratic
irrational numbers belonging to the discriminant $D$, and
if we develop them in continued fractions, then by
\S.\,124 we shall at last come upon closing numbers $\omega_1$, $\omega_2$ which belong to the reduced
numbers, and which must therefore be contained in the periods
discussed above.

It remains to determine whether these closing numbers $\omega_1$, $\omega_2$ are
contained in the same period or in different periods.

By 6. they are contained in the same period if and only
if $\Omega$ and $\Omega'$ are equivalent to each other. If we call
such quadratic irrational numbers as are equivalent to their conjugates
\gesperrt{two-sided numbers}\footnote{``\gesperrt{Two-sided}'' we use, following \gesperrt{Dedekind}'s proposal, in place of the
\gesperrt{Gauss}ian \gesperrt{anceps} or the otherwise usual \gesperrt{ambig}: \gesperrt{Dirichlet-Dedekind}, Vorlesungen über Zahlentheorie. 4th edition, 1894, p. 139.}, then we can
say that the continued-fraction development of the roots of
equation (16) leads to one or to two different periods,
according as the roots are two-sided or not.
Which of the two cases occurs can be recognised from the period
itself.

If, namely, $\omega$ is a reduced number, and $\omega'$ conjugate
to $\omega$, then $-1 : \omega'$ is also reduced; and if $\Omega$ is equivalent
to $\omega$, then $\Omega'$ is equivalent to $\omega'$, and hence also to $-1 : \omega'$ (\S.\,120).

If
\[
\omega = \{a, b, c\},
\]
then by (12):
\[
-\frac{1}{\omega'} = \{c, b, a\}.
\]

By \S.\,124, (9), $-1 : \omega'$ therefore follows $-1 : \omega_1'$
in the same way as $\omega_1$ follows $\omega$. If therefore the period of $\omega$ is
\[
\omega, \omega_1, \omega_2 \ldots \omega_{\nu-1},
\]
then the period of $-1 : \omega'$ is:
\[
\frac{-1}{\omega'_{\nu-1}}, \frac{-1}{\omega'_{\nu-2}} \cdots \frac{-1}{\omega_1'}, \frac{-1}{\omega'}.
\]


% ===== scan leaf 416, printed 391 =====
% !PRINTED 391
\weberpage{391}
If the period of the continued fraction development of $\omega$ is
\weq{17}{[\alpha_0, \alpha_1 \ldots \alpha_{\nu-1}],}
then for $-1 : \omega'$ it is
\weq{18}{[\alpha_{\nu-1}, \alpha_{\nu-2} \ldots \alpha_0],}
that is, the reverse.

Now if $\omega$ is equivalent to $\omega'$, and hence also to $-1 : \omega'$, then
by 6. these two periods must agree with one another when one
begins at a suitable term, or the period of
$\omega$ must be \gesperrt{reversible}.

This means: in the continued fraction development
for $\omega$ two elements must admit of being so chosen that the development
reads the same whether one reads it proceeding from the one of these numbers to
the left or from the other to the right.
In symbols: for some suitably determined $k$ and
for $i = 0, 1, 2 \ldots \nu - 1$ we must have
\[
\alpha_i = \alpha_{\nu+k-i}
\]
If the period is reversible, then conversely
(by 6.) $\omega$ is equivalent to $-1 : \omega'$, and hence also to $\omega'$. From this,
then, the result:

8. \gesperrt{Two-sided numbers have in their continued
 fraction development a reversible period, and conversely
 numbers which have a reversible period
 in their continued fraction development
are two-sided.}

Finally we wish to remark further that by the property
of admitting development into a periodic continued fraction,
the real quadratic irrational numbers are distinguished from all other
numbers. For if $\omega$ is a number developed into a continued
fraction, and $\omega_n$, $\omega_m$ are two closing numbers
of this continued fraction, then
\[
\omega = \frac{P_n\omega_n + P_{n-1}}{Q_n\omega_n + Q_{n-1}} = \frac{P_m\omega_m + P_{m-1}}{Q_m\omega_m + Q_{m-1}}.
\]

Now if the continued fraction is periodic, no matter whether
the periodicity begins right from the start or only in the course
of the development, then for two different values of $n$
and $m$ we have $\omega_n = \omega_m$, and one thus obtains, by elimination of
%CONT>>


% ===== scan leaf 417, printed 392 =====
% !PRINTED 392
\weberpage{392}
\pdfbookmark[2]{§ 126. Examples}{bm140}
%<<CONT
$\omega_n$ a quadratic equation for $\omega$, which can be put in the form:
\[
\begin{vmatrix} P_n - \omega Q_n, & P_{n-1} - \omega Q_{n-1} \\ P_m - \omega Q_m, & P_{m-1} - \omega Q_{m-1} \end{vmatrix} = 0.
\]

\begin{center}
\S.\ 126.

Examples.
\end{center}

We shall here carry out some examples of the determination of the reduced
numbers and periods, in order to show the application
of the method; the examples can of course be multiplied
quite at pleasure.

1. $D = 29$.

We first determine, according to \S.\,125, 4., all the
reduced numbers belonging to $D$. The greatest integer $\lambda$
contained in $\sqrt{D}$ is here $= 5$; hence $b$ can have only the values:
\[
b = 1, 3, 5,
\]
whence
\[
ac = \frac{D-b^2}{4} = 7, 5, 1.
\]

The limits within which $a$ and $c$ must lie are, by
\S.\,125, (7), for the three values of $b$:
\[
3, 3; \quad 2, 4; \quad 1, 5.
\]
One therefore obtains only the single reduced number:
\[
\{1, 5, 1\}.
\]

The period consists of a single term, and the
continued fraction for $\omega$ is:
\[
\omega = \frac{5+\sqrt{29}}{2} = (5, 5, 5 \ldots)
\]

2. $D = 116 = 4.29$.

Here $\lambda = 10$ and $b$ has one of the values
\[
b = 2, 4, 6, 8, 10,
\]
hence
\[
ac = 28, 25, 20, 13, 4;
\]
the limits for $a$ and $c$ are
\[
5, 6; \quad 4, 7; \quad 3, 8; \quad 2, 9; \quad 1, 10,
\]
hence, since the cases in which $a$, $b$, $c$ have the common factor
%CONT>>


% ===== scan leaf 418, printed 393 =====
% !PRINTED 393
\weberpage{393}
%<<CONT
2, are still to be omitted, the reduced numbers are found
to be:
\[
\{5, 4, 5\}, \quad \{4, 6, 5\}, \quad \{5, 6, 4\}, \quad \{1, 10, 4\}, \quad \{4, 10, 1\},
\]
which arrange themselves into a period as follows, the first term
being appended once more at the end:
\[
\{4, 10, 1\}, \quad \{1, 10, 4\}, \quad \{4, 6, 5\}, \quad \{5, 4, 5\}, \quad \{5, 6, 4\}, \quad \{4, 10, 1\},
\]
and for $\omega$ the continued-fraction period is found:
\[
5 + \sqrt{29} = [10, 2, 1, 1, 2].
\]
If one begins with the fourth term, that is, considers the continued
fraction $[1, 2, 10, 2, 1]$, which is the expansion of $\dfrac{2 + \sqrt{29}}{5}$,
one sees that the period is \gesperrt{reversible}.

3. $D = 76 = 4 \ldotp 19, \quad \lambda = 8$
% [resolved 2026-08-06, 5x zoom strip 2/5] Confirmed. The line reads "3. D = 76 = 4. 19,
% lambda = 8": the mark is Weber's raised multiplication period, and the arithmetic settles
% it -- 4 x 19 = 76. Transcription stands. Same strip confirms b = 2, 4, 6, 8 and
% ac = 18, 15, 10, 3.
\[
b = \quad 2, \quad 4, \quad 6, \quad 8
\]
\[
a\,c = 18, \quad 15, \quad 10, \quad 3.
\]
\gesperrt{Limits} for $a$ and $c$:
\[
4, 5; \quad 3, 6; \quad 2, 7; \quad 1, 8;
\]
hence the reduced numbers are:
\[
\{3, 4, 5\}, \quad \{5, 4, 3\}, \quad \{2, 6, 5\}, \quad \{5, 6, 2\}, \quad \{1, 8, 3\}, \quad \{3, 8, 1\}.
\]
We obtain a period:
\[
\{3, 8, 1\}, \quad \{1, 8, 3\}, \quad \{3, 4, 5\}, \quad \{5, 6, 2\}, \quad \{2, 6, 5\}, \quad \{5, 4, 3\}, \quad \{3, 8, 1\}
\]
and the period of the continued fraction for $4 + \sqrt{19}$ becomes
\[
[8, 2, 1, 3, 1, 2].
\]
This continued fraction too has a reversible period, since
one obtains the same succession of numbers whether one reads
from the first 2 to the right or from the second 2 to the left.

4. $D = 37, \qquad \lambda = 6,$
\[
b = 1, \quad 3, \quad 5
\]
\[
a\,c = 9, \quad 7, \quad 3;
\]
\gesperrt{Limits} for $a$, $c$:
\[
3, 3; \quad 2, 4; \quad 1, 5;
\]
\gesperrt{reduced numbers}:
\[
\{3, 1, 3\}, \quad \{1, 5, 3\}, \quad \{3, 5, 1\};
\]
\gesperrt{Period}:
\[
\{3, 5, 1\}, \quad \{1, 5, 3\}, \quad \{3, 1, 3\}, \quad \{3, 5, 1\};
\]
\gesperrt{Continued fraction} with reversible period $[1, 5, 1]$:
%CONT>>


% ===== scan leaf 419, printed 394 =====
% !PRINTED 394
\weberpage{394}
%<<CONT
\[
\frac{5 + \sqrt{37}}{2} = [5, 1, 1].
\]
5. $D = 148 = 4 \ldotp 37 \quad \lambda = 12$
\[
b = \quad 2, \quad 4, \quad 6, \quad 8, \quad 10, \quad 12
\]
\[
a\,c = 36, \quad 33, \quad 28, \quad 21, \quad 12, \quad 1;
\]
\gesperrt{limits} for $a$, $c$:
\[
6, 7; \quad 5, 8; \quad 4, 9; \quad 3, 10; \quad 2, 11; \quad 1, 12;
\]
\gesperrt{reduced numbers}:
\[
\{4, 6, 7\}, \quad \{7, 6, 4\}, \quad \{3, 8, 7\}, \quad \{7, 8, 3\}
\]
\[
\{3, 10, 4\}, \quad \{4, 10, 3\}, \quad \{1, 12, 1\}.
\]
Here we obtain three periods of 1, 3, and 3 terms:
\[
\{1, 12, 1\}, \quad \{1, 12, 1\},
\]
\[
\{3, 10, 4\}, \quad \{4, 6, 7\}, \quad \{7, 8, 3\}, \quad \{3, 10, 4\},
\]
\[
\{4, 10, 3\}, \quad \{3, 8, 7\}, \quad \{7, 6, 4\}, \quad \{4, 10, 3\}
\]
and the continued-fraction periods
\[
[12], \quad [2, 1, 3], \quad [3, 1, 2];
\]
of these the first is reversible, and the other two pass into one
another under reversal.

6. $D = 136 = 4 \ldotp 34, \quad \lambda = 11,$
\[
b = \quad 2, \quad 4, \quad 6, \quad 8, \quad 10,
\]
\[
a\,c = 33, \quad 30, \quad 25, \quad 18, \quad 9;
\]
\gesperrt{limits} for $a$, $c$:
\[
5, 6; \quad 4, 7; \quad 3, 8; \quad 2, 9; \quad 1, 10;
\]
\gesperrt{reduced numbers}:
\[
\{5, 4, 6\}, \quad \{6, 4, 5\}, \quad \{5, 6, 5\}, \quad \{2, 8, 9\}, \quad \{9, 8, 2\}
\]
\[
\{3, 8, 6\}, \quad \{6, 8, 3\}, \quad \{1, 10, 9\}, \quad \{9, 10, 1\}, \quad \{3, 10, 3\};
\]
\gesperrt{periods}:
\[
\{5, 4, 6\}, \quad \{6, 8, 3\}, \quad \{3, 10, 3\}, \quad \{3, 8, 6\}, \quad \{6, 4, 5\}, \quad \{5, 6, 5\}, \quad \{5, 4, 6\},
\]
\[
\{9, 10, 1\}, \quad \{1, 10, 9\}, \quad \{9, 8, 2\}, \quad \{2, 8, 9\}, \quad \{9, 10, 1\};
\]
\gesperrt{continued-fraction periods}:
\[
[1, 3, 3, 1, 1, 1], \quad [10, 1, 4, 1].
\]
Here, then, we have two different periods, both of which are
reversible.


% ===== scan leaf 420, printed 395 =====
% !PRINTED 395
\weberpage{395}
\pdfbookmark[2]{§ 127. The Pell equation}{bm141}
\begin{center}
\S.\,127.

The Pell Equation.
\end{center}

If $\omega$ is a reduced number belonging to the discriminant $D$, then
by the continued fraction development one obtains
\weq{1}{\omega = \dfrac{P_n \omega_n + P_{n-1}}{Q_n \omega_n + Q_{n-1}},}
and $\omega_n = \omega$ holds always when $n$ is a multiple of the
number of terms $\nu$ of the period. But then it follows from (1):
\weq{2}{Q_n \omega^2 = (P_n - Q_{n-1}) \omega + P_{n-1}.}
Now let $\omega$ satisfy the quadratic equation \S.\,122, (5):
\weq{3}{c \omega^2 = a + b \omega, \quad D = b^2 + 4 a c,}
and from this, since $a$, $b$, $c$ are without common divisor,
equation (2) must follow by multiplication with an integer.
If we denote this integer by $u$, we thus have
\weq{4}{Q_n = u c, \quad P_{n-1} = u a, \quad P_n - Q_{n-1} = u b.}
If we further set
\weq{5}{P_n + Q_{n-1} = t,}
so that $t$ too is a positive integer, then from (4) and (5) it follows that
\[
P_n = \frac{t + u b}{2}, \quad Q_n = u c,
\]
\weq{6}{P_{n-1} = u a, \quad Q_{n-1} = \frac{t - u b}{2},}
and hence, having regard to the equation
\[
P_n Q_{n-1} - Q_n P_{n-1} = (-1)^n\text{:}
\]
\weq{7}{t^2 - D u^2 = (-1)^n 4.}
Equation (7) is called the \gesperrt{Pell equation}. The
problem is this: for a given $D$, to find all its integral
solutions $t$, $u$. Since, if $t$, $u$ satisfy the equation,
$\pm t$, $\pm u$ also satisfy it, we may restrict ourselves to the determination
of its positive solutions.

The formulae (4), (5) give us an infinite number of such
solutions; we have only to choose for $n$ an arbitrary multiple of the
period number $\nu$, and for $u$ the greatest common
divisor of $Q_n$, $P_{n-1}$, $P_n - Q_{n-1}$, or also simply the number
%CONT>>


% ===== scan leaf 421, printed 396 =====
% !PRINTED 396
\weberpage{396}
%<<CONT
$Q_n : c$, and to determine $t$ from (5), or also, after $u$ has
been determined, directly from (7).

If $\nu$ is odd and $n$ is an odd multiple of $\nu$, then
\weq{8}{t^2 - D u^2 = -4,}
but if $\nu$ is even or $n$ is an even multiple of $\nu$, then
\weq{9}{t^2 - D u^2 = 4.}
Equation (9) has the self-evident solution $t = 2$,
$u = 0$; in all other solutions of (8) or (9), $t$ and
$u$ are different from zero, and we consider here only the positive
solutions.

We shall now show that from an arbitrary
equation of the form (3), by the formulae (4), (5), one obtains \gesperrt{all
solutions} of (8) and (9).

So let us assume that $t$, $u$ is any positive solution
of equation (8) or (9), and that $\omega = \{a, b, c\}$ is a reduced irrational number belonging to the discriminant $D$. By the equations (6) we determine $P_n$, $Q_n$, $P_{n-1}$, $Q_{n-1}$, which evidently become
integers, since $t$ and $u b$ are either both even or both
odd. From equation (3), which $\omega$ satisfies, (2) then results,
from which in turn
\weq{10}{\omega = \dfrac{P_n \omega + P_{n-1}}{Q_n \omega + Q_{n-1}}}
is to be inferred; and from (8) or (9) it results that
\weq{11}{P_n Q_{n-1} - Q_n P_{n-1} = \mp 1,}
in which, as also in the following formulae, the upper or the
lower sign holds according as $t$, $u$ satisfies equation (8) or (9).
But now from (6) it follows that
\[
Q_n - Q_{n-1} = \frac{(2 c + b) u - t}{2},
\]
and because $\{a, b, c\}$ is a reduced number, we have $2 c + b > \sqrt{D}$
[\S.\,125, (4)], hence
\[
Q_n - Q_{n-1} > \frac{u \sqrt{D} - t}{2} = \frac{u^2 D - t^2}{2(u \sqrt{D} + t)} = \frac{\pm 2}{u \sqrt{D} + t};
\]
but since $u$ and $t$ are positive integers and $D > 1$, we have
$u \sqrt{D} + t > 2$, hence $Q_n - Q_{n-1} > 0$ for the upper and $> -1$
for the lower sign.

Now because $Q_n - Q_{n-1}$ must be an integer, it is
accordingly equal to or greater than zero, and we have
%CONT>>


% ===== scan leaf 422, printed 397 =====
% !PRINTED 397
\weberpage{397}
%<<CONT
\[
Q_{n-1} \leqq Q_n,
\]
where the equality sign is possible only in the case of the lower
sign, that is, in the case of equation (9).

Further (since $b < \sqrt{D}$) we have
\[
Q_{n-1} = \frac{t-ub}{2} > \frac{t-u\sqrt{D}}{2} = \frac{t^2-u^2D}{2(t+u\sqrt{D})} = \frac{\mp 2}{t+u\sqrt{D}},
\]
so $Q_{n-1}$ is zero or positive, and zero can occur only in the case
of the upper sign, that is, in the case of equation (8).
From this we conclude
\weq{12}{0 \leqq Q_{n-1} \leqq Q_n,}
where the equality sign can occur at the lower limit only in the
case of equation (8), that is, of the upper signs, and at the upper
limit only in the case of equation (9), that is, of the lower signs.

We now develop the rational fraction $P_n : Q_n$ into a finite
continued fraction
\[
\frac{P_n}{Q_n} = (\alpha_0, \alpha_1, \alpha_2 \ldots \alpha_{n-1}),
\]
by taking $n$, according to \S.\,115, I, so that
\[
(-1)^n = \mp 1
\]
If then $P' : Q'$ denotes the next-to-last convergent, we have
\[
P_n Q' - Q_n P' = \mp 1
\]
and, by \S.\,117, II:
\[
0 \leqq Q' \leqq Q_n,
\]
where the equality sign can occur at the lower limit only for $n=1$,
that is, in the case of the upper signs, and at the upper limit only
for $n=2$, that is, of the lower signs, just as in (12). From this
it follows by \S.\,118, III:
\[
Q' = Q_{n-1}, \quad P' = P_{n-1},
\]
and hence, by (10),
\[
(\alpha_0, \alpha_1, \alpha_2 \ldots \alpha_{n-1}, \omega) = \omega,
\]
that is, $[\alpha_0, \alpha_1, \alpha_2 \ldots \alpha_{n-1}]$ is a period of the continued fraction for
$\omega$, from which, by the formulae (4), (5), the solution $t, u$ of (7)
is derived. It is thereby proved that in this way we find
\gesperrt{all solutions} of equation (7).

From this some further important consequences can be drawn.


% ===== scan leaf 423, printed 398 =====
% !PRINTED 398
\weberpage{398}
From (7) it follows that when $u$ increases, $t$ must increase
as well, and hence that when $u$ has the smallest possible positive
value, $t$ is also as small as possible. One may therefore speak of a
smallest positive solution, which we shall denote by $T$, $U$. It is
obtained by taking $n$ in the formulae (4), (5) as small as possible,
that is, equal to the number of terms $\nu$ of the period. Hence if $\nu$
is even, only the equation (9), and not the equation (8), is
solvable; but if $\nu$ is odd, both (8) and (9) are solvable. Since
this can depend only on the discriminant $D$, and not on the particular
number $\omega$, it follows

\gesperrt{9. that for a given discriminant the various
periods contain either all an even or all
an odd number of terms.}

If $D$ is even, and hence divisible by 4, then $t$ must also be
even, and the equation (7) can be divided term by term by 4.
If $D$ is odd, then $t$, $u$ are likewise either both even
or both odd; if they are both odd, then their squares,
like all odd square numbers, are congruent to
1 modulo 8, and hence we must have $D \equiv 5\ (\mathrm{mod}\ 8)$. But it is not
always possible, when $D \equiv 5\ (\mathrm{mod}\ 8)$, to solve the equation (7)
in odd numbers\footnote{Cf. \gesperrt{Cayley}, Note sur l'équation $x^2 - D y^2 = \pm 4$, $D \equiv 5\ (\mathrm{mod}\ 8)$
Crelle's Journal, Bd. 53. Mathematical papers, Vol. IV. Tables for the
solutions of the Pell equation in the form $y^2 = a x^2 + 1$ have been
computed by \gesperrt{Degen} (Canon Pellianus, Havniae 1817). Such a table is also to be found in \gesperrt{Legendre}'s
number theory (German translation by \gesperrt{Maser}).}. If $D \equiv 1\ (\mathrm{mod}\ 8)$, then
$t$, $u$ must be even.

The examples of \S.\,126, treated in this way, yield
the following smallest positive solutions of the Pell
equation, where the factor 4 has been cancelled wherever possible.
\[
\begin{aligned}
5^2 - 29.1^2 &= -4\\
70^2 - 29.13^2 &= -1\\
170^2 - 19.39^2 &= 1\\
6^2 - 37.1^2 &= -1\\
35^2 - 34.6^2 &= 1.
\end{aligned}
\]


% ===== scan leaf 424, printed 399 =====
% !PRINTED 399
\weberpage{399}
\pdfbookmark[2]{§ 128. Derivation of all solutions of the Pell equation from the smallest positive one}{bm142}
\begin{center}
\S.\,128.

\gesperrt{Derivation of all solutions of the Pell equation
from the least positive one.}
\end{center}

If $t$, $u$ is any solution of the Pell equation, then
the expression
\[
\frac{t^2 - D u^2}{4},
\]
which has the value $\pm 1$, is the norm of the two conjugate numbers
\[
\frac{t + u\sqrt{D}}{2}, \quad \frac{t - u\sqrt{D}}{2}.
\]
We shall call them the \gesperrt{units} belonging to $\sqrt{D}$.
If, then, $(t_1, u_1)$, $(t_2, u_2)$ are any two (positive or negative)
solutions of the Pell equation
\[
t^2 - D u^2 = \pm 4,
\]
then the two numbers
\[
\Theta_1 = \frac{t_1 + u_1 \sqrt{D}}{2}, \quad \Theta_2 = \frac{t_2 + u_2 \sqrt{D}}{2}
\]
have the norm $\pm 1$, and the same therefore holds also of their product
\[
\Theta_3 = \Theta_1 \Theta_2 = \frac{t_3 + u_3 \sqrt{D}}{2}.
\]
Here
\[
t_3 = \frac{t_1 t_2 + u_1 u_2 D}{2}, \quad u_3 = \frac{t_1 u_2 + t_2 u_1}{2},
\]
and $t_3$ and $u_3$ are, as one sees at once, integers.

For if $D$ is even, then $t_1$, $t_2$ are even, and if $D$
is odd, then $t_1$, $u_1$ are either both even or both
odd, and likewise $t_2$, $u_2$. Thus $\Theta_3$ is likewise a unit
belonging to $\sqrt{D}$, and it follows \gesperrt{that the product of two}
(and consequently also of several) \gesperrt{units is again a unit.}

If $\Theta$ is an arbitrary unit, then $\pm \Theta^{-1}$ is the unit conjugate
to $\Theta$. Among the units different from $\pm 1$
there always belong four together,
\[
\pm \Theta, \quad \pm \Theta^{-1},
\]
among which \gesperrt{one} is positive and greater than 1. This one corresponds
%CONT>>


% ===== scan leaf 425, printed 400 =====
% !PRINTED 400
\weberpage{400}
%<<CONT
the positive values of $t$, $u$; a second is positive and
smaller than 1, and the two others are negative. If $T$, $U$
is the smallest positive solution of the Pell equation, then
\[
\Theta = \frac{T + U\sqrt{D}}{2}
\]
is the smallest improperly fractional positive unit belonging to $\sqrt{D}$,
and we can show that all units belonging to $\sqrt{D}$ are contained
in the form $\pm \Theta^n$, in which $n$ denotes an integer. For this
it suffices to show that the expression $\Theta^n$, for a positive $n$,
yields all the units that are greater than 1.
But this is very simple; for if $\Theta_1$ is any one of these
units, then it will lie between two successive
powers of $\Theta$, since these powers grow to infinity with the
exponent; thus
\[
\Theta^n \leqq \Theta_1 < \Theta^{n+1},
\]
hence
\[
1 \leqq \Theta_1 \Theta^{-n} < \Theta.
\]
But since $\Theta_1 \Theta^{-n}$ is likewise a unit, it follows, unless
the equality sign holds, that $\Theta_1 \Theta^{-n}$ is greater than 1 and smaller than
$\Theta$, which is contrary to the assumption that $\Theta$ is the smallest
improperly fractional unit; hence we must have
\[
\Theta_1 = \Theta^n
\]
which was to be proved.

\gesperrt{Remark concerning the \gesperrt{Gauss} theory of quadratic
forms}\footnote{Cf. Disq. ar. art. 183 ff., \gesperrt{Dirichlet}-\gesperrt{Dedekind}, Vorlesungen
über Zahlentheorie. Fourth edition. \S.\,72 ff.}.

In order to establish the connection of our considerations with
the otherwise well-known \gesperrt{Gauss} theory of quadratic
forms, we add a few remarks, which in our
context are not strictly required, and which may therefore
also be passed over.

Two quadratic forms,
\[
(A, B, C), \quad (-A, B, -C)
\]
in the notation of \gesperrt{Gauss}, are equivalent to one another in
the sense of \gesperrt{Gauss} then and only then, when the period
%CONT>>


% ===== scan leaf 426, printed 401 =====
% !PRINTED 401
\weberpage{401}
\pdfbookmark[2]{§ 129. Approximate computation of the real roots of a numerical equation by continued fractions}{bm143}
%<<CONT
of the continued fraction development consists of an odd number
of terms. Now a reduced number $\omega = \{a, b, c\}$ is the first
root of
\[
(a, \tfrac{1}{2}b, -c) \quad \text{or} \quad (2a, b, -2c)
\]
(according as $D$ is even or odd) and the second root of
\[
(-a, \tfrac{1}{2}b, c) \quad \text{or} \quad (-2a, b, 2c).
\]
If these two forms are equivalent, that is, if the periods
of $\omega$ consist of an odd number of terms, then to each
two-sided class of forms there corresponds one, and to each pair of
opposite classes two systems of equivalent numbers $\omega$; hence
the class number in the \gesperrt{Gauss}ian sense is just as great as the
number of periods of the reduced numbers $\omega$.

But when the number of terms of the period is even, the
forms $(A, B, C), (-A, B, -C)$ are not equivalent, and the
\gesperrt{Gauss}ian class number is twice as great as the number of
periods of the $\omega$.

\begin{center}
\S.\,129.

\gesperrt{Approximate calculation of the real roots of a
numerical equation by continued fractions.}
\end{center}

On the theory of continued fractions is founded a procedure
due to \gesperrt{Lagrange} for calculating the real roots
of a numerical equation with any desired degree of approximation,
a procedure which can now be set forth in a few words.

Let $f(x) = 0$ be a real equation which has at least one
real root. By one of the various methods of Sections
VIII and IX we can then always find an integer
$a_0$ such that one or more of the roots of $f(x)$ lie between
$a_0$ and $a_0 + 1$. If, then, we transform $f(x)$ by
the substitution
\[
x = a_0 + \frac{1}{x_1}
\]
into $f_1(x_1)$, then $f_1(x_1)$ will have one or more roots which are
greater than 1. We therefore determine a \gesperrt{positive integer}
$a_1$ such that one or more of the roots
of $f_1(x_1)$ lie between $a_1$ and $a_1 + 1$, and transform $f_1(x_1)$ by the
substitution
%CONT>>


% ===== scan leaf 427, printed 402 =====
% !PRINTED 402
\weberpage{402}
%<<CONT
\[
x_1 = a_1 + \frac{1}{x_2}
\]
in $f_2(x_2)$. The equation $f_2 = 0$ then again has at least
one root that is greater than 1, and which therefore lies between
the two positive integers $a_2$ and $a_2 + 1$, and so on.
Finally one will in this way come to isolate a single
root, and one obtains in the sequence of convergents
\[
(a_0), \quad (a_0, a_1), \quad (a_0, a_1, a_2), \quad (a_0, a_1, a_2, a_3) \ldots
\]
a sequence of rational numbers which lie alternately above and below
one of the roots $x$ and approach this root without limit.
If one has advanced in this way to the convergent
\[
(a_0, a_1, a_2 \ldots a_{\nu-1}) = \frac{P_\nu}{Q_\nu}
\]
then this differs from the true value
of $x$ by less than $1 : Q_\nu^2$. The size of the denominator thus gives
directly a measure of the accuracy already attained.

In the ordinary cases one will obtain the numbers $a_0, a_1, a_2 \ldots$
simply by attending to the changes of sign of the
functions $f, f_1, f_2 \ldots$

To test this procedure on an example, we take
the cubic equation which we have already considered
earlier:
\[
x^3 - 2x - 2 = 0.
\]
It has a root between 1 and 2, so we have to set $a_0 = 1$.

If one actually carries out the transformation indicated above,
one obtains the following chain of equations:
\[
\begin{aligned}
x^3 - 2x - 2 &= 0, & a_0 &= 1\\
3x^3 - x^2 - 3x - 1 &= 0, & a_1 &= 1\\
2x^3 - 4x^2 - 8x - 3 &= 0, & a_2 &= 3\\
9x^3 - 22x^2 - 14x - 2 &= 0, & a_4 &= 1\\ % [sic] printed order a_4 before a_3
46x^3 - 6x^2 - 32x - 9 &= 0, & a_3 &= 2\\
x^3 - 94x^2 - 132x - 46 &= 0, & a_5 &= 95\\
3561x^3 - 9183x^2 - 191x - 1 &= 0, & a_6 &= 2,
\end{aligned}
\]
hence
\[
x = (1, 1, 3, 2, 1, 95, 2 \ldots)
\]
and the convergents:
%CONT>>


% ===== scan leaf 428, printed 403 =====
% !PRINTED 403
\weberpage{403}
\pdfbookmark[2]{§ 130. Rational roots of equations with integer coefficients. Reducible equations}{bm144}
%<<CONT
\[
\frac{1}{1}, \frac{2}{1}, \frac{7}{4}, \frac{16}{9}, \frac{23}{13}, \frac{2201}{1244}, \frac{4425}{2501} \cdots
\]
If we convert the two last fractions into decimal fractions, we obtain for $x$ the two limits:
\[
1{,}7692926; \quad 1{,}76929228,
\]
and the error in the latter, a number slightly too small, is less than
\[
0{,}00000016.
\]
That one here obtains a good result after few steps rests on the circumstance that a larger partial denominator 95 occurs fairly early. In most cases one obtains small numbers for the partial denominators; the denominators of the convergents then grow slowly, and one only laboriously obtains a result of greater accuracy.

The finding of larger partial denominators, e.g. 95, is made easier by the remark that, when a root of an equation has a large value, the negative coefficient of the second-highest power of the unknown gives an approximation to this root, even if only a rough one, as e.g. in the next to last of the above equations 94 to 95.

\begin{center}
\S.\,130.\\[0.3em]
\gesperrt{Rational Roots of Equations with Integral Coefficients.}\\
\gesperrt{Reducible Equations.}
\end{center}

We close this section with a consideration of equations with integral coefficients which, even though it does not stand in immediate relation to the theory of continued fractions, yet most fittingly finds a place here.

When one has to do with equations whose coefficients are rational numbers, one will first look for possible rational roots. If the equation
\[
f(x) = a_0 x^n + a_1 x^{n-1} + a_2 x^{n-2} + \cdots + a_{n-1} x + a_n = 0
\]
has rational coefficients, we can always assume that these coefficients are integers; to reduce the case of fractional coefficients to this one, one need only multiply the whole equation by a common multiple of all the denominators.


% ===== scan leaf 429, printed 404 =====
% !PRINTED 404
\weberpage{404}
We can, however, further assume that $a_0 = 1$; for if in the product $a_0^{n-1} f(x)$ we set
\[
a_0 x = x_1,
\]
the equation becomes
\[
x_1^n + a_1 x_1^{n-1} + a_0 a_2 x_1^{n-2} + \cdots + a_0^{n-1} a_n = 0.
\]
We shall therefore now assume that in the equation
\weq{1}{x^n + a_1 x^{n-1} + a_2 x^{n-2} + \cdots + a_{n-1} x + a_n = 0}
the coefficients $a_1, a_2 \ldots a_n$ are integers. A \gesperrt{rational} root of (1) cannot be a fractional number; for if we assume that (1) is satisfied by the fraction
\[
x = \frac{p}{q},
\]
where $p$ and $q$ are integers without a common divisor, and $q$ is positive and greater than 1, then multiplication by $q^{n-1}$ gives
\[
\frac{p^n}{q} + a_1 p^{n-1} + a_2 q p^{n-2} + \cdots + a_n q^{n-1} = 0.
\]
Thus $p^n : q$ would have to be an integer, which is impossible.

Now if the equation (1) can be satisfied by putting an integer for $x$, then this integer, as the equation (1) shows, must necessarily be a divisor of $a_n$, and one has therefore only to substitute the various divisors of $a_n$, taken with positive and negative sign, into the equation by way of trial.

To the question of rational roots of an equation with integral coefficients there also comes back the general question whether an integral rational function
\[
f(x) = x^n + a_1 x^{n-1} + a_2 x^{n-2} + \cdots + a_{n-1} x + a_n,
\]
whose coefficients $a_1, a_2 \ldots a_n$ are integers, can be divisible by a rational function of lower degree $\nu$:
\[
\varphi(x) = x^\nu + \alpha_1 x^{\nu-1} + \alpha_2 x^{\nu-2} + \cdots + \alpha_{\nu-1} x + \alpha_\nu
\]
with rational coefficients $\alpha_1, \alpha_2 \ldots \alpha_\nu$. If $f(x)$ is divisible by $\varphi(x)$, then the quotient $\psi(x)$ is again an integral rational function with rational coefficients of degree $n - \nu$:
\[
\psi(x) = x^{n-\nu} + \beta_1 x^{n-\nu-1} + \cdots + \beta_{n-\nu-1} x + \beta_{n-\nu},
\]
and since $f(x)$ is also divisible by $\psi(x)$, we can, in
%CONT>>


% ===== scan leaf 430, printed 405 =====
% !PRINTED 405
\weberpage{405}
the decision of our question assume $\nu\leqq\frac{1}{2}n$. Moreover it
follows from \S.\,2 that the coefficients $\alpha$ and $\beta$ of $\varphi$ and $\psi$
must be \gesperrt{integers}.

Now in order to decide on the possibility of such a division, one
takes the function $\varphi(x)$ with indeterminate
coefficients $\alpha$, and divides $f(x)$ by $\varphi(x)$ (according to \S.\,3). There
then results a remainder of degree $\nu-1$, and if one sets its
coefficients, which are all rational functions of the $\alpha$,
equal to zero, one obtains $\nu$ equations which these
coefficients must satisfy. By elimination (\S.\,50) one can
produce an equation which contains only one of these unknowns,
and which must then have an integral root. It
will often be more expedient, however, not actually to carry out this
elimination,. but to attempt directly to satisfy the system of the % [sic] stray comma-period as printed
$\nu$ equations by integral values of the $\alpha$. The
selection of the numbers to be tried can be greatly restricted by many
an artifice resting on theorems of number theory.

As an example we shall take an equation of the tenth degree
from the theory of elliptic functions, where
such problems very frequently occur in such a way that the
decomposability into factors of a definite degree is theoretically established,
and where it is then a matter of finding these
factors $^1$). Such an equation is
\[
x^{10} - 36\,x^8 + 528\,x^6 - 3897\,x^4 + 14354\,x^2 - 21025 = 0.
\]
We ask whether the left side is decomposable into two factors of the
fifth degree. If one of the two factors is
\[
x^5 + \alpha x^4 + \beta x^3 + \gamma x^2 + \delta x + \varepsilon
\]
then, since no odd powers occur in the given equation,
the other is
\[
x^5 - \alpha x^4 + \beta x^3 - \gamma x^2 + \delta x - \varepsilon,
\]
and hence we must have:
\[
\begin{aligned}
x^{10} - 36\,x^8 + 528\,x^6 - 3897\,x^4 + 14354\,x^2 - 21025 \\
= (x^5 + \beta x^3 + \delta x)^2 - (\alpha x^4 + \gamma x^2 + \varepsilon)^2.
\end{aligned}
\]

\noindent\rule{4cm}{0.4pt}

$^1$) Cf.\ the author's work: ``Elliptische Functionen und algebraische
Zahlen'', Braunschweig 1891. The example treated here, from the theory
of the transformation of degree 47, is taken from the paper: ``Ein
Beitrag zur Transformationstheorie der elliptischen Functionen etc.'' by
\gesperrt{H.\ Weber} (1893). Mathematische Annalen, Bd.\ 43.


% ===== scan leaf 431, printed 406 =====
% !PRINTED 406
\weberpage{406}
If one sets corresponding coefficients equal to one another,
it follows that
\[
\begin{aligned}
1.&\quad \alpha^2 - 2\beta = 36,\\
2.&\quad \beta^2 + 2\delta - 2\alpha\gamma = 528,\\
3.&\quad \gamma^2 + 2\alpha\varepsilon - 2\beta\delta = 3897,\\
4.&\quad \delta^2 - 2\gamma\varepsilon = 14354,\\
5.&\quad \varepsilon^2 = 21025,
\end{aligned}
\]
and it is now a matter of seeking integral values of $\alpha, \beta, \gamma, \delta, \varepsilon$
that satisfy these five equations. First of all, from 5. one obtains
\[
\varepsilon = \sqrt{21025} = 145 = 5\,.\,29;
\]
and one may take $\varepsilon$ positive, because the term independent of $x$
is in any case positive in one of the two factors. Equation 4.
yields, if one takes on the right the residue with respect to the modulus
145,
\weq{4}{\delta^2 \equiv -1 \pmod{145},}
and this congruence has the four solutions
\[
\delta \equiv \pm 12,\ \pm 17 \pmod{145}.
\]
Since $\delta$ must moreover be even, by 4., $\delta$ could have the following
values:
\[
\delta = \pm 12,\ \pm 128,\ \pm 162.
\]
For the time being one will go no further, since large numerical values for
$\delta$ are improbable from the outset. Corresponding to these values of $\delta$
one obtains from 4. for $\gamma$:
\[
\gamma = -49,\ 7,\ 41.
\]
If we apply in addition the modulus 5, it follows that:
\[
\gamma^2 \equiv 1,\ -1,\ 1 \pmod{5}
\]
and from 3.:
\[
\beta \equiv \pm 1,\ \pm 2,\ \pm 1 \pmod{5}
\]
and from 1.:
\[
\alpha^2 \equiv 1 \pm 2,\ 1 \mp 1,\ 1 \pm 2 \pmod{5}.
\]
From this it follows, since only 0, 1 and 4, not 2 and 3, can be
congruent to a square with respect to the modulus 5, that
the signs in $\delta$ must be taken thus:
\[
\begin{aligned}
\delta &= -12,\ +128,\ -162\\
\alpha &\equiv \pm 2,\ 0,\ \pm 2 \pmod{5}.
\end{aligned}
\]
Equation 2. shows further that the middle case is to be excluded, since the right-hand side is not divisible by 5.


% ===== scan leaf 432, printed 407 =====
% !PRINTED 407
\weberpage{407}
That $\delta = -12$ must also be excluded follows from equation 2., which for $\delta = -12$, $\gamma = -49$ has as a consequence the congruence
\[
\beta^2 \equiv 6 \pmod{7}
\]
which, however, is not solvable.

There thus remains only:
\[
\delta = -162,\quad \gamma = 41.
\]
From 3. there follows for $\beta$ the congruence:
\[
17\,\beta \equiv 93 \pmod{145}
\]
or
\[
\beta \equiv -1 \pmod{5},\quad 2\,\beta \equiv -1 \pmod{29},
\]
hence
\[
\beta \equiv 14 \pmod{145}.
\]

If one takes $\beta = 14$, then $\alpha = -8$ follows, and all the equations 1. to 5. are satisfied. With this, then, the decomposition has been obtained, which is easily verified afterwards.
\[
\begin{aligned}
x^{10} - 36\,x^8 + 528\,x^6 - 3897\,x^4 + 14354\,x^2 - 21025 = \\
(x^5 - 8\,x^4 + 14\,x^3 + 41\,x^2 - 162\,x + 145)\\
(x^5 + 8\,x^4 + 14\,x^3 - 41\,x^2 - 162\,x - 145).
\end{aligned}
\]


% ===== scan leaf 433, printed 408 =====
\weberpage{408}% page number not printed on this page
\pdfbookmark[1]{Twelfth Section. Theory of the Roots of Unity.}{bm145}
\pdfbookmark[2]{§ 131. The roots of unity}{bm146}
% !PRINTED none
% Section-opening page: no running page number visible in the scan; part title, section
% heading, and body text follow.

\begin{center}
\gesperrt{Twelfth Section.}

\vspace{1em}
\gesperrt{Theory of the Roots of Unity.}

\vspace{0.5em}
\noindent\rule{4cm}{0.4pt}

\vspace{1em}
\S.\,131.

\vspace{0.5em}
\gesperrt{The Roots of Unity.}
\end{center}

\vspace{1em}
By a root of unity one understands in general a real or
imaginary number of such a nature that some power of it with
an integral exponent is equal to 1. We have already become
acquainted with these roots of unity in \S.\,33 and represented
them by trigonometric and exponential functions. We have,
moreover, made use of special cases, e.\,g. in \S.\,35 the third
roots of unity, which were there expressed without trigonometric
functions.

In all deeper investigations of algebraic equations a more
precise knowledge of the roots of unity and of their properties
is now indispensable. In this section we shall therefore occupy
ourselves more thoroughly with the elementary part of the
algebraic theory of the roots of unity, without making use of
their representation by trigonometric functions. On account of
the geometric application to the construction of the regular
polygons, to which we have already pointed in \S.\,33, the theory
of the roots of unity is also called \gesperrt{the theory of the division
of the circle}. It is in essence a creation of \gesperrt{Gauss} $^1$).

\noindent\rule{4cm}{0.4pt}

$^1$) \gesperrt{Gauss}, Disq. ar. Sectio VII.


% ===== scan leaf 434, printed 409 =====
% !PRINTED 409
\weberpage{409}
When the $n$th power of a number $r$ is equal to 1, that is, when the equation

\weq{1}{r^n = 1}
is satisfied for a positive integer $n$, then $r$ is called an $n$th
\gesperrt{root of unity}, or a \gesperrt{root of unity of degree $n$}. Thus all the $n$th
roots of unity, and only these, are roots of the equation of degree $n$

\weq{2}{f(x) = x^n - 1 = 0.}
We have

\weq{3}{f'(x) = n\,x^{n-1},}
and consequently $f(x)$ has no divisor in common with $f'(x)$; hence $f(x)$ has no
multiple roots, and there are $n$, and no more, distinct
$n$th roots of unity.

If $r$ is an $n$th root of unity, then so is every integral power of $r$;
for from $r^n = 1$ it follows that $r^{kn} = 1$, if $k$ is any positive or negative
integer ($k = 0$ not being excluded); hence $r^k$ is also an $n$th
root of unity. But since there are only $n$ roots of unity of degree $n$, the
powers $r^k$ are not all distinct from one another. On this point the following now holds:

If two numbers $k, k'$ differ by a multiple of $n$, that is, if

\weq{4}{k' \equiv k \pmod{n}}
then also

\weq{5}{r^k = r^{k'}.}
For if $k' = k + h\,n$, then

\[
r^{k'} = r^k r^{h n},
\]
from which, since $r^n = 1$, equation (5) follows.

Thus in the sequence of numbers

\weq{6}{1, r, r^2, r^3 \ldots r^{n-1}}
all the $r^k$ that are distinct from one another are certainly contained; but conversely
the quantities (6) need not all be distinct from one another.

Let us assume that $k$ and $k' = k + \mu$ are two numbers of the sequence

\[
0, 1, 2 \ldots n-1
\]
and

\[
r^k = r^{k+\mu},
\]
%CONT>>


% ===== scan leaf 435, printed 410 =====
% !PRINTED 410
\weberpage{410}
%<<CONT
it follows that $r^\mu = 1$ must hold. Thus no term that has already occurred can recur in
the sequence (6) before the first term 1 occurs for the second time, and if $\mu$ is the
least positive number for which $r^\mu = 1$, then the numbers:

\weq{7}{1, r, r^2 \ldots r^{\mu - 1}}
are all different from one another.

\gesperrt{Then $\mu$ must be a divisor of $n$.} For by division the integers $h, \mu'$ can be
determined so that

\[
n = h\mu + \mu'; \quad 0 \leqq \mu' < \mu.
\]
But then, as follows from

\[
r^n = r^{h\mu} r^{\mu'}
\]
we also have $r^{\mu'} = 1$, that is, since $\mu$ is to be the least positive number for
which $r^\mu = 1$, $\mu' = 0$, or $n$ is divisible by $\mu$.

\gesperrt{Thus $r$ is at the same time a $\mu$th root of unity, but not a root of unity of any
still lower degree.}

\gesperrt{The number $r$ is called a primitive $n$th root of unity if it is not at the
same time a root of unity of a lower degree.}

From this definition it follows that the numbers of the sequence (6), if $r$ is a
primitive $n$th root of unity, are all different from one another, and that all the
$n$th roots of unity are contained among them, but that they make up only a part of the
$n$th roots of unity if $r$ is an imprimitive $n$th root of unity.

Every root of unity whose degree is a divisor of $n$ different from $n$ is at the same
time an imprimitive $n$th root of unity.

\gesperrt{If $r$ is at the same time an $n$th and an $m$th root of unity, then it is also a
$\mu$th root of unity, if $\mu$ is the greatest common divisor of $n$ and $m$.}

For by \S.\,118 the integers $x, y$ can be determined so that

\[
m x + n y = \mu
\]
and consequently

\[
r^\mu = r^{mx} r^{ny} = 1.
\]
To this corresponds the other theorem:

\gesperrt{If $r_1, r_2 \ldots$ are roots of unity of the degrees $n_1, n_2 \ldots$, then they are
all at the same time roots of unity of the degree $m$, if $m$ denotes any common
multiple of $n_1, n_2 \ldots$.}


% ===== scan leaf 436, printed 411 =====
% !PRINTED 411
\weberpage{411}
\pdfbookmark[2]{§ 132. Primitive roots of unity}{bm147}
\begin{center}
§. 132.

\gesperrt{Primitive Roots of Unity.}
\end{center}

That primitive roots of unity exist for every degree $n$ we have not yet proved in the
preceding section. We must first make good this omission, and in doing so we shall also
determine the exact number of the primitive roots of unity.

Let the degree $n$ be decomposed into two factors $a, b$ which are relatively prime to one
another, so that

\[
n = ab,
\]
and let $\alpha$ be an $a$th and $\beta$ a $b$th root of unity. Then the product

\weq{1}{r = \alpha\beta}
is an $n$th root of unity. If $\alpha', \beta'$ are two other $a$th and $b$th roots of
unity, then $r' = \alpha'\beta'$ is also an $n$th root of unity, and it is to be shown
that $r'$ is different from $r$ unless at the same time $\alpha = \alpha'$ and
$\beta = \beta'$.

For since $a, b$ are relatively prime, one can by \S.\,118 determine the integers $x, y$
so that

\weq{2}{ax + by = 1}
holds; and then it follows from (1) that

\[
\alpha = r^{by}, \quad \beta = r^{ax}.
\]
Accordingly $\alpha$ and $\beta$ are completely determined by $r$, and if $\alpha'\beta'$
is also to be equal to $r$, then we must have $\alpha = \alpha'$, $\beta = \beta'$.

If, therefore, we let $\alpha$ in (1) run through all the $a$th, and $\beta$ through all
the $b$th roots of unity, then $r$ takes exactly $ab = n$ different values, and it
follows:

I.~\gesperrt{that in the form $\alpha\beta$ all the $n$th roots of unity can be
represented,}

and further:

II.~\gesperrt{that $r$ is a primitive $n$th root of unity if and only if $\alpha$ is a
primitive $a$th and $\beta$ a primitive $b$th root of unity.}

For, first, let $\mu$ be the smallest positive exponent for which $r^\mu = 1$; then also
$\alpha^\mu \beta^\mu = 1$, and from this it follows by (2), if we raise both sides to the
powers $by$ and $ax$,

\[
\alpha^\mu = 1, \quad \beta^\mu = 1.
\]


% ===== scan leaf 437, printed 412 =====
% !PRINTED 412
\weberpage{412}
Now if $\mu < n$, it cannot be divisible at once by $a$ and by $b$,
and hence $a$ and $b$ cannot both be the least positive exponents of the
powers of $\alpha$ and $\beta$ that become equal to 1; that is to say, if $r$ is not
a primitive $n$th root of unity, then $\alpha$ and $\beta$ are not at the same time
primitive $a$th and $b$th roots of unity; or, if
$\alpha$ and $\beta$ are primitive $a$th and $b$th roots of unity,
then their product $r$ is a primitive $n$th root of unity. On the other
hand it is clear that if $\alpha$ or $\beta$ is at the same time a root of unity of lower
degree than $a$ or $b$, then $r$ too will be a root of unity of degree lower than
the $n$th.

If $n$ decomposes into several factors $a, b, c \ldots$, of which every two are
relatively prime to one another, and if $\alpha, \beta, \gamma \ldots$ are roots of unity of the degrees
$a, b, c \ldots$, and one sets

\weq{3}{r = \alpha\beta\gamma \ldots,}
then one concludes, by repeated application of the preceding theorems, that in (3) all
the $n$th roots of unity, and each only once, are contained, and further that
$r$ is a \gesperrt{primitive} $n$th root of unity then and only then, when
$\alpha, \beta, \gamma \ldots$ are primitive roots of unity of the degrees $a, b, c \ldots$.

If we now denote the number of primitive $n$th roots of unity by
$\varphi(n)$, it follows from what has here been proved that

\weq{4}{\varphi(n) = \varphi(a)\varphi(b)\varphi(c) \ldots,}
if

\[
n = abc \ldots
\]
and $a, b, c \ldots$ are numbers which, two by two, are relatively prime to one another.

Now every number $n$ can be decomposed in one and only one way into a product of
prime powers

\[
n = p^\pi p_1^{\pi_1} p_2^{\pi_2} \ldots,
\]
in which $p, p_1, p_2 \ldots$ are distinct primes and $\pi, \pi_1, \pi_2 \ldots$
are positive exponents, so that from the formula (4) it follows:

\weq{5}{\varphi(n) = \varphi(p^\pi)\, \varphi(p_1^{\pi_1})\, \varphi(p_2^{\pi_2}) \ldots,}
and that it therefore only remains to decide whether and how many primitive
roots of unity of the degree $p^\pi$ exist.

But this question is very simple to decide. For if $\varrho$ is a
root of unity of the degree $p^\pi$, then the lowest degree to which $\varrho$ belongs as a
root of unity is a divisor of $p^\pi$, hence a power of $p$, and if it is
therefore not equal to
%CONT>>


% ===== scan leaf 438, printed 413 =====
% !PRINTED 413
\weberpage{413}
%<<CONT
$p^\pi$, a divisor of $p^{\pi-1}$; that is, every non-primitive root of unity of degree $p^\pi$ is at the same time a root of unity of degree $p^{\pi-1}$. But since there are $p^\pi$ roots of unity of degree $p^\pi$ and only $p^{\pi-1}$ roots of unity of degree $p^{\pi-1}$, there must be
\[
p^\pi - p^{\pi-1} = p^\pi \left(1 - \frac{1}{p}\right)
\]
primitive roots of unity of degree $p^\pi$. From this, by (5), one obtains the number of all primitive $n$th roots of unity
\weq{6}{\varphi(n) = n\,\Pi\left(1 - \frac{1}{p}\right),}
in which the product sign $\Pi$ extends over all the mutually distinct primes dividing $n$. Only in the case $n = 1$ does the formula (6) no longer apply; in this case $\varphi(1) = 1$ is to be set. The number $\varphi(n)$ is therefore never equal to zero. Since, by this, for every degree $n$ there is at least one primitive root of unity $r$, all $n$th roots of unity can be represented by the powers of $r$:
\weq{7}{1, r, r^2, r^3 \ldots r^{n-1}.}
If $r^k$ is any power of $r$, then $r^{km} = 1$ will hold if and only if $km$ is divisible by $n$. If, then, $n = n'n''$ and $n'$ is the greatest common divisor of $k$ and $n$, then $m$ must be divisible by $n''$, and $n''$ is the exponent of the lowest power of $r^k$ which becomes equal to 1; that is, $r^k$ is a primitive $n''$th root of unity. From this follows the theorem:

\gesperrt{III. If $r$ is a primitive $n$th root of unity, then $r^k$ is a primitive $n$th root of unity if and only if $k$ is relatively prime to $n$.}

If we take $k$ from the sequence of numbers $1, 2, 3 \ldots n$, there \gesperrt{results the theorem of number theory that $\varphi(n)$ is equal to the number of the numbers which are not greater than $n$ and relatively prime to $n$.} This is the original definition of the symbol $\varphi(n)$ in general use in number theory.

If $k$ is relatively prime to $n$, one can determine an integer $x$ so that $kx \equiv 1\ (\mathrm{mod}\ n)$. If then $r, r_1$ are two $n$th roots of unity, $r^k = r_1^k$ can hold only when $r = r_1$, as follows by raising to the power $x$. From this it follows, by III.:
%CONT>>


% ===== scan leaf 439, printed 414 =====
% !PRINTED 414
\weberpage{414}
\pdfbookmark[2]{§ 133. Equations for the primitive nth roots of unity}{bm148}
%<<CONT
\gesperrt{IV. If $k$ is relatively prime to $n$ and $r$ runs through the sequence of the primitive $n$th roots of unity, then $r^k$ runs through the same sequence of numbers, though in a different order.}

Finally we cite the further theorem: --

\gesperrt{V. If $n$ is a prime, then every $n$th root of unity, with the exception of 1, is a primitive $n$th root of unity.}

\begin{center}
\S.\,133.\\[0.3em]
\gesperrt{Equations for the primitive roots of unity of degree $n$.}
\end{center}

All $n$th roots of unity are, as we have seen, roots of an equation $f_n(x) = 0$, where
\weq{1}{f_n(x) = x^n - 1}
is understood. If now we free the function $f_n(x)$ from all factors which it has in common with other functions of the same form $f_{n_1}(x)$, which is done by rational operations (\S.\,6), we obtain an equation $X_n = 0$, which is satisfied by the primitive $n$th roots of unity and by these only, and $X_n$ has the form
% [mirrored 2026-08-07] canon ruled the v/nu question for this page: Greek \nu
\weq{2}{X_n = x^\nu + a_1 x^{\nu-1} + \cdots + a_\nu.}
The degree $\nu$ is equal to $\varphi(n)$, and the coefficients $a_1, a_2 \ldots a_\nu$ are \gesperrt{rational numbers}.

In seeking the common factors of $f_n$ and $f_{n_1}$ we may, for $n_1$, restrict ourselves to the divisors of $n$. How the function $X_n$ may be formed in a simple way we shall presently see more closely. But first we prove a general theorem concerning these functions.

The $n$th roots of unity comprise all primitive $\mu$th roots of unity, where $\mu$ is any divisor of $n$, $n$ itself and 1 included, since it is precisely the unity 1 itself that is to be regarded as the primitive first root of unity. If, then, we let $\mu$ run through all divisors of $n$, and observe that two different $X_\mu$ never have a common divisor, and that neither $f_n(x)$ nor $X_\mu$ contains multiple factors, it follows that
\weq{3}{f_n(x) = \Pi X_\mu,}
%CONT>>


% ===== scan leaf 440, printed 415 =====
% !PRINTED 415
\weberpage{415}
in which the product sign $\Pi$ extends over all divisors $\mu$ of $n$.

From this we obtain, incidentally, a proof of the number-theoretic
theorem
\weq{4}{n = \overset{\mu}{\Sigma}\varphi(\mu),}
if we set the degrees of the functions on the right and the
left side of (3) equal to one another.

On the other hand we conclude, by \gesperrt{Gauss}'s theorem
(\S.\,2), that the coefficients of all the $X_\mu$ are \gesperrt{integers}.
For if fractional numbers occurred among them, the
product could not consist solely of integral coefficients, as by (3) and (1) it nevertheless must.

All the functions $f_n(x)$ have the divisor $x-1$, and if
we carry out the division, we obtain
\weq{5}{\frac{x^n-1}{x-1} = x^{n-1} + x^{n-2} + \cdots + x + 1.}
From this we conclude that, if $r$ is any primitive
or non-primitive $n$th root of unity, with the sole exception of $r=1$, always
\weq{6}{1 + r + r^2 + \cdots + r^{n-1} = 0,}
while for $r=1$ the sum on the left side of (6)
obviously has the value $n$.

If $n$ is a prime, then apart from 1 there are no imprimitive $n$th roots of unity, and therefore, when $n$ is a prime, which we shall indicate by writing $p$ for it,
we have
\weq{7}{X_p = x^{p-1} + x^{p-2} + \cdots + x + 1.}
Just as simply, $X_n$ can be formed when $n$ is a \gesperrt{prime power}, that is,
\[
n = p^{\pi},
\]
or, if for brevity we set $p^{\pi-1} = p'$:
\[
n = p'p.
\]
The $n$th roots of unity consist in this case of the
primitive $n$th roots of unity and of the $p'$th roots of unity,
and thus:
\weq{8}{X_n = \frac{x^{pp'}-1}{x^{p'}-1} = x^{p'(p-1)} + x^{p'(p-2)} + \cdots + x^{p'} + 1}


% ===== scan leaf 441, printed 416 =====
% !PRINTED 416
\weberpage{416}
If $p' > 1$, and hence $\pi > 1$, then the term with the
$(\nu-1)$th power of the unknown, whose coefficient is equal to the negative sum of the roots, is missing from this equation. We therefore
have the theorem:

\gesperrt{The sum of all primitive $n$th roots of unity,
when $n$ is a higher power of a prime,
is always equal to zero.}

For the general formation of $X_n$ we shall apply a recurrent
procedure; we assume $X_n$ as already formed,
denote by $p$ a prime which does not divide $n$, by $p'$,
as above, the $(\pi-1)$th power of $p$, and now form $X_{np'p}$,
in which of course $p'$ may also be equal to 1.

If we denote by $r$ the \gesperrt{primitive} $n$th roots of unity,
by $\alpha$ every root of unity of degree $pp'$, and by $\alpha'$ every root of unity of degree $p'$, then one obtains all the primitive roots of unity of degree $npp'$ by taking away from all the
$r\alpha$ the $r\alpha'$ (\S.\,132, II). But the $r\alpha$ are
the roots of the equation:
\[
X_n(x^{pp'}) = 0,
\]
because the quantities
\[
(r\alpha)^{pp'} = r^{pp'},
\]
agree, apart from their order, with the $r$ themselves
(\S.\,132, IV). Likewise the $r\alpha'$ are the roots of the equation
\[
X_n(x^{p'}) = 0;
\]
and from this it follows that
\weq{9}{X_{npp'}(x) = \frac{X_n(x^{pp'})}{X_n(x^{p'})}.}
The value $n=1$ forms no exception to this,
if we understand by $X_1$ the function $x-1$. According to this
we can form $X_n$ comparatively simply in all cases.
We consider some special cases.

Let us suppose that $n$ contains only two distinct primes $p$, $q$, so that
\[
n = pp'qq',
\]
then the formulae (8) and (9) yield:
\[
X_n = \frac{(x^n-1)\left(x^{\frac{n}{pq}}-1\right)}{\left(x^{\frac{n}{p}}-1\right)\left(x^{\frac{n}{q}}-1\right)},
\]
%CONT>>


% ===== scan leaf 442, printed 417 =====
% !PRINTED 417
\weberpage{417}
\pdfbookmark[2]{§ 134. Irreducibility}{bm149}
%<<CONT
a formula which can easily be generalised by complete induction as follows.

If one denotes by $\mu_1$ all the numbers that arise from $n$ when one divides
$n$ by an even number of distinct prime divisors of $n$ ($n$ itself included),
and by $\mu_2$ the numbers that arise from $n$ when one divides $n$ by an odd
number of such prime divisors, then

\weq{10}{X_n = \dfrac{\prod\limits_{\mu_1}(x^{\mu_1} - 1)}{\prod\limits_{\mu_2}(x^{\mu_2} - 1)}.}

The proof results from (9) for $npp'$, if one assumes that the correctness has
already been proved for $n$.

If $r$ denotes each of the primitive $n$th roots of unity, then for \gesperrt{odd} $n$
the $-r$ are the primitive $2n$th roots of unity (\S.\,132, I, II).
Accordingly, for odd $n$,

\weq{11}{X_{2n}(x) = X_n(-x).}

If $n$ is a power of $2$, then by (8):

\weq{12}{X_n = \dfrac{x^n - 1}{x^{\frac{n}{2}} - 1} = x^{\frac{n}{2}} + 1.}

If we set $x = 1$ in formula (8), then $X_n$ takes the value $p$ if $n$ is a
power of $p$. On the other hand, formula (9) yields, when $n > 1$, the value
$1$ for $X_{npp'}$ (for $n = 1$ the numerator and denominator on the
right-hand side would be $= 0$).

We thus obtain the theorem:

\gesperrt{V. The function $X_n$ takes for $x = 1$ the value $p$ if $n$ is a
power of the prime $p$, and the value $1$ if $n$ has more than one prime as
divisor.}

\begin{center}
\S.\,134.

\gesperrt{Irreducibility.}
\end{center}

The functions $X_n$, which are defined in the preceding section, have the
property of \gesperrt{irreducibility}, which is very important for their
deeper study and all applications, i.e.

\gesperrt{It is not possible to decompose a function $X_n$ of $x$ into factors of lower
degree which themselves}
%CONT>>


% ===== scan leaf 443, printed 418 =====
% !PRINTED 418
\weberpage{418}
%<<CONT
\gesperrt{integral rational functions of $x$ with rational numerical coefficients.}

This important theorem was set up and first proved by \gesperrt{Gauss} for
the case that $n$ is a prime. Later many other proofs have been given, some
of them more general and simpler. We follow here a proof by
\gesperrt{Kronecker}.

In the function $X_n$, as we have seen above, the highest power of $x$ has
the coefficient $1$, while the remaining coefficients are integers. If we
now assume that $X_n$ can be decomposed into two factors with rational
coefficients

\weq{1}{X_n = \varphi(x)\,\psi(x)}

then the product of the coefficients of the highest powers of $x$ in
$\varphi$ and $\psi$ must be equal to $1$, and we may therefore assume that
these highest coefficients in $\varphi$ and $\psi$ themselves have the value
$1$ (since otherwise we would divide by the product of the two
coefficients). But then it follows from the theorem of \gesperrt{Gauss}
(\S.\,2) that the remaining coefficients of $\varphi(x)$ and $\psi(x)$ must
be \gesperrt{integers}, and hence that $\varphi(x)$ and $\psi(x)$, when an integer is
put for $x$, must themselves pass into integers.

If we therefore set $x = 1$, it follows that:
\[
\varphi(1)\,\psi(1) = X_n(1),
\]
and so by Theorem V. this is equal to $1$ or to a prime $p$. Accordingly,
of the two factors $\varphi(1)$, $\psi(1)$, one must certainly have the
value $\pm 1$, while the other is equal to $\pm 1$ or $\pm p$. We therefore
assume that

\weq{2}{\varphi(1) = \pm 1.}

The roots $r$ of the equation $X_n = 0$ now distribute themselves between
the two equations $\varphi = 0$, $\psi = 0$, and since neither of the two
factors $\varphi$, $\psi$ is independent of $x$, there must certainly be
one among the primitive $n$th roots of unity $r$, which we shall denote by
$\rho$, for which
\[
\varphi(\rho) = 0
\]
holds.

Now if $r$ denotes any primitive $n$th root of unity, then among the powers
\[
r,\ r^2,\ r^3\ \dots\ r^{n-1}
\]
all $n$th roots of unity with the exception of $1$, and therefore also
$\rho$, are contained; and it follows therefore that the product
%CONT>>


% ===== scan leaf 444, printed 419 =====
% !PRINTED 419
\weberpage{419}
%<<CONT
\weq{3}{\varPhi(x) = \varphi(x)\,\varphi(x^2)\,\varphi(x^3)\dots\varphi(x^{n-1})}

vanishes for $x = r$. But since $r$ may be any root of $X_n$, and
these roots are all different from one another, $\varPhi(x)$ must be
divisible by $X_n$, hence

\weq{4}{\varPhi(x) = X_n\,\varPsi(x).}

The function $\varPhi(x)$, as appears from its mode of formation, has the
coefficient $1$ for the highest power of $x$ and otherwise integer
coefficients; likewise $X_n$. Hence $\varPsi(x)$ must also have the
coefficient $1$ for the highest power of $x$ and otherwise integer
coefficients. But it follows from (2) and (3) that $\varPhi(x)$ takes
the value $\pm 1$ for $x = 1$, hence

\weq{5}{\pm 1 = X_n(1)\,\varPsi(1),}

wherein $\varPsi(1)$ must be an integer.

Now if $n$ is a prime power $p^{\pi}$, then $X_n(1) = p$, and this
contradicts equation (5). Thus the irreducibility of $X_n$ is proved
for the case in which $n$ is a prime or a power of a prime.

The general proof therefore still requires the following: \gesperrt{If $a$ and
$b$ are two integers without a common divisor, and if it is further assumed
that $X_a$ and $X_b$ are irreducible, it follows that $X_{ab}$
is also irreducible.}

To prove this, we set $n = ab$ and assume that $\varphi(x)$ is
a rational factor of $X_n$. The roots of $\varphi(x) = 0$ are then
contained among the primitive $n$th roots of unity, and are
therefore of the form

\weq{6}{\varrho = \alpha\beta,}

wherein $\alpha$ and $\beta$ denote certain primitive $a$th and
$b$th roots of unity (\S.\,132).

Now (by a special case of the Tschirnhausen transformation)
a function $\varPhi(x)$ with rational coefficients can be derived
from $\varphi(x)$, whose roots are the $a$th powers of the roots of
$\varphi(x)$. The roots of $\varPhi(x)$ are therefore
\[
\varrho^a = \beta^a,
\]
that is, since $a$ is relatively prime to $b$, all primitive $b$th
roots of unity; and if therefore any $\varrho$ exists at all,
$\varPhi(x)$ must have a non-constant factor in common with $X_b$. But now
we have assumed, as already proved, that $X_b$ has no rational
factor other than itself; hence $\varPhi(x)$ must be divisible by
%CONT>>


% ===== scan leaf 445, printed 420 =====
% !PRINTED 420
\weberpage{420}
%<<CONT
$X_b$, and the $\beta^a$, and consequently also the $\beta$
themselves, which occur in (6), \gesperrt{together comprise all the primitive $b$th
roots of unity.} In the same way it is now proved that $\alpha$ in (6) runs
through all the primitive $a$th roots of unity, and that therefore $\rho$
runs through all the primitive $n$th roots of unity.

The roots of $\varphi(x)$ thus agree completely with the totality of the
roots of $X_n$, and consequently, if for the determination of a numerical
factor we further stipulate, say, that the highest power of $x$ in
$\varphi(x)$ shall have the coefficient $1$, $\varphi(x)$ must \gesperrt{be
identical} with $X_n$.

\gesperrt{Under the assumptions made, therefore, $X_n$ is not decomposable into
factors of lower degree with rational coefficients.}

By the same means we can prove the following theorem of
\gesperrt{Kronecker}, which expresses the irreducibility of the functions $X_n$
in a wider sense\footnote{Mém. sur les facteurs
irréductibles de l'expression $(x^n - 1)$. Liouville's Journal, Vol. XIV.}.

\gesperrt{The function $X_n$ cannot be decomposed into factors which are
integral and rational with respect to $x$, and which contain in their
coefficients only rational numbers and roots of unity whose degree is
relatively prime to $n$.}

Suppose that $a$ is relatively prime to $n$, and that $r$ is a primitive
$n$th, $\alpha$ a primitive $a$th root of unity,
and let $\varphi(x,\alpha)$ be a divisor of $X_n$, so that at least for
one $r$

\weq{7}{\varphi(r,\alpha) = 0.}

Now let $\rho$ be a primitive $na$th root of unity, hence
a root of $X_{na} = 0$; then both $r$ and $\alpha$ are powers
of $\rho$ (\S.\,131), say
\[
r = \rho^{\xi},\quad \alpha = \rho^{\eta},
\]
where $\xi$ is relatively prime to $n$ but divisible by $a$, and $\eta$
relatively prime to $a$ but divisible by $n$. Hence

\weq{8}{\varphi(\rho^{\xi}, \rho^{\eta}) = 0.}

Now, as we have just proved, $X_{na}$ is irreducible, and consequently
the function $\varphi(x^{\xi}, x^{\eta})$, which has a divisor in common
with $X_{na}$, must be divisible by this function; that is, for
every integer $h$ that is relatively prime to $na$:

\weq{9}{\varphi(\rho^{h\xi}, \rho^{h\eta}) = 0.}


% ===== scan leaf 446, printed 421 =====
% !PRINTED 421
\weberpage{421}
\pdfbookmark[2]{§ 135. The discriminant of the cyclotomic equation}{bm150}

Now if $s$ is an arbitrary number coprime to $n$, we can determine $h$ so that
\[
h \equiv s \ (\mathrm{mod}\ n), \quad h \equiv 1 \ (\mathrm{mod}\ a)
\]
and accordingly it follows from (9) that
\[
\varphi(r^s, \alpha) = 0,
\]
that is, equation (7) is satisfied for \gesperrt{all} roots of $X_n$. Hence $\varphi(x,\alpha)$ is divisible by $X_n$, and since on the other hand $X_n$ was assumed to be divisible by $\varphi(x,\alpha)$, it must be identical with $X_n$ up to a numerical factor.

\begin{center}
§. 135.

\gesperrt{The Discriminant of the Cyclotomic Equation.}
\end{center}

The roots of unity of degree $n$ can be represented in transcendental form by
\weq{1}{r = e^{\frac{2\pi i k}{n}} = \cos\frac{2\pi k}{n} + i \sin\frac{2\pi k}{n}.}
This root of unity is primitive or not primitive according as $k$ is coprime to $n$ or not. The simplest among the primitive ones is obtained for $k = 1$, namely
\weq{2}{r_0 = e^{\frac{2\pi i}{n}} = \cos\frac{2\pi}{n} + i \sin\frac{2\pi}{n} = a + bi.}
The angle $2\pi : n$ is the $n$th part of the whole circumference of the circle. The various points of division, which one obtains by laying off this angle on the circumference from an arbitrary initial point, determine the regular $n$-gon inscribed in the circle.

% Registered Bd. I scan, PDF page 447 / printed p. 421, rendered at 600 dpi.
% Exact source crop x=100,y=3150,w=810,h=550; SHA-256 27F4BCA9D93A7A0A2DED6313E3840AD29B999168A364F0FDE053D0091B6B361B.
\weberfigureimage{weber-b1-fig28.png}{0.28}

The points of division can be constructed if one knows $r$, and with it $a$ and $b$ (or even only one of these two quantities). In particular, for the side $s$ of the regular $n$-gon, if the radius $c$ of the circumscribed circle is taken equal to 1, we obtain
\[
s = 2 \sin\frac{\pi}{n} = \sqrt{2\left(1 - \cos\frac{2\pi}{n}\right)}.
\]
The remaining quantities $r$ stand in the same relation to the other points of division. On account of this geometric
%CONT>>


% ===== scan leaf 447, printed 422 =====
% !PRINTED 422
\weberpage{422}

%<<CONT
meaning we call the equation $X_n = 0$, whose roots are the quantities (1), the \gesperrt{cyclotomic equation}.

We now determine the discriminant of the cyclotomic equation $X_n = 0$, but restrict ourselves to the case in which $n$ is an \gesperrt{odd prime}.

We then have
\weq{3}{X_n = x^{n-1} + x^{n-2} + \cdots + x + 1,}
and since only the one imprimitive $n$th root of unity 1 exists, we have
\weq{4}{f_n(x) = x^n - 1 = (x-1) X_n,}
hence
\weq{5}{f_n'(x) = n x^{n-1} = X_n + (x-1) X_n'.}
Now if $D$ is the discriminant of $X_n$, then by \S.\,46
\weq{6}{D = (-1)^{\frac{n-1}{2}} \Pi X_n'(r),}
where the product sign $\Pi$ extends over all roots $r$ of the equation $X_n = 0$, and where, since $n$ is odd, the equality
\[
(-1)^{\frac{n(n-1)}{2}} = (-1)^{\frac{n-1}{2}}
\]
could be used. But by (5) we have
\weq{7}{n r^{n-1} = (r-1) X_n'(r),}
and since the product of all the $r$ is equal to 1 (because it is equal to the term of $X_n$ independent of $x$), and since
\[
\Pi(x-r) = \Pi(r-x) = X_n,
\]
hence
\[
\Pi(r-1) = X_n(1) = n, \qquad (\text{\S.\,133})
\]
it follows from (6) and (7) that
\weq{8}{D = (-1)^{\frac{n-1}{2}} n^{n-2}.}
The roots of $X_n$, if $r$ is one of them, are all contained in the form
\[
r, r^2, r^3 \ldots r^{n-1},
\]
and if we therefore form the difference product
\begin{equation}
\tag*{9)}
\begin{gathered}
P = (r - r^2)(r - r^3) \cdots (r - r^{n-1}) \\
(r^2 - r^3) \cdots (r^2 - r^{n-1}) \\
\cdot\ \cdot\ \cdot\ \cdot\ \cdot\ \cdot\ \cdot\ \cdot\ \cdot\ \cdot \\
(r^{n-2} - r^{n-1})
\end{gathered}
\end{equation} % [sic] source prints the equation tag as $9)$, without an opening parenthesis
%CONT>>


% ===== scan leaf 448, printed 423 =====
% !PRINTED 423
\weberpage{423}

%<<CONT
introduce, then (\S.\,46)
\[
D = P^2.
\]
By (8) it is therefore also
\weq{10}{P = n^{\frac{n-3}{2}} \sqrt{(-1)^{\frac{n-1}{2}} n},}
whereby $P$ is determined, apart from the sign. $P$ is real when $n \equiv 1$, and imaginary when $n \equiv -1 \ (\mathrm{mod}\ 4)$, hence real for $n = 5, 13, 17, 29, 37 \ldots$, imaginary for $n = 3, 7, 11, 19, 23, 31 \ldots$ The sign that we have to give the root in (10) depends on which $r$ we have chosen in the expression (9).

If we choose a definite $r$, e.g. $r = r_0$, then we can still determine the sign in (10). To carry this out, we divide the binomial factors of $P$,
\[
r^\nu - r^\mu, \quad \nu < \mu,
\]
whose number amounts to $\frac{n(n-1)}{2}$, into two classes. The differences of the one class form the product
\weq{11}{Q = (r - r^{n-1})(r^2 - r^{n-2}) \cdots \left(r^{\frac{n-1}{2}} - r^{\frac{n+1}{2}}\right),}
which thus contains all those factors in which $\mu + \nu = n$.

The remaining factors can be gathered into pairs of the following form
\weq{12}{R = (r^\nu - r^\mu)(r^{n-\mu} - r^{n-\nu}).}
The number of factors in $Q$ is $\frac{n-1}{2}$, and hence the number of pairs $R$ is
\[
\frac{1}{2}\left(\frac{n(n-1)}{2} - \frac{n-1}{2}\right) = \left(\frac{n-1}{2}\right)^2.
\]
Now if
\[
r = e^{\frac{2\pi i k}{n}}
\]
then
\[
R = -2 + r^{\mu-\nu} + r^{\nu-\mu} = -2\left(1 - \cos\frac{2\pi(\mu-\nu)}{n}\,k\right),
\]
that is, $R$ is always negative. Consequently the product of all the $R$ has the sign
\[
(-1)^{\frac{n-1}{2}}.
\]
Further,
\weq{13}{Q = (2i)^{\frac{n-1}{2}} \sin\frac{2\pi k}{n} \sin\frac{4\pi k}{n} \cdots \sin\frac{(n-1)\pi k}{n},}
%CONT>>


% ===== scan leaf 449, printed 424 =====
% !PRINTED 424
\weberpage{424}

%<<CONT
and the sign of this depends on $k$. But if we take $k = 1$, and hence $r = r_0$, then all the angles
\[
\frac{2\pi}{n}, \frac{4\pi}{n} \ldots \frac{(n-1)\pi}{n}
\]
lie between zero and $\pi$, and the sines are all positive. We conclude from this that in this case
\weq{14}{P = (-i)^{\frac{n-1}{2}} n^{\frac{n-2}{2}} \sqrt{n}}
and that $\sqrt{n}$ is to be taken positive.

The sign of the product of the $R$ does not depend on the choice of $r$, but only on the nature of $n$. It is therefore of interest to determine by itself the product $Q$, whose sign does depend on the choice of $r$. We obtain it in the following way:

If we multiply the expression (11) by
\[
r \cdot r^2 \cdot r^3 \cdots r^{\frac{n-1}{2}} = r^{\frac{n^2-1}{8}}
\]
and then by
\[
r^{-1} \cdot r^{-2} \cdot r^{-3} \cdots r^{-\frac{n-1}{2}} = r^{-\frac{n^2-1}{8}},
\]
we obtain
\[
r^{\frac{n^2-1}{8}} Q = (r^2 - 1)(r^4 - 1) \cdots (r^{n-1} - 1)
\]
\[
r^{-\frac{n^2-1}{8}} Q = (1 - r^{n-2})(1 - r^{n-4}) \cdots (1 - r).
\]
Now since the exponents $2, 4 \ldots n-1, n-2, n-4 \ldots 1$ together comprise all the numbers $1, 2 \ldots n-1$, multiplication of these two expressions yields:
\[
Q^2 = (-1)^{\frac{n-1}{2}} \Pi(r-1) = (-1)^{\frac{n-1}{2}} n,
\]
hence
\weq{15}{Q = \pm i^{\frac{n-1}{2}} \sqrt{n}.}
Comparison with (13) yields
\weq{16}{2^{\frac{n-1}{2}} \sin\frac{2\pi k}{n} \sin\frac{4\pi k}{n} \cdots \sin\frac{(n-1)\pi k}{n} = \pm \sqrt{n}.}
The sign in this formula, as in (15), still depends on $k$; but it is the positive one when $k = 1$. For the rest, we shall derive one more important theorem concerning this sign, which will be investigated more closely later on.


% ===== scan leaf 450, printed 425 =====
% !PRINTED 425
\weberpage{425}
\pdfbookmark[2]{§ 136. Primitive congruence roots}{bm151}

By the definition (11), $Q$ is an integral rational function
of $r$ which, when one orders it according to powers of $r$, has
integral rational numerical coefficients. If we therefore set
\[
r = r_0^k,
\]
then
\[
Q = F(r_0^k),
\]
where $F$ is the symbol for a rational function whose
coefficients do not depend on $k$. If we now set for
$k$ in turn the values
\[
k = 1,\ 2 \ldots n-1,
\]
then $r_0^k$ represents all the roots of the cyclotomic equation $X_n = 0$,
and the sum
\[
\overset{k}{\Sigma} Q = F(r_0) + F(r_0^2) + \cdots + F(r_0^{n-1})
\]
is a \gesperrt{symmetric function} of the roots of this equation; it
can therefore be expressed rationally in terms of the coefficients of the equation,
that is, in terms of rational numbers, and is consequently itself a
\gesperrt{rational number}. On the other hand, by the formula (15),
\[
\overset{k}{\Sigma} Q = h\, i^{\frac{n-1}{2}} \sqrt{n},
\]
where $h$ denotes an integer, namely the number of cases
in which the positive sign is to be taken in (15), diminished by
the number of cases in which the negative sign holds. But the two
are compatible with one another only when $h = 0$; and
with this the following theorem is proved:

\gesperrt{1. If $k$ runs through the sequence of numbers 1, 2, 3, \dots\ $n-1$,
then in the formula (16) the positive sign holds just as often as the
negative.}

It hardly needs special mention that one may also take for $k$
another complete system of residues of $n$, excluding the number
divisible by $n$.

\begin{center}
\S.\,136.\\[2pt]
\gesperrt{Primitive congruence roots.}
\end{center}

Parallel to the theory of the roots of unity runs a theory of the
so-called binomial congruences, without which, in the theory of the
cyclotomic equations, further steps cannot
%CONT>>


% ===== scan leaf 451, printed 426 =====
% !PRINTED 426
\weberpage{426}

%<<CONT
can be made, the fundamentals of which may fittingly be inserted
here, where the analogy with the theory of the roots of unity
shows itself clearly.

So let $n$ now be a \gesperrt{prime} and
\weq{1}{f(x) = a_0 x^m + a_1 x^{m-1} + \cdots + a_{m-1} x + a_m}
an integral function of $x$, whose coefficients $a$ are integers,
and $a_0$ not divisible by $n$. If we put for $x$ an integer
$\alpha$ such that $f(\alpha)$ becomes divisible by $n$, then we
say (by analogy with equations) that $\alpha$ is a root of the
congruence of degree $m$:
\weq{2}{f(x) \equiv 0 \quad (\mathrm{mod}\ n).}
If $\alpha$ is increased by a multiple of $n$, it remains a root
of the congruence (2). Such roots, congruent with respect to the
modulus $n$, are not counted as distinct. Under this assumption
we can state the theorem:

\gesperrt{1. A congruence of degree $m$ for a prime
modulus cannot have more than $m$ distinct roots.}

The theorem is true for $m = 1$, for the congruence $a_0 x + a
\equiv 0\ (\mathrm{mod}\ n)$ has only one root, namely, if
% [sic] The source prints modulus m here although the inverse is being taken modulo n; preserved diplomatically.
$a_0'$ is determined so that $a_0 a_0' \equiv 1\ (\mathrm{mod}\ m)$
holds, $\alpha \equiv -a a_0'\ (\mathrm{mod}\ n)$. We therefore
now assume our theorem to be proved for the degree $m-1$ and
deduce its validity for the degree $m$ from this. By \S.\,4 we
can, if at first $\alpha$ is arbitrary, set
\weq{3}{\frac{f(x) - f(\alpha)}{x - \alpha} = f_1(x),}
in which $f_1(x)$ is an integral function of degree $(m-1)$,
which, if $\alpha$ is an integer, has integral
coefficients. If therefore $f(\alpha) \equiv 0$, it follows
from (3) that
\weq{4}{f(x) \equiv (x - \alpha)\, f_1(x) \quad (\mathrm{mod}\ n).}
Every root $\beta$ of the congruence (2) must therefore satisfy
the condition $(\beta - \alpha)\, f_1(\beta) \equiv 0\ (\mathrm{mod}\ n)$.
Hence either $\beta - \alpha$ or $f_1(\beta)$ is
divisible by $n$.

Now by hypothesis there are at most $m-1$ values $\beta$
for which $f_1(\beta)$ becomes divisible by $n$; if therefore
there is still an $m$th root of (4), it must be equal to $\alpha$.

If therefore a congruence of the form $f(x) \equiv 0\ (\mathrm{mod}\
n)$ has more roots than its degree amounts to, then we conclude
that
%CONT>>


% ===== scan leaf 452, printed 427 =====
% !PRINTED 427
\weberpage{427}
%<<CONT
the congruence is identical, i.e. that all coefficients of $f(x)$
must be divisible by $n$.

This theorem is, as one sees, entirely analogous to the algebraic theorem that an equation cannot have more roots
than its degree indicates. But the other theorem, that every
equation also actually has that many roots, cannot be carried
over. A congruence of degree $m$ can have fewer roots, indeed
none at all. All the more remarkable is a particular
congruence in which the number of roots always equals the
degree, on the basis of a theorem which is called \gesperrt{Fermat's theorem}, and which we formulate here as follows.

\gesperrt{2. The congruences}
\weq{5}{x^n - x \equiv 0, \quad x^{n-1} - 1 \equiv 0 \quad (\mathrm{mod}\ n)}
\gesperrt{have, if $n$ is a prime, as many roots
as their degree amounts to, namely $n$ and $n-1$.}

The two assertions are not essentially different, for
the first of the congruences (5) has all the roots of the second and
in addition the root 0, which does not satisfy the second. Likewise
the second has all the roots of the first, with the exception of the
root 0.

Now since for the modulus $n$ there are altogether only $n$ different
numbers, it is thus to be proved that for every integer
$\alpha$ the congruence holds
\weq{6}{\alpha^n \equiv \alpha \quad (\mathrm{mod}\ n).}
This congruence is correct for $\alpha = 0$ and $\alpha = 1$. We
therefore prove it again by complete induction, deriving the congruence
\weq{7}{(\alpha + 1)^n \equiv \alpha + 1 \quad (\mathrm{mod}\ n)}
from congruence (6), assumed to be correct. This follows from the binomial theorem,
if one observes that all binomial coefficients, with the exception of the first and the last, which are equal to 1, namely
\[
n, \quad \frac{n(n-1)}{1\,.\,2}, \quad \frac{n(n-1)(n-2)}{1\,.\,2\,.\,3}, \dots
\]
are integers divisible by $n$, because the numerator, but not
the denominator, is divisible by $n$. Accordingly
\[
(\alpha + 1)^n \equiv \alpha^n + 1 \quad (\mathrm{mod}\ n),
\]
hence the formula (7) follows from the formula (6).


% ===== scan leaf 453, printed 428 =====
% !PRINTED 428
\weberpage{428}
The difference
\[
x(x-1)(x-2)\dots(x-n+1) - x^n + x
\]
is a function of degree at most $n-1$. But it is congruent
to zero for $n$ values of $x$, namely for $x = 0, 1, 2, \dots n-1$, and hence we have identically
\[
x(x-1)(x-2)\dots(x-n+1) \equiv x^n - x \quad (\mathrm{mod}\ n).
\]
From this there results \gesperrt{Wilson's theorem}:
\[
1\,.\,2\,.\,3\dots(n-1) \equiv -1 \quad (\mathrm{mod}\ n),
\]
if one compares the coefficients of the first power of $x$ on both
sides.

We now restrict ourselves to the second of the congruences (5)
\weq{8}{x^{n-1} \equiv 1 \quad (\mathrm{mod}\ n)}
and first prove the theorem:

\gesperrt{3. If $a$ is a divisor of $n-1$, then the congruence}
\weq{9}{x^a \equiv 1 \quad (\mathrm{mod}\ n)}
\gesperrt{has exactly $a$ different roots.}

For, if $n-1 = ab$, we have
\[
x^{n-1} - 1 = (x^a - 1)(x^{a(b-1)} + x^{a(b-2)} + \cdots + x^a + 1),
\]
and since every root of the left-hand side must be a root of
one or the other factor on the right-hand side, it follows that
if $x^a - 1 \equiv 0$ had fewer than $a$ roots, the second factor
would have to have more roots than its degree indicates,
contrary to Theorem 1.

If $\alpha$ is a number not divisible by $n$, and $a$ is the least
positive number for which $\alpha^a \equiv 1 \pmod n$, and if further $h$ is any
positive integer for which $\alpha^h \equiv 1 \pmod n$, then $a$ is
necessarily a divisor of $h$, and so in particular $a$ is a divisor
of $n - 1$.

For if we set $h = qa + a'$, where $q$ is an integer and
$a' < a$, then also $\alpha^{a'} \equiv 1$; hence we must have $a' = 0$, because $a$ was to be the
least positive number for which this congruence is satisfied.

\gesperrt{A root $\alpha$ of the congruence (9) with the
property that no power lower than the $a$th
becomes congruent to unity is called a
primitive root of this congruence, or a primitive $a$th congruence root of the prime $n$, and a
primitive root $g$ of the congruence (8) we
also call, for short, a primitive root of the prime $n$.}


% ===== scan leaf 454, printed 429 =====
% !PRINTED 429
\weberpage{429}

We wish to prove that primitive roots exist for every prime $n$, and to
determine their number; and in doing so we follow the same path as with
the roots of unity.

Let $a$ and $b$ be two divisors of $n-1$ which are relatively prime
to each other, and let $\alpha$ be an $a$th, $\beta$ a $b$th
primitive congruence root of $n$. If we set $ab=c$, then
$\gamma=\alpha\beta$ is a primitive $c$th congruence root, and
conversely every primitive $c$th congruence root can also be represented
in the form $\alpha\beta$.

For, first: $(\alpha\beta)^c = \alpha^c\beta^c \equiv 1 \pmod{n}$;
and secondly: if $(\alpha\beta)^h = \alpha^h\beta^h \equiv 1 \pmod{n}$,
then also $\alpha^{hb}\beta^{hb}\equiv 1$, hence $\alpha^{hb}\equiv 1$,
hence $hb$ is divisible by $a$, hence $h$ too is divisible by $a$; and in
the same way one concludes that $h$ must be divisible by $b$, and hence
also by $ab=c$; that is, $\alpha\beta$ is a primitive $c$th
congruence root.

If conversely $\gamma$ is a primitive $c$th congruence root, then we
determine the positive integers $x, y$ so that
\weq{10}{a y + b x \equiv 1 \pmod{c}}
holds, by first determining $x$ from $bx \equiv 1 \pmod{a}$ and then $y$ from
$y \equiv \dfrac{1-bx}{a} \pmod{b}$. Then $y$ is relatively prime
to $b$ and $x$ relatively prime to $a$. Accordingly
\[
\gamma \equiv \gamma^{bx}\gamma^{ay} = \alpha\beta \pmod{n}
\]
and $\gamma^{bx} = \alpha$ is a primitive $a$th
congruence root, $\gamma^{ay} = \beta$ a primitive $b$th
congruence root (because $\alpha^h = \gamma^{bxh}$ can be $\equiv 1$
only when $h$ is divisible by $a$).

It remains to be shown that the products $\alpha\beta$ are all distinct
from one another, that is, that the congruence
\weq{11}{\alpha\beta \equiv \alpha'\beta' \pmod{n},}
when $\alpha, \alpha'$; $\beta, \beta'$ are primitive $a$th and
$b$th congruence roots, can be satisfied only when
\[
\alpha \equiv \alpha', \quad \beta \equiv \beta' \pmod{n}
\]
holds. But this follows by (10), if one raises (11) to the power
\[
bx \equiv 1 \pmod{a}
\]
from which $\alpha \equiv \alpha'$ results, and then also
$\beta \equiv \beta'$.

If, then, we denote the number of primitive $a$th
congruence roots by $\varphi(a)$, it follows from what has been proved that
\weq{12}{\varphi(c) = \varphi(a)\varphi(b).}


% ===== scan leaf 455, printed 430 =====
% !PRINTED 430
\weberpage{430}

It still remains, when $p$ is a prime and $pp_1$ is a power of $p$
dividing $n-1$, to determine $\varphi(pp_1)$.

If $\alpha$ is a non-primitive congruence root of degree $p^{\pi}$ and
$p^{\pi-1} = p_1$, then, if $\alpha^h$ is the lowest power of $\alpha$
congruent to 1 with respect to the modulus $n$, $h$ is a divisor of $pp_1$,
i.\,e. a power of $p$; but since $h$ is to be smaller than $pp_1$, it
is also a divisor of $p_1$, and consequently $\alpha^{p_1} \equiv 1
\pmod{n}$; thus: a non-primitive congruence root of degree $pp_1$
is at the same time a congruence root of degree $p_1$. Conversely, every
congruence root of degree $p_1$ is at the same time an imprimitive
congruence root of degree $pp_1$. Now since the two congruences
$x^{pp_1} \equiv 1$, $x^{p_1} \equiv 1$ have as many roots as their
degree amounts to, there remain $pp_1 - p_1$ \gesperrt{primitive congruence roots}
of the first, and thus we have
\weq{13}{\varphi(pp_1) = p_1(p-1).}

The function $\varphi(a)$ therefore has the same meaning as in \S.\,132,
and with this the number of primitive roots of the prime
$n$ is at the same time found to be equal to $\varphi(n-1)$. We therefore state the further theorem:

4.\ \gesperrt{A prime $n$ has $\varphi(n-1)$ primitive roots.}

With this the existence of primitive roots for every prime is
proved, and at the same time their number is determined. For finding the
primitive roots we have, to be sure, no other general method
than trial, which can be somewhat shortened by a few
artifices.

If $g$ is a primitive root of the prime $n$, then the residues of the
powers
\weq{14}{1, g, g^2, g^3 \ldots g^{n-2}}
are all different from one another, since, if $g^m \equiv g^{m+\mu}$ held,
$g^{\mu} \equiv 1$ would have to hold, which is not possible so long as $\mu$
is smaller than $n-1$. Hence among the residues of the sequence (14),
among which the residue 0 does not occur, each of the numbers
\weq{15}{1, 2, 3 \ldots n-1}
must be contained, and each only once, or in other words:

5.\ \gesperrt{If $a$ is a number not divisible by $n$, there is
one and only one number $\alpha$ from the sequence of the numbers
$0, 1, 2 \ldots n-2$, by which the congruence}
\weq{16}{g^{\alpha} \equiv a \pmod{n}}
%CONT>>


% ===== scan leaf 456, printed 431 =====
% !PRINTED 431
\weberpage{431}
%<<CONT
\gesperrt{is satisfied. But the same congruence is also satisfied if a number
congruent to $\alpha$ modulo $n-1$ is put as exponent, and in every
complete system of residues one finds only one such number
$\alpha$.}

The exponent $\alpha$ is called the \gesperrt{index} of $a$ with respect to the
\gesperrt{base} $g$, and one writes
\weq{17}{\alpha \equiv \operatorname{ind} a \pmod{n-1}.}

One thus has, for every number not divisible by $n$,
\weq{18}{g^{\operatorname{ind} a} \equiv a \pmod{n}.}

From the two formulae
\weq{19}{g^{\operatorname{ind} a + \operatorname{ind} b} \equiv ab, \quad g^{m \operatorname{ind} a} \equiv a^m \pmod{n}}
there follows the theorem:

6.\ \gesperrt{The index of a product is equal to the sum of the indices
of the factors; the index of the $m$th power of $a$ is equal
to $m$ times the index of $a$.}

These theorems are quite analogous to the corresponding theorems that
relate to computation with logarithms, and the analogy can be pursued
still further. Thus, for example, the transition from one base $g$ to
another base $g'$ is expressed by the relation
\weq{20}{\overset{g}{\operatorname{ind}}\, a \equiv \overset{g}{\operatorname{ind}}\, g' \; \overset{g'}{\operatorname{ind}}\, a \pmod{n-1}}
between the two systems of indices. We shall further adduce the easily proved theorem that among the
numbers $a$ not divisible by $n$, \gesperrt{those and only those} are
primitive roots whose indices are relatively prime to $n-1$.

For practical computations one conveniently makes use of so-called
index tables, which correspond to the tables of logarithms.

In the congruence (17) one calls $\alpha$ the index and $a$ the numerus,
and it is best to set up two tables, of which the one gives for each
index the numerus, the other for each numerus the index, where the
second is obtained from the first by rearrangement. Here it is to be
observed that for the indices the modulus $n-1$, and for the numeri the
modulus $n$, comes into consideration, and congruent numbers count
as equivalent. Thus one obtains, e.g., for $n = 13$ and the base 2
%CONT>>


% ===== scan leaf 457, printed 432 =====
% !PRINTED 432
\weberpage{432}
\pdfbookmark[2]{§ 137. Multiplication and division of the trigonometric functions}{bm152}
%<<CONT
\[
\begin{array}{c|cccccccccccc}
I & 0 & 1 & 2 & 3 & 4 & 5 & 6 & 7 & 8 & 9 & 10 & 11 \\ \hline
N & 1 & 2 & 4 & 8 & 3 & 6 & 12 & 11 & 9 & 5 & 10 & 7
\end{array}
\]

\[
\begin{array}{c|cccccccccccc}
N & 1 & 2 & 3 & 4 & 5 & 6 & 7 & 8 & 9 & 10 & 11 & 12 \\ \hline
I & 0 & 1 & 4 & 2 & 9 & 5 & 11 & 3 & 8 & 10 & 7 & 6
\end{array}.
\]

In the Canon Arithmeticus of \gesperrt{Jacobi} (1839) index tables
are compiled for all primes in the first thousand. On a smaller
scale, after computations of \gesperrt{Ostrogradsky}, in the ``Theory of
Congruences'' of \gesperrt{Tschebyscheff} (German by
\gesperrt{Schapira}, Berlin 1889).

It is still to be remarked that the index of 1 is always equal to 0,
and that the index of $-1$ (or of $n-1$) is always equal to
$\dfrac{n-1}{2}$. For since $(-1)^2$ has the index 0 or $n-1$, $-1$
must have the index $\dfrac{n-1}{2}$, since it does not have the
index 0, and twice its index is equal to 0 or $n-1$; in formulae:
\weq{21}{\operatorname{ind} 1 \equiv 0, \quad \operatorname{ind}(-1) \equiv \frac{n-1}{2} \pmod{n-1}.}

\begin{center}
\S.\,137.

\gesperrt{Multiplication and division of the trigonometric functions.}
\end{center}

We have already pointed out the connection that exists between the
theory of the roots of unity and the geometric problem of the
division of the circumference of the circle. The generalisation of
this problem leads to the division of an arbitrary angle. This
problem too leads to algebraic equations, of which we have already
used the one for trisection earlier, for the general solution of the
cubic equations (\S.\,112).

We now begin with the setting up of the general formulae.

By the theorem of \gesperrt{Moivre} (\S.\,33), if $n$ is an
arbitrary positive integer and $\varphi$ an arbitrary angle, we have
%CONT>>


% ===== scan leaf 458, printed 433 =====
% !PRINTED 433
\weberpage{433}
\weq{1}{\cos n\varphi + i\sin n\varphi = (\cos\varphi + i\sin\varphi)^n,}
and if we apply the binomial theorem on the right-hand side, then by separation of the real from the imaginary it follows that
\[
\begin{aligned}
\cos n\varphi &= \cos^n\varphi - B_2^{(n)}\cos^{n-2}\varphi\sin^2\varphi\\
&\quad + B_4^{(n)}\cos^{n-4}\varphi\sin^4\varphi - \cdots
\end{aligned}
\]
\weq{2}{
\begin{aligned}
\frac{\sin n\varphi}{\sin\varphi} &= n\cos^{n-1}\varphi - B_3^{(n)}\cos^{n-3}\varphi\sin^2\varphi\\
&\quad + B_5^{(n)}\cos^{n-5}\varphi\sin^4\varphi - \cdots
\end{aligned}
}
wherein $B_\nu^{(n)}$ are, as before, the binomial coefficients, and the sums on the right-hand side are to be continued as far as no negative exponents of $\cos\varphi$ occur. Since now $\sin^2\varphi = 1 - \cos^2\varphi$, these two expressions can be represented rationally in terms of $\cos\varphi$; we set
\weq{3}{2\cos\varphi = x,}
and introduce the notation
\weq{4}{2\cos n\varphi = A_n(x), \qquad \frac{\sin n\varphi}{\sin\varphi} = B_n(x),}
wherein, then, $A_n(x)$ and $B_n(x)$ are integral rational functions of $x$, $A_n(x)$ of degree $n$, $B_n(x)$ of degree $n-1$. We wish to determine the general law of these functions. We set them up for the first values $n = 1, 2, 3\ldots$, and then derive a recursion formula for the computation of the higher functions. One finds
\weq{5}{
\begin{aligned}
A_1(x) &= x & B_1(x) &= 1\\
A_2(x) &= x^2 - 2, & B_2(x) &= x\\
A_3(x) &= x^3 - 3x, & B_3(x) &= x^2 - 1.
\end{aligned}
}
But now, by well-known formulae,
\[
\begin{aligned}
\cos(n+1)\varphi + \cos(n-1)\varphi &= 2\cos\varphi\cos n\varphi\\
\sin(n+1)\varphi + \sin(n-1)\varphi &= 2\cos\varphi\sin n\varphi,
\end{aligned}
\]
from which the recursion formulae for $A_n$ and $B_n$ follow:
\weq{6}{
\begin{aligned}
A_{n+1}(x) &= x A_n(x) - A_{n-1}(x),\\
B_{n+1}(x) &= x B_n(x) - B_{n-1}(x).
\end{aligned}
}
From this one can compute $A_n$, $B_n$ for any arbitrary $n$, if these functions are known for $n = 1$ and $n = 2$. Thus one finds
%CONT>>


% ===== scan leaf 459, printed 434 =====
% !PRINTED 434
\weberpage{434}

\[
\begin{gathered}
A_4(x) = x^4 - 4x^2 + 2\\
A_5(x) = x^5 - 5x^3 + 5x\\
A_6(x) = x^6 - 6x^4 + 9x^2 - 2
\end{gathered}
\]
\[
\begin{gathered}
B_4(x) = x^3 - 2x\\
B_5(x) = x^4 - 3x^2 + 1\\
B_6(x) = x^5 - 4x^3 + 3x,
\end{gathered}
\]
and so one may continue. By induction one arrives at the following general law:
\weq{7}{A_n(x) = \overset{\nu}{\Sigma}(-1)^\nu \dfrac{n}{n-\nu} B_\nu^{(n-\nu)} x^{n-2\nu}, \quad 0 \leqq \nu \leqq \dfrac{n}{2},}
\weq{8}{B_n(x) = \overset{\nu}{\Sigma}(-1)^\nu B_\nu^{(n-\nu-1)} x^{n-2\nu-1}, \quad 0 \leqq \nu \leqq \dfrac{n-1}{2},}
where the sums begin with $\nu=0$, and are to be continued so long as the exponents of $x$ do not become negative.

Since these formulae prove to be correct in the first cases, in order to prove them generally it is only necessary to show that they verify the formulae (6). If for this purpose we form
\[
A_{n-1}(x) = \overset{\nu}{\Sigma}(-1)^\nu \dfrac{n-1}{n-\nu-1} B_\nu^{(n-\nu-1)} x^{n-2\nu-1},
\]
and put $\nu-1$ in place of $\nu$, it follows that
\weq{9}{\begin{gathered} A_{n-1}(x) = -\overset{\nu}{\Sigma}(-1)^\nu \dfrac{n-1}{n-\nu} B_{\nu-1}^{(n-\nu)} x^{n-2\nu+1}, \\ 1 \leqq \nu \leqq \dfrac{n+1}{2}. \end{gathered}}
If we assume $B_{-1}^{(n)} = 0$, we can let $\nu$ increase from 0 here as well.

The values that $\nu$ takes in formula (7) then differ from those in formula (9) only when $n$ is odd, because then the value $\nu = (n+1):2$ occurs in (9), but not in (7). We can, however, add this value in (7), provided that
\[
B_{\frac{n+1}{2}}^{\left(\frac{n-1}{2}\right)} = 0
\]
is assumed. Then there results
%CONT>>


% ===== scan leaf 460, printed 435 =====
% !PRINTED 435
\weberpage{435}
%<<CONT
\weq{10}{\begin{gathered} x A_n(x) - A_{n-1}(x) \\ = \overset{\nu}{\Sigma}(-1)^\nu x^{n-2\nu+1} \dfrac{n B_\nu^{(n-\nu)} + (n-1) B_{\nu-1}^{(n-\nu)}}{n-\nu}, \quad 0 \leqq \nu \leqq \dfrac{n+1}{2}. \end{gathered}}
Now, as follows easily from the expressions for the binomial coefficients (\S.\,8),
\[
\dfrac{n B_\nu^{(n-\nu)} + (n-1) B_{\nu-1}^{(n-\nu)}}{n-\nu} = \dfrac{n+1}{n-\nu+1} B_\nu^{(n-\nu+1)},
\]
also in the two limiting cases $\nu=0$, $\nu=(n+1):2$, and thus (10) agrees with what results from the right-hand side of (7) by interchanging $n$ with $n+1$.

In just the same way, formula (8) yields
\[
x B_n(x) - B_{n-1}(x) = \overset{\nu}{\Sigma}(-1)^\nu x^{n-2\nu}\left(B_\nu^{(n-\nu-1)} + B_{\nu-1}^{(n-\nu-1)}\right),
\]
and since
\[
B_\nu^{(n-\nu-1)} + B_{\nu-1}^{(n-\nu-1)} = B_\nu^{(n-\nu)}
\]
holds, (8) is thereby also proved in general.

For the sake of a clearer view we set out the formulae somewhat more fully:
\weq{11}{A_n(x) = x^n - n x^{n-2} + \dfrac{n(n-3)}{1\,.\,2} x^{n-4} - \dfrac{n(n-4)(n-5)}{1\,.\,2\,.\,3} x^{n-6} + \cdots}
\weq{12}{B_n(x) = x^{n-1} - (n-2) x^{n-3} + \dfrac{(n-3)(n-4)}{1\,.\,2} x^{n-5} - \dfrac{(n-4)(n-5)(n-6)}{1\,.\,2\,.\,3} x^{n-7} + \cdots}
If we take $\cos n\varphi$ and $\sin n\varphi$ as given, then the equation of degree $n$
\weq{13}{2 \cos n\varphi = A_n(x)}
serves to determine the unknown $x = 2\cos\varphi$. The meaning of the $n$ roots of this equation follows from the fact that the cosine does not change when the angle increases by a multiple of $2\pi$. Accordingly, equation (13) is satisfied by every value
\weq{14}{x = 2\cos\left(\varphi + \dfrac{2\nu\pi}{n}\right),}
where $\nu$ denotes an arbitrary positive or negative integer. But one obtains all the mutually distinct values of (14) by letting $\nu$ run through a complete system of residues with respect to $n$, thus setting, for example, $\nu = 0, 1, 2 \ldots n-1$, and one
%CONT>>


% ===== scan leaf 461, printed 436 =====
% !PRINTED 436
\weberpage{436}
%<<CONT
one also sees easily that these $n$ values are all different from one another, if one leaves aside the special case that $\varphi$ is a multiple of $\pi$. For two cosines are equal to one another only when the sum or the difference of the angles is a multiple of $2\pi$.

The equation
\weq{15}{\sin n\varphi = \sin\varphi\, B_n(x)}
determines $\sin\varphi$ uniquely in terms of $x$, when $B_n(x)$ is not equal to zero.

If we seek in equation (13) the term independent of $x$, then with the use of (6) there results for an even $n$
\[
2(-1)^{\frac{n}{2}} - 2\cos n\varphi,
\]
and for an odd $n$
\[
-2\cos n\varphi,
\]
and if one sets this equal to the positive or negative product of the roots, one obtains for an even $n$
\weq{16}{2^{n-1} \overset{\nu}{\Pi} \cos\left(\varphi + \dfrac{2\nu\pi}{n}\right) = (-1)^{\frac{n}{2}} - \cos n\varphi,}
and for an odd $n$
\weq{17}{2^{n-1} \overset{\nu}{\Pi} \cos\left(\varphi + \dfrac{2\nu\pi}{n}\right) = \cos n\varphi,}
and these formulae hold, since on the right and on the left stand continuous functions of $\varphi$, also for the case at first excluded, where $\varphi$ is a multiple of $\pi$.

Herein we may put $\nu m$ for $\nu$ under the product sign, if $m$ is any number coprime to $n$, because then $\nu m$ runs, together with $\nu$, through a complete system of residues to the modulus $n$. If we separate off the value $\nu=0$, and then let $\nu$ run through both the positive and the negative numbers that are absolutely smaller than $n:2$, we obtain from (17) the formula, valid for every odd $n$,
\weq{18}{2^{n-1} \overset{\nu}{\underset{1,\,\frac{n-1}{2}}{\Pi}} \cos\left(\varphi + \dfrac{2\nu m\pi}{n}\right) \cos\left(\varphi - \dfrac{2\nu m\pi}{n}\right) = \dfrac{\cos n\varphi}{\cos\varphi}.}
If we set $\varphi = 0$ herein and take the square root, there results
%CONT>>


% ===== scan leaf 462, printed 437 =====
% !PRINTED 437
\weberpage{437}
%<<CONT
\[
2^{\frac{n-1}{2}} \prod\limits_{1,\,\frac{n-1}{2}}^{\nu} \cos\frac{2\nu m\pi}{n} = \pm 1, \tag{19}
\]
in which for the time being the sign still remains undetermined.

But if in (18) we let $\varphi$ pass into $\pi:2$, the right-hand side, which presents itself under the indeterminate form $0:0$, acquires the limiting value $(-1)^{\frac{n-1}{2}} n$, and we obtain again, generally for every odd $n$, the formula which we derived in the preceding section for a prime $n$
\[
2^{\frac{n-1}{2}} \prod\limits_{1,\,\frac{n-1}{2}}^{\nu} \sin\frac{2\nu m\pi}{n} = \pm \sqrt{n}, \tag{20}
\]
and for $m=1$ the positive sign is to be taken.

The function $B_n$, which with respect to $x$ is of degree $(n-1)$, vanishes when $n\varphi$, but not $\varphi$, is a multiple of $\pi$, that is, for
\[
x = 2\cos\frac{h\pi}{n}, \tag{21}
\]
in which $h$ is an integer not divisible by $n$. If we first assume $n$ even, we obtain all the \gesperrt{distinct} values contained in (21) by taking for $m$ an arbitrary number coprime to $n$ and setting
\[
h = m,\ 2m,\ 3m \ldots (n-1)m \tag{22}
\]
for
\[
\cos\frac{h\pi}{n} = \cos\frac{h'\pi}{n},
\]
holds only when either $h+h'$ or $h-h'$ is a multiple of $2n$, which cannot occur when $h$ and $h'$ are both taken from the sequence (22).

Among the values (21) there is one, namely the one corresponding to the value $h = nm:2$, which is equal to zero, and the rest are opposite in pairs. Accordingly, for an \gesperrt{even} $n$ we obtain
\[
B_n(x) = x \prod\limits_{1,\,\frac{n}{2}-1}^{\nu} \left(x^2 - 4\cos^2\frac{\nu m\pi}{n}\right). \tag{23}
\]

If $n$ is odd, then $B_n(x)$ depends only on $x^2$; hence the values (21) are equal and opposite in pairs (for $h$ and $h+n$). If, again supposing $m$ relatively prime to $n$, one
%CONT>>


% ===== scan leaf 463, printed 438 =====
% !PRINTED 438
\weberpage{438}
%<<CONT
\[
h = 2m,\ 4m,\ 6m,\ \ldots (n-1)m
\]
assumes, one obtains in (21) neither two equal nor two opposite values, and all the mutually distinct values of $x^2$ result.

Thus for an \gesperrt{odd} $n$ we have
\[
B_n(x) = \prod\limits_{1,\,\frac{n-1}{2}}^{\nu} \left(x^2 - 4\cos^2\frac{2\nu m\pi}{n}\right). \tag{24}
\]

If in this formula we set
\[
x = 2\cos\varphi,\quad x^2 = 4(1-\sin^2\varphi),\quad \cos^2\frac{2\nu m\pi}{n} = 1 - \sin^2\frac{2\nu m\pi}{n},
\]
we can represent it thus:
\[
\frac{\sin n\varphi}{\sin\varphi} = (-1)^{\frac{n-1}{2}} 2^{n-1} \prod\limits_{1,\,\frac{n-1}{2}}^{\nu} \left(\sin^2\varphi - \sin^2\frac{2\nu m\pi}{n}\right), \tag{25}
\]
from which again, for $\varphi=0$, formula (20) follows.

Since $\cos n\varphi$ vanishes for odd $n$ when
\[
\cos\varphi = \pm \sin\frac{2\nu m\pi}{n}
\]
holds, we have thereby found the decomposition into factors of the function $A_n(x)$ for this case also. For \gesperrt{odd} $n$ we have
\[
A_n(x) = x \prod\limits_{1,\,\frac{n-1}{2}}^{\nu} \left(x^2 - 4\sin^2\frac{2\nu m\pi}{n}\right). \tag{26}
\]

The quantities
\[
\alpha = \sin^2\frac{2\nu m\pi}{n} \tag{27}
\]
are thus the roots of an algebraic equation of degree $(n-1):2$
\[
\Phi_n(x) = 0,
\]
and we obtain the function $\Phi_n(x)$ if, in place of the variable $x^2$ in the function $A_n(x):x$, we put a new variable $4x$ and divide by $2^{n-1}$. It therefore becomes [formula (11)]
\[
\Phi_n(x) = \overset{\alpha}{\Pi}(x-\alpha)
= x^{\frac{n-1}{2}} - \frac{n}{4} x^{\frac{n-3}{2}} + \frac{n(n-3)}{1.2}\,\frac{1}{16}\, x^{\frac{n-5}{2}} - \cdots, \tag{28}
\]
and for (25) we can write
\[
\frac{\sin n\varphi}{\sin\varphi}
= (-1)^{\frac{n-1}{2}} 2^{n-1} \overset{\alpha}{\Pi}(\sin^2\varphi - \alpha). \tag{29}
\]


% ===== scan leaf 464, printed 439 =====
% !PRINTED 439
\weberpage{439}
\pdfbookmark[2]{§ 138. Determination of the sign. Quadratic residues}{bm153}
\begin{center}
\S.\,138.

\gesperrt{Determination of signs. Quadratic residues.}
\end{center}

In the preceding section we derived two formulae in which signs still remained to be determined, namely
\[
2^{\frac{n-1}{2}} \overset{\nu}{\Pi} \cos\frac{2\nu m\pi}{n} = \pm 1, \tag{1}
\]
\[
2^{\frac{n-1}{2}} \overset{\nu}{\Pi} \sin\frac{2\nu m\pi}{n} = \pm \sqrt{n}, \tag{2}
\]
where $m$ is relatively prime to the odd number $n$ and $\nu$ runs through the following sequence.
\[
1,\ 2,\ \cdots\ \frac{n-1}{2} \tag{$\nu$}
\]
We shall here retain the symbol $\nu$ for the numbers of this sequence.

Every integer whatsoever gives, upon division by $n$, as residue one of the numbers $\nu$ or $0$ or $n-\nu$. Instead of $n-\nu$ we may, if we also admit negative residues, choose $-\nu$; and if we accordingly set aside the numbers divisible by $n$, there remains one of the numbers $\pm\nu$ as residue. We call this the \gesperrt{absolutely least residue} (in distinction from the least residue in the ordinary sense, which is taken from the numbers $1, 2, \ldots n-1$).

Now let $m$ be a number coprime to $n$; we consider the sequence of numbers
\[
m,\ 2m,\ 3m \cdots \frac{n-1}{2}\, m \tag{$m\nu$}
\]
and form for each the absolutely least residue $\varrho$:
\[
\pm \nu_1,\ \pm \nu_2,\ \pm \nu_3 \cdots \pm \nu_{\frac{1}{2}(n-1)}; \tag{$\varrho$}
\]
among these numbers $\varrho$ no two equal ones occur, and also no two opposite ones; for if the sum or the difference of two numbers $(\varrho)$, that is $\varrho+\varrho'$ or $\varrho-\varrho'$, were equal to zero, then for two distinct numbers $\nu, \nu'$ from $(\nu)$
\[
m(\nu \pm \nu')
\]
would have to be divisible by $n$, and hence $\nu \pm \nu'$ would also have to be divisible by $n$, which is impossible, since each of the numbers $\nu, \nu'$ is less than $n:2$,
%CONT>>


% ===== scan leaf 465, printed 440 =====
% !PRINTED 440
\weberpage{440}
%<<CONT
\gesperrt{The totality of the numbers $(\varrho)$ thus agrees, apart from sign and from order, with the totality of the numbers $(\nu)$.}

In the formulae (1) and (2) we can, however, on account of the periodicity of sine and cosine, replace $\nu m$ by $\varrho$, and since $\cos(-\varphi) = \cos\varphi$, we can in formula (1) also replace $m\nu$ by $\nu$, i.e. the sign does not depend on $m$. To determine it, we take into account that the cosine of an angle in the first quadrant is positive, in the second quadrant negative, and thus that the sign in (1) is positive or negative according as an even or an odd number of the angles
\[
\frac{2\nu\pi}{n}
\]
lies in the second quadrant, or according as an even or an odd number of the numbers $\nu$ lies between $\dfrac{n}{4}$ and $\dfrac{n}{2}$.

If $x$ is any number that is not an integer, we denote by $E(x)$ the greatest integer that is less than $x$, so that $x$ lies between $E(x)$ and $E(x)+1$. The number of integers lying between two numbers $x$ and $y$ is then equal to $E(y) - E(x)$, and we therefore have to investigate whether
\[
E\left(\frac{n}{2}\right) - E\left(\frac{n}{4}\right)
\]
is an even or an odd number. We must distinguish four cases, as appears from the following compilation, in which $k$ denotes a non-negative integer:
\[
\begin{array}{lll}
n = 8k+1, & E\left(\dfrac{n}{2}\right) = 4k, & E\left(\dfrac{n}{4}\right) = 2k, \\[2ex]
n = 8k+3, & E\left(\dfrac{n}{2}\right) = 4k+1, & E\left(\dfrac{n}{4}\right) = 2k, \\[2ex]
n = 8k+5, & E\left(\dfrac{n}{2}\right) = 4k+2, & E\left(\dfrac{n}{4}\right) = 2k+1, \\[2ex]
n = 8k+7, & E\left(\dfrac{n}{2}\right) = 4k+3, & E\left(\dfrac{n}{4}\right) = 2k+1.
\end{array}
\]
Thus in (1) the positive sign is to be taken when $n$ is of the form $8k+1$ or $8k+7$, the negative when $n$ is of the form $8k+3$ or $8k+5$. In the first cases
%CONT>>


% ===== scan leaf 466, printed 441 =====
% !PRINTED 441
\weberpage{441}
%<<CONT
$\dfrac{n^2-1}{8}$ is an even number, in the last two an odd one,
and accordingly we obtain the exact formula (1):

\weq{3}{2^{\frac{n-1}{2}} \overset{\nu}{\Pi} \cos \frac{2\nu m \pi}{n} = (-1)^{\frac{n^2-1}{8}}.}

In formula (2) the sign depends on $m$; it is positive or
negative according as, among the numbers $(\varrho)$, an even or an
odd number of negative ones occurs.

Now, in order to express in the notation the dependence of
the sign on $m$ and $n$, we write for (2)

\weq{4}{2^{\frac{n-1}{2}} \overset{\nu}{\underset{1,\,\frac{n-1}{2}}{\Pi}} \sin \frac{2\nu m \pi}{n} = \left(\frac{m}{n}\right)\sqrt{n},}

in which

\[
\left(\frac{m}{n}\right) = \pm 1
\]

holds.

The symbol $\left(\dfrac{m}{n}\right)$ has been introduced into number theory
from another side, and in a more special form, by \gesperrt{Legendre},
and has been generalised by \gesperrt{Jacobi}. We shall
call it the \gesperrt{Legendre symbol}. Its significance for
number theory will appear presently; first we have
the theorem:

1. \gesperrt{We have} $\left(\dfrac{m}{n}\right)$ \gesperrt{equal to} $+1$ \gesperrt{or equal to} $-1$,
\gesperrt{according as the number of negative ones among the absolutely
least residues of} $m\nu$ \gesperrt{is even or odd.}

For $m = 1$, as we have already seen, the positive sign
holds in formula (2); hence we have, as formula
(4) also shows directly,

\[
2.\quad \left(\frac{1}{n}\right) = +1.
\]

If we change $m$ into $-m$, the sign changes in all factors
of the product on the left-hand side of (4); since
the number of factors is $\dfrac{n-1}{2}$, it follows that

\[
3.\quad \left(\frac{-m}{n}\right) = (-1)^{\frac{n-1}{2}} \left(\frac{m}{n}\right)
\]

\gesperrt{and for} $m = 1$
%CONT>>


% ===== scan leaf 467, printed 442 =====
% !PRINTED 442
\weberpage{442}
%<<CONT
\[
4.\quad \left(\frac{-1}{n}\right) = (-1)^{\frac{n-1}{2}}.
\]

If we put $2m$ for $m$ and use the formula

\[
\sin \frac{4\nu m \pi}{n} = 2 \sin \frac{2\nu m \pi}{n} \cos \frac{2\nu m \pi}{n},
\]

then from (3) it follows that

\[
5.\quad \left(\frac{2m}{n}\right) = (-1)^{\frac{n^2-1}{8}} \left(\frac{m}{n}\right)
\]

\gesperrt{and for} $m = 1$

\[
6.\quad \left(\frac{2}{n}\right) = (-1)^{\frac{n^2-1}{8}}.
\]

If $m$ is altered by a multiple of $n$, all the
factors of the product (4) remain unchanged, and from this it follows:

7. \gesperrt{We have}

\[
\left(\frac{m}{n}\right) = \left(\frac{m'}{n}\right),
\]

\gesperrt{if} $m \equiv m'$ \gesperrt{(mod} $n$\gesperrt{).}

We now assume not only $n$ but also $m$ to be odd
and positive, and apply the formula (29) of \S.\,137,
in which we replace $n$ by $m$:

\weq{5}{\frac{\sin m \varphi}{\sin \varphi} = (-1)^{\frac{m-1}{2}} 2^{m-1} \overset{\beta}{\Pi} (\sin^2 \varphi - \beta),}

in which $\beta$ runs through all the roots of the equation $\Phi_m(x) = 0$.

\weq{6}{\beta = \sin^2 \frac{2\mu\pi}{m}}

Now if $m'$ is a number coprime to $n$, then

\weq{7}{\alpha = \sin^2 \frac{2\nu m' \pi}{n}}

runs through all the roots of the equation $\Phi_n(x) = 0$, and if
we set

\weq{8}{\varphi = \frac{2\nu m' \pi}{n}}

in (5) and take the product with respect to all $\nu$, it follows that:

\weq{9}{\frac{\overset{\nu}{\Pi} \sin \dfrac{2\nu m m' \pi}{n}}{\overset{\nu}{\Pi} \sin \dfrac{2\nu m' \pi}{n}} = (-1)^{\frac{(m-1)(n-1)}{4}} 2^{\frac{(m-1)(n-1)}{2}} \overset{\alpha}{\Pi}\overset{\beta}{\Pi} (\alpha - \beta).}


% ===== scan leaf 468, printed 443 =====
% !PRINTED 443
\weberpage{443}
But the right-hand side of this formula is entirely independent
of $m'$. The value of the left-hand side is therefore the same as if
$m' = 1$, and if one substitutes for the products the values from
(4), it follows that

\[
\left(\frac{mm'}{n}\right) : \left(\frac{m'}{n}\right) = \left(\frac{m}{n}\right) : \left(\frac{1}{n}\right),
\]

or

\[
8.\quad \left(\frac{mm'}{n}\right) = \left(\frac{m}{n}\right)\left(\frac{m'}{n}\right).
\]

Here it was at first assumed that $m$ and $m'$ are positive and
odd. But by 3. and 5., 8. still remains valid
when $m$ or $m'$ is negative or even, provided only that $m$ and $m'$
are relatively prime to $n$.

If we again take $m$ odd and positive, it now follows
from (4) and (9)

\weq{10}{\left(\frac{m}{n}\right) = (-1)^{\frac{(m-1)(n-1)}{4}} 2^{\frac{(m-1)(n-1)}{2}} \overset{\alpha}{\Pi}\overset{\beta}{\Pi} (\alpha - \beta).}

If we interchange $m$ with $n$ here, then in the
double product on the right-hand side every factor changes its sign;
everything else remains unchanged. But the number of these factors
is

\[
\frac{n-1}{2} \cdot \frac{m-1}{2}
\]

and from this it follows

\[
9.\quad \left(\frac{m}{n}\right) = (-1)^{\frac{(m-1)(n-1)}{4}} \left(\frac{n}{m}\right),
\]

\gesperrt{or} $\left(\dfrac{m}{n}\right)$ \gesperrt{is equal to} $\left(\dfrac{n}{m}\right)$\gesperrt{, if of the two
odd numbers} $m, n$ \gesperrt{at least one is of the form}
$4k+1$ \gesperrt{, and} $\left(\dfrac{m}{n}\right)$ \gesperrt{is opposite
 to} $\left(\dfrac{n}{m}\right)$\gesperrt{, if both numbers are of the
form} $4k+3$\gesperrt{.}

This famous theorem is known under the name of the \gesperrt{law of reciprocity}.
 The proof given here is due to
\gesperrt{Eisenstein}. With its help, and by 5. and
6., one can determine the value of the symbol $\left(\dfrac{m}{n}\right)$ very quickly, by
%CONT>>


% ===== scan leaf 469, printed 444 =====
% !PRINTED 444
\weberpage{444}
%<<CONT
one proceeds as if it were a matter of determining the greatest
common divisor of $m$ and $n$.

From 8. and 9. there follows yet one last theorem, which holds when
$n$ and $n'$ are two positive odd numbers coprime to $m$:

\[
10.\quad \left(\frac{m}{n}\right)\left(\frac{m}{n'}\right) = \left(\frac{m}{nn'}\right).
\]

Its correctness follows, when $m$ is odd and positive,
on the ground of the congruence

% [sic] Within §138 the source reuses equation label (7), already printed on p.442 for the alpha formula; preserved diplomatically.
\weq{7}{(n-1)(n'-1) = (nn'-1) - (n-1) - (n'-1) \equiv 0 \ (\mathrm{mod}\ 4),}

if one applies formula 9. and then 8. to both sides of 10.;
and, when $m$ is even or negative, from 4. and 6.

By these theorems one can reduce $\left(\dfrac{m}{n}\right)$ to a product of
values $\left(\dfrac{q}{p}\right)$, in which $p, q$ are primes. The
computation of $\left(\dfrac{m}{n}\right)$, however, is carried out more easily by 9. without this
decomposition.

We now wish to learn another meaning of the symbol
$\left(\dfrac{m}{p}\right)$ in the case where $p$ is an odd prime,
for which \gesperrt{Legendre} originally introduced the symbol.

If we raise to the square all the integers $x$ not divisible by $p$
and seek out the residues of the division by $p$,
we obtain in all only $\dfrac{p-1}{2}$ different
residues; for, first, the squares of congruent numbers give the same
residues, and secondly the squares of opposite
numbers also give the same residues. We therefore certainly obtain all the
residues if we take the $\dfrac{p-1}{2}$ numbers $\nu$

\[
(\nu)\qquad 1,\ 2,\ 3\cdots\frac{p-1}{2}
\]

and raise them to the square. On the other hand, the
squares of these numbers $\nu$ also yield residues that are all different; for,
if $\nu$ and $\nu'$ are different ones among them, the quantity

\[
\nu^2 - \nu'^2 = (\nu - \nu')(\nu + \nu')
\]
can never
%CONT>>


% ===== scan leaf 470, printed 445 =====
% !PRINTED 445
\weberpage{445}
be divisible by $p$, because both $\nu - \nu'$ and $\nu + \nu'$ are smaller than
$p$. Among the numbers

\weq{s}{1, 2, 3 \ldots p - 1}

half therefore occur among the residues of $x^2$, and half do not.
The former numbers, and all those congruent to them with respect to $p$,
are called the \gesperrt{quadratic residues of $p$}, the others the
\gesperrt{quadratic non-residues of $p$}.

Now if $m$ is a quadratic residue, the congruence
\[
x^2 \equiv m \ (\mathrm{mod}\ p)
\]
is possible, and from 7. and 8. it follows that
\[
\left(\frac{m}{p}\right) = \left(\frac{x^2}{p}\right) = \left(\frac{x}{p}\right)^2 = +1.
\]
Thus $\left(\dfrac{m}{p}\right) = +1$ when $m$ is a quadratic residue of
$p$.

But in \S.\,135, 1. the theorem was stated that, if for $m$ the
sequence of numbers 1, 2, \ldots p - 1 is put,
$\left(\dfrac{m}{p}\right)$ has the positive sign just as often as the
negative one; and from this one sees that the negative sign holds
when $m$ is a quadratic non-residue. It therefore follows:

11. \gesperrt{If $p$ is an odd prime and $m$ is not divisible
by $p$, then $\left(\dfrac{m}{p}\right) = +1$ or $= -1$,
according as $m$ is a quadratic residue or a quadratic non-residue of
$p$.}
In this all the principal theorems on the quadratic residues
are contained, e.\,g. the theorem:

The product of two quadratic residues or of two non-residues
is a residue; the product of a residue and a non-residue is
a non-residue.

\begin{center}
\rule{0.25\linewidth}{0.4pt}
\end{center}


% ===== scan leaf 471, printed 446 =====
\weberpage{446}% page number not printed on this page
% !PRINTED none
% Blank verso after printed 445. No direct black text, running head, folio, or
% rule. Faint offset/show-through covers the top, middle, and lower text fields,
% including fraction-display shapes; it is not page content.


% ===== scan leaf 472, printed 447 =====
\weberpage{447}% page number not printed on this page
% !PRINTED none
% Part-title recto opening Book Three: no leading rule. The letterspaced book
% heading comes first, followed by a short inter-title rule, the letterspaced
% title, and a longer lower rule.

\begin{center}
\gesperrt{THIRD BOOK.}

\rule{0.15\linewidth}{0.4pt}

\bigskip
\gesperrt{ALGEBRAIC QUANTITIES.}
\bigskip

\rule{0.25\linewidth}{0.4pt}
\end{center}


% ===== scan leaf 473, printed 448 =====
\weberpage{448}% page number not printed on this page
% !PRINTED none
% Blank verso of the Book Three part-title leaf. No direct-printed text, rule,
% running head, or folio. Both title lines are visible only as faint,
% horizontally reversed show-through from the recto at i0472.


% ===== scan leaf 474, printed 449 =====
\weberpage{449}% page number not printed on this page
\pdfbookmark[0]{Third Book. Algebraic Quantities.}{bm154}
\pdfbookmark[1]{Thirteenth Section. The Galois Theory.}{bm155}
\pdfbookmark[2]{§ 139. The concept of a field}{bm156}
% !PRINTED none
% Abschnitt-opening / part-title page: no page number in running head. Footer furniture
% "Weber, Algebra. I." and signature "29" are ignored per standing policy.

\begin{center}
\gesperrt{Thirteenth Section.}

\gesperrt{The Galois Theory.}
\end{center}

\begin{center}
\rule{3cm}{0.4pt}
\end{center}

\begin{center}
\S.\,139.

\gesperrt{The Concept of a Field.}
\end{center}

A system of numbers is called a \gesperrt{number field} when it is so complete and closed in itself that the four fundamental operations of reckoning (the four species) — \gesperrt{addition}, \gesperrt{subtraction}, \gesperrt{multiplication} and \gesperrt{division} — carried out with any numbers of the system, excepting division by zero, always lead to numbers that belong to the same system. This concept, which gives a classification of the kinds of numbers according to a natural point of view, was introduced by \gesperrt{Dedekind} (\gesperrt{Dirichlet-Dedekind}, Vorlesungen \"uber Zahlentheorie, 2nd ed. 1871, \S.\,159). It is of the greatest importance for algebra, and it is not a matter of indifference to have a distinctive and expressive name for it. The word \gesperrt{Zahlk\"orper} (number field) was formed by \gesperrt{Dedekind} on the model of numerous analogies, in which the word K\"orper (corpus, corps) in a similar way signifies a union of things belonging together, to which a certain completeness attaches.

The concept of the number field can be extended and carried over to all magnitudes with which one can reckon according to the rules of the four species; in particular, therefore, to rational functions of any variables.

We accordingly call a \gesperrt{function field} a system of functions of one or more variables of such a nature that in this system the reckonings with the four species can be carried out without limit and always lead to a definite function of the same system (always with the exception of division by zero). The variables
%CONT>>


% ===== scan leaf 475, printed 450 =====
% !PRINTED 450
\weberpage{450}
%<<CONT
may be entirely independent of one another; but the case is not excluded that certain relations are fixed between them, which are to be taken into account in calculation.

A function of a function field counts as zero only when it vanishes identically, that is, for all values of the variables that come into consideration.

Since for the present we do not wish to restrict our considerations, we shall now speak of \gesperrt{fields} simply, and shall designate the objects with which the calculations are to be carried out, which may be numbers as well as functions, as \gesperrt{quantities} or also as the \gesperrt{elements} of the field $^1$).

A \gesperrt{field} is then a system of quantities of such completeness that in it the quantities can be added, subtracted, multiplied and divided.

When all the quantities of one field are contained in a second field, the first field is called a \gesperrt{divisor} of the second.

If $\Omega$ is a divisor of $\Omega'$, we shall also say that $\Omega'$ \gesperrt{is a field over} $\Omega$.

% [sic] The German source prints the anomalous hyphenated phrase "ein- für allemal"; the idiomatic translation is retained.
Strictly speaking, the single number ``zero'' forms a field by itself. But for the sake of brevity we shall exclude this once and for all.

Then the most obvious example of a field is the \gesperrt{totality of all rational numbers}.

This field is a divisor of every other field; for every field contains at least one quantity $\omega$ different from zero, hence also the quotient $\omega:\omega=1$, that is, the number 1; hence also, since all integers arise by addition and subtraction of units, all integers; and from the integers one can again derive all fractions by division.

Other examples of number fields are: the totality of all complex numbers of the form $x+yi$, in which $i=\sqrt{-1}$, and $x$ and $y$ denote all rational numbers; further the totality of all (rational and irrational) \gesperrt{real numbers}, or the totality of all complex numbers $x+yi$ that exist at all.

As an example of a function field may serve the totality of all integral and fractional rational functions of a
%CONT>>

\par\noindent\rule{4cm}{0.4pt}\par
{\footnotesize $^1$) Cf. the author's paper: ``Die allgemeinen Grundlagen der \gesperrt{Galois}'schen Gleichungstheorie.'' Mathematische Annalen, Bd.\,43.}


% ===== scan leaf 476, printed 451 =====
% !PRINTED 451
\weberpage{451}
\pdfbookmark[2]{§ 140. Adjunction}{bm157}
%<<CONT
variables (with the inclusion of the constants), in which one may restrict the constant coefficients to an arbitrary number field, say to that of the rational numbers, or may also take them entirely unrestricted.

We do not wish to multiply examples further here, since the exact investigation of fields of special kinds will be one of the principal tasks of our later expositions.

\begin{center}
\S.\,140.

\gesperrt{Adjunction.}
\end{center}

If to a field of any magnitudes whatever, which we shall denote by $\Omega$, there is added any magnitude $\alpha$ that is not contained in it, there arises a new system of magnitudes
\[
\Omega,\ \alpha,
\]
which, however, will not be a field; in order to derive a field from it, one must add all magnitudes that can be derived by combining $\alpha$ with the numbers of $\Omega$ by means of the fundamental operations of arithmetic. One thus obtains an extended field $\Omega'$, which contains both $\alpha$ and $\Omega$, and which is completely determined by $\alpha$ and $\Omega$. We shall call it the field $\Omega, \alpha$.

This adding of a new magnitude $\alpha$ to a field $\Omega$ is called \gesperrt{adjoining}, and one says that the field $\Omega'$ arises from $\Omega$ by \gesperrt{adjunction of} $\alpha$.

To $\Omega'$ one can again adjoin a magnitude $\alpha'$, and obtains a third field $\Omega''$, which then contains $\Omega$, $\alpha$ and $\alpha'$, and so on. But one can also adjoin two or more numbers to $\Omega$ simultaneously, and thus obtains the same field $\Omega''$ if one adjoins $\alpha$ and $\alpha'$ to $\Omega$ simultaneously. One can also adjoin to one field a second field, and thus obtains a field that contains the numbers of both fields at once. In the above example one would, by adjunction of the field $\Omega, \alpha'$ to the field $\Omega' = \Omega, \alpha$, again obtain the same field $\Omega''$.

If one adjoins to a field one of its parts, no new field arises.

Thus, for example, by adjoining to the field of rational numbers, which may be denoted by $\mathfrak{R}$, the number $i=\sqrt{-1}$
%CONT>>


% ===== scan leaf 477, printed 452 =====
% !PRINTED 452
\weberpage{452}
\pdfbookmark[2]{§ 141. Functions in a field}{bm158}
%<<CONT
one obtains the field of complex numbers $x+yi$, with rational $x, y$, which we shall call $\mathfrak{J}$. By the adjunction of a variable $u$ to the field $\mathfrak{R}$ one obtains the field which consists of all integral and fractional rational functions of $u$ with rational coefficients; by the adjunction of $u$ and $i$ to $\mathfrak{R}$, the field of the rational functions of $u$ whose coefficients are numbers in $\mathfrak{J}$, and so on.

Let it finally be remarked that a \gesperrt{field} of quantities is often also called, following \gesperrt{Kronecker}'s example, a \gesperrt{domain of rationality}. It is sometimes useful to employ the expression domain of rationality for a field when it is to be expressed thereby that, in a particular problem, the quantities of this field are to be regarded as known or rational.

\begin{center}
\S.\,141.

\gesperrt{Functions in a Field.}
\end{center}

We often have occasion to consider integral rational functions of a variable $x$
\weq{1}{f(x) = a_0 x^m + a_1 x^{m-1} + \cdots + a_{m-1} x + a_m}
whose coefficients belong to a particular field $\Omega$, and which we shall call \gesperrt{functions contained in} $\Omega$, or briefly \gesperrt{functions in} $\Omega$. They are not to be confused with functions which, say, \gesperrt{belong} to the field $\Omega$, when $\Omega$ is a function field.

In the third section we saw that, when the coefficients $a_0, a_1 \ldots a_m$ have definite numerical values, there exist $m$ roots $\alpha_1, \alpha_2 \ldots \alpha_m$ of $f(x)$, which likewise have definite numerical values, among which, however, several equal to one another occur when the discriminant $\Delta(a)$ of the function (1) vanishes. If the coefficients $a$ of $f(x)$ are independent variables, then for every definite system of values of the $a$ there is a definite system of values of the $\alpha$, and the $\alpha$ vary continuously with the $a$ (\S.\,40). So long, therefore, as we restrict the variation of the $a$ to a sufficiently narrow region, in which no such values occur for which the discriminant $\Delta(a)$ vanishes, the $m$ quantities $\alpha$, as functions of the $a$, are completely determined and distinguished from one another as soon as the value of each of them is given for some particular system of values of the $a$.
%CONT>>


% ===== scan leaf 478, printed 453 =====
% !PRINTED 453
\weberpage{453}

If the $a$ are not independent variables, but are themselves in turn functions of other variables, all this remains valid, so long only as the changes of the $a$ take place continuously.

Accordingly, in all these cases we may speak of the $m$ roots of an equation of degree $m$, which are to be understood as numbers or as algebraic functions of their coefficients.

Now if $\Omega$ is an arbitrary field and $f(x)$ is a function in $\Omega$, two cases are possible.

The function $f(x)$ is either decomposable or not decomposable into factors of lower degree whose coefficients belong to the field $\Omega$. In the first case the function $f(x)$ is called \gesperrt{reducible}, in the second \gesperrt{irreducible} in $\Omega$. When the field $\Omega$ may be regarded as known and fixed, for instance when it is the field of rational numbers, one sometimes also speaks simply of reducible or irreducible functions, without expressly adding the designation ``in the field $\Omega$''. That in itself, however, the precise specification of the field to which the reducibility refers, or which is regarded as the domain of rationality, is very essential, follows from the following remarks.

A linear function is of course irreducible in every field. If one multiplies several functions of the form (1) with one another, there arises a function of the same form, which, however, is then naturally reducible, since it can again be decomposed into the factors from which it arose.

A function $f(x)$ which is irreducible in $\Omega$ can become reducible in an extended field $\Omega'$ which arises from $\Omega$ by any adjunction; thus every function $f(x)$ becomes reducible in the field $\Omega'$ which arises from $\Omega$ by the adjunction of a root $\alpha$ of $f(x) = 0$; for a linear factor $x - \alpha$ can be separated off from $f(x)$. In the field; which consists of \gesperrt{all} numbers, every function with numerical coefficients is reducible. % [sic] semicolon as printed

I. \gesperrt{An irreducible function $f(x)$ can have no common divisor with another function $F(x)$ whose coefficients belong to the same field, unless $F(x)$ is divisible by $f(x)$.}


% ===== scan leaf 479, printed 454 =====
% !PRINTED 454
\weberpage{454}

This theorem, which will render us very important services in what follows, is almost self-evident; for the greatest common divisor of two functions $F(x)$ and $f(x)$ can be found by rational computation, and is therefore also contained in $\Omega$. But since $f(x)$ has no divisor contained in $\Omega$ except itself or a constant (i.\,e. a quantity contained in $\Omega$), this greatest common divisor must therefore be either a constant or $f(x)$ itself.

It follows from this theorem that an irreducible function can never have multiple factors; for otherwise it would have to have a common divisor with its derivative $f'(x)$; thus $f'(x)$ would have to be divisible by $f(x)$, which is impossible, since the degree of $f'(x)$ is lower than that of $f(x)$.

We can also give Theorem I the following expression:

\gesperrt{II. If $f(x)$ is irreducible, and $F(x)$ vanishes for one root of $f(x) = 0$, then it also vanishes for all other roots of $f(x) = 0$.}

In particular we can conclude from this:

\gesperrt{III. If $F(x)$ is of lower degree than the irreducible function $f(x)$, and $F(x)$ vanishes for a root of $f(x)$, then $F(x)$ must vanish identically, i.\,e. all coefficients of $F(x)$ must be zero.}

Integral rational functions of several variables, $f(x, y, z \ldots)$, whose coefficients belong to the field $\Omega$, are likewise called functions in $\Omega$. These too are designated as \gesperrt{reducible} and \gesperrt{irreducible} functions, according as they are decomposable or not decomposable into several factors which are themselves functions in $\Omega$, and actually contain at least some of the variables $x, y, z \ldots$.

Here, however, attention must be called to a distinction.

Among the functions of more than one independent variable there are some that are not decomposable at all into factors of lower degree which depend rationally on the variables; and again others that are indeed decomposable, but only into factors whose coefficients are not contained in $\Omega$; and finally, thirdly, functions that are decomposable into several functions in $\Omega$; only the last kind shall we call reducible in $\Omega$, while those of the second kind are described as \gesperrt{decompos-}
%CONT>>


% ===== scan leaf 480, printed 455 =====
% !PRINTED 455
\weberpage{455}
\pdfbookmark[2]{§ 142. Algebraic fields}{bm159}

%<<CONT
\gesperrt{able, those of the first kind being designated as indecomposable functions}. If, for example, $\Omega$ is the field of rational numbers, then $x^2 - y^2 = (x - y)(x + y)$ is reducible in $\Omega$. The function $x^2 - 2y^2$ is indeed decomposable into the factors $x + y\sqrt{2}, x - y\sqrt{2}$, but these factors are not functions in $\Omega$. The function $x^2 + y^2 + 1$, finally, is not decomposable at all. A function irreducible in a field $\Omega$, but decomposable, becomes reducible in an extended field.

The fundamental theorems on irreducible functions we came to know in \S.\,51. In those developments, it is true, the field of all numbers was at first assumed as the domain of rationality, in accordance with the standpoint there taken, in which field the distinction between indecomposable and irreducible functions disappears. But only the reducibility or irreducibility of a function is ever made use of, and the particular nature of the field $\Omega$ does not come into consideration. Thus there is contained in it the proof of the theorem:

\gesperrt{IV. A reducible function in $\Omega$ can be decomposed into irreducible factors in only one way, if rational functions which differ only by a constant factor are regarded as not distinct.}

\begin{center}
§. 142.\\
\gesperrt{Algebraic Fields.}
\end{center}

If $\Omega$ is an arbitrary field, and $F(x)$ a function in $\Omega$, then the equation
\weq{1}{F(x) = 0}
is called an \gesperrt{equation in} $\Omega$. This equation is called reducible or irreducible according as the function $F(x)$ is reducible or irreducible. By the adjunction of a root $\alpha$ of such an equation to $\Omega$ there arises (if $\alpha$ does not itself already belong to $\Omega$) a new field $\Omega'$, which we shall call an \gesperrt{algebraic field} ``\gesperrt{over}'' $\Omega$, or also, when doubt as to the meaning is excluded, briefly, an \gesperrt{algebraic field}. For such a field we use the sign
\[ \Omega(\alpha). \]


% ===== scan leaf 481, printed 456 =====
% !PRINTED 456
\weberpage{456}

If $\beta, \gamma \ldots$ are roots of other equations in $\Omega$, or also other roots of the same equation, then by simultaneous adjunction of $\alpha, \beta, \gamma \ldots$ we likewise obtain algebraic fields over $\Omega$, which we denote by
\[ \Omega(\alpha, \beta, \gamma \ldots) \]
We shall see presently that this extension of the concept of algebraic fields is only an apparent one, and so for the time being we stop at the adjunction of \gesperrt{one} quantity $\alpha$, that is, at the consideration of $\Omega(\alpha)$.

The equation (1), whose root is $\alpha$, may be reducible. But among the irreducible factors of $F(x)$ there is at least one that vanishes for $x = \alpha$, which we shall denote by $f(x)$; this function $f(x)$ is, if we take the coefficient of the highest power of $x$ equal to 1, completely determined, because by the preceding section $x = \alpha$ cannot be the root of two different irreducible equations.

The equation $f(x) = 0$ thus has the form
\weq{2}{x^n + a_1x^{n-1} + a_2x^{n-2} + \cdots + a_{n-1}x + a_n = 0,}
in which $a_1, a_2 \ldots a_n$ are quantities in $\Omega$. If $n$ is the degree of this equation, then we also call $n$ the \gesperrt{degree of the field} $\Omega(\alpha)$.

All quantities $\Theta$ of the field $\Omega(\alpha)$ can be derived by addition, subtraction, multiplication and division from $\alpha$ and from numbers in $\Omega$; they can therefore be represented as rational functions of $\alpha$ with coefficients in $\Omega$, or in the form
\weq{3}{\Theta = \dfrac{\chi(\alpha)}{\psi(\alpha)},}
in which $\chi(x)$ and $\psi(x)$ are functions in $\Omega$. But since $\psi(\alpha)$ must be different from zero, $\psi(x)$ is not divisible by $f(x)$, and since $f(x)$ is assumed irreducible, $\psi(x)$ is relatively prime to $f(x)$.

Accordingly, by \S.\,6, Theorem II, we can determine the functions $Q(x)$, $\varphi(x)$ in $\Omega$, and indeed $\varphi(x)$ at most of degree $n - 1$, so that
\weq{4}{Q(x) f(x) + \varphi(x) \psi(x) = \chi(x)}
holds, and if we then set $x = \alpha$ herein, so that $f(\alpha) = 0$, it follows that
\weq{5}{\dfrac{\chi(\alpha)}{\psi(\alpha)} = \varphi(\alpha).}


% ===== scan leaf 482, printed 457 =====
% !PRINTED 457
\weberpage{457}
\pdfbookmark[2]{§ 143. Simultaneous adjunction of several algebraic quantities}{bm160}
If we denote the coefficients of $\varphi(x)$, which, as we have seen, are quantities in $\Omega$, by $c_0, c_1 \ldots c_{n-1}$, then every quantity $\Theta$ in $\Omega(\alpha)$ can accordingly be represented in the form:

\[
\Theta = c_0 + c_1\alpha + c_2\alpha^2 + \cdots + c_{n-1}\alpha^{n-1}.
\]

This representation is possible in only \gesperrt{one way}; for if at the same time

\weq{6}{\Theta = c_0' + c_1'\alpha + c_2'\alpha^2 + \cdots + c_{n-1}'\alpha^{n-1},}

then the function of degree $(n-1)$

\[
c_0 - c_0' + (c_1 - c_1')x + (c_2 - c_2')x^2 + \cdots + (c_{n-1} - c_{n-1}')x^{n-1}
\]

vanishes for $x = \alpha$. But by Theorem III of the preceding section this is possible only if

\[
c_0 = c_0', \quad c_1 = c_1', \ldots c_{n-1} = c_{n-1}'
\]

holds.

If we regard the quantities of the field $\Omega$ as the given ones, its quantities as rational, and the field $\Omega(\alpha)$ as the one derived from them, then we also call a rational function in $\Omega$ briefly a \gesperrt{rational function}, and can accordingly also say:

\gesperrt{The field $\Omega(\alpha)$ consists of all rational functions of $\alpha$.}

\begin{center}
\S.\,143.

\gesperrt{Simultaneous adjunction of several algebraic quantities.}
\end{center}

To a more exact investigation of the field $\Omega(\alpha)$ we prefix a general theorem, very simple but yet rich in consequences.

1. \gesperrt{If $\Phi_1(x,y,z\ldots)$, $\Phi_2(x,y,z\ldots)$, $\Phi_3(x,y,z\ldots)\ldots$ are integral rational functions of the variables $x,y,z\ldots$ with arbitrary coefficients, which do not all vanish at once in any of the functions, then one can substitute for the variables, in infinitely many ways, such rational numerical values that none of the functions $\Phi_1, \Phi_2, \Phi_3\ldots$ vanishes.}


% ===== scan leaf 483, printed 458 =====
% !PRINTED 458
\weberpage{458}
The theorem is first of all evident when the functions $\Phi_1, \Phi_2, \Phi_3\ldots$ depend on only one variable; for then there are altogether only a finite number of numerical values for this variable for which one of these functions becomes zero.

But then we can easily see the correctness of the theorem for functions of $n+1$ variables, provided we regard it as proved for functions of $n$ variables. For if we order the functions according to the $(n+1)$th variable $t$, then by hypothesis we can put for the remaining $n$ variables such rational values that in none of the functions do all the coefficients of the powers of $t$ vanish; we then have functions of the single variable $t$, and for it we can put such a rational value that none of the functions vanishes. The theorem is thus proved in general.

It shall now first be shown that the simultaneous adjunction of \gesperrt{several} algebraic quantities can be replaced by the adjunction of \gesperrt{a single one}; in other words, that every field $\Omega(\alpha, \beta, \gamma\ldots)$ can be regarded as a field $\Omega(\alpha)$.

Let then $\alpha, \beta, \gamma\ldots$ be algebraic quantities in arbitrary number, and let

\weq{1}{A(x) = 0, \quad B(x) = 0, \quad C(x) = 0\ldots}

be equations in $\Omega$ whose roots are $x=\alpha$, $x=\beta$, $x=\gamma\ldots$. None of these equations is to have a multiple root, an assumption by which the generality is not restricted.

We have seen in the preceding section that every quantity of a field $\Omega(\alpha)$ can be represented as an \gesperrt{integral rational function} in $\Omega$ of $\alpha$. If in this theorem we put in place of the field $\Omega$ the field $\Omega(\beta, \gamma\ldots)$, it follows that every number in $\Omega(\alpha, \beta, \gamma\ldots)$ can be represented as an integral rational function of $\alpha$ with coefficients from $\Omega(\beta, \gamma\ldots)$; and if we repeat the same mode of inference with respect to the coefficients, we obtain:

2. \gesperrt{Every quantity of the field $\Omega(\alpha, \beta, \gamma\ldots)$ can be represented as an integral rational function in $\Omega$ of $\alpha, \beta, \gamma\ldots$.}


% ===== scan leaf 484, printed 459 =====
% !PRINTED 459
\weberpage{459}

We form a linear function of the $\alpha, \beta, \gamma \ldots$
\[
\xi = x\alpha + y\beta + z\gamma + \cdots,
\]
and observe that each of the equations (1) has not merely one root, but several, whose number is equal to the degree of the equation. If then $\alpha', \beta', \gamma' \ldots$ is any combination, different from $\alpha, \beta, \gamma \ldots$, of one root from each of the equations (1), we set:
\[
\xi' = x\alpha' + y\beta' + z\gamma' + \ldots
\]
and form in the same way $\xi'', \xi''' \ldots$. The number of the quantities $\xi$ so formed is, if $a$ is the degree of $A(x)$, $b$ the degree of $B(x)$, $c$ the degree of $C(x)$, and so on,
\[
m = abc \ldots
\]
The differences $\xi - \xi'$, $\xi - \xi''$, $\xi' - \xi'' \ldots$ are linear functions of $x, y, z \ldots$, and none of them is identically zero, since we have assumed that none of the equations (1) has equal roots. By Theorem 1 we can therefore put for $x, y, z \ldots$ such \gesperrt{rational numbers} that none of these differences vanishes, so that the $m$ values $\xi, \xi', \xi'' \ldots$ are all different from one another.

Now every rational function that is symmetric with respect to the roots of each of the equations (1) is rationally expressible in terms of the coefficients of these equations, i.e. it is a number in $\Omega$. To these belong also the coefficients of the function of degree $m$ in $t$:
\weq{2}{F(t) = (t - \xi)(t - \xi')(t - \xi'') \ldots,}
and $F(t) = 0$ is therefore an equation in $\Omega$ that has no equal roots, and one of whose roots is $\xi$.

If, further, $\Theta$ is a quantity in $\Omega(\alpha, \beta, \gamma \ldots)$, and hence an integral rational function of $\alpha, \beta, \gamma \ldots$, and if we denote by $\Theta', \Theta'' \ldots$ the quantities that arise from this rational function when the variables are replaced by $\alpha', \beta', \gamma' \ldots$ or $\alpha'', \beta'', \gamma'' \ldots$ and so on, then
\[
F(t) \left( \frac{\Theta}{t - \xi} + \frac{\Theta'}{t - \xi'} + \frac{\Theta''}{t - \xi''} + \cdots \right) = \psi(t)
\]
is, as a symmetric function of the $\alpha$ and of the $\beta \ldots$, an integral rational function of degree $(m-1)$ in $t$ in $\Omega$, and if one sets $t = \xi$, it follows
%CONT>>


% ===== scan leaf 485, printed 460 =====
% !PRINTED 460
\weberpage{460}
\pdfbookmark[2]{§ 144. Primitive and imprimitive fields}{bm161}

\weq{3}{\Theta = \frac{\psi(\xi)}{F'(\xi)},}

and $\Theta$ is therefore contained in $\Omega(\xi)$. Since conversely every
quantity in $\Omega(\xi)$ is at the same time contained in $\Omega(\alpha, \beta, \gamma \dots)$, the
two fields are identical, i.\,e. we have

\[
\Omega(\xi) = \Omega(\alpha, \beta, \gamma \dots).
\]

But with this our assertion is proved.

\begin{center}
§. 144.

\gesperrt{Primitive and imprimitive fields.}
\end{center}

By the theorem last proved, we now restrict ourselves to
the consideration of algebraic fields $\Omega(\alpha)$, where $\alpha$
is the root of an equation in $\Omega$. Let this equation,
freed from multiple factors, be denoted by

\weq{1}{F(x) = 0}

and be of degree $m$. Its roots are as follows.

\weq{2}{\alpha, \alpha_1, \alpha_2 \dots \alpha_{m-1}.}

It may be that $F(x)$ is reducible in $\Omega$; then there will be one and
only one \gesperrt{irreducible factor} $f(x)$ of $F(x)$, so
that $\alpha$ is a root of the following equation.

\weq{3}{f(x) = 0}

Let the degree of $f(x)$ be $n$. The roots of (3), which are all
contained among the roots (2), are as follows.

\weq{4}{\alpha, \alpha_1, \alpha_2 \dots \alpha_{n-1}.}

The degree of the field $\Omega(\alpha)$ is then, by the stipulation
made in §.\,142, likewise $n$.

Each of the quantities (4) gives rise to an algebraic field,
and thus arise the fields

\weq{5}{\Omega(\alpha), \Omega(\alpha_1) \dots \Omega(\alpha_{n-1}),}

which we call the \gesperrt{fields conjugate to} $\Omega(\alpha)$. It may
be that these fields are all or in part identical, but they
may also all be different from one another, of which more
later.

By §.\,142 we obtain every quantity $\Theta$ in $\Omega(\alpha)$ if, in an
integral rational function $\chi(t)$ in $\Omega$, of degree at most
$(n-1)$, we put the number $\alpha$ for the variable $t$.


% ===== scan leaf 486, printed 461 =====
% !PRINTED 461
\weberpage{461}

If we substitute for $t$ in $\chi(t)$ the various quantities (4), we
obtain $n$ quantities, one from each of the conjugate fields
\weq{6}{\Theta = \chi(\alpha),\ \Theta_1 = \chi(\alpha_1) \ldots \Theta_{n-1} = \chi(\alpha_{n-1}),}
and these are called the \gesperrt{quantities conjugate to} $\Theta$.

Every symmetric function of these quantities is at the same time a
symmetric function of the roots of the equation (3), and hence
contained in $\Omega$.

Among these symmetric functions we wish to single out two, which
occur frequently, by special names and notations; there is the sum
\weq{7}{S(\Theta) = \Theta + \Theta_1 + \Theta_2 + \cdots + \Theta_{n-1},}
which we call the \gesperrt{trace} of $\Theta$, and the product
\weq{8}{N(\Theta) = \Theta \Theta_1 \Theta_2 \ldots \Theta_{n-1},}
which is called the \gesperrt{norm} of $\Theta$. Conjugate numbers accordingly
have the same traces and norms.

The product
\weq{9}{(t - \Theta)(t - \Theta_1) \ldots (t - \Theta_{n-1}) = \Phi(t)}
is an integral rational function of degree $n$ of $t$ in
$\Omega$, and its roots are $\Theta$ and the quantities conjugate to
$\Theta$. From this results the theorem:

\gesperrt{1. Every quantity in $\Omega(\alpha)$ is a root of an
equation of degree $n$ in $\Omega$, whose remaining roots
are the quantities conjugate to $\Theta$.}

The computation of the function $\Phi(t)$ from $f(x)$ is nothing
other than the \gesperrt{Tschirnhausen transformation} considered
from another point of view in the fourth and sixth sections.

We now have to examine the function $\Phi(t)$ as to its
irreducibility.

If the function $\Phi(t)$ is reducible, then it has an
irreducible factor $\varphi(t)$, in which we may assume the coefficient of the
highest power of $t$ equal to 1, and every such
factor vanishes for at least one of the values $\Theta,
\Theta_1 \ldots \Theta_{n-1}$.

Let us then have
\[
\varphi(\Theta) = \varphi[\chi(\alpha)] = 0,
\]
that is, the equations $\varphi[\chi(x)] = 0$ and $f(x) = 0$ have
a common root. But since $f(x)$ is assumed irreducible,
$\varphi[\chi(x)]$ must, by the Theorem II, \S.\,141, for all roots
%CONT>>


% ===== scan leaf 487, printed 462 =====
% !PRINTED 462
\weberpage{462}
%<<CONT
$\alpha, \alpha_1, \alpha_2 \ldots \alpha_{n-1}$ of $f(x) = 0$ vanish, i.\,e. we have
\[
\varphi(\Theta) = 0,\quad \varphi(\Theta_1) = 0,\quad \varphi(\Theta_2) = 0 \ldots \varphi(\Theta_{n-1}) = 0.
\]

If, then, the values conjugate to $\Theta$ are all different from
one another, $\varphi(t)$ is identical with $\Phi(t)$, i.\,e.
$\Phi(t)$ is irreducible. But if among the values conjugate to
$\Theta$ only $n_1$ different ones are present,
say $\Theta, \Theta_1, \Theta_2 \ldots \Theta_{n_1-1}$, then
\weq{10}{\varphi(t) = (t - \Theta)(t - \Theta_1) \ldots (t - \Theta_{n_1-1}),}
and every other irreducible factor of $\Phi(t)$, since it must
vanish for at least one of the conjugate values $\Theta$ and
consequently for all of them, is identical with $\varphi(t)$. Hence
$\Phi(t)$ is a power of $\varphi(t)$, say
\weq{11}{\Phi(t) = \varphi(t)^{n_2}}
and
\weq{12}{n = n_1 n_2.}
From this, then, follows the theorem:

\gesperrt{2. The function $\Phi(t)$ is either irreducible or it
is a power of an irreducible function. The $n$ numbers conjugate
to a number in $\Omega(\alpha)$ are either all different from
one another, or they fall into $n_1$ systems of
$n_2$ numbers each, equal among themselves. In the first case $\Phi(t)$
is irreducible, in the second it is the $n_2$th power of an
irreducible function of degree $n_1$ in $\Omega$.}

A quantity $\Theta$ in $\Omega(\alpha)$ that is different from all its
conjugate numbers, and that therefore satisfies an
irreducible equation of degree $n$, is called a
\gesperrt{primitive quantity of the field}. By the theorem \S.\,143, 1. infinitely
many such primitive quantities can be determined, indeed in such a way that
the coefficients of $\Theta = \chi(\alpha)$ are rational numbers.
One need only dispose of the coefficients of $\chi$ in such a way
that among the conjugate quantities $\chi(\alpha)$ no equal ones
occur.

\gesperrt{3. Every quantity $\omega$ of the field $\Omega(\alpha)$ can
be expressed rationally in terms of an arbitrary primitive quantity $\Theta$
of the field.}

For if $\omega, \omega_1, \omega_2 \ldots \omega_{n-1}$ are the numbers
conjugate to $\omega$, and likewise $\Theta, \Theta_1, \Theta_2
\ldots \Theta_{n-1}$ those conjugate to $\Theta$, and
%CONT>>


% ===== scan leaf 488, printed 463 =====
% !PRINTED 463
\weberpage{463}
%<<CONT
\[
\Phi(t) = (t - \Theta)(t - \Theta_1) \ldots (t - \Theta_{n-1}),
\]
then
\[
\Phi(t) \left( \frac{\omega}{t - \Theta} + \frac{\omega_1}{t - \Theta_1} + \cdots + \frac{\omega_{n-1}}{t - \Theta_{n-1}} \right) = \Psi(t)
\]
is an integral rational function of degree $n - 1$ of $t$,
whose coefficients are numbers in $\Omega$, and from this it
follows, on setting $t = \Theta$,
\[
\omega = \frac{\Psi(\Theta)}{\Phi'(\Theta)},
\]
in which $\Phi'(\Theta)$ is different from zero.

\gesperrt{Accordingly the field $\Omega(\Theta)$ is identical with the field
$\Omega(\alpha)$.}

\gesperrt{4. If $\Theta$ is not primitive, then not every quantity
in $\Omega(\alpha)$ can be expressed rationally in terms of $\Theta$. The
field $\Omega(\Theta)$ is a divisor of the field $\Omega(\alpha)$,
and the degree of $\Omega(\Theta)$ is a divisor of the degree of
$\Omega(\alpha)$.}

For every number of the field $\Omega(\Theta)$ satisfies an equation
in $\Omega$ of degree $n_1$, where $n_1$ is a divisor of $n$ and
less than $n$. Hence a primitive quantity of the field
$\Omega(\alpha)$, which satisfies an irreducible equation of degree
$n$, cannot be contained in $\Omega(\Theta)$.

\gesperrt{The field $\Omega(\alpha)$ is called primitive if, apart from the
quantities in $\Omega$, it contains no imprimitive quantities; imprimitive,
if it contains still other imprimitive quantities.}

From this definition it follows first of all that a field whose
degree is a prime is necessarily primitive; for an imprimitive
quantity $\Theta$ in $\Omega(\alpha)$ satisfies an equation whose
degree is a divisor of the degree $n$ of $\Omega(\alpha)$ different
from $n$. If this divisor is equal to 1, then $\Theta$ is contained
in $\Omega$.

We shall now become acquainted with some of the most important
properties of the imprimitive fields.

If $\Omega(\alpha)$ has a divisor $\Omega'$ different from
$\Omega$, which in its turn contains $\Omega$ as a divisor (so that
$\Omega'$ is a field over $\Omega$), and if $\beta$ is a quantity that
belongs to the field $\Omega'$ but not to $\Omega$, then the
field $\Omega(\beta)$ is an algebraic field over $\Omega$ and
at the same time a divisor of $\Omega'$ and of $\Omega(\alpha)$, and
according to what we have just
%CONT>>


% ===== scan leaf 489, printed 464 =====
% !PRINTED 464
\weberpage{464}
\pdfbookmark[2]{§ 145. Normal fields. Galois resolvent}{bm162}
%<<CONT
proved, the degree $n_1$ of the field $\Omega(\beta)$ is a
divisor of $n$. If now $\Omega'$ is not exhausted by $\Omega(\beta)$,
we take a quantity $\gamma$ in $\Omega'$, but not in
$\Omega(\beta)$; then the field $\Omega(\beta,
\gamma)$ is a divisor of $\Omega'$ and of $\Omega(\alpha)$ and has,
for its part, $\Omega(\beta)$ as a divisor. The degree of $\Omega(\beta,
\gamma)$ is therefore greater than that of $\Omega(\beta)$ and smaller
than that of $\Omega(\alpha)$. If $\Omega'$ is not yet exhausted
by this, we continue in the same way, but must finally come to an
end, since the degrees of the fields $\Omega(\beta), \Omega(\beta,
\gamma) \ldots$ always increase and yet remain smaller than $n$.
From this it follows:

\gesperrt{5. Every divisor of $\Omega(\alpha)$ which contains the
field $\Omega$ is an algebraic field $\Omega(\beta)$ over
$\Omega$. The degree of $\Omega(\beta)$ is a divisor of the degree
of $\Omega(\alpha)$, and if both fields have the same degree,
they are identical.}

We can therefore also state our definition thus:

\gesperrt{6. An algebraic field over $\Omega$ is primitive
if it has no field over $\Omega$ except $\Omega$ and itself
as a divisor.}

\begin{center}
\S.\,145.

\gesperrt{Normal fields. Galois resolvent.}
\end{center}

If $\Omega(\alpha)$ is an algebraic field, then the
conjugate fields
% [sic] The source prints alpha_{m-1}; §144 defines the conjugate fields only through alpha_{n-1}, where n is the irreducible-factor degree. Preserved diplomatically.
\[
\Omega(\alpha), \Omega(\alpha_1) \ldots \Omega(\alpha_{m-1})
\]
are all of the same degree, and hence if one of them is contained
in another, the two are identical.

A field which is identical with all its conjugate fields
is called a \gesperrt{normal field}. In the normal fields
much simpler laws prevail, and the great advance which
algebra owes to \gesperrt{Galois} rests essentially
on the fact that arbitrary fields are reduced to normal
fields. The normal fields are therefore also called \gesperrt{Galois
fields}.

If a field $\Omega(\varrho)$ of degree $\mu$ has the
property that the numbers $\varrho_1,
\varrho_2 \ldots \varrho_{\mu-1}$ conjugate to $\varrho$ are all
contained in $\Omega(\varrho)$, then it is a normal field; for then the
fields $\Omega(\varrho_1), \Omega(\varrho_2) \ldots
\Omega(\varrho_{\mu-1})$ are also all contained in $\Omega(\varrho)$
and hence all identical with $\Omega(\varrho)$.


% ===== scan leaf 490, printed 465 =====
% !PRINTED 465
\weberpage{465}

We shall call an equation a \gesperrt{normal equation if it is irreducible and has the property that all its roots can be expressed rationally (in $\Omega$) by one of them}. Then a primitive element of a normal field of degree $\mu$ is a root of a normal equation of degree $\mu$. But the converse also holds, namely, that if $\varrho$ is a root of a normal equation, then $\Omega(\varrho)$ is a normal field of the same degree. For if $\varrho_0$ is the root of the normal equation by which all the others are rationally expressed, then $\Omega(\varrho_0)$ is certainly a normal field. But $\Omega(\varrho)$ is conjugate to this field, and hence identical with it. From this it follows further that in a normal equation every root is rationally expressible not only by one determinate root among them, \gesperrt{but by any one whatever}.

One can now reduce arbitrary algebraic fields to normal fields in the following way.

Let
\weq{1}{F(x) = 0}
be an arbitrary reducible or irreducible equation in $\Omega$ of degree $m$, of which we shall only assume \gesperrt{that it has no multiple roots}. Its roots are as follows.
\weq{2}{\alpha,\ \alpha_1 \ldots \alpha_{m-1}.}
From these one obtains $m$ algebraic fields
\weq{3}{\Omega(\alpha),\ \Omega(\alpha_1) \ldots \Omega(\alpha_{m-1}),}
and one can evidently choose the function $F(x)$ so that among the fields (3) any given algebraic fields over $\Omega$ occur. We call the field over $\Omega$ derived from all the quantities (2), i.\,e. from all the roots of the equation (1),
\weq{4}{N = \Omega(\alpha, \alpha_1 \ldots \alpha_{m-1})}
the \gesperrt{Galois field of the equation} $F(x) = 0$.

If $F(x)$ is irreducible, the fields (3) are the fields conjugate to $\Omega(\alpha)$. In this case the field $N$ shall be called the \gesperrt{norm of each of the fields} $\Omega(\alpha), \Omega(\alpha_1) \ldots \Omega(\alpha_{m-1})$.

If $\Omega(\alpha)$ is a normal field, it is identical with its norm. In the general case it is to be shown \gesperrt{that $N$ is a normal field}. We choose a primitive element $\varrho$ of the field $N$ (according to \S.\,143), and can then also denote the field $N$
%CONT>>


% ===== scan leaf 491, printed 466 =====
% !PRINTED 466
\weberpage{466}

%<<CONT
denote by $\Omega(\varrho)$. If $\mu$ is the degree of $N$, then $\varrho$ satisfies an irreducible equation of degree $\mu$
\weq{5}{g(t) = 0,}
of which it is to be shown that it is a normal equation. For this purpose we first remark that the one root $\varrho$ of this equation is a rational function of $\alpha, \alpha_1 \ldots \alpha_{m-1}$, because it was contained in $N$. To indicate this, let us set
\[
\varrho = \varrho(\alpha, \alpha_1 \ldots \alpha_{m-1}),
\]
and now form all the different arrangements of the digits
\[
0,\ 1,\ 2 \ldots m - 1,
\]
whose number amounts to $\Pi(m)$:
\weq{6}{(0, 1, 2 \ldots m - 1),\ (0', 1', 2' \ldots (m-1)'),\ (0'', 1'', 2'' \ldots (m-1)'')\,\ldots,}
in which the digits with one, two, etc. accents are the same as those without accent, only in a different order, and from these form the functions
\weq{7}{\varrho = \varrho(\alpha_0, \alpha_1 \ldots \alpha_{m-1}),\ \varrho' = \varrho(\alpha_{0'}, \alpha_{1'} \ldots \alpha_{(m-1)'}),\ \varrho'' = \varrho(\alpha_{0''}, \alpha_{1''} \ldots \alpha_{(m-1)''})\,\ldots,}
without concern as to whether some among them may be equal to one another or not.

If in all the arrangements (6) we carry out one and the same permutation, e.g. 0 with 1, the totality of these arrangements does not change; only their order becomes a different one. For, first, such a permutation can produce nothing other than arrangements of the digits, and secondly, two different arrangements cannot pass into the same arrangement by one and the same permutation.

Hence by every such permutation the functions (7) will likewise only be permuted among one another. If, then, we form the product
\[
G(t) = (t - \varrho)\,(t - \varrho')\,(t - \varrho'')\,\ldots
\]
for a variable $t$, its coefficients, which are certainly functions of $\alpha, \alpha_1 \ldots \alpha_{m-1}$, remain unchanged when these quantities are permuted in any way; that is, they are symmetric functions of the roots of the function $F(x)$, and hence $G(t)$ is a function of $t$ in $\Omega$. All the roots of $G(t)$ are quantities in $N$, since they are rationally expressed by the $\alpha$.


% ===== scan leaf 492, printed 467 =====
% !PRINTED 467
\weberpage{467}
\pdfbookmark[2]{§ 146. The substitutions of a normal field}{bm163}

Now $G(t)$ and $g(t)$ have a root in common. But since $g(t)$ is irreducible, $G(t)$ must be divisible by $g(t)$; hence all the roots of $g(t)$ are also contained in $N$, i.\,e. $N$ is a normal field, q.\,e.\,d.

Every equation $g(t) = 0$ is called a \gesperrt{Galois resolvent} of the equation $F(x) = 0$, and a \gesperrt{Galois resolvent} is accordingly defined by the following specification:

\gesperrt{An equation $g(t) = 0$ is a Galois resolvent of a given equation $F(x) = 0$ in $\Omega$ if 1) $g(t)$ is irreducible, if 2) all the roots of $F(x)$ are rationally expressible through a root $\varrho$ of $g(t)$, and 3) one root of $g(t)$ is rationally expressible through the roots of $F(x)$.}

For by 2), $N$ is contained in $\Omega(\varrho)$, and by 3) one of the fields conjugate to $\Omega(\varrho)$, namely $\Omega(\varrho_1)$, is contained in $N$. The degrees of $N$ and $\Omega(\varrho)$ therefore cannot be different, and consequently $N = \Omega(\varrho)$.

\gesperrt{Every Galois resolvent is a normal equation, and a normal equation is its own Galois resolvent.}

\begin{center}
\S.\,146.

\gesperrt{The Substitutions of a Normal Field.}
\end{center}

Now let $\Omega(\varrho)$ be a normal field of degree $\mu$ and $\varrho$ one of its primitive numbers, and further let
\weq{1}{g(t) = 0}
be the irreducible equation of degree $\mu$ of which one root is $t = \varrho$. The elements conjugate to $\varrho$ are as follows.
\weq{2}{\varrho,\ \varrho_1,\ \varrho_2 \ldots \varrho_{\mu-1}.}
Since by the definition of the normal field the quantities (2) are all contained in $\Omega(\varrho)$, we can determine $\Theta_1(t), \Theta_2(t) \ldots \Theta_{\mu-1}(t)$ as integral rational functions in $\Omega$, of degree at most $\mu - 1$, such that
\weq{3}{\varrho_1 = \Theta_1(\varrho),\ \varrho_2 = \Theta_2(\varrho) \ldots \varrho_{\mu-1} = \Theta_{\mu-1}(\varrho)}
%CONT>>


% ===== scan leaf 493, printed 468 =====
% !PRINTED 468
\weberpage{468}

%<<CONT
becomes. If $\omega$ is an arbitrary quantity in $\Omega(\varrho)$, then the quantities conjugate to $\omega$ can be represented thus:
\weq{4}{\omega = \varphi(\varrho),\ \omega_1 = \varphi(\varrho_1) \ldots \omega_{\mu-1} = \varphi(\varrho_{\mu-1}),}
in which $\varphi(t)$ is a rational function in $\Omega$.

Now since $g(t)$ is irreducible, the theorem of \S.\,141, II holds, which we now state as follows:

1. \gesperrt{If a rational function $\Phi(t)$ in $\Omega$ has a root in common with $g(t)$, then there vanish all the conjugate quantities}
\[
\Phi(\varrho),\ \Phi(\varrho_1) \ldots \Phi(\varrho_{\mu-1}).
\]

If in one of the functions $\Theta_k(\varrho)$, by which according to (3) the roots of (1) are expressed, $\varrho$ is replaced by another root $\varrho_h$, then there arises from it again one of the roots; for if
\[
g[\Theta_k(\varrho)] = 0,
\]
then by 1. also
\[
g[\Theta_k(\varrho_h)] = 0,
\]
and the sequence of quantities
\weq{5}{\varrho_h,\ \Theta_1(\varrho_h),\ \Theta_2(\varrho_h) \ldots \Theta_{\mu-1}(\varrho_h)}
agrees, apart from the arrangement, with the following sequence.
\weq{6}{\varrho,\ \Theta_1(\varrho),\ \Theta_2(\varrho) \ldots \Theta_{\mu-1}(\varrho)}
This will be proved when we show that no two equal values occur in (5). If, for the sake of uniformity, we denote by $\Theta_0(\varrho)$ or $\Theta(\varrho)$ the root $\varrho$ itself, then from the equality of two of the quantities (5)
\weq{7}{\Theta_i(\varrho_h) = \Theta_k(\varrho_h)}
Theorem 1 gives
\weq{8}{\Theta_i(\varrho) = \Theta_k(\varrho),}
which, however, if $i$ is different from $k$, is not possible. Hence, together with the quantities (6), the quantities (5) are also different from one another.

We can sum this up as a theorem:

2. \gesperrt{If one interchanges $\varrho$ with $\varrho_h$, then at the same time every quantity conjugate to $\varrho$ passes over into a definite other one, and never two different ones into the same.}

If in all functions of $\varrho$ we put another root $\varrho_a$ in place of $\varrho$, we carry out a \gesperrt{substitution}. We denote this substitution by
%CONT>>


% ===== scan leaf 494, printed 469 =====
% !PRINTED 469
\weberpage{469}

\[
\sigma_a = (\varrho, \varrho_a), \qquad a = 0, 1, 2 \ldots \mu - 1,
\]
where by $\sigma_0$ or $\sigma$ is understood the so-called \gesperrt{identical substitution} $(\varrho, \varrho)$, which consists in $\varrho$ being replaced by itself, so that all numbers in $\Omega(\varrho)$ remain unchanged.

If in any one of the roots (3)
\[
\varrho_h = \Theta_h(\varrho)
\]
we carry out the substitution $\sigma_a$, then $\varrho_h$ passes into another root $\varrho_k$, which is determined by
\[
\varrho_k = \Theta_h(\varrho_a) = \Theta_h \Theta_a(\varrho) = \Theta_k(\varrho).
\]
Thus $\varrho_k$ is the same function of $\varrho_a$ as $\varrho_h$ is of $\varrho$.

An arbitrary quantity $\omega = \varphi(\varrho)$ of the field $\Omega(\varrho)$ passes, by the substitution $\sigma_a$, into $\omega_a = \varphi(\varrho_a)$. If we express $\omega$ in terms of $\varrho_h$ and $\omega_a$ in terms of $\varrho_k$:
\[
\omega = \psi(\varrho_h), \qquad \omega_a = \psi(\varrho_k),
\]
then one sees that $\omega$ passes into $\omega_a$ by the substitution
\weq{9}{\sigma_a = (\varrho_h, \varrho_k) = (\varrho_h, \Theta_h \Theta_a(\varrho)),}
and this is only another expression for the substitution $(\varrho, \varrho_a)$.

Here $\varrho_h$ is an arbitrary root of $g(t)$, and the associated $\varrho_k$ is determined by $\sigma_a$. But for a given $\sigma_a$ one may also take $\varrho_k$ arbitrarily and determine the associated $\varrho_h$, by letting the index $h$ in $\Theta_h(\varrho_a)$ run through the values $0, 1, 2 \ldots \mu - 1$, whereby each root $\varrho_k$ makes its appearance once. Hence:

\gesperrt{3. Every substitution $\sigma_a$ can be represented in the form $(\varrho_h, \varrho_k)$, in which either $\varrho_h$ or $\varrho_k$ is an arbitrarily given one of the $\mu$ roots $\varrho$, while the other is completely determined by $\sigma_a$.}

% [sic] The German source breaks off grammatically after „in die gleiche“; the implied completion is „Grösse übergehen“ (pass into the same quantity). The English supplies the implied noun while the German transcription remains diplomatic.
The number of all mutually distinct substitutions $\sigma_a$ is therefore, the identical substitution included, equal to the degree of the field $\Omega(\varrho)$, which we have already denoted above by $\mu$. Each of these substitutions carries the totality of the quantities of the field $\Omega(\varrho)$ into itself, in such a way that each passes into a definite other one, and never two distinct quantities into the same one [formulae (7), (8)].

We therefore call the $\sigma_a$ the \gesperrt{substitutions of the field} $\Omega(\varrho)$.

If $\omega = \varphi(\varrho)$ remains unchanged in the field $\Omega(\varrho)$ when $\varrho$ is replaced by $\varrho_a$, that is, if
%CONT>>


% ===== scan leaf 495, printed 470 =====
% !PRINTED 470
\weberpage{470}
\pdfbookmark[2]{§ 147. Composition of substitutions}{bm164}

%<<CONT
\[
\varphi(\varrho) = \varphi(\varrho_a)
\]
then we say that \gesperrt{$\omega$ permits or admits the substitution $(\varrho, \varrho_a)$ or $\sigma_a$}.

The quantities $\omega$ which, apart from the identical substitution, admit no substitution $\sigma_a$, are the primitive elements of the field $\Omega(\varrho)$.

\gesperrt{4. A quantity which admits all the substitutions $\sigma_a$ is necessarily an element of $\Omega$.}

For such a quantity is identical with all its conjugate quantities, and hence by \S.\,144 satisfies an equation of the first degree in $\Omega$.

\begin{center}
\S.\ 147.
\end{center}

\begin{center}
\gesperrt{Composition of the Substitutions.}
\end{center}

If in any function of $\varrho$ we replace the root $\varrho$ first by $\varrho_a$ and then $\varrho_a$ by $\varrho_b$, the result is the same as if we had from the outset interchanged $\varrho$ with $\varrho_b$. If therefore we set
\weq{1}{\sigma = (\varrho, \varrho_a), \quad \sigma' = (\varrho_a, \varrho_b), \quad \sigma'' = (\varrho, \varrho_b),}
then it is indifferent for the result whether in all the numbers of the field $\Omega(\varrho)$ we first carry out the substitution $\sigma$ and then the substitution $\sigma'$, or whether we make the substitution $\sigma''$ in a single step.

\gesperrt{1. We therefore call $\sigma''$ compounded or composed of $\sigma$ and $\sigma'$, and denote this relation by the symbolic equation}
\weq{2}{\sigma \sigma' = \sigma''.}
Since by Theorem 3.\ of the preceding section we may, in the designation of a substitution $(\varrho_h, \varrho_k)$, choose the first or the second element arbitrarily, $\varrho_a, \varrho_b$ can be so selected that $\sigma, \sigma'$ are any two arbitrarily given ones among the $\mu$ substitutions of the field $\Omega(\varrho)$; but $\sigma''$ is thereby completely determined. But likewise, if $\sigma$ and $\sigma''$ are given, $\sigma'$ is uniquely determined. For if we choose $\varrho$ arbitrarily, then $\varrho_a$ is completely determined by $\sigma$, and $\varrho_b$ by $\sigma''$, whereby $\sigma' = (\varrho_a, \varrho_b)$ is also given. Finally, if $\sigma'$ and $\sigma''$ are given, then $\sigma$ is uniquely determined, since first $\varrho_b$ is determined by $\sigma''$ and then $\varrho_a$ by $\sigma'$. We therefore have:
%CONT>>


% ===== scan leaf 496, printed 471 =====
% !PRINTED 471
\weberpage{471}
%<<CONT

\gesperrt{2. Of the three substitutions of the field $\Omega(\varrho)$ connected by the symbolic equation $\sigma \sigma' = \sigma''$, any two may be given arbitrarily, while the third is thereby uniquely determined.}

We have expressed the composition of the substitutions symbolically by the sign of multiplication. Here, however, one must pay careful attention to the order of the factors or components, because $\sigma \sigma'$ can be different from $\sigma' \sigma$; that is, for the symbolic multiplication expressing composition, the commutative law of ordinary multiplication does not hold; but the so-called associative law does hold, which is expressed in the following theorem:

\gesperrt{3. If $\sigma, \sigma', \sigma''$ are any three of the $\mu$ substitutions of the field $\Omega(\varrho)$, then}
\weq{3}{(\sigma \sigma') \sigma'' = \sigma (\sigma' \sigma'').}
The proof is very simple; for by Theorem 3., \S.\ 146, we can determine $\varrho_a, \varrho_b, \varrho_c$ so that
\weq{4}{\sigma = (\varrho, \varrho_a), \quad \sigma' = (\varrho_a, \varrho_b), \quad \sigma'' = (\varrho_b, \varrho_c),}
and then
\[
\sigma \sigma' = (\varrho, \varrho_b), \qquad (\sigma \sigma') \sigma'' = (\varrho, \varrho_b)(\varrho_b, \varrho_c) = (\varrho, \varrho_c),
\]
\[
\sigma' \sigma'' = (\varrho_a, \varrho_c), \qquad \sigma (\sigma' \sigma'') = (\varrho, \varrho_a)(\varrho_a, \varrho_c) = (\varrho, \varrho_c).
\]
We therefore denote briefly the substitution composed of the three substitutions $\sigma, \sigma', \sigma''$ by
\[
\sigma \sigma' \sigma'',
\]
and can in general compose from arbitrarily many components a definite substitution
\[
\sigma \sigma' \sigma'' \sigma''' \ldots
\]
by uniting, one after another, two adjacent components at a time into a single substitution, until the whole symbolic product has contracted to a single substitution. For by \S.\ 146, 3.\ we can set
\[
\sigma = (\varrho, \varrho_a), \quad \sigma' = (\varrho_a, \varrho_b), \quad \sigma'' = (\varrho_b, \varrho_c), \quad \sigma''' = (\varrho_c, \varrho_e),
\]
and always obtain
\[
(\varrho, \varrho_a)(\varrho_a, \varrho_b)(\varrho_b, \varrho_c)(\varrho_c, \varrho_e) = (\varrho, \varrho_e),
\]
which evidently can be extended to an arbitrary number of components.


% ===== scan leaf 497, printed 472 =====
% !PRINTED 472
\weberpage{472}
\pdfbookmark[2]{§ 148. Permutation groups}{bm165}

On the basis of the properties 1., 2., 3.\ we call the totality of the $\mu$ substitutions of the normal field $\Omega(\varrho)$ a \gesperrt{group}, and indeed a \gesperrt{finite group} of \gesperrt{degree} $\mu$. It is also called the \gesperrt{group of the substitutions of the field} $\Omega(\varrho)$.

With the general theory of groups we shall occupy ourselves in detail in later sections.

\begin{center}
\S.\ 148.
\end{center}

\begin{center}
\gesperrt{Permutation groups.}
\end{center}

In \S.\ 145 we have seen how, for an arbitrary equation in $\Omega$ of degree $m$
\weq{1}{F(x) = 0,}
of which it is only assumed that it has no equal roots, one can determine a normal field $\Omega(\varrho)$ to which all the roots of (1) belong. If, then, the roots of (1) are denoted by
\weq{2}{\alpha, \alpha_1, \alpha_2 \ldots \alpha_{m-1}}
then
\weq{3}{\alpha = \chi(\varrho), \quad \alpha_1 = \chi_1(\varrho), \ldots \alpha_{m-1} = \chi_{m-1}(\varrho),}
and the $\chi$ are rational functions in $\Omega$.

But if we set in place of $\varrho$ a quantity $\varrho_1$ conjugate to $\varrho$, that is, if we carry out a substitution $\sigma = (\varrho, \varrho_1)$, then every root $\alpha$ passes over into a definite other one, and never do two different ones pass into the same; for if
\[
F(\alpha) = F[\chi(\varrho)] = 0
\]
then by \S.\ 146, 1:
\[
F[\chi(\varrho_1)] = 0,
\]
that is, $\chi(\varrho_1)$ is also a root of the equation (1). But that two different quantities of the field $\Omega(\varrho)$ cannot pass over into the same quantity through a substitution $\sigma$ we have already shown in general in \S.\ 146.

\gesperrt{1. A substitution $\sigma$ therefore calls forth a certain permutation among the roots $\alpha$, in that each of the indices $0, 1, 2 \ldots m-1$ is replaced by a definite other one.}


% ===== scan leaf 498, printed 473 =====
% !PRINTED 473
\weberpage{473}

This leads us, then, to study more closely the permutations of $m$ figures
\[
0, 1, 2 \ldots m-1
\]
and their properties in general.

If the figure 0 passes into $a_0$, 1 into $a_1$, 2 into $a_2$, and so on, $m-1$ into
$a_{m-1}$, then

\weq{4}{a_0, a_1, a_2 \ldots a_{m-1}}

must be some other arrangement of the figures $0, 1, 2 \ldots m-1$. The
transition from the one arrangement to the other we denote by
a symbol, such as

\weq{5}{\pi_a = \begin{pmatrix} 0, & 1, & 2 & \ldots & m-1 \\ a_0, & a_1, & a_2 & \ldots & a_{m-1} \end{pmatrix}}

and we call this transition a \gesperrt{permutation}.

The number of all possible permutations of $m$ elements is
\[
\Pi(m) = 1 \,.\, 2 \,.\, 3 \ldots m,
\]
among which the \gesperrt{identical permutation}
\[
\begin{pmatrix} 0, & 1, & 2 & \ldots & m-1 \\ 0, & 1, & 2 & \ldots & m-1 \end{pmatrix},
\]
which leaves every element in its place, is also counted.

The symbol (5) does not change its meaning if one arranges the individual
pairs
\[
\begin{array}{ccccc}
0, & 1, & 2, & \ldots & m-1 \\
a_0, & a_1, & a_2, & \ldots & a_{m-1}
\end{array}
\]
differently. If then
\[
b_0, b_1, b_2 \ldots b_{m-1}
\]
is any arrangement of the figures $0, 1, 2 \ldots m-1$, we may
also put, in place of (5):

\weq{6}{\pi_a = \begin{pmatrix} b_0, & b_1, & b_2 & \ldots & b_{m-1} \\ a_{b_0}, & a_{b_1}, & a_{b_2} & \ldots & a_{b_{m-1}} \end{pmatrix},}

which is entirely equivalent to (5).

From this the very important concept of the \gesperrt{composition of
permutations} can be derived. If we replace an index $a$ by an index
$b$, and then $b$ by $c$, the effect is the same as if $a$ were replaced
at once by $c$. If, then, we carry out two permutations one after
the other, the result is a definite permutation,
which we can represent in the following way. If
%CONT>>


% ===== scan leaf 499, printed 474 =====
% !PRINTED 474
\weberpage{474}

%<<CONT
\[
\pi_a = \begin{pmatrix} 0, & 1 & \ldots & m-1 \\ a_0, & a_1 & \ldots & a_{m-1} \end{pmatrix}, \quad
\pi_b = \begin{pmatrix} 0, & 1 & \ldots & m-1 \\ b_0, & b_1 & \ldots & b_{m-1} \end{pmatrix}
\]
are any two permutations, then by (6) we can also write the second
thus:
\[
\pi_b = \begin{pmatrix} a_0, & a_1 & \ldots & a_{m-1} \\ b_{a_0}, & b_{a_1} & \ldots & b_{a_{m-1}} \end{pmatrix},
\]
and if we therefore carry out both one after the other, first $\pi_a$,
then $\pi_b$, there arises a third permutation, which we call composed
(componirt) from $\pi_a$ and $\pi_b$, and which we represent as the
symbolic product of $\pi_a$ and $\pi_b$:

\weq{7}{\pi_a \pi_b = \begin{pmatrix} 0, & 1 & \ldots & m-1 \\ b_{a_0}, & b_{a_1} & \ldots & b_{a_{m-1}} \end{pmatrix}.}

In this composition the factors usually cannot be interchanged,
since $\pi_a \pi_b$ does not in general agree with
$\pi_b \pi_a$. The following symbolic equation expresses the associative
principle.

\weq{8}{(\pi_a \pi_b)\pi_c = \pi_a(\pi_b \pi_c) = \pi_a \pi_b \pi_c}

One convinces oneself of this by the following simple
consideration.

If an index $p$ is changed by $\pi_a$ into $q$, the resulting index is changed
by $\pi_b$ into $r$, and $r$ by $\pi_c$ into $s$, then $p$ passes by
$\pi_a \pi_b$ into $r$, and $q$ by $\pi_b \pi_c$ into $s$; hence
$p$ is changed into $s$ both by $(\pi_a \pi_b)\pi_c$ and by
$\pi_a(\pi_b \pi_c)$. But formula (8) says nothing else.

To every permutation $\pi_a$ there is a definite
\gesperrt{opposite} (inverse) permutation, which consists in this,
that the change produced by $\pi_a$ is undone again, and which
therefore, composed with $\pi_a$, produces the identical
permutation. We denote this permutation by
$\pi_a^{-1}$, and we have, if $\pi_a$ has the expression (5),

\weq{9}{\pi_a^{-1} = \begin{pmatrix} a_0, & a_1 & \ldots & a_{m-1} \\ 0, & 1 & \ldots & m-1 \end{pmatrix}.}

For by the rule of composition in (7) we have
\[
\pi_a \pi_a^{-1} = \begin{pmatrix} 0, & 1 & \ldots & m-1 \\ 0, & 1 & \ldots & m-1 \end{pmatrix}.
\]
But in the same way we also have
\[
\pi_a^{-1}\pi_a = \begin{pmatrix} a_0, & a_1 & \ldots & a_{m-1} \\ a_0, & a_1 & \ldots & a_{m-1} \end{pmatrix},
\]
%CONT>>


% ===== scan leaf 500, printed 475 =====
% !PRINTED 475
\weberpage{475}

%<<CONT
which is again the identical permutation. \gesperrt{Thus the permutation
opposite to $\pi_a^{-1}$ is again the original permutation $\pi_a$.}

If $\pi_a \pi_b = \pi_c$, then
\[
\pi_c^{-1} = \pi_b^{-1} \pi_a^{-1},
\]
for $\pi_a \pi_b \pi_b^{-1} \pi_a^{-1}$ is the identical
permutation.

If we denote the identical permutation by $\pi_0$, then

\weq{10}{\pi_a \pi_a^{-1} = \pi_a^{-1} \pi_a = \pi_0,}

and $\pi_0$ is opposite to itself (which does not exclude the
possibility that there are still other permutations opposite to
themselves). By composition with the identical permutation no other
permutation is changed; thus

\weq{11}{\pi_0 \pi_a = \pi_a \pi_0 = \pi_a.}

\gesperrt{2. If $\pi_c$ is composed of $\pi_a$ and $\pi_b$, that is,}

\weq{12}{\pi_a \pi_b = \pi_c,}

\gesperrt{then not only is $\pi_c$ completely and uniquely determined
by $\pi_a$ and $\pi_b$, but also $\pi_a$ by $\pi_b$ and
$\pi_c$, and $\pi_b$ by $\pi_a$ and $\pi_c$.}

For it follows from (12) that

\weq{13}{\pi_a = \pi_c \pi_b^{-1}, \quad \pi_b = \pi_a^{-1} \pi_c.}

\gesperrt{3. A system
\[
Q = \pi, \pi', \pi'', \pi''' \ldots,
\]
selected from the totality of all $\Pi(m)$ permutations of $m$
elements, which is so closed and complete in itself that any two, and
consequently also arbitrarily many, permutations of $Q$ always yield
by composition a permutation of the same system, is called a
permutation group. By the degree of a permutation group we understand
the number of permutations that it contains.}

The totality of all $\Pi(m)$ permutations forms such a group.
Likewise the single identical permutation is a group by itself. What
may still lie between these two extreme cases of groups is what will
occupy us in the following sections.


% ===== scan leaf 501, printed 476 =====
% !PRINTED 476
\weberpage{476}
\pdfbookmark[2]{§ 149. Galois group}{bm166}

As a simple example we mention here only the \gesperrt{group of cyclic
permutations, briefly the cyclic group}, which one obtains when, in any
arrangement of digits $(1, 2, 3 \ldots n)$, one replaces each digit by
the following one and the last by the first, and when one repeats this
permutation until one comes back to the original arrangement. The
degree of the cyclic group is equal to the number of digits, e.\,g.
with three digits:
\[
\begin{pmatrix} 1, & 2, & 3 \\ 2, & 3, & 1 \end{pmatrix}, \quad
\begin{pmatrix} 1, & 2, & 3 \\ 3, & 1, & 2 \end{pmatrix}, \quad
\begin{pmatrix} 1, & 2, & 3 \\ 1, & 2, & 3 \end{pmatrix}.
\]

There are special groups, which will occupy us still further later
on, that consist only of such permutations in whose composition,
besides the laws mentioned above, the commutative law also holds, for
which therefore always $\pi_1 \pi_2 = \pi_2 \pi_1$, and in consequence
of which the analogy with multiplication is a much more complete one.
Such groups, to which the cyclic groups belong, are called
\gesperrt{commutative} or \gesperrt{Abelian groups}. When all the
elements of a permutation group $Q_1$ are contained in a group $Q$,
then $Q_1$ is called a \gesperrt{divisor of $Q$, and $Q$ divisible by
$Q_1$}.

\begin{center}
\S.\,149.

\gesperrt{Galois Group.}
\end{center}

If, in the roots $\alpha, \alpha_1, \alpha_2 \ldots \alpha_{m-1}$ of
an equation $F(x) = 0$, as they are expressed by \S.\,148, (3) as
quantities of the normal field $\Omega(\varrho)$ of degree $\mu$, we
carry out one of the $\mu$ substitutions $\sigma$ of the normal field,
then a certain permutation takes place among these roots, as has
already been shown above. To every substitution $\sigma$ there
corresponds therefore a certain permutation $\pi$ of the $m$ digits
$0, 1, 2 \ldots m-1$. \gesperrt{Two different substitutions $\sigma_h, \sigma_k$
lead also to two different permutations $\pi_h, \pi_k$.}

For by \S.\,145, $\varrho$ is a rational function of $\alpha,
\alpha_1 \ldots \alpha_{m-1}$; we use the following notation.
\[
\varrho = \Phi(\alpha, \alpha_1 \ldots \alpha_{m-1})
\]
If we here set
\[
\alpha_s = \chi_s(\varrho)
\]
it follows that

\weq{1}{\varrho = \Phi[\chi(\varrho), \chi_1(\varrho) \ldots \chi_{m-1}(\varrho)],}
%CONT>>


% ===== scan leaf 502, printed 477 =====
% !PRINTED 477
\weberpage{477}

%<<CONT
and by \S.\,146, 1. we thus have, for every index $k$,
\weq{2}{\varrho_k = \Phi[\chi(\varrho_k), \chi_1(\varrho_k) \ldots \chi_{m-1}(\varrho_k)].}
Now if the substitutions $\sigma_h, \sigma_k$ produce the same permutation among the $\alpha$, then for every index $s = 0, 1 \ldots m-1$
\[
\chi_s(\varrho_h) = \chi_s(\varrho_k),
\]
and from this it follows by (2) that we must also have $\sigma_h = \sigma_k$.

If the substitution $\sigma_h$ produces the permutation $\pi_h$ among the roots of $F(x) = 0$, then we say that \gesperrt{the permutation $\pi_h$ corresponds to the substitution $\sigma_h$}.

To every substitution there corresponds a definite permutation, and the number of these permutations — counting in the identical permutation, which corresponds to the identical substitution — is therefore equal to $\mu$.

If the permutations $\pi_h$ and $\pi_k$ correspond to the substitutions $\sigma_h$ and $\sigma_k$, then to the composed substitution $\sigma_h \sigma_k$ there corresponds the composed permutation $\pi_h \pi_k$. For the permutations are the result of the corresponding substitutions on certain quantities of the field $\Omega(\varrho)$, and here it makes no difference whether we carry out two substitutions one after the other or carry out the substitution composed of both at one stroke (by \S.\,147). We therefore have the theorem:

\gesperrt{I. To the $\mu$ substitutions of the field $\Omega(\varrho)$}
\weq{3}{\sigma, \sigma_1, \sigma_2 \ldots \sigma_{\mu-1},}
\gesperrt{there correspond $\mu$ permutations of $m$ figures}
\weq{4}{\pi, \pi_1, \pi_2 \ldots \pi_{\mu-1}}
\gesperrt{in such a way that to the substitution composed of two substitutions there corresponds the permutation composed of the corresponding permutations.}

But from this it follows that the permutations (4) form a group. We denote this group by $P$ and call it the \gesperrt{Galois group of the equation $F(x) = 0$, or also the Galois group of each of the fields $\Omega(\alpha), \Omega(\alpha_1) \ldots \Omega(\alpha_{m-1})$}.

This group stands to the group of substitutions of the field $\Omega(\varrho)$ in the relation expressed by Theorem I., for which reason the two groups are called \gesperrt{isomorphic}.


% ===== scan leaf 503, printed 478 =====
% !PRINTED 478
\weberpage{478}

To find the characteristic properties of the Galois group, we have only to express the theorems \S.\,146, 1. to 4. somewhat differently. Of a function of the $m$ roots which remains unchanged under any permutation we say that it \gesperrt{admits this permutation}, and thus obtain:

a) \gesperrt{Every rational equation in $\Omega$ that holds between the $m$ roots of $F(x)$ remains valid when the roots are subjected to any permutation of the Galois group.}

b) \gesperrt{Every rational function in $\Omega$ of the $m$ roots of $F(x)$ that admits all the permutations of the Galois group is a number in $\Omega$.}

For if we express the roots of $F(x)$ as rational functions of $\varrho$, then a function of the roots $\alpha$ passes into a function $\varphi(\varrho)$. If a permutation of the Galois group has first been carried out, one obtains one of the conjugate numbers $\varphi(\varrho_a)$, which arises from $\varphi(\varrho)$ by a substitution $\sigma_a = (\varrho, \varrho_a)$. Now if $\varphi(\varrho) = 0$, then by the theorem \S.\,146, 1. all the conjugates $\varphi(\varrho_a) = 0$ as well, whereby a) is proved; and if the conjugate quantities $\varphi(\varrho_a)$ are all equal to one another, then their common value is, by \S.\,146, 4., contained in $\Omega$, whereby b) is proved.

To a) and b) a third statement is added:

c) \gesperrt{If any permutation $\pi$ of the roots of $F(x)$ is applicable to all rational equations in $\Omega$ that hold between the roots, then $\pi$ belongs to the Galois group, and the Galois group can therefore also be defined as the totality of all permutations that are applicable to all the rational equations between the roots.}

For by \S.\,143 one can take $\varrho$ to be a rational (e.\,g. linear) function of the $m$ quantities $\alpha, \alpha_1, \alpha_2 \ldots \alpha_{m-1}$ in such a way that all the $\Pi(m)$ values which one obtains from it by the $\Pi(m)$ permutations of the $\alpha$ are different from one another. If then $g(t) = 0$ is the Galois resolvent of $F(x) = 0$ whose root is this $\varrho$, then
\[
g(\varrho) = 0.
\]


% ===== scan leaf 504, printed 479 =====
% !PRINTED 479
\weberpage{479}

Hereupon, if $\varrho$ is expressed in terms of the roots $\alpha$, we can by hypothesis apply the permutation $\pi$, and thus, if $\varrho$ thereby passes into $\varrho_a$, we obtain $g(\varrho_a) = 0$; that is, $\varrho_a$ is also a root of the resolvent, and the permutation $\pi$ therefore corresponds to a substitution $(\varrho, \varrho_a)$ of the field $\Omega(\varrho)$, i.\,e. $\pi$ belongs to the \gesperrt{Galois group}.

From this we conclude further the following important theorem:

d) \gesperrt{If $P$ is a group of permutations of the $m$ roots $\alpha$ which possesses the properties a) and b), then $P$ is the Galois group of the equation $F(x) = 0$.}

For, in the first place, by c) every permutation of $P$ belongs to the \gesperrt{Galois group}, and $P$ is therefore certainly a divisor of the latter.

But if $P$ comprises only $\nu$ permutations, we use the following notation.
\weq{5}{\pi_1, \pi_2 \ldots \pi_\nu}
If we apply these permutations to $\varrho$, let the result be:
\weq{6}{\varrho_1, \varrho_2 \ldots \varrho_\nu.}
If we apply one of the permutations (5), say $\pi_k$, to one of the quantities (6), say to $\varrho_i$, the effect is the same as if $\pi_i \pi_k$ had been applied to $\varrho$; let the result of this permutation be $\varrho_i'$. Now it lies in the hypothesis that $P$ is a group that the composed permutation $\pi_i \pi_k$ also belongs to $P$, and hence that $\varrho_i'$ is contained among the quantities (6). At the same time, by \S.\,146, 2., $\varrho_i'$ is different from $\varrho_h'$ as soon as $\varrho_i$ is different from $\varrho_h$. From this it follows that the quantities
\weq{7}{\varrho_1', \varrho_2' \ldots \varrho_\nu'}
agree with the quantities (6), apart from the order.

The product
\[
g'(t) = (t - \varrho_1)(t - \varrho_2) \ldots (t - \varrho_\nu)
\]
therefore remains unchanged by the permutations of the group $P$, and is consequently, since we presuppose the property b) of $P$, a rational function of $t$ in $\Omega$. At the same time $g'(t)$ is a divisor of $g(t)$, and must therefore, since $g(t)$ is irreducible, agree with $g(t)$; hence $\nu$ is not smaller than the degree of the \gesperrt{Galois group}, and $P$ is identical with the \gesperrt{Galois group}.


% ===== scan leaf 505, printed 480 =====
% !PRINTED 480
\weberpage{480}

If, as above, we choose the quantity $\varrho$ as a rational function of the $\alpha$ in such a way that the $\Pi(m)$ quantities arising through the permutations of the $\alpha$,
\[
\varrho, \varrho', \varrho'' \ldots
\]
are all different from one another, then
\[
G(t) = (t - \varrho)(t - \varrho')(t - \varrho'') \ldots
\]
is a function in $\Omega$. Now each of the quantities $\varrho, \varrho', \varrho'' \ldots$ is a primitive quantity of the normal field $N = \Omega(\alpha, \alpha_1 \ldots \alpha_{m-1})$, and each of them is therefore the root of a \gesperrt{Galois} resolvent of degree $\mu$. These quantities, $\mu$ at a time, are the roots of one such equation. Hence $G(t)$, which has no equal roots, must split into nothing but irreducible factors of degree $\mu$, and from this it follows further \gesperrt{that $\mu$ is a divisor of $\Pi(m)$. At the same time $\mu$ is the degree of the Galois group} of $F(x)$.

If the degree of the \gesperrt{Galois} resolvent of an equation of degree $m$ has the greatest value $\Pi(m)$, we say, with an expression originating with \gesperrt{Kronecker}, that the equation has \gesperrt{no affect}. It has an affect that is the higher, the lower the degree $\mu$ of the \gesperrt{Galois} resolvent is. The quotient $\Pi(m) : \mu$, which must always be an integer, at most equal to $\Pi(m)$ and at least equal to 1, we shall call the \gesperrt{degree of the affect}; for an equation without affect it accordingly has the value 1. If the affect has the highest possible degree $\Pi(m)$, then the roots of the equation are themselves contained in the domain of rationality $\Omega$, and the equation is thus solved.

If by the adjunction of an algebraic quantity to $\Omega$ the \gesperrt{Galois} resolvent becomes reducible, a new resolvent of lower degree arises, and the degree of the affect of the equation $F(x) = 0$ increases. In this way, therefore, one comes nearer to the solution of the equation.

The \gesperrt{Galois} conception of the problem of solving an equation $F(x) = 0$ consists in this: by the successive adjunction of algebraic quantities of the simplest possible kind, the group is to be gradually diminished, or the affect raised, until it has reached its highest degree.

\begin{quote}
\gesperrt{The general equation of degree $m$ has no affect in the field $\Omega$ which consists of the rational functions of the coefficients $a_1, a_2 \ldots a_m$.}
\end{quote}


% ===== scan leaf 506, printed 481 =====
% !PRINTED 481
\weberpage{481}
\pdfbookmark[2]{§ 150. Transitive and intransitive groups}{bm167}
%<<CONT
For the coefficients $a_1, a_2 \ldots a_m$, regarded as independent
variables, can also be represented as the symmetric fundamental functions of
the roots $\alpha, \alpha_1, \alpha_2 \ldots \alpha_{m-1}$. If then $g(t)$ is
an irreducible factor of $G(t)$, expressed rationally in terms of the $a$,
which vanishes for $t = \varrho$, then, on representing $a_i$ and $\varrho$
by the $\alpha$, one obtains in $g(\varrho) = 0$ an identical equation. In
this, however, the $\alpha$ may be permuted arbitrarily, whereby the $a$ do
not change, while $\varrho$ can pass into any of the roots
$\varrho, \varrho', \varrho'' \ldots$ of $G(t)$. Hence $g(t)$ vanishes for
all the roots of $G(t)$ and must consequently coincide with
$G(t)$\footnote{\'Evariste \gesperrt{Galois} fell in a duel in the
year 1832, scarcely 20 years old. The first indication of the theory that
today bears his name is found in a memoir ``Analyse d'un m\'emoire sur la
r\'esolution alg\'ebrique des \'equations'', which appeared in 1830 in the
``Bulletin des sciences math\'em.''\ of \gesperrt{F\'erussac}. More detailed
communications are contained in the famous letter to \gesperrt{Auguste Chevalier},
written on the eve of his death, which is published in the ``Revue
encyclop\'edique'' of September 1832. Only in the year 1846 did \gesperrt{Liouville}
have all the already published papers of \gesperrt{Galois}, together with some
investigations from his literary remains, among them the most important,
``M\'em. sur les conditions de r\'esolubilit\'e des \'equations par
radicaux'', printed in Vol.~XI of his Journal. A German edition by \gesperrt{Maser}
appeared in 1889 (Berlin, published by Springer). Interesting biographical
communications are found in the ``Revue encyclop\'edique'' of September
1832, which \gesperrt{Liouville} did not reprint along with them and which are also
lacking in the German edition. Of further writings that have contributed to
the understanding or the further development of the \gesperrt{Galois} theory, let there
still be mentioned: J. A. \gesperrt{Serret} (``Cours d'alg\`ebre sup\'erieure'',
II. edition, 1854, IV. 1879); \gesperrt{Betti} (Annali di scienze fisiche e matematiche
[1853]); C. \gesperrt{Jordan}, ``Trait\'e des substitutions'' (1870). E. \gesperrt{Netto},
``Substitutionentheorie'' (Leipzig 1882).}.

\begin{center}
\S.\,150.

\gesperrt{Transitive and intransitive groups.}
\end{center}

From the \gesperrt{Galois} group $P$ we can derive a very
simple criterion for whether the equation $F(x) = 0$ is reducible or
irreducible.

Let us assume that $F(x)$, of degree $m$, is reducible, and that $f(x)$ is a
factor of $F(x)$ in $\Omega$ of degree $n$, with $n < m$. Let the roots of
$f(x)$
%CONT>>


% ===== scan leaf 507, printed 482 =====
% !PRINTED 482
\weberpage{482}
%<<CONT
\weq{1}{\alpha, \alpha_1, \alpha_2 \ldots \alpha_{n-1}}
and the remaining roots of $F(x)$ be
\weq{2}{\alpha_n, \alpha_{n+1} \ldots \alpha_{m-1}.}

If $\alpha'$ is any one of the quantities (1), and $\alpha''$ one of the
quantities (2), then no permutation can occur in the group $P$ by which
$\alpha'$ is carried into $\alpha''$; for by hypothesis
\[
f(\alpha') = 0,
\]
and if now a permutation occurred in $P$ by which $\alpha'$ were replaced by
$\alpha''$, then on account of the property \S.\,149, a) we should also have
$f(\alpha'') = 0$, which contradicts our assumption.

\gesperrt{Thus by the permutations of} $P$ \gesperrt{the roots (1) of} $f(x)$
\gesperrt{are permuted only among themselves.}

If conversely the group $P$ has the property that its permutations permute a
part of the $m$ roots $\alpha$, such as the system (1), only among
themselves, then the product
\[
(x - \alpha)(x - \alpha_1) \ldots (x - \alpha_{n-1}) = f(x)
\]
admits all the permutations of $P$, and is therefore by \S.\,149, b)
contained in $\Omega$; that is, $F(x)$ is reducible.

A permutation group is called \gesperrt{transitive} if it contains at least
one permutation which carries an arbitrary element into any other arbitrary
element; in the opposite case, that is, when the elements can be so
decomposed into two or more portions that by no permutation of the group
does an element of the one portion pass into an element of the other,
the group is called
\gesperrt{intransitive}. Accordingly we may sum up what has been proved in
the theorem:

% [sic] The printed theorem reverses the correspondence proved immediately above: irreducible corresponds to transitive, and reducible to intransitive. The source wording is preserved.
\gesperrt{1. The equation $F(x) = 0$ is reducible or irreducible, according
as its Galois group is transitive or intransitive.}

If a part of the elements is so connected that by permutations from $P$
every element of this part can pass into every other, we call the elements
of this part \gesperrt{transitively connected}. The various transitively
connected systems, which by all the permutations of $P$ are permuted only
among themselves, are called the \gesperrt{systems of
intransitivity}.


% ===== scan leaf 508, printed 483 =====
% !PRINTED 483
\weberpage{483}
\pdfbookmark[2]{§ 151. Primitive and imprimitive groups}{bm168}
\begin{center}
\S.\,151.

\gesperrt{Primitive and imprimitive groups.}
\end{center}

Now let $f(x) = 0$ be an irreducible equation of degree $n$, so that its
\gesperrt{Galois} group $P$ is transitive. Let the roots of $f(x)$ be
\weq{1}{\alpha, \alpha_1, \alpha_2 \ldots \alpha_{n-1}.}

If the field $\Omega(\alpha)$ is imprimitive (\S.\,144) and
$\Theta = \chi(\alpha)$ is an imprimitive element whose conjugate values fall
into $s$ systems of $r$ mutually equal values each, so that
\[
n = rs
\]
with $r > 1$, $s < n$, then the values (1) can be decomposed into $s$
sequences of $r$ elements each, which we shall denote by
\weq{2}{
\begin{array}{l}
A = \alpha, \alpha_1 \ldots \alpha_{r-1} \\
B = \beta, \beta_1 \ldots \beta_{r-1} \\
\cdots\cdots\cdots\cdots \\
S = \sigma, \sigma_1 \ldots \sigma_{r-1}
\end{array}
}
so that
\weq{3}{
\begin{array}{l}
\Theta = \chi(\alpha) = \chi(\alpha_1) \ldots = \chi(\alpha_{r-1}) \\
\Theta_1 = \chi(\beta) = \chi(\beta_1) \ldots = \chi(\beta_{r-1}) \\
\cdots\cdots\cdots\cdots\cdots\cdots\cdots\cdots \\
\Theta_{s-1} = \chi(\sigma) = \chi(\sigma_1) \ldots = \chi(\sigma_{r-1})
\end{array}
}
are the conjugate values of $\Theta$. By \S.\,144,
\weq{4}{(t - \Theta)(t - \Theta_1) \ldots (t - \Theta_{s-1}) = \varphi(t)}
is then an irreducible function in $\Omega$ of degree $s$ with respect to $t$,
whose roots are the values (3).

It now follows from (3) that the group $P$ must be so constituted that
all its permutations interchange the elements of the individual sequences
$A, B \ldots S$ only among themselves, and moreover permute the whole
sequences $A, B \ldots S$ with one another, so that two elements of one
sequence are never replaced by two elements of different sequences. For if,
say, a permutation $\pi$ of $P$ carried $\alpha$ and $\alpha_1$ into $\beta$
and $\sigma$, it would follow, since one may perform the permutation $\pi$
[by a), \S.\,149] in the equation $\chi(\alpha) = \chi(\alpha_1)$, that
$\chi(\beta) = \chi(\sigma)$ would also have to hold, which contra-
%CONT>>


% ===== scan leaf 509, printed 484 =====
% !PRINTED 484
\weberpage{484}
%<<CONT
dicts the assumption that the values (3) are different from one another. Permutation
groups $P$ which have this property, namely that the permuted
elements can be decomposed into rows $A, B \dots S$ of equally
many elements in such a way that no permutation of $P$ carries
two elements of the same row into elements of different rows,
are called \gesperrt{imprimitive}. The individual rows
$A, B \dots S$ are called the \gesperrt{systems of imprimitivity}.
Permutation groups for which such a decomposition of the elements
is not possible are called \gesperrt{primitive groups}.

It may very well happen that a group is imprimitive in more
than one way, that it possesses entirely different systems of
imprimitivity. Thus, for example, the group of cyclic permutations
of six elements $(1, 2, 3, 4, 5, 6)$ is imprimitive in a twofold
manner and has the systems of imprimitivity

\[
A = 1, 3, 5, \quad B = 2, 4, 6
\]
and
\[
A = 1, 4, \quad B = 2, 5, \quad C = 3, 6.
\]

Now that the concept of primitive and imprimitive groups has
thus been established, we can give the following expression to
what was proved above:

1. \gesperrt{An imprimitive field has an imprimitive
group.}

From the imprimitivity of the group there results an important
consequence for the imprimitive fields.

We will denote by
\[
\psi(x, x_1, \dots x_{r-1})
\]
a rational \gesperrt{symmetric function} of the $r$ variables
$x, x_1 \dots x_{r-1}$, and set

\weq{4}{
\begin{aligned}
\omega &= \psi(\alpha, \alpha_1 \dots \alpha_{r-1}) \\
\omega_1 &= \psi(\beta, \beta_1 \dots \beta_{r-1}) \\
& \cdot \ \cdot \ \cdot \ \cdot \ \cdot \ \cdot \ \cdot \ \cdot \ \cdot \ \cdot \\
\omega_{s-1} &= \psi(\sigma, \sigma_1 \dots \sigma_{r-1});
\end{aligned}
}
then it is to be proved that $\omega$ can be expressed rationally
in terms of $\Theta$, that is, is contained in the field $\Omega(\Theta)$. For
we have
\weq{5}{
\varphi(t)\left(\dfrac{\omega}{t - \Theta} + \dfrac{\omega_1}{t - \Theta_1} + \cdots + \dfrac{\omega_{s-1}}{t - \Theta_{s-1}}\right) = \Phi(t)
}
%CONT>>


% ===== scan leaf 510, printed 485 =====
% !PRINTED 485
\weberpage{485}

an integral rational function of $t$, whose coefficients are rational
functions of the $\alpha, \alpha_1 \ldots \alpha_{n-1}$, which remain unchanged
when any permutation of the group $P$ is carried out,
because under these permutations the $\omega$ either remain unchanged,
or are permuted among one another in the same way as the $\Theta$. By \S.\,149 b) these coefficients are therefore contained
in $\Omega$. If we then substitute in (5) the value $\Theta$ for the
indeterminate $t$, it follows, since $\varphi'(\Theta)$ is different from zero,
\weq{6}{\omega = \frac{\Phi(\Theta)}{\varphi'(\Theta)}.}

If we apply this to the coefficients of the product
\weq{7}{(u - \alpha)(u - \alpha_1) \ldots (u - \alpha_{r-1}) = \varphi(u, \Theta),}
where $u$ denotes a variable, it follows that this function
of degree $r$, whose roots are the $\alpha, \alpha_1 \ldots \alpha_{r-1}$, can be
expressed rationally in terms of $\Theta$.

Thus the quantity $\alpha$, which is originally the root of an
equation of degree $n$ in $\Omega$, is at the same time the root of an equation
of degree $r$, whose coefficients depend rationally on the root of an
equation of degree $s$. An equation $f(x) = 0$
whose root $\alpha$ has this property is called an \gesperrt{imprimitive equation}.

We can also express what has been proved thus:

\gesperrt{2. The imprimitive field $\Omega(\alpha)$ of degree $n$ passes,
by the adjunction to $\Omega$ of the field of degree $s$, $\Omega' = \Omega(\Theta)$,
into a field $\Omega'(\alpha)$ of degree $r$.}

We wish further to investigate whether the imprimitivity of the
group is a sufficient criterion for the imprimitivity of the
field, that is, whether one can always find imprimitive elements in a field $\Omega(\alpha)$
with imprimitive group.

Let now, therefore, $f(x) = 0$ be an irreducible equation of degree $n$
with imprimitive group, and let (1) be the roots
of this equation, decomposed into the sequences (2) in such a way that the
elements of these sequences are not separated from one another
by the permutations of the group.

We choose any symmetric function $\psi(x, x_1 \ldots x_{r-1})$
in such a way that the values
%CONT>>


% ===== scan leaf 511, printed 486 =====
% !PRINTED 486
\weberpage{486}
%<<CONT
\[
\begin{aligned}
y &= \psi(\alpha, \alpha_1 \ldots \alpha_{r-1}) \\
y_1 &= \psi(\beta, \beta_1 \ldots \beta_{r-1}) \\
& \phantom{=} \; \cdots\cdots\cdots\cdots \\
y_{s-1} &= \psi(\sigma, \sigma_1 \ldots \sigma_{r-1})
\end{aligned}
\tag{8}
\]
are all different from one another.

To see the possibility of this, we can, for example, set
\[
\begin{aligned}
\psi(t, A) &= (t - \alpha)(t - \alpha_1) \ldots (t - \alpha_{r-1}), \\
\psi(t, B) &= (t - \beta)(t - \beta_1) \ldots (t - \beta_{r-1}), \\
& \phantom{=} \; \cdots\cdots\cdots\cdots \\
\psi(t, S) &= (t - \sigma)(t - \sigma_1) \ldots (t - \sigma_{r-1})
\end{aligned}
\tag{9}
\]
and then (by \S.\,143, 1.) a rational value can be found for $t$ such
that these quantities all turn out different from one another. These
values can therefore be taken for the $y$ in (8).

If we now understand by $u$ an independent variable and set
\weq{10}{\varphi(u) = (u - y)(u - y_1) \ldots (u - y_{s-1})}
then $\varphi(u)$, since its coefficients remain unchanged under the permutations
of the group, is a function in $\Omega$ whose roots are the quantities $y$. \gesperrt{It is moreover irreducible}, for
from the assumed transitivity of the group it follows that
if a rational function of $y$ vanishes for one of the values (8),
it must also vanish for all the others.

If $\omega$ is any symmetric function of the $\alpha, \alpha_1 \ldots \alpha_{r-1}$,
then from the consideration of the expression
\[
\varphi(u) \left( \frac{\omega}{u - y} + \frac{\omega_1}{u - y_1} + \cdots + \frac{\omega_{s-1}}{u - y_{s-1}} \right),
\]
which is an integral rational function of $u$ in $\Omega$, we conclude, just as above,
by setting $u = y$, that $\omega$ can be represented rationally
in terms of $y$; and accordingly the function
\weq{11}{\psi(t, A) = \psi(t, y),}
whose roots are the quantities $\alpha, \alpha_1 \ldots \alpha_{r-1}$, can also be expressed
rationally in terms of $y$.

\gesperrt{The equation $\psi(t, A) = 0$ is also irreducible
in the field $\Omega(y)$}; for if $\psi_1(t, y)$ is a rational divisor
of $\psi(t, y)$, and if
%CONT>>


% ===== scan leaf 512, printed 487 =====
% !PRINTED 487
\weberpage{487}
%<<CONT
\[
\psi_1(\alpha, y) = 0,
\]
then, on account of the transitivity of the group, we can apply to this equation a permutation by which $\alpha$ passes into an arbitrary quantity of the sequence $A$, while $y$ remains unchanged; hence every root of $\psi(t, y)$ is at the same time a root of $\psi_1(t, y)$, and therefore the two functions are identical.

Hence the equation $f(x) = 0$ of degree $n$ is imprimitive in the sense fixed above, i.e.\ a root $\alpha$ of this equation can be regarded as a root of an equation of degree $r$, whose coefficients depend on the root $y$ of an equation of degree $s$.

To prove that then the field $\Omega(\alpha)$ (in the sense of \S.\,144) is also imprimitive, we have still to show that there are quantities in $\Omega(\alpha)$ which satisfy an irreducible equation of degree $s$.

We shall show that $y$ itself is such a quantity, and hence that $y$ can be represented rationally by $\alpha$ alone.

But this is very simple to conclude. By (11) we have
\weq{12}{\psi(\alpha, y) = 0,}
while $\psi(\alpha, y_1), \ldots \psi(\alpha, y_{s-1})$ are different from zero. For if, say, $\psi(\alpha, y_1) = 0$, then $\alpha$ would be a root of the equation of degree $r$, $\psi(t, y_1) = 0$, whose roots are precisely the quantities $\beta, \beta_1 \ldots \beta_{r-1}$, which are different from $\alpha$. Hence the two equations
\weq{13}{\psi(\alpha, u) = 0, \quad \varphi(u) = 0}
have one and only one common root $u = y$. The greatest common divisor of these two functions is therefore linear with respect to $u$, and it gives $y$ rationally in terms of the coefficients of $\varphi(u)$ and $\psi(\alpha, u)$, and \gesperrt{hence rationally in terms of} $\alpha$. With this the converse of 1.\ is proved:

3. \gesperrt{A primitive field has a primitive group.}

\begin{center}
\rule{3cm}{0.4pt}
\end{center}


% ===== scan leaf 513, printed 488 =====
\weberpage{488}% page number not printed on this page
\pdfbookmark[1]{Fourteenth Section. Application of the permutation groups to equations.}{bm169}
\pdfbookmark[2]{§ 152. Action of the permutation groups on functions of independent variables}{bm170}
% !PRINTED none
% Abschnitt (part) title page: "Vierzehnter Abschnitt." / section title, rule, then opening of §.152 with its heading and the start of the section's running text and first footnote. No running-head page number is printed on this title page.
\begin{center}
\gesperrt{Fourteenth Section.}
\end{center}

\begin{center}
\gesperrt{Application of the Permutation Groups to Equations.}
\end{center}

\begin{center}
\rule{2cm}{0.4pt}
\end{center}

\begin{center}
\S.\,152.
\end{center}

\begin{center}
\gesperrt{Effect of the Permutation Groups on Functions of Independent Variables.}
\end{center}

For a deeper penetration into algebra, a more exact study of the permutation groups is, after the results of the preceding section, necessary. We first set up some general theorems about them, which form the foundation for what follows\footnote{For the theory of permutations the following are especially to be noted: \gesperrt{Cauchy}, ``Journal de l'École polytechn. X. cah. (1815)'' (several memoirs). ``C. \gesperrt{Jordan}, Traité des substitutions et des équations algébriques'', Paris 1870. \gesperrt{Netto}, ``Substitutionentheorie und ihre Anwendung auf die Algebra.'' Leipzig 1882.}.

Let
\weq{1}{u,\ u_1,\ u_2 \ldots u_{m-1}}
be a system of mutually independent symbols (variables), and
\weq{2}{\psi(u,\ u_1,\ u_2 \ldots u_{m-1})}
an integral rational function of them with coefficients from an arbitrary field.

Two such functions $\psi$ are to be regarded as equal only when, in the expressions ordered according to powers and products of the $u$, corresponding terms have the same coefficients.


% ===== scan leaf 514, printed 489 =====
% !PRINTED 489
\weberpage{489}
If we subject the indices of the variables $u$, or, what comes to the same thing, the $u$ themselves to a permutation $\pi$, then the function $\psi$ may change, or it may remain unchanged. If it remains unchanged, we say that it \gesperrt{admits the permutation $\pi$}. If it admits all $\Pi(m)$ permutations that are possible at all, then it is symmetric and can be expressed by the symmetric fundamental functions (\S.\,44). The other extreme case is that in which the function takes nothing but distinct values for all $\Pi(m)$ arrangements. In general, certain permutations will leave the function $\psi$ unchanged, and others will change it. We now set up the theorem:

1. \gesperrt{The totality of all those permutations of $u, u_1, u_2 \ldots u_{m-1}$ which leave an integral function $\psi(u, u_1, u_2 \ldots u_{m-1})$ unchanged is a group of permutations.}

To prove it, assume that $\pi$ is one of the permutations which leaves $\psi$ unchanged. We express this by the following identity.
\[
\psi = \psi_\pi
\]
Now since the $u$ are independent variables, this identical equation must remain correct when the variables are subjected to any permutation $\pi'$ whatever. But when $\pi'$ is applied to $\psi_\pi$, the result is the same as when the composed permutation $\pi\pi'$ is applied to $\psi$. It follows therefore that
\[
\psi_{\pi'} = \psi_{\pi\pi'}
\]
and if now
\[
\psi = \psi_{\pi'}
\]
then it follows that
\[
\psi = \psi_{\pi\pi'},
\]
i.\,e. if the function $\psi$ remains unchanged under the permutations $\pi$ and $\pi'$, then it also remains unchanged under $\pi\pi'$, whereby, according to the definition of the group, the theorem we have stated is proved.

We have in this a means of forming groups of permutations, by taking any function of $m$ variables and seeking out all the permutations that leave such a function unchanged. That all groups of permutations can be formed in this way will appear later.


% ===== scan leaf 515, printed 490 =====
% !PRINTED 490
\weberpage{490}
2. \gesperrt{If}
\[
P = \pi,\ \pi',\ \pi'',\ \pi''' \ldots
\]
\gesperrt{is a group of permutations of $m$ elements and $\pi_1$ is any one among them, then the system}
\[
P_1 = \pi\pi_1,\ \pi'\pi_1,\ \pi''\pi_1,\ \pi'''\pi_1 \ldots,
\]
\gesperrt{agrees completely with $P$, apart from the order.}

For the permutations contained in $P_1$, whose number is just as great as that of the permutations $P$, are, by the group property of $P$ (\S.\,148, 3.), in any case all contained among the permutations of $P$. Moreover, the permutations of $P_1$ are (by \S.\,148, 2.) all different from one another, and hence $P$ is identical with $P_1$, apart from the order. In the same way one can also show that
\[
P_1' = \pi_1\pi,\ \pi_1\pi',\ \pi_1\pi'',\ \pi_1\pi''', \ldots
\]
is identical with $P$.

From this the following theorems result:

3. \gesperrt{Every permutation group contains the identical permutation.}

For among the permutations of $P_1$, $\pi_1$ itself must also occur, and if $\pi_0\pi_1 = \pi_1$, then (by \S.\,148, 2.) $\pi_0$ is the identical permutation.

4. \gesperrt{A permutation group $P$ contains, for every permutation, also the opposite one.}

For by 3. the identical permutation must occur in $P_1$. But if $\pi\pi_1$ is the identical permutation, then $\pi = \pi_1^{-1}$. Here $\pi_1$ may be any permutation from $P$.

If in place of the variables $u, u_1 \ldots u_{m-1}$ definite quantities $\alpha, \alpha_1, \alpha_2 \ldots \alpha_{m-1}$ are put, be they numerical values or quantities of any field, then Theorem 1. loses its general validity. If we take, for example, $m = 3$ and fix between the three quantities $\alpha, \alpha_1, \alpha_2$ the relation $2\alpha_2 = \alpha + \alpha_1$
while otherwise not restricting them further, then the function $\alpha_1 - \alpha_2$ remains unchanged only under the identical permutation and under the permutation
\[
\begin{pmatrix} 0, & 1, & 2 \\ 1, & 2, & 0 \end{pmatrix}.
\]


% ===== scan leaf 516, printed 491 =====
% !PRINTED 491
\weberpage{491}

But these two permutations do not form a group, because the repetition of the latter,
\[
\begin{pmatrix} 0, & 1, & 2 \\ 2, & 0, & 1 \end{pmatrix}
\]
does not occur among them.

For these functions the following theorem now holds:

5. \gesperrt{If $P$ is a permutation group of $m$ figures, and if $\alpha, \alpha_1 \dots \alpha_{m-1}$ are arbitrary quantities different from one another, then there are rational functions of the $\alpha$, even with rational coefficients, which do not change when a permutation of the group $P$ is applied to the indices of $\alpha$, and which do change when a permutation not belonging to $P$ is applied.}

To prove it, we take a function
\[
\varrho = \Phi(\alpha, \alpha_1 \dots \alpha_{m-1}),
\]
such as we have already considered in \S.\,143, which takes $\Pi(m)$ different values when the $\alpha$ are permuted in all possible ways.

If we denote the values which $\varrho$ takes through the application of the permutations of a group $P$ by
\weq{1}{\varrho, \varrho_1 \dots \varrho_{\mu-1},}
then by Theorem 2. these quantities are only interchanged among one another when a permutation $\pi_1$ from the group $P$ is applied to all of them at once. For the application of the permutation $\pi_1$ to all the functions (1) has the same effect as the application of the permutations contained in $P_1$ to $\varrho$. But by 2., $P_1$ is identical with $P$, apart from the order.

If, then, $t$ denotes a variable, the function
\weq{2}{\psi = (t - \varrho)(t - \varrho_1) \dots (t - \varrho_{\mu-1})}
remains unchanged when any permutation from $P$ is applied. Now if, by a permutation not contained in $P$, the quantities (1) pass into
\weq{3}{\varrho', \varrho_1' \dots \varrho_{\mu-1}'}
then at least $\varrho'$ certainly does not occur among the quantities (1), and consequently
\[
\psi_1 = (t - \varrho')(t - \varrho_1') \dots (t - \varrho_{\mu-1}')
\]
%CONT>>


% ===== scan leaf 517, printed 492 =====
% !PRINTED 492
\weberpage{492}
\pdfbookmark[2]{§ 153. Decomposition of permutations into transpositions and cycles}{bm171}
%<<CONT
not identical with $\psi$. One can therefore (by \S.\,143, 1.) give $t$ such a rational value that $\psi$ is different from all functions formed after the manner of $\psi_1$.

It is self-evident that these considerations also hold when independent variables take the place of the $\alpha$.

A function which admits all the permutations of a group $P$, while it changes under all permutations not belonging to $P$, is called a function \gesperrt{belonging to the group $P$}.

To the group $P_0$, which comprises \gesperrt{all} permutations of $m$ figures, belong the symmetric functions. This group is therefore also called the \gesperrt{symmetric group}.

A function $\varrho$ which belongs to the group consisting of the single identical permutation can serve as an unambiguous way of denoting the permutations; for there is only one permutation $\pi'$ by which $\varrho$ passes into $\varrho'$. We can therefore use for this permutation the symbol
\[
\pi' = (\varrho, \varrho')
\]
so that by
\[
(\varrho, \varrho), \quad (\varrho, \varrho'), \quad (\varrho, \varrho'') \dots
\]
all permutations of the $m$ figures are unambiguously denoted.

\begin{center}
\S.\,153.

\gesperrt{Decomposition of Permutations into Transpositions\\
and into Cycles.}
\end{center}

We have already mentioned, in the second section, on the occasion of the determinants, the composition of arbitrary permutations from a sequence of transpositions, that is, the permutation of only two figures. Such transpositions we denote by placing the figures in question side by side, e.g.\ by $(0,1)$ the permutation of the figures 0 and 1.

A transposition, repeated twice, leads to the identity, so that every transposition is opposite to itself. Now if
\[
\pi = \begin{pmatrix} 0, & 1, & 2 \dots m-1 \\ a_0, & a_1, & a_2 \dots a_{m-1} \end{pmatrix}
\]
is an arbitrary permutation of $m$ figures, then
\[
\pi(m-1, a_{m-1}) = \pi_1
\]
is a permutation which leaves $m-1$ unchanged, and which therefore by
%CONT>>


% ===== scan leaf 518, printed 493 =====
% !PRINTED 493
\weberpage{493}
%<<CONT
\[
\begin{pmatrix} 0, & 1, & 2 \dots m-2 \\ b_0, & b_1, & b_2 \dots b_{m-2} \end{pmatrix}
\]
can be represented, and is therefore a permutation of at most $m-1$ figures. Since now also
\[
\pi = \pi_1(m-1, a_{m-1})
\]
it follows from this by complete induction

1. \gesperrt{that every permutation can be decomposed (and indeed in infinitely many different ways) into transpositions.}

And from this it follows further:

2. \gesperrt{If a permutation group of $m$ figures contains all transpositions with one fixed figure, e.g.:
\[
(0,1), (0,2) \dots (0, m-1)
\]
then it is identical with the symmetric group.}

For such a group, as one recognises from the composition
\[
(1,2) = (0,1)(0,2)(0,1)
\]
contains also all the other transpositions, and by 1. all the permutations of the symmetric group can be composed from these.

A permutation $\pi$ is called \gesperrt{cyclic} if the figures can be so arranged in a sequence that by $\pi$ each figure passes into the following one, and the last again into the first, thus e.g.
\[
\begin{pmatrix} 0, & 1, & 2 \dots m-2, & m-1 \\ 1, & 2, & 3 \dots m-1, & 0 \end{pmatrix};
\]
such cyclic permutations are denoted more simply, by setting the figures of the cycle next to one another, by
\[
(0, 1, 2 \dots m-1).
\]
Here it is indifferent with which figure one begins; one could therefore also write
\[
(1, 2 \dots m-1, 0)
\]
in its place. The following theorem now holds:

3. \gesperrt{Every permutation $\pi$ can be decomposed, and indeed in only one way, into a sequence of cyclic permutations, so that no two of these cycles have a figure in common.}


% ===== scan leaf 519, printed 494 =====
% !PRINTED 494
\weberpage{494}

For if
% [sic] The print omits the comma between $a_1$ and $a_2$ in the lower row of (1); preserved source-faithfully.
\weq{1}{\pi = \begin{pmatrix} 0, & 1, & 2 \dots m-1 \\ a_0, & a_1 & a_2 \dots a_{m-1} \end{pmatrix},}
one begins with an arbitrary figure, say with 0, and continues the sequence
\weq{2}{0, \quad a_0 = b, \quad a_b = c \dots}
until one strikes upon some figure for the second time. The first figure to recur must be 0, since for every figure the one preceding it in the sequence (2) is uniquely determined by (1). The figures
\weq{3}{(0, b, c \dots)}
then form a first cycle. If not all $m$ figures of (1) are thereby exhausted, one picks out one of the remaining figures and proceeds in the same way, until all $m$ figures of (1) have been accommodated in the cycles. Since in every such cycle, for each figure, the preceding one as well as the following one is completely determined by (1), the cycles themselves are also uniquely determined. In the \gesperrt{designation} of $\pi$ by the cycles, however, not only may the various cycles be arranged arbitrarily, but in each cycle one may also begin with any one of its figures. A figure that is not changed forms by itself a one-term cycle.

We choose an entirely arbitrary example, by which the procedure at once becomes clear:
\[
\pi_1 = \begin{pmatrix} 0, & 1, & 2, & 3, & 4, & 5, & 6, & 7 \\ 2, & 3, & 6, & 7, & 1, & 4, & 5, & 0 \end{pmatrix} = (0, 2, 6, 5, 4, 1, 3, 7)
\]
\[
\pi_2 = \begin{pmatrix} 0, & 1, & 2, & 3, & 4, & 5, & 6, & 7 \\ 4, & 2, & 1, & 3, & 5, & 7, & 0, & 6 \end{pmatrix} = (0, 4, 5, 7, 6)(1, 2)(3)
\]
\[
\pi_3 = \begin{pmatrix} 0, & 1, & 2, & 3, & 4, & 5, & 6, & 7 \\ 7, & 2, & 3, & 6, & 5, & 0, & 1, & 4 \end{pmatrix} = (0, 7, 4, 5)(1, 2, 3, 6).
\]
The setting of the cycles side by side is equivalent to a composition in the sense used hitherto; only it is to be remarked that permutations which contain no common figure whatever can always be interchanged in the composition.

One-term cycles, which change nothing, are also often omitted in the designation.


% ===== scan leaf 520, printed 495 =====
% !PRINTED 495
\weberpage{495}

From the simultaneous consideration of the decomposition of permutations into cycles and into transpositions there results a proof of the division of the permutations into two classes, with which we have already become acquainted in the second section, in connection with determinants.

If $\tau$ denotes a transposition and $\pi$ an arbitrary permutation, then in the compounded permutation $\tau\pi$ the number of cycles is greater by one or smaller by one than in $\pi$. The two figures that are interchanged by $\tau$ can occur either in the same cycle of $\pi$ or in two different cycles. Let us assume that a cycle occurring in $\pi$ is $\gamma = (1, 2 \dots a, a + 1 \dots b)$, and that $\tau$ contains two figures that occur in $\gamma$, say $\tau = (1, a)$; then
\[
\tau\gamma = (1, a + 1 \dots b)\ (2, 3 \dots a),
\]
that is, $\gamma$ is decomposed by composition with $\tau$ into two cycles, while the remaining cycles are not touched by $\tau$.

But if the two figures of $\tau$ occur in two different cycles of $\pi$, then let these two cycles be
\[
\gamma = (1, 2, 3 \dots a), \quad \gamma' = (1', 2', 3' \dots a')
\]
and $\tau = (1, 1')$. Then
\[
\tau\gamma\gamma' = (1, 2', 3' \dots a', 1', 2, 3 \dots a),
\]
that is, the two cycles $\gamma$, $\gamma'$ are united by $\tau$ into a single cycle, while again the remaining cycles stay unchanged.

If, then, we have represented a permutation $\pi$ of $m$ figures, consisting of $\nu$ cycles (the one-term cycles being counted along with the rest), by $\mu$ transpositions, then by composition with these $\mu$ transpositions in reversed order it can be converted into the identical permutation, which consists of $m$ one-term cycles. In all, therefore, $m - \nu$ cycles have been gained, and since to every transposition there corresponds a gained or a lost cycle, we must have $\mu \equiv m - \nu\ (\mathrm{mod}.\ 2)$; that is, $\mu$ is an even or an odd number according as $m - \nu$ is even or odd. But the last number depends only on the permutation $\pi$, not on the manner of the decomposition into transpositions. We have thereby the theorem:
%CONT>>


% ===== scan leaf 521, printed 496 =====
% !PRINTED 496
\weberpage{496}

\gesperrt{4. The permutations of $m$ figures fall into two kinds, of which the first is decomposable into an even, the second into an odd number of transpositions.}

Each of these two kinds comprises equally many permutations, namely $\frac{1}{2}\Pi(m)$; for by the addition of a fixed transposition every permutation of the first kind passes over into one of the second kind, and conversely.

The composition of two permutations of the same kind always gives a permutation of the first kind, while a permutation of the first and one of the second kind, composed, give a permutation of the second kind, as is at once to be seen from the decomposition into transpositions. From this it follows:

\gesperrt{5. The permutations of the first kind form a group.}

If we call a function of $m$ variables $u_0, u_1 \dots u_{m-1}$, or also of $m$ mutually distinct quantities $\alpha_0, \alpha_1 \dots \alpha_{m-1}$, which changes its sign when two of the variables are interchanged with one another, as e.\,g. the product of all the differences:
\[
(u_0 - u_1)(u_0 - u_2) \dots (u_0 - u_{m-1})
\]
\[
(u_1 - u_2) \dots (u_1 - u_{m-1})
\]
\[
\cdots\cdots\cdots\cdots
\]
an \gesperrt{alternating function}, then such a function is not changed by the permutations of the first kind, while by a permutation of the second kind it is changed, namely transformed into the opposite value. The alternating functions therefore \gesperrt{belong} to the group of the permutations of the first kind.

As we earlier called the group to which the symmetric functions belong the symmetric group, so we call the group of the permutations of the first kind, to which the alternating functions belong, the \gesperrt{alternating group}. The group which consists of the single identical permutation we call the \gesperrt{identical group} or also the \gesperrt{unit group}.

A cyclic permutation of $n$ terms can be decomposed into $n - 1$ transpositions, as one sees from the composition
\[
(1, 2, 3 \dots n) = (1, 2)(1, 3) \dots (1, n)
\]
%CONT>>


% ===== scan leaf 522, printed 497 =====
% !PRINTED 497
\weberpage{497}
%<<CONT
recognises, and from this it follows that a cyclic permutation belongs to the first or to the second kind, according as the number of figures is odd or even.

Just as every permutation can be composed of transpositions, so too every permutation of the first kind can be composed of three-term cycles. It suffices if this is proved for every pair of transpositions, since every permutation of the first kind can be compounded of such pairs. But we have
\[
(0, 1)(0, 2) = (0, 1, 2)
\]
\[
(0, 3)(1, 2) = (0, 1, 2)(0, 1, 3),
\]
from which the correctness of the assertion is to be seen, since the two transpositions of a pair have either one figure in common or none. Hence:

\gesperrt{6. All permutations of the first kind can be decomposed (in infinitely many ways) into cyclic permutations of three elements.}

From 6. there follows further

\gesperrt{7. A permutation group which contains all three-term cyclic permutations with two fixed figures must contain the whole alternating group.}

For from the cycles $(0, 1, 2), (0, 1, 3), (0, 1, 4) \dots (0, 1, m - 1)$ all three-term cycles, and hence by 6. all permutations of the first kind, can be compounded, as from the compositions
\[
(1, 0, 2) = (0, 1, 2)(0, 1, 2),
\]
\[
(0, 2, 3) = (0, 1, 2)(1, 0, 3)(1, 0, 2),
\]
\[
(1, 3, 4) = (1, 0, 3)(0, 1, 4)(0, 1, 3)
\]
\[
(2, 3, 4) = (1, 0, 2)(1, 3, 4)(0, 1, 2)
\]
may be seen.

\gesperrt{8. If $m$ is greater than 4, then a permutation group which contains all permutations $(0, 1)(2, 3)$ composed of any two transpositions without a common figure is divisible by the alternating group.}

For we have
\[
(0, 1, 2) = (0, 1)(3, 4)(3, 4)(0, 2)
\]
%CONT>>


% ===== scan leaf 523, printed 498 =====
% !PRINTED 498
\weberpage{498}
%<<CONT
and one can therefore, if besides the four figures 0, 1, 2, 3 a fifth, 4, is present, compose all three-term cycles from pairs of transpositions of the form $(0, 1)(3, 4)$.

To this we append the proofs of two further important theorems relating to transitive permutation groups:

\gesperrt{9. If a transitive permutation group of $m$ figures contains a single transposition, then it is either the symmetric group or it is imprimitive.}

Let the given group be $P$, and let the $m$ figures that are permuted by it be $0, 1, 2 \dots m - 1$, and let the transposition $(0, 1)$ occur in it. Suppose that it contains altogether the transpositions
\[
(0, 1), (0, 2), (0, 3), \dots (0, \mu - 1),
\]
but no other transposition with the figure 0. Now if $\mu = m$, then by 2. $P$ is the symmetric group. But if $\mu < m$, then by the same theorem $P$ contains all permutations of the figures
\[
M = 0, 1, 2 \dots \mu - 1.
\]
From this it follows that $P$ can contain no transposition of one of the figures $0, 1, 2, 3 \dots \mu - 1$ with another figure, say with $\mu$. For if $(1, \mu)$ occurs in $P$, then $(0, \mu) = (0, 1)(1, \mu)(0, 1)$ also occurs in it, contrary to the hypothesis.

Now if $P$ is transitive, there is in it a permutation $\pi$ by which the figure 0 passes into a figure not contained in the system $M$, say into $\mu$; and then by $\pi$ none of the figures of $M$ can pass into a figure of $M$; for if, e.\,g., 1 passes into $r$, then
\[
\pi^{-1}(0, 1)\pi = (\mu, r)
\]
and $(\mu, r)$, if $r$ belongs to $M$, cannot, as we have seen, occur in $P$. By every permutation of $P$, therefore, the figures of the system $M$ are either permuted only among themselves, or they are carried over into the figures of a system $M'$ entirely distinct from it, and in $P$ all permutations of the figures $M'$ also occur.

If with $M$ and $M'$ not all the figures are yet exhausted, there is a permutation in $P$ by which $M$ is carried over into a third system $M''$, which again can have no figure in common with $M$ or with $M'$.


% ===== scan leaf 524, printed 499 =====
% !PRINTED 499
\weberpage{499}
For if, when $\alpha$ denotes any one of the figures
$0, 1, 2, \ldots n - 1$, we denote by $\alpha_\pi$ and $\alpha_{\pi'}$ the figures into which $\alpha$
passes under $\pi$ and under $\pi'$, then $\alpha_{\pi\pi'}$ is the figure into which $\alpha_\pi$
passes under $\pi'$, or $\alpha$ under $\pi\pi'$. If $\alpha, \beta, \gamma \ldots$ are the figures
from the system $M$, then
\[
\begin{gathered}
\pi'^{-1}(\alpha, \beta) \pi' = (\alpha_{\pi'}, \beta_{\pi'}) \\
\pi\pi'^{-1}(\alpha, \beta) \pi'\pi^{-1} = (\alpha_{\pi'\pi^{-1}}, \beta_{\pi'\pi^{-1}}) = (\gamma, \delta).
\end{gathered}
\]
Now if $\alpha_{\pi'} = \gamma_\pi$ is contained in $M'$, then $\alpha_{\pi'\pi^{-1}} = \gamma$ is contained
in $M$, and consequently $\beta_{\pi'\pi^{-1}} = \delta$ must likewise be contained in $M$,
since the transposition $(\gamma, \delta)$ occurs in $P$. Consequently $\alpha_{\pi'} = \gamma_\pi$
and $\beta_{\pi'} = \delta_\pi$, i.e. if one of the figures from $M''$ is contained in $M'$, then
$M'$ and $M''$ must be identical.
The totality of the figures $0, 1 \ldots n - 1$ therefore falls apart
into systems $M, M', M'' \ldots$ of $m$ figures each, which are permuted imprimitively
by $P$; that is, the group $P$ is imprimitive.
If $n$ is a prime, there are no imprimitive permutation groups;
hence in this case a transitive group which contains a transposition
must be the symmetric group.
In exactly the same way we conclude from Theorem 7:

10. \gesperrt{If a transitive permutation group of
$n$ figures contains a three-term cycle,
then it contains the whole alternating group
or it is imprimitive.}

For if the group $P$ contains the three-term cycles
\[
(0, 1, 2), (0, 1, 3) \ldots (0, 1, m - 1),
\]
but no other with the figures $0, 1$, then by 7.
it contains the whole alternating group of the figures $M = 0, 1, 2 \ldots m - 1$.
Now if $m < n$, no three-term cycle can occur in $P$ which joins one
or two of the figures of $M$ with other figures. For if $P$ contains
e.g. $(0, m, m + 1)$, then, as one sees from the compositions
\[
(0, 2, 1) (0, m, m + 1) (0, 1, 2) = (1, m, m + 1)
\]
it also contains $(1, m, m + 1), \ldots (m - 1, m, m + 1)$ and hence
by 7. the whole alternating group of the figures $0, 1, 2 \ldots m + 1$,
which is contrary to the hypothesis. The $n$ figures of $P$ therefore
fall apart, as before, into several systems $M, M', M'' \ldots$ of imprimitivity. We may add further,
%CONT>>


% ===== scan leaf 525, printed 500 =====
% !PRINTED 500
\weberpage{500}
%<<CONT
11. \gesperrt{that, when the case of imprimitivity occurs, the group $P$ contains, in case 9. the whole symmetric group, in case 10. the whole alternating group, for each single system of imprimitivity.}

Repeated composition of a permutation with itself is denoted by
powers. If $\pi$ is any permutation, then among the successive
powers
\[
\pi^0, \pi^1, \pi^2, \pi^3 \ldots,
\]
in which $\pi^0 = 1$ denotes the identical permutation, the same
permutation must necessarily recur. But from $\pi^u = \pi^{u+e}$ it follows
that $\pi^e = 1$. If $e$ is the least positive exponent for which $\pi^e = 1$,
then the permutations
\weq{4}{1, \pi, \pi^2 \ldots \pi^{e-1}}
are all different from one another, while, as the sequence is continued with
$\pi^e, \pi^{e+1}, \ldots \pi^{2e-1} \ldots$, they repeat periodically in the same
order. The sequence (4) is called the \gesperrt{period of the permutation}
$\pi$, and $e$ is called the \gesperrt{degree of} $\pi$. The period
of $\pi$ forms a group, since composing any two of its elements
yields another element of the same period.
If $\pi$ is a cyclic permutation
\[
\pi = (0, 1, 2 \ldots n - 1),
\]
then for even $n$
\[
\pi^2 = (0, 2, 4 \ldots n - 2) (1, 3 \ldots n - 1),
\]
and for odd $n$
\[
\pi^2 = (0, 2, 4 \ldots n - 1, 1, 3 \ldots n - 2),
\]
and the degree of $\pi$ is in both cases equal to $n$.
If $\pi$ is decomposed into several cycles, the degree of $\pi$
is equal to the least common multiple of the number of terms of the single cycles, e.\,g.
\[
\pi = (0, 1) (2, 3, 4), \; \pi^2 = (0) (1) (2, 4, 3), \; \pi^3 = (0, 1) (2) (3) (4),
\]
% [sic] The print separates the disjoint cycles in $\pi^4$ and $\pi^5$ by commas, unlike $\pi^2$ and $\pi^3$ immediately above; the printed notation is preserved.
\[
\pi^4 = (0), (1), (2, 3, 4), \; \pi^5 = (0, 1), (2, 4, 3), \; \pi^6 = 1.
\]
If $\pi$ occurs in any permutation group $P$,
then $P$ contains the whole period of $\pi$.


% ===== scan leaf 526, printed 501 =====
% !PRINTED 501
\weberpage{501}
\pdfbookmark[2]{§ 154. Divisors of the group, cosets and conjugate groups}{bm172}
\begin{center}
\S. 154.\\[4pt]
\gesperrt{Divisors of Groups. Subsidiary Groups and conjugate Groups.}
\end{center}

We saw in \S.\,152 that a function of $m$ independent variables always determines a permutation group of $m$ figures, but that this does not always hold when definite quantities are put in place of the independent variables. By the \gesperrt{Galois} theory, however, this difference is evened out, or at least reduced to its kernel. Now let $F(x) = 0$ be any equation of degree $m$ in a given field $\Omega$ with the mutually distinct roots $\alpha, \alpha_1, \alpha_2 \ldots \alpha_{m-1}$.

The normal field
\[
\Omega(\varrho) = \Omega(\alpha, \alpha_1, \alpha_2 \ldots \alpha_{m-1})
\]
we shall denote by $N$.

By the \gesperrt{Galois} group $P$ of the field $N$, or of the equation $F(x) = 0$, we may understand (by \S.\,146, 149) either the group of substitutions of the field $N$, or also the permutation group of the indices of the $\alpha$ isomorphic with it.

Let $p$ be the degree of $P$, and let us denote the operations of $P$ (be they now substitutions of $N$ or permutations of the $\alpha$) by
\weq{1}{P = \pi,\ \pi_1,\ \pi_2, \ldots \pi_{p-1}}
If the system $Q$ of the permutations contained in $P$
\weq{2}{\varkappa,\ \varkappa_1,\ \varkappa_2, \ldots \varkappa_{q-1}.}
% Source check: the faint terminal mark in (2) is the printed period.
forms a group by itself, that is, if any two of the elements of $Q$, compounded with one another, again yield an element of $Q$, then the group $Q$ is called a \gesperrt{divisor} or \gesperrt{factor of} $P$\footnote{Also called \gesperrt{subgroup}.}).

We first derive an important general theorem on the divisors of a group.

If the degree $q$ of the divisor is smaller than the degree $p$ of $P$, we take an element $\pi_1$ of $P$ not contained in $Q$ and form the system
\weq{3}{Q\,\pi_1 = \varkappa\,\pi_1,\ \varkappa_1\,\pi_1 \ldots \varkappa_{q-1}\,\pi_1,}
%CONT>>


% ===== scan leaf 527, printed 502 =====
% !PRINTED 502
\weberpage{502}
%<<CONT
which we represent as the symbolic product $Q\pi_1$. The elements
of (3) are distinct both from one another and from the elements of
$Q$; for from $\varkappa_1\pi_1 = \varkappa_2\pi_1$ would follow $\varkappa_1 = \varkappa_2$, and from
$\varkappa = \varkappa_1\pi_1$ would follow $\pi_1 = \varkappa_1^{-1}\varkappa$, which is not possible if
$\pi_1$ does not belong to $Q$. The elements of $Q\pi_1$ form no
group, for otherwise $\varkappa_1\pi_1\varkappa_2\pi_1 = \varkappa_3\pi_1$ or $\pi_1 = \varkappa_1^{-1}\varkappa_3\varkappa_2^{-1}$
would have to hold, which would again require that $\pi_1$ be contained in $Q$.

We call the system $Q\pi_1$ a \gesperrt{subsidiary group} to $Q$
(within $P$). The subsidiary group remains unaltered in its
totality when $\pi_1$ is replaced by any element $\varkappa\pi_1$ of $Q\pi_1$,
because $Q\varkappa$ is identical with $Q$.

If we call every system $Q\pi$ generally a subsidiary group,
so that $Q$ itself is counted among them, then we can show
that two subsidiary groups are either \gesperrt{entirely identical}, or
\gesperrt{contain no common element.} For if $Q\pi_1$,
$Q\pi_2$ are two subsidiary groups, and if $\varkappa_1\pi_1 = \varkappa_2\pi_2$, in which $\varkappa_1, \varkappa_2$ are two
elements from $Q$, then, if we set $\varkappa = \varkappa_2^{-1}\varkappa_1$, it follows that
% [sic] The source prints a bare \pi here; the derivation requires \pi_1. The source wording is preserved.
$\pi_2 = \varkappa\pi$. But since $\varkappa$ belongs to $Q$, and consequently $Q\varkappa = Q$,
it follows that
\[
Q\pi_2 = Q\varkappa\pi_1 = Q\pi_1.
\]
If the group $P$ is not yet completely exhausted by (2) and (3), we can form a further subsidiary group $Q\pi_2$,
and can continue in the formation of the subsidiary groups until $P$ is decomposed into nothing but subsidiary groups. Since each of these
% [sic] The source gives the number of elements in each subsidiary group as \nu rather than the previously defined group degree q; preserved source-faithfully.
subsidiary groups contains $\nu$ elements, there follows the important fundamental theorem of \gesperrt{Cauchy}:

1. \gesperrt{The degree of a divisor of a group $P$ is
a divisor of the degree of $P$.}

As a special case there is contained in this the fact that the \gesperrt{degree
of an element of the group} (\S.\,153) \gesperrt{is always a divisor
of the degree of the group}.

The decomposition of $P$ into its subsidiary groups we can
(following \gesperrt{Galois}) express very fittingly by the symbolic equation
\weq{4}{P = Q + Q\pi_1 + Q\pi_2 + \cdots + Q\pi_{j-1},}
if
\weq{5}{p = jq}
holds. The quotient $j$ we call the \gesperrt{index of the divisor}
$Q$. This number expresses no peculiarity of the group $Q$
in itself, but only a relation of $Q$ to $P$.


% ===== scan leaf 528, printed 503 =====
% !PRINTED 503
\weberpage{503}

Since, as we have already seen above, two subsidiary groups of $Q$ are either wholly identical, or contain no common element, we can further state the theorem:

2. \gesperrt{If $\pi$ is an arbitrary permutation from $P$, then the system of subsidiary groups}
\[
Q\pi,\ Q\pi_1\pi,\ Q\pi_2\pi\ldots Q\pi_{j-1}\pi
\]
\gesperrt{is distinguished from the system}
\[
Q,\ Q\pi_1,\ Q\pi_2\ldots Q\pi_{j-1}
\]
\gesperrt{only by the arrangement.}

We extend the designation subsidiary groups also to systems of the form $\pi Q$, and can decompose $P$ into such systems as well. From the decomposition (4) there follows in particular the second
\[
P = Q + \pi_1^{-1}Q + \pi_2^{-1}Q + \cdots \pi_{j-1}^{-1}Q.
\]
For if $\pi_1$ does not occur in $Q$, then $\pi_1^{-1}$ does not occur in it either, and we obtain a first subsidiary group $\pi_1^{-1}Q$. If then $\pi_2$ does not occur in $Q\pi_1$, then $\pi_2^{-1}$ does not occur in $\pi_1^{-1}Q$ either; for from $\pi_2^{-1} = \pi_1^{-1}\varkappa$ would follow $\pi_2 = \varkappa^{-1}\pi_1$, and so forth.

We here give the preference to the decomposition (4) on account of the significance which, as we shall see presently, it has for the algebraic problem.

We call a quantity $\psi$ of the field $N$ \gesperrt{belonging to the group $Q$} (when necessary with the addition ``within $P$''), if $\psi$ does not change when one of the operations of $Q$ is applied to it, but does change when any other operation from $P$ is applied. As one sees, this concept of belonging is somewhat extended as against the one fixed in \S.\,152, inasmuch as the permutations that lie outside $P$ do not come into consideration here at all. For the case that $P$ is the symmetric group, however, the two definitions coincide.

With this definition of belonging, the following theorems now hold generally.

3. \gesperrt{To every divisor $Q$ of $P$ there belong quantities in $N$, and every quantity in $N$ belongs to a determinate divisor $Q$ of $P$.}

The first part of the theorem is already proved in \S.\,152, 5. For if a function of the quantities $\alpha$ within the sym-
%CONT>>


% ===== scan leaf 529, printed 504 =====
% !PRINTED 504
\weberpage{504}

%<<CONT
metric group belongs to a group $Q$, then it belongs to $Q$ also within every other group $P$ of which $Q$ is a divisor.

The second part, however, follows from the fundamental properties of the \gesperrt{Galois} group.

To see this, we shall introduce the following notation. If any permutation $\pi$ of the group $P$ is applied to a function $\psi$ of the field $N$, let $\psi$ pass over into $\psi(\pi)$; if a permutation $\pi'$ is applied to $\psi(\pi)$, the result is the same as if the composed permutation $\pi\pi'$ had been applied to $\psi$, i.\,e.\ we have
\weq{6}{\psi(\pi)(\pi') = \psi(\pi\pi').}
If now $\psi$ remains unchanged by the permutation $\varkappa$, then
\[
\psi = \psi(\varkappa),
\]
and to this, by \S.\,149, a), any permutation $\pi$ from the group $P$ may be applied. From this it follows that
\[
\psi(\pi) = \psi(\varkappa\pi).
\]
If therefore $\psi(\pi)$ is also equal to $\psi$, then $\psi(\varkappa\pi) = \psi$ as well, i.\,e.\ the permutations which leave $\psi$ unchanged form a group, q.\,e.\,d.

If $\psi$ is a function in $N$ belonging to the group $Q$, so that
\weq{7}{\psi(\varkappa) = \psi(\varkappa_1) = \cdots = \psi(\varkappa_{q-1}) = \psi,}
and $\pi_1$ is any permutation of $P$ not belonging to $Q$, then we can apply $\pi_1$ to the equations (7) and obtain
\[
\psi(\varkappa\pi_1) = \psi(\varkappa_1\pi_1) = \cdots = \psi(\varkappa_{q-1}\pi_1) = \psi_1,
\]
and $\psi_1$ is different from $\psi$. If conversely, for some $\pi$,
\weq{8}{\psi_1 = \psi(\pi) = \psi(\varkappa\pi_1)}
holds, then $\pi$ must be contained in the sequence of substitutions
\weq{9}{Q\pi_1 = \varkappa\pi_1,\ \varkappa_1\pi_1,\ldots \varkappa_{\nu-1}\pi_1}
% [sic] The source uses the terminal index \nu-1, although Q is indexed through q-1 in (2) and (7); preserved source-faithfully.
for by applying $\pi^{-1}$ to (8) it follows that $\varkappa\pi_1\pi^{-1}$ must be contained in $Q$, say equal to $\varkappa_1$; but from this it follows that
\[
\pi = \varkappa_1^{-1}\varkappa\pi_1,
\]
which is indeed contained in (9).

We can accordingly state the following theorem:

4. \gesperrt{A function $\psi$ belonging to $Q$ passes, under all the permutations of a coset $Q\pi$, and}
%CONT>>


% ===== scan leaf 530, printed 505 =====
% !PRINTED 505
\weberpage{505}

%<<CONT
\gesperrt{by no other operation of $P$ into a determinate quantity $\psi_1$ different from $\psi$.}

Accordingly, to the $j$ subsidiary groups there correspond just as many functions $\psi$, namely
\weq{10}{\psi,\ \psi_1,\ \psi_2\ldots\psi_{j-1},}
and from Theorem 2. there results the theorem:

5. \gesperrt{The quantities (10) undergo a permutation when one and the same operation $\pi$ of $P$ is applied to all of them simultaneously.}

For we can represent the quantities (10) thus
\[
\psi(\varkappa),\ \psi(\varkappa\pi_1),\ \psi(\varkappa\pi_2),\ldots \psi(\varkappa\pi_{j-1}),
\]
and if now $\pi$ is applied to them,
\[
\psi(\varkappa\pi),\ \psi(\varkappa\pi_1\pi),\ \psi(\varkappa\pi_2\pi),\ldots \psi(\varkappa\pi_{j-1}\pi),
\]
in which $\varkappa$ and $\pi$ are arbitrary permutations of $Q$ and $P$. But in each of the two sequences
\[
\begin{aligned}
&\varkappa,\ \varkappa\pi_1,\ \varkappa\pi_2\ldots\varkappa\pi_{j-1}\\
&\varkappa\pi,\ \varkappa\pi_1\pi,\ \varkappa\pi_2\pi\ldots\varkappa\pi_{j-1}\pi
\end{aligned}
\]
there occurs one and only one permutation from each of the subsidiary groups (4).

The functions (10), which, as we have seen, are all different from one another, are called \gesperrt{conjugate functions}.

By what precedes we can easily form the groups to which the individual functions (10) belong. We have, e.g., $\psi_1 = \psi(\pi_1)$, and if now $\pi$ is a permutation of $P$ by which $\psi_1$ remains unaltered, then
\[
\psi(\pi_1\pi) = \psi(\pi_1)
\]
must hold. But we can apply to this equation the permutation $\pi_1^{-1}$ contained in $P$, whereby
\[
\psi(\pi_1\pi\pi_1^{-1}) = \psi
\]
results; that is, $\pi_1\pi\pi_1^{-1}$ must belong to $Q$. If we set it equal to $\varkappa$, it follows that
\weq{11}{\pi = \pi_1^{-1}\varkappa\pi_1.}
If conversely $\pi$ has this form, then
\[
\psi_1(\pi) = \psi(\pi_1\pi) = \psi(\pi_1\pi_1^{-1}\varkappa\pi_1) = \psi(\varkappa\pi_1) = \psi_1,
\]
and $\psi_1$ therefore admits this permutation. Accordingly we can denote the group to which $\psi_1$ belongs by the symbol
\[
\pi_1^{-1}Q\pi_1
\]
%CONT>>


% ===== scan leaf 531, printed 506 =====
% !PRINTED 506
\weberpage{506}

%<<CONT
represent, which is in fact a group, as is to be seen from $\pi_1^{-1}\varkappa\pi_1\pi_1^{-1}\varkappa_1\pi_1 = \pi_1^{-1}\varkappa\varkappa_1\pi_1$. From this it follows at the same time that $\pi_1^{-1}Q\pi_1$ is \gesperrt{isomorphic} with $Q$ (\S.\,149).

The groups to which the conjugate functions (10) belong, namely
\weq{12}{Q,\quad \pi_1^{-1}Q\pi_1,\quad \pi_2^{-1}Q\pi_2,\ldots \pi_{j-1}^{-1}Q\pi_{j-1}}
we call \gesperrt{conjugate divisors} of $P$, or briefly conjugate groups\footnote{The name ``equally entitled subgroup'' is also in use for this.}.

If $\pi$ is an arbitrary permutation from $P$, and hence of the form $\varkappa\pi_r$, then the group $\pi^{-1}Q\pi$ is always contained among the conjugate groups (12), namely $= \pi_r^{-1}\varkappa^{-1}Q\varkappa\pi_r = \pi_r^{-1}Q\pi_r$.

The derivation of $\pi^{-1}Q\pi$ from $Q$ is also called a \gesperrt{transformation of the group $Q$ by $\pi$}, and $\pi^{-1}Q\pi$ a \gesperrt{group transformed from $Q$}.

For the formation of the permutations of the conjugate groups $\pi^{-1}Q\pi$ the following simple rule serves:

6. \gesperrt{One obtains the permutation $\pi^{-1}\varkappa\pi$ by carrying out the permutation $\pi$ in the cycles of $\varkappa$.}

To prove the rule, let, decomposed into its cycles (\S.\,153),
\[
\varkappa = (\alpha,\beta,\gamma\ldots)(\alpha',\beta',\gamma'\ldots)\ldots
\]
and let
\[
\pi =
\begin{pmatrix}
\alpha, & \beta, & \gamma & \ldots \alpha', & \beta', & \gamma' & \ldots \\
\alpha_\pi, & \beta_\pi, & \gamma_\pi & \ldots \alpha'_\pi, & \beta'_\pi, & \gamma'_\pi & \ldots
\end{pmatrix}.
\]
By $\pi^{-1}$, $\alpha_\pi$ passes into $\alpha$; by $\varkappa$, $\alpha$ passes into $\beta$; and by $\pi$, $\beta$ is carried into $\beta_\pi$. By $\pi^{-1}\varkappa\pi$, therefore, $\alpha_\pi$ passes into $\beta_\pi$. Since the same consideration is applicable to $\beta_\pi, \gamma_\pi\ldots$ and so forth, it follows that:
\[
\pi^{-1}\varkappa\pi = (\alpha_\pi,\beta_\pi,\gamma_\pi\ldots)(\alpha'_\pi,\beta'_\pi,\gamma'_\pi\ldots),
\]
by which the rule 6. is proved.

Of the decomposition into subsidiary groups we can make yet another application, to the proof of the theorem:

7. \gesperrt{The degree of a transitive permutation group of $m$ figures is always a multiple of $m$.}


% ===== scan leaf 532, printed 507 =====
% !PRINTED 507
\weberpage{507}
\pdfbookmark[2]{§ 155. Reduction of the Galois resolvent by adjunction. Normal divisors of a group}{bm173}

For all the permutations in $P$ that leave the figure $0$ unchanged
form a divisor $Q$ of $P$. Moreover, on account of the transitivity,
there are in $P$ permutations $\pi_1, \pi_2 \ldots \pi_{m-1}$ which carry the
figure $0$ into $1, 2 \ldots m-1$, and we have

\[
P = Q + Q\pi_1 + Q\pi_2 + \cdots + Q\pi_{m-1},
\]

hence the degree of $P$ is equal to $m$ times the degree of $Q$.

\begin{center}
§. 155.

\gesperrt{Reduction of the Galois resolvent by adjunction.\\
Normal divisor of a group.}
\end{center}

We have now to consider the gradual reduction of the group
of a given equation, which can come about through the adjunction
of certain algebraic quantities.

Let, as in the preceding section, $P$ be the \gesperrt{Galois}
group of degree $p$ and $Q$ one of its divisors of degree $q$
and of index $j$; further let $\psi$ be a function of the roots
belonging to $Q$, and let

\weq{1}{\psi, \psi_1, \psi_2 \ldots \psi_{j-1}}

be the conjugate quantities. We place the following theorem
at the head.

1. \gesperrt{The conjugate quantities (1) are the roots
of an irreducible equation of degree $j$
in $\Omega$.}

For if $t$ denotes a variable, then

\weq{2}{\varphi(t) = (t - \psi)(t - \psi_1) \ldots (t - \psi_{j-1})}

remains unchanged (by §.154, 5.) when a permutation from $P$ is
applied; hence the coefficients of the function $\varphi(t)$, whose
roots are the quantities (1), are contained in $\Omega$ [§. 149, b)].

To prove the irreducibility of $\varphi(t)$, we assume
that $\Phi(t)$ is any function in $\Omega$ which vanishes
for $t = \psi$, so that $\Phi(\psi) = 0$. Since all the permutations
of $P$ may be applied to this equation, it follows that
$\Phi(\psi_1), \Phi(\psi_2) \ldots \Phi(\psi_{j-1})$ must also be zero, and hence that $\Phi(t)$
must be divisible by $\varphi(t)$. But in this the irreducibility
is contained.


% ===== scan leaf 533, printed 508 =====
% !PRINTED 508
\weberpage{508}

To this is joined the \gesperrt{theorem of Lagrange}\footnote{\gesperrt{Lagrange}, Réflexions sur la résolution algébrique des équations. Mémoires de l'Académie de Berlin, années 1770, 1771. Oeuvres de Lagrange. Tome III. The theorem was, to be sure, given by \gesperrt{Lagrange} only in a more special form. The general formulation is due to \gesperrt{Galois}.}:

2. \gesperrt{Every quantity of the field $N$ which admits the permutations
 of the group $Q$ is contained in the
field $\Omega(\psi)$, if $\psi$ is a function belonging
 to $Q$.}

A function $\omega$ which admits the permutations of $Q$ passes,
under the permutations of a subsidiary group $Q\pi_1$, into one and
the same function. Hence to the conjugate
values

\weq{3}{\psi, \psi_1, \psi_2 \ldots \psi_{j-1}}

there correspond the values

\weq{4}{\omega, \omega_1, \omega_2 \ldots \omega_{j-1},}

which, however, are not necessarily all different from one another.

If one applies to the sequences of quantities (3), (4) one of the
permutations of $P$, a certain permutation takes place, and
indeed the same one in both sequences, since, if for example $\psi_1$ passes into $\psi_2$,
then $\omega_1$ must also pass into $\omega_2$ (§. 154, 5.).

Let us now consider the sum

\[
\varphi(t) \left( \frac{\omega}{t - \psi} + \frac{\omega_1}{t - \psi_1} + \cdots + \frac{\omega_{j-1}}{t - \psi_{j-1}} \right) = \chi(t),
\]

which is an integral function of $t$ of degree $(j-1)$; we then find
that it remains unchanged under all the permutations $P$ and
is consequently contained in $\Omega$. If one then sets $t = \psi$ and observes
that $\varphi(t)$ has no equal roots, it follows that

\weq{5}{\omega = \frac{\chi(\psi)}{\varphi'(\psi)}.}

Thus $\omega$ is expressed rationally in terms of $\psi$, and this is the
content of theorem 2.

3. \gesperrt{If we adjoin to the field $\Omega$ a function $\psi$ belonging
to $Q$, the group of the
field $N$ reduces to $Q$.}

% [page alignment mirrored 2026-08-07] Printed page 508 contains the complete
% line through "because $\psi$ remains".  Page 509 begins "unchanged by $Q$";
% the English boundary now follows the source-confirmed German recut.
For if we denote the field $\Omega(\psi)$ by $\Omega'$, then \gesperrt{first},
every equation in $\Omega'$ between $\alpha, \alpha_1 \ldots \alpha_{m-1}$ admits the permutations
of $Q$, because $Q$ is contained in $P$, and because $\psi$ remains
%CONT>>


% ===== scan leaf 534, printed 509 =====
% !PRINTED 509
\weberpage{509}
%<<CONT
unchanged by $Q$; and \gesperrt{secondly}, every function that remains
unchanged by $Q$ is, by \gesperrt{Lagrange}'s theorem, contained
in $\Omega'$. But these [by §. 149, a), b)] are the characteristic
marks of the \gesperrt{Galois} group in the field $\Omega'$.
In order, therefore, to reduce the \gesperrt{Galois} group from degree $p$ to degree $q$,
\gesperrt{one} root of an auxiliary equation of degree $j$
must be adjoined to $\Omega$.

Theorem 3. we can also express thus:

\gesperrt{The normal field $N = \Omega(\alpha, \alpha_1 \ldots \alpha_{m-1})$ is a field
of degree $p$ over $\Omega$ and of degree $q$ over $\Omega' = \Omega(\psi)$.}

To the lowering of the group by the adjunction of $\psi$ there corresponds,
as already follows from the general principles,
a splitting of the \gesperrt{Galois} resolvent. Let us assume
that (as in §. 145) $g(t) = 0$ is the \gesperrt{Galois} resolvent and
$\varrho$ one of its roots, and that by the permutations of $Q$,
$\varrho$ passes over into

\[
\varrho, \varrho_1, \varrho_2 \ldots \varrho_{q-1},
\]

and by the permutations of the subsidiary group $Q\pi_i$ into

% [sic] The source omits the comma after \varrho_{0,i} in this parallel list.
\[
\varrho_{0,i}\ \varrho_{1,i}, \varrho_{2,i} \ldots \varrho_{q-1,i},
\]

then each of the factors of $g(t)$

\[
g_i(t) = (t - \varrho_{0,i})(t - \varrho_{1,i}) \ldots (t - \varrho_{q-1,i})
\]

is unchanged by the substitutions of $Q$ and is therefore contained in $\Omega(\psi)$.

If we set

\[
g(t, \psi) = (t - \varrho)(t - \varrho_1) \ldots (t - \varrho_{q-1}),
\]

then it results, if we carry out in it a permutation from $Q\pi_i$,
that

\[
g_i(t) = g(t, \psi_i),
\]

and $g(t)$ is thus decomposable in the following manner:

\weq{6}{g(t) = g(t, \psi)\, g(t, \psi_1) \ldots g(t, \psi_{j-1}),}

in which each factor is of degree $q$, and no two of these
factors have a common divisor.

If $Q$ and $Q'$ are two different divisors of the group $P$, then
$Q$ and $Q'$ will have certain permutations in common, among
which the identical permutation is always found. It is
possible that this is the only common permutation of $Q$
and $Q'$, and in that case these two groups are called \gesperrt{coprime}.
But still more common elements may
%CONT>>


% ===== scan leaf 535, printed 510 =====
% !PRINTED 510
\weberpage{510}
%<<CONT
be present. The totality $R$ of all common permutations of $Q$
and $Q'$ we call the \gesperrt{greatest common divisor},
or also the \gesperrt{intersection}, of $Q$ and $Q'$. This $R$ is always a
group, for if $\pi_1$ and $\pi_2$ both occur in $Q$ as well as in $Q'$,
then $\pi_1 \pi_2$ must also occur in $Q$ and in $Q'$, and hence also in $R$.
The concept carries over at once to several
groups $Q, Q', Q'' \ldots$

If $\psi$ is a function belonging to $Q$ and $\psi'$ one belonging to $Q'$,
we can determine the rational numbers $x, x'$ (by §. 143, 1.)
so that $\omega = x\psi + x'\psi'$ becomes a function belonging to $R$.
For $\omega$ certainly admits the permutations of $R$, since
$\psi$ and $\psi'$ admit them. If then $\pi$ is a permutation of $P$ not
contained in $R$, then certainly $\psi = \psi(\pi)$ and
$\psi' = \psi'(\pi)$ do not both hold; hence we can determine $x, x'$ so that
$\omega = \omega(\pi)$ likewise does not hold. If therefore we adjoin $\psi$ and $\psi'$,
and consequently $\omega$, simultaneously, the group of the field $N$
reduces by 3. to $R$. In the same way we may conclude for more than two groups $Q$,
and obtain the theorem:

4. \gesperrt{If $Q, Q', Q'' \ldots$ are divisors of $P$, $R$ their intersection,
$\psi, \psi', \psi'' \ldots$ functions belonging to $Q, Q', Q'' \ldots$,
then by the simultaneous adjunction of
$\psi, \psi', \psi'' \ldots$ the group of the
field $N$ reduces to $R$.}

If we adjoin not merely one root $\psi$ of the auxiliary equation
$\varphi(t) = 0$, but all the quantities conjugate to $\psi$,
namely $\psi, \psi_1, \psi_2 \ldots \psi_{j-1}$ — if, then, we derive from $\Omega$ the field

\[
\Omega'' = \Omega(\psi, \psi_1, \psi_2 \ldots \psi_{j-1})
\]

— then by this theorem the outcome is that the
group of $N$ in $\Omega''$ becomes the greatest common divisor $R$
of all the groups conjugate to $Q$, so that (by §. 154)
$R$ is the greatest common divisor of all the groups $\pi^{-1}Q\pi$,
as $\pi$ runs through the permutations of $P$.

\gesperrt{Of special interest now is the case in which all the conjugate
groups $\pi^{-1}Q\pi$ are identical with one another.} In this
case, by Theorem 2., each of the conjugate quantities
$\psi, \psi_1, \psi_2 \ldots \psi_{j-1}$ is contained in $\Omega(\psi)$, the fields
$\Omega(\psi), \Omega(\psi_1) \ldots \Omega(\psi_{j-1})$ and $\Omega(\psi, \psi_1 \ldots \psi_{j-1})$ are identical,
and $\Omega(\psi)$ is a \gesperrt{normal field} over $\Omega$.


% ===== scan leaf 536, printed 511 =====
% !PRINTED 511
\weberpage{511}
\pdfbookmark[2]{§ 156. The group of the resolvents}{bm174}

We therefore call a divisor $Q$ of the group $P$ which has this
property a \gesperrt{normal divisor} (or normal factor)\footnote{\gesperrt{Galois} speaks of the proper decomposition (d\'ecomposition propre) of a group; hence the expression ``proper divisor'' which is also in use; more recent writers designate the normal divisors as ``distinguished or invariant subgroups''.}.

But if $Q$ is a normal divisor, then the adjunction of a
single root $\psi$ of the auxiliary equation is equivalent to the
simultaneous adjunction of all these roots.

5. \gesperrt{If $Q$ itself is not normal, then the greatest
common divisor $R$ of all the conjugate
divisors $\pi^{-1}Q\pi$ is certainly normal.}

For if a permutation $\varkappa$ is contained in $R$, and hence in all the groups
$\pi^{-1}Q\pi$, then the same holds also of every $\pi^{-1}\varkappa\pi$;
that is, $\pi^{-1}R\pi$ is a divisor of $R$, and consequently, since it is of the same
degree, $= R$.

A group which has no normal divisor other than itself and the identical
group is called a \gesperrt{simple group}.
The intersection of all the groups conjugate to a simple group
is either the group itself or the identical group.

If $Q$ is a normal divisor of $P$, and $R$ a normal divisor
of $Q$, then $R$ is indeed a divisor of $P$, but by no means always
a normal one. On the other hand, if $R$ is a normal divisor of $P$, then $R$
is also a normal divisor of every divisor $Q$ of $P$ of which $R$
is a divisor.

\begin{center}
\S.\,156.

\gesperrt{The Group of the Resolvents.}
\end{center}

The auxiliary equation $\varphi(t) = 0$, on which the determination of the
function $\psi$ depends, passes over into the \gesperrt{Galois} resolvent
when the divisor $Q$ to which $\psi$ belongs is the identical group.
For then, by the theorem of \gesperrt{Lagrange} (\S.\,155, 2.),
every function of the field $N$, and hence also the roots $\alpha$ themselves,
can be expressed rationally in terms of $\psi$, and $N$ is identical with $\Omega(\psi)$.

We shall therefore call these equations $\varphi(t) = 0$, in a general
sense, \gesperrt{resolvents}. But two cases are here
to be distinguished.


% ===== scan leaf 537, printed 512 =====
% !PRINTED 512
\weberpage{512}

1. If the conjugate groups $\pi^{-1}Q\pi$ are coprime, then
the simultaneous adjunction of all the roots of the
resolvent $\varphi(t) = 0$ is (by Theorem 4, \S.\,155) equivalent to
the adjunction of a function belonging to the unit group,
and the solution of the given equation is reduced to the complete
solution of the resolvent. We have
\[ N = \Omega(\psi, \psi_1 \ldots \psi_{j-1}), \]
and a \gesperrt{Galois} resolvent of the equation $\varphi(t) = 0$ is
at the same time a \gesperrt{Galois} resolvent of the original equation.
In this case we call $\varphi(t) = 0$ a \gesperrt{total resolvent}
of the given equation.

2. If the conjugate groups $\pi^{-1}Q\pi$ have a divisor $R$
of degree $r$ different from the unit group, which is then
a normal divisor of $P$, then the simultaneous adjunction
of all the functions conjugate to $\psi$ is equivalent
to the adjunction of a quantity belonging to $R$. The \gesperrt{Galois}
resolvent of the given equation is not yet completely solved
by this adjunction; rather it is only decomposed into factors
of degree $r$. In this case $\varphi(t) = 0$ is called a
\gesperrt{partial resolvent}.

If $P$ is a \gesperrt{simple group}, then \gesperrt{only} total resolvents
exist, whereas, if $P$ has normal divisors, a partial
resolvent can be found for every such normal
divisor.

The same distinction also comes to light when we investigate the
\gesperrt{Galois} group of the resolvent $\varphi(t)$. This
group consists of all the permutations that are produced in the sequence of the
quantities
\weq{1}{\psi, \psi_1, \psi_2 \ldots \psi_{j-1}}
by the application of all the operations $\pi$ of $P$; for every equation in $\Omega$ between the quantities (1)
remains valid when such a permutation is carried out,
since one can apply the operations $\pi$ to such an equation;
and if a function in $\Omega$ of the quantities (1) admits all these
permutations, then it also admits all the permutations $\pi$
and is therefore equal to a quantity in $\Omega$.

The degree of this group is now to be determined. Among
the operations of $P$, those and only those will produce no
change among the quantities (1) which occur simultaneously in
%CONT>>


% ===== scan leaf 538, printed 513 =====
% !PRINTED 513
\weberpage{513}
\pdfbookmark[2]{§ 157. Reduction of the Galois group by adjunction of arbitrary irrationalities}{bm175}
%<<CONT
the groups $\pi^{-1}Q\pi$ of all these functions, and hence also in their
greatest common divisor $R$. Let the operations
of $R$ be denoted by $\sigma$. If then $\pi$ and $\pi'$ are two
operations that produce the same permutation among the quantities (1),
then $\pi'\pi^{-1}$ will restore the original arrangement of the $\psi$,
and hence be equal to one of the quantities $\sigma$, or
\[ \pi' = \sigma\pi. \]
We therefore have the result:

\begin{quote}
\gesperrt{The permutations of the subsidiary group $R\pi$, and only
these, produce one and the same permutation
among the quantities (1).}
\end{quote}

The degree of the \gesperrt{Galois} group of the resolvent $\varphi(t) = 0$
is therefore equal to the number of these subsidiary groups, i.e.\ equal to the
quotient $p : r$, or the index of the divisor $R$ of $P$.

If $R$ is the identical group, so that $r = 1$, and consequently $\varphi(t) = 0$
is a total resolvent, then the degree of its group is just as high
as the degree of the group of the original equation, and the two
groups are moreover isomorphic. With regard to the group, then,
nothing is gained. The given equation is equivalent to the
resolvent, however different their degrees may be.

If, on the other hand, $Q$ is a normal divisor, so that $R$ is
identical with $Q$, then $\varphi(t) = 0$ is a partial resolvent and the degree
of its group is equal to the index $j$ of the divisor $Q$ of $P$.
After the adjunction of a root of this resolvent, the group
of the original equation reduces to $Q$, and hence to
degree $q$. A splitting of the group has therefore taken place.

If one wishes to form resolvents of the lowest possible
degree, one has to seek out divisors of the group $P$ of the smallest
possible index, and hence of the highest possible degree; but if
at the same time a reduction of the group is to occur, then the
divisors must be normal divisors.

\begin{center}
\S.\,157.

\gesperrt{Reduction of the Galois Group by Adjunction
of Arbitrary Irrationalities.}
\end{center}

According to the \gesperrt{Galois} conception, the problem of solving
an algebraic equation is replaced by another,
%CONT>>


% ===== scan leaf 539, printed 514 =====
% !PRINTED 514
\weberpage{514}
%<<CONT
namely, by the adjunction of algebraic quantities of the simplest
possible nature, to bring about a splitting of the \gesperrt{Galois} resolvent,
and thus a lowering of the degree of the group.

We now put the question thus. In the original
field $\Omega$ the \gesperrt{Galois} resolvent $g(t)$ is irreducible. The
field $\Omega$ is to be extended to another field $\Omega'$ in such a way
that $g(t)$ becomes reducible. Here $\Omega'$ is to be an algebraic field
over $\Omega$, and so by \S.\,143 there must be an algebraic
quantity $\varepsilon$ such that $\Omega' = \Omega(\varepsilon)$. This quantity $\varepsilon$ will be the
root of a certain irreducible equation in $\Omega$, which
we shall denote by
\weq{1}{\chi(\varepsilon) = 0}

Let us assume that in the field $\Omega(\varepsilon)$ the irreducible
factor
\weq{2}{g_1(t) = g_1(t, \varepsilon)}
separates off from $g(t)$, having the root $t = \varrho$. The degree of $g_1(t, \varepsilon)$ we shall
denote by $q$, and the roots by
\weq{3}{\varrho, \varrho_1, \varrho_2 \ldots \varrho_{q-1};}
so that
\weq{4}{g_1(t, \varepsilon) = (t - \varrho)(t - \varrho_1) \ldots (t - \varrho_{q-1}).}
It is now first to be shown that the substitutions
\[ (\varrho, \varrho), (\varrho, \varrho_1), (\varrho, \varrho_2) \ldots (\varrho, \varrho_{q-1}) \]
contained in the group $P$ form a group $Q$. To show this, we set, as
before (\S.\,146),
\[ \varrho_i = \Theta_i(\varrho) \]
and observe that then the equation
\[ g_1[\Theta_i(t), \varepsilon] = 0 \]
has a root, namely $t = \varrho$, in common with $g_1(t)$, and hence,
since $g_1(t)$ is irreducible, is divisible by $g_1(t)$. From this it follows
that, if $\varrho_1$ and $\varrho_2$ are any two of the roots (3), then
$\Theta_i(\varrho_2) = \varrho_3$ is also contained among these roots. But the
substitution $(\varrho, \varrho_3)$ is composed of $(\varrho, \varrho_1)$ and $(\varrho, \varrho_2)$, and
with this the group property of $Q$ is proved. We denote the index of $Q$ by $j$ and set
\weq{5}{p = jq.}


% ===== scan leaf 540, printed 515 =====
% !PRINTED 515
\weberpage{515}

The product (4) now admits the substitutions of the group $Q$,
and hence, if $\psi$, as above, denotes a function in $N$
belonging to this group, then by \S.\,155, 2. the function
$g_1(t,\varepsilon)$ can be expressed rationally in terms of
$\psi$; that is, $g_1(t,\varepsilon)$ is a function of $t$ in
the field of degree $j$, $\Omega(\psi)$, and shall accordingly
be denoted by $g(t,\psi)$.

The quantity $\psi$ is one of the roots of an irreducible
equation of degree $j$

\weq{6}{\varphi(u) = 0,}

whose remaining roots are $\psi_1, \psi_2 \dots \psi_{j-1}$,
and by \S.\,155, (6) $g(t)$ is decomposable into the factors

\weq{7}{g(t) = g(t,\psi)\, g(t,\psi_1) \dots g(t,\psi_{j-1}).}

Now

\[ g(t,\psi) = g_1(t,\varepsilon), \]

and from (7) it follows that this equation does not continue to
hold if, on the left-hand side, one of the other roots of (6) is
put in place of $\psi$. One can therefore assign to $t$ a
rational value such that the equation

\weq{8}{g(t,u) - g_1(t,\varepsilon) = 0}

has only the one root $u = \psi$ in common with equation (6).
The greatest common divisor of (6) and (8) is then linear with
respect to $u$, and if one sets $u = \psi$ in it, it results
\gesperrt{that $\psi$ is rationally expressible in terms of
$\varepsilon$}. The field $\Omega(\psi)$ is therefore certainly a
divisor of the field $\Omega(\varepsilon)$. The degree of the
field $\Omega(\varepsilon)$, i.\,e. the degree of the auxiliary
equation $\chi(\varepsilon) = 0$, is therefore a multiple of $j$
and never lower than $j$. If the degree of the auxiliary
equation is equal to $j$, which, for example, always occurs when
the degree of $\chi$ is a prime, then $\Omega(\varepsilon)$ is
identical with $\Omega(\psi)$; that is, $\varepsilon$ is also
rationally expressible in terms of $\psi$. In this case the
conjugate values $\varepsilon, \varepsilon_1 \dots
\varepsilon_{j-1}$ are all the roots of $\chi = 0$, and one
obtains them by replacing, in the rational expression of
$\varepsilon$ in terms of $\psi$, the function $\psi$ by each of
the conjugate functions $\psi, \psi_1 \dots \psi_{j-1}$. Hence
$\varepsilon$ itself can be taken for $\psi$, and one finds a
decomposition

\[ g(t) = g(t,\varepsilon)\, g(t,\varepsilon_1) \dots g(t,\varepsilon_{j-1}). \]

The quantities $\varepsilon, \varepsilon_1, \varepsilon_2 \dots
\varepsilon_{j-1}$ all belong to the field $\Omega(\varrho)$.


% ===== scan leaf 541, printed 516 =====
% !PRINTED 516
\weberpage{516}
\pdfbookmark[2]{§ 158. Imprimitive groups}{bm176}

These results are of great importance and deserve to be
specially emphasised.

Following \gesperrt{Kronecker}, if $F(x) = 0$ is the given
equation and $N$ the associated \gesperrt{Galois} field, the quantities of
$N$ are called \gesperrt{the natural irrationalities of the equation
$F(x) = 0$}.

We then have the theorem:

\gesperrt{Every possible reduction of the Galois group is
brought about by the adjunction of a natural irrationality. If
$j$ is the index of the reduced group with respect to the
original one, then the reduction cannot come about through the
adjunction of a field of degree lower than $j$. But the degree
can be equal to $j$, and then the adjoined irrationality is a
natural one. If the degree of the adjoined field is greater
than $j$, then it is a multiple of $j$.}

\begin{center}
\S.\,158.

\gesperrt{Imprimitive Groups.}
\end{center}

We make application of our general theorems to groups of a
special nature. First we take the \gesperrt{imprimitive groups}. We have
seen in \S.\,151 that, if $f(x) = 0$ is an irreducible equation
of degree $n$ with imprimitive group, then $n$ can be decomposed
into two factors $r \boldsymbol{.} s$ in such a way that the
roots split into $s$ rows of $r$ terms each:

\weq{1}{\begin{aligned}
A &= \alpha,\ \alpha_1 \dots \alpha_{r-1},\\
B &= \beta,\ \beta_1 \dots \beta_{r-1},\\
&\cdots\cdots\cdots\cdots\cdots\cdots\cdots\cdots\\
S &= \sigma,\ \sigma_1 \dots \sigma_{r-1},
\end{aligned}}

and that by the permutations of the group $P$ the quantities of
the individual rows are not separated from one another, but are
only permuted among themselves, while on the other hand the
rows in turn are permuted among one another.

Among the permutations of the \gesperrt{Galois} group $P$ of $f(x)$
there will now be certain ones, whose totality we shall denote by $Q$,
%CONT>>


% ===== scan leaf 542, printed 517 =====
% !PRINTED 517
\weberpage{517}

%<<CONT
which leave the individual rows $A, B \dots S$ in their
place and only permute the indices within each row. These
permutations form a group by themselves, since by the
composition of two of them the rows $A, B \dots
S$ cannot be interchanged. This group $Q$ is a
\gesperrt{normal divisor} of $P$. For if $\varkappa$ is a permutation of
$Q$, and $\pi$ one from $P$, then the composed
permutation $\pi^{-1} \varkappa \pi$ belongs again to $Q$, because the
permutation among the $A, B \dots S$ that is produced by $\pi^{-1}$
is not altered by $\varkappa$ and is undone again by
$\pi$. The group $Q$ is
intransitive, and by adjunction of a function $\psi$ belonging
to $Q$ not only does the \gesperrt{Galois} resolvent
become reducible, but the function $f(x)$ itself splits into $s$
factors of degree $r$,

\weq{2}{f(x) = f_\alpha(x,\psi)\, f_\beta(x,\psi) \dots f_\sigma(x,\psi).}

The roots of $f_\alpha(x,\psi) = 0$ are the quantities $\alpha$,
and the \gesperrt{Galois} group of this equation consists of all the
permutations that are produced by $Q$ among the $\alpha$.

The quantity $\psi$ is, if $j$ is the index of the divisor $Q$ of
$P$, a root of a partial resolvent of degree $j$,
$\chi(u) = 0$, which is a normal equation and hence its
own \gesperrt{Galois} resolvent (\S.\,155). If we denote
its roots by

\weq{3}{\psi,\ \psi_1,\ \psi_2 \dots \psi_{j-1},}

then the \gesperrt{Galois} group of this resolvent consists of the
substitutions

\weq{4}{(\psi,\psi),\ (\psi,\psi_1),\ (\psi,\psi_2) \dots (\psi,\psi_{j-1}),}

which are produced by the permutations of the subsidiary groups

\weq{5}{Q,\ Q\pi_1,\ Q\pi_2 \dots Q\pi_{j-1}}

Each of these subsidiary groups produces a certain
permutation among the rows $A, B \dots S$, and
indeed these permutations are all different from one another;
for if, say, the same permutation among the rows were produced
by $\pi_1$ and $\pi_2$, then $\pi_2 \pi_1^{-1}$ belongs
to $Q$, and $Q\pi_1$, $Q\pi_2$ are not different
subsidiary groups. The permutation group of the partial resolvent
$\chi(u) = 0$ is therefore isomorphic with the group of the
permutations that are produced by $P$ among the rows $A, B \dots S$.

If we understand by $Q_\alpha$ the totality of the permutations
of $P$ by which the row $A$ is not displaced, then
$Q_\alpha$ is
%CONT>>


% ===== scan leaf 543, printed 518 =====
% !PRINTED 518
\weberpage{518}

%<<CONT
certainly a group. If $\pi_\beta$ is a permutation by which the
row $A$ passes into the row $B$, which by the hypothesis of
transitivity must exist in $P$, then
$Q_\alpha \pi_\beta$ is a subsidiary group of $Q_\alpha$ by which
$A$ passes into $B$. $\pi_\beta^{-1} Q_\alpha \pi_\beta =
Q_\beta$ is a group conjugate to $Q_\alpha$, and indeed the one
by whose permutations the row $B$ is not displaced.
In this way we obtain the conjugate groups

\[ Q_\alpha,\ Q_\beta \dots Q_\sigma, \]

whose greatest common divisor is the group $Q$ of index $j$
considered above.

A function $y_\alpha$ belonging to $Q_\alpha$ is the root
of an equation of degree $s$ in $\Omega$,
$\varphi(y) = 0$, and the remaining roots $y_\beta \dots y_\sigma$
arise from $y_\alpha$ by application of the permutations
of the subsidiary groups $Q_\alpha \pi_\beta, \dots Q_\alpha \pi_\sigma$.
The equation $\chi(u) = 0$ of degree $j$ found above is
the Galois resolvent of the equation $\varphi(y) = 0$. We
have remarked above that the group of the equation $\chi(u) = 0$
of degree $j$ is isomorphic with the group of the permutations
that are produced by $P$ among the systems of the $A, B \dots S$.
But since, on account of the transitivity of $P$,
these systems also pass transitively into one another, $j$
must, by the theorem of \S.\,154, 7., be divisible by $s$, and in the
particular case where $j$ is a prime, we must have $j = s$.

The simultaneous adjunction of all the roots of $\varphi(y)$
has the same effect as the adjunction of the function $\psi$,
namely, to reduce the group of the equation $f(x) = 0$ from $P$ to $Q$.

By the adjunction of $y_\alpha$ a rational factor
$f(x, y_\alpha)$ splits off from $f(x)$, whose roots are the
quantities of the row $A$, and by adjunction of the remaining $y$
one obtains the decomposition

\weq{6}{f(x) = f(x,y_\alpha)\, f(x,y_\beta) \dots f(x,y_\sigma),}

The functions $f(x,y_\alpha), f(x,y_\beta) \dots f(x,y_\sigma)$
are precisely the factors $f_\alpha(x,\psi), f_\beta(x,\psi) \dots
f_\sigma(x,\psi)$ in (2).

It is further to be remarked here that $f(x,y_\alpha)$ is
irreducible in the field $\Omega(y_\alpha)$. For on account of the
assumed irreducibility of $f(x)$ there are in $P$
permutations by which $\alpha$ is carried into every $\alpha_i$,
and these permutations belong to
$Q_\alpha$. If therefore a rational function $\varphi(x,
y_\alpha)$ vanishes for $x = \alpha$,
%CONT>>


% ===== scan leaf 544, printed 519 =====
% !PRINTED 519
\weberpage{519}
%<<CONT
then each $\varphi(\alpha_i, y_\alpha)$ must be equal to zero, that is, $\varphi(x, y_\alpha)$ must
be divisible by $f(x, y_\alpha)$. But it does not follow from this that
$f(x, y_\alpha) = f_\alpha(x, \psi)$ is also irreducible in the field $\Omega(\psi)$.

These are the same theorems as in \S.\,151. It is further shown
there that for $y_\alpha$ one can take a function of $\alpha$ alone,
which we shall not go into again.

It is of interest, however, to follow the connection between the \gesperrt{Galois}
groups of the individual factors $f(x, y_\alpha), f(x, y_\beta), \ldots f(x, y_\sigma)$.
We have already seen above that one obtains the
group of the factor $f(x, y_\alpha)$ in the field $\Omega(\psi)$ by seeking out
all the permutations of the $\alpha$ that are produced by $Q$.
We shall denote this group by $P_\alpha$, and likewise
use the signs $P_\beta, \ldots P_\sigma$ for the other factors.

We select from $P$ any $s - 1$ permutations $\pi_\beta \ldots \pi_\sigma$
of such a kind that by $\pi_\beta$ the row $A$ passes into $B$ etc.,
by $\pi_\sigma$ the row $A$ into $S$; and for the sake of
simplicity we shall now — which is evidently still entirely open to us — assume the
designation of the $\beta, \ldots \sigma$ so chosen that $\pi_\beta, \ldots \pi_\sigma$
produce the following permutations:
\[
\begin{pmatrix} \alpha_0, & \alpha_1 & \ldots & \alpha_{r-1} \\ \beta_0, & \beta_1 & \ldots & \beta_{r-1} \end{pmatrix}, \ldots
\begin{pmatrix} \alpha_0, & \alpha_1 & \ldots & \alpha_{r-1} \\ \sigma_0, & \sigma_1 & \ldots & \sigma_{r-1} \end{pmatrix}.
\]
Now let $\varkappa$ be any element of $Q$ which among the $\alpha$ effects the
permutation
\[
(\alpha) = \begin{pmatrix} \alpha_0, & \alpha_1 & \ldots & \alpha_{r-1} \\ \alpha_{a_0}, & \alpha_{a_1} & \ldots & \alpha_{a_{r-1}} \end{pmatrix}
\]
then
\[
\varkappa_\beta = \pi_\beta^{-1} \varkappa \pi_\beta
\]
is likewise contained in $Q$, and by $\varkappa_\beta$ there is produced among the $\beta$ the
permutation
\[
(\beta) = \begin{pmatrix} \beta_0, & \beta_1 & \ldots & \beta_{r-1} \\ \beta_{a_0}, & \beta_{a_1} & \ldots & \beta_{a_{r-1}} \end{pmatrix},
\]
and hence among the indices of the $\beta$ exactly the same permutation
as by $\varkappa$ among the indices of $\alpha$. If we take
$(\beta)$ and a corresponding $\varkappa_\beta$ as given, then $\pi_\beta \varkappa_\beta \pi_\beta^{-1}$
yields $\varkappa$ again, so that the permutations of the $\beta$
produced by $Q$ agree exactly with those which are produced among the
$\alpha$; and the same holds for the remaining
rows. But with this the theorem is proved.


% ===== scan leaf 545, printed 520 =====
% !PRINTED 520
\weberpage{520}
\gesperrt{1. In the case of an irreducible imprimitive equation
the roots in the individual rows can be
so designated that the partial equations
$f(x, y_\alpha) = 0, f(x, y_\beta) = 0, \ldots f(x, y_\sigma) = 0$ in the field
$\Omega(\psi)$ all receive the same group.}

If one does not choose to designate the roots in the manner
indicated, the groups will at least be
isomorphic.

We further raise the question, which is in a certain measure
answered by the converse of this theorem: \gesperrt{When does a
transitive group have an intransitive normal divisor?}

If $P$ is a transitive permutation group and $Q$ an intransitive
normal divisor of $P$, let
\[
A = \alpha, \alpha_1 \ldots \alpha_{r-1}
\]
be one of the systems whose elements are interchanged by $Q$ only
among themselves, but in such a way that the elements of $A$
are transitively connected by $Q$, that is, that by permutations
from $Q$ every $\alpha$ can pass into every other $\alpha$. Now if
by an element $\pi$ from $P$ the quantities $A$ pass into
\[
B = \beta, \beta_1 \ldots \beta_{r-1}
\]
then these elements must either all coincide with the $\alpha$,
or all be different from the $\alpha$; for
by the group $\pi^{-1} Q \pi$, which by hypothesis is identical
with $Q$, the $\beta$ are interchanged only among themselves, and indeed
in the same manner as the $\alpha$, hence transitively. If therefore only a part
of the $\alpha$ occurred in $B$, then, contrary to the hypothesis,
these $\alpha$ would be interchanged by $Q$ with other quantities not
occurring in $A$.

If the $A$ do not exhaust all the elements to which
the permutations of $P$ extend, then, since $P$
is transitive, there is a permutation $\pi$ by which $A$ is carried over
into a system $B$ entirely different from it. If $\varkappa$ is a permutation
from $Q$, then by the permutation $\pi^{-1} \varkappa \pi$, which likewise
belongs to $Q$, the same permutation is produced among the indices of the $\beta$
as by $\varkappa$ among the indices of the $\alpha$. From this
it follows that by $Q$ the $\beta$ also are permuted only
among themselves. If with $A$ and $B$ the elements are not yet exhausted,
one can form in the same way a third system $C$,
which arises from $A$ by $\pi'$ and from $B$ by $\pi^{-1}\pi'$,
%CONT>>


% ===== scan leaf 546, printed 521 =====
% !PRINTED 521
\weberpage{521}
%<<CONT
and which has no element in common with either $A$ or $B$. One
continues in this way until all elements are exhausted, and comes to
the following theorem:

\gesperrt{2. A transitive group has, only when it is imprimitive,
an intransitive normal divisor different
from the identical group. The systems of
intransitivity of the normal divisor are
systems of imprimitivity of the given
group.} % The German source prints "Imprimivit\"at" for "Imprimitivit\"at"; the intended technical term is translated here.

With respect to the equations, having regard to 1.,
we can also express theorem 2. as follows:

\gesperrt{3. An irreducible equation $f(x) = 0$ splits,
when it becomes reducible through the adjunction of all
the roots of a resolvent or of one root of a normal equation
of degree $j$, into several
irreducible factors of equal degree and
of equal group. The number of these factors
is a divisor of the degree $j$ and is equal to $j$
when $j$ is a prime.}

If, then, an irreducible primitive equation is reduced by the
adjunction of the roots of a resolvent, it splits
into linear factors, i.\,e. it is completely solved.
This case always occurs when the degree of the equation
$f(x) = 0$ is a prime, because such an equation cannot
be imprimitive.

\begin{center}
\rule{3cm}{0.4pt}
\end{center}


% ===== scan leaf 547, printed 522 =====
\weberpage{522}% page number not printed on this page
\pdfbookmark[1]{Fifteenth Section. Cyclic equations.}{bm177}
\pdfbookmark[2]{§ 159. Cubic equations}{bm178}
% !PRINTED none
% The source page between printed pp. 521 and 523 is unnumbered; its complete
% section-opening content is translated below.
\begin{center}
\gesperrt{Fifteenth Section.}

\gesperrt{Cyclic Equations.}

\rule{3cm}{0.4pt}

\S. 159.

\gesperrt{Cubic Equations.}
\end{center}

From the standpoint now gained we shall first consider the
solution of the equations of the third and fourth degree.

The group of the permutations of three figures consists of
six elements, namely of the identical permutation 1, of two
three-term cycles, and three transpositions:
\weq{1}{\begin{array}{ccc} 1 & (0,\,1,\,2) & (0,\,2,\,1) \\ (1,\,2) & (2,\,0) & (0,\,1). \end{array}}
The first three
\[
1, \; (0,\,1,\,2) \quad (0,\,2,\,1)
\]
form the alternating group; it consists of the powers of the
cyclic permutation $\pi = (0, 1, 2)$:
\weq{2}{\pi^0 = 1, \; \pi, \; \pi^2}
and is therefore called the \gesperrt{cyclic group}.

Now let $\alpha, \alpha_1, \alpha_2$ be the three roots of the cubic
equation:
\weq{3}{f(x) = x^3 - ax^2 + bx - c = 0.}
If we take as basis the field $\Omega$, which consists of all rational
functions of the independent variables $a, b, c$,
then (1) is the group of this equation; but if we
\weq{4}{\sqrt{D} = (\alpha - \alpha_1)(\alpha - \alpha_2)(\alpha_1 - \alpha_2)}
%CONT>>


% ===== scan leaf 548, printed 523 =====
% !PRINTED 523
\weberpage{523}

adjoin, in which
\weq{5}{D = + a^2b^2 + 18\,abc - 4a^3c - 4b^3 - 27c^2}
is the discriminant of the equation (3), then the group reduces
to (2).

In the field $\Omega'$, which arises from $\Omega$ by this
adjunction, every root can be expressed rationally in terms of every
other, for we have
\[
\begin{gathered}
\alpha_1 + \alpha_2 = a - \alpha\\
\alpha_1 - \alpha_2 = \dfrac{\sqrt{D}}{f'(\alpha)},
\end{gathered}
\]
hence
\weq{6}{\begin{gathered}
2\alpha_1 = a - \alpha + \dfrac{\sqrt{D}}{f'(\alpha)}\\
2\alpha_2 = a - \alpha - \dfrac{\sqrt{D}}{f'(\alpha)},
\end{gathered}}
and in these the permutations $\pi$, $\pi^2$ can be carried out.
The cubic equation is therefore, after adjunction of $\sqrt{D}$,
its own \gesperrt{Galois} resolvent.

All this remains valid if the domain of rationality is any
special field $\Omega$ in which $a$, $b$, $c$ and
$\sqrt{D}$ are contained, unless $f(x)$ is itself reducible in
$\Omega$, and so has a rational root. For besides itself
and the cyclic group (2), the group (1) has only
intransitive divisors, namely the unit group and three groups
of the type (1, 2).

If we now wish to solve the cubic equation, that is, to reduce it
to a pure cubic equation, we must seek a
function $v$ of the roots which does not itself admit the cyclic
permutation $\pi$, but whose cube does. Such a function
can exist, however, only if the third roots of unity, or, what
is the same thing, $\sqrt{-3}$, is adjoined to the field $\Omega$; for
$v$ must, on application of $\pi$, acquire a third root of unity as
a factor.

By this adjunction no reduction of the group can
occur, because the group (2) has no divisor besides the
unit group, and the reduction to the unit group can
take place not through a quadratic, but only through a cubic
equation (\S.\,157). We therefore set
\[
\varepsilon = \dfrac{-1+\sqrt{-3}}{2}, \qquad
\varepsilon^2 = \dfrac{-1-\sqrt{-3}}{3}
\]
% [sic] The source prints 3 in the denominator of \varepsilon^2; algebra requires 2. The printed value is preserved source-faithfully.
%CONT>>


% ===== scan leaf 549, printed 524 =====
% !PRINTED 524
\weberpage{524}
\pdfbookmark[2]{§ 160. Permutation groups of four elements}{bm179}
%<<CONT
and
\weq{7}{\begin{gathered}
v = \alpha + \varepsilon\alpha_1 + \varepsilon^2\alpha_2 \\
v' = \alpha + \varepsilon^2\alpha_1 + \varepsilon\alpha_2;
\end{gathered}}
then the effect of $\pi$ is that $v$ passes into $\varepsilon^2 v$, and $v'$ into $\varepsilon v'$, while $v^3$, $v'^3$ and $vv'$ remain unchanged, and hence are contained in the field $\Omega'$. There is no difficulty in computing these quantities. If for brevity we set
\[
\begin{gathered}
A = \alpha^2\alpha_1 + \alpha_1^2\alpha_2 + \alpha_2^2\alpha, \\
A' = \alpha_1^2\alpha + \alpha_2^2\alpha_1 + \alpha^2\alpha_2,
\end{gathered}
\]
then
\[
\begin{gathered}
A + A' = ab - 3c \\
A - A' = \sqrt{D},
\end{gathered}
\]
and for $v^3$ one obtains
\[
v^3 = \alpha^3 + \alpha_1^3 + \alpha_2^3 + 6\alpha\alpha_1\alpha_2 + 3\varepsilon A + 3\varepsilon^2 A',
\]
hence
\weq{8}{v^3 = a^3 - \frac{9}{2}ab + \frac{27}{2}c + \frac{3}{2}\sqrt{-3D},}
and likewise
\weq{9}{v'^3 = a^3 - \frac{9}{2}ab + \frac{27}{2}c - \frac{3}{2}\sqrt{-3D}.}

But we also have
\weq{10}{vv' = \alpha^2 + \alpha_1^2 + \alpha_2^2 - \alpha\alpha_1 - \alpha\alpha_2 - \alpha_1\alpha_2 = a^2 - 3b,}
and from (8), (9), (10) one obtains, in agreement with an earlier formula [\S. 47, (8)],
\[
27D = 4(a^2 - 3b)^3 - (2a^3 - 9ab + 27c)^2.
\]
If to (7) we further add the equation
\[
a = \alpha + \alpha_1 + \alpha_2,
\]
then one finds, in agreement with the \gesperrt{Cardan} formula,
\[
\begin{gathered}
3\alpha = a + v + v' \\
3\alpha_1 = a + \varepsilon^2 v + \varepsilon v' \\
3\alpha_2 = a + \varepsilon v + \varepsilon^2 v'.
\end{gathered}
\]

\begin{center}
§. 160.

\gesperrt{Permutation groups of four elements.}
\end{center}

A good example of the \gesperrt{Galois} theory is offered by the equations of the fourth degree, where the circumstances are so simple
%CONT>>


% ===== scan leaf 550, printed 525 =====
% !PRINTED 525
\weberpage{525}
%<<CONT
that everything can easily be surveyed, and yet the most important phenomena of group formation come to light in the process.

From four figures, 0, 1, 2, 3, one can form 24 permutations, and so high, therefore, is the degree of the \gesperrt{Galois} group of the \gesperrt{general} equation of the fourth degree in the field of the rational functions of the coefficients (cf. the end of §. 149). We represent these permutations by their cycles in the following way, where 1 denotes the identical permutation:
\weq{1}{\begin{gathered}
P = 1, \quad (0,1), \quad (0,2), \quad (0,3), \quad (1,2), \quad (1,3), \quad (2,3) \\
(0,1)(2,3), \quad (0,2)(1,3), \quad (0,3)(1,2) \\
(0,1,2), \quad (0,1,3), \quad (0,2,3), \quad (1,2,3) \\
(0,2,1), \quad (0,3,1), \quad (0,3,2), \quad (1,3,2) \\
(0,1,2,3), \quad (0,1,3,2), \quad (0,2,3,1) \\
(0,2,1,3), \quad (0,3,1,2), \quad (0,3,2,1).
\end{gathered}}

We thus have in $P$, besides the identical permutation, six transpositions, eight three-term cycles, six four-term cycles, and three permutations of the form $(0,1)(2,3)$, which we shall call \gesperrt{transposition pairs}.

In $P$ the alternating group of degree 12 is contained as a normal divisor. This group can contain no single transposition and no four-term cyclic permutation, and there thus remain the following:
\weq{2}{\begin{gathered}
Q = 1, \quad (0,1)(2,3), \quad (0,2)(1,3), \quad (0,3)(1,2) \\
(0,1,2), \quad (0,1,3), \quad (0,2,3), \quad (1,2,3) \\
(0,2,1), \quad (0,3,1), \quad (0,3,2), \quad (1,3,2).
\end{gathered}}

We now wish to seek out all the existing divisors of the group $P$. In doing so we leave aside the intransitive groups, which all consist either of three or of six permutations of only three elements, or of the permutations of two pairs of two elements. We therefore seek out only the \gesperrt{transitive divisors} of $P$. Moreover, in the case of a system of conjugate divisors we restrict ourselves to one representative, from which the others can indeed be derived by permutation of the figures (§. 154, 6.).

An essential simplification is brought about by the following two theorems:

1. \gesperrt{If a transitive permutation group of four figures contains a three-term cycle,}
%CONT>>


% ===== scan leaf 551, printed 526 =====
% !PRINTED 526
\weberpage{526}
%<<CONT
\gesperrt{it is the symmetric or the alternating group, and}

2. \gesperrt{if it contains two transpositions with a common element, it is the symmetric group.}

For if the cycle $(0,1,2)$ occurs in a group, and the group is moreover transitive, then it contains a permutation $\pi$ by which 0 is carried into 3. In the three-term cycle $\pi^{-1}(0,1,2)\pi$ the figure 3 therefore certainly occurs; it is thus, say, $(0,1,3)$, and accordingly we have in the group
\[
\begin{gathered}
(0,1,2), \quad (0,1,2)^2 = (0,2,1) \\
(0,1,3), \quad (0,1,3)^2 = (0,3,1) \\
(0,1,2)(0,3,1) = (1,2,3) \\
(1,2,3)^2 = (1,3,2) \\
(0,3,1)(0,1,2) = (0,3,2) \\
(0,3,2)^2 = (0,2,3),
\end{gathered}
\]
that is, all the three-term cycles that exist at all. Accordingly all permutations of the alternating group occur in the group in question (§. 153), and it is therefore either exhausted therewith, or it contains in addition a permutation of the second kind and is then the symmetric group.

But if a transitive group contains the two transpositions $(0,1)$, $(0,2)$, then it also contains
\[
(0,1)(0,2) = (0,1,2),
\]
and with it the whole alternating group, and since it also contains a transposition, it is the symmetric group.

If now, then, a group different from $P$ and $Q$ contains first a transposition $(0,2)$, then, if it is transitive, it must also contain $(1,3)$, and accordingly the whole intransitive group of the fourth degree:
\weq{3}{1, \quad (0,2), \quad (1,3), \quad (0,2)(1,3).}
If, then, the group is to be transitive, further permutations must be added. These can be sought only among the four-term cycles or among the pairs of transpositions. But a four-term cycle in which 0 and 2 adjoin one another cyclically cannot occur, because, combined with $(0,2)$, it would give a three-term cycle, such as:
\[
(0,2)(0,2,3,1) = (0,3,1).
\]


% ===== scan leaf 552, printed 527 =====
% !PRINTED 527
\weberpage{527}
Likewise 1 and 3 cannot adjoin one another cyclically, and there remain therefore only the following:
\weq{4}{(0,1,2,3), \quad (0,3,2,1), \quad (0,3)(1,2), \quad (0,1)(2,3).}
But if one of these permutations is added to (3), then all the others must also be admitted, for we have
\[
\begin{gathered}
(0,2)(0,1,2,3) = (0,3)(1,2) \\
(1,3)(0,1,2,3) = (0,1)(2,3) \\
(0,2)(1,3)(0,1,2,3) = (0,3,2,1),
\end{gathered}
\]
and one can therefore, by means of the transpositions $(0,2)$ and $(1,3)$, derive each of the permutations (4) from each of the others. From the assumption, then, that a transposition occurs at all in a transitive group different from $P$ and $Q$, it follows that this must be the group of the eighth degree
\weq{5}{\begin{gathered}
P_1 = 1, \quad (0,2), \quad (1,3), \quad (0,2)(1,3), \quad (0,1)(2,3), \quad (0,3)(1,2) \\
(0,1,2,3), \quad (0,3,2,1)
\end{gathered}}
or at least a group conjugate with it.

That $P_1$ really is a group one easily recognises by composing any two of its permutations.

There are three distinct groups conjugate with $P_1$, which one obtains by starting, instead of from $(0,2)$, $(1,3)$, from the transpositions $(0,1)$, $(2,3)$ or $(0,3)$, $(1,2)$. The intersection of these three groups is a normal divisor of $P$ and at the same time of $Q$, namely
\weq{6}{Q_1 = 1, \quad (0,2)(1,3), \quad (0,1)(2,3), \quad (0,3)(1,2).}
If a transitive group different from $P$ and $Q$ contains no transposition, then it either contains no four-term cycle and is then identical with (6), or it contains a four-term cycle and then contains also the whole cyclic group, such as
\weq{7}{P_2 = 1, \quad (0,1,2,3), \quad (0,2)(1,3), \quad (0,3,2,1).}
If such a group contains yet a second four-term cycle, say $(0,1,3,2)$, then on account of $(0,1,2,3)(0,1,3,2) = (0,3,1)$ it is the symmetric group. But if, besides $P_2$, it contains also one of the permutations of (6), say $(0,1)(2,3)$, then on account of
\[
(0,1,2,3)(0,1)(2,3) = (1,3)
\]
it is identical with $P_1$.


% ===== scan leaf 553, printed 528 =====
% !PRINTED 528
\weberpage{528}
\pdfbookmark[2]{§ 161. Solution of the biquadratic equations}{bm180}
From $P_2$ likewise three conjugate groups arise, which however have no common divisor other than 1.

\gesperrt{Thus, apart from $P$, $Q$, $P_1$, $Q_1$, $P_2$ and the divisors of $P$ conjugate with $P_1$ and $P_2$, there are no transitive permutation groups of four elements.}

Among the intransitive groups, the group (3):
\[
P_3 = 1, \ (0,2), \ (1,3), \ (0,2)(1,3),
\]
still deserves to be singled out, which likewise gives rise to a system of three conjugate groups.

\begin{center}
\S.\,161.

\gesperrt{Solution of the biquadratic equations.}
\end{center}

The various methods of solution of the biquadratic equations differ from one another in the order in which the various divisors of the group $P$ are employed.

Let $\alpha$, $\alpha_1$, $\alpha_2$, $\alpha_3$ be the roots of the biquadratic equation
\weq{1}{f(x) = x^4 - a_1x^3 + a_2x^2 - a_3x + a_4 = 0.}
We denote the discriminant by $D$ and set
\weq{2}{\sqrt{D} = (\alpha-\alpha_1)(\alpha-\alpha_2)(\alpha-\alpha_3)(\alpha_1-\alpha_2)(\alpha_1-\alpha_3)(\alpha_2-\alpha_3).}
Moreover we recall the expressions for the invariants (\S.\,64):
\weq{3}{A = a_2^2 - 3a_1a_3 + 12a_4}
\weq{4}{B = 27a_1^2a_4 + 27a_3^2 + 2a_2^3 - 72a_2a_4 - 9a_1a_2a_3,}
by means of which $D$ is expressed in the form:
\weq{5}{27D = 4A^3 - B^2.}
If one first adjoins $\sqrt{D}$, the group of the equation reduces to $Q$, and if one then seeks a function belonging to $Q_1$, the group is reduced by it to $Q_1$. Since $Q_1$ has index 3 within $Q$, such a function will depend on a cubic equation, and indeed, since $Q_1$ is a normal divisor of $Q$, on a normal equation. Let us take, for example,
\weq{6}{y = (\alpha-\alpha_1)(\alpha_2-\alpha_3)}
%CONT>>


% ===== scan leaf 554, printed 529 =====
% !PRINTED 529
\weberpage{529}
%<<CONT
as a function belonging to $Q_1$, then the two values conjugate to it are
\weq{7}{\begin{aligned}
y_1 &= (\alpha-\alpha_2)(\alpha_3-\alpha_1)\\
y_2 &= (\alpha-\alpha_3)(\alpha_1-\alpha_2).
\end{aligned}}
These three quantities are the roots of a cubic partial resolvent, whose coefficients are easily formed as invariants of $f$ (cf.\ \S.\,64, 65, where the quantities $y$, $y_1$, $y_2$ were denoted by $U$, $V$, $W$). The equation becomes
\weq{8}{y^3 - Ay + \sqrt{D} = 0.}
Its discriminant is $B^2$, and for the product of the root differences the sign can easily be determined (by comparison of the terms of highest order according to the \gesperrt{Gauss} definition, \S.\,45, or also from \S.\,64). Thus one obtains:
\weq{9}{(y-y_1)(y-y_2)(y_1-y_2) = B.}
The equation (8) is, as already remarked, a normal equation, and it is easy to represent $y_1$, $y_2$ actually in terms of $y$. If one multiplies (9) by $(y_1-y_2)$, one finds, by making use of the formulae $(y-y_1)(y-y_2) = 3y^2 - A$, $(y_1-y_2)^2 = y^2 - 4y_1y_2 = 4A - 3y^2$:
\weq{10}{
\begin{aligned}
B(y_1-y_2) &= -9y^4 + 15Ay^2 - 4A^2\\
 &= 6Ay^2 + 9y\sqrt{D} - 4A^2\\
y_1+y_2 &= -y,
\end{aligned}
}
whereby $y_1$, $y_2$ are determined. If therefore one further adjoins $y$, the group of the biquadratic equation reduces to
\[
Q_1 = 1, \ (0,2)(1,3), \ (0,1)(2,3), \ (0,3)(1,2).
\]
This group is imprimitive (and indeed in three ways), and the equation can now be solved by two square roots. This is best reached in the following way. The three quantities
\weq{11}{
\begin{aligned}
v_1 &= (\alpha+\alpha_1-\alpha_2-\alpha_3)^2\\
v_2 &= (\alpha-\alpha_1+\alpha_2-\alpha_3)^2\\
v_3 &= (\alpha-\alpha_1-\alpha_2+\alpha_3)^2
\end{aligned}
}
all admit the permutations of the group $Q_1$ and are therefore rationally representable in terms of $y$. One obtains this representation as follows: We have
%CONT>>


% ===== scan leaf 555, printed 530 =====
% !PRINTED 530
\weberpage{530}
%<<CONT
\[
\begin{aligned}
v_1 &= a_1^2 - 4(\alpha+\alpha_1)(\alpha_2+\alpha_3)\\
y_1-y_2 &= (\alpha+\alpha_1)(\alpha_2+\alpha_3) - 2(\alpha\alpha_1+\alpha_2\alpha_3)\\
a_2 &= (\alpha+\alpha_1)(\alpha_2+\alpha_3) + \alpha\alpha_1 + \alpha_2\alpha_3,
\end{aligned}
\]
and thus, if one forms $v_2$, $v_3$ in the same way:
\weq{12}{
\begin{aligned}
3v_1 &= 3a_1^2 - 8a_2 - 4y_1 + 4y_2\\
3v_2 &= 3a_1^2 - 8a_2 - 4y_2 + 4y\\
3v_3 &= 3a_1^2 - 8a_2 - 4y + 4y_1.
\end{aligned}
}
But there is yet another relation, which results from the fact that the product
\[
(\alpha+\alpha_1-\alpha_2-\alpha_3)(\alpha-\alpha_1+\alpha_2-\alpha_3)(\alpha-\alpha_1-\alpha_2+\alpha_3)
\]
is a symmetric function of the $\alpha$, which one easily finds to be equal to
\[
a_1^3 - 4a_1a_2 + 8a_3
\]
We therefore make the signs of the three root quantities depend on one another by the relation
\weq{13}{\sqrt{v_1}\sqrt{v_2}\sqrt{v_3} = a_1^3 - 4a_1a_2 + 8a_3}
and set:
\weq{14}{
\begin{aligned}
\sqrt{v_1} &= \alpha+\alpha_1-\alpha_2-\alpha_3\\
\sqrt{v_2} &= \alpha-\alpha_1+\alpha_2-\alpha_3\\
\sqrt{v_3} &= \alpha-\alpha_1-\alpha_2+\alpha_3\\
a_1 &= \alpha+\alpha_1+\alpha_2+\alpha_3.
\end{aligned}
}
From this it then follows:
\weq{15}{
\begin{aligned}
4\alpha &= a_1+\sqrt{v_1}+\sqrt{v_2}+\sqrt{v_3}\\
4\alpha_1 &= a_1+\sqrt{v_1}-\sqrt{v_2}-\sqrt{v_3}\\
4\alpha_2 &= a_1-\sqrt{v_1}+\sqrt{v_2}-\sqrt{v_3}\\
4\alpha_3 &= a_1-\sqrt{v_1}-\sqrt{v_2}+\sqrt{v_3}.
\end{aligned}
}
If one gives $\sqrt{D}$ in (8) the opposite sign, and replaces $y$ by $-y$, then by (10) one must interchange $y_1$, $y_2$ with $-y_2$, $-y_1$. Thus $v_1$ remains unchanged, and $v_2$ and $v_3$ are interchanged with one another. If one takes for $y$ another root of the equation (8), say $y_1$, then on account of (9) $y_1$, $y_2$ must be replaced by $y_2$, $y$, and the quantities $v_1$, $v_2$, $v_3$ undergo a cyclic permutation, whereby the quantities $\alpha_1$, $\alpha_2$, $\alpha_3$ in (15) are also cyclically permuted. If, finally, one chooses the signs of $\sqrt{v_1}$, $\sqrt{v_2}$, $\sqrt{v_3}$ differently, but in such a way that (13) always remains satisfied, then certain permutations likewise occur among the quantities (15).


% ===== scan leaf 556, printed 531 =====
% !PRINTED 531
\weberpage{531}

\gesperrt{Whatever values one may assign to the many-valued algebraic quantities that occur in the solution, the expressions (15) thus always represent the roots of our biquadratic equation in some order.}

With regard to the remaining ways of solving the biquadratic equation we can be briefer.

If we first adjoin not $\sqrt{D}$ but a function belonging to the group $P_1$, we obtain a cubic resolvent which is not a normal equation. For this function we can choose, say, $y^2$, which satisfies the following cubic equation in that quantity:
\weq{16}{y^6 - 2Ay^4 + A^2y^2 - D = 0}
Better still, one takes as roots of the cubic resolvent
\weq{17}{z = y_1 - y_2, \quad z_1 = y_2 - y, \quad z_2 = y - y_1.}
Then $z$ belongs to the group $P_1$, as already follows from the second equation (10), or as may also be seen directly. From (8) and (9) one obtains for $z$ the cubic equation
\weq{18}{z^3 - 3Az + B = 0.}
If we further set
\weq{19}{u = \alpha\alpha_1 + \alpha_2\alpha_3}
then
\[
z = a_2 - 3u
\]
and hence obtain the resolvent for $u$:
\weq{20}{u^3 - a_2u^2 + (a_1a_3 - 4a_4)u + 4a_2a_4 - a_1^2a_4 - a_3^2 = 0.}

But if one adjoins not one but all the roots of equation (18) or (20), one arrives again at the group $Q_1$. The quantities $v_1, v_2, v_3$ are by (12) expressed immediately in terms of $z, z_1, z_2$, and one can therefore also apply the formulae (15) for the representation of the $\alpha$ by this procedure.

One can further arrive at the solution of the biquadratic equation by first adjoining a function of the roots belonging to the cyclic group $P_2$ (a cyclic function), as for example
\weq{21}{w = \alpha\alpha_1^2 + \alpha_1\alpha_2^2 + \alpha_2\alpha_3^2 + \alpha_3\alpha^2.}
This function is six-valued and is therefore the root of an equation of the sixth degree. But this equation of the sixth degree is, as one easily sees, imprimitive, and can be
%CONT>>


% ===== scan leaf 557, printed 532 =====
% !PRINTED 532
\weberpage{532}

%<<CONT
to an equation of the third degree and to two square roots. This
equation of the sixth degree is a total resolvent, and two of its
roots suffice for the rational representation of the roots $\alpha$. The
form of these resolvents does not turn out simple.

Finally, one may also aim, by the adjunction of a function $\xi$
belonging to the group $P_3$, to make the function $f(x)$ directly
reducible and to decompose it into two quadratic factors. Such a
function $\xi$ is likewise the root of an equation of the sixth degree,
which, however, also reduces by imprimitivity to the third and
second degree. If one takes, e.\,g.,
\weq{22}{\xi = \alpha + \alpha_2 - \tfrac{1}{2}a_1 = \tfrac{1}{2}\sqrt{v_2},}
then of the six roots two at a time are equal and opposite, so
that a cubic equation for $\xi^2$ results. By means of one of these
quantities $\xi$ the quadratic factors of $f(x)$ can then be represented
rationally.

If, to simplify the formulae, we assume $a_1 = 0$ and
\weq{23}{f(x) = x^4 + ax^2 - bx + c,}
then
\[
\begin{aligned}
(\alpha\alpha_2 - \alpha_1\alpha_3)\,\xi &= b \\
\alpha\alpha_2 + \alpha_1\alpha_3 &= \xi^2 + a,
\end{aligned}
\]
hence
\weq{24}{2\,\alpha\alpha_2 = \frac{b}{\xi} + \xi^2 + a, \qquad 2\,\alpha_1\alpha_3 = -\frac{b}{\xi} + \xi^2 + a,}
so that the two quadratic factors of $f(x)$ become the following:
\weq{25}{
\begin{aligned}
x^2 - \xi x + \tfrac{1}{2}\Bigl(\dfrac{b}{\xi} + \xi^2 + a\Bigr) &= 0 \\
x^2 + \xi x + \tfrac{1}{2}\Bigl(-\dfrac{b}{\xi} + \xi^2 + a\Bigr) &= 0,
\end{aligned}
}
and for $\xi^2$ there results from (24), by multiplication, the cubic
equation
\[
4c = (\xi^2 + a)^2 - \frac{b^2}{\xi^2}
\]
or
\weq{26}{\xi^6 + 2a\xi^4 + (a^2 - 4c)\xi^2 - b^2 = 0.}
If one takes for $\xi$ any root of this equation of the sixth
degree, then each of the equations (25) gives a pair of roots
of $f(x) = 0$.


% ===== scan leaf 558, printed 533 =====
% !PRINTED 533
\weberpage{533}
\pdfbookmark[2]{§ 162. Abelian equations}{bm181}

\begin{center}
\S.\,162.

\gesperrt{Abelian Equations.}
\end{center}

The degree of a transitive permutation group of $m$ figures
is always divisible by $m$, and therefore never smaller than $m$.
For if $P$ is such a group and $Q_0$ is the totality of the permutations
of $P$ that leave the figure 0 unchanged, then $Q_0$ is also
a group, and therefore a divisor of $P$. Now, on account of the
assumed transitivity, one can find in $P$ a system of
permutations $\pi_1, \pi_2 \ldots \pi_{m-1}$ that carry 0 into 1, 0 into 2, \ldots
% [sic] authority prints Q\pi_1 here, omitting the subscript 0 expected from Q_0.
0 into $m-1$, and then $Q\pi_1$ is the system of all those
permutations of $P$ that change 0 into 1. Accordingly

\weq{1}{P = Q_0 + Q_0\pi_1 + Q_0\pi_2 + \cdots + Q_0\pi_{m-1},}

and so the degree of $P$ is equal to the product of $m$ and the degree
of $Q_0$.

\gesperrt{The Galois group of an irreducible equation is therefore never of lower degree than the equation itself.}

A normal equation we have earlier defined by the stipulation
that it be irreducible and that each of its
roots should be rationally expressible in terms of every other
(\S.\,145). From this it follows that the group of a normal equation
must reduce to the identical group when
one root is adjoined; and conversely, an equation
whose group is so constituted is, provided it is at the same time irreducible,
always a normal equation.

Now $P$ reduces, by adjunction of the root $\alpha_0$, to
$Q_0$, by $\alpha_1$ to $\pi_1^{-1}Q_0\pi_1$, etc., and so if $P$ is to be the group
of a normal equation, it is necessary and sufficient
that $Q_0$ be the unit group, and hence that the group be exhausted by
% [sic] authority prints an unindexed \pi before \pi_1,...,\pi_{m-1}; none is introduced here, and context calls for the identity representative.
$\pi, \pi_1, \pi_2, \ldots \pi_{m-1}$. We therefore have:

\gesperrt{In order that an irreducible equation be a normal equation, it is necessary and sufficient that the degree of the group agree with the degree of the equation.}

We first consider here the special class of equations that includes
the cyclotomic equations first solved by \gesperrt{Gauss},
%CONT>>


% ===== scan leaf 559, printed 534 =====
% !PRINTED 534
\weberpage{534}

%<<CONT
which \gesperrt{Abel} taught how to solve in general, and which we shall therefore, after him, call \gesperrt{Abelian equations}\footnote{\gesperrt{Abel}, Mémoire sur une classe d'équations résolubles algébriquement. Crelle's Journal f. Mathematik, Bd.~4, 1829. Oeuvres complètes, nouvelle édition, 1881, Bd.~1, S.~418.}.

\gesperrt{An equation of degree $m$, $F(x) = 0$, with the roots $\alpha, \alpha_1, \alpha_2 \ldots \alpha_{m-1}$, is called an Abelian equation if every root is rationally expressible in terms of one of them, $\alpha$, and if, when}
\weq{2}{\alpha_1 = \Theta_1(\alpha), \quad \alpha_2 = \Theta_2(\alpha) \ldots \alpha_{m-1} = \Theta_{m-1}(\alpha)}
\gesperrt{are these rational expressions, the condition}
\weq{3}{\Theta_h\Theta_k(\alpha) = \Theta_k\Theta_h(\alpha)}
\gesperrt{holds for every pair of these functions.}

Here the symbol $\Theta_h\Theta_k(\alpha)$ means that the function $\Theta_h(x)$ is to be formed for the argument $x = \Theta_k(\alpha)$. This whole definition refers, of course, to a definite field $\Omega$ assumed as rational. To these equations certainly belong, when $\Omega$ is the field of rational numbers, the equations by which the roots of unity are determined. For every $m$th root of unity is, if $r$ is a primitive one among them, contained in the form $\Theta_h(r) = r^h$, and we have
\[
\Theta_h\Theta_k(r) = \Theta_k\Theta_h(r) = r^{hk}.
\]
Likewise, when the field $\Omega$ contains the $m$th roots of unity $r$, the \gesperrt{pure equation}
\[
x^m - a = 0,
\]
in which $a$ belongs to the field $\Omega$, belongs to the \gesperrt{Abelian equations}. For if $\alpha$ is one of its roots, then all the roots are contained in the form
\[
\alpha_h = r^h\alpha
\]
we have therefore $\Theta_h(x) = r^hx$, and hence
\[
\Theta_h\Theta_k(x) = \Theta_k\Theta_h(x) = r^{h+k}x.
\]
If the function $F(x)$ is not irreducible, then it has a definite irreducible factor $\varphi(x)$ which has the root $\alpha$, and among the roots of $\varphi(x) = 0$ the relations expressed by (2) and (3) likewise hold. Hence $\varphi(x) = 0$ is, by \S.\,145, a \gesperrt{Galois} resolvent of
%CONT>>


% ===== scan leaf 560, printed 535 =====
% !PRINTED 535
\weberpage{535}
%<<CONT
$F(x) = 0$, and when $\varphi(x) = 0$ is solved, all the roots
of $F(x)$ are thereby known. It therefore suffices if we confine
ourselves to the consideration of \gesperrt{irreducible Abelian equations}.

\quad The \gesperrt{Galois} group of an \gesperrt{Abelian} equation, whether it
be irreducible or not (provided only that its roots are distinct
from one another), has the property that in the composition
of its permutations the commutative law holds; it is
therefore a \gesperrt{commutative group}.

\quad Namely, let
\weq{4}{\alpha,\ \alpha_1 = \Theta_1(\alpha),\ \alpha_2 = \Theta_2(\alpha) \ldots \alpha_{m-1} = \Theta_{m-1}(\alpha)}
be the roots of $F(x)$, and
\weq{5}{\alpha,\ \alpha' = \Theta'(\alpha),\ \alpha'' = \Theta''(\alpha) \ldots}
the roots of the irreducible factor $\varphi(x)$ of $F(x)$
contained among them.

\quad Since, as already remarked, $\varphi(x) = 0$ is a \gesperrt{Galois} resolvent,
the group of the equation consists of the substitutions
of the field $\Omega(\alpha)$, that is, of the substitutions
\weq{6}{\sigma = (\alpha, \alpha),\ \sigma' = (\alpha, \alpha'),\ \sigma'' = (\alpha, \alpha'') \ldots}

\quad But this group obeys the commutative law; for
let
\[
\begin{aligned}
\sigma' &= (\alpha, \alpha') = [\alpha, \Theta'(\alpha)]\\
\sigma'' &= (\alpha, \alpha'') = [\alpha, \Theta''(\alpha)];
\end{aligned}
\]
then by \S.\,147
\[
\begin{aligned}
\sigma'\sigma'' &= [\alpha, \Theta'(\alpha)][\Theta'(\alpha), \Theta'\Theta''(\alpha)] = [\alpha, \Theta'\Theta''(\alpha)]\\
\sigma''\sigma' &= [\alpha, \Theta''(\alpha)][\Theta''(\alpha), \Theta''\Theta'(\alpha)] = [\alpha, \Theta''\Theta'(\alpha)].
\end{aligned}
\]
On account of (3) we therefore have $\sigma'\sigma'' = \sigma''\sigma'$, where $\sigma', \sigma''$ may be any
two of the substitutions $\sigma$. Consequently the group of the
substitutions of the field $\Omega(\alpha)$, and with it also the isomorphic
permutation group of the equation $F(x) = 0$, is \gesperrt{commutative}.

\quad The converse now also holds, but here the \gesperrt{irreducibility}
must be presupposed.

\quad \gesperrt{An irreducible equation $\varphi(x) = 0$ with commutative group is an Abelian equation.}

\quad Let $\alpha, \alpha_1, \alpha_2 \ldots \alpha_{m-1}$ be the roots of $\varphi(x) = 0$ and
$P$ the group of this equation, which is transitive because of the irreducibility of
$\varphi(x)$, and which we now moreover assume to be commutative.
Let $Q$ be the divisor of $P$ that leaves the element $0$
%CONT>>


% ===== scan leaf 561, printed 536 =====
% !PRINTED 536
\weberpage{536}
%<<CONT
fixed. If then $\pi_i$ is a permutation that carries the figure 0 into $i$, then $\pi_i^{-1}\,Q\,\pi_i$ is the group of the permutations in $P$ that do not alter the figure $i$. But since in every composition of permutations from $P$ the components may be interchanged, we have
\[
\pi_i^{-1}\,Q\,\pi_i = \pi_i^{-1}\pi_i\,Q = Q,
\]
that is, the group $Q$ also leaves the figure $i$ unchanged. On account of the transitivity of $P$, however, $i$ may denote any of the figures $1,2\ldots m-1$; consequently $Q$ consists of the single identical permutation, and in $P$ there is, besides the identical one, no permutation that leaves a figure unchanged. But if one adjoins \gesperrt{a} root $\alpha$, the group $P$ reduces to the unit group, or the equation is solved; $\varphi(x)=0$ is a normal equation and hence its own \gesperrt{Galois} resolvent. If $\alpha_k=\Theta_k(\alpha)$, then the \gesperrt{Galois} group of this equation consists of the substitutions:
\[
\sigma_k = [\alpha,\Theta_k(\alpha)],
\]
and we have
\[
% [sic] authority prints x in the second composite; the parallel expression and the next line require \alpha.
\sigma_h\sigma_k = [\alpha,\Theta_h\Theta_k(\alpha)],\quad \sigma_k\sigma_h = [\alpha,\Theta_k\Theta_h(x)].
\]
Since the group is to be commutative, we must have
\[
\Theta_h\Theta_k(\alpha) = \Theta_k\Theta_h(\alpha)
\]
that is, $\varphi(x)=0$ is an \gesperrt{Abelian equation}. This is the reason why the commutative groups are also called \gesperrt{Abelian} groups (\S.\,148). It therefore also follows from the theorem given above on the group of normal equations that in a transitive \gesperrt{Abelian} permutation group the number of permutations agrees with the number of the permuted figures.

Here we have presupposed the irreducibility of the given equation. If we wish to take into consideration also reducible equations with commutative group, we need only put in place of the equation itself one of its \gesperrt{Galois} resolvents.

It remains, finally, to remark that in a commutative group every divisor is normal, since we always have $\pi^{-1}\varkappa\pi = \varkappa$ when $\pi$ and $\varkappa$ are arbitrary elements of a commutative group.


% ===== scan leaf 562, printed 537 =====
% !PRINTED 537
\weberpage{537}
\pdfbookmark[2]{§ 163. Reduction of the Abelian equations to cyclic ones}{bm182}
\begin{center}
\S.\,163.

\gesperrt{Reduction of Abelian equations to cyclic ones.}
\end{center}

We have already seen above that in a transitive \gesperrt{Abelian} group no permutation other than the identical one occurs which leaves a figure unchanged. But from this there follows yet another theorem. Let us assume that a permutation $\pi$ of such a group is decomposed into its cycles, and that the cycle which contains the fewest terms is an $r$-term one. Then $\pi^r$ will leave the terms of this cycle unchanged, and must therefore be the identical permutation. But from this it follows that all the remaining cycles of $\pi$, if there are any others, must likewise consist of $r$ terms; hence:

\gesperrt{1. A permutation of a transitive Abelian group contains only cycles with the same number of terms.}

The number $r$ of the terms of a cycle must be a divisor of $m$, if $m$ is the degree of the group, and if $m=rs$, then $s$ is the number of the $r$-term cycles of which $\pi$ consists.

If now $P$ is the group of an \gesperrt{Abelian} equation $F(x)=0$, we take any non-identical permutation $\pi$ out of $P$ and decompose it into its cycles
\weq{1}{\pi = \gamma\gamma_1\gamma_2\cdots\gamma_{s-1},}
in which each of the cycles $\gamma$ relates to $r$ roots of the equation $F(x)=0$. We shall arrange these roots in such a way that
\weq{2}{\begin{aligned}
\gamma &= (\alpha,\alpha_1,\alpha_2\ldots\alpha_{r-1})\\
\gamma_1 &= (\beta,\beta_1,\beta_2\ldots\beta_{r-1})\\
&\ \,\vdots\\
\gamma_{s-1} &= (\sigma,\sigma_1,\sigma_2\ldots\sigma_{r-1})
\end{aligned}}
holds.

If $\pi_1$ is any permutation of $P$, then on account of the permutability
\weq{3}{\pi_1^{-1}\pi\pi_1 = \pi.}

By the theorem on the formation of the permutations $\pi_1^{-1}\pi\pi_1$ (\S.\,154, 6.), $\pi$ must therefore not be changed when the permutations $\pi_1$ are carried out in the cycles of $\pi$. But since the cycles are completely determined, apart from their ar-
%CONT>>


% ===== scan leaf 563, printed 538 =====
% !PRINTED 538
\weberpage{538}
%<<CONT
rangement and their initial element, it follows that by the application of any permutation $\pi_1$ from $P$ the elements of the individual cycles $\gamma$ are not separated from one another, but can only be interchanged among themselves (cyclically), and moreover the cycles can be interchanged with one another. The group is therefore, if $s>1$, \gesperrt{imprimitive}.

A rational function of the arguments $\alpha,\alpha_1,\ldots\alpha_{r-1}$ which does not change its value when the $\alpha$ are subjected to the cyclic permutation $\gamma$ and its repetitions is called a \gesperrt{cyclic function} of the $\alpha$. If one denotes by
\weq{4}{\omega = \psi(\alpha,\alpha_1\ldots\alpha_{r-1})}
such a cyclic function, and by
\weq{5}{\omega,\ \omega_1,\ \omega_2\ldots\omega_{s-1}}
the conjugate values of $\omega$, setting thus e.\,g.
\weq{6}{\omega_1 = \psi(\beta,\beta_1\ldots\beta_{r-1}),}
then these quantities are, by \S.\,158, the roots of an irreducible equation of degree $s$:
\weq{7}{\Phi(t) = 0,}
whose group one obtains by seeking out the permutations of the quantities (5) produced by $P$.

But one obtains the permutation effected by $\pi_1\pi_2$ among the quantities (5) by composing the two permutations produced by $\pi_1$ and $\pi_2$ singly, and the group of the permutations of (5) is therefore also commutative. Hence $\Phi(t)=0$ is an \gesperrt{Abelian} equation of degree $s$.

If one adjoins $\omega$, then $F(x)$ splits into $s$ factors of degree $r$:
\[
F(x) = F(x,\omega)F(x,\omega_1)\ldots F(x,\omega_{s-1}),
\]
of which the first, $F(x,\omega)$, has the roots $\alpha,\alpha_1\ldots\alpha_{r-1}$, and the group of the equation $F(x,\omega)=0$ consists solely of the period of the cyclic permutation $\gamma$.

We shall call an equation whose group consists of a single cycle and its repetitions a \gesperrt{cyclic equation}, so that the cyclic equations are the simplest special case of the \gesperrt{Abelian} equations. We have then proved that the solution of every \gesperrt{Abelian} equation is reduced to the solution of an \gesperrt{Abelian} equation of lower degree and to the solution of a series
%CONT>>


% ===== scan leaf 564, printed 539 =====
% !PRINTED 539
\weberpage{539}
of cyclic equations. This theorem can again be applied to the auxiliary equation of degree $s$, and one can continue in this way until this auxiliary equation reduces to the first degree. With this, then, one obtains the result:

\gesperrt{The solution of an Abelian equation can always be reduced to the solution of a sequence of cyclic equations whose degrees are divisors of the degree of the given equation.}

It is not necessary to include irreducibility in the definition of cyclic equations. We may therefore frame the definition generally as follows:

\gesperrt{An equation of degree $m$, $F(x) = 0$, with $m$ distinct roots is called a cyclic equation in the field $\Omega$ if its roots $\alpha, \alpha_1, \alpha_2 \ldots \alpha_{m-1}$ are not rational, but can be so arranged that the cyclic functions of the roots are rational in $\Omega$.}

If
\weq{7}{\pi = (\alpha, \alpha_1, \alpha_2 \ldots \alpha_{m-1})}
holds, then every function that admits the permutations of the cyclic group
\weq{8}{C = 1, \pi, \pi^2, \pi^3 \ldots \pi^{m-1}}
must be contained in $\Omega$.
The \gesperrt{Galois} group of a cyclic equation is either the period $C$ itself, and then the equation is irreducible, or it is a divisor of $C$; in that case, if $e$ and $f$ are two integer factors of $m$ and
\[
m = ef
\]
it consists of the permutations
\weq{9}{C_e = 1, \pi^e, \pi^{2e} \ldots \pi^{(f-1)e},}
and the cyclic equation splits into $e$ factors of degree $f$. For if we take a function $\psi(\alpha, \alpha_1 \ldots \alpha_{m-1})$ \gesperrt{belonging} to the group $C$, then by hypothesis it is equal to a quantity in $\Omega$, which we denote by $a$. No permutation not contained in $C$ is then applicable to the rational equation $\psi = a$, and consequently the group of the equation can contain no substitutions other than those which occur in $C$. Now if $\pi$ occurs in the group of the equation, then that group is identical with $C$, and since $C$ is transitive, the equation is irreducible. But if $\pi^e$ is the lowest power of $\pi$ which occurs in the
%CONT>>


% ===== scan leaf 565, printed 540 =====
% !PRINTED 540
\weberpage{540}
%<<CONT
group of the equation, then $C_e$ is this group. But $C_e$ is intransitive when $e < m$, and the equation is reducible.

The reduction carried out above can also be applied to the cyclic equations whose degree is not a prime.

For if $m = ef$ is any decomposition of $m$ into two factors, the permutation $\pi^e$ splits into $e$ cycles $\gamma$ of $f$ terms each:
\weq{10}{
\begin{aligned}
\gamma &= (\alpha, \alpha_e, \alpha_{2e} \ldots \alpha_{(f-1)e}), \\
\gamma_1 &= (\alpha_1, \alpha_{e+1}, \alpha_{2e+1} \ldots \alpha_{(f-1)e+1}), \\
&\cdots\cdots\cdots\cdots\cdots\cdots\cdots\cdots \\
% [sic] authority prints c-1 here; mathematical cycle indexing requires e-1.
\gamma_{c-1} &= (\alpha_{e-1}, \alpha_{2e-1}, \alpha_{3e-1} \ldots \alpha_{m-1}).
\end{aligned}
}
Now we determine a function $\psi$ belonging to the first of these cycles, and set
\weq{11}{
\begin{aligned}
\eta &= \psi(\alpha, \alpha_e, \alpha_{2e} \ldots \alpha_{(f-1)e}), \\
\eta_1 &= \psi(\alpha_1, \alpha_{e+1}, \alpha_{2e+1} \ldots \alpha_{(f-1)e+1}), \\
&\cdots\cdots\cdots\cdots\cdots\cdots\cdots\cdots \\
\eta_{e-1} &= \psi(\alpha_{e-1}, \alpha_{2e-1}, \alpha_{3e-1} \ldots \alpha_{m-1}).
\end{aligned}
}
These quantities are all different from one another, but under a cyclic permutation of their arguments they remain unchanged. By application of the permutation $\pi$ the quantities
\[
\eta, \eta_1, \eta_2 \ldots \eta_{e-1}
\]
pass cyclically into one another, and by the hypothesis their cyclic functions, and consequently also their symmetric functions, are therefore rational. They are thus the roots of a cyclic equation of degree $e$, while $F(x)$ splits into $e$ factors of degree $f$
\[
F(x) = F(x, \eta)\, F(x, \eta_1) \ldots F(x, \eta_{e-1})
\]
each of which yields a cyclic equation of degree $f$ for the roots of one of the cycles $\gamma$.

For if we set, say,
\[
F_1 = (x - \alpha)(x - \alpha_e) \ldots (x - \alpha_{(f-1)e}),
\]
then this function admits the permutations of the period of $\gamma$. Hence it also admits the permutations of the group which, after adjunction of $\eta$, becomes the \gesperrt{Galois} group of our equation, and which indeed can consist only of powers of $\gamma, \gamma_1 \ldots \gamma_{e-1}$. Hence, by the theorem of \gesperrt{Lagrange} [\S.\,155, (2)], $F_1$ is rationally expressible in terms of $\eta$, so $F_1 = F(x, \eta)$. The equation
%CONT>>


% ===== scan leaf 566, printed 541 =====
% !PRINTED 541
\weberpage{541}

%<<CONT
$F(x, \eta) = 0$ is, however, again cyclic in $\Omega(\eta)$, since the cyclic functions of its roots belong to this field.

The cyclic equations, like all \gesperrt{Abel}ian equations, have the property that every root is rationally expressible in terms of every other. Here these expressions can be arranged cyclically in the following manner. If $\alpha, \alpha_1, \alpha_2 \ldots \alpha_{m-1}$ are the roots of the cyclic equation $F(x) = 0$, then the polynomial of degree $(m-1)$ in $x$
\[
F(x)\left(\frac{\alpha_1}{x-\alpha} + \frac{\alpha_2}{x-\alpha_1} + \cdots + \frac{\alpha}{x-\alpha_{m-1}}\right) = \Psi(x)
\]
is unchanged by the permutation $\pi$ and therefore contained in $\Omega$. If we set $x = \alpha, \alpha_1 \ldots \alpha_{m-1}$ therein and introduce the sign
\[
\frac{\Psi(x)}{F'(x)} = \Theta(x)
\]
it follows that
\weq{11}{\alpha_1 = \Theta(\alpha), \alpha_2 = \Theta(\alpha_1), \ldots \alpha_{m-1} = \Theta(\alpha_{m-2}), \alpha = \Theta(\alpha_{m-1}),}
and this holds, whether $F(x)$ be reducible or irreducible, provided only that $F(x)$ and $F'(x)$ have no common divisor.

The solution of the cyclic equations of degree $m$ is hereby made dependent on the solution of cyclic equations whose degrees are the prime factors of $m$. If, for example, $m$ is a power of $2$, the solution is accomplished by a sequence of square roots. In this way \gesperrt{Gauss} first treated the cyclotomic equations\footnote{\gesperrt{Gauss}, Disquisitiones arithmeticae, Sectio VII.}.

We finally attach here the remark, important for what follows, that for a prime degree the concepts of the \gesperrt{normal equation and of the cyclic equation} coincide. For in any case a cyclic equation of prime degree, since it is irreducible, is a normal equation. And if, conversely, the degree $n$ of a normal equation, which is at the same time the degree of the group of this equation, is a prime, and $\pi$ is a non-identical permutation of this group, then the degree of $\pi$, which must indeed be a divisor of $n$, is equal to $n$, and the group of the equation is $1, \pi, \pi^2 \ldots \pi^{n-1}$, hence cyclic.


% ===== scan leaf 567, printed 542 =====
% !PRINTED 542
\weberpage{542}
\pdfbookmark[2]{§ 164. Resolvents of Lagrange}{bm183}

\begin{center}
\S.\,164.\\
\gesperrt{Resolvents of Lagrange.}
\end{center}

The method of solving cyclic equations with which we now wish to become acquainted is applicable uniformly to prime degrees and to composite degrees. For this purpose one makes use of certain expressions, known under the name of the \gesperrt{resolvents of Lagrange}\footnote{\gesperrt{Lagrange}, Réflexions etc., see p. 508. Earlier we understood by resolvents solving equations; here they are solving functions.}, which are of great use in all investigations concerning the algebraic solution of equations.

Let $F(x) = 0$ be an equation with the roots $\alpha, \alpha_1, \alpha_2 \ldots \alpha_{m-1}$. We denote by $\varepsilon$ any $m$th root of unity and introduce the notation
\weq{1}{(\varepsilon, \alpha) = \alpha + \varepsilon\alpha_1 + \varepsilon^2\alpha_2 + \cdots + \varepsilon^{m-1}\alpha_{m-1}.}
The sums so defined are those which one calls the \gesperrt{Lagrange resolvents}. When these functions are known for all $m$th roots of unity $\varepsilon$, the equation itself is also solved, for by \S.\,133 (6) we have
\weq{2}{\overset{\varepsilon}{\Sigma}\,\varepsilon^k = m \ \text{or} = 0,}
according as $k$ is divisible by $m$ or not, and from this it follows
\weq{3}{m\alpha = \overset{\varepsilon}{\Sigma}(\varepsilon, \alpha),}
in which the sums extend over all $m$th roots of unity $\varepsilon$. One can also express the other roots in the same way:
\weq{4}{m\alpha_k = \overset{\varepsilon}{\Sigma}\varepsilon^{-k}(\varepsilon, \alpha),}
so that in fact everything is reduced to the knowledge of $\varepsilon$ and of the functions $(\varepsilon, \alpha)$.

The sums (2) and (3) may also be represented in a somewhat different manner, since all the $m$th roots of unity are powers of a primitive one among them. If, then, one understands by $\varepsilon$ a fixed primitive $m$th root of unity, and by $\lambda$ an index which runs through a complete system of residues with respect to the modulus $m$, say the
%CONT>>


% ===== scan leaf 568, printed 543 =====
% !PRINTED 543
\weberpage{543}

%<<CONT
sequence of numbers $0, 1 \ldots m-1$, then for the equations (2), (3) we can set
\[
m\alpha = \overset{\lambda}{\Sigma}(\varepsilon^\lambda, \alpha), \quad m\alpha_k = \overset{\lambda}{\Sigma}\varepsilon^{-\lambda k}(\varepsilon^\lambda, \alpha).
\]
We investigate these resolvents $(\varepsilon, \alpha)$ first as functions of the $m$ \gesperrt{independent variables} $\alpha$, and, for the sake of a simpler presentation of the formulae, we shall agree that $\alpha_m = \alpha_0 = \alpha$, and in general that $\alpha_h = \alpha_k$ shall hold whenever $h \equiv k \pmod{m}$. We then obtain the whole system of the variables $\alpha$ when in $\alpha_h$ we let the index run through a complete system of residues to the modulus $m$.

For these resolvents the following theorems now hold.

1. \gesperrt{If one applies to the indices of the $\alpha$ the cyclic permutation $\pi = (0, 1, \ldots m-1)$, then $(\varepsilon, \alpha)$ passes into $\varepsilon^{-1}(\varepsilon, \alpha)$, and by the permutation $\pi^k$, $(\varepsilon, \alpha)$ passes into $\varepsilon^{-k}(\varepsilon, \alpha)$.}

This is shown immediately by the definition (1) of the resolvents.

We further understand by $\nu$ an arbitrary positive exponent, and form $(\varepsilon, \alpha)^\nu$ according to the polynomial theorem. In the developed power we set $\varepsilon^m = 1$ and then arrange according to powers of $\varepsilon$. There then results an expression of the form
\weq{5}{\begin{split}(\varepsilon, \alpha)^\nu &= A_0^{(\nu)} + \varepsilon A_1^{(\nu)} + \varepsilon^2 A_2^{(\nu)} + \cdots + \varepsilon^{m-1}A_{m-1}^{(\nu)}\\ &= \sum_{0,\,m-1}^{h} \varepsilon^h A_h^{(\nu)},\end{split}}
in which the $A_0^{(\nu)}, A_1^{(\nu)} \ldots A_{m-1}^{(\nu)}$ are forms of degree $\nu$ with integral coefficients and the variables $\alpha$, but independent of $\varepsilon$. Here too let $A_h = A_k$ hold whenever $h \equiv k \pmod{m}$. We then prove the theorem:

2. \gesperrt{If one applies to the indices of the $\alpha$ the cyclic permutation $\pi$, then the indices of the coefficients of $(\varepsilon, \alpha)^\nu$ undergo the cyclic permutation $\pi^\nu$, i.\,e. $A_h^{(\nu)}$ passes into $A_{h+\nu}^{(\nu)}$.}

To prove this, we remark that the formula (5) remains correct for every arbitrary $m$th root of unity $\varepsilon$, including $1$, and accordingly there follows from (2)
\weq{6}{mA_k^{(\nu)} = \overset{\varepsilon}{\Sigma}\varepsilon^{-k}(\varepsilon, \alpha)^\nu.}


% ===== scan leaf 569, printed 544 =====
% !PRINTED 544
\weberpage{544}

If one carries out on the right-hand side of this formula the permutation $\pi$ in the indices of the $\alpha$, then by 1. there results
\[
\sum^{\varepsilon} \varepsilon^{-k-\nu}(\varepsilon,\alpha)^{\nu},
\]
that is, $A_k^{(\nu)}$ passes into $A_{k+\nu}^{(\nu)}$, as is asserted in Theorem 2.

Theorem 2. is a special case of a general theorem. If we develop a product of arbitrarily many factors
\[
(\varepsilon,\alpha)^{\nu}(\varepsilon^{\lambda_1},\alpha)^{\nu_1}(\varepsilon^{\lambda_2},\alpha)^{\nu_2}\dots,
\]
in which $\nu, \nu_1, \nu_2\dots$ are positive and $\lambda_1, \lambda_2\dots$ arbitrary integers, and order it, as before $(\varepsilon,\alpha)^{\nu}$, according to powers of $\varepsilon$, the result may be
\weq{7}{(\varepsilon,\alpha)^{\nu}(\varepsilon^{\lambda_1},\alpha)^{\nu_1}(\varepsilon^{\lambda_2},\alpha)^{\nu_2}\dots = \sum_{0,\,m-1}^{h}\varepsilon^h B_h,}

in which the $B_h$ are forms of the variables $\alpha$ independent of $\varepsilon$. Since this development also holds for all $m$th roots of unity $\varepsilon$, one can apply the formula (2) and obtains
% [sic] authority prints B_k rather than mB_k and puts each \lambda_j on \alpha in the
% factors below; coefficient extraction from (7) and the product displayed above require
% mB_k and (\varepsilon^{\lambda_j},\alpha), respectively.
\[
B_k = \sum^{\varepsilon} \varepsilon^{-k}(\varepsilon,\alpha)^{\nu}(\varepsilon,\alpha^{\lambda_1})^{\nu_1}(\varepsilon,\alpha^{\lambda_2})^{\nu_2}\dots
\]
from which one can conclude, by 1., that $B_k$ passes by the substitution $\pi$ into
\[
B_{k+\nu+\lambda_1\nu_1+\lambda_2\nu_2+\cdots}
\]
We have therefore:

\gesperrt{3. The permutation $\pi$, applied to the indices of the $\alpha$, calls forth among the indices of the coefficients of the product, ordered according to $\varepsilon$,}
\weq{8}{(\varepsilon,\alpha)^{\nu}(\varepsilon^{\lambda_1},\alpha)^{\nu_1}(\varepsilon^{2},\alpha)^{\nu_2}\dots} % [sic] authority prints \varepsilon^2 here, omitting the expected \lambda_2.
\gesperrt{the permutation $\pi^{\nu+\lambda_1\nu_1+\lambda_2\nu_2\cdots}$.}

In these theorems the permutation $\pi$ may be repeated; they remain correct when $\pi$ is replaced by any power of $\pi$, and they hold whatever $m$th root of unity $\varepsilon$ may be. But if $\varepsilon$ is an imprimitive root of unity, certain simplifications occur, with which we must still become acquainted. If $e$ is a divisor of $m$,
\[
m = ef
\]
and $\varepsilon$ is any $e$th root of unity, then the expression (1) becomes
\begin{align*}
(\varepsilon,\alpha) &= \alpha + \varepsilon\alpha_1 + \cdots + \varepsilon^{e-1}\alpha_{e-1}\\
&\quad + \alpha_e + \varepsilon\alpha_{e+1} + \cdots + \varepsilon^{e-1}\alpha_{2e-1}\\
&\quad \cdots\cdots\cdots\cdots\cdots\cdots\cdots\cdots\cdots\cdots\\
&\quad + \alpha_{e(f-1)} + \varepsilon\alpha_{e(f-1)+1} + \cdots + \varepsilon^{e-1}\alpha_{m-1}
\end{align*}


% ===== scan leaf 570, printed 545 =====
% !PRINTED 545
\weberpage{545}

We now introduce the notation:
\weq{9}{\begin{aligned}
\eta &= \alpha + \alpha_e + \alpha_{2e} + \cdots + \alpha_{(f-1)e}\\
\eta_1 &= \alpha_1 + \alpha_{e+1} + \alpha_{2e+1} + \cdots + \alpha_{(f-1)e+1},\\
&\cdots\cdots\cdots\cdots\cdots\cdots\cdots\cdots\cdots\cdots\\
\eta_{e-1} &= \alpha_{e-1} + \alpha_{2e-1} + \alpha_{3e-1} + \cdots + \alpha_{m-1},
\end{aligned}}

and call these quantities, as \gesperrt{Gauss} did in the special case of the division of the circle, the $f$-term periods of the quantities $\alpha$. We then have
\weq{10}{(\varepsilon,\alpha) = \eta + \varepsilon\eta_1 + \varepsilon^2\eta_2 + \cdots + \varepsilon^{e-1}\eta_{e-1},}
for which we may also write $(\varepsilon,\eta)$.

The application of the permutation $\pi$ to the indices of $\alpha$ brings forth, among the indices of the quantities $\eta$, which are to be taken with respect to the modulus $e$, the cyclic permutation
\weq{11}{\gamma = (0, 1, 2\dots e-1)}

Corresponding to formula (10), we can now also write the formulas (5) and (7) thus:
\[
(\varepsilon,\eta)^{\nu} = E_0^{(\nu)} + \varepsilon E_1^{(\nu)} + \varepsilon^2 E_2^{(\nu)} + \cdots + \varepsilon^{e-1}E_{e-1}^{(\nu)}
\]
\[
(\varepsilon,\eta)^{\nu}(\varepsilon^{\lambda_1},\eta)^{\nu_1}(\varepsilon^{\lambda_2},\eta)^{\nu_2}\dots = G_0 + \varepsilon G_1 + \varepsilon^2 G_2 + \cdots + \varepsilon^{e-1}G_{e-1},
\]
in which then
\[
E_k^{(\nu)} = A_k^{(\nu)} + A_{e+k}^{(\nu)} + \cdots + A_{(f-1)e+k}^{(\nu)}
\]
\[
G_k = B_k + B_{e+k} + \cdots + B_{(f-1)e+k}
\]
and $E_k^{(\nu)}$ and $G_k$ are homogeneous polynomials in the $\alpha$, which can also be represented as functions of the $\eta$.

The quantities $\eta_k$, $E_k^{(\nu)}$, $G_k$ remain unchanged when the permutation $\pi^e$ is applied to their indices, that is, when the index is altered by a multiple of $e$; and we can therefore now complete Theorems 2. and 3. as follows:

\gesperrt{4. If $\varepsilon$ is an arbitrary $e$th root of unity and $e$ a divisor of $m$, then, when the permutation $\pi$ is applied to the indices of $\alpha$, the indices $k$ of the coefficients $E_k^{(\nu)}$ of $(\varepsilon,\eta)^{\nu}$ undergo the permutation $\gamma^{\nu}$, and the indices of the coefficients $G$ of the product $(\varepsilon,\eta)^{\nu}(\varepsilon^{\lambda_1},\eta)^{\nu_1}(\varepsilon^{\lambda_2},\eta)^{\nu_2}\dots$ the permutation $\gamma^{\nu+\lambda_1\nu_1+\lambda_2\nu_2\dots}$}


% ===== scan leaf 571, printed 546 =====
% !PRINTED 546
\weberpage{546}
\pdfbookmark[2]{§ 165. Solution of the cyclic equations}{bm184}

\begin{center}
§. 165.

\gesperrt{Solution of the Cyclic Equations.}
\end{center}

The \gesperrt{Lagrange} resolvents, by application of the theorems now proved, lead to the solution of the cyclic equations — more precisely, to the reduction to \gesperrt{pure equations}.

We now understand by the $\alpha$ no longer arbitrary variables, but the roots of a cyclic equation, so that the cyclic functions of the $\alpha$ are to be regarded as known quantities.

By the theorem of \S.\,164, 2., the coefficients of $(\varepsilon,\alpha)^m$ are cyclic functions of the $\alpha$.

If we understand by $a_0, a_1\dots a_{m-1}$ quantities in $\Omega$ and set
\weq{1}{\psi_\lambda = a_0 + a_1\varepsilon^\lambda + a_2\varepsilon^{2\lambda} + \cdots + a_{m-1}\varepsilon^{(m-1)\lambda},}
then from this theorem it follows that
\weq{2}{(\varepsilon^\lambda,\alpha) = \sqrt[m]{\psi_\lambda}.}

If $\varepsilon$ denotes therein a primitive $m$th root of unity, then all the resolvents are contained in the form $(\varepsilon^\lambda,\alpha)$.

If we observe further that $(1,\alpha) = a$, as the sum of the roots, belongs to the known quantities, then by \S.\,164 (3) we have
\weq{3}{m\alpha = a + \sqrt[m]{\psi_1} + \sqrt[m]{\psi_2} + \cdots + \sqrt[m]{\psi_{m-1}},}
and with this $\alpha$ is expressed by radicals of degree $m$, which under the root signs contain, besides the quantities that occur in $\Omega$ from the outset, also $m$th roots of unity.

Each of these radicals has, considered by itself, $m$ different values, which differ from one another by $m$th roots of unity as factors. If we give each $m$th root all its values, we obtain from (3) many different values of $\alpha$, among which, by \S.\,164, (4), all the roots $\alpha, \alpha_1, \alpha_2\dots\alpha_{m-1}$ occur. But the number of the expressions thus derived from (3) is much greater, and it remains to distinguish those which are to be retained from those which are to be discarded. The following way leads to this most simply.


% ===== scan leaf 572, printed 547 =====
% !PRINTED 547
\weberpage{547}

If we apply the theorem of \S.\,164, 3. to only two factors, it follows that
% [sic] authority prints the Latin u as the exponent; the surrounding formula requires \mu.
\[
(\varepsilon,\alpha)^{\nu}(\varepsilon^{\lambda},\alpha)^{u}
\]
is a function in $\Omega(\varepsilon)$ when $\nu + \lambda\mu \equiv 0 \pmod{m}$. If we therefore set $\mu = 1$, $\nu = m - \lambda$, it follows, if
\[
\chi_\lambda = b_0^{(\lambda)} + b_1^{(\lambda)}\varepsilon + \cdots + b_{m-1}^{(\lambda)}\varepsilon^{m-1}
\]
denotes a quantity in $\Omega(\varepsilon)$, that
\[
(\varepsilon,\alpha)^{m-\lambda}(\varepsilon^{\lambda},\alpha) = \chi_\lambda,
\]
and hence by (2)
\weq{4}{\sqrt[m]{\psi_\lambda} = \frac{\chi_\lambda}{\left(\sqrt[m]{\psi_1}\right)^{m-\lambda}} = \frac{\chi_\lambda\left(\sqrt[m]{\psi_1}\right)^{\lambda}}{\psi_1},}
and by this, if $\varepsilon$ denotes a fixed primitive $m$th root of unity, all the radicals occurring in (3) are expressed rationally in terms of one of them, $\sqrt[m]{\psi_1}$.

If one gives this one its $m$ different values, one obtains from (3) precisely the $m$ different values $\alpha$.

There is only one exceptional case in which this procedure is not applicable, namely when $\psi_1 = 0$. But by a small modification of the procedure we can free ourselves from such an exceptional case. To this we premise the following.

Let $p$ be a prime dividing $m$, and $m = pn$; as above, let $\varepsilon$ be any fixed primitive $m$th root of unity. Then there is always a $\lambda$ not divisible by $p$ such that $(\varepsilon^{\lambda},\alpha)$ is different from zero. For if we form the difference $\alpha_n - \alpha_0$ according to the formula \S.\,164, (4), we obtain
\[
m(\alpha_n - \alpha_0) = \sum_{0,\,m-1}^{\lambda}(\varepsilon^{-n\lambda} - 1)(\varepsilon^{\lambda},\alpha).
\]
Now, however, $\varepsilon^{-n\lambda} - 1$ is always $= 0$ whenever $\lambda$ is divisible by $p$, and if $(\varepsilon^{\lambda},\alpha)$ vanishes in all the other cases, that is, where $\lambda$ is not divisible by $p$, then $\alpha_n = \alpha_0$, contrary to the hypothesis that the $\alpha$ are all to be different. There is therefore at least one $\lambda$ not divisible by $p$ such that $(\varepsilon^{\lambda},\alpha)$ is different from zero.

We now decompose $m$ into its prime factors and set
\[
m = p_1 p_2 \dots,
\]
%CONT>>


% ===== scan leaf 573, printed 548 =====
% !PRINTED 548
\weberpage{548}
%<<CONT
in which $p_1, p_2 \ldots$ are powers of \gesperrt{different} primes.
We further set
\[
m = p_1 m_1 = p_2 m_2 \ldots,
\]
and choose, as is always possible by what has just been proved,
$\lambda_1$ relatively prime to $p_1$, $\lambda_2$ relatively prime to $p_2$ etc., so that
\[
(\varepsilon^{\lambda_1}, \alpha), \quad (\varepsilon^{\lambda_2}, \alpha) \ldots
\]
are different from zero. Then by the Theorem \S.\,164, 3.:
\weq{5}{(\varepsilon^{\lambda}, \alpha)(\varepsilon^{\lambda_1}, \alpha)^{m_1 \nu}(\varepsilon^{\lambda_2}, \alpha)^{m_2 \nu} \ldots = \chi_{\lambda}}
is a quantity contained in $\Omega(\varepsilon)$, if
\weq{6}{\lambda \equiv -\nu(\lambda_1 m_1 + \lambda_2 m_2 + \cdots) \ (\mathrm{mod}\ m).}
But $(\varepsilon^{\lambda_1}, \alpha)^{m_1}$ is a root of degree $p_1$ of a function
in $\Omega(\varepsilon)$, and we therefore set
% [sic] authority omits the exponent m_2 from the second factor; symmetry with the first factor requires it.
\weq{7}{(\varepsilon^{\lambda_1}, \alpha)^{m_1} = \sqrt[p_1]{\varphi_1}, \quad (\varepsilon^{\lambda_2}, \alpha) = \sqrt[p_2]{\varphi_2} \ldots, \quad (\varepsilon^{\lambda}, \alpha) = \sqrt[m]{\psi_{\lambda}}.}
Then by (5) we have
\weq{8}{\sqrt[m]{\psi_{\lambda}} = \frac{\chi_{\lambda}}{\left(\sqrt[p_1]{\varphi_1}\sqrt[p_2]{\varphi_2} \ldots \right)^{\nu}}.}
Now $(\lambda_1 m_1 + \lambda_2 m_2 + \cdots)$ is relatively prime to $m$, since $m_2 \ldots$
are divisible by $p_1$, while $\lambda_1 m_1$ is relatively prime to $p_1$; and thus from (6)
one obtains for each $\lambda$ a number $\nu$ completely determined
with respect to the modulus $m$.
If, then, we substitute the expressions (8) in (3) and assign to the
radicals $\sqrt[p_1]{\varphi_1}, \sqrt[p_2]{\varphi_2} \ldots$ all their values, we obtain
exactly $m$ different values for $\alpha$, and no more.

The last results we can use to derive a form of the
representation of the roots $\alpha$ in a somewhat generalised shape,
which \gesperrt{Abel} communicates in the place cited, and which
relates to the case where the field $\Omega$ is real, i.e. consists of
nothing but real numbers. The functions $\varphi, \psi, \chi$, as we
have used them above, are then composed of real
numbers and of the root of unity $\varepsilon$, which one finds by the division
of the circumference of the circle into $m$ equal parts; one may, for instance, set
\[
\varepsilon = e^{\frac{2\pi i}{m}} = \cos\frac{2\pi}{m} + i \sin\frac{2\pi}{m}
\]

The function $\varphi_1$ passes over, when $\varepsilon$ is changed into $\varepsilon^{-1}$, into
the conjugate imaginary value, which we denote by $\varphi_1'$.


% ===== scan leaf 574, printed 549 =====
% !PRINTED 549
\weberpage{549}
We shall choose a positive quantity $\varrho_1$ and an angle $\Theta_1$ in such
a way that
\weq{9}{\varphi_1 = \varrho_1 e^{i\Theta_1}, \qquad \varphi_1' = \varrho_1 e^{-i\Theta_1},}
or
\weq{10}{\begin{aligned} \varphi_1 &= \varrho_1(\cos\Theta_1 + i\sin\Theta_1) \\ \varphi_1' &= \varrho_1(\cos\Theta_1 - i\sin\Theta_1), \end{aligned}}
from which it further follows that
\weq{11}{\varrho_1^2 = \varphi_1 \varphi_1'.}
Now
\weq{12}{(\varepsilon^{\lambda_1}, \alpha)(\varepsilon^{-\lambda_1}, \alpha) = \pm a_1}
is a quantity of the field $\Omega(\varepsilon)$ by (\S.\,164, 3.), and indeed,
since it does not change on passing to the conjugate imaginary value,
that is, on the interchange of $\varepsilon$ and $\varepsilon^{-1}$, it is a
real quantity. We shall determine the sign so that $a_1$ is
positive. Hence by (11), since $\varrho_1^2$ is positive,
\[
\varrho_1^2 = (\pm a_1)^m = a_1^m,
\]
and from this it follows that for odd $m$ the upper sign
holds in any case; for even $m$ the lower can also occur. We then have
\weq{13}{\varrho_1 = \sqrt{a_1^m},}
where the square root is to be taken positive.

Further, for the same reason,
\weq{14}{\frac{\varphi_1 + \varphi_1'}{2} = b_1, \qquad \frac{\varphi_1 - \varphi_1'}{2i} = c_1}
are real quantities in $\Omega(\varepsilon)$, and it follows from (10) that
\weq{15}{\cos\Theta_1 = \frac{b_1}{\sqrt{a_1^m}}, \qquad \sin\Theta_1 = \frac{c_1}{\sqrt{a_1^m}},}
from which the relation further follows:
\[
a_1^m = b_1^2 + c_1^2.
\]

Accordingly we obtain
\[
\sqrt[p_1]{\varphi_1} = \sqrt{a_1^{m_1}}\, e^{\frac{i\Theta_1}{p_1}};
\]
and one now proceeds in the same way with the functions $\varphi_2$ and so on.
One thus determines $\Theta_2 \ldots$ from a system of equations
like (15), and, if one further puts
\weq{16}{m_1\Theta_1 + m_2\Theta_2 + \cdots = \Theta}
one obtains, by (8)
%CONT>>


% ===== scan leaf 575, printed 550 =====
% !PRINTED 550
\weberpage{550}
%<<CONT
\weq{17}{\sqrt[m]{\psi_{\lambda}} = \chi_{\lambda}\left(\sqrt{a_1^{m_1} a_2^{m_2} \ldots}\right)^{-\nu} e^{-\frac{i\Theta}{m}\nu}.}

By (15) we can set
\[
e^{i\Theta_1} = \frac{b_1 + ic_1}{\sqrt{a_1^m}}, \qquad e^{i\Theta_2} = \frac{b_2 + ic_2}{\sqrt{a_2^m}} \cdots,
\]
and if, then, by decomposition into the real and imaginary part we
obtain
\weq{18}{(b_1 + ic_1)^{m_1}(b_2 + ic_2)^{m_2} \cdots = B + iC}
and set
\weq{19}{a_1^{m_1} a_2^{m_2} \cdots = A}
it follows by (16) that
\[
e^{i\Theta} = \frac{B + iC}{\sqrt{A^m}},
\]
hence
\weq{20}{\sqrt{A^m}\cos\Theta = B, \qquad \sqrt{A^m}\sin\Theta = C,}
and
\weq{21}{\sqrt[m]{\psi_{\lambda}} = \chi_{\lambda}\sqrt{A}^{-\nu}\left(\cos\frac{\Theta\nu}{m} - i\sin\frac{\Theta\nu}{m}\right).}

The angle $\Theta$ is determined by (20) only up to a multiple of
$2\pi$, and if one therefore puts $\Theta + 2h\pi$ for $\Theta$ and lets $h$
run from $0$ to $m - 1$, one obtains from (21) the $m$ different values of the root quantity. The functions $\cos\dfrac{\Theta\nu}{m}$
and $\sin\dfrac{\Theta\nu}{m}$ can further be expressed rationally in terms of $\cos\dfrac{\Theta}{m}$, $\sin\dfrac{\Theta}{m}$. The solution of the cyclic equations in
the real field $\Omega$ is thus reduced to the following.

\gesperrt{One first adjoins to the field $\Omega$ the $m$th root of unity $\varepsilon$ (division of the circumference of the circle into $m$ equal
parts). Thereupon $A, B, C$ are known. One further adjoins
the positive square root $\sqrt{A}$; then
$\cos\Theta$ and $\sin\Theta$ are known by (20). Finally one adjoins
$\cos\dfrac{\Theta}{m}$, $\sin\dfrac{\Theta}{m}$ (division of the angle $\Theta$ into $m$ parts).
Then the cyclic equation is solved by (21).}

These considerations lead further to an interesting
result on the reality relations of the roots of cyclic
equations.


% ===== scan leaf 576, printed 551 =====
% !PRINTED 551
\weberpage{551}
\pdfbookmark[2]{§ 166. Division of the angle}{bm185}
If, as in \S.\,163, we represent each root in the sequence of roots
$\alpha, \alpha_1, \ldots \alpha_{m-1}$ rationally by the preceding one:
\[
\alpha_1 = \Theta(\alpha), \quad \alpha_2 = \Theta(\alpha_1), \ldots \alpha_{m-1} = \Theta(\alpha_{m-2}), \quad \alpha = \Theta(\alpha_{m-1}),
\]
in which, since $\Omega$ is here assumed to be real, $\Theta(x)$ denotes a real rational function of $x$, then it follows first that if
one of the roots is real, all the others must be real as well.
This is always the case for odd $m$, since a real equation of odd degree must always have at least one real
root.

For even $m$ imaginary roots may also be present,
and if one root is imaginary, then all must be so,
since, if a real root occurs, all the others are also
real. If we denote by $\Theta^{\nu}(x)$ the $\nu$-fold repetition of the
function $\Theta$, then
% [sic] authority prints alpha^k; the later parallel occurrence alpha_k confirms k should be a subscript.
\[
\alpha_{\nu+k} = \Theta^{\nu}(\alpha^k).
\]
If then $\alpha = \alpha_0$ is conjugate imaginary with $\alpha_k$, then for
every $\nu$ the functions $\Theta^{\nu}(\alpha_0)$ and $\Theta^{\nu}(\alpha_k)$, that is, $\alpha_{\nu}$ and $\alpha_{\nu+k}$,
are also conjugate imaginary.

From this it follows that $2k = m$ must hold, and we conclude
that in the case of imaginary roots
\[
\alpha_k \quad \text{and} \quad \alpha_{k + \frac{m}{2}}
\]
form, for every index $k$, a pair of conjugate imaginary roots.

The cubic equations become cyclic equations by the adjunction of the
square root of the discriminant.
When the discriminant is positive, the roots are, in
agreement with this theorem, real.

\begin{center}
\S.\,166.\\
\gesperrt{Division of the angle.}
\end{center}

Cyclic equations also include those
on which the division of an angle into $m$ equal parts depends; to these we have reduced the general cyclic equations in
a real field.

The problem can be formulated thus:

\gesperrt{When $\cos m\varphi$ and $\sin m\varphi$ are given, then
$\cos\varphi$ and $\sin\varphi$ are to be found from them.}


% ===== scan leaf 577, printed 552 =====
% !PRINTED 552
\weberpage{552}
We suppose that the given values of $\cos m\varphi$ and $\sin m\varphi$ are less than 1 and that the sum of their squares is $= 1$.

We further adjoin $m$th roots of unity, with which we shall occupy ourselves more thoroughly in the next section, but which in any case depend on equations whose degree is lower than $m$. In order, however, to keep the field real, we shall adjoin, instead of the roots of unity themselves, the two real quantities
\weq{1}{\sin\frac{2\pi}{m}, \qquad \cos\frac{2\pi}{m}}

Let the field $\Omega$, then, consist of all rational numbers and of the rational functions of $\cos m\varphi$, $\sin m\varphi$, $\cos\frac{2\pi}{m}$, $\sin\frac{2\pi}{m}$.

The equation of degree $m$ on which $x = 2\cos\varphi$ depends we have set up in \S.\,137 in the form
\weq{2}{2\cos m\varphi = A_m(x),}
where $A_m(x)$ is a polynomial of degree $m$ in $x$; but there we have also found the equation
\weq{3}{\sin\varphi = \frac{\sin m\varphi}{B_m(x)}}
by which $\sin\varphi$ is expressed rationally in terms of $x$. In these formulae the roots of unity are not yet contained.

The $m$ roots of (2) now have the following meaning:
\[
x_0 = 2\cos\varphi, \qquad x_1 = 2\cos\left(\varphi + \frac{2\pi}{m}\right)\dots,
\]
\[
x_{m-1} = 2\cos\left(\varphi + \frac{2(m-1)\pi}{m}\right),
\]
and with the aid of (2) and (3), and upon adjunction of the quantities (1), each of them can be represented as a rational function of another, namely thus:
\[
x_1 = f(x_0), \qquad x_2 = f(x_1)\dots, \qquad x_0 = f(x_{m-1}).
\]

Any function of the $x_0, x_1 \dots x_{m-1}$ can therefore be represented as a rational function $F(x_0)$; under a cyclic permutation the following transformations occur:
\[
F(x_0) \text{ into } F(x_1), \ F(x_1) \text{ into } F(x_2) \dots F(x_{m-1}) \text{ into } F(x_0)
\]
For a \gesperrt{cyclic function}, therefore,
%CONT>>


% ===== scan leaf 578, printed 553 =====
% !PRINTED 553
\weberpage{553}
%<<CONT
\[
F(x_0) = F(x_1) \cdots = F(x_{m-1}) = \frac{1}{m}\left[F(x_0) + F(x_1) \cdots + F(x_{m-1})\right],
\]
and consequently $F(x_0)$ is rational, since it is represented as a symmetric function of the roots.

\gesperrt{The division of the angle thus depends on a cyclic equation.}

This cyclic equation is irreducible; for if in any rational equation $\Phi(x_0, \cos m\varphi, \sin m\varphi) = 0$ one replaces $\varphi$ by $\varphi + \frac{2\pi k}{m}$, it goes over into
\[
\Phi(x_k, \cos m\varphi, \sin m\varphi) = 0
\]
and is therefore satisfied for \gesperrt{all} roots of (2).

If we wished to adjoin only $\cos m\varphi$, and not at the same time $\sin m\varphi$, then the equation would no longer be a cyclic one, which one easily confirms by the example of the trisection, where precisely $\sin m\varphi$ becomes the square root of the discriminant (\S.\,112).

\begin{center}
\rule{3cm}{0.4pt}
\end{center}

% Below the direct-printed closing rule, only reversed/faint show-through is visible; excluded by policy.


% ===== scan leaf 579, printed 554 =====
\weberpage{554}% page number not printed on this page
\pdfbookmark[1]{Sixteenth Section. Cyclotomy.}{bm186}
\pdfbookmark[2]{§ 167. The cyclotomic periods and the period equations}{bm187}
% !PRINTED none
% Abschnitt-opening (part-title) page: no folio is printed on this page in the scan
% (section-opening pages in this book suppress the running head/page number).
% Content: "Sechzehnter Abschnitt." / "Kreistheilung." (both letter-spaced, centred),
% a horizontal divider rule, then "§. 167." and its letter-spaced title, followed
% immediately by the start of ordinary body text (no page break before the body begins).
\begin{center}
\gesperrt{Sixteenth Section.}

\gesperrt{Division of the Circle.}

\rule{3cm}{0.4pt}

\S.\,167.

\gesperrt{The cyclotomic periods and the period equations.}
\end{center}

The most important among the Abelian equations are those on which the determination of the roots of unity depends, whose elementary properties we have already come to know in the twelfth section. On account of the relation to the geometric problem of dividing the circumference of the circle into a definite number of equal parts, all the equations connected with this are also called the \gesperrt{cyclotomic equations}.

The field $\Omega$, which contains the quantities regarded as known, is here simply the \gesperrt{field of rational numbers}, which we shall call the field $R$.

We shall first occupy ourselves with the $n$th roots of unity under the assumption that $n$ is an odd prime (the prime 2 gives rise to no further questions, since the only second roots of unity, $\pm 1$, are contained in $R$).

Besides the rational $n$th root of unity 1 there exist $n - 1$ more primitive ones, which can be represented by the transcendental expressions
\[
e^{\frac{2\pi i}{n}}, \quad e^{\frac{4\pi i}{n}}, \quad e^{\frac{6\pi i}{n}} \dots e^{\frac{2(n-1)\pi i}{n}}
\]
If we denote one of them by $r$, then the whole system is also represented by the following list:
\weq{1}{r, \ r^2, \ r^3 \dots r^{n-1}}
%CONT>>


% ===== scan leaf 580, printed 555 =====
% !PRINTED 555
\weberpage{555}
%<<CONT
In the exponent of $r$ only the residue with respect to the modulus $n$ matters, since $r^h = r^k$ holds always and only when $h \equiv k \pmod{n}$.

As we have already seen in the twelfth section, the quantities (1) are the roots of an \gesperrt{irreducible equation} of degree $(n-1)$
\weq{2}{X = x^{n-1} + x^{n-2} + x^{n-3} + \cdots + x + 1 = 0,}
and we wish first to show that this equation is \gesperrt{cyclic}. This property results from the existence of primitive roots of the prime $n$, with which we occupied ourselves in \S.\,136. It was proved there that for every prime $n$ there are certain numbers $g$, called primitive roots of $n$, which are characterised by the property that among the residues of the powers
\[
1, \ g, \ g^2, \ g^3 \dots g^{n-2},
\]
upon division by $n$, each of the numbers
% [sic] authority includes 0 here; powers of a primitive root modulo n cover only 1,...,n-1.
\[
0, \ 1, \ 2, \dots n-1
\]
occurs once and only once. The residue of a power $g^h$ remains the same when $h$ is changed by a multiple of $n-1$. If
\[
g^\alpha \equiv a \pmod{n},
\]
then $\alpha$ is called the index of $a$ ($\alpha = \operatorname{ind} a$) and $a$ the number or the numerus. The index is taken with respect to the modulus $n-1$, the numerus with respect to the modulus $n$ (\S.\,136).

If, then, we assume such a primitive root of $n$, we can, apart from the order, represent the quantities (1) thus:
\[
r, \ r^g, \ r^{g^2} \dots r^{g^{n-2}},
\]
or, if we set
\weq{3}{r^{g^h} = r_h}
then
\weq{4}{r, \ r_1, \ r_2 \dots r_{n-2}.}
Every number of the field $R(r)$, i.\,e. every rational function of $r$, can be brought into the form
\weq{5}{\varphi(r) = b_0 + b_1 r + b_2 r^2 + \cdots + b_{n-2} r^{n-2},}
or, since by (2)
\[
1 + r + r^2 + \cdots + r^{n-1} = 0
\]
holds,
%CONT>>


% ===== scan leaf 581, printed 556 =====
% !PRINTED 556
\weberpage{556}

\weq{6}{\varphi(r) = (b_1 - b_0)\,r + (b_2 - b_0)\,r^2 + \cdots + (b_{n-2} - b_0)\,r^{n-2} - b_0 r^{n-1};}

or, since the powers $r, r^2, \ldots r^{n-1}$, apart from their order, agree with $r, r_1, r_2 \ldots r_{n-2}$, into the form

\weq{7}{\varphi(r) = a\,r + a_1 r_1 + a_2 r_2 + \cdots + a_{n-2} r_{n-2},}

in which the coefficients $b$ and $a$ are rational numbers.

\gesperrt{I. Such a function $\varphi(r)$ can be equal to zero only when all its coefficients $a, a_1 \ldots a_{n-2}$ vanish;}

for the expression (5) shows that $\varphi(r)$ cannot vanish otherwise than when the $b_0, b_1 \ldots b_{n-2}$ are zero. But if these are zero, then (6) shows that the coefficients $a$ must also all be zero. Hence $\varphi(r)$ can be brought into the form (7) in only one way.

If we make the stipulation that $a_h$ and $r_h$ shall not change when the index $h$ increases by a multiple of $n - 1$, we can also set

\weq{8}{\varphi(r) = \sum^h a_h r_h}

and let $h$ run through a complete system of residues with respect to the modulus $n - 1$.

We now have to prove that the group of the equation $X = 0$ is no other than the period of the cyclic permutation
\[
\pi = (r, r_1, r_2 \ldots r_{n-2}).
\]

This period or cyclic group we denote by

\weq{9}{C = 1, \pi, \pi^2 \ldots \pi^{n-2}.}

The equation $X = 0$ is in any case a \gesperrt{normal equation}, because it is irreducible, and because all its roots are, by (3), rationally expressible in terms of one among them; it is therefore its own \gesperrt{Galois} resolvent (\S.\,145). Its substitutions are

\weq{10}{(r, r), (r, r_1), (r, r_2) \ldots (r, r_{n-2}),}

and by (3) we have
\[
(r, r_h) = (r_k, r_{k+h}).
\]

From this it follows that $(r, r_1)$ produces among the roots (4) the cyclic permutation $\pi$, and that therefore the substitution group
%CONT>>


% ===== scan leaf 582, printed 557 =====
% !PRINTED 557
\weberpage{557}

%<<CONT
(10) is isomorphic with the permutation group $C$. Hence $C$ is the \gesperrt{Galois} group of $X = 0$.

By \S.\,163 the equation $X = 0$ can be reduced to a sequence of cyclic equations of lower degree, if one seeks functions that belong to the divisors of the group $C$. But the divisors of $C$, if $n - 1$ is decomposed into two factors $e, f$, so that
\[
n - 1 = ef
\]
is set, are the cyclic groups

\weq{11}{C_e = 1, \pi^e, \pi^{2e} \ldots \pi^{(f-1)e},}

and the permutation $\pi^e$ is therein equivalent to the substitution $(r, r_e)$.

In order to arrive at a function of the roots belonging to $C_e$ of the simplest possible form, we consider, following \gesperrt{Gauss}, the $f$-\gesperrt{term periods}

\weq{12}{\begin{aligned}
\eta &= r + r_e + r_{2e} + \cdots + r_{(f-1)e} \\
\eta_1 &= r_1 + r_{e+1} + r_{2e+1} + \cdots + r_{(f-1)e+1} \\
&\cdots\cdots\cdots\cdots\cdots\cdots\cdots\cdots \\
\eta_{e-1} &= r_{e-1} + r_{2e-1} + r_{3e-1} + \cdots + r_{n-2},
\end{aligned}}

which arise from the first of them by the substitutions $(r, r), (r, r_1), (r, r_2) \ldots (r, r_{e-1})$. We also designate the quantities (12) as a system of \gesperrt{conjugate periods}.

Each of these functions remains unchanged by the substitution $(r, r_e)$, and they are all different from one another, since an equation $\eta_h - \eta_k = 0$, if $h$ is different from $k$, could be written in the form
\[
a r + a_1 r_1 + \cdots + a_{n-2} r_{n-2} = 0
\]
where the $a, a_1 \ldots a_{n-2}$ are integers not all vanishing at once ($+ 1$, $0$ or $- 1$), which by Theorem I is not possible. Each of the periods (12) is therefore a function \gesperrt{belonging} to the group $C_e$, and these quantities are therefore the roots of an irreducible equation of degree $e$ with integer coefficients. If one adjoins one of these quantities, the determination of $r$ still depends on an equation of degree $f$.

By the general theorems of the fourteenth section, every function of the $r$ that admits the substitutions of the group $C_e$ can be represented rationally by any one of the quantities $\eta$. But here we can also set up the following theorem:
%CONT>>


% ===== scan leaf 583, printed 558 =====
% !PRINTED 558
\weberpage{558}
%<<CONT

\gesperrt{II. Every number in $R(r)$ which admits the substitution $(r, r_e)$, and hence also every number of the field $R(\eta)$, can be represented as a homogeneous linear function of the periods $\eta, \eta_1 \ldots \eta_{e-1}$.}

For by I. we can put every number of $R(r)$ in one way into the form

\weq{13}{\varphi(r) = \sum a_h r_h}

where the coefficients $a_h$ are rational numbers, and $h$ runs through a complete system of residues with respect to the modulus $n - 1$. But we have
\[
\varphi(r_e) = \sum a_h r_{h+e} = \sum a_{h-e} r_h,
\]
and therefore, if $\varphi(r) = \varphi(r_e)$ is to hold, we must have $a_{h-e} = a_h$ for every $h$; hence also $a_h = a_{h+e} = a_{h+2e} \ldots$ Accordingly (13) can be put into the form
\[
\varphi(r) = \sum_{0,\,e-1}^{h} a_h r_h + \sum_{0,\,e-1}^{h} a_h r_{h+e} + \cdots + \sum_{0,\,e-1}^{h} a_h r_{h+(f-1)e},
\]
or

\weq{14}{\varphi(r) = a\eta + a_1\eta_1 + \cdots + a_{e-1}\eta_{e-1},}

which was to be proved. It is further to be remarked here that, if $\varphi$ has integer coefficients, the coefficients of the $\eta_h$ also become integers.

Since every number of the field $R(\eta)$ can be represented in one and only one way in the linear form (14), we call the system of quantities $\eta, \eta_1, \eta_2 \ldots \eta_{e-1}$ a \gesperrt{basis of the field $R(\eta)$}.

If one knows the representation of the product of each two periods $\eta_h \eta_k$ in the form (14), one can from it, by repeated application, represent all rational functions of the periods $\eta$ in the same form. Likewise one can also easily form the equation whose roots are the $e$ quantities $\eta$, if one applies the procedure to the product

\weq{15}{F_e(x) = (x - \eta)(x - \eta_1) \cdots (x - \eta_{e-1})}

This equation can also be represented in determinant form.

Namely, let us put

\weq{16}{\eta\eta_h = a_{0,h}\eta + a_{1,h}\eta_1 + \cdots + a_{e-1,h}\eta_{e-1},}

in which the $a_{0,h}, a_{1,h} \ldots a_{e-1,h}$ are integers, about whose calculation further details will follow later. If one sets up this equation for $h = 0, 1, \ldots e - 1$ in the form:
%CONT>>


% ===== scan leaf 584, printed 559 =====
% !PRINTED 559
\weberpage{559}

%<<CONT
\[
\begin{aligned}
(a_{0,0} - \eta)\,\eta + &\ a_{1,0}\eta_1 + \cdots \hspace{2em} a_{e-1,0}\eta_{e-1} = 0 \\
a_{0,1}\eta \ + &\ (a_{1,1} - \eta)\eta_1 + \cdots \hspace{2em} a_{e-1,1}\eta_{e-1} = 0 \\
&\cdots\cdots\cdots\cdots\cdots\cdots\cdots\cdots\cdots\cdots \\
a_{0,e-1} \ + &\ a_{1,e-1}\eta_1 + \cdots (a_{e-1,e-1} - \eta)\eta_{e-1} = 0,
\end{aligned}
\]
% [sic] The printed last row omits the factor \eta after a_{0,e-1}; retained diplomatically although the parallel rows require it.
one sees, since the $\eta, \eta_1 \ldots \eta_{e-1}$ do not vanish, that the determinant of this system must be zero, hence

\weq{17}{\begin{vmatrix} a_{0,0}-\eta, & a_{1,0}, & \ldots & a_{e-1,0} \\ a_{0,1}, & a_{1,1}-\eta, & \ldots & a_{e-1,1} \\ \cdots & \cdots & \cdots & \cdots \\ a_{0,e-1}, & a_{1,e-1}, & \ldots & a_{e-1,e-1}-\eta \end{vmatrix} = 0,}

which, if one writes $x$ for $\eta$, must agree, apart from the sign, with $F_e(x)$.

\gesperrt{The function $F_e(x)$ therefore has not merely rational, but integer coefficients.}

% [sic] authority says that r is indeterminate here, although the displayed polynomial has indeterminate x and r is one of its roots.
One can now also form the equation of degree $f$ whose roots are $r, r_e, r_{2e} \ldots r_{(f-1)e}$; for, for an indeterminate $r$,

\weq{18}{\Phi_e(x) = (x - r)(x - r_e) \cdots (x - r_{(f-1)e})}

is a function in the field $R(\eta)$. Its coefficients can easily be computed from the \gesperrt{Newton} formulae, by which the coefficients of an equation are expressed by means of the power sums of the roots (\S.\,42). For by (12) and (3) we have
\[
r^{g^h} + r_e^{g^h} + \cdots + r_{(f-1)e}^{g^h} = \eta_h,
\]
and this is thus directly equal to one of the power sums, and every power sum is equal to one of the periods.

But the equation (18) can be treated in just the same way as the equation (2); for the group of $\Phi_e(x) = 0$ is precisely $C_e$, and if $e'$ is a divisor of $f$, say $f = e'f'$, then $C_{ee'}$ is a divisor of $C_e$. But to $C_{ee'}$ belongs the period
\[
\eta' = r + r_{ee'} + r_{2ee'} + \cdots + r_{(f'-1)ee'},
\]
which therefore depends on an equation of degree $e'$ in the field $R(\eta)$.

If one takes the numbers $e, e' \ldots$ to be primes, then the sequence of resolvents consists of a sequence of equations whose degrees are the prime factors of $n - 1$.

As for the group of the equation of degree $e$ whose roots are the $e$ quantities $\eta$, we find this very simply
%CONT>>


% ===== scan leaf 585, printed 560 =====
% !PRINTED 560
\weberpage{560}
\pdfbookmark[2]{§ 168. Gauss's method for the computation of the resolvents}{bm188}
%<<CONT
from the remark that the substitution $(r, r_h)$ calls forth among the $\eta_k$
the substitution $(\eta_k, \eta_{k+h})$, when the index of $\eta$
is taken modulo $e$. The group of $F_e(x) = 0$
therefore consists of the powers of the cyclic permutation

\[
(\eta, \eta_1, \eta_2 \ldots \eta_{e-1}).
\]

\begin{center}
\S.\,168.

\gesperrt{The Gauss Method for the Computation of the\\
Resolvents.}
\end{center}

In order actually to carry out the solution of the cyclotomic
equation in concrete cases, all that matters, after what we
have developed in the preceding section, is to represent the
products of any two conjugate periods as linear functions
of these same periods. For this \gesperrt{Gauss} has given a simple
procedure, which we shall now become acquainted with. To this
end we must alter the notation of the periods a little.

We write any two of the periods \S.\,167, (12) in
the form

\weq{1}{\begin{aligned}
\eta^{(\lambda)} &= r^\lambda + r^{\lambda'} + r^{\lambda''} + \cdots\\
\eta^{(\mu)} &= r^\mu + r^{\mu'} + r^{\mu''} + \cdots
\end{aligned}}

so that, if $\lambda \equiv g^h \ (\mathrm{mod}\ n)$, then $\eta^{(\lambda)} = \eta_h$ or also
$\eta^{(\lambda)} = \eta_{\mathrm{ind}\,\lambda}$; we then have

\weq{2}{\begin{aligned}
\lambda' &\equiv \lambda g^e, & \lambda'' &\equiv \lambda g^{2e} \cdots\\
\mu' &\equiv \mu g^e, & \mu'' &\equiv \mu g^{2e} \cdots
\end{aligned}\ (\mathrm{mod}\ n);}

if, then, we let the symbols $s$ and $t$ each run through a complete system of residues
modulo $f$, we have

\[
\eta^{(\lambda)} = \sum^{s} r^{\lambda g^{se}}, \qquad \eta^{(\mu)} = \sum^{t} r^{\mu g^{te}}.
\]

The product of these is

\[
\eta^{(\lambda)} \eta^{(\mu)} = \sum^{s}\sum^{t} r^{\lambda g^{se} + \mu g^{te}}.
\]

If, in forming this sum, we first hold $s$ fixed,
and sum with respect to $t$, we may replace $t$ by $t + s$, because both run simultaneously through a complete system of residues
modulo $f$. Accordingly we obtain
\[
\eta^{(\lambda)}\eta^{(\mu)} = \sum^{s}\sum^{t} r^{(\lambda + \mu g^{te})g^{se}}.
\]
%CONT>>


% ===== scan leaf 586, printed 561 =====
% !PRINTED 561
\weberpage{561}
%<<CONT
In this, however, the order of summation may be interchanged,
so that we also have

\weq{3}{\eta^{(\lambda)}\eta^{(\mu)} = \sum^{t}\sum^{s} r^{(\lambda + \mu g^{te})g^{se}}}

The sum taken with respect to $s$,

\[
\sum^{s} r^{(\lambda + \mu g^{te})g^{se}}
\]

is, however, itself one of the periods, namely, according to the notation (1), the period

\[
\eta^{(\lambda + \mu g^{te})}.
\]

From (2) and (3) it therefore follows that

\weq{4}{\eta^{(\lambda)}\eta^{(\mu)} = \eta^{(\lambda+\mu)} + \eta^{(\lambda+\mu')} + \eta^{(\lambda+\mu'')} + \cdots}

The right-hand side of this expression contains $f$ terms; among them the same period may of course occur several times;
there may also occur among them the improper period $\eta^{(0)}$,
which is to be set equal to the integer $f$. If one wishes to restore the homogeneous form of the expression, one makes use of
the relation

\[
\eta + \eta_1 + \eta_2 + \cdots + \eta_{e-1} = -1,
\]

which is indeed only another way of writing the equation \S.\,167, (2).

When one has computed according to (4) the expression for a product $\eta\eta_h$,
to which we have already given the form above,

\weq{5}{\eta\eta_h = a_{0,h}\eta + a_{1,h}\eta_1 + \cdots + a_{e-1,h}\eta_{e-1},}

one obtains the product of any two of the periods
by the substitution $(\eta, \eta_k)$,

\weq{6}{\eta_k\eta_{h+k} = a_{0,h}\eta_k + a_{1,h}\eta_{k+1} + \cdots + a_{e-1,h}\eta_{k-1}.}

In order to illustrate these rules by a simple example,
we take $n = 13$. For 13, 2 is a primitive root, and
we can apply the index table given in \S.\,136:

\weq{7}{
\begin{array}{c|cccccccccccc}
I & 0, & 1, & 2, & 3, & 4, & 5, & 6, & 7, & 8, & 9, & 10, & 11\\
\hline
N & 1, & 2, & 4, & 8, & 3, & 6, & 12, & 11, & 9, & 5, & 10, & 7
\end{array}.
}

If we first assume $e = 3$, $f = 4$, we obtain
a cubic resolvent whose roots are

\weq{8}{\begin{aligned}
\eta &= r + r_3 + r_6 + r_9\\
\eta_1 &= r_1 + r_4 + r_7 + r_{10}\\
\eta_2 &= r_2 + r_5 + r_8 + r_{11}
\end{aligned}}

From the table (7) this also gives
%CONT>>


% ===== scan leaf 587, printed 562 =====
% !PRINTED 562
\weberpage{562}
%<<CONT
\weq{9}{\begin{aligned}
\eta &= r + r^{-5} + r^{-1} + r^5 = \eta^{(1)}\\
\eta_1 &= r^2 + r^3 + r^{-2} + r^{-3} = \eta^{(2)}\\
\eta_2 &= r^4 + r^6 + r^{-4} + r^{-6} = \eta^{(4)}
\end{aligned}}

If one applies formula (4), one obtains, e.g.,

\[
\begin{aligned}
\eta\eta &= \eta^{(2)} + \eta^{(-4)} + \eta^{(0)} + \eta^{(6)} = 4 + \eta^{(2)} + 2\eta^{(4)}\\
&= -4\eta - 3\eta_1 - 2\eta_2,
\end{aligned}
\]

and thus the formulae result

\[
\begin{aligned}
\eta^2 &= -4\eta - 3\eta_1 - 2\eta_2,\\
\eta\eta_1 &= \eta + 2\eta_1 + \eta_2,\\
\eta\eta_2 &= 2\eta + \eta_1 + \eta_2,
\end{aligned}
\]

from which, by \S.\,167, (17), the equation for $\eta$

\[
\begin{vmatrix}
-4-\eta, & -3, & -2\\
1, & 2-\eta, & 1\\
2, & 1, & 1-\eta
\end{vmatrix} = 0,
\]

or

\weq{10}{\eta^3 + \eta^2 - 4\eta + 1 = 0}

follows. The discriminant of this equation is found to be equal to
$169 = 13^2$, hence positive. The equation consequently has three real
% [sic] The print attributes the sign pattern to the positive constant term alone. That
% term shows only that an odd number of roots is negative; the negative coefficient of
% eta supplies the additional exclusion of three negative roots. The conclusion is correct.
roots, and indeed, since the last term is positive, two positive and one
negative root.

If we set $r = e^{\frac{2\pi i}{13}}$, then

\[
\begin{aligned}
\eta &= 2\cos\frac{2\pi}{13} + 2\cos\frac{10\pi}{13} = 2\left(\cos\frac{2\pi}{13} - \cos\frac{3\pi}{13}\right),\\
\eta_1 &= 2\cos\frac{4\pi}{13} + 2\cos\frac{6\pi}{13} = 2\left(\cos\frac{4\pi}{13} + \cos\frac{6\pi}{13}\right),\\
\eta_2 &= 2\cos\frac{8\pi}{13} + 2\cos\frac{12\pi}{13} = -2\left(\cos\frac{\pi}{13} + \cos\frac{5\pi}{13}\right),
\end{aligned}
\]

or also

\[
\begin{aligned}
\eta &= 4\cos\frac{4\pi}{13}\cos\frac{6\pi}{13}, \qquad \eta_1 = 4\cos\frac{\pi}{13}\cos\frac{5\pi}{13},\\
\eta_2 &= -4\cos\frac{2\pi}{13}\cos\frac{3\pi}{13}.
\end{aligned}
\]

Thus $\eta_2$ is the negative root, $\eta_1$ the greater, and $\eta$ the smaller of the
two positive roots.

The equation of degree $4$, whose roots are $r$, $r^{-1}$, $r^5$, $r^{-5}$,
can now easily be formed, if one
%CONT>>


% ===== scan leaf 588, printed 563 =====
% !PRINTED 563
\weberpage{563}
%<<CONT
\[
\xi = r + r^{-1} = 2\cos\frac{2\pi}{13}, \qquad \xi' = r^5 + r^{-5} = 2\cos\frac{10\pi}{13}
\]

sets. For we have

\[
\xi + \xi' = \eta, \qquad \xi\xi' = \eta_2,
\]

hence

\weq{11}{\xi^2 - \eta\xi + \eta_2 = 0}

is the quadratic equation for $\xi$, from which one obtains the biquadratic
equation for $r$,

\weq{12}{r^4 - \eta r^3 + (\eta_2 + 2)r^2 - \eta r + 1 = 0.}

The quadratic equation (11) has one positive and one negative
root. The positive root is $2\cos\frac{2\pi}{13}$. Instead of computing $r$
itself, one computes further $2\sin\frac{2\pi}{13}$, which one obtains as
the positive square root $\sqrt{4 - \xi^2}$.

In place of the three four-term periods $\eta$ one could also
have first considered two periods of six terms.

\weq{13}{\begin{aligned}
\xi &= r + r^3 + r^4 + r^{-1} + r^{-3} + r^{-4}\\
\xi_1 &= r^2 + r^5 + r^6 + r^{-2} + r^{-5} + r^{-6}.
\end{aligned}}

For the product $\xi\xi_1$ one finds by formula (4) the
value $3(\xi + \xi_1)$ or $-3$, so that $\xi$, $\xi_1$ are the roots of the quadratic equation

\weq{14}{\xi^2 + \xi - 3 = 0}

from which, since it is easy to see from the expression by the
cosines that $\xi$ is positive:

\weq{15}{\xi = \frac{-1 + \sqrt{13}}{2}, \qquad \xi_1 = \frac{-1 - \sqrt{13}}{2}.}

After adjunction of the values $\xi$, $\xi_1$ one can form a cubic
equation for the two-term periods. For if one
sets

\[
\begin{aligned}
\xi &= r + r^{-1}, & \xi_1 &= r^2 + r^{-2},\\
\xi_2 &= r^4 + r^{-4}, & \xi_3 &= r^5 + r^{-5},\\
\xi_4 &= r^3 + r^{-3}, & \xi_5 &= r^6 + r^{-6},
\end{aligned}
\]

then it follows that

\[
\zeta = \xi + \xi_2 + \xi_4, \qquad \zeta_1 = \xi_1 + \xi_3 + \xi_5, \qquad \zeta + \zeta_1 = -1,
\]

further

\[
\xi\xi_2 = \xi_3 + \xi_4, \qquad \xi\xi_4 = \xi_1 + \xi_2, \qquad \xi_2\xi_4 = \xi + \xi_5,
\]

\[
\xi\xi_2\xi_4 = 2 + \zeta_1 = 1 - \zeta,
\]
%CONT>>


% ===== scan leaf 589, printed 564 =====
% !PRINTED 564
\weberpage{564}
\pdfbookmark[2]{§ 169. Reduction of the cyclotomic equation to pure equations. Division into seventeen parts}{bm189}
%<<CONT
hence $\xi$, $\xi_2$, $\xi_4$ are the roots of the cubic equation
\weq{16}{x^3 - \zeta x^2 - x - 1 + \zeta = 0.}

\begin{center}
\S.\,169.

\gesperrt{Reduction of the cyclotomic equation to}\\
\gesperrt{pure equations.}
\end{center}

\gesperrt{Gauss}, already in his first exposition in the disq.
arithm., reduced the cyclotomic equations directly to \gesperrt{pure equations}
by the use of the resolvents\footnote{\gesperrt{Gauss}, disq. arithm. art. 359, 360; disq. circa aequationes puras ulterior evolutio, Werke Bd. II. \gesperrt{Lagrange}, r\'es. des \'equations num\'eriques. \gesperrt{Jacobi}, ``Ueber die Kreistheilung und ihre Anwendung in der Zahlentheorie''. Werke Bd. 6. \gesperrt{Kummer}, Crelle's Journal, Bd. 35 and Abhandl. d. Berl. Akademie 1856. \gesperrt{Eisenstein} and \gesperrt{Cauchy} are also to be mentioned. The theory of the division of the circle is presented in connected form, accompanied by the historical references, in the book by \gesperrt{Bachmann}, ``Die Lehre von der Kreistheilung''. Leipzig 1872.}. To
carry this out, we set
\[ n - 1 = m, \]
and denote, as in the preceding section, by $\varepsilon$ a primitive
$m$th root of unity, and by $r, r_1, r_2 \ldots r_{m-1}$ the $n$th roots of unity
 in the order fixed above. The \gesperrt{Lagrange}
resolvents are then
\[ (\varepsilon^\lambda, r) = \sum_{0,\,n-2}^{h} \varepsilon^{\lambda h} r_h = \sum_{0,\,n-2}^{h} \varepsilon^{\lambda h} r^{g^h}, \]
in which $\lambda$ is an arbitrary exponent taken with respect to the modulus $m$.
If we set $g^h = s$, so that $h \equiv \mathrm{ind}\,s \pmod{m}$, we may
also write
\weq{1}{(\varepsilon^\lambda, r) = \sum_{1,\,n-1}^{s} \varepsilon^{\lambda\,\mathrm{ind}\,s} r^s.}

The basis for the further reduction of these expressions
is the computation of the product of two such resolvents.

We denote by $\mu$ a second number like $\lambda$, and understand
by $s$ and $t$ two symbols which, independently of one another, each
run through a system of residues with respect to the modulus $n$, with the exclusion
of zero.
Then
\[ (\varepsilon^\lambda, r)(\varepsilon^\mu, r) = \sum^{s} \sum^{t} r^{(s+t)} \varepsilon^{\lambda\,\mathrm{ind}\,s + \mu\,\mathrm{ind}\,t}. \]


% ===== scan leaf 590, printed 565 =====
% !PRINTED 565
\weberpage{565}
If the summation with respect to $t$ is carried out first,
then $st$ may be put in place of $t$, since $s$ is different from zero,
and hence $st$ and $t$ run at the same time through a complete system of residues
with respect to the modulus $n$. Thus
\weq{2}{(\varepsilon^\lambda, r)(\varepsilon^\mu, r) = \sum^{s} \sum^{t} r^{s(t+1)} \varepsilon^{(\lambda+\mu)\,\mathrm{ind}\,s} \varepsilon^{\mu\,\mathrm{ind}\,t}.}
We first dispose of the special case $\mu = -\lambda$, which
we obtain from (2) when we carry out the summation with respect to $s$
first:
\[ (\varepsilon^\lambda, r)(\varepsilon^{-\lambda}, r) = \sum^{t} \varepsilon^{-\lambda\,\mathrm{ind}\,t} \sum^{s} r^{s(t+1)}. \]
Herein
\[
\begin{aligned}
\sum^{s} r^{s(t+1)} &= -1, \quad \text{if } t = 1, 2 \ldots n-2\\
&= n-1, \quad \text{if } t = n-1,
\end{aligned}
\]
hence, since by \S.\,136 $\mathrm{ind}(n-1) = \frac{1}{2}(n-1)$ and
\[ \varepsilon^{\frac{n-1}{2}} = -1 \]
holds,
\[ (\varepsilon^\lambda, r)(\varepsilon^{-\lambda}, r) = -\sum_{1,\,n-1}^{t} \varepsilon^{-\lambda\,\mathrm{ind}\,t} + (-1)^\lambda n. \]
In the sum taken with respect to $t$, $\mathrm{ind}\,t$ runs through a complete
system of residues with respect to the modulus $n-1$, and hence, if $\lambda$ is
assumed \gesperrt{not divisible} by $n-1$,
\[ \sum_{1,\,n-1}^{t} \varepsilon^{-\lambda\,\mathrm{ind}\,t} = 0, \]
and therefore
\weq{3}{(\varepsilon^\lambda, r)(\varepsilon^{-\lambda}, r) = (-1)^\lambda n,}
while, if $\lambda$ is divisible by $n-1$, $(\varepsilon^\lambda, r) = (1, r) = -1$.

We now, then, treat the sum (2) further under the
assumption that $\lambda + \mu$ is not divisible by $(n-1)$.
Then $\sum^{s} \varepsilon^{(\lambda+\mu)\,\mathrm{ind}\,s} = 0$, and in the sum on the
right-hand side of (2) the terms corresponding to the value $t = n-1$
may be omitted. One thus obtains
\weq{4}{
\begin{aligned}
(\varepsilon^\lambda, r)(\varepsilon^\mu, r) &= \sum_{1,\,n-2}^{t} \varepsilon^{\mu\,\mathrm{ind}\,t} \sum_{1,\,n-1}^{s} \varepsilon^{(\lambda+\mu)\,\mathrm{ind}\,s} r^{s(t+1)}\\
&= \sum_{1,\,n-2}^{t} \varepsilon^{\mu\,\mathrm{ind}\,t - (\lambda+\mu)\,\mathrm{ind}(t+1)} \sum^{s} \varepsilon^{(\lambda+\mu)\,\mathrm{ind}\,s(t+1)} r^{s(t+1)}.
\end{aligned}
}
Now, for fixed $t$, $s(t+1)$ runs at the same time as $s$
through a complete system of residues with respect to the modulus $n$, and hence
%CONT>>


% ===== scan leaf 591, printed 566 =====
% !PRINTED 566
\weberpage{566}
%<<CONT
\[ \sum^{s} \varepsilon^{(\lambda+\mu)\,\mathrm{ind}\,s(t+1)} r^{s(t+1)} = \sum^{s} \varepsilon^{(\lambda+\mu)\,\mathrm{ind}\,s} r^s = (\varepsilon^{\lambda+\mu}, r), \]
and formula (4) yields
\weq{5}{\frac{(\varepsilon^\lambda, r)(\varepsilon^\mu, r)}{(\varepsilon^{\lambda+\mu}, r)} = \psi_{\lambda,\mu}(\varepsilon),}
if for abbreviation
\weq{6}{\sum_{1,\,n-2}^{t} \varepsilon^{\mu\,\mathrm{ind}\,t - (\lambda+\mu)\,\mathrm{ind}(t+1)} = \psi_{\lambda,\mu}(\varepsilon)}
is set, and $\lambda + \mu$ is assumed not divisible by $n - 1 = m$.
These functions $\psi$ are then composed of the $(n-1)$th roots of unity
$\varepsilon$ with integral coefficients, and are therefore
known if the $(n-1)$th roots of unity are assumed to be known.
They are numbers of the field $R(\varepsilon)$.

We now assume $\lambda$ divisible by $\mu$, and accordingly replace $\lambda$
by $\mu\lambda$ and set
\weq{7}{\alpha = \varepsilon^\mu,}
so that $\alpha$ may also be a non-primitive $(n-1)$th root of unity,
e.g.\ an $e$th root of unity (if, as above, $m = ef$
and $\mu$ is divisible by $f$, say taken $= f$), where
however the case $\alpha = 1$, that is $e = 1$, is to be excluded.
Then we set
\weq{8}{\sum_{1,\,n-2}^{t} \alpha^{\mathrm{ind}\,t - (\lambda+1)\,\mathrm{ind}(t+1)} = \psi_\lambda(\alpha);}
here $\psi_\lambda(\alpha)$ is a number of the field $R(\alpha)$, which can easily
be computed with the aid of an index table; $\lambda$ can here be taken
with respect to the modulus $e$. Formula (5) gives for this
assumption
\weq{9}{(\alpha^\lambda, r)(\alpha, r) = (\alpha^{\lambda+1}, r)\psi_\lambda(\alpha),}
and this formula holds so long as $\mu(\lambda+1)$ is not divisible by $n - 1$;
thus, if $\mu$ is coprime to $m$, and hence $\alpha$ is a
primitive $m$th root of unity, for $\lambda = 1, 2 \ldots n-3$; and
if $f$ is the greatest common divisor of $\mu$ and $m$ and
$m = ef$, for $\lambda = 1, 2 \ldots e-2$.
Formula (3) now gives, if $\alpha^\lambda$ is not $= 1$,
\weq{10}{(\alpha^\lambda, r)(\alpha^{-\lambda}, r) = (-1)^{\mu\lambda} n.}
Formulae (9), (10) suffice to determine not only the $r$ themselves,
but also all the periods $\eta$, by radicals
whose degrees are the primes dividing $m$,
and among which the functions $\psi_\lambda$ occur as the known
quantities.


% ===== scan leaf 592, printed 567 =====
% !PRINTED 567
\weberpage{567}
Since the $f$-term periods (for $f = 1$) include the quantities $r$
among them, we shall pass at once to the determination of these
periods.

We shall therefore understand by $\alpha$ a primitive $e$th root of unity
and set $\mu = f$.
Then
\[ (\alpha^\lambda, r) = r + \alpha^\lambda r_1 + \alpha^{2\lambda} r_2 + \cdots + \alpha^{(n-2)\lambda} r_{n-2}, \]
or by \S.\,167, (12)
\weq{11}{(\alpha^\lambda, r) = (\alpha^\lambda, \eta) = \eta + \alpha^\lambda \eta_1 + \alpha^{2\lambda} \eta_2 + \cdots + \alpha^{(e-1)\lambda} \eta_{e-1},}
and the formulae (9) and (10) yield
\weq{12}{(\alpha^\lambda, \eta)(\alpha, \eta) = (\alpha^{\lambda+1}, \eta)\psi_\lambda(\alpha) \qquad \lambda = 1, 2 \ldots (e-2),}
\weq{13}{(\alpha^\lambda, \eta)(\alpha^{-\lambda}, \eta) = (-1)^{f\lambda} n.}
Equation (13) shows that none of the numbers $(\alpha^\lambda, \eta)$ vanishes,
and by (12) the $\psi_\lambda(\alpha)$ also are different from zero;
and if we form equation (12) for $\lambda = 1, 2 \ldots \lambda - 1$,
it follows that
% [source-original mathematical notation, verified 2026-08-29] The print uses lambda
% both for the running values and for the fixed upper index in "lambda-1". The displayed
% sequence below makes the intended range clear; the diplomatic meaning is preserved.
\[
\begin{aligned}
(\alpha, \eta)(\alpha, \eta) &= (\alpha^2, \eta)\psi_1(\alpha)\\
(\alpha^2, \eta)(\alpha, \eta) &= (\alpha^3, \eta)\psi_2(\alpha)\\
&\cdots\cdots\cdots\cdots\cdots\cdots\cdots\cdots\cdots\cdots\\
(\alpha^{\lambda-1}, \eta)(\alpha, \eta) &= (\alpha^\lambda, \eta)\psi_{\lambda-1}(\alpha).
\end{aligned}
\]
From this we obtain, by multiplication and cancellation
of the factor $(\alpha^2, \eta) \ldots (\alpha^{\lambda-1}, \eta)$,
\weq{14}{(\alpha, \eta)^\lambda = (\alpha^\lambda, \eta)\psi_1(\alpha)\psi_2(\alpha) \ldots \psi_{\lambda-1}(\alpha).}
By this, $(\alpha^\lambda, \eta)$ is rationally expressed in terms of $(\alpha, \eta)$. But if
we set $\lambda = e - 1$ in (14), there results
\[ (\alpha, \eta)^{e-1} = (\alpha^{-1}, \eta)\psi_1(\alpha)\psi_2(\alpha) \ldots \psi_{e-2}(\alpha), \]
and, by multiplication with $(\alpha, \eta)$, if (13) is used in the form
\[ (\alpha, \eta)(\alpha^{-1}, \eta) = (-1)^f n \]
we obtain
\weq{15}{(\alpha, \eta)^e = (-1)^f n\, \psi_1(\alpha)\psi_2(\alpha) \ldots \psi_{e-2}(\alpha),}
whereby the formation of $(\alpha, \eta)$ is reduced to an $e$th
root. If we set $e = m$, we obtain $(\varepsilon, r)$.

We shall apply this procedure to the interesting case $n = 17$,
which leads to the geometrically remarkable result, first found by
\gesperrt{Gauss}, that the division of the circumference of the circle
into 17 parts depends on a sequence of quadratic equations,
and that therefore the regular seventeen-gon can be
constructed with compass and ruler.


% ===== scan leaf 593, printed 568 =====
% !PRINTED 568
\weberpage{568}

For the prime 17, 3 is a primitive root, and one finds
the index table

\[
\begin{array}{c|cccccccccccccccc}
I & 0 & 1 & 2 & 3 & 4 & 5 & 6 & 7 & 8 & 9 & 10 & 11 & 12 & 13 & 14 & 15\\
\hline
N & 1 & 3 & 9 & 10 & 13 & 5 & 15 & 11 & 16 & 14 & 8 & 7 & 4 & 12 & 2 & 6
\end{array}
.
\]

The two-term periods are, if we set $r = e^{\frac{2\pi i}{17}}$,

\weq{16}{
\begin{aligned}
\eta &= 2\cos\frac{2\pi}{17}, & \eta_1 &= 2\cos\frac{6\pi}{17} \\
\eta_2 &= 2\cos\frac{18\pi}{17} = -2\cos\frac{\pi}{17}, & \eta_3 &= 2\cos\frac{20\pi}{17} = -2\cos\frac{3\pi}{17}, \\
\eta_4 &= 2\cos\frac{26\pi}{17} = 2\cos\frac{8\pi}{17}, & \eta_5 &= 2\cos\frac{10\pi}{17} = -2\cos\frac{7\pi}{17}, \\
\eta_6 &= 2\cos\frac{30\pi}{17} = 2\cos\frac{4\pi}{17}, & \eta_7 &= 2\cos\frac{22\pi}{17} = -2\cos\frac{5\pi}{17}.
\end{aligned}
}

For the computation of the functions $\psi$ one most simply
applies the following table:

\[
\begin{array}{r|ccccccccccccccc}
t & 1 & 2 & 3 & 4 & 5 & 6 & 7 & 8 & 9 & 10 & 11 & 12 & 13 & 14 & 15\\
\hline
\operatorname{ind} t & 0 & 14 & 1 & 12 & 5 & 15 & 11 & 10 & 2 & 3 & 7 & 13 & 4 & 9 & 6 \\
\hline
\operatorname{ind}(t+1) & 14 & 1 & 12 & 5 & 15 & 11 & 10 & 2 & 3 & 7 & 13 & 4 & 9 & 6 & 8
\end{array}
,
\]

from which, if $\alpha$ denotes an 8th root of unity, one finds by
formula (8):

\weq{17}{
\begin{aligned}
\psi_1(\alpha) &= \psi_6(\alpha) = 2\alpha + 2\alpha^2 + 3\alpha^4 + 4\alpha^5 + 2\alpha^6 + 2\alpha^7 \\
\psi_2(\alpha) &= \psi_5(\alpha) = 2 + 3\alpha + \alpha^3 + \alpha^4 + 3\alpha^5 + 4\alpha^6 + \alpha^7 \\
\psi_3(\alpha) &= \psi_4(\alpha) = 3 + 3\alpha + 2\alpha^2 + 3\alpha^3 + \alpha^5 + 2\alpha^6 + \alpha^7.
\end{aligned}
}

If one takes

\[
\alpha = \frac{1+i}{\sqrt{2}}, \qquad \alpha^2 = i, \qquad \alpha^4 = -1,
\]

there results

\weq{18}{
\begin{aligned}
\psi_1(\alpha) &= \psi_6(\alpha) = -3 - i\sqrt{8} \\
\psi_2(\alpha) &= \psi_5(\alpha) = 1 - 4i \\
\psi_3(\alpha) &= \psi_4(\alpha) = 3 + i\sqrt{8} \\
\psi_1(i) &= -1 + 4i,
\end{aligned}
}

hence by (15)

\weq{19}{
(\alpha, \eta)^8 = 17\,(3 + i\sqrt{8})^4\,(1 - 4i)^2.
}

But from (13) and (14) one obtains
%CONT>>


% ===== scan leaf 594, printed 569 =====
% !PRINTED 569
\weberpage{569}
%<<CONT
\weq{20}{
\begin{aligned}
(-1,\eta)^2 &= 17 \\
(i,\eta)^2 &= (-1,\eta)(-1+4i) \\
(\alpha,\eta)^2 &= -(i,\eta)(3+i\sqrt{8});
\end{aligned}
}

in addition one has $(1,\eta) = -1$,

\weq{21}{
\begin{aligned}
(i,\eta)(-i,\eta) &= 17 \\
(\alpha,\eta)(\alpha^{-1},\eta) &= 17 \\
(\alpha^3,\eta)(\alpha^{-3},\eta) &= 17
\end{aligned}
}

and by (9)

\weq{22}{
(\alpha,\eta)(\alpha^3,\eta) = (-1,\eta)(3+i\sqrt{8}),
}

whereby $(1,\eta)$, $(-1,\eta)$, $(\pm i,\eta)$, $(\alpha,\eta)$, $(\alpha^{-1},\eta)$, $(\alpha^{-3},\eta)$ are
determined by square roots.
% [source-original mathematical omission, verified 2026-08-29] Formula (22), together
% with the already determined (alpha,eta), also determines (alpha^3,eta), but the printed
% enumeration omits that quantity while retaining its reciprocal partner (alpha^-3,eta).

With the determination that we have made concerning $r$ and $\alpha$,
it follows easily from a consideration of the values of the cosines
(16) that $(-1,\eta)$ is positive, and hence equal to the positive square root
$\sqrt{17}$, and further that $(i,\eta)$ and $(\alpha,\eta)$ have positive real
parts, whereby these quantities too are completely determined.
We have therefore

\[
\begin{aligned}
(-1,\eta) &= \sqrt{17} \\
(i,\eta) &= \sqrt[4]{17}\left\{\sqrt{\frac{\sqrt{17}-1}{2}} + i\sqrt{\frac{\sqrt{17}+1}{2}}\right\},
\end{aligned}
\]

if all the roots are taken positive; and by the last
equation (20), $(\alpha,\eta)$ can be represented, by a formula which is,
to be sure, somewhat long, through square roots of real quantities.
We shall not enter into the geometric constructions that
attach themselves here\footnote{Cf. on this v. \gesperrt{Staudt}, Construction des regulären Siebzehnecks. Crelle's Journ. Bd. 24, 1842.}.

The determination of signs employed is obtained simply
if, following (16), one sets

\[
\frac{1}{2}(-1,\eta) = \cos\frac{2\pi}{17} - \cos\frac{\pi}{17} + \cos\frac{8\pi}{17} + \cos\frac{4\pi}{17}
- \cos\frac{6\pi}{17} + \cos\frac{3\pi}{17} + \cos\frac{7\pi}{17} + \cos\frac{5\pi}{17}.
\]

But by a well-known trigonometric formula,

\[
\begin{aligned}
\cos\frac{2\pi}{17} - \cos\frac{\pi}{17} &= -2\sin\frac{\pi}{34}\sin\frac{3\pi}{34} \\
-\cos\frac{6\pi}{17} + \cos\frac{5\pi}{17} &= 2\sin\frac{\pi}{34}\sin\frac{11\pi}{34},
\end{aligned}
\]
%CONT>>


% ===== scan leaf 595, printed 570 =====
% !PRINTED 570
\weberpage{570}
\pdfbookmark[2]{§ 170. Properties of the numbers ψ}{bm190}
%<<CONT
and since $\sin\frac{11\pi}{34}$ is greater than $\sin\frac{3\pi}{34}$, the sum of these
two expressions, and with it the whole sum $(-1,\eta)$, is positive % [sic] the printed authority ends the sentence without punctuation

There results further, for the real part of $(i,\eta)$,

\[
\eta - \eta_2 + \eta_4 - \eta_6 = 2\cos\frac{2\pi}{17} + 2\cos\frac{\pi}{17} + 2\cos\frac{8\pi}{17} - 2\cos\frac{4\pi}{17},
\]

which is immediately recognised as positive; finally, for the real
part of $(\alpha,\eta)$,

\[
\eta - \eta_4 + \frac{1}{\sqrt{2}}(\eta_1 - \eta_3 - \eta_5 + \eta_7)
= 2\cos\frac{2\pi}{17} - 2\cos\frac{8\pi}{17} + \sqrt{2}\left(\cos\frac{6\pi}{17} + \cos\frac{3\pi}{17} + \cos\frac{7\pi}{17} - \cos\frac{5\pi}{17}\right),
\]

which likewise proves to be positive.

\begin{center}
\S.\,170.

\gesperrt{Properties of the Numbers} $\psi$.
\end{center}

The formulae for the resolvents derived in the preceding
section admit a wealth of applications, even without an exact
determination of the prime $n$, for special assumptions
about the number $e$; they thus give number-theoretic theorems for
whole categories of primes.

For algebra they provide arbitrarily many solutions of cyclic
equations of prescribed degree.

We precede the derivation of the most important of these applications
for algebra with some general considerations on
the functions $\psi$ which we defined in the preceding
section.

For the computation of the numbers $\psi_{\lambda,\mu}(\varepsilon)$ [\S.\,169, (5)] one makes
use, as we have seen in the case $n = 17$, of an index table.
But between these functions there exist relations by
which they can be reduced to a smaller number, on which
\gesperrt{Jacobi} has made remarkable investigations\footnote{Cf. \gesperrt{Kronecker}, Zur Theorie der \gesperrt{Abel'}schen Gleichungen. Journ. f. Mathem. Bd. 93.}.

The function $\psi_{\lambda,\mu}(\varepsilon)$ does not change, according to its definition,
when $\lambda,\mu$ are interchanged, although its expression
%CONT>>


% ===== scan leaf 596, printed 571 =====
% !PRINTED 571
\weberpage{571}
%<<CONT
[\S.\,169, (6)] thereby takes on another form. We thus obtain
the same function $\psi_\lambda(\alpha)$ when we replace $\lambda,\mu$ by $\lambda f, f$ or
by $f, \lambda f$, hence [\S.\,169, (8)]

% [resolved 2026-08-06; confirmed 2026-08-07 at 600 dpi] CORRECTED THROUGHOUT THIS LEAF.
% All four sigmas in the two displays carry "t" ABOVE only; what sits beneath each is a short
% horizontal RULE under the sigma, not a second "t". The original transcription "\sum_t^t"
% duplicated the index. Reproduced here as "t" above, with the printed rule recorded in this
% note rather than faked in LaTeX.
\[
\psi_\lambda(\alpha) = \sum^{t} \alpha^{\operatorname{ind} t - (\lambda+1)\operatorname{ind}(t+1)}
= \sum^{t} \alpha^{\lambda \operatorname{ind} t - (\lambda+1)\operatorname{ind}(t+1)}.
\]

If, then, two numbers $\lambda,\lambda'$ are connected with one another by the
congruence

\weq{1}{
\lambda\lambda' \equiv 1 \pmod{e}
}

it follows that

\[
\psi_\lambda(\alpha^{\lambda'}) = \sum^{t} \alpha^{\lambda\lambda' \operatorname{ind} t - (\lambda\lambda'+\lambda')\operatorname{ind}(t+1)}
= \sum^{t} \alpha^{\operatorname{ind} t - (\lambda'+1)\operatorname{ind}(t+1)},
\]

hence the formula

\weq{2}{
\psi_\lambda(\alpha^{\lambda'}) = \psi_{\lambda'}(\alpha).
}

If we then put $-\lambda-\mu$ for $\lambda$ in \S.\,169, (5), and
hence $-\lambda$ for $\lambda+\mu$ (where it must then be assumed that
$\lambda$ is not divisible by $m$), it follows that

\[
\frac{(\varepsilon^{-\lambda-\mu}, r)(\varepsilon^\mu, r)}{(\varepsilon^{-\lambda}, r)} = \psi_{-\lambda-\mu,\mu}(\varepsilon).
\]

We multiply the left side by

\[
\frac{(\varepsilon^{\lambda+\mu}, r)(\varepsilon^\lambda, r)}{(\varepsilon^{\lambda+\mu}, r)(\varepsilon^\lambda, r)} = 1,
\]

and apply in the numerator and denominator the formula

\[
(\varepsilon^\lambda, r)(\varepsilon^{-\lambda}, r) = (-1)^\lambda n
\]

[\S.\,169, (3)]; then it follows that

\[
(-1)^\mu \frac{(\varepsilon^\lambda, r)(\varepsilon^\mu, r)}{(\varepsilon^{\lambda+\mu}, r)} = \psi_{-\lambda-\mu,\mu}(\varepsilon),
\]

or

\[
\psi_{-\lambda-\mu,\mu}(\varepsilon) = (-1)^\mu \psi_{\lambda,\mu}(\varepsilon),
\]

and if we replace $\lambda,\mu$ by $\lambda f, f$ and determine $\lambda''$ from the
congruence

\weq{3}{
\lambda + \lambda'' + 1 \equiv 0 \pmod{e}
}

we obtain

\weq{4}{
\psi_\lambda(\alpha) = (-1)^f \psi_{\lambda''}(\alpha).
}

This formula can then in turn be combined with formula (2),
and the number of functions $\psi_\lambda$ to be computed for a given
$e$, whose number amounts to $e-2$, thereby be reduced to about
the sixth part.

A third formula results when one puts $\varepsilon^{-1}$ for
$\varepsilon$ in \S.\,169, (5) and then multiplies, thus
%CONT>>


% ===== scan leaf 597, printed 572 =====
% !PRINTED 572
\weberpage{572}

\[
\psi_{\lambda,\mu}(\varepsilon)\,\psi_{\lambda,\mu}(\varepsilon^{-1})
= \frac{(\varepsilon^\lambda,\,r)(\varepsilon^{-\lambda},\,r)(\varepsilon^\mu,\,r)(\varepsilon^{-\mu},\,r)}
       {(\varepsilon^{\lambda+\mu},\,r)(\varepsilon^{-\lambda-\mu},\,r)}
\]
is formed. And with the help of \S.\,169, (3) one thus finds the formula
\weq{5}{\psi_{\lambda,\mu}(\varepsilon)\,\psi_{\lambda,\mu}(\varepsilon^{-1}) = n,}
which is now again specialised by the substitution $\lambda f, f$ for $\lambda, \mu$:
\weq{6}{\psi_\lambda(\alpha)\,\psi_\lambda(\alpha^{-1}) = n.}
Finally we wish to mention one further property of the functions $\psi$, which can also be read off immediately from their definition \S.\,169, (5). According to this, namely, if $\lambda, \mu, \nu$ are three numbers such that $\lambda + \mu$, $\lambda + \mu + \nu$ is not divisible by $m$,
\weq{7}{\psi_{\lambda,\mu}(\varepsilon)\,\psi_{\lambda+\mu,\nu}(\varepsilon) = \frac{(\varepsilon^\lambda,\,r)(\varepsilon^\mu,\,r)(\varepsilon^\nu,\,r)}{(\varepsilon^{\lambda+\mu+\nu},\,r)},}
i.\,e. the product $\psi_{\lambda,\mu}(\varepsilon)\,\psi_{\lambda+\mu,\nu}(\varepsilon)$ remains unchanged when the three numbers $\lambda, \mu, \nu$ are permuted arbitrarily (provided that $\lambda + \nu$ and $\mu + \nu$ are also not divisible by $m$). From this we obtain, e.\,g.,
\[
\psi_{\lambda,\mu}(\varepsilon)\,\psi_{\lambda+\mu,\nu}(\varepsilon) = \psi_{\mu,\nu}(\varepsilon)\,\psi_{\mu+\nu,\lambda}(\varepsilon),
\]
and if one replaces $\lambda, \mu, \nu$ in it by $2\lambda f, f, f$, and observes that $\psi_{2f,2\lambda f}(\varepsilon) = \psi_\lambda(\alpha^2)$,
\weq{8}{\psi_{2\lambda}(\alpha)\,\psi_{2\lambda+1}(\alpha) = \psi_1(\alpha)\,\psi_\lambda(\alpha^2),}
provided that $2, 2\lambda + 1, 2\lambda + 2$ are not divisible by $e$.

\gesperrt{Jacobi} made the attempt, with the help of these formulae, to express, for a given $e$ assumed to be a prime, all the $e - 2$ functions $\psi_\lambda(\alpha)$ by the conjugate values of a single one, i.\,e. by $\psi_1(\alpha), \psi_1(\alpha^2), \psi_1(\alpha^3) \ldots$, and carried the computation through as far as $e = 23$.

His conjecture, however, that such a representation is possible in general, is, according to the investigations of \gesperrt{Kronecker}, probably not correct.

We wish to carry through the case $e = 7$ as an example. For this case
\[
\begin{aligned}
\lambda &= 1, 2, 3, 4, 5\\
\lambda' &= 1, 4, 5, 2, 3\\
\lambda'' &= 5, 4, 3, 2, 1,
\end{aligned}
\]
so the formulae (2) and (4) yield
\[
\begin{aligned}
&\psi_4(\alpha) = \psi_2(\alpha^4), \quad \psi_5(\alpha) = \psi_3(\alpha^5)\\
&\psi_1(\alpha) = \psi_5(\alpha), \quad \psi_2(\alpha) = \psi_4(\alpha),
\end{aligned}
\]
and from this further
%CONT>>


% ===== scan leaf 598, printed 573 =====
% !PRINTED 573
\weberpage{573}

%<<CONT
\[
\begin{aligned}
\psi_1(\alpha) &= \psi_3(\alpha^5)\\
\psi_3(\alpha) &= \psi_1(\alpha^3),
\end{aligned}
\]
and if in (8) one sets $\lambda = 1$
\[
\psi_2(\alpha)\,\psi_3(\alpha) = \psi_1(\alpha)\,\psi_1(\alpha^2);
\]
finally, by (6)
\[
n = \psi_3(\alpha)\,\psi_3(\alpha^6) = \psi_3(\alpha)\,\psi_1(\alpha^4).
\]
If one now computes accordingly the formula \S.\,169, (15)
\[
(\alpha, \eta)^7 = n\,\psi_1(\alpha)\,\psi_2(\alpha)\,\psi_3(\alpha)\,\psi_4(\alpha)\,\psi_5(\alpha),
\]
there results
\weq{9}{(\alpha, \eta)^7 = \psi_1(\alpha)^4\,\psi_1(\alpha^2)^2\,\psi_1(\alpha^4).}
By this formula the relation from \S.\,169, (14)
\[
(\alpha, \eta)^{14} = (\alpha^2, \eta)^7\,\psi_1(\alpha)^7
\]
can be verified immediately.

We shall conclude these considerations on the functions $\psi$ with the proof of a theorem that will be useful to us later.

If we determine a positive integer $\nu < m$ from the congruence
\weq{10}{\lambda + \mu + \nu \equiv 0 \pmod{m},}
then we can represent the function $\psi_{\lambda,\mu}(\varepsilon)$, as it is defined by \S.\,169, (6), in the following way
\weq{11}{\psi_{\lambda,\mu}(\varepsilon) = \sum_{1,\,n-2}^{t} \varepsilon^{\mu\,\mathrm{ind}\,t + \nu\,\mathrm{ind}(t+1)}.}
In this we now wish to replace the $m$th root of unity $\varepsilon$ by the primitive congruence root $g$ of the prime $n$ on which the index system is based, i.\,e.\ we wish to form the integer
\weq{12}{\psi_{\lambda,\mu}(g) = \sum_{1,\,n-2}^{t} g^{\mu\,\mathrm{ind}\,t + \nu\,\mathrm{ind}(t+1)} \equiv \sum_{1,\,n-2}^{t} t^\mu(t+1)^\nu \pmod{n}}
Considering this number with respect to the modulus $n$, we can extend the summation to $t = n - 1$, since for this value $(t + 1)$ is congruent to zero. We thus obtain
\weq{13}{\psi_{\lambda,\mu}(g) \equiv \sum_{1,\,n-1}^{t} t^\mu(t+1)^\nu \pmod{n}.}
To the power $(t + 1)^\nu$ we can apply the binomial theorem, and we obtain, if, as before, $B_h^{(\nu)}$ denote the binomial coefficients,
\weq{14}{\psi_{\lambda,\mu}(g) \equiv \sum_{0,\,\nu}^{h} B_h^{(\nu)} \sum_{1,\,n-1}^{t} t^{\mu+h} \pmod{n}.}


% ===== scan leaf 599, printed 574 =====
% !PRINTED 574
\weberpage{574}
\pdfbookmark[2]{§ 171. The Gaussian sums}{bm191}

If we set $t \equiv g^s \pmod{n}$, then $s$ runs through a system of residues with respect to the modulus $m$, and we obtain
\[
\sum_{1,\,n-1}^{t} t^{\mu+h} \equiv \sum_{0,\,m-1}^{s} g^{s(\mu+h)} \pmod{n}.
\]
But by the summation formula for the geometric series
\[
(g^{\mu+h} - 1) \sum_{0,\,m-1}^{s} g^{s(\mu+h)} = g^{m(\mu+h)} - 1 \equiv 0 \pmod{n},
\]
and consequently
\weq{15}{\sum_{0,\,m-1}^{s} g^{s(\mu+h)} \equiv 0 \pmod{n},}
except when $\mu + h \equiv 0 \pmod{m}$, and in this case
\weq{16}{\sum_{0,\,m-1}^{s} g^{s(\mu+h)} \equiv m \equiv -1 \pmod{n}.}
If we take $\lambda, \mu$ between $0$ and $m$, then
\begin{quote}
1) when $\lambda + \mu < m$, \quad $\nu = m - \lambda - \mu$,\\
2) when $\lambda + \mu > m$, \quad $\nu = 2m - \lambda - \mu$.
\end{quote}
The exponent $\mu + h$ in the sum (14) therefore runs, in case 1), through the values $\mu + h = \mu, \mu + 1, \ldots m - \lambda$, and in case 2) through $\mu, \mu + 1, \ldots 2m - \lambda$.

In case 1), therefore, the case $\mu + h \equiv 0 \pmod{m}$ assumed in (16) does not occur at all; in case 2) it occurs once, for $h = m - \mu$.

Accordingly, from (14), having regard to (15) and (16), if for the binomial coefficient $B_h^{(\nu)}$ the value $\dfrac{\Pi(\nu)}{\Pi(h)\,\Pi(\nu-h)}$ is put, we have
\weq{17}{
\begin{aligned}
&\text{1)}\quad \psi_{\lambda,\mu}(g) \equiv 0, \quad \lambda + \mu < m\\
&\text{2)}\quad \psi_{\lambda,\mu}(g) \equiv -\frac{\Pi(2m-\lambda-\mu)}{\Pi(m-\lambda)\,\Pi(m-\mu)}, \quad m < \lambda + \mu < 2m.
\end{aligned}
\quad \pmod{n}}

\begin{center}
\S.\ 171.

\gesperrt{The Gauss Sums.}
\end{center}

In \S.\,169 it is shown that, if $n - 1 = ef$, the periods $\eta, \eta_1, \ldots \eta_{e-1}$ of an equation of degree $e$ with integral coefficients
%CONT>>


% ===== scan leaf 600, printed 575 =====
% !PRINTED 575
\weberpage{575}

%<<CONT
satisfy. This equation has only real roots if $f$ is even, because then $\tfrac{1}{2}(n - 1)$ is a multiple of $e$, and consequently $r$ and $r^{-1}$ occur in the same period. But if $f$ is odd, all the roots are imaginary.

It is now of the highest interest to investigate this equation of degree $e$ for individual particular values of $e$, without any special assumption about the prime $n$ beyond the requirement that $n - 1$ shall be divisible by $e$.

We consider the first special cases, and assume to begin with $e = 2$, which is admissible for every odd prime $n$; here $\eta, \eta_1$ are thus the two periods of $\tfrac{1}{2}(n - 1)$ terms, which we now wish to denote by $A, B$, so that
\[
\begin{aligned}
A &= r + r_2 + r_4 + \cdots + r_{n-3}\\
B &= r_1 + r_3 + r_5 + \cdots + r_{n-2}.
\end{aligned}
\]
The $r, r_2, r_4 \ldots$ have the exponents $g^0, g^2 \ldots g^{n-3}$, i.\,e. the exponents of $g$ are even numbers. The exponents of $r$ are therefore the \gesperrt{quadratic residues} of $n$ (\S.\,138). Likewise the exponents of $r$ in the sum $B$ are the non-residues. If, then, we denote the residues by $a$ and the non-residues by $b$, we have
\weq{1}{A = \Sigma r^a, \quad B = \Sigma r^b,}
and $(-1, r) = A - B$. These expressions $A, B$, are called the \gesperrt{Gauss sums}.

If in the formula \S.\,169, (10) we make the assumption
\[
e = 2, \quad \mu = f = \frac{n-1}{2}, \quad \lambda = 1, \quad \alpha = -1,
\]
there results
\weq{2}{A - B = \pm \sqrt{(-1)^{\frac{n-1}{2}}\, n},}
while on the other hand
\[
A + B = -1,
\]
hence
\weq{3}{2A = -1 \pm \sqrt{(-1)^{\frac{n-1}{2}}\, n}, \quad 2B = -1 \mp \sqrt{(-1)^{\frac{n-1}{2}}\, n}.}
The sign that is to be given to the root depends on the choice of $r$. But once $r$ has been fixed upon, the sign is completely determined. Its ascertainment presents peculiar difficulties, which gave occasion to a separate memoir by \gesperrt{Gauss}\textsuperscript{1}).

\noindent\rule{3cm}{0.4pt}

\noindent\textsuperscript{1}) Summatio quarundam serierum singularium. \gesperrt{Gauss} Werke, Bd.\ II. In another way, by application of the higher analysis, the sign
%CONT>>


% ===== scan leaf 601, printed 576 =====
% !PRINTED 576
\weberpage{576}

Let $\mu$ and $\nu$ be two positive integers and $x$ a
variable. We define a rational function of $x$ by
the equation

\weq{4}{(\mu, \nu) = \frac{(1 - x^{\mu})(1 - x^{\mu - 1}) \ldots (1 - x^{\mu - \nu + 1})}{(1 - x)(1 - x^2) \ldots (1 - x^{\nu})}.}

If $\nu = \mu + 1$, this function vanishes
identically, on account of the factor $1 - x^{\mu - \nu + 1}$, and likewise when $\nu$
becomes still greater.

We set up for these functions some recursion formulae,
which follow immediately from an inspection of formula (4):

\weq{5}{(1 - x^{\mu - \nu + 1})(\mu + 1, \nu) = (1 - x^{\mu + 1})(\mu, \nu)}

\weq{6}{(1 - x^{\nu + 1})\ (\mu, \nu + 1) = (1 - x^{\mu - \nu})(\mu, \nu),}

and by the combination of (5) and (6), after $\mu$ has been
replaced by $\mu + 1$ in (6),

\weq{7}{(1 - x^{\nu + 1})(\mu + 1, \nu + 1) = (1 - x^{\mu + 1})(\mu, \nu).}

If one subtracts (6) from (7) and divides by the non-vanishing $1 - x^{\nu + 1}$, it follows further

\weq{8}{(\mu + 1, \nu + 1) = (\mu, \nu + 1) + x^{\mu - \nu}(\mu, \nu).}

If we set, by way of definition, $(\mu, 0) = 1$, and $(\mu, \nu) = 0$ for
negative $\nu$, then these formulae all hold generally for every
integer $\nu$.

From the last formula we conclude by complete
induction that $(\mu, \nu)$ is an \gesperrt{integral function} of $x$. For
if this property is proved for any one $\mu$, it follows for
the next greater $\mu$ from formula (8). But, for example,
$(1, \nu)$ is equal to 1 or equal to 0, and so the theorem is generally correct.

If we put $\mu + 1$ in place of $\mu$ in (8), we can
with its help transform formula (5) thus,

\weq{9}{(1 - x^{\mu + 1})(\mu, \nu) = (\mu + 1, \nu) + (\mu + 1, \nu + 1) - (\mu + 2, \nu + 1).}

We now consider an integral function

\weq{10}{f(x, \mu) = 1 - (\mu, 1) + (\mu, 2) - (\mu, 3) + \cdots = \overset{\nu}{\Sigma} (-1)^{\nu} (\mu, \nu),}

determined by \gesperrt{Dirichlet}, \gesperrt{Cauchy}, \gesperrt{Kronecker}. The path that
\gesperrt{Gauss} takes, which we follow in the text, uses only algebraic
means. \gesperrt{Kronecker} has given in Liouville's Journal Ser.\,2, Bd.\,1 a
way of determining the sign which, with the use of a
remark of \gesperrt{Dedekind} (Schlömilch's Ztschr. Bd.\,15, Literaturzeitung
S.\,21), likewise becomes a purely algebraic one.
%CONT>>


% ===== scan leaf 602, printed 577 =====
% !PRINTED 577
\weberpage{577}

%<<CONT
in which the summation over $\nu$ extends from 0 to $\mu$, but may
also be extended arbitrarily further, since $(\mu, \nu)$ vanishes
when $\nu$ is negative or greater than $\mu$. Accordingly we also have

\[ f(x, \mu) = -\overset{\nu}{\Sigma} (-1)^{\nu} (\mu, \nu + 1), \]

and if we therefore multiply formula (9) by $(-1)^{\nu}$ and
sum with respect to $\nu$, there results

\weq{11}{(1 - x^{\mu + 1}) f(x, \mu) = f(x, \mu + 2).}

From this it follows that $f(x, \mu)$ vanishes for every odd $\mu$, since $f(x, 1) = 0$. On the other hand, there results directly
from (10)

\[ f(x, 2) = 1 - x, \]

and therefore from (11), for every even $\mu$,

\weq{12}{f(x, \mu) = (1 - x)(1 - x^3)(1 - x^5) \ldots (1 - x^{\mu - 1}).}

For if this formula is correct for $\mu$, its correctness for
$\mu + 2$ follows from (11).

We now apply this identical transformation to our
problem, by setting $\mu = m = n - 1$ and $x = r$. Then

\[ \frac{1 - x^m}{1 - x} = \frac{1 - r^{-1}}{1 - r} = -r^{-1}, \quad \frac{1 - x^{m - 1}}{1 - x^2} = \frac{1 - r^{-2}}{1 - r^2} = -r^{-2} \cdots, \]

hence, so long as $0 \leqq \nu \leqq m$,

\[ (m, \nu) = (-1)^{\nu} r^{-(1 + 2 + \cdots + \nu)} = (-1)^{\nu} r^{-\frac{\nu(\nu + 1)}{2}}, \]

or, replacing $r$ by $r^{-2}$,

\[ (m, \nu) = (-1)^{\nu} r^{\nu^2 + \nu}. \]

Now

\[ \nu^2 + \nu \equiv \nu^2 - (n - 1) \nu \equiv \left( \nu - \frac{n - 1}{2} \right)^2 - \left( \frac{n - 1}{2} \right)^2 \ (\mathrm{mod}\ n), \]

hence

\[ (m, \nu) = (-1)^{\nu} r^{-\frac{m^2}{4}} r^{\left( \nu - \frac{m}{2} \right)^2}, \]

and it therefore becomes, by the defining equation (10),

\[ f(r^{-2}, m) = r^{-\frac{m^2}{4}} \Sigma r^{\left( \nu - \frac{m}{2} \right)^2}, \]

in which $\nu$ runs through a complete system of residues with respect to the modulus $n$.
Now simultaneously with $\nu$, $\nu - {}^1/_2 m$ also runs through such a
system of residues, so that we may also write

\[ f(r^{-2}, m) = r^{-\frac{m^2}{4}} \Sigma r^{\nu^2} \]
%CONT>>


% ===== scan leaf 603, printed 578 =====
% !PRINTED 578
\weberpage{578}

%<<CONT
can set. Now since in the sequence of numbers $\nu^2$ the number 0
occurs, and moreover every quadratic residue occurs twice [for
$\nu^2$ and $(-\nu)^2$], we have

\[ r^{\frac{m^2}{4}} f(r^{-2}, m) = 1 + 2A = A - B, \]

and from formula (12) it follows that

\[ A - B = r^{\frac{m^2}{4}} (1 - r^{-2})(1 - r^{-6}) \ldots (1 - r^{-2n + 4}). \]

Now

\[ r^{\frac{m^2}{4}} = r\, r^3 r^5 \ldots r^{n - 2}, \]

and so we finally obtain the formula

\weq{13}{A - B = (r - r^{-1})(r^3 - r^{-3}) \ldots (r^{n - 2} - r^{-n + 2}).}

If $\nu$ denotes an odd and $\mu$ an even number between
0 and ${}^1/_2 n$, this product consists of the factors

\[ r^{\nu} - r^{-\nu}, \quad r^{n - \mu} - r^{-n + \mu} = -(r^{\mu} - r^{-\mu}). \]

The number of factors of the second category is as great
as the number of even numbers between 0 and ${}^1/_2 n$, or
equal to the number of positive integers that are smaller than
${}^1/_4 n$. If therefore, understanding by $l$ an integer, we set
$n = 4l + 1$ or $4l + 3$, then $l$ is this number, and
we can set

\[ A - B = (-1)^{l} \overset{\nu}{\Pi} (r^{\nu} - r^{-\nu}) \overset{\mu}{\Pi} (r^{\mu} - r^{-\mu}), \]

for which, if $h$ runs through the sequence of numbers from 1 to ${}^1/_2 (n - 1)$,
we may also set

\[ A - B = (-1)^{l} \overset{h}{\underset{1,\, \frac{n - 1}{2}}{\Pi}} (r^{h} - r^{-h}). \]

If we now set

\[ r = e^{\frac{2 \pi i k}{n}}, \quad r^{h} - r^{-h} = 2i \sin \frac{2 \pi k h}{n}, \]

then

\[ A - B = (-1)^{l} (2i)^{\frac{n - 1}{2}} \overset{h}{\Pi} \sin \frac{2 \pi h k}{n}. \]

The product that still occurs here we have already
determined earlier [\S.\,138, (4)]:

\[ 2^{\frac{n - 1}{2}} \overset{h}{\Pi} \sin \frac{2 \pi k h}{n} = \left( \frac{k}{n} \right) \sqrt{n}, \]

wherein the $\sqrt{n}$ is to be taken positive.


% ===== scan leaf 604, printed 579 =====
% !PRINTED 579
\weberpage{579}
\pdfbookmark[2]{§ 172. The periods of 1!/3,(n-1) and 1!/4,(n-1) terms}{bm192}

If we further set ${}^1/_2 (n - 1) = 2l$ or $= 2l + 1$, according as
$n = 4l + 1$ or $4l + 3$, we now obtain, with
exact determination of sign,

\weq{14}{\begin{aligned} A - B &= \left( \frac{k}{n} \right) \sqrt{n}, & n \equiv 1 \ (\mathrm{mod}\ 4) \\ &= i \left( \frac{k}{n} \right) \sqrt{n}, & n \equiv 3 \ (\mathrm{mod}\ 4). \end{aligned}}

\begin{center}
\S.\,172.
\end{center}

\begin{center}
\gesperrt{The periods of} ${}^1/_3 (n - 1)$ \gesperrt{and} ${}^1/_4 (n - 1)$ \gesperrt{terms.}
\end{center}

We now pass to the case $e = 3$, in which $n - 1$
must be assumed divisible by 3, so that $n = 7, 13, 19,
31, 37, 43 \ldots$ We denote by $\varrho$ an imaginary third
root of unity:

\[ \varrho = \frac{-1 + \sqrt{-3}}{2}, \]

and determine the three periods $\eta, \eta_1, \eta_2$ of ${}^1/_3 (n - 1)$
terms each, which, as we have seen in \S.\,168, are the roots
of a cyclic cubic equation. By \S.\,169 we have

\[ \psi_1(\varrho) = \overset{t}{\underset{1,\, n - 2}{\Sigma}} \varrho^{\mathrm{ind}\, t - 2\, \mathrm{ind}(t + 1)}, \]

for which, since $-2 \equiv 1 \ (\mathrm{mod}\ 3)$, we may also set

\weq{1}{\psi_1(\varrho) = \overset{t}{\underset{1,\, n - 2}{\Sigma}} \varrho^{\mathrm{ind}(t^2 + t)}}

This number $\psi_1(\varrho)$ can be represented in one of the two
forms

\weq{2}{\psi_1(\varrho) = a + b\varrho = \frac{A + b\sqrt{-3}}{2}}

in which $a, b, A$ are integers and $A = 2a - b$.

We then have

\weq{3}{\psi_1(\varrho^2) = a + b\varrho^2 = \frac{A - b\sqrt{-3}}{2},}

and the formulae (10), (15) of \S.\,169 yield

\weq{4}{\eta + \eta_1 + \eta_2 = -1}

\weq{5}{(\varrho, \eta)^3 = n \psi_1(\varrho), \quad (\varrho^2, \eta)^3 = n \psi_1(\varrho^2)}


% ===== scan leaf 605, printed 580 =====
% !PRINTED 580
\weberpage{580}
\weq{6}{(\varrho,\eta)(\varrho^2,\eta) = n, \quad \psi_1(\varrho)\psi_1(\varrho^2) = n,}
in which

\weq{7}{(\varrho,\eta) = \eta + \varrho\eta_1 + \varrho^2\eta_2, \quad (\varrho^2,\eta) = \eta + \varrho^2\eta_1 + \varrho\eta_2.}

For example, for $n = 7, 13$, if one takes the primitive roots 3, 2 as basis and applies the index tables, one obtains:

\[
n = 7
\]
\[
\begin{array}{c|c|c|c|c|c|c}
N & 1 & 2 & 3 & 4 & 5 & 6 \\
\hline
I & 0 & 2 & 1 & 4 & 5 & 3
\end{array}
\]
\[
n = 13
\]
\[
\begin{array}{c|c|c|c|c|c|c|c|c|c|c|c}
1 & 2 & 3 & 4 & 5 & 6 & 7 & 8 & 9 & 10 & 11 & 12 \\
\hline
0 & 1 & 4 & 2 & 9 & 5 & 11 & 3 & 8 & 10 & 7 & 6
\end{array}
\]

\[
\text{for } n = 7, \quad a + b\varrho = -(1 + 3\varrho),
\]
\[
\text{for } n = 13, \quad a + b\varrho = -(4 + 3\varrho).
\]

From these formulae we can easily derive the cubic equation whose roots are the $\eta, \eta_1, \eta_2$. On account of (4) it has the form
\weq{8}{\eta^3 + \eta^2 + \beta\eta + \gamma = 0,}

and the integers $\beta, \gamma$ are to be determined.

If we carry out the multiplication in (6), there results
\[
\begin{aligned}
n &= \eta^2 + \eta_1^2 + \eta_2^2 - \eta\eta_1 - \eta\eta_2 - \eta_1\eta_2 \\
  &= (\eta + \eta_1 + \eta_2)^2 - 3\beta,
\end{aligned}
\]
hence
\weq{9}{3\beta = -(n - 1).}

Carrying out the cubes in (5) yields, if for the moment we set
\[
s_3 = \eta^3 + \eta_1^3 + \eta_2^3, \quad s = \eta^2\eta_1 + \eta_1^2\eta_2 + \eta_2^2\eta,
\]
\[
s' = \eta_1^2\eta + \eta_2^2\eta_1 + \eta^2\eta_2
\]
the relation
\[
n\psi_1(\varrho) = s_3 - 6\gamma + 3s\varrho + 3s'\varrho^2.
\]
But from (8) and (9) one obtains $s_3 = -n - 3\gamma$, hence
\weq{10}{n[\psi_1(\varrho) + 1] = -9\gamma + 3s\varrho + 3s'\varrho^2,}
and likewise
\weq{11}{n[\psi_1(\varrho^2) + 1] = -9\gamma + 3s\varrho^2 + 3s'\varrho,}
to which is added, if one raises (4) to the cube,
\[
n - 1 = -9\gamma + 3s + 3s'
\]
as well. If one adds the last three equations, it follows that
\[
n[\psi_1(\varrho) + \psi_1(\varrho^2)] + 3n - 1 = -27\gamma,
\]
%CONT>>


% ===== scan leaf 606, printed 581 =====
% !PRINTED 581
\weberpage{581}
%<<CONT
or finally, by (2) and (3),
\weq{12}{-27\gamma = nA + 3n - 1.}

From this relation it follows, since $\gamma$ is an integer,
\weq{13}{A \equiv 1\ (\mathrm{mod}\ 3)}
and from this $A^3 \equiv 1\ (\mathrm{mod}\ 9)$. Now by (2), (3) and (6)
\weq{14}{n = a^2 - ab + b^2}
or
\weq{15}{4n = A^2 + 3b^2}
and if we substitute this value in (12),
\[
-4\cdot 27\gamma = A^3 + 3Ab^2 + 3A^2 + 9b^2 - 4,
\]
and from this
\[
3b^2 \equiv 0\ (\mathrm{mod}\ 9).
\]

Thus $b$ is divisible by 3, and if we set $b = 3B$,
it follows from (15) that
\weq{16}{4n = A^2 + 27B^2.}

We can also easily form the square root of the dis\-criminant, namely
\weq{17}{\sqrt{D} = (\eta - \eta_1)(\eta - \eta_2)(\eta_1 - \eta_2) = s - s',}
if we subtract (10) and (11) from one another. From this one
obtains
\weq{18}{\sqrt{D} = nB.}

The cubic equation whose roots are the three quantities $\eta$
therefore has, as indeed already follows from what was remarked at the
beginning of \S.\,171, three real roots. It is simplified
if one sets
\[
3\eta + 1 = \xi
\]
and then yields
\weq{19}{\xi^3 - 3n\xi - nA = 0.}

If one substitutes on the left-hand side for $\xi$ in succession
\weq{20}{-2\sqrt{n}, \quad -\sqrt{n}, \quad +\sqrt{n}, \quad +2\sqrt{n},}
there result the values
\[
-n(2\sqrt{n} + A), \quad +n(2\sqrt{n} - A), \quad -n(2\sqrt{n} + A),
\]
\[
+n(2\sqrt{n} - A),
\]
which, because by (16) $A$ is absolutely smaller than $2\sqrt{n}$, have alternating
signs. Hence in each of the three intervals bounded by the
values (20) there lies \gesperrt{one} root of equation (19).


% ===== scan leaf 607, printed 582 =====
% !PRINTED 582
\weberpage{582}
On the question which of the three roots of the cubic
equation (19) is to be set for $\xi = 3\eta + 1$, which for $3\eta_1 + 1$ or for
$3\eta_2 + 1$, this much may be said in general: the sign of the
product (17) must agree with the sign of (18), which is determined
by the definition (1), (2), though it still depends on the choice
of the primitive root $g$. By this, three of the six possible
assignments are excluded. Which assignment among the three
remaining is to be made depends on the choice of $r$, and would
give occasion to investigations similar to those carried out in the
previous section for the determination of the sign of the square root\footnote{\gesperrt{Kummer}, Journ. f. Mathematik, Bd. 32.}.

We emphasise further the theorem proved by the formulae
(14), (16):

\gesperrt{If} $n$ \gesperrt{is a prime of the form} $3k + 1$, \gesperrt{then} $n$ \gesperrt{is}
\gesperrt{representable by the form} $a^2 - ab + b^2$ \gesperrt{and} $4n$ \gesperrt{by the form}
$A^2 + 27B^2$, \gesperrt{where} $a, b, A, B$ \gesperrt{are integers.}

From this it follows further that $n$ is also representable in
the form $x^2 + 3y^2$. For if of the two numbers $a, b$ one,
say $a$, is even, then
\[
n = \left(\frac{a}{2} - b\right)^2 + \tfrac{3}{4}a^2,
\]
and if $a$ and $b$ are both odd, then
\[
n = \tfrac{1}{4}(a + b)^2 + \tfrac{3}{4}(a - b)^3. % [sic] printed exponent 3, expected 2
\]

The two-term periods of the $9$th roots of unity also
satisfy a cubic equation, and because this equation plays an
important role in the general theory of the cubic cyclotomic
fields to be treated later, we shall consider it here for the
sake of completeness, although it properly belongs to a more
general region, in which the degrees of the roots of unity
are no longer primes.

A $9$th root of unity $r$ satisfies the equation of degree $6$ % [sic] the German prints "Einheitswuzel" for "Einheitswurzel"
\weq{21}{r^6 + r^3 + 1 = 0,}
and we have three conjugate two-term periods
\weq{22}{\eta = r + r^{-1}, \quad \eta' = r^2 + r^{-2}, \quad \eta'' = r^4 + r^{-4},}
which are the roots of the cubic equation
\weq{23}{\eta^3 - 3\eta + 1 = 0}
%CONT>>


% ===== scan leaf 608, printed 583 =====
% !PRINTED 583
\weberpage{583}
%<<CONT
are. If $\varrho$ is a third root of unity, we can set $\varrho = r^3$,
and obtain the resolvents
\weq{24}{(\varrho,\eta) = \eta + \varrho\eta' + \varrho^2\eta'', \quad (\varrho^{-1},\eta) = \eta + \varrho^{-1}\eta' + \varrho^{-2}\eta'',}
for which, by (22), making use of (21), one obtains
\[
(\varrho, r) = 3r, \quad (\varrho^{-1}, r) = 3r^{-1},
\]
hence
\weq{25}{(\varrho, r)^3 = 27\varrho,}
\weq{26}{(\varrho, r)(\varrho^{-1}, r) = 9,}
% [source-original mathematical notation, verified 2026-08-29] Immediately after defining
% the resolvents (varrho,eta) in (24), the print switches to (varrho,r) in the two displayed
% evaluations and in (25)-(26). Direct substitution from (21)-(24) gives (varrho,eta)=3r
% and (varrho^-1,eta)=3r^-1, so eta is the expected second argument; print preserved.
which is entirely analogous to the formulae (5) and (6).

We enter briefly also on the case $e = 4$, which occurs for
the primes $n$ of the form $4f + 1$ ($f$ an integer),
thus for $n = 5, 13, 17, 29, 37, 41 \dots$

Let the four periods of $f$ terms each again be $\eta, \eta_1, \eta_2, \eta_3$.

For $\alpha$ in \S.\,169 (8), we have to set $i$, and thus obtain,
since $-2 \equiv +2\ (\mathrm{mod}\ 4)$,
\weq{27}{\begin{aligned}
\psi_1(i) &= \sum i^{\operatorname{ind}t + 2\operatorname{ind}(t+1)}\\
\psi_2(i) &= \sum i^{\operatorname{ind}(t^2+t)} = a + bi,
\end{aligned}}
wherein $a$ and $b$ are integers. $\psi_1(i)$ has the same form
as $\psi_2(i)$. But we shall presently reduce the one function to the
other.

By the results of \S.\,171 we have
\weq{28}{(-1,\eta) = \eta + \eta_2 - \eta_1 - \eta_3 = \sqrt{n},}
and hence
\weq{29}{2(\eta + \eta_2) = -1 + \sqrt{n}, \quad 2(\eta_1 + \eta_3) = -1 - \sqrt{n},}
where the sign of $\sqrt{n}$ is to be determined according to \S.\,171 and
may, on a suitable assumption about $r$, be taken positive.

We have further, by \S.\,169, (14), (15)
\weq{30}{\begin{aligned}
(i,\eta)(-i,\eta) &= (-1)^f n\\
(i,\eta)^2 &= (-1,\eta)\psi_1(i),
\end{aligned}}
hence by (28)
\weq{31}{(i,\eta)^2 = \psi_1(i)\sqrt{n}, \quad (-i,\eta)^2 = \psi_1(-i)\sqrt{n}}
\weq{32}{(i,\eta)^4 = (-1)^f n\,\psi_1(i)\psi_2(i) = n\,\psi_1(i)^2,}
from which it follows
\weq{33}{\psi_1(i) = (-1)^f \psi_2(i) = (-1)^f(a + bi),}
and from this by (30)
\weq{34}{\psi_1(i)\psi_1(-i) = \psi_2(i)\psi_2(-i) = a^2 + b^2 = n.}


% ===== scan leaf 609, printed 584 =====
% !PRINTED 584
\weberpage{584}

To form the biquadratic equation whose roots are $\eta$,
$\eta_1$, $\eta_2$, $\eta_3$, we first seek the quadratic equation
with the roots $\eta$, $\eta_2$, and obtain it from

\[
\begin{aligned}
2(\eta+\eta_2)&=-1+\sqrt{n}\\
2(\eta-\eta_2)&=(i,\eta)+(-i,\eta).
\end{aligned}
\]

If we square these two equations and subtract the
second from the first, then by (30), (31) and (33) it follows that

\weq{35}{16\,\eta\,\eta_2=1+n-2\sqrt{n}-2n(-1)^f-2a(-1)^f\sqrt{n}.}

But by (29) and (35) the coefficients of the quadratic
equation sought are determined. It reads:
\[
\eta^2+\frac{\eta}{2}+\frac{1+n-2n(-1)^f}{16}=\left(\frac{\eta}{2}+\frac{a(-1)^f+1}{8}\right)\sqrt{n},
\]
and the biquadratic equation for the four periods is obtained
by squaring both sides:

\weq{36}{\left(\eta^2+\frac{\eta}{2}+\frac{1+n-2n(-1)^f}{16}\right)^2-n\left(\frac{\eta}{2}+\frac{a(-1)^f+1}{8}\right)^2=0.}

If we seek the coefficient of the first power of $\eta$, which
must be an integer, we find for it

\weq{37}{\frac{1+n-2n(-1)^f}{16}-n\,\frac{a(-1)^f+1}{8};}

$f$ is even when $n\equiv 1$, odd when $n\equiv 5$ (mod 8),
and accordingly
\[
1+n-2n(-1)^f\equiv 0\ (\mathrm{mod}\ 8).
\]
But from this it follows by (37) that $a(-1)^f+1$ must be
divisible by 4, hence

\weq{38}{a\equiv -(-1)^f\ (\mathrm{mod}\ 4),}

and from this it is further to be concluded, by (34), that $b$ must
be even, hence

\weq{39}{b\equiv 0\ (\mathrm{mod}\ 2).}

The biquadratic equation takes a simpler form under the substitution

\weq{40}{4\eta+1=\xi}

and by (36) becomes:

\weq{41}{[\xi^2+n(1-2(-1)^f)]^2-4n[\xi+(-1)^f a]^2=0,}

or, if we separate the two cases $n\equiv 1$ or $\equiv 5$ (mod 8):

\weq{42}{\begin{aligned}(\xi^2-n)^2-4n(\xi+a)^2&=0\\(\xi^2+3n)^2-4n(\xi-a)^2&=0.\end{aligned}}


% ===== scan leaf 610, printed 585 =====
% !PRINTED 585
\weberpage{585}
\pdfbookmark[2]{§ 173. The complex numbers of Gauss}{bm193}

We wish here also to emphasise the theorem expressed in
formula (34):

\gesperrt{Every prime of the form} $4f+1$ \gesperrt{can be decomposed into
the sum of two squares.}

\begin{center}
\S.~173.\\[6pt]
\gesperrt{The complex numbers of Gauss.}
\end{center}

The theorem last proved forms the foundation for the
theory of the number field which arises from the field of rational
numbers by adjunction of the imaginary
unit $i=\sqrt{-1}$, and which by our stipulation we have to denote
by $R(i)$\textsuperscript{1}).

We have here to distinguish the primes of the form $4f+1$
from those of the form $4f+3$, and for the sake of brevity we
wish to combine the former with the prime 2 and denote them
by $p$, the latter by $q$; thus we set
\[
p=4f+1,2;\quad q=4f+3
\]
The following theorems then hold:

1. \gesperrt{For every $p$ there are two integers $a$ and $b$
such that}
\[
p=a^2+b^2
\]
\gesperrt{holds.}

For the primes of the form $4f+1$ this has been proved in the
concluding theorem of the last section, and for $p=2$ it is
immediately evident from $2=1^2+1^2$.

Against this stands a second theorem:

2. \gesperrt{A prime $q$ is never representable in the form
$a^2+b^2$,} or the still more general one:

3. \gesperrt{The sum of two squares of integers $a^2+b^2$
is divisible by a prime $q$ only when $a$ and $b$
are divisible by $q$.}

Theorem 2. follows simply from the fact that, when $a^2+b^2$
is odd, one of the two numbers $a$, $b$ is even, and hence its
square divisible by 4, while the other is odd, and hence its square
%CONT>>

\textsuperscript{1}) \gesperrt{Gauss}, Theoria residuorum biquadraticorum, commentatio
secunda. Werke, Bd.~II.


% ===== scan leaf 611, printed 586 =====
% !PRINTED 586
\weberpage{586}

%<<CONT
$\equiv 1$ (mod 4), and hence $a^2+b^2\equiv 1$ (mod 4). Theorem 3.,
which moreover includes the second, is proved as follows.

Suppose that $a^2+b^2$, but not $b$, is divisible by $q$;
we then determine $b'$ from the congruence $bb'\equiv 1$ (mod $q$) and
obtain from $a^2+b^2\equiv 0$ (mod $q$)
\[
(ab')^2\equiv -1\ (\mathrm{mod}\ q).
\]
But this is impossible, since by \S.\,138, 4. the number $-1$ is a
quadratic non-residue for every prime $q$.

The field $R(i)$ is the totality of all numbers of the
form $x+yi$, where $x$, $y$ are integral or fractional rational
numbers.

\gesperrt{The numbers} $a+bi$, \gesperrt{in which} $a$ \gesperrt{and} $b$ \gesperrt{are integers
in} $R$, \gesperrt{are called the integers of the field} $R(i)$.

The product of two conjugate numbers
\[
\xi=x+yi,\quad \xi'=x-yi,
\]
that is
\[
\xi\xi'=x^2+y^2
\]
is called the norm of $\xi$ and is denoted by $N(\xi)$. The norm
of a non-vanishing number $\xi$ is a positive rational
number, and the norm of an integer is an integer. The
norm of a product or of a quotient is equal to the product
or the quotient of the norms.

\gesperrt{An integer whose norm is equal to 1 is called
a unit.}

Since $a^2+b^2$ (for integers $a$, $b$) can be $=1$ only
when $a=\pm 1$, $b=0$ or $a=0$, $b=\pm 1$, there are in
$R(i)$ only the four units
\[
+1,\quad -1,\quad +i,\quad -i.
\]
The reciprocal value of a unit is also a unit.

Two numbers, of which one arises from the other by
multiplication with a unit, are called \gesperrt{associated
numbers}.

Every complex number belongs to a system of four associated
numbers
\[
a+bi,\quad -a-bi,\quad -b+ai,\quad b-ai.
\]
The sum, difference and product of two integers are
again integers.

\gesperrt{An integer} $\alpha$ \gesperrt{is called divisible by an integer} $\beta$
\gesperrt{if a third integer} $\gamma$ \gesperrt{exists, such
that} $\alpha=\beta\gamma$.


% ===== scan leaf 612, printed 587 =====
% !PRINTED 587
\weberpage{587}

If $\alpha$ is divisible by $\beta$, then $N(\alpha)$ is divisible by $N(\beta)$.
For from $\alpha=\beta\gamma$ it follows that $N(\alpha)=N(\beta)N(\gamma)$.

If $\alpha$, $\beta$, $\gamma$ are integers in $R(i)$ and $\alpha$ is divisible by $\gamma$,
then $\alpha\beta$ is also divisible by $\gamma\beta$; and if $\alpha$ and $\beta$ are divisible by $\gamma$,
then $\alpha\pm\beta$ is also divisible by $\gamma$.

Every integer is divisible by every unit.

% [source-original mathematical ellipsis, verified 2026-08-29] From
% $N(\alpha)N(\beta)=1$ with positive integral norms, both norms equal $1$;
% the print omits ``$=1$'' after the first $N(\alpha)$.
The units, however, are divisible only by units; for
if $\alpha\beta$ is a unit, then $N(\alpha)N(\beta)=1$, hence $N(\alpha)$ and
$N(\beta)=1$, i.\,e. $\alpha$ and $\beta$ are units.

Associated numbers are mutually divisible by one another,
and when two numbers are mutually divisible, they are
associated. For if $\alpha:\beta$ and $\beta:\alpha$ are both integers, then
the product of these two numbers is $=1$; hence both are units.

For the numbers of the field $R(i)$ the same
laws hold concerning the decomposition into prime factors as for the real
integers.

\gesperrt{An integer $\alpha$ of the field $R(i)$ which is not a unit
is called composite if it can be decomposed into several
integer factors, none of which is a unit.}

\gesperrt{If it cannot be so decomposed, it shall be called a
prime in the field $R(i)$.}

If $\xi=x+yi$ is a fractional number of the field $R(i)$, then
an integer $\mu=m+ni$ can be determined so that the
norm of the difference $\xi-\mu$, that is $(x-m)^2+(y-n)^2$, is smaller
than, or at least not greater than, $1/2$; for one need only
choose the integral rational numbers $m$, $n$ so that $x-m$
and $y-n$, taken absolutely, are not greater than $1/2$, and hence their squares
not greater than $1/4$.

If, then, $\alpha$ and $\alpha_1$ are two integers in $R(i)$, and $\alpha_1$ is
different from zero, one can accordingly determine the integers $\mu$
and $\alpha_2$, so that
\[
\frac{\alpha}{\alpha_1}-\mu=\frac{\alpha_2}{\alpha_1},\quad \frac{N(\alpha_2)}{N(\alpha_1)}\leqq \frac{1}{2},
\]
hence
\[
\alpha=\mu\alpha_1+\alpha_2,\quad N(\alpha_2)<N(\alpha_1).
\]
If $\alpha_2$ is not zero, one can proceed in the same way with the numbers
$\alpha_1$, $\alpha_2$ and obtains
\[
\alpha_1=\mu_1\alpha_2+\alpha_3,\quad N(\alpha_3)<N(\alpha_2),
\]
%CONT>>


% ===== scan leaf 613, printed 588 =====
% !PRINTED 588
\weberpage{588}
%<<CONT
and in this way one determines a sequence of integers $\alpha, \alpha_1, \alpha_2, \alpha_3 \dots$
with steadily decreasing norm, which can be continued so long
as zero does not occur among them. Since these norms are positive
rational integers which always decrease, zero must appear
after a finite number of steps, and one thus obtains
a system of equations

\weq{1}{\begin{aligned}
\alpha &= \mu\alpha_1 + \alpha_2,\\
\alpha_1 &= \mu_1\alpha_2 + \alpha_3,\\
&\cdots\cdots\cdots\cdots\\
\alpha_{h-2} &= \mu_{h-2}\alpha_{h-1} + \alpha_h,\\
\alpha_{h-1} &= \mu_{h-1}\alpha_h.
\end{aligned}}

From this one concludes that the whole sequence of numbers
$\alpha, \alpha_1, \alpha_2 \dots$, and thus in particular also $\alpha$ and $\alpha_1$, are divisible
by $\alpha_h$, and conversely, that every common divisor of $\alpha$
and $\alpha_1$ is a divisor of all the following $\alpha$, and hence also of $\alpha_h$.
Every other number that has these two properties must
be associated with $\alpha_h$, and one therefore calls $\alpha_h$ (and every number
associated with $\alpha_h$) the \gesperrt{greatest common divisor} of
$\alpha$ and $\alpha_1$.

If from the algorithm (1) we eliminate the intermediate $\alpha$,
by substituting $\alpha_2$ from the first equation into the second, then $\alpha_3$
from the second into the third, there results
a result of the form $\alpha\varkappa + \alpha_1\lambda = \alpha_h$, and we can
therefore, if instead of $\alpha, \alpha_1, \alpha_h$ we write $\alpha, \beta, \delta$, state the theorem:

4. \gesperrt{If $\alpha, \beta$ are any two integers in $R(i)$,
and $\delta$ their greatest common divisor,
then the integers $\varkappa, \lambda$
can be so determined that}

\weq{2}{\alpha\varkappa + \beta\lambda = \delta}
\gesperrt{holds.}

If the greatest common divisor $\delta$ of two numbers
$\alpha, \beta$ is a unit, then $\alpha$ and $\beta$ are called relatively prime, and in this
case we can satisfy the equation

\weq{3}{\alpha\varkappa + \beta\lambda = 1}
by integral $\varkappa, \lambda$.

If $\beta$ is a prime, then either $\alpha$ is divisible by $\beta$,
or $\alpha$ and $\beta$ are relatively prime. In the latter case
equation (3) holds. If then $\gamma$ is another number in $R(i)$, it follows from (3) that
\[
\alpha\gamma\varkappa + \beta\gamma\lambda = \gamma.
\]


% ===== scan leaf 614, printed 589 =====
% !PRINTED 589
\weberpage{589}
Accordingly, if $\alpha\gamma$ is divisible by $\beta$, then $\gamma$ too must be
divisible by $\beta$. We therefore have the following fundamental theorem:

5. \gesperrt{A product of two or more integers
in $R(i)$ is divisible by a prime in $R(i)$
only if at least one of its factors
is divisible by this prime.}

If one now further makes the remark that a prime can be
divisible by another only if the quotient
is a unit, that is, if the two are associated, and regards
associated primes as not essentially different, it follows:

6. \gesperrt{An integer of the field $R(i)$ can always
and in only one way be decomposed into a product
 of primes.}

For let $\alpha$ be any integer in $R(i)$. If $\alpha$ is not
itself a prime, it is decomposable, say into $\gamma\beta$; but the norms
of $\gamma$ and of $\beta$ are smaller than the norm of $\alpha$. If $\beta$
is not yet a prime, it is again decomposable, say into $\delta\varepsilon$; the
norms of these factors, which are all positive integers,
decrease continually, and one must therefore necessarily at some point come upon
a factor which is a prime. Hence every number $\alpha$
is certainly divisible by a prime $\pi$. If accordingly $\alpha = \pi\alpha'$,
then the same holds of $\alpha'$, and we have, say, $\alpha' = \pi'\alpha'' \dots$ Since the
norms of $\alpha, \alpha', \alpha'' \dots$ again all decrease, one must,
in continuing this sequence, finally come upon a unit.
Hence $\alpha$ is decomposable into a finite number of prime factors,
thus
\[
\alpha = \pi\pi'\pi'' \dots,
\]
where a unit factor remaining over at the end can be united with one of the
primes $\pi$. If now $\varkappa, \varkappa', \varkappa'' \dots$
are likewise primes and $\alpha = \varkappa\varkappa'\varkappa'' \dots$, then it follows from
\[
\varkappa\varkappa'\varkappa'' \dots = \pi\pi'\pi'' \dots
\]
by 5.\ that one of the primes $\pi$, e.g.\ the first, is divisible by $\varkappa$,
and hence is not essentially different from $\varkappa$. If one cancels
on both sides by $\varkappa = \pi$, one can repeat the same inference with $\varkappa'$,
and so forth, and thus finds that each of the primes $\varkappa$
must also occur among the $\pi$, and if among the $\varkappa$ a
prime occurs several times, it must occur at least as
often among the $\pi$. But the same inference can also be
made in the reverse direction, and from this it follows that the totality
%CONT>>


% ===== scan leaf 615, printed 590 =====
% !PRINTED 590
\weberpage{590}
%<<CONT
the $\varkappa$, apart from unit factors, must agree with the totality
of the $\pi$.

It now remains to actually determine the primes of the
field $R(i)$.

Every prime is a divisor of infinitely many rational
integers, e.g. of the norm and all its multiples.
Among the rational positive integers that are divisible by a prime
$\pi$ there is one, which we shall denote by $n$,
that is the smallest, and this must be a prime in the field $R$.
For if it were decomposable in $R$, then by 5. one of its factors,
and hence a still smaller number, would have to be divisible by $\pi$. Likewise,
conversely, every real prime $n$ is also divisible by at least one
prime $\pi$. Let, then, $n = \pi\alpha$, from which, by forming
the norm, $n^2 = N(\pi)N(\alpha)$ follows. Since $n$ is a prime and
$N(\pi)N(\alpha)$ integers, two cases are possible:
\[
\begin{aligned}
1.\quad & N(\pi) = n, \quad N(\alpha) = n,\\
2.\quad & N(\pi) = n^2, \quad N(\alpha) = 1.
\end{aligned}
\]

In the first case, if $\pi'$ denotes the number conjugate to $\pi$
and $\pi = a+bi$ is set, then $n = \pi\pi' = a^2+b^2$,
$\alpha = \pi'$, and one sees that this case occurs only when
$n$ belongs to the primes $p$ that are decomposable into the sum of two
squares. Conversely, every such number $p$
can be decomposed into two conjugate factors $\pi\pi'$, of which neither
is a unit. \gesperrt{The primes} $p$ \gesperrt{are therefore not primes in the field}
$R(i)$.

In the second case, which accordingly occurs with the primes $q$,
$\alpha$ is a unit, hence $n$ is associated with $\pi$; that is, the \gesperrt{real
primes} $q$ \gesperrt{are also primes in the field} $R(i)$.

The prime 2 belongs, as already remarked, to the first
kind, and $2 = (1+i)(1-i)$. But it occupies a special
position, because $1-i = -i(1+i)$, and hence $2 = -i(1+i)^2$.
The real prime 2 is therefore (apart from the factor $-i$)
\gesperrt{the square of a prime} in the field $R(i)$. This
case, as one easily sees, occurs with none of the remaining
primes $p$.

% [sic] German authority prints "des Körper R(i)" for grammatical "des Körpers R(i)"; diplomatic DE preserved.
The totality of the primes of the field $R(i)$ thus consists
of the real primes $q$ and of the factors of the primes
$p$. If $p$ is decomposed into the sum of two squares, $a^2+b^2$, then
one also knows the conjugate factors $a\pm bi$ of $p$. To this
%CONT>>


% ===== scan leaf 616, printed 591 =====
% !PRINTED 591
\weberpage{591}
%<<CONT
come also the associated numbers. From 5. it then further follows
that a prime $p$ can be decomposed in only \gesperrt{one way} into the sum of two
squares. By this decomposition, however, the numbers $a, b$ are
not yet completely determined; rather, they can still be interchanged
with one another and furnished with two signs.
This comes down to saying that one can replace $\pi$ by
any of the four associated numbers $a+bi, -a-bi, -b+ai,$
$b-ai$. If $p$ is odd, then of the two
numbers $a, b$ one must be even, the other odd. If we choose,
say, for $b$ the even one of the two numbers, and determine
the sign so that $a \equiv 1 \pmod{4}$, then among the
four associated numbers a definite one is selected, which one may
\gesperrt{call the primary one}\footnote{Gauss gives at the place mentioned two different determinations
 of the primary numbers to choose from, of which this is the first.
He afterwards retains the second.}.

This definition of the primary numbers can be carried over to all
integers of the field $R(i)$ whose norm is
odd, and one then has the law that the product
of two primary numbers again yields a primary number.

The system of four associated numbers in the field $R(i)$ is
analogous to the pair of opposite numbers in the field $R$. In $R$
one regards the positive numbers as the primary ones.

By the formula (27) of the preceding section, for
any given prime $p$ one of the factors $a+bi$ is derived
from the division of the circle. By the formulae \S.\,172, (38), (39)
this number is primary if $p \equiv 5 \pmod{8}$, but opposite to
the primary one if $p \equiv 1 \pmod{8}$. But as to
which of the two conjugate numbers $a \pm bi$ is represented by these
formulae, we have no general criterion.

For example, the complex primes in $R(i)$ whose
norms are less than 200 are:
\[
\begin{aligned}
&1+i,\ 1+2i,\ 3+2i,\ 1+4i,\ 5+2i,\ 1+6i,\ 5+4i,\\
&7+2i,\ 5+6i,\ 3+8i,\ 5+8i,\ 9+4i,\ 1+10i,\ 3+10i,\\
&7+8i,\ 11+4i,\ 7+10i,\ 11+6i,\ 13+2i,\ 9+10i,\ 7+12i,\\
&1+14i.
\end{aligned}
\]


% ===== scan leaf 617, printed 592 =====
% !PRINTED 592
\weberpage{592}
\pdfbookmark[2]{§ 174. The field of the third roots of unity}{bm194}

\begin{center}
\S.\,174.

\gesperrt{The Field of the Third Roots of Unity.}
\end{center}

The principal theorem of the preceding section, that every integer of the field $R(i)$ can be decomposed in only one way into indecomposable factors, and that in consequence the concept of the indecomposable number coincides with that of the prime, rests, as one sees, in the main on the algorithm of the greatest common divisor contained in the formulae \S.\,173, (1); and thus if in another field one has such an algorithm, terminating after finitely many steps, the same conclusions can be drawn.

This is the case for the field which can be derived from $R$ by means of the third root of unity $\varrho = \dfrac{-1+\sqrt{-3}}{2}$, and which may be denoted by $R(\varrho)$ or also by $R(\sqrt{-3})$. This field consists of the totality of all numbers of the form $\xi = x+\varrho y$, in which $x$, $y$ are rational numbers, and the norm of such a number is
\[
\begin{aligned}
N(\xi) &= (x+\varrho y)(x+\varrho^2 y) = x^2 - xy + y^2\\
&= \frac{2x-y+y\sqrt{-3}}{2}\cdot\frac{2x-y-y\sqrt{-3}}{2} = \frac{(2x-y)^2+3y^2}{4}.
\end{aligned}
\]
A number $\alpha = a+b\varrho$ is called an integer in $R(\varrho)$ if $a$, $b$ are rational integers. The units in this field result from
\[
N(\alpha) = a^2 - ab + b^2 = 1,
\]
or
\[
(2a-b)^2 + 3b^2 = 4,
\]
which is satisfied only for $b=0$, $a=\pm 1$ or $b=\pm 1$, $a=\pm 1$, so that here one has six units
% [source-original mathematical omission/ambiguity, verified 2026-08-29] Solving
% a^2-ab+b^2=1 gives (a,b)=(+-1,0),(0,+-1),(1,1),(-1,-1). The printed
% condition omits the a=0,b=+-1 cases and leaves the signs in its second pair ambiguous;
% the following six-unit list itself is correct and the printed meaning is preserved.
\[
\pm 1,\quad \pm\varrho,\quad \pm\varrho^2.
\]
A system of associated numbers is here
\[
\begin{aligned}
&\pm(a+b\varrho),\quad \pm(a\varrho+b\varrho^2) = \pm[-b+(a-b)\varrho]\\
&\pm(a\varrho^2+b) = \pm(b-a-a\varrho).
\end{aligned}
\]
If $\xi$ is an arbitrary fractional number in $R(\varrho)$, one can determine the integer $\mu$ so that in $\xi-\mu = x+\varrho y$ the
%CONT>>


% ===== scan leaf 618, printed 593 =====
% !PRINTED 593
\weberpage{593}
%<<CONT
components $x$ and $y$ are absolutely not greater than $1/2$, and hence that
\[
N(\xi-\mu) = x^2 - xy + y^2 \leqq \frac{3}{4}.
\]
On the basis of this property one can, just as in the preceding section, determine the greatest common divisor of two numbers $\alpha$, $\beta$, and thereby prove that every integer in $R(\varrho)$ can be decomposed into primes of this field in only one way, if associated numbers are not regarded as essentially different.

As in the preceding section one concludes that one obtains all the primes of the field $R(\varrho)$ by decomposing all real primes, and that a real prime either splits into two conjugate prime factors in $R(\varrho)$, $p=\pi\pi'$, or is a prime in the field $R(\varrho)$ as well.

The prime 3 splits into the two factors $(1-\varrho)$, $(1-\varrho^2)$; but these are associated, and so $3=-\varrho^2(1-\varrho)^2$ is associated with the square of a prime in $R(\varrho)$. One may choose $\sqrt{-3}$ for this prime. All primes of the form $3f+1$ can, by \S.\,172, be decomposed into two conjugate factors in $R(\varrho)$, while the primes of the form $3f+2$ are never decomposable, because an expression of the form $a^2-ab+b^2$ cannot be $\equiv 2\ (\mathrm{mod}\ 3)$. Thus one obtains all the primes in $R(\varrho)$ by seeking out the real primes of the form $3f+2$ and the conjugate factors of the remaining real primes.

If $a+b\varrho$ is a complex integer whose norm $a^2-ab+b^2$ is not divisible by 3, then of the two numbers $a$, $b$ either one must be divisible by 3 and the other not, or we must have $a\equiv b$ and consequently $a^2-ab+b^2\equiv 1\ (\mathrm{mod}\ 3)$. Among the associated numbers
\[
a+b\varrho,\quad b+a\varrho^2,\quad b-(a-b)\varrho,
\]
there is therefore always one and only one in which the coefficient of $\varrho$ (or of $\varrho^2$) is divisible by 3. Let us assume that $b=3B$ is divisible by 3, and set $N(a+b\varrho)=p$; then it follows, if in addition $2a-b=A$ is set,
\[
4p = A^2 + 27B^2,
\]
in agreement with \S.\,172. If $p$ is a real prime, and $4p$ has been represented in this form, one obtains from it the complex factors of $p$
%CONT>>


% ===== scan leaf 619, printed 594 =====
% !PRINTED 594
\weberpage{594}
%<<CONT
% [page alignment mirrored 2026-08-07] The formula/list tail below is direct
% print on page 594.  The former section opener belongs to the following
% unnumbered leaf, i0620; page 594 ends after the centered closing rule.
\[
p = \left(\frac{A+3B}{2}+3\varrho B\right)\left(\frac{A-3B}{2}-3\varrho B\right)\!.
\]
The complex primes whose norm is under 200 are here $1-\varrho$, $1+3\varrho$, $4+3\varrho$, $5+3\varrho$, $5+6\varrho$, $7+3\varrho$, $7+6\varrho$, $5+9\varrho$, $7+9\varrho$, $8+9\varrho$, $10+3\varrho$, $11+3\varrho$, $11+9\varrho$, $7+12\varrho$, $13+6\varrho$, $13+3\varrho$, $14+9\varrho$, $13+12\varrho$, $14+3\varrho$, $11+15\varrho$, $16+9\varrho$, $13+15\varrho$.

\begin{center}
\rule{3cm}{0.4pt}
\end{center}


% ===== scan leaf 620, printed 595 =====
\weberpage{595}% page number not printed on this page
\pdfbookmark[1]{Seventeenth Section. Algebraic solution of equations.}{bm195}
\pdfbookmark[2]{§ 175. Reduction of the groups by pure equations}{bm196}
% !PRINTED none
% Abschnitt-opening (part-title) page: no running-head folio visible on this
% scan. Carries the Abschnitt heading, §.175 heading, and body text.

\begin{center}
\gesperrt{Seventeenth Section.}

\gesperrt{Algebraic Solution of Equations.}

\rule{3cm}{0.4pt}

\S.\,175.

\gesperrt{Reduction of the Group by Pure Equations.}
\end{center}

One of the oldest questions, in connection with which the more
recent algebra has chiefly developed, is that of the so-called
\gesperrt{algebraic solution} of equations, by which one understands a
representation of the roots of an equation by a series of
radicals, or their computation by a finite chain of root
extractions. On this question the theory of groups sheds the
brightest light.

If we first make precise the question at issue, it is evidently
this: whether and how one can extend the field $\Omega$ by
successive adjunction of root quantities (radicals) in such a way
that either all or at least a part of the roots
are contained in the extended field. A root quantity, however,
is one which is not itself contained in $\Omega$, but of which some
integral power is; that is, if $a$ is a quantity in $\Omega$,
it is the root of an equation of the form

\[
y^m - a = 0,
\]

that is, of a \gesperrt{pure equation}.

If an \gesperrt{irreducible} equation is to be algebraically solvable,
or at least to have one or several algebraically representable roots,
then after a successive adjunction of roots of
pure equations in finite number the given equation
must become reducible, since a part of its roots is to be contained
in the extended field. Since the initial group $P$ is transitive,
this group must finally become intransitive, or
%CONT>>


% ===== scan leaf 621, printed 596 =====
% !PRINTED 596
\weberpage{596}

%<<CONT
reduce to the unit group. Hence the case must at all events
occur at some point that the group $P$ is reduced by the adjunction
of a root of a pure equation.

The investigation of this question is extraordinarily simplified
if one transforms it somewhat further\,\footnote{This form of the question was first pointed out by C. \gesperrt{Jordan} (Trait\'e des substitutions p. 386).}.

We have seen in \S.\,162 that the pure equations belong to the
\gesperrt{Abel}ian ones, that all \gesperrt{Abel}ian equations are solvable through a
chain of cyclic equations of prime degree, and these latter
by radicals.

We therefore replace the question of solvability by radicals
by the equivalent one of solvability by roots of \gesperrt{Abel}ian
equations. If we replace the \gesperrt{Abel}ian equation by a chain of
cyclic equations of prime degree, then the first reduction of the
group will occur through the adjunction of the root of one such
equation, and so we stand first of all before the question:

\gesperrt{Under what conditions is the group $P$
of an equation of degree $n$, $f(x) = 0$, reduced by the adjunction
of a root of a cyclic equation of prime degree $m$?}

We do not restrict ourselves here to irreducible equations
$f(x) = 0$, but discuss the question in general, always under the
self-evident assumption that $f(x)$ has no multiple
roots.

It is conceivable that it may be necessary first to adjoin to
the field $\Omega$ various roots of cyclic equations (e.\,g. roots of unity)
which do not alter the group $P$. But since we are now
concerned only with ascertaining the necessary properties of the
group $P$, we assume that all the necessary preparations
have been made, and the field $\Omega$ so constituted
that by the adjunction of a root of the cyclic equation
$\varphi(x) = 0$, of prime degree $m$, the group $P$ reduces to
one of its divisors, $Q$.

If we denote the roots of $\varphi(x) = 0$ by

\[
\varepsilon,\ \varepsilon_1,\ \varepsilon_2 \ldots \varepsilon_{m-1},
\]
then all these quantities are expressible rationally (in $\Omega$)
through an arbitrary one among them, and if $P$ is the group of
$f(x) = 0$
%CONT>>


% ===== scan leaf 622, printed 597 =====
% !PRINTED 597
\weberpage{597}
\pdfbookmark[2]{§ 176. Metacyclic equations}{bm197}
%<<CONT
in $\Omega$, then $Q$ is the group of the same equation in $\Omega(\varepsilon)$, or,
what is the same thing, in $\Omega(\varepsilon_1), \Omega(\varepsilon_2) \ldots \Omega(\varepsilon_{m-1})$. But now we can
apply the concluding theorem in \S.\,157. By this theorem the
index $j$ of the divisor $Q$ of $P$ must be a divisor of $m$, and
since $m$ is assumed to be a prime, we have $m = j$. Moreover, by
the same theorem, $\varepsilon$ is rationally representable by the roots of the
equation $f(x) = 0$:

\[
\varepsilon = \psi(x_0, x_1 \ldots x_{m-1}).
\]

This function belongs to the group $Q$, and if we apply to it
all the permutations of the group $P$, we obtain
the functions $\varepsilon, \varepsilon_1, \varepsilon_2 \ldots \varepsilon_{m-1}$ and no others. These
functions belong to the conjugate groups $\pi^{-1}Q\pi$. But since
each of these functions is rationally expressible by every other,
they must all belong to the same group,
that is, $Q$ is a normal divisor of $P$.

We have thus proved the first theorem:

\begin{quote}
\noindent I.\ \gesperrt{If the group $P$ of an equation is reduced by
the adjunction of the roots of an Abelian equation, then $P$ has a normal divisor $Q$ of prime index.}
\end{quote}

This theorem can also be reversed.

For if the group $P$ has a normal divisor $Q$ of
index $m$, then we can choose a function $\psi$ belonging to $Q$,
and the functions conjugate with it, $\psi, \psi_1, \psi_2 \ldots \psi_{m-1}$,
all belong to the same group. The field $\Omega(\psi)$ is a
normal field and $\psi$ the root of a normal equation. In the
field $\Omega(\psi)$, $Q$ is the group of $f(x) = 0$ (\S.\,155). But if
$m$ is a prime, then by \S.\,163 $\psi$ is the root of a
cyclic equation, and thus it is proved:

\begin{quote}
\noindent II.\ \gesperrt{If the group $P$ of $f(x) = 0$ possesses a normal divisor $Q$ of prime index $m$, then by
the adjunction of the root of a cyclic
equation of degree $m$ the group $P$ is
reduced to $Q$.}
\end{quote}

\begin{center}
\S.\,176.

\gesperrt{Metacyclic Equations.}
\end{center}

We shall call an equation whose complete solution can be
reduced to a chain of cyclic equations
%CONT>>


% ===== scan leaf 623, printed 598 =====
% !PRINTED 598
\weberpage{598}
%<<CONT
call a \gesperrt{metacyclic equation}. The cyclic equations
themselves are contained among them as a special case,
and by what was remarked in \S.\,175 the metacyclic
equations are the same as the equations solvable by radicals. If $P$ is the group of such an equation, then by
the preceding section it must possess a normal divisor of prime index $j_1$, which we now wish to denote by $P_1$.
If $P_1$ consists of the single identical permutation, then $P$
is itself cyclic and of prime degree. If $P_1$ is not the unit group, then $P_1$ must in turn contain a normal divisor $P_2$ of prime index $j_2$, and so on, until we finally arrive at the unit
group.

That this composition of its group $P$ is also the sufficient condition for a metacyclic equation follows from the preceding section. We
therefore pronounce the theorem:

\begin{quote}
\noindent III.\ \gesperrt{The necessary and sufficient condition for
a metacyclic equation is that there is
a sequence of groups}

\[
P, P_1, P_2, P_3 \ldots
\]

\gesperrt{whose first is the Galois group of the
equation, whose last is the unit group,
and of which each following one is a normal divisor
of the one next preceding, of prime index.}
\end{quote}

Accordingly we call a permutation group $P$ which has this
property, for which therefore a chain of groups

\[
P, P_1, P_2 \ldots P_{\mu-1}, 1
\]

can be so determined that each term is a normal divisor of the preceding, of prime index, a \gesperrt{metacyclic
group}\footnote{The expression ``metacyclic groups'' was first used by \gesperrt{Kronecker},
though in a more restricted sense. I should like here to propose this
slight generalisation of an already known expression in place of the
expression ``solvable groups'' used by \gesperrt{Frobenius} and \gesperrt{Hölder}.}.

We have here obtained the condition for the \gesperrt{complete} solvability of an equation by a chain of cyclic
equations. But it is still a question of the further inquiry,
%CONT>>


% ===== scan leaf 624, printed 599 =====
% !PRINTED 599
\weberpage{599}
%<<CONT
whether one or several of the roots can be represented in this way,
while others admit no such representation. This question is justified
only for irreducible equations, since for reducible equations all
conceivable combinations can occur; and here the following theorem now holds:

\begin{quote}
\noindent IV.\ \gesperrt{If a root of an irreducible equation is
determinable by the solution of cyclic equations, then the equation is metacyclic.}
\end{quote}

If an irreducible equation of degree $n$, $f(x) = 0$, has even
one root which becomes rational through successive adjunction of roots
of cyclic equations, then it must necessarily become reducible
through this adjunction, since in the end a linear
factor must separate off. Let $P$ then be the group
of our equation, after everything necessary has been adjoined so far
that $f(x)$ is not yet reducible, but splits into factors through the next
adjunction of the root $\varepsilon$ of a cyclic equation of prime
degree $m$. Then by \S.\,175, $P$ must have
a normal divisor $Q$ of index $m$, to which the
group of the equation reduces after the adjunction of $\varepsilon$, and the
permutation group $Q$ must be intransitive. If $Q$ is the unit
group, then $f(x) = 0$ is completely solved by the adjunction of $\varepsilon$. But if $Q$ is still different from the unit group, then the theorems of \S.\,158 come into force. According to these,
$f(x)$ splits after the adjunction of $\varepsilon$ — since the index of $Q$ is here a
prime — into $m$ factors, $\varphi(x, \varepsilon), \varphi_1(x, \varepsilon), \ldots \varphi_{m-1}(x, \varepsilon)$, of
the same degree $\mu$, and we have $\mu m = n$.

If, for example, $n$ is a prime, then $\mu = 1$ must hold; the functions $\varphi$ are linear, and the equation $f(x) = 0$ is completely solved; hence Theorem IV is correct for equations of prime degree. In the general case, one of the factors
of degree $\mu$, $\varphi(x, \varepsilon), \varphi_1(x, \varepsilon), \ldots \varphi_{m-1}(x, \varepsilon)$, say $\varphi(x, \varepsilon)$, must have a
root which becomes rational through the adjunction of the roots of cyclic
equations; and if we therefore assume that
our theorem is already proved for equations of degree $\mu$, it
follows that $\varphi(x, \varepsilon) = 0$ itself, and hence also its group, is metacyclic.

But since by \S.\,158 the different equations $\varphi = 0$,
$\varphi_1 = 0, \ldots \varphi_{m-1} = 0$ have the same group, they are all,
and consequently also $f(x) = 0$, metacyclic.


% ===== scan leaf 625, printed 600 =====
% !PRINTED 600
\weberpage{600}
\pdfbookmark[2]{§ 177. Simplicity of the alternating group}{bm198}

Under the assumption, then, that Theorem IV is correct for equations of degree $\mu$, its correctness follows for equations of degree $\mu m$; and since it holds for equations of prime degree, it is proved in general.

\begin{center}
\S.\,177.

\gesperrt{Simplicity of the alternating group.}
\end{center}

We have seen earlier (\S.\,149) that if we regard the coefficients of an equation of degree $n$ as independent variables, and the field of all rational functions of these coefficients as the domain of rationality, the \gesperrt{Galois} group of the equation is the symmetric group. In the symmetric group there is always contained a normal divisor of index 2, the \gesperrt{alternating group}, to which the group of the equation reduces when the square root of the discriminant is adjoined. For four figures the alternating group, which consists of the identical permutation, eight three-term cycles, and three pairs of transpositions, has the normal divisor of index 3:

\[
1, \quad (0,\,1)(2,\,3), \quad (0,\,2)(1,\,3), \quad (0,\,3)(1,\,2).
\]

This group has three different normal divisors of index 2, of which we prefer one, 1, $(0,\,1)(2,\,3)$, of which in turn the unit group is a normal divisor of index 2. The group of the 24 permutations of four figures is therefore metacyclic, and on this rests every method of solution of the biquadratic equation (\S.\,160, 161).

We now wish to prove that if $n$ is greater than 4, the alternating group has no normal divisors at all besides the unit group, or, in the designation introduced earlier, is \gesperrt{simple}. From this it then follows that the condition which we found to be necessary for the algebraic solvability of an equation is not fulfilled for equations of higher than the fourth degree whose group is the symmetric or the alternating group, and that therefore \gesperrt{equations of higher than the fourth degree, so long as the coefficients are independent variables, are no longer algebraically solvable.}


% ===== scan leaf 626, printed 601 =====
% !PRINTED 601
\weberpage{601}

The proof that the alternating group is simple may be conducted as follows:

Let $A$ be the alternating group of the permutations of $n$ figures $0, 1, 2 \ldots n-1$, and $Q$ a normal divisor of $A$. We have seen in \S.\,153, 6. and \S.\,154, 6. that all permutations of $A$ can be composed of three-term cycles, and that, if $\varkappa$ and $\pi$ are any permutations whatever, one obtains the transformed permutation of $\varkappa$, $\pi^{-1}\varkappa\pi$, by carrying out in the cycles of $\varkappa$ the permutations $\pi$. Now if $\varkappa$ is a three-term cycle, say $(0,\,1,\,2)$, then $\pi$ can be chosen from $A$ so that $\pi^{-1}\varkappa\pi$ represents any arbitrary three-term cycle of the $n$ figures; for in
% [sic] the German column prints "zusammmensetzen", with three m, on this page

\[
\pi = \begin{pmatrix} 0,\,1,\,2,\,3 \ldots n-1 \\ a_0,\,a_1,\,a_2,\,a_3 \ldots a_{n-1} \end{pmatrix}
\]

one may choose the first three figures $a_0, a_1, a_2$ arbitrarily, and, if it is necessary in order that $\pi$ belong to $A$, further interchange $a_1$ and $a_2$. In this way there arises from $\varkappa$ one of the two cycles $(a_0,\,a_1,\,a_2)$, $(a_0,\,a_2,\,a_1)$, of which each is the second power of the other. Now if $Q$ is a normal divisor of $A$, and $\varkappa$ a permutation. from $Q$, then $\pi^{-1}\varkappa\pi$ is also contained in $Q$ whenever $\pi$ is contained in $A$, and from this it follows that, if a three-term cycle occurs in $Q$, then $Q$ is identical with $A$.
% [sic] source prints a full stop after "Permutation" before "aus Q"

Our proof now rests on the fact that, if $\varkappa$ is any permutation in $Q$, then $\pi^{-1}\varkappa\pi$, and consequently also

\weq{1}{\lambda = \varkappa^{-1}\pi^{-1}\varkappa\pi}

must occur in $Q$; and it is then to be shown that, if $\varkappa$ is any non-identical permutation, $\pi$ can always be so chosen from $A$ that the permutation $\lambda$ becomes a three-term cycle, and consequently $Q$ identical with $A$.

For this purpose we suppose both $\varkappa$ and $\pi$ decomposed into their cycles, and remark that in the formation of $\lambda$ one need not take into account at all those cycles of $\varkappa$ whose figures remain unchanged by $\pi$, because they cancel each other in $\varkappa^{-1}$ and $\varkappa$. We must now consider the various possible forms of $\varkappa$ one by one.

1. Let $\varkappa$ contain a cycle of more than three figures, say $(1,\,2,\,3 \ldots m)$; we take $\pi = (1,\,2,\,3)$ and obtain

\[
\lambda = \varkappa^{-1}\pi^{-1}\varkappa\pi = (1,\,m,\,m-1 \ldots 2)(2,\,3,\,1,\,4 \ldots m) = (1,\,2,\,4)
\]


% ===== scan leaf 627, printed 602 =====
% !PRINTED 602
\weberpage{602}

Thus a three-term cycle occurs in $Q$.

2. Let $\varkappa$ contain two three-term cycles $(1,\,2,\,3)(4,\,5,\,6)$. We take $\pi = (1,\,3,\,4)$ and obtain

\[
\lambda = (1,\,3,\,2)(4,\,6,\,5)(3,\,2,\,4)(1,\,5,\,6) = (1,\,2,\,5,\,3,\,4).
\]

This permutation $\lambda$, which is contained in $Q$, falls, however, under case 1.

3. Let $\varkappa$ contain a three-term and a two-term cycle $(1,\,2,\,3)(4,\,5)$. (That in $\varkappa$, if it belongs to $A$, a second cycle of even number of terms must also occur, is immaterial here). For $\pi = (1,\,2,\,4)$ we get

\[
\lambda = (1,\,3,\,2)(4,\,5)(2,\,4,\,3)(1,\,5) = (1,\,2,\,5,\,3,\,4),
\]

which again falls under case 1.

4. Let $\varkappa$ contain three transpositions, $(1,\,2)(3,\,4)(5,\,6)$. For $\pi = (1,\,3,\,5)$ it follows that

\[
\lambda = (1,\,2)(3,\,4)(5,\,6)(3,\,2)(5,\,4)(1,\,6) = (1,\,3,\,5)(2,\,6,\,4),
\]

which falls under case 2.

5. Let $\varkappa$ contain two transpositions and one unchanged element $(1,\,2)(3,\,4)(5)$. One sets $\pi = (1,\,2,\,5)$ and obtains

\[
\lambda = (1,\,2)(3,\,4)(5)(2,\,5)(3,\,4)(1) = (1,\,5,\,2).
\]

With this all cases are exhausted when $n > 4$. For $n = 4$ there still remains the one case in which $\varkappa$ consists of two transpositions, and it is just this that brings about the special behaviour at $n = 4$, by which the algebraic solution of the equation of degree $4$ is made possible\footnote{The first complete proof that the general equation of degree higher than the $4$th is not solvable by radicals is due to \gesperrt{Abel} (Crelle's Journal, Bd I, 1826). On the older attempts at a proof by \gesperrt{Ruffini} (1799 to 1806) and their relation to the \gesperrt{Abel} proof, compare the memoir of \gesperrt{Burkhardt}, ``Die Anfänge der Gruppentheorie und \gesperrt{Paolo Ruffini}'' (Abhandlungen zur Geschichte der Mathematik VI. Supplement zu Schlömilch's Zeitschrift. Leipzig 1892).}.

It follows further from this theorem that the symmetric group has no normal divisors other than itself, the alternating group and the unit group. For if $S$ is the symmetric group, $A$ the alternating group, and $Q$ a normal divisor of $S$, then the greatest common divisor $Q'$ of $A$ and $Q$ is a normal divisor of $A$, and is therefore equal to $A$ or equal to 1.


% ===== scan leaf 628, printed 603 =====
% !PRINTED 603
\weberpage{603}
\pdfbookmark[2]{§ 178. Non-metacyclic equations in the field of rational numbers}{bm199}

If $Q'$ is equal to $A$, then $Q$ is either likewise equal to $A$ or equal to $S$. But if $Q'$ is $=1$, then $Q$ contains no permutation of the first kind besides the unit. If, then, $\varkappa$ and $\lambda$ are two permutations of $Q$, different from one another and different from the unit, then $\varkappa^2$ and $\varkappa\lambda$, being of the first kind, must be $=1$; that is, $\lambda$ must be $=\varkappa$. Thus $Q$ can contain at most one permutation $\varkappa$ different from the identical one. But since $Q$ is to be a normal divisor of $S$, for every permutation $\pi$ of the symmetric group we must have $\pi^{-1}\varkappa\pi = \varkappa$; that is, $\varkappa$ may not change when any permutation whatever is carried out in its cycles. But this is possible only when there are only two figures 1, 2 present at all, and $\varkappa = (1,\,2)$. But then $1, (1,\,2)$ is the whole group $S$.

\begin{center}
\S.\,178.

\gesperrt{Non-metacyclic equations in the field of the rational numbers.}
\end{center}

By the theorem of the preceding section the proof has been given that an equation of degree $n$, when $n$ is greater than 4, can no longer be solved algebraically if the coefficients are regarded as independent variables. But of still greater interest is the question whether there are, in the field of the rational numbers, equations of degree $n$ that are not algebraically solvable. The question can be put somewhat more generally, namely thus: whether there are equations with integral coefficients whose group is the symmetric group, and which therefore, according to a mode of expression explained earlier, have no affect.

If one forms the \gesperrt{Galois} resolvent $G(t)=0$ of degree $\Pi(n)$ of a general equation of degree $n$, $f(x)=0$, with indeterminate coefficients $a$, then $G(t)$ is an integral function of the variables $t$ and $a$ which cannot be decomposed into factors that are again rational functions of $t$ and $a$. If one substitutes for the variables $a$ quantities of any field $\Omega$, and $G(t)$ then becomes reducible in this field, the equation $f(x)=0$ has an affect in the domain of rationality $\Omega$. Our question thus comes down to whether one can put for the variables $a$ such rational numbers in $G(t)$ that $G(t)$ remains irreducible in the field of the rational numbers.


% ===== scan leaf 629, printed 604 =====
% !PRINTED 604
\weberpage{604}

This question has found a quite general answer in a memoir by \gesperrt{Hilbert}\footnote{Ueber die Irreducibilit\"at ganzer rationaler Functionen mit ganzzahligen Co\"efficienten. Journal f\"ur Mathematik, Bd. 110.}, in which the general theorem is proved that in an irreducible function of arbitrarily many variables one can substitute for an arbitrary part of the variables rational numbers such that an irreducible function of the remaining variables results. We cannot enter upon this general theorem here.

We shall, however, answer the question posed much more simply, though by far not so generally, by showing that equations without affect can be found for every prime degree.

We have proved in \S.\,153, 9. that a transitive permutation group of $n$ figures which is not the symmetric group, if $n$ is a prime, can contain no single transposition.

By an equation with an affect we have understood one whose group is not the symmetric group. If, then, the irreducible equation $f(x) = 0$ has an affect and its degree is a prime, its group $P$ must be transitive, and it can contain no transposition of two roots. If therefore we adjoin any $n - 2$ of the roots to the domain of rationality, the group must reduce to the unit group, since besides this only the permutation of the two last roots could still remain, which does not occur in $P$. From this it follows, then, that the two last roots are contained in the extended field $\Omega$, or the theorem:

1. \gesperrt{If an irreducible equation whose degree $n$ is a prime has an affect, then any two of its roots can be expressed rationally in terms of the remaining ones.}

From this follows as a corollary:

2. \gesperrt{If the field $\Omega$ contains only real numbers, then an irreducible equation of prime degree $n$ with an affect in $\Omega$ cannot have two imaginary and $n - 2$ real roots.}

Now there are, however, innumerable equations of every arbi-
%CONT>>


% ===== scan leaf 630, printed 605 =====
% !PRINTED 605
\weberpage{605}
%<<CONT
trary degree $n$, with real coefficients, which have two conjugate imaginary and $n - 2$ real roots; for in
\[
f(x) = (x - \alpha_1)(x - \alpha_2) \ldots (x - \alpha_n)
\]
\[
= x^n + a_1 x^{n-1} + a_2 x^{n-2} + \cdots + a_n
\]
one may take $\alpha_1, \alpha_2$ conjugate imaginary at will and the remaining $\alpha$ real. But the number of real and imaginary roots of $f(x)$ does not change when the coefficients are varied continuously within certain finite limits (cf. the seventh section), and consequently there are also such equations whose coefficients are rational numbers, since one can always find rational numbers in arbitrary nearness to any given real numbers.

It remains, then, only to show that these rational coefficients can be so chosen that $f(x)$ becomes irreducible. But this follows from the following theorem:

3. \gesperrt{If $p$ is a prime, $c_0, c_1, c_2 \ldots c_n$ a sequence of integers, of which $c_0$ is not divisible by $p$, while $c_1, c_2 \ldots c_n$ are divisible by $p$, but $c_n$ is not divisible by $p^2$, then the function}
\[
\varphi(x) = c_0 x^n + c_1 x^{n-1} + \cdots + c_n
\]
\gesperrt{is irreducible.}

For if $\varphi(x)$ splits into two factors with rational coefficients, then (by \S.\,2) the factors may also be assumed to have integral coefficients. So let
\[
\varphi(x) = (\alpha_0 x^h + \alpha_1 x^{h-1} + \cdots + \alpha_h)(\beta_0 x^k + \beta_1 x^{k-1} + \cdots + \beta_k),
\]
$h$ and $k$ both greater than zero and their sum equal to $n$.

Since $\alpha_h \beta_k = c_n$, one of the two factors $\alpha_h, \beta_k$ must be divisible by $p$, and the other not. Suppose that $\beta_k$ is divisible by $p$ and $\alpha_h$ not divisible; and since not all the $\beta$ can be divisible by $p$, because otherwise $c_0$ too would be divisible by $p$, let $\beta_\nu$ be not divisible by $p$, and $\beta_{\nu+1}, \beta_{\nu+2} \ldots \beta_k$ divisible by $p$. The coefficient of $x^{k-\nu}$ in the product of the two factors is then $\alpha_h \beta_\nu + \alpha_{h-1}\beta_{\nu+1} + \cdots$, and hence not divisible by $p$. It would therefore have to be that $k - \nu = n$, which is not possible, since already $k < n$ must hold.

If, then, one sets for the coefficients of $f(x)$:
\[
a_1 = \frac{c_1}{c_0}, \quad a_2 = \frac{c_2}{c_0} \cdots a_n = \frac{c_n}{c_0},
\]
% [page alignment mirrored 2026-08-07] Printed 605 ends immediately after
% the coefficient display above.  The following theorem/prose belongs to
% printed 606 and has been moved to i0631.
%CONT>>


% ===== scan leaf 631, printed 606 =====
% !PRINTED 606
\weberpage{606}
\pdfbookmark[2]{§ 179. Solution by real radicals}{bm200}
%<<CONT
% [page alignment mirrored 2026-08-07] Printed 606 begins with this
% continuation from i0630 and ends after the §179 footnote.  The former
% "We assume ..." tail is printed on 607 and already belongs to i0632.
then $f(x)$ is irreducible, and one still has freedom enough in the choice of the integers $c$ to bring the rational fractions $a_1, a_2 \ldots a_n$ arbitrarily near to any arbitrarily given system of values.

But with this the theorem is proved:

4. \gesperrt{Of every arbitrary prime degree there are infinitely many equations with rational coefficients without affect.}

The proof that has been given here shows that such affect-free equations can be found whose coefficients fill a finite region everywhere densely — namely the region in which the real coefficients of the equations with only two
conjugate imaginary roots lie. The proof is therefore, apart from the fact that it is restricted to prime degrees, also not exhaustive in this respect, that it gives us no information about the remaining regions, in which in all probability the matter stands just the same.

\begin{center}
\S.\,179.

\gesperrt{Solution by real radicals.}
\end{center}

In the solution of the cubic equations with real coefficients, the two cases in which the discriminant is negative or positive have been distinguished from ancient times. In the second case, although the equation then has precisely three real roots, one has, in the application of the \gesperrt{Cardan formula}, to extract the cube root of an imaginary expression, and the attempt to represent the roots in real form leads always back to a cubic equation of the same character. For this reason this case has been called the \gesperrt{casus irreducibilis}. The cubic equations of the casus irreducibilis are a special case of the \gesperrt{cyclic equations with real roots}, which we came to know in \S.\,165, and in whose solution likewise roots of imaginary quantities, or what comes to the same thing, divisions of angles, occur. That these roots of imaginary quantities or divisions of angles can in no way be avoided, we can now prove.\footnote{Cf. \gesperrt{H\"older}, Mathematische Annalen, Bd. 38; \gesperrt{Kneser}, ebendas., Bd. 41.}


% ===== scan leaf 632, printed 607 =====
% !PRINTED 607
\weberpage{607}

% [German fragment carries a correction note: the running head of printed 607
% ("§. 179. | Reelle Radicale. | 607") had been transcribed as a section heading and was
% removed; running heads are excluded page furniture. Mirrored here for symmetry.]
\quad We assume a \gesperrt{real domain of rationality} $\Omega$,
and take in it a normal equation $g(t) = 0$,
i.\,e. an irreducible equation whose roots are all rationally
expressible in terms of an arbitrary one of them, $\varrho$. If one
of the roots is real, then all the others must be real
as well, and $g(t)$ thus has either only real or only pairwise
conjugate imaginary roots. $g(t) = 0$ can be the \gesperrt{Galois}
resolvent of any given equation $F(x) = 0$, whether irreducible
or reducible, and if therefore $g(t)$ has only
real roots, then $F(x)$ too can have only real roots, % [sic: the German page prints "relle" for "reelle"]
because the roots of $F(x)$ are rationally representable
in terms of a root of $g(t)$. If among the roots of $F(x)$ there are
imaginary ones, then $g(t)$ has only imaginary roots. But if $F(x)$
has only real roots, then the roots of $g(t)$ are also real,
because they can indeed be expressed rationally in terms of the
roots of $F(x)$.

\quad Now if $g(t)$ is reduced by the adjunction of a root $\varepsilon$ of an
irreducible equation $\chi = 0$ \gesperrt{whose degree is a prime},
then, as we have seen in \S.\,157, all the
roots of $\chi$ are contained in $\Omega(\varrho)$, and if therefore
$\varrho$ is real, then all the roots of $\chi$ are real.

\quad 1. \gesperrt{A normal equation with real roots can
therefore be reduced only by such irreducible
equations of prime degree as have only
real roots.}

\quad We now ask whether a reduction of the normal equation
can be effected by the adjunction of a real radical $\sqrt[p]{a}$.
We may here assume that the degree $p$ of the radical
is a prime, because every radical can be reduced to a series of
extractions of roots of prime degree. We can also
presuppose that $a$ is not the $p$th power of a rational
quantity, because otherwise the real root $\sqrt[p]{a}$ would be rational.

\quad Under this assumption, as we now wish to prove, % [sic: the German page prints "Vorraussetzung" for "Voraussetzung"]
$x^p - a$ is irreducible. For if we denote one, e.\,g. the real
value $\sqrt[p]{a}$, by $\alpha$, and an imaginary $p$th root of unity by $\varepsilon$,
then
\weq{1}{\alpha,\ \varepsilon\alpha,\ \varepsilon^2\alpha \ldots \varepsilon^{p-1}\alpha}
%CONT>>


% ===== scan leaf 633, printed 608 =====
% !PRINTED 608
\weberpage{608}
%<<CONT
the roots of the equation
\weq{2}{x^p - a = 0.}

If this equation is reducible, then
\weq{3}{x^p - a = f_1(x) f_2(x),}
and $f_1$, $f_2$ are functions in $\Omega$ of degree lower than $p$.
One part of the roots (1) will fall on $f_1 = 0$, another part on $f_2 = 0$. If
then $\mu$ is the degree of $f_1$, then for some exponent $\lambda$
\[ \varepsilon^\lambda \alpha^\mu = b \]
must be a quantity in $\Omega$ [the term in $f_1(x)$ independent of $x$]; hence, if one
raises to the $p$th power,
\weq{4}{a^\mu = b^p.}

Now $0 < \mu < p$, and therefore $\mu$ and $p$ are relatively prime, so that two integers
$h$ and $k$ can be determined from the equation
\[ \mu h + p k = 1 \]
Then by (4)
\[ a = a^{\mu h} a^{p k} = (b^h a^k)^p; \]
thus, contrary to the hypothesis, $a$ would be the $p$th power of a quantity in $\Omega$.

This theorem, as one sees, is independent of the assumption that $\Omega$ be a
real field. But if we take it to be real, and $p$ is not equal to 2, then the
equation (2) has imaginary roots, and by 1. cannot reduce a normal equation $g(t) = 0$.
But if $p = 2$, then $g(t)$ can be reduced by this equation (2)
only if $a$ is positive and the degree of $g(t)$ is an \gesperrt{even} number; hence:

2.~\gesperrt{A normal equation of odd degree cannot be reduced by the adjunction of a
real radical.}

This case occurs in the casus irreducibilis of the cubic equation, when one has adjoined the
square root of the discriminant to the domain of rationality of the coefficients.
For since the discriminant is positive, the domain of rationality remains
real, and it also remains real however many real radicals one adjoins. If
the equation is thus to be solvable by real radicals, it must at last split under such
adjunctions, which by 2. is impossible.


% ===== scan leaf 634, printed 609 =====
% !PRINTED 609
\weberpage{609}
\pdfbookmark[2]{§ 180. Metacyclic equations of prime degree}{bm201}
In the same case are found the cyclic and in general all irreducible
\gesperrt{Abel}ian equations of odd degree, which therefore can never be solved by real
radicals.

Accordingly we can, for example, distinguish the cubic equations in the field of the
rational numbers into three or four kinds, which all occur in certain geometric
problems famous from antiquity. We here leave the reducible equations out of account.
The degree of the \gesperrt{Galois} group must then always be divisible by 3
(\S.\,154, 7.) and is therefore either equal to 3 or equal to 6. By \S.\,157 the group
can therefore never be reduced by adjunction of mere square roots to a lower than
the third degree; hence the equation also cannot be solved by
square roots. The corresponding geometric problems cannot be solved with
compass and ruler.

If the group is of degree 3, we have a cyclic equation. Here belong the
equations arising from the division of the circle, e.g. the one on which the construction
of the regular heptagon depends. The three roots of such an equation are real
and cannot be expressed by real radicals. Among the cubic
equations with a group of degree $6$ we must distinguish those with
positive and those with negative discriminant.

The first belong to the casus irreducibilis and can be reduced to the
equation on which the trisection of an arbitrary angle depends. To the
equations with negative discriminant belong also, as special cases, the pure
cubic equations $x^3 = a$, when $a$ is not a cube number. For $a = 2$ there results
the equation on which the \gesperrt{Delian problem} of the duplication of the cube depends.

{\centering\S.\,180.\par}
{\centering \gesperrt{Metacyclic equations of prime degree.}\par}

The general conditions which we found in \S.\,176 are not yet
simple enough to permit an immediate application to the finding of metacyclic
equations or to the decision on the solvability of a proposed
%CONT>>


% ===== scan leaf 635, printed 610 =====
% !PRINTED 610
\weberpage{610}
%<<CONT
equation to admit solution by radicals. We therefore derive, at first under the
assumption that the degree $n$ of the equation is a prime, another criterion first set
up by \gesperrt{Galois}.

Now let $f(x) = 0$ be an irreducible equation of degree $n$, and $n$ a prime greater
than 2. If $f(x)$ is to be algebraically solvable, then by \S.\,176 its group $P$ must
be metacyclic, i.e. there must be a chain of groups
\weq{1}{P,\ P_1,\ P_2 \ldots P_{\mu-1},\ 1}
each of which is a normal divisor of the one immediately preceding, with prime index.
By successive adjunction of roots of cyclic equations of prime degree the group of the
equation is reduced from $P$ to $P_1, P_2 \ldots P_{\mu-1}, 1$.

The function $f(x)$ itself, as was already shown in \S.\,176, cannot be reduced before
the last adjunction has been made, whereupon it splits into $n$ linear factors. The
degree of the last cyclic equation, and hence the degree of the next-to-last group
$P_{\mu-1}$, must therefore by \S.\,158, 3. be equal to $n$.

The degree of an element of a group is always a divisor of the degree of the group, and
hence $P_{\mu-1}$ can contain, besides the identical one, only permutations of order
$n$. It is therefore identical with the period $1, \pi, \pi^2 \ldots \pi^{n-1}$ of one
of its elements. From this it follows by \S.\,153 that $\pi$ must be a cyclic
permutation of all $n$ figures, and we can choose the designation of the roots of
$f(x)$ so that
\[ \pi = (0,\ 1,\ 2 \ldots n-1) \]
holds. If we denote the roots by $x_z$, and stipulate that $x_z = x_{z'}$ shall hold
when $z \equiv z' \, (\mathrm{mod}\ n)$, then by $\pi$ every $z$ passes into $z+1$, by
$\pi^2$ into $z+2$, and so on, and we can therefore say that the group $P_{\mu-1}$
consists of the substitutions for $z$
\weq{2}{(z,\ z+b), \qquad b = 0,\ 1,\ 2 \ldots n-1}
Every irreducible metacyclic equation of prime degree can therefore be transformed by
adjunction of radicals into a cyclic equation.

The substitutions $(z, z+b)$ form a special case of the more general linear
substitution
\weq{3}{(z,\ az+b),}
in which $a$ and $b$ are fixed numbers, which may also be reduced according to the
modulus $n$, but where for $a$ the value 0 of course
%CONT>>


% ===== scan leaf 636, printed 611 =====
% !PRINTED 611
\weberpage{611}
%<<CONT
is to be excluded, because otherwise by (3) all the different $z$
would pass into the same $b$, and thus no permutation of the quantities $x_z$ would be
expressed at all. But if $a$ is different from 0 $(\mathrm{mod}\ n)$, then a permutation
will always be represented by (3), because then $az+b \equiv az'+b$ can hold only when
$z \equiv z'$. We call (3) a substitution for $z$, which produces a permutation among
the $x$. The number of the different linear substitutions of the form (3)
is $n(n-1)$. To them correspond just as many different permutations among the $x$;
for it is only then that, for every $z$,
\[ az + b \equiv a'z + b' \, (\mathrm{mod}\ n) \]
can hold, when $b \equiv b'$ (to be concluded from $z = 0$) and $a \equiv a'$ (to be
concluded from $z = 1$).

The totality of the permutations that are represented by (3) forms a group;
for let
\[ \lambda = (z,\ az+b), \qquad \lambda' = (z,\ a'z+b') \]
be two of these permutations; then
\[
\lambda \lambda' =
\begin{pmatrix} x_z \\ x_{az+b} \end{pmatrix}
\begin{pmatrix} x_z \\ x_{a'z+b'} \end{pmatrix}
=
\begin{pmatrix} x_z \\ x_{a'(az+b)+b'} \end{pmatrix},
\]
hence
\weq{4}{\lambda \lambda' = \lambda'' = (z,\ aa'z + a'b + b'),}
which is again of the form $(z, a''z+b'')$.

The group formed by all substitutions of the form (3), which is a
generalisation of the cyclic group (2), we shall call a
\gesperrt{linear group}\footnote{\gesperrt{Kronecker} calls only this group
metacyclic.}; the permutations contained in it shall be called
linear permutations.

In the linear group various divisors are contained. Among them we find a
group of degree $n-1$, namely $(z, az)$, which may itself in turn have divisors.

The composition rule (4) shows how we can form all divisors of the linear
group. If $\lambda = (z, az+b)$ is a linear substitution, then by (4), through
repetition, we obtain for every exponent $h$
\[ \lambda^h = [z,\ a^h z + (1 + a + \cdots + a^{h-1})b], \]
and for $a = 1$
\[ \lambda^h = (z,\ z + hb). \]


% ===== scan leaf 637, printed 612 =====
% !PRINTED 612
\weberpage{612}

If, then, in a group $L$ there occurs a substitution $\lambda$ in which $a = 1$ and $b$ is different from zero, then by the above $L$ contains the whole \gesperrt{cyclic group} $(z, z + b)$ $(b = 0, 1, \ldots n - 1)$.

If $a$ is not equal to 1, we set
\[
1 + a + \cdots + a^{h-1} = \frac{a^h - 1}{a - 1},
\]
and conclude from this that, if $e$ is the least positive exponent for which $a^e \equiv 1\ (\mathrm{mod}\ n)$, then $\lambda$ is of degree $e$.

The period of the substitution $\lambda$ is, unless $a = 1$, an intransitive group. There is a figure $z$ which is not altered by $\lambda$ and its repetitions, which is determined from the congruence $z \equiv az + b\ (\mathrm{mod}\ n)$, and can be denoted by
\[
z_0 \equiv \frac{b}{a - b}\ (\mathrm{mod}\ n)
\]
% [sic] source denominator is "a-b"; fixed-point algebra would give "1-a"
if by $\dfrac{\mu}{\nu}\ (\mathrm{mod}\ n)$ is understood an integer which, by multiplication with the factor $\nu$, becomes congruent to $\mu$ with respect to the modulus $n$.

The cyclic group $(z, z + b)$ is transitive and of degree $n$. It can be shown that every transitive linear group $(\mathrm{mod}\ n)$ must contain the cyclic group.

Let $g$ be a primitive root of $n$ and $\alpha$ the index of $a$, i.\,e. $g^\alpha \equiv a\ (\mathrm{mod}\ n)$ (\S.\,136). For brevity we call $\alpha$, for the moment, also the \gesperrt{index of the substitution} $\lambda = (z, az + b)$. The index may be assumed to be less than $n - 1$ and equal to or greater than zero. If it is equal to zero, then $\lambda$ is either the identical substitution, or it belongs to the cyclic group. By (4) the rule holds that the index of a compounded substitution is equal to the sum of the indices of the components.

From this it follows that all the indices of the substitutions of a linear group $L$ are multiples of the least positive one among them, and that consequently, if $\alpha_0$ is this least positive index and $a^{\alpha_0} = a_0$ is set, all the substitutions of $L$ can be represented in the form
% [sic] The source prints a^{\alpha_0}=a_0 here; the preceding definition uses g^{\alpha}=a. Retained source-faithfully.
\[
\lambda = (z, a_0^h z + b)
\]
in which $a_0$ is held fixed, and only $h$ and $b$ run through certain sequences of numbers.

A substitution of the form $\lambda_0 = (z, a_0 z + b_0)$ must occur in the group $L$.


% ===== scan leaf 638, printed 613 =====
% !PRINTED 613
\weberpage{613}

If we now form, according to (4), the composition
\[
\lambda \lambda_0^{-h} = \left(z,\, z + a_0^{-h} b + \frac{a_0^{-h} - 1}{a_0 - 1} b_0\right),
\]
then we obtain the identical substitution when, and only when,
\[
\frac{b_0}{a_0 - 1} \equiv \frac{b}{a_0^h - 1}\ (\mathrm{mod}\ n).
\]

If this congruence is satisfied for \gesperrt{all} substitutions $\lambda$ of $L$, then the element $\dfrac{b_0}{a_0 - 1}$ remains unchanged by the whole group, and $L$ is intransitive. Consequently, in a transitive group $L$ there is always an element $\lambda$ for which $\lambda \lambda_0^{-h}$ is not the identical substitution, although the index is zero, and $L$ therefore contains the whole cyclic group.

If one forms the various powers of $\lambda_0$, it turns out that the exponent $h$ in the substitutions $\lambda$ of the group $L$ can take every value. If $e$ is the least positive exponent that satisfies the condition $a_0^e \equiv 1\ (\mathrm{mod}\ n)$, then the values $h = 0, 1, 2 \ldots e - 1$ therefore occur, and here $e$ is a divisor of $n - 1$. From the composed substitution
\[
(z, a_0^h z + b)(z, z + 1) = (z, a_0^h z + b + 1),
\]
which after all must also occur in $L$, one sees further that with every value of $h$ each of the values $b = 0, 1, 2, \ldots n - 1$ occurs in combination, and the whole group $L$ therefore consists of the substitutions
\weq{5}{\lambda = (z, a_0^h z + b), \quad \begin{array}{l} h = 0, 1 \ldots e - 1 \\ b = 0, 1 \ldots n - 1' \end{array}}
% [sic] source prints a raised apostrophe-like mark after "n-1" in the b-row
and is of degree $en$. Conversely, every system of substitutions of this form constitutes a group.

Of the linear groups the following theorem now holds:

I. \gesperrt{If a transitive linear group $L$ is a normal divisor of another permutation group $P$ of the same $n$ figures, then $P$ is also a linear group.}

The theorem can be proved as follows:
If $\pi$ is an arbitrary permutation:
\[
\pi = \begin{pmatrix} 0, & 1, & 2 \ldots n - 1 \\ a_0, & a_1, & a_2 \ldots a_{n-1} \end{pmatrix},
\]
%CONT>>


% ===== scan leaf 639, printed 614 =====
% !PRINTED 614
\weberpage{614}

%<<CONT
which can also be represented by $(z, a_z)$, then one can determine an integral function $\varphi(z)$ of $z$ with rational coefficients, whose degree is not higher than $n - 1$, so that $a_z = \varphi(z)$ can be set. One need only determine the $n$ coefficients in $\varphi(z)$ from the $n$ linear equations
\[
a_0 = \varphi(0), \quad a_1 = \varphi(1), \ldots a_{n-1} = \varphi(n - 1)
\]
For this one may use the \gesperrt{Lagrange} interpolation formula derived in \S.\,29, which, however, can be given a still simpler form, since here only congruences with respect to the modulus $n$ come into question. Namely, if one sets
\[
\psi(z) = z(z - 1)(z - 2) \ldots (z - n + 1),
\]
the interpolation formula just mentioned gives
\[
\varphi(z) = \psi(z)\left(\frac{a_0}{\psi'(0) z} + \frac{a_1}{\psi'(1)(z - 1)} + \cdots + \frac{a_{n-1}}{\psi'(n-1)(z - n + 1)}\right).
\]
Now, as we have seen in \S.\,136, the congruence
\[
\psi(z) \equiv z^n - z\ (\mathrm{mod}\ n)
\]
is identical; hence also
\[
\psi'(z) \equiv n z^{n-1} - 1 \equiv -1\ (\mathrm{mod}\ n),
\]
and therefore we can represent $\varphi(z)$ with integral coefficients as follows:
\weq{6}{\varphi(z) \equiv -a_0 \frac{\psi(z)}{z} - a_1 \frac{\psi(z)}{z - 1} - \cdots - a_{n-1}\frac{\psi(z)}{z - n + 1}\ (\mathrm{mod}\ n).}

If, as in our problem, the numbers $a_0, a_1 \ldots a_{n-1}$ agree, apart from the order, with the numbers $0, 1, 2 \ldots n - 1$, then, if $n > 2$, we have $a_0 + a_1 + \cdots + a_{n-1} \equiv 0\ (\mathrm{mod}\ n)$, and hence $\varphi(z)$ is at most of degree $n - 2$, which, however, is of no significance for our proof.

Accordingly, then, every arbitrary permutation $\pi$ can be represented by a substitution $[z, \varphi(z)]$.

Now if $L$ is a transitive linear group and at the same time a normal divisor of another group $P$, and if $\lambda$ is an arbitrary permutation of $L$, and $\pi$ a likewise arbitrary permutation of $P$, then $\pi^{-1}\lambda\pi = \lambda'$ is contained in $L$, and $\lambda\pi = \pi\lambda'$. If, by what has just been proved, we set $\pi = [z, \varphi(z)]$ and choose for $\lambda$ the cyclic substitution $(z, z + 1)$, which by the hypothesis of transitivity occurs in $L$, then $a'$ and $a$ must admit of being determined so that
\[
(z, z + 1)[z, \varphi(z)] = [z, \varphi(z)](z, a'z + a),
\]
%CONT>>


% ===== scan leaf 640, printed 615 =====
% !PRINTED 615
\weberpage{615}

%<<CONT
or, if one carries out the composition,
\weq{7}{\varphi(z + 1) \equiv a' \varphi(z) + a\ (\mathrm{mod}\ n)}
for every integer $z$. Since $\varphi(z)$ does not reach the degree $n$, the coefficients of the individual powers of $z$ on the two sides of (7) must (\S.\,136) be congruent, and from the comparison of the coefficients of the highest powers of $z$ it follows that $a' \equiv 1\ (\mathrm{mod}\ n)$, and hence
\[
\varphi(z + 1) \equiv \varphi(z) + a\ (\mathrm{mod}\ n).
\]
If one here puts $z + 1, z + 2 \ldots z + h - 1$ for $z$, it follows that
\[
\varphi(z + h) = \varphi(z) + ah,
\]
where $h$ can be congruent to any arbitrary integer with respect to the modulus $n$. In this one may now set $z = 0$, and obtains, on writing $z$ again for $h$ and setting $\varphi(0) = b$,
\weq{8}{\varphi(z) = az + b.}

The group $P$ therefore consists entirely of linear permutations, and our theorem is proved.

We add further the remark, which follows immediately from an inspection of the formulae, that the \gesperrt{cyclic group is a normal divisor of every linear group in which it is contained at all}.

If we apply the theorem just proved to the group of the metacyclic equation, we obtain the theorem of \gesperrt{Galois}:

II. \gesperrt{The group of a metacyclic equation of prime degree is linear.}

For if we return to the chain of groups (1), we have seen that $P_{\mu - 1}$ is the cyclic group. If $P_{\mu - 1}$ is not identical with $P$, then $P_{\mu - 1}$ is a normal divisor of $P_{\mu - 2}$, and hence $P_{\mu - 2}$ is linear. $P_{\mu - 2}$ is in turn a normal divisor of $P_{\mu - 3}$, hence $P_{\mu - 3}$ too is linear, and so on, until we reach the conclusion that $P$ itself must also be linear. At the same time it follows that $P_{\mu - 1}$ is a normal divisor of all the preceding groups, and hence also of $P$ itself.

The converse theorem also holds:

III. \gesperrt{Every irreducible equation of prime degree whose group is linear is metacyclic.}

To prove it, it suffices to show that every transitive linear group $L$ has a normal divisor $L'$ whose index
%CONT>>


% ===== scan leaf 641, printed 616 =====
% !PRINTED 616
\weberpage{616}
%<<CONT
a prime, and which itself, if it is not the unit group, is a transitive linear group.

But this is shown immediately by the representation of the substitutions of the group $L$ by the formula (5).

If in it $e=1$, then $L$ is the cyclic group of degree $n$ with the normal divisor 1. But if $e>1$, let $p$ be a prime dividing $e$, and $e=pe'$. The linear group $L'$ of degree $ne'$, which consists of the substitutions:

\[
\lambda' = (z, a_0^{ph} z + b), \quad
\begin{aligned}
h &= 0, 1, 2 \ldots e'-1\\
b &= 0, 1, 2 \ldots n-1,
\end{aligned}
\]

is certainly a divisor of $L$ of index $p$. But it is also normal, as one recognises without difficulty from the rule of composition (4).

The concepts of the transitive linear and of the metacyclic groups thus coincide completely for the equations of prime degree, and we can therefore in what follows also use the two expressions synonymously.

If we wish to pass from a group $P$ to a conjugate group $P' = \pi^{-1}P\pi$, this transition is effected by carrying out the permutation $\pi$ in the cycles of all the permutations of $P$. $P'$ will therefore agree with $P$ up to the designation of the roots, and is therefore, if $P$ is linear, also to be called linear, even though a change in the numbering of the roots is necessary in order to represent it by linear substitutions.

To simplify the application, let it further be observed that one can derive the \gesperrt{full linear group} by repetition and composition of the two substitutions

\weq{9}{s = (z, z+1), \quad t = (z, gz),}

and every divisor $L$ of the full linear group likewise from

\weq{10}{s = (z, z+1), \quad t^{\alpha_0} = (z, a_0 z)}

One therefore calls the two substitutions (9) or (10) the \gesperrt{generating substitutions of these groups}.

If we set $a_0 = g^2$, we obtain a frequently occurring linear group $(z, az+b)$, in which $a$ runs only through the quadratic residues of $n$, and which has been called by \gesperrt{Kronecker} the \gesperrt{half-metacyclic group}.

This representation by the generating substitutions admits a simple inference concerning the relation of the meta-
%CONT>>


% ===== scan leaf 642, printed 617 =====
% !PRINTED 617
\weberpage{617}
%<<CONT
cyclic groups to the symmetric and the alternating group.

The permutation arising from $s$ consists of a single cycle with an odd number of terms $(0, 1, 2 \ldots n-1)$ and therefore belongs to the alternating group. The substitution $t$ leaves the index 0 unchanged and yields for the remaining figures again a single cycle. For by $t$, 1 passes into $g$, $g$ into $g^2$, $g^2$ into $g^3$, etc., and since the powers $1, g, g^2 \ldots g^{n-2}$, apart from their order, coincide with the figures $1, 2 \ldots n-1$, $t$ corresponds to the cycle $(1, g, g^2 \ldots g^{n-2})$, which consists of an even number of terms and consequently belongs to the permutations of the second kind, and is therefore not contained in the alternating group. On the other hand $t^2$ is again contained in it.

From this the result follows:

\gesperrt{IV. The full metacyclic group is not a divisor of the alternating group. The greatest common divisor of the two groups is the half-metacyclic group.}

Every transitive divisor of the metacyclic group is itself metacyclic.

The condition for the metacyclic equations can be given various forms, which can be derived from what has gone before.

The full linear group is contained in the symmetric group as a divisor of index $\nu = 1 . 2 . 3 \ldots (n-2)$. A function $y$ of the $n$ variables $x_0, x_1 \ldots x_{n-1}$ belonging to the full linear group, which we may call a \gesperrt{metacyclic function}, therefore satisfies a resolvent $F(y)=0$ of degree $\nu$, whose coefficients are symmetric functions of the $x$.

If now for $x$ one substitutes the roots of a metacyclic equation $f(x)=0$, then $y$ becomes rational.

If conversely the function $y$ becomes rational by the substitution of the roots of an irreducible equation $f(x)=0$ for the $x$, while $F'(y)$ remains different from zero, then $f(x)$ is metacyclic; for then the group of $f(x)=0$ is certainly a divisor of the full metacyclic group, and therefore, since $f(x)$ is irreducible, itself metacyclic. But it also suffices for the algebraic solubility of $f(x)=0$ if the resolvent $F(y)=0$ has any rational root at all, which
%CONT>>


% ===== scan leaf 643, printed 618 =====
% !PRINTED 618
\weberpage{618}
%<<CONT
is not a double root. For the different roots of this resolvent belong to conjugate groups, and if therefore one of these roots is rational, then one of the conjugate groups is metacyclic. We therefore have the theorem:

\gesperrt{V. The necessary and sufficient condition for the solvability of the equation $f(x)=0$ by radicals is that the resolvent of degree $\nu$, $F(y)=0$, which one has arranged, by a suitable choice of the function $y$, so that it receives no double roots, has a rational root.}

Another form of this condition results in the following way.

Among the permutations of a linear group, only the identical one leaves any two figures unchanged. For if $(z, az+b)$ does not alter two figures, then the congruence $az+b \equiv z \pmod n$ must have two distinct solutions. But this is possible only if $a \equiv 1$, $b \equiv 0 \pmod n$.

If, then, the group $P$ of $f(x)$ is metacyclic, it reduces itself, by the adjunction of two arbitrary roots, to the unit, and consequently all the roots are rationally expressible in terms of any two among them. Hence:

\gesperrt{VI. If an irreducible equation of prime degree is metacyclic, then all the roots are rationally expressible in terms of any two among them.}

But this theorem also holds conversely, which we can prove simply in the following manner.

Suppose an irreducible equation of prime degree $n$ has the property that all the roots are rationally expressible in terms of any two among them. If, say, $x_\nu = \psi(x_0, x_1)$ is such a representation, then all the permutations of the group $P$ of our equation can be applied to this relation, and if therefore one of these permutations leaves $x_0$ and $x_1$ unchanged, then it also leaves all the $x_\nu$ unchanged, i.e. it is the identical permutation. \gesperrt{Hence $P$ contains, besides the identical permutation, none that leaves two figures unaltered.}

If now one of the permutations $\pi = cc_1 \ldots$ of $P$ contains two or more different cycles $c, c_1 \ldots$, and $c$ is of degree $h$, but $c_1$ of a higher degree, then $\pi^h = c_1 \ldots$, and $\pi^h$ leaves
%CONT>>


% ===== scan leaf 644, printed 619 =====
% !PRINTED 619
\weberpage{619}
%<<CONT
the figures of $c$ unchanged, while it is nevertheless not the identical substitution, because $c_1^h$ does not leave all figures unchanged. But this is not possible when $h>1$. Hence either $\pi$ must leave one figure unchanged, or it must consist of cycles of equal degree. But since $n$ is a prime, $\pi$ can in this case form only a single cycle of degree $n$.

Thus $P$ contains, besides the identical permutation, only cyclic permutations of degree $n$, which we shall denote by $\gamma$, and permutations $\varkappa$ which leave one figure unchanged and moreover consist of cycles of equal degree.

Each of the permutations $\varkappa$ by which the figure 0 is left unchanged we shall denote by $\varkappa_0$. Likewise $\varkappa_1, \varkappa_2 \ldots \varkappa_{n-1}$ denote the permutations $\varkappa$ which leave the figures 1, 2 \ldots $n-1$ unchanged. Since $P$ is transitive, there must be just as many $\varkappa_0$ as $\varkappa_1$, as $\varkappa_2$, etc.

For if $\pi$ is a permutation which carries 0 into 1, then every $\pi^{-1}\varkappa_0\pi$ is a $\varkappa_1$, and conversely every $\pi\varkappa_1\pi^{-1}$ is a $\varkappa_0$, and likewise for the remaining figures. If then $\mu$ is the number of the $\varkappa_0$, $\nu$ the number of the $\gamma$, $m$ the degree of $P$, then, since the identical permutation is also to be added, we have

\weq{11}{m = \mu n + \nu + 1.}

Now the $\varkappa_0$, together with the identical permutation, form a group $Q$ of degree $\mu+1$ and a divisor of $P$, namely the totality of all permutations of $P$ which leave 0 unchanged. If then $\pi_1, \pi_2 \ldots \pi_{n-1}$ are permutations in $P$ which carry 0 into $1, 2 \ldots n-1$, then $Q\pi_1$ is the totality of all those permutations of $P$ which carry 0 into 1, and so forth.

We thus obtain the decomposition of $P$ into the subsidiary groups

\[
P = Q + Q\pi_1 + Q\pi_2 + \cdots + Q\pi_{n-1},
\]

from which follows

\weq{12}{m = n(\mu+1),}

and hence from (11) and (12) $\nu = n-1$.

There are therefore $n-1$ and no more cyclic permutations of degree $n$ in $P$, and these consequently form, again together with the unit, a cyclic group

\[
C = 1, \gamma, \gamma^2 \ldots \gamma^{n-1},
\]

since with $\gamma$ all powers of $\gamma$ likewise occur in $P$.

Now if $\gamma$ is a cyclic permutation of degree $n$, then every $\pi^{-1}\gamma^h\pi$ is also cyclic, and hence, if $\pi$ belongs to $P$
%CONT>>


% ===== scan leaf 645, printed 620 =====
% !PRINTED 620
\weberpage{620}
%<<CONT
belongs, must also be contained in $C$. $C$ is therefore a normal
divisor of $P$, and consequently, by Theorem I., the group $P$
must be linear.

VII.~\gesperrt{An irreducible equation of prime degree,
in which all the roots are rationally expressible in terms of
any two of them, is metacyclic.}

From these theorems we further infer a remarkable property
of metacyclic equations, first noticed by \gesperrt{Kronecker},
for the case of a real domain of rationality\footnote{Ueber algebraisch auflösbare Gleichungen. Monatsbericht der Berliner
Akademie, 14. April 1856.}.
If in such an equation two roots are real, then it follows
from VI that all its roots are real. But such an equation,
since it is of odd degree, must always have one real
root. Hence follows

VIII.~\gesperrt{An irreducible metacyclic equation of
odd prime degree with real coefficients has either
all its roots real or only one.}

If all the roots are real, the discriminant, as a product
of nothing but squares of real quantities $(x_h - x_k)^2$, is positive. If only
one root is real, then to every complex factor of the
discriminant $(x_h - x_k)^2$ there corresponds a conjugate one, and their product is
positive. Only when $x_h$ and $x_k$ are conjugate imaginary, so that $x_h - x_k$
is purely imaginary, is $(x_h - x_k)^2$ negative, and the discriminant
thus receives the factor $-1$ for every pair of conjugate imaginary
roots. The sign of the discriminant is therefore $(-1)^{\frac{n-1}{2}}$.
From this follows:

IX.~\gesperrt{When $n \equiv 1\ (\mathrm{mod}\ 4)$, the
discriminant is always positive; but if $n \equiv 3\ (\mathrm{mod}\ 4)$,
the sign of the discriminant decides which of the two
cases of Theorem VIII. occurs.}


% ===== scan leaf 646, printed 621 =====
% !PRINTED 621
\weberpage{621}
\pdfbookmark[2]{§ 181. Application to the metacyclic equations of the 5th degree}{bm202}
% Faint lines above §181 are reversed verso show-through; no direct recto body text is present.

\begin{center}
\S.\,181.

\gesperrt{Application to the metacyclic equations}
\gesperrt{of degree} $5$.
\end{center}

\quad We make one further application of the general theorems
to the equations of degree $5$. If $n = 5$, the cyclic
group has $5$, the half-metacyclic $10$, and the full metacyclic
group $20$ permutations. For a general equation of degree $5$, a metacyclic
function satisfies a resolvent of degree $6$,
and if this equation of degree $6$ has a rational
root, then the given equation of degree $5$ is metacyclic.
The half-metacyclic functions, i.\,e. the functions
that admit the permutations of the half-metacyclic group,
satisfy a resolvent of degree $12$, which however, after
adjunction of the square root of the discriminant, splits into two
factors of degree $6$.

\quad If one applies the generating substitutions for the modulus $5$
\[
s = (z, z+1),\quad t = (z, 2z),\quad t^2 = (z, 4z)
\]
to the figures $0, 1, 2, 3, 4$, there result the permutations
\weq{1}{
\begin{array}{lccccc}
 & 0, & 1, & 2, & 3, & 4\\
(s) & 1, & 2, & 3, & 4, & 0\\
(t) & 0, & 2, & 4, & 1, & 3\\
(t^2) & 0, & 4, & 3, & 2, & 1.
\end{array}
}

\quad If we apply the substitutions $s, t^2, t$ to the pairs of
figures $(0,1), (1,2), (2,3), (3,4), (4,0)$, it follows that
\weq{2}{
\begin{array}{lccccc}
 & (0,1), & (1,2), & (2,3), & (3,4), & (4,0)\\
(s) & (1,2), & (2,3), & (3,4), & (4,0), & (0,1)\\
(t^2) & (0,4), & (4,3), & (3,2), & (2,1), & (1,0);
\end{array}
}
these pairs therefore remain unchanged, in their totality, by $s$ and
by $t^2$, and consequently by the half-metacyclic group,
and a symmetric function of the corresponding
pairs of roots is half-metacyclic. The substitution $t$ effects
the permutations
\weq{3}{
\begin{array}{lccccc}
 & (0,1), & (1,2), & (2,3), & (3,4), & (4,0)\\
(t) & (0,2), & (2,4), & (4,1), & (1,3), & (3,0),
\end{array}
}
%CONT>>


% ===== scan leaf 647, printed 622 =====
% !PRINTED 622
\weberpage{622}
%<<CONT
thus carries the five pairs over into five others; and since there
are only ten pairs of five things, all pairs occur in the two
sequences (3).

One can now form half-metacyclic functions in very different
ways. The simplest is
\weq{4}{u = x_0x_1 + x_1x_2 + x_2x_3 + x_3x_4 + x_4x_0,}
which by $t$ passes into
\weq{5}{u' = x_0x_2 + x_2x_4 + x_4x_1 + x_1x_3 + x_3x_0}
and $u'$ itself belongs to the half-metacyclic functions.

If
\weq{6}{f(x) = x^5 - ax^4 + bx^3 - cx^2 + dx - e = 0}
is the equation whose roots are $x_0, x_1, x_2, x_3, x_4$, then accordingly
\weq{7}{u + u' = b.}

The function
\weq{8}{y = u - u'}
is likewise half-metacyclic, but the square $y^2$ is full metacyclic
and therefore the root of a resolvent of degree 6,
as \gesperrt{Lagrange} already found.

If $\sqrt{\Delta}$ is the product of the ten root differences $(x_0 - x_1)$,
$(x_0 - x_2) \dots$, so that $\Delta$ is the discriminant of $f(x)$, then
\weq{9}{Y = \dfrac{u - u'}{\sqrt{\Delta}}}

is likewise full metacyclic and the root of an equation of degree 6 $^{1)}$.

In order to form the conjugate functions to $u$, we denote
by $M$ the full metacyclic, by $N$ the half-metacyclic,
by $A$ the alternating group, and decompose $A$ into the subsidiary groups:
\weq{10}{\begin{aligned}A &= N + N(1,2)(3,4) + Nt(0,1) + Nt(0,2)\\ &+ Nt(0,3) + Nt(0,4).\end{aligned}}

These subsidiary groups are in fact all different from one another.
For if, e.\,g., $N = N(1,2)(3,4)$, then $(1,2)(3,4)$ would have to
occur in $N$, and hence
\[ (1,2)(3,4)t = (1,2)(3,4)(1,2,4,3) = (1,4) \]
would have to occur in $M$, and this is not possible, because in $M$ no
permutation appears that leaves two figures unchanged; and in

\par\noindent\rule{4cm}{0.4pt}\par
\noindent $^{1)}$ \gesperrt{Jacobi}, ``observatiunculae ad theoriam aequationum pertinentes'',
Crelle's Journ. Bd. 13. \gesperrt{Jacobi's Werke}, Bd. 3. \gesperrt{Cayley}, philos. transactions
 1861, Collected math. papers, Vol. IV, p. 309.
%CONT>>


% ===== scan leaf 648, printed 623 =====
% !PRINTED 623
\weberpage{623}

%<<CONT
in a similar way one shows that no two others of these
subsidiary groups are equal to one another.

According to this, the values of $u$ conjugate within $A$ are
obtained by applying the permutations

\[
(1,\,2)(3,\,4), \quad t(0,1) = (0,1,2,4,3), \quad t(0,2) = (0,2,4,3,1)
\]
\[
t(0,3) = (0,3,1,2,4), \quad t(0,4) = (0,4,3,1,2)\text{:}
\]

\weq{11}{
\begin{aligned}
u_1 &= x_0x_1 + x_1x_2 + x_2x_3 + x_3x_4 + x_4x_0, \\
u_2 &= x_0x_2 + x_2x_1 + x_1x_4 + x_4x_3 + x_3x_0, \\
u_3 &= x_1x_2 + x_2x_4 + x_4x_0 + x_0x_3 + x_3x_1, \\
u_4 &= x_2x_0 + x_0x_4 + x_4x_1 + x_1x_3 + x_3x_2, \\
u_5 &= x_3x_2 + x_2x_4 + x_4x_1 + x_1x_0 + x_0x_3, \\
u_6 &= x_4x_2 + x_2x_0 + x_0x_1 + x_1x_3 + x_3x_4,
\end{aligned}
}

and the corresponding $u_1', u_2', u_3', u_4', u_5', u_6'$ are found by
applying the permutation $t$ to the $u_1$, and then proceeding with $u_1'$
just as was done with $u_1$, or also simply by uniting the pairs
missing from each $u$ into the corresponding $u'$,
so that for all these $u$ the condition (7) is satisfied.

\weq{12}{
\begin{aligned}
u_1' &= x_0x_2 + x_0x_3 + x_1x_3 + x_1x_4 + x_2x_4, \\
u_2' &= x_0x_1 + x_0x_4 + x_1x_3 + x_2x_3 + x_2x_4, \\
u_3' &= x_0x_1 + x_0x_2 + x_2x_3 + x_3x_4 + x_1x_4, \\
u_4' &= x_0x_1 + x_0x_3 + x_1x_2 + x_2x_4 + x_3x_4, \\
u_5' &= x_0x_2 + x_0x_4 + x_1x_2 + x_1x_3 + x_3x_4, \\
u_6' &= x_0x_3 + x_0x_4 + x_1x_2 + x_1x_4 + x_2x_3.
\end{aligned}
}

The six quantities (8): $y_1, y_2, y_3, y_4, y_5, y_6$ are the roots of an
equation of the sixth degree, whose coefficients depend rationally on $a, b, c, d, e$
and $\sqrt{\varDelta}$. Since the $y$ all change their sign when
$\sqrt{\varDelta}$ changes its sign, this equation has the form

\weq{12}{
y^6 + a_2y^4 + a_4y^2 + a_6 - \sqrt{\varDelta}(a_1y^5 + a_3y^3 + a_5y) = 0, % [sic] duplicate equation number (12); printed page repeats (12) rather than continuing to (13)
}

in which the $a$ are rational functions of the coefficients of $f(x)$.
But they must also be integral functions of these coefficients;
for $a_2, a_4, \alpha_6, a_1\sqrt{\varDelta}, a_3\sqrt{\varDelta}, a_5\sqrt{\varDelta}$ are integral functions of the % [sic] printed as $\alpha_6$ (lowercase alpha) rather than $a_6$
$x$, which can indeed be regarded here as independent variables, and
therefore the three last must be divisible by the difference product $\sqrt{\varDelta}$,
and then $a_2, a_4, a_6, a_1, a_3, a_5$ must prove to be integral
symmetric functions of the $x$. If one determines the
degrees with respect to the $x$, then, since the $y$ are of the second
degree, there result for
%CONT>>


% ===== scan leaf 649, printed 624 =====
% !PRINTED 624
\weberpage{624}

\[
\begin{array}{lcccccc}
 & a_1\sqrt{\Delta}, & a_2, & a_3\sqrt{\Delta}, & a_4, & a_5\sqrt{\Delta}, & a_6\\
\text{the degrees} & 2 & 4 & 6 & 8 & 10 & 12.
\end{array}
\]
But $\sqrt{\Delta}$ is of degree 10 in the $x$, and consequently $a_1=0$,
$a_3=0$ must hold, and $a_5$ must be a number.

One can compute the coefficients $a_2$, $a_4$, $a_5$, $a_6$ by actual
formation of the expressions according to the rules for the representation of the symmetric
functions, which has been done by \gesperrt{Cayley}. But the calculation and the formulae are very
lengthy.

We shall content ourselves here with forming the resolvent for
a special case. The value found in doing so
for the number $a_5$ is then of course generally valid.

Let the equation of degree 5 have the \gesperrt{Bring-Jerrard}
form\footnote{\gesperrt{Runge}, Acta mathematica, Bd.\,7.}
\weq{13}{x^5+\alpha x+\beta=0.}
Then $a_2$, $a_4$, $a_6$ are integral rational functions of $\alpha$ and $\beta$.
Of the coefficients of the general equation (6), $d$ is of
the fourth, $e$ of the fifth order in the $x$. From this it follows that
in $a_2$, $a_4$, $a_6$ the coefficient $e$ can occur only multiplied by one of the coefficients
$a$, $b$, $c$, and hence that, if $a$, $b$, $c$
are set to zero, $e$ must drop out of these expressions.
For equation (13), therefore, $a_2$, $a_4$, $a_6$ cannot depend
on $\beta$, and, having regard to the degree, one obtains, if
$m$, $m_1$, $m_2$, $m_3$ denote numbers,
\weq{14}{a_2=m_1\alpha,\quad a_4=m_2\alpha^2,\quad a_6=m_3\alpha^3,\quad a_5=m.}

The numbers $m$, $m_1$, $m_2$, $m_3$ can be determined from a special
assumption. If we set $\beta=0$, then
\weq{15}{\begin{aligned}
x_0&=0,\quad x_1=\sqrt[4]{-\alpha},\quad x_2=i\sqrt[4]{-\alpha},\\
&\hphantom{{}=0,\quad{}} x_3=-\sqrt[4]{-\alpha},\quad x_4=-i\sqrt[4]{-\alpha},
\end{aligned}}
and from this there results
\weq{16}{\begin{aligned}
\sqrt{\Delta} &= (x_0-x_1)(x_0-x_2)(x_0-x_3)(x_0-x_4)(x_1-x_2)\\
&\quad (x_1-x_3)(x_1-x_4)(x_2-x_3)(x_2-x_4)(x_3-x_4)\\
&= 16i\sqrt{-\alpha}^5 = -16\sqrt{\alpha}^5\\
\Delta &= 256\alpha^5.
\end{aligned}}


% ===== scan leaf 650, printed 625 =====
% !PRINTED 625
\weberpage{625}

Since here $b=0$, by (7) and (8) we have $y=2u$, and
from (11) and (15) the values of the $y$ result:
\[
\begin{aligned}
y_1&=y_2=y_3=y_6=-2\sqrt{\alpha}\\
y_4&=(4-2i)\sqrt{\alpha},\quad y_5=(4+2i)\sqrt{\alpha}.
\end{aligned}
\]
From this one obtains for $\beta=0$ the following equation,
identical in $y$:
\weq{17}{\begin{aligned}
&y^6+a_2y^4+a_4y^2+a_6-a_5\sqrt{\Delta}\,y\\
&=(y+2\sqrt{\alpha})^4(y^2-8y\sqrt{\alpha}+20\alpha)\\
&=y^6-20\alpha y^4+240\alpha^2y^2+512\sqrt{\alpha}^5y+320\alpha^3.
\end{aligned}}
But from this one finds the general resolvent for the
equation (13), if according to (16) one puts $\sqrt{\Delta}$ for $-16\sqrt{\alpha}^5$, thus
\weq{18}{y^6-20\alpha y^4+240\alpha^2y^2-32\sqrt{\Delta}\,y+320\alpha^3=0.}

The discriminant $\Delta$ here has, as follows from formula (3)
in \S.\,74, the expression
\weq{19}{\Delta=5^5\beta^4+2^8\alpha^5.}

Simpler still is the equation for $u$, namely
\weq{20}{u^6-5\alpha u^4+15\alpha^2u^2-\sqrt{\Delta}\,u+5\alpha^3=0,}
and if one puts $u^2=v$, there results for $v$ the rational
equation of degree 6
\weq{21}{(v^3-5\alpha v^2+15\alpha^2v+5\alpha^3)^2=\Delta v.}

Another form of this equation is obtained from (17),
if one multiplies the second of the three expressions by the one
obtained from it by changing $\sqrt{\alpha}$ into $-\sqrt{\alpha}$.
In this way one obtains
\[
(v-\alpha)^4(v^2+6\alpha v+25\alpha^2)=0,
\]
and (21) must pass over into this when $\beta=0$ is set.
But since in (21) the part depending on $\beta$ is determined
by (19), the sought form of the resolvent follows:
\weq{22}{(v-\alpha)^4(v^2+6\alpha v+25\alpha^2)=5^5\beta^4v.}

One may further raise the question whether the function $v$
introduced here is always really metacyclic, whether it might not
perhaps, for special equations, remain unchanged under still
other permutations.

If this case should occur, then the equation
(21) or (22) would have to have equal roots, or (20) equal or
opposite roots.


% ===== scan leaf 651, printed 626 =====
% !PRINTED 626
\weberpage{626}

The value $u=0$ or $v=0$ occurs only when
$\alpha=0$. In that case we have in fact in (13) the well-known
metacyclic equation $x^5+\beta=0$; this case, then, we
now set aside. The equation (20) can only have
two opposite roots when $\Delta=0$, that is,
when the equation (13) has equal roots. This too is
to be excluded, since we have assumed (13) to be irreducible.

There remains, then, only the question whether (20) can have
equal roots. In that case, together with (20), the derived
equation
\[
6u^5-20\alpha u^3+30\alpha^2u-\sqrt{\Delta}=0
\]
would have to be satisfied, and if we eliminate $\sqrt{\Delta}$ from this
and from (20),
\[
5u^6-15\alpha u^4+15\alpha^2u^2-5\alpha^3=5(u^2-\alpha)^3=0.
\]
We should therefore have to have $v=\alpha$, and hence by (22) $\beta=0$.
But then the equation of degree 5 would again be reducible, since
it has the factor $x$.

If one wishes, then, to test a given equation of the form (13)
with regard to its solvability, one has first to establish the
irreducibility, and then to investigate whether (21)
or (22) has a rational root.

If $\alpha$ and $\beta$ are integers, then a rational $v$ must also
be an integer, which is to be sought among the factors
of $25\alpha^6$.

If, e.\,g., $\alpha=5$, $\beta=5t$, then, if $t$ is an integer not
divisible by 5, $x^5+\alpha x+\beta$ is irreducible by \S.\,178, 3.
For $v$ one has to substitute powers of 5, with positive or
negative sign, but for no such number
can (22) be satisfied, because on the left-hand side the power of
5 does not become high enough.

\gesperrt{Thus no equation of the form}
$x^5+5x+5t=0$ \gesperrt{is metacyclic}.

If one wishes to determine metacyclic equations, one sets
in (22)
\[
\beta=\alpha\mu,\quad v=\alpha\lambda,
\]
whereby one obtains:
\weq{23}{\begin{aligned}
\alpha&=\dfrac{5^5\mu^4\lambda}{(\lambda-1)^4(\lambda^2+6\lambda+25)}\\
\beta&=\dfrac{5^5\mu^5\lambda}{(\lambda-1)^4(\lambda^2+6\lambda+25)},
\end{aligned}}
%CONT>>


% ===== scan leaf 652, printed 627 =====
% !PRINTED 627
\weberpage{627}
\pdfbookmark[2]{§ 182. The group of the resolvent}{bm203}

%<<CONT
and if one here puts rational quantities for $\lambda$ and $\mu$ (from
an arbitrary domain of rationality, e.g. also from the field
of rational functions of $\lambda$ and $\mu$), then the equation
of degree 5
\weq{24}{x^5+\alpha x+\beta=0}
is algebraically solvable, since indeed the case that this equation
becomes reducible also leads to algebraically solvable equations.

Conversely we can say that no irreducible equation
of degree 5 of the form (24) is algebraically solvable in
which the coefficients $\alpha$, $\beta$ are not representable in the form (23).

If we set, e.g., $\lambda=-1$, $\mu=1$, we obtain $64\alpha=-5^4$,
$64\beta=-5^4$. The equation (24) is simplified, if we introduce for
$x$ a new unknown $\xi$, by the equation $x\xi=5$.
We thus obtain the equation
\[
\xi^5+5\xi^4-5\cdot 64=0
\]
as an example of an \gesperrt{algebraically solvable} equation which, moreover,
is \gesperrt{irreducible} by the theorem mentioned above.

\begin{center}
\S.\,182.

\gesperrt{The Group of the Resolvent.}
\end{center}

The theorems that we have derived in the preceding section
admit a remarkable conclusion about the equations
of degree 6.

We now again take the $x_0$, $x_1$, $x_2$, $x_3$, $x_4$ of the preceding
section as independent variables. Then the six
quantities
\weq{1}{\begin{aligned}
v_1&=(u_1-u_1')^2, & v_2&=(u_2-u_2')^2, & v_3&=(u_3-u_3')^2\\
v_4&=(u_4-u_4')^2, & v_5&=(u_5-u_5')^2, & v_6&=(u_6-u_6')^2
\end{aligned}}
are the roots of an equation of degree 6, $F(v)=0$, whose
coefficients are rationally expressible in terms of the symmetric fundamental functions
$a$, $b$, $c$, $d$, $e$ of the $x$, and which is irreducible in the field
of the rational functions of these quantities. This
equation is a resolvent of the equation of degree 5, $f(x)=0$,
whose roots are the $x$.

If $M$ is the full metacyclic group of the $x$, then the
groups $\pi^{-1}M\pi$ conjugate to $M$ have no other common
%CONT>>


% ===== scan leaf 653, printed 628 =====
% !PRINTED 628
\weberpage{628}
%<<CONT
% [sic] The source defines the resolvent as F(v)=0 but subsequently prints F(z) here; notation retained source-faithfully.
common divisor than the identical permutation. For every common divisor of all these groups would have to be a normal divisor of the symmetric, and consequently also of the alternating group; but since it cannot be the alternating group itself, it must by \S.\,177 reduce to the unit group.

The resolvent $F(v) = 0$ is therefore, according to the designation introduced in \S.\,156, a \gesperrt{total resolvent} of $f(x) = 0$, and its group must be of the same degree as the group of $f(x) = 0$, whose degree is $1 . 2 . 3 . 4 . 5 = 120$. Since $F(z)$ is also irreducible, the permutation group among the indices of the six quantities $v$ is transitive. It is of degree 120, while the degree of the symmetric group of the permutations of six figures is $6 . 120$. With this we have the theorem:

\gesperrt{The symmetric permutation group of six figures has a transitive divisor of index 6.}

In order to find this remarkable group, which we shall denote by $C$, we have only to investigate the influence which all the 120 permutations of the $x$ exert on the $v$, or on the quantities (11) of the preceding section.

Now we have seen earlier (\S.\,153, 2.) that the transpositions $(0, 1)$, $(0, 2)$, $(0, 3)$, $(0, 4)$ generate, by repeated composition, the whole symmetric group of the permutations of five figures. It therefore suffices to determine the influence of these four transpositions on the indices of the $v$. By the application of a transposition $(0, 1)$ each of the expressions (11), \S.\,181 passes into one of the expressions (12) and conversely, and if one determines this change for the four transpositions named, one obtains the following four generating permutations $\pi_1, \pi_2, \pi_3, \pi_4$ of the group $C$

\[
\begin{aligned}
(0, 1)&: \quad \pi_1 = (1, 3) (2, 5) (4, 6) \\
(0, 2)&: \quad \pi_2 = (1, 4) (2, 3) (5, 6) \\
(0, 3)&: \quad \pi_3 = (1, 5) (2, 6) (3, 4) \\
(0, 4)&: \quad \pi_4 = (1, 6) (2, 4) (3, 5),
\end{aligned}
\]

where the transpositions $(1, 3) \ldots$ occurring in the $\pi$ refer to the indices of the $u$ or of the $v$.

Among the 120 permutations of the group $C$ there occur, among others, also
%CONT>>


% ===== scan leaf 654, printed 629 =====
% !PRINTED 629
\weberpage{629}
%<<CONT
\[
\begin{aligned}
\pi_1 \pi_2 &= (1, 2, 6) (3, 4, 5) \\
\pi_1 \pi_2 \pi_3 &= (1, 6, 5, 4) \\
\pi_1 \pi_2 \pi_3 \pi_4 &= (2, 4, 6, 3, 5) \\
(\pi_1 \pi_2 \pi_3 \pi_4)^3 &= (2, 3, 4, 5, 6).
\end{aligned}
\]
According to this, a function of the six quantities $v$ belonging to the group $C$ can easily be formed.

A function such as
\[
v_1 v_3 + v_2 v_4 + v_5 v_6,
\]
remains unchanged by $\pi_1$ and by $\pi_1 \pi_2$; but if we apply to it the powers of the cyclic permutation $(2, 3, 4, 5, 6)$, five functions arise from it, which we shall denote as follows:
\weq{2}{
\begin{aligned}
% [sic] Print has a prime after $v_4$ in this row; the parallel combinatorial pattern requires unprimed $v_4$.
w_0 &= v_1 v_2 + v_4' v_5 + v_3 v_6 \\
w_1 &= v_1 v_4 + v_2 v_6 + v_3 v_5 \\
w_2 &= v_1 v_6 + v_2 v_5 + v_3 v_4 \\
w_3 &= v_1 v_3 + v_2 v_4 + v_5 v_6 \\
w_4 &= v_1 v_5 + v_2 v_3 + v_4 v_6.
\end{aligned}
}

If we apply the permutations $\pi_1, \pi_2, \pi_3, \pi_4$ to the functions (2), these functions merely pass into one another, and indeed in the following way:
\[
\pi_1 = (w_0, w_1), \quad \pi_2 = (w_0, w_2), \quad \pi_3 = (w_0, w_3), \quad \pi_4 = (w_0, w_4),
\]
i.e. the $\pi_1, \pi_2, \pi_3, \pi_4$ correspond to the transpositions $(0, 1)$, $(0, 2)$, $(0, 3)$, $(0, 4)$ among the indices of the $w$, and thus by \S.\,153, 2. the symmetric group of the indices of the $w$ corresponds to the whole group $C$. A symmetric function of the five quantities $w$ which is not at the same time a symmetric function of the $v$, such as the sum of the $w$, is therefore a function belonging to $C$.

One can, for example, take for this the sum of the squares
\[
W = w_0^2 + w_1^2 + w_2^2 + w_3^2 + w_4^2,
\]
or a function
\[
(\lambda - w_0) (\lambda - w_1) (\lambda - w_2) (\lambda - w_3) (\lambda - w_4)
\]
for an arbitrary rational $\lambda$.

Such a quantity $W$ is the root of an irreducible equation of degree $6$, whose coefficients are symmetric functions of the six quantities $v$.

\begin{center}
\rule{3cm}{0.4pt}
\end{center}


% ===== scan leaf 655, printed 630 =====
\weberpage{630}% page number not printed on this page
\pdfbookmark[1]{Eighteenth Section. Roots of metacyclic equations.}{bm204}
\pdfbookmark[2]{§ 183. Statement of the problem. Lemma}{bm205}
% !PRINTED none
% Abschnitt-opening page: no running head / page number visible on this page.
\begin{center}
\gesperrt{Eighteenth Section.}

\gesperrt{Roots of Metacyclic Equations.}

\rule{0.3\linewidth}{0.4pt}

\S.\,183.

\gesperrt{Statement of the Problem. Lemma.}
\end{center}

In the preceding section we have found general criteria by which one can decide whether a given equation of prime degree is metacyclic or not. But with this the subject is by no means exhausted. Rather, according to the form of the statement of the problem as it derives from \gesperrt{Abel}, the task is to give a procedure by which one can find all metacyclic equations. In this form the problem, on which \gesperrt{Abel} left behind only brief indications without proofs, was taken up by \gesperrt{Kronecker}, and completely solved in such a way that at first not the equations themselves, but their roots, are found\footnote{\gesperrt{Abel}, Sur la résolution algébrique des équations (Oeuvres complètes, ed. Sylow, tome II, p. 217), letter to \gesperrt{Crelle} of 14 March 1826 (Oeuvres tome II, p. 266). --- \gesperrt{Kronecker}, Ueber die algebraisch auflösbaren Gleichungen. Monatsberichte der Berl. Akademie 1853, 1856. --- \gesperrt{H. Weber}, Ueber algebraisch auflösbare Gleichungen von Primzahlgrad. Sitzungsberichte der Gesellschaft zur Beförderung der Naturwissenschaften zu Marburg 1892. The following presentation is simplified as against the one given there, and corrected in one point.}. Thus \gesperrt{Kronecker} has given, for every prime $n$, an expression which is derived from a given field $\Omega$ by repeated extraction of roots, and which has the double
%CONT>>


% ===== scan leaf 656, printed 631 =====
% !PRINTED 631
\weberpage{631}
%<<CONT
property that every root of an irreducible metacyclic equation of degree $n$ in $\Omega$ is contained in this expression, and that conversely every such expression also satisfies an irreducible equation of degree $n$ in $\Omega$.

The foundation for this investigation is furnished by the resolvents of \gesperrt{Lagrange}.

Now let $n$ be a prime, let $x_0, x_1, x_2 \ldots x_{n-1}$ be a system of independent variables, and let the notation be so chosen that $x_n = x_0$, $x_{n+1} = x_1 \ldots$ Finally, let $\varepsilon$ be an $n$th root of unity. We then set
\weq{1}{
(\varepsilon, x) = x_0 + \varepsilon x_1 + \varepsilon^2 x_2 + \cdots + \varepsilon^{n-1} x_{n-1} = \sum_{0, n-1}^{h} \varepsilon^h x_h.
}
These expressions we have already come to know in the fourteenth section under the name of the \gesperrt{resolvents of Lagrange}, and we have there (\S.\,164) investigated their behaviour with respect to the \gesperrt{cyclic permutations}.

It is now a matter of determining the influence of linear permutations, of the kind considered in \S.\,180, on these functions.

We first send ahead a lemma, which we really need only reproduce from what has gone before. The \gesperrt{imaginary} $n$th roots of unity $\varepsilon$ are the roots of an equation $X = 0$, if
\weq{2}{
X = \xi^{n-1} + \xi^{n-2} + \cdots + \xi + 1
}
is set, and $\xi$ denotes a variable.

The function $X$ is, as we have seen earlier (\S.\,134), irreducible in the field $R$ of the rational numbers.

If we take as basis any field $\Omega$ in which $X$ is irreducible, and derive from it the field $\Omega(\varepsilon)$, which consists of all rational functions of $\varepsilon$ with coefficients in $\Omega$, then it follows by \S.\,142 that every quantity of this field can be put in one and only one way into the form
\weq{3}{
c_0 + c_1 \varepsilon + c_2 \varepsilon^2 + \cdots + c_{n-2} \varepsilon^{n-2}
}
in which the $c$ are quantities in $\Omega$.

For brevity we shall call this form of the quantities in $\Omega(\varepsilon)$ the \gesperrt{normal form}, and shall accordingly state the theorem:

1. \gesperrt{If $\Omega$ is a field in which the cyclotomic equation $X = 0$ is irreducible, then every}
%CONT>>


% ===== scan leaf 657, printed 632 =====
% !PRINTED 632
\weberpage{632}
%<<CONT
\gesperrt{quantity of the field $\Omega(\varepsilon)$ can be brought in only one way into the normal form
\[
\Phi(\varepsilon) = c_0 + c_1\varepsilon + c_2\varepsilon^2 + \cdots + c_{n-2}\varepsilon^{n-2}
\]
in which the coefficients $c$ are contained in $\Omega$.}

In this theorem we can take for $\Omega$ the field that arises from $R$ by the adjunction of $n$ independent variables $x_0, x_1, x_2 \ldots x_{n-1}$; for if $X$ were reducible in this field, so that
\[
X = X_1 X_2,
\]
where $X_1, X_2$ are integral functions of $\varepsilon$ and rational functions of the $x$, then one could put for the variables $x_0, x_1 \ldots x_{n-1}$ such rational numbers that $X_1, X_2$, without changing their degree with respect to $\varepsilon$, pass into functions in $R$ (\S.\,143), and this contradicts the irreducibility of $X$ in $R$.

This field $\Omega$ has various divisors in which $X$ is likewise irreducible, and which can therefore all be taken for $\Omega$ in Theorem 1. Such divisors are obtained if one fixes any \gesperrt{permutation group} $P$ of the $n$ figures $0, 1, 2 \ldots n-1$ and considers the totality of all rational functions of $x_0, x_1 \ldots x_{n-1}$ in $R$ that \gesperrt{admit} the permutations of this group.

The totality of all these functions evidently forms a field, and thus to every permutation group of the $n$ figures one obtains a definite associated field $\Omega$. The quantities of such a field can be represented rationally by the \gesperrt{symmetric functions of the $x$ and by a function belonging to the group $P$}.

We call the field $\Omega$ so determined \gesperrt{the field belonging to the group $P$}. Every such field can take the place of $\Omega$ in 1. (Cf. \S.\,152.)

We can therefore state the second theorem:

\gesperrt{2. If a rational function $\Phi(x_0, x_1 \ldots x_{n-1}, \varepsilon)$ of the field $\Omega(\varepsilon)$ remains unchanged when the permutations of a group $P$ are applied to the indices of the $x$, then the coefficients $c_0, c_1, c_2 \ldots c_{n-2}$ of the normal form of $\Phi$ also admit the permutations of $P$.}

If $\Phi(\varepsilon)$ is any element of the field $\Omega(\varepsilon)$, one obtains the conjugate quantities by putting for $\varepsilon$ any other
%CONT>>


% ===== scan leaf 658, printed 633 =====
% !PRINTED 633
\weberpage{633}
\pdfbookmark[2]{§ 184. Theorems on the resolvents}{bm206}
%<<CONT
root of $X = 0$, that is, carries out the substitution $(\varepsilon, \varepsilon^h)$ for $h = 1, 2 \ldots n-1$. All these fields are again contained in $\Omega(\varepsilon)$; they are therefore identical with one another, i.e. $\Omega(\varepsilon)$ is a \gesperrt{normal field}. If a quantity $\Phi(\varepsilon)$ has the property of remaining unchanged when $\varepsilon$ is replaced by all the $\varepsilon^h$, then it is itself contained in $\Omega$. This already follows from the general theorems of \S.\,144. But we see its correctness at once if, under the hypothesis $\Phi(\varepsilon^h) = \Phi(\varepsilon)$, we set
\[
\Phi(\varepsilon) = \frac{1}{n-1}\sum_{1,\,n-1}^{h}\Phi(\varepsilon^h) = c_0 - \frac{c_1 + c_2 + \cdots + c_{n-2}}{n-1}
\]
from which it is then to be concluded, by 2., that $c_1, c_2 \ldots c_{n-2}$ are equal to zero. Hence:

\gesperrt{3. If a quantity in $\Omega(\varepsilon)$ admits all the substitutions $(\varepsilon, \varepsilon^h)$, then it is contained in $\Omega$.}

\begin{center}
\S.\,184.

\gesperrt{Theorems on the Resolvents.}
\end{center}

We now consider the \gesperrt{Lagrange} resolvents $(\varepsilon, x)$ under the hypothesis that $\varepsilon$ is an \gesperrt{imaginary} $n$th root of unity and the $x$ are independent variables, and apply to them the theorems of the preceding section.

We first investigate the influence of the linear permutations on $(\varepsilon, x)$; for this it suffices to determine the influence of the two generating substitutions
\weq{1}{s = (h, h+1), \quad t = (h, gh)}
where $h$, taken modulo $n$, runs through the indices of $x$, and $g$ denotes a primitive root of the prime $n$ (\S.\,180).

But if we apply to
\weq{2}{(\varepsilon, x) = x_0 + \varepsilon x_1 + \varepsilon^2 x_2 + \cdots + \varepsilon^{n-1}x_{n-1}}
the substitution $s$, this function goes over into
\[
x_1 + \varepsilon x_2 + \varepsilon^2 x_3 + \cdots + \varepsilon^{n-1}x_0 = \varepsilon^{-1}(\varepsilon, x).
\]
The effect of the substitution $s$, then, is that
\[
% the German sentence closes with "übergeht." after the display; the verb is absorbed into the
% \text and the display copies the German character for character, without a final stop
(\varepsilon, x) \quad \text{goes over into} \quad \varepsilon^{-1}(\varepsilon, x)
\]


% ===== scan leaf 659, printed 634 =====
% !PRINTED 634
\weberpage{634}
We now pass to the substitution $t$, and in
\weq{2}{(\varepsilon, x) = x_0 + \sum_{1,\,n-1}^{h}\varepsilon^h x_h}
put $x_{gh}$ for $x_h$; this function then goes over into
\weq{3}{x_0 + \sum_{1,\,n}^{h}\varepsilon^h x_{gh}.} % [sic] Print has lower limits $1,n$ although the residue system excludes $0$; preserved.
In this sum $h$ may run through an arbitrary system of residues with respect to the modulus $n$, with the exclusion of $0$, and we can therefore replace $h$ by $g^{-1}h$; thereby the sum (3) becomes
\[
x_0 + \sum^{h}\varepsilon^{g^{-1}h}x_h = (\varepsilon^{g^{-1}}, x).
\]
If one applies the permutations $s$ and $t$ repeatedly, there results the influence of $s^\lambda, t^\lambda$, which consists in the permutations
\[
(\varepsilon, x), \varepsilon^{-\lambda}(\varepsilon, x)
\]
\[
(\varepsilon, x), (\varepsilon^{g^{-\lambda}}, x)
\]
and so in particular, for $\lambda = -1$,
\[
(\varepsilon, x), \varepsilon(\varepsilon, x)
\]
\[
(\varepsilon, x), (\varepsilon^{g}, x),
\]
and we obtain the following theorem:

\gesperrt{4. The generating permutations of the metacyclic group, $s, t$, effect the permutations
\[
(s)\colon (\varepsilon, x), \varepsilon^{-1}(\varepsilon, x); \qquad (s^{-1})\colon (\varepsilon, x), \varepsilon(\varepsilon, x)
\]
\[
(t)\colon (\varepsilon, x), (\varepsilon^{g^{-1}}, x); \qquad (t^{-1})\colon (\varepsilon, x), (\varepsilon^{g}, x).
\]}

But from 2. and 4., because $s$ is the generating permutation of the cyclic group, there results the consequence:

\gesperrt{5. If one represents the functions
\weq{4}{(\varepsilon, x)^n = \Phi(\varepsilon), \quad (\varepsilon^\lambda, x)(\varepsilon, x)^{-\lambda} = F(\varepsilon)}
in the normal form [\S.\,183, (3)], then the coefficients are cyclic functions of $x_0, x_1 \ldots x_{n-1}$ (\S.\,163).}

Here $\lambda$ may denote any arbitrary integer that is not divisible by $n$.

From the function $F(\varepsilon)$ there arise, if we take $\lambda = g$ and put for $\varepsilon$ its $n-1$ different values, $n-1$ functions, which all remain unchanged under $s$, and which we shall denote as follows:
%CONT>>


% ===== scan leaf 660, printed 635 =====
% !PRINTED 635
\weberpage{635}
%<<CONT
\weq{5}{\begin{aligned}
(\varepsilon^{g}, x)(\varepsilon, x)^{-g} &= f_0\\
(\varepsilon^{g^2}, x)(\varepsilon^{g}, x)^{-g} &= f_1\\
\cdot\ \cdot\ \cdot\ \cdot\ \cdot\ \cdot\ \cdot\ \cdot\ \cdot\ \cdot & \\
(\varepsilon^{g^{n-1}}, x)(\varepsilon^{g^{n-2}}, x)^{-g} &= f_{n-2}.
\end{aligned}}
so that, if $(\varepsilon^{g}, x)(\varepsilon, x)^{-g} = f(\varepsilon)$ is set, we have
\weq{6}{f_h = f(\varepsilon^{g^h})}
If one applies the permutation $t^{-1}$ to (5), then, as is to be seen from Theorem 4., the functions
\[
f_0, f_1, f_2 \ldots f_{n-2}
\]
undergo a cyclic permutation $(0, 1, 2 \ldots n-2)$.

Moreover, so long as the $x$ are variable, the functions $f_\nu$ are, as one sees from (5), all \gesperrt{different from one another} since they are decomposed into different linear factors.

If, then, we form a cyclic function of these $n-1$ quantities
\weq{7}{C(f_0, f_1, f_2 \ldots f_{n-2}),}
this remains unchanged both by $s$ and by $t^{-1}$, hence by $t$ and the whole metacyclic group, and we obtain the theorem:

\gesperrt{6. If a cyclic function of the quantities $f_0, f_1 \ldots f_{n-1}$ is represented in the normal form, then the coefficients are metacyclic functions of the variables $x_0, x_1 \ldots x_{n-1}$.} % [sic] Print says $f_{n-1}$ although (5) defines the sequence only through $f_{n-2}$; preserved.

Here a metacyclic function is to be understood as any function that remains unchanged by the permutations of the linear group.

But the functions $f_0, f_1 \ldots f_{n-2}$ also pass cyclically into one another if one leaves the $x$ unchanged and makes the substitution $(\varepsilon, \varepsilon^g)$ for $\varepsilon$. If, then, the cyclic function $C$ of the $f$ does not contain $\varepsilon$ in its coefficients, then $C$ remains unchanged by all the substitutions $(\varepsilon, \varepsilon^h)$, and we obtain from 3. the theorem:

\gesperrt{7. If $C(f_0, f_1 \ldots f_{n-2})$ is a cyclic function of the quantities $f_0, f_1 \ldots f_{n-2}$ with rational coefficients, then $C$ is a metacyclic function of the variables $x_0, x_1 \ldots x_{n-1}$ with rational coefficients.}

If we raise the functions $f_0, f_1 \ldots f_{n-2}$, as they are given by (5), in order to the powers $g^{n-2}, g^{n-3} \ldots 1$
%CONT>>


% ===== scan leaf 661, printed 636 =====
% !PRINTED 636
\weberpage{636}
%<<CONT
and then multiply everything, and if we further observe that
$\varepsilon^{g^{n-1}} = \varepsilon$, then in the product of the
left-hand sides of the equations (5) all the resolvents except
$(\varepsilon, x)$ cancel out, and it follows that

\weq{8}{(\varepsilon, x)^{1-g^{n-1}} = f_0^{g^{n-2}} f_1^{g^{n-3}} \cdots f_{n-2}.}

If we now denote by $\lambda$ a still undetermined integer
exponent, we derive from (8) the following relation:

\weq{9}{(\varepsilon^\lambda, x)^n = \left[ (\varepsilon, x)^{\frac{g^{n-1}-1}{n}} (\varepsilon^\lambda, x) \right]^n f_0^{g^{n-2}} f_1^{g^{n-3}} \cdots f_{n-2}.}

If we understand by $g_1$ any fixed primitive root of $n$, we
can set $g = g_1 + l n$, where $l$ denotes an arbitrary integer,
since the primitive roots are defined only up to multiples of $n$.
Then

\[
g^{n-1} \equiv g_1^{n-1} + n(n-1) l g_1^{n-2} \pmod{n^2},
\]

from which

\[
\frac{g^{n-1}-1}{n} \equiv \frac{g_1^{n-1}-1}{n} - l g^{n-2} \pmod{n},
\]

and $l$ can be so determined that

\[
\frac{g^{n-1}-1}{n} \equiv 1 \pmod{n}
\]

holds.

In (9) we now set, under this assumption,

\weq{10}{\lambda = \frac{g^{n-1}-1}{n} \equiv 1 \pmod{n},}

and, as in (4),

\weq{11}{F(\varepsilon) = (\varepsilon^\lambda, x)(\varepsilon, x)^{-\lambda} = (\varepsilon, x)^{nk},}

where $\varepsilon^\lambda = \varepsilon$, and $1 - \lambda$, which is
divisible by $n$, is set $= nk$. From this we obtain

\weq{12}{(\varepsilon, x)^n = [F(\varepsilon)]^n f_0^{g^{n-2}} f_1^{g^{n-3}} \cdots f_{n-2}.}

These formulae can be generalised if one observes that the
functions $f_0, f_1 \ldots f_{n-2}$ undergo a cyclic permutation
when the substitution $(\varepsilon, \varepsilon^g)$ is carried out.
If we therefore stipulate that $f_h = f_k$ is to hold when $h \equiv k
\pmod{n-1}$, and set

\weq{13}{F(\varepsilon^{g^\nu}) = F_\nu = (\varepsilon^{g^\nu}, x)^{nk},}

then the substitution $(\varepsilon, \varepsilon^{g^\nu})$ in (12) yields

\weq{14}{(\varepsilon^{g^\nu}, x)^n = F_\nu^n f_\nu^{g^{n-2}} f_{\nu+1}^{g^{n-3}} \cdots f_{\nu+n-2}.}


% ===== scan leaf 662, printed 637 =====
% !PRINTED 637
\weberpage{637}
These formulae hold for every $\nu$, if we further stipulate with regard
to the $F_\nu$ that $F_h = F_k$ shall hold whenever $h \equiv k
\pmod{n-1}$. It then follows from 4. that the functions $F_0,
F_1 \ldots F_{n-2}$ remain unchanged by the substitution $s$, and
undergo a cyclic permutation by $t^{-1}$.

If we wish here to reduce the exponents to their least positive residues
with respect to the modulus $n$, we set

\weq{15}{g^\nu = n q_\nu + r_\nu, \quad 0 < r_\nu < n, \quad r_0 = 1,}

and the numbers $r_\nu$ coincide, in some order, with the numbers
$1, 2 \ldots n-1$, so that $r_0$ is always $= 1$.

Then by (5) we have

\weq{16}{f_{\nu-1} = (\varepsilon^{r_\nu}, x)(\varepsilon^{r_{\nu-1}}, x)^{-g},}

and we further set

\weq{17}{\Phi_\nu = F_\nu f_\nu^{q_{n-2}} f_{\nu+1}^{q_{n-3}} \cdots f_{\nu+n-3},}

and obtain from (14)

\weq{18}{(\varepsilon^{r_\nu}, x)^n = \Phi_\nu^n f_\nu^{r_{n-2}} f_{\nu+1}^{r_{n-3}} \cdots f_{\nu+n-2}^{r_0}.}

The exponents $r_\nu$ remain unchanged when $g$ is replaced by $g - l n$,
and are therefore independent of the condition (10) and \gesperrt{can be derived
from an arbitrary primitive root $g$}.

The characteristic properties of the functions $f_\nu, F_\nu, \Phi_\nu$
occurring in these formulae, whose totality we shall denote by $\omega_\nu$,
are, to repeat them once more, the following:

$\alpha)$ \gesperrt{By the cyclic permutation $s = (h, h+1)$ among the
indices of $x$ the functions $\omega_\nu$ are not changed.}

$\beta)$ \gesperrt{By the linear permutation $t^{-1} = (h, g^{-1} h)$ among the
indices of the $x$ the cyclic permutation $(0, 1, 2 \ldots n-2)$
is produced among the indices of the $\omega$.}

$\gamma)$ \gesperrt{By the substitution $\sigma = (\varepsilon, \varepsilon^g)$
the same cyclic permutation $(0, 1, 2 \ldots n-2)$
is effected among the indices of $\omega$.}

$\delta)$ \gesperrt{The cyclic functions of $\omega_\nu$ are metacyclic
functions of the $x$ and free of $\varepsilon$.}


% ===== scan leaf 663, printed 638 =====
% !PRINTED 638
\weberpage{638}
\pdfbookmark[2]{§ 185. Roots of metacyclic equations}{bm207}
By the theorem of \gesperrt{Lagrange} (§.155) every function
$\Theta_\nu$ possessing the four properties $\alpha)$, $\beta)$,
$\gamma)$, $\delta)$ can be expressed rationally, in the field $\Omega$
of the metacyclic functions of $x$, by one of them, $\omega_\nu$,
provided only that the $\omega_0, \omega_1 \ldots \omega_{n-2}$ are
different from one another; one may therefore choose for these functions
the $f_\nu$.

For if $u$ denotes a variable, and if

\weq{19}{(u - \omega_0)(u - \omega_1) \cdots (u - \omega_{n-2}) = \varphi(u),}

then the function $\varphi(u)$ admits the permutations $s, t$ and also
the substitution $\sigma = (\varepsilon, \varepsilon^g)$; and hence, by
Theorem 3. of §.183, it is an integral function of $u$ with coefficients
metacyclic in regard to $x$, and independent of $\varepsilon$.

If now $\Theta_0, \Theta_1 \ldots \Theta_{n-2}$ is a system of
functions possessing the properties $\alpha) \beta) \gamma) \delta)$,
then the sum

\weq{20}{\sum_{0,\, n-2}^{\nu} \frac{\Theta_\nu\, \varphi(u)}{u - \omega_\nu} = \chi(u)}

is an integral function of degree $(n-2)$ of $u$, which likewise admits
the substitutions $s, t, \sigma$ and which therefore has coefficients in
$\Omega$. If one sets

\weq{21}{\frac{\chi(u)}{\varphi'(u)} = \Theta(u)}

and then sets $u = \omega_\nu$ in (20), it follows that

\[
\Theta_\nu = \Theta(\omega_\nu),
\]

where $\Theta$ denotes a rational function whose coefficients are
metacyclic functions of the $x$, and no longer contain $\varepsilon$.

\begin{center}
§. 185.

\gesperrt{Roots of metacyclic equations.}
\end{center}

In the formulae of the preceding section, the roots $\xi_0, \xi_1
\ldots \xi_{n-1}$ of a metacyclic equation of degree $n$, irreducible
in some field $\mathfrak{K}$, are now to be introduced in place of the
variables $x_0, x_1 \ldots x_{n-1}$, and indeed in such an order that
the metacyclic functions of the $\xi$ are rational (in $\mathfrak{K}$),
which by the theorems of the preceding section is always possible. But
in doing so we first make two further restrictive as-
%CONT>>


% ===== scan leaf 664, printed 639 =====
% !PRINTED 639
\weberpage{639}
%<<CONT
sumptions, from which we shall afterwards free the result again.
These assumptions are:

\gesperrt{1. that by this substitution of the $\xi$ for the $x$ none of
the resolvents $(\varepsilon, \xi)$, in which $\varepsilon$ is an
imaginary $n$th root of unity, vanishes.}

By (13), (16), (17), §.184 the functions $f_\nu, F_\nu,
\Phi_\nu$ then receive definite values different from zero, and we make
the second assumption,

% [sic] Print runs the list through $f_{n-1}$, although the system defined in
% §.184 (5) and used immediately below runs only from $f_0$ through $f_{n-2}$.
\gesperrt{2. that by the same substitution no two of the functions
$f_0, f_1 \ldots f_{n-1}$ become equal to one another.}

We assume that by the substitution of the $\xi$ for the $x$

\[
f_0, f_1 \ldots f_{n-2}
\]

pass into

\[
k_0, k_1 \ldots k_{n-2},
\]

so that the $k_0, k_1 \ldots k_{n-2}$ are therefore different from one
another, and

\[
\Phi_0, \Phi_1 \ldots \Phi_{n-2}
\]

into

\[
K_0, K_1 \ldots K_{n-2}
\]

where, however, among the $K$ equal ones may also occur.

On account of the property $\delta)$ of the functions $f_\nu$, the
quantities $k_0, k_1 \ldots k_{n-2}$ are the roots of a cyclic equation
of degree $(n-1)$ (hence even), $\psi(u) = 0$, so that, if $u$ denotes a
variable,

\weq{1}{\psi(u) = (u - k_0)(u - k_1) \cdots (u - k_{n-2})}

is a function of degree $(n-1)$ in $u$ with rational coefficients. The
quantities $k_0, k_1 \ldots k_{n-2}$ can be expressed rationally in
terms of one another in the form

\weq{2}{k_1 = \Theta(k_0), \quad k_2 = \Theta(k_1), \ldots k_{n-1} = \Theta(k_{n-2}), \quad k_0 = \Theta(k_{n-1})} % [sic] Print refers to $k_{n-1}$ although the sequence is introduced only through $k_{n-2}$; preserved.
(§.\,163).

Likewise the quantities $K_\nu$ can, by the concluding theorem of
§.184, be expressed in terms of the $k_\nu$, and indeed in the form

\weq{3}{K_0 = \Phi(k_0), \quad K_1 = \Phi(k_1), \ldots K_{n-2} = \Phi(k_{n-2}),}

in which $\Phi$ denotes a rational function (in $\mathfrak{K}$).

Accordingly the formula (18), §.184 yields us the following result: We
put, as an abbreviation,

\weq{4}{\tau_\nu = \sqrt[n]{k_\nu},}
%CONT>>


% ===== scan leaf 665, printed 640 =====
% !PRINTED 640
\weberpage{640}
%<<CONT

and obtain

\weq{5}{(\varepsilon^{r_\nu}, \xi) = K_\nu \tau_\nu^{r_{n-2}} \tau_{\nu+1}^{r_{n-3}} \ldots \tau_{\nu+n-2}^{r_0},}

and if one adds all these formulae, and further sets the rational quantity $(1, x) = A$:

\weq{6}{n\xi_0 = A + \sum_{0,\,n-2}^{\nu} K_\nu \tau_\nu^{r_{n-2}} \tau_{\nu+1}^{r_{n-3}} \ldots \tau_{\nu+n-2}^{r_0}.}

By \S.\,184, (16) and (17) the quantities $K_\nu$ and $k_\nu$ are all different from zero.

By (6) one of the roots $\xi$ is represented.

If we set, for abbreviation,

\[
R_\nu = k_\nu^{r_{n-2}} k_{\nu+1}^{r_{n-3}} \ldots k_{\nu+n-2}
\]

and write (6) in the form

\weq{7}{n\xi_0 = A + K_0 \sqrt[n]{R_0} + K_1 \sqrt[n]{R_1} + \cdots + K_{n-2} \sqrt[n]{R_{n-2}},}

then $\xi_0$ is thus expressed by $n-1$ radicals of degree $n$, each of which can in itself have $n$ different values. But these values cannot be independent of one another, since otherwise the number of values of $\xi$ that result from (7) would be too great. One obtains in fact from (5), if by \S.\,184, (16) one sets

\[
(\varepsilon^{r_\nu}, \xi) (\varepsilon^{r_{\nu-1}}, \xi)^{-g} = k_{\nu-1}
\]

\weq{8}{K_\nu \sqrt[n]{R_\nu} = K_{\nu-1}^g k_{\nu-1} \left( \sqrt[n]{R_{\nu-1}} \right)^g,}

and can accordingly express all these radicals rationally by one of them and by $k_\nu$.

But the expression (6) has over (7) the great advantage \gesperrt{that, however one may determine the $n-1$ radicals $\sqrt[n]{k_\nu}$, it nevertheless represents only $n$ different values, namely the $n$ roots $\xi$}.

To prove this, we denote by $\varepsilon_0, \varepsilon_1 \ldots \varepsilon_{n-2}$ any arbitrary system of $n$th roots of unity, and replace in (5)

\[
\sqrt[n]{k_0}, \quad \sqrt[n]{k_1} \ldots \sqrt[n]{k_{n-2}}
\]

by

\[
\varepsilon_0 \sqrt[n]{k_0}, \; \varepsilon_1 \sqrt[n]{k_1} \ldots \varepsilon_{n-2} \sqrt[n]{k_{n-2}}:
\]
%CONT>>


% ===== scan leaf 666, printed 641 =====
% !PRINTED 641
\weberpage{641}
\pdfbookmark[2]{§ 186. Removal of the restrictive assumptions}{bm208}
%<<CONT

Then (5) passes into another form, which we represent thus

\weq{9}{n\xi = A + \sum_{0,\,n-2}^{\nu} E_\nu K_\nu \tau_\nu^{r_{n-2}} \tau_{\nu+1}^{r_{n-3}} \ldots \tau_{\nu+n-2}^{r_0},}

in which $E_\nu$ is an $n$th root of unity, which is expressed by means of $\varepsilon$ as follows:

\weq{10}{E_\nu = \varepsilon_\nu^{r_{n-2}} \varepsilon_{\nu+1}^{r_{n-3}} \ldots \varepsilon_{\nu+n-2}^{r_0}.}

Now, by the definition of $r_\nu$ [\S.\,184, (15)]

\weq{11}{r_\nu \equiv g r_{\nu-1} \pmod{n},}

and by (10), since the index of $r$ is to be taken with respect to the modulus $n-1$,

\[
E_{\nu-1} = \varepsilon_{\nu-1}^{r_{n-2}} \varepsilon_\nu^{r_{n-3}} \ldots \varepsilon_{\nu+n-3}^{r_0}
\]
\[
E_{\nu-1}^g = \varepsilon_\nu^{r_{n-2}} \varepsilon_{\nu+1}^{r_{n-3}} \ldots \varepsilon_{\nu-1}^{r_0},
\]

hence

\weq{12}{E_\nu = E_{\nu-1}^g = E_0^{r_\nu}.}

If, then, the $\varepsilon$ are determined in any manner whatever, we obtain

\weq{13}{E_0 = \varepsilon_0^{r_{n-2}} \varepsilon_1^{r_{n-3}} \ldots \varepsilon_{n-2}^{r_0},}

and thereby, by (12), the remaining $E_\nu$ are completely determined. Thus in fact only $n$ values result from (6), which one can obtain, for example, by assigning to \gesperrt{one} of the radicals $\tau_\nu$ its various values. (Compare with this the \gesperrt{Cayley} solution of the cubic equations, \S.\,36.)

If one wishes to determine the quantities $\xi_0, \xi_1 \ldots \xi_{n-1}$ in the order which underlies the formation of the cyclic and metacyclic group, one must, in (5), multiply by $\varepsilon^{-h r_\nu}$ before the summation, and finds (\S.\,133)

\weq{14}{n\xi_h = A + \sum_{0,\,n-2}^{\nu} \varepsilon^{-h r_\nu} K_\nu \tau_\nu^{r_{n-2}} \tau_{\nu+1}^{r_{n-3}} \ldots \tau_{\nu+n-2}^{r_0},}

in which an altered determination of the radicals occasions only a cyclic permutation of the $\xi_h$.

\begin{center}
\S.\,186.

\gesperrt{Removal of the Restrictive Assumptions.}
\end{center}

It would be an essential restriction of these investigations on solvable equations if the assumptions 1., 2. of the preceding section had to be maintained.


% ===== scan leaf 667, printed 642 =====
% !PRINTED 642
\weberpage{642}

But this is not necessary, as we now wish to show.

Now let $\eta_0, \eta_1 \ldots \eta_{n-1}$ be the $n$ roots of \gesperrt{any} irreducible metacyclic equation of degree $n$.

We introduce the new unknowns $\xi_0, \xi_1 \ldots \xi_{n-1}$ by a system of equations

\weq{1}{\xi_0 = \psi(\eta_0), \quad \xi_1 = \psi(\eta_1), \ldots \xi_{n-1} = \psi(\eta_{n-1})}

in which $\psi(y)$ is to denote an integral function of degree $(n-1)$

\weq{2}{\psi(y) = a_0 + a_1 y + a_2 y^2 + \cdots + a_{n-1} y^{n-1}}

with indeterminate coefficients from the field $\mathfrak{K}$. We assume, however, that the $n$ quantities $\xi_0, \xi_1 \ldots \xi_{n-1}$ are distinct from one another; then the $\xi$ are likewise the roots of an irreducible metacyclic equation, and by (1) a \gesperrt{Tschirnhausen} transformation is expressed (\S.\,52); $\xi_0$ is a primitive element of the field $\mathfrak{K}(\eta_0)$; consequently the fields $\mathfrak{K}(\xi_0)$ and $\mathfrak{K}(\eta_0)$, and likewise the conjugate fields $\mathfrak{K}(\xi_r)$ and $\mathfrak{K}(\eta_r)$, are identical with one another. Hence also, if we set

\weq{3}{\chi(x) = b_0 + b_1 x + b_2 x^2 + \cdots + b_{n-1} x^{n-1}}

in which the coefficients $b$ are likewise quantities in $\mathfrak{K}$,

\weq{4}{\eta_0 = \chi(\xi_0), \quad \eta_1 = \chi(\xi_1) \ldots \eta_{n-1} = \chi(\xi_{n-1}).}

It is now easy to see that one can dispose of the coefficients $a_0, a_1 \ldots a_{n-1}$ in (1) in such a way that the hypotheses 1. and 2. of the preceding section are fulfilled for the $\xi$. For if we regard the $a$ as independent variables, then the $\xi$ determined by (1) likewise become variables, which are connected with the $a$ by a linear substitution with non-vanishing determinant. The determinant of this substitution is namely

\[
\begin{vmatrix}
1, & \eta_0, & \eta_0^2 & \ldots & \eta_0^{n-1} \\
1, & \eta_1, & \eta_1^2 & \ldots & \eta_1^{n-1} \\
\cdot & \cdot & \cdot & \cdot & \cdot \\
1, & \eta_{n-1}, & \eta_{n-1}^2 & \ldots & \eta_{n-1}^{n-1}
\end{vmatrix},
\]

that is, equal to the difference product of the $\eta$, which by hypothesis is different from zero.

The functions $(\varepsilon, x)$ and the differences $f_\alpha - f_\beta$, which by \S.\,184 are non-vanishing functions of the variables $x_0, x_1 \ldots x_{n-1}$
%CONT>>


% ===== scan leaf 668, printed 643 =====
% !PRINTED 643
\weberpage{643}

%<<CONT
therefore, when the substitution $x_h = \psi(\eta_h)$ is made for the $x$, they pass over into functions of the variables $a$ which likewise cannot vanish identically. By the theorem of \S.\,143, 1. one can therefore set for the quantities $a_0, a_1 \ldots a_{n-1}$ such rational numbers that the functions $(\varepsilon, x)$ and the differences $f_\alpha - f_\beta$, which then pass over into $(\varepsilon, \xi)$ and into $k_\alpha - k_\beta$, become different from zero, and that the $\xi_h$ also remain different from one another, which was to be proved.

Having thus established this, we introduce the system of variables $x_0, x_1 \ldots x_{n-1}$ and a system of variables $y$ depending on it,

\weq{5}{y_0 = \chi(x_0), \quad y_1 = \chi(x_1), \ldots y_{n-1} = \chi(x_{n-1})}

in which $\chi$ denotes the function (3), and set

\weq{6}{\dfrac{(\varepsilon^{r_\nu}, y)}{(\varepsilon^{r_\nu}, x)} = \Theta_\nu.}

If any permutation is applied to the indices of the $x$, the indices of the $y$ undergo the same permutation. If, then, we apply to (6) the cyclic permutation $s$, $\Theta_\nu$ does not change (by \S.\,184, 4.).

If, then, we represent $\Theta_\nu$ in the normal form of \S.\,183, its coefficients are cyclic functions of $x_0, x_1 \ldots x_{n-1}$.

If one applies to $\Theta_\nu$ the permutation $t^{-1} = (h, g^{-1}h)$, then $\Theta_\nu$ passes over, by \S.\,184, 4., into $\Theta_{\nu+1}$, or the $\Theta_\nu$ undergo a cyclic permutation.

But the substitution $\sigma = (\varepsilon, \varepsilon^g)$ has the same effect as well, and therefore the system of functions $\Theta_0, \Theta_1 \ldots \Theta_{n-2}$ has the properties $\alpha)\ \beta)\ \gamma)\ \delta)$ (\S.\,184).

By the theorem proved at the end of \S.\,184, therefore, when we replace the $x$ by the $\xi$ and consequently the $y$ by the $\eta$, $\Theta_\nu$ passes over into a rational function of $k_\nu$

\weq{7}{Q_\nu = Q(k_\nu)}

and from (6) it follows that

\weq{8}{(\varepsilon^{r_\nu}, \eta) = Q_\nu (\varepsilon^{r_\nu}, \xi).}

If one substitutes herein for $(\varepsilon^{r_\nu}, \xi)$ the expression (5), \S.\,185, one obtains for $(\varepsilon^{r_\nu}, \eta)$ an expression of entirely the same form, only with the difference that in place of $K_\nu$ there has entered $Q_\nu K_\nu$. The functions $Q_\nu K_\nu$ have essentially the same properties as the function $K_\nu$, only that they also
%CONT>>


% ===== scan leaf 669, printed 644 =====
% !PRINTED 644
\weberpage{644}
%<<CONT
may in part be zero. By the same alteration, \S.\,185 (6) then yields
 the value of $\eta_0$.

With this, then, the restrictive assumptions of \S.\,185 are
removed, and we have the general theorem:

\gesperrt{I. Every root $\xi$ of a metacyclic equation
of prime degree $n$ can be represented in the
form}
\weq{9}{\xi = A + \sum_{0,\,n-2}^{\nu} K_\nu \tau_\nu^{r_{n-2}} \tau_{\nu+1}^{r_{n-3}} \cdots \tau_{\nu+n-2}^{r_0},}
\gesperrt{where $A$ is a rational quantity, $k_0, k_1 \dots k_{n-2}$ are the
roots, different from one another and from zero, of a cyclic
 equation of degree $(n-1)$, and $K_\nu$ is a rational function
of $k_\nu$, whose form is the same for all $\nu$. The exponents
$r_0, r_1 \dots r_{n-2}$ are the least positive residues of the
numbers $1, g, g^2 \dots g^{n-2}$, where $g$ is a primitive
root of $n$. The $n$ values that one obtains from (9)
when one assigns to the radicals $\tau_\nu = \sqrt[n]{k_\nu}$ their various
 values are the $n$ roots
of one and the same rational equation.}

The formula (9) results from \S.\,185 (6) when $A$ and $K_\nu$
are replaced by $nA$ and $nK_\nu$, which has been done for simplification.

But the converse of Theorem I can now also be stated.

To prove this, we denote by $\varepsilon$ any
imaginary $n$th root of unity and set
\weq{10}{\begin{aligned}\xi_h &= A + \sum_{0,\,n-2}^{\nu} \varepsilon^{h r_\nu} K_\nu \tau_\nu^{r_{n-2}} \tau_{\nu+1}^{r_{n-3}} \cdots \tau_{\nu+n-2}^{r_0}\\ h &= 0, 1, 2 \dots n-1,\end{aligned}}
or, abbreviated,
\weq{11}{\xi_h = A + \sum_{0,\,n-2}^{\nu} \varepsilon^{h r_\nu} K_\nu \sqrt[n]{R_\nu},}
where, as before,
\weq{12}{\sqrt[n]{R_\nu} = \tau_\nu^{r_{n-2}} \tau_{\nu+1}^{r_{n-3}} \cdots \tau_{\nu+n-2}^{r_0}}
is set. We now have to investigate the changes
that result for $\xi_h$ when we

1. determine one of the radicals $\tau_\nu = \sqrt[n]{k_\nu}$ differently, and

2. permute the $k_0, k_1 \dots k_{n-2}$ cyclically.


% ===== scan leaf 670, printed 645 =====
% !PRINTED 645
\weberpage{645}
We wish to give one of the radicals, say $\tau_\alpha$, a different sign, that is, to make the permutation

\weq{13}{(\tau_\alpha,\ \varepsilon^\beta \tau_\alpha)}
where $\beta$ can be an arbitrary exponent, and $\alpha$ is one of the indices $0, 1 \ldots n-2$.

In (12) the radical $\tau_\alpha$ has the exponent $r_{n+\nu-\alpha-2} \equiv r_{\nu-\alpha-1}$, and hence to the permutation (13) there corresponds the permutation

\weq{14}{\left(\sqrt[n]{R_\nu},\ \varepsilon^{\beta r_{\nu-\alpha-1}} \sqrt[n]{R_\nu}\right).}
Now, by the meaning of the numbers $r$, we have generally $r_{\alpha+\beta} \equiv r_\alpha r_\beta \pmod{n}$, and if therefore one introduces the permutation (13) into (11), $h$ passes over into $h + \beta r_{-\alpha-1}$.

1. Thus the permutation (13) calls forth among the indices of $\xi$ the cyclic permutation
\[
(h,\ h + \beta r_{-\alpha-1})
\]

If, secondly, we make the cyclic permutation

\weq{15}{(\tau_0,\ \tau_1, \ldots \tau_{n-2}),}
then this corresponds to the permutations

\weq{16}{\left(\sqrt[n]{R_\nu},\ \sqrt[n]{R_{\nu+1}}\right), \quad (K_\nu, K_{\nu+1}), \quad \nu = 0, 1, 2 \ldots n-2,}
and $\xi_h$ passes over into

\[
A + \sum^{\nu} \varepsilon^{h r_\nu} K_{\nu+1} \sqrt[n]{R_{\nu+1}} = A + \sum \varepsilon^{h r_{\nu-1}} K_\nu \sqrt[n]{R_\nu}
\]
\[
= A + \sum \varepsilon^{h g^{-1} r_\nu} K_\nu \sqrt[n]{R_\nu}.
\]

2. Accordingly the indices of $\xi$ suffer, through the permutation (15), the permutation
\[
(h,\ g^{-1} h).
\]

If we now consider any rational symmetric (or even merely metacyclic) function of the $\xi$:

\[
S(\xi_0,\ \xi_1 \ldots \xi_{n-1}),
\]
then, when we introduce the values (10), we obtain from it a rational function of the radicals

\[
\tau_0,\ \tau_1, \ldots \tau_{n-2}.
\]
But this function does not change when one of these radicals is given another of its $n$ values, and consequently the function must depend rationally on $k_0, k_1 \ldots k_{n-2}$. On account of
%CONT>>


% ===== scan leaf 671, printed 646 =====
% !PRINTED 646
\weberpage{646}
%<<CONT
% [sic] Print changes the argument list here to $\xi_0,\ldots,\xi_{n-2}$;
% the function introduced on printed p. 645 has arguments through $\xi_{n-1}$.
2. but neither does this function change when the cyclic permutation $(k_0, k_1 \ldots k_{n-2})$ is performed among the $k_\nu$, and since now the quantities $k_\nu$ are the roots of a cyclic equation in the field $\mathfrak{K}$, the function $S(\xi_0, \xi_1 \ldots \xi_{n-2})$ is a quantity in $\mathfrak{K}$, i.e. rational. If one applies this to the coefficients of the equation whose roots are the quantities (10), it follows that these quantities are the roots of an equation of degree $n$ in $\mathfrak{K}$.

As for the irreducibility of this equation, the following is to be remarked on that point:

The radicals $\sqrt[n]{R_\nu}$ are, if one of them is rational, all rational, and if one is irrational, so are the others [§.\,185, (8)]. Which of the two cases occurs therefore depends, by §.\,179, on whether the function $x^n - R_0$ is reducible or irreducible in $\mathfrak{K}$. If it is reducible, then one of the quantities $\xi_h$ is rational in the field $\mathfrak{K}$, and the equation for $\xi$ is reducible. All the quantities $\xi$ are contained in the field $\mathfrak{K}(\varepsilon)$, whose degree is at most equal to $n-1$.

But if $\sqrt[n]{R_0}$ is irrational, then $x^n - R_0$ is irreducible in $\mathfrak{K}$. Every equation $\Phi(\xi_0) = 0$ passes, on substitution of the value of $\xi_0$, into an equation $\Psi\!\left(\sqrt[n]{R_0}\right) = 0$ whose coefficients are contained in $\mathfrak{K}$, and which must therefore continue to hold when $\sqrt[n]{R_0}$ is replaced by any one of its $n$ values.

But thereby $\xi_0$ passes into each of the values $\xi_0, \xi_1 \ldots \xi_{n-1}$, i.e. the equation of degree $n$ whose roots are the $n$ quantities $\xi$ is irreducible in $\mathfrak{K}$.

We can therefore state the converse of Theorem I as follows:

\gesperrt{II. Every quantity $\xi$ contained in the form (9) is the root of an equation of degree $n$ in $\mathfrak{K}$, and this equation is irreducible if one leaves aside the special case that one of its roots is rational and the others are contained in $\mathfrak{K}(\varepsilon)$.}


% ===== scan leaf 672, printed 647 =====
% !PRINTED 647
\weberpage{647}
\pdfbookmark[2]{§ 187. Conditions of reality}{bm209}
\begin{center}
\S.\,187.

\gesperrt{Reality Relations.}
\end{center}

If the field $\mathfrak{K}$ is a real one, then, as we have seen in §.\,180, there are two kinds of metacyclic equations of prime degree $n$: those with one real and $n-1$ imaginary roots, and those with all real roots. Likewise we have seen in §.\,165 that there are two kinds of cyclic equations of an even degree, namely those with all real and those with all imaginary roots.

% [sic] Print runs both the $k$ and $\tau$ lists through index $n-1$ here;
% the cyclic equation and radical system defined in §.185 run through $n-2$.
If now the cyclic equation whose roots are the quantities $k_0, k_1 \ldots k_{n-1}$ belongs to the first kind, that is, if $k_0, k_1 \ldots k_{n-1}$ are real and the radicals $\tau_0, \tau_1, \ldots \tau_{n-1}$ are also taken to be real, then formula (10) of §.\,186 shows that $\xi_0$ is real, and that $\xi_h$ and $\xi_{-h}$ are conjugate imaginary.

If on the other hand the $k$ are imaginary, then by §.\,165 $k_\nu$ and $k_{\nu+\frac{n-1}{2}}$ are conjugate imaginary, and by formula (8) in §.\,185 we can set, if $\Phi$ denotes a rational function,

\[
\sqrt[n]{R_{\nu+\frac{n-1}{2}}} = \Phi(k_\nu) \left(\sqrt[n]{R_\nu}\right)^{g^{\frac{n-1}{2}}},
\]
or also

\[
\sqrt[n]{R_{\nu+\frac{n-1}{2}}}\ \sqrt[n]{R_\nu} = \Phi(k_\nu) \left(\sqrt[n]{R_\nu}\right)^{g^{\frac{n-1}{2}}+1}.
\]

Now because $g$ is a primitive root of $n$, $g^{\frac{n-1}{2}}+1$ is divisible by $n$, and the right side therefore becomes rational. If then $\Psi$ denotes a rational function, we have

\[
\sqrt[n]{R_{\nu+\frac{n-1}{2}}}\ \sqrt[n]{R_\nu} = \Psi(k_\nu),
\]
and this formula shows that

\[
\Psi(k_\nu) = \Psi\!\left(k_{\nu+\frac{n-1}{2}}\right),
\]
hence that $\Psi(k_\nu)$ is real. From this it follows that

\[
\sqrt[n]{R_{\nu+\frac{n-1}{2}}} \quad \text{and} \quad \sqrt[n]{R_\nu}
\]
%CONT>>


% ===== scan leaf 673, printed 648 =====
% !PRINTED 648
\weberpage{648}
\pdfbookmark[2]{§ 188. Metacyclic equations of the 5th degree}{bm210}
%<<CONT
are conjugate imaginary; for, first, their $n$th powers are conjugate
imaginary, and hence they themselves are so, at first up to an $n$th root of unity as
factor. But since their product is real, this $n$th root of unity
must be $= 1$. Since, moreover,
\[
r_\nu \equiv - r_{\nu + \frac{n-1}{2}} \pmod{n}
\]
formula (11) shows that in this case the roots $\xi_z$ are all real.

With this, then, the theorem is proved for a real field $\mathfrak{K}$:

\gesperrt{If the cyclic equation whose roots are the $k_0$, $k_1 \ldots k_{n-2}$ has real
roots, then of the roots $\xi$ one is real and the rest imaginary; if this
cyclic equation has imaginary roots, then all the $\xi$ are real.}

\begin{center}
§. 188.

\gesperrt{Metacyclic Equations of the Fifth Degree.}
\end{center}

Through the theorems of the preceding section the roots of a metacyclic
equation of prime degree have been brought into a generally valid form, which permits
all these quantities to be actually formed, if one presupposes a knowledge of the roots of cyclic
equations. Whole systems of cyclic equations of every degree are furnished to us,
e.g., by the theory of the division of the circle, whose roots are the cyclotomic periods. Thus
we can, for example, form arbitrarily many
metacyclic equations for the field of rational numbers. That all these equations in the field of
rational numbers are contained among them is a very remarkable theorem due to \gesperrt{Kronecker},
which we shall come to know in the second volume. Here, as an
application and illustration of what precedes, we wish further to treat the special problem
of finding, in an arbitrary field $\mathfrak{K}$, the roots of all metacyclic
equations of the fifth degree.

For this it is necessary and sufficient that we determine generally the roots $k_0, k_1, k_2, k_3$ of a
cyclic equation of the fourth degree, and indeed under the
assumption that \gesperrt{$k_0, k_1, k_2, k_3$ are different from one another and none
of them is zero}.


% ===== scan leaf 674, printed 649 =====
% !PRINTED 649
\weberpage{649}
As a foundation we here use the function
\weq{1}{w = (k_0 - k_2)(k_1 - k_3),}
which must be different from zero, and which for a real field $\mathfrak{K}$
is always real, since either $k_0, k_1, k_2, k_3$ are real or $k_0, k_2$ and
$k_1, k_3$ form two conjugate imaginary pairs.

Under the cyclic permutation $(0, 1, 2, 3)$, $w$ changes its sign. Hence the
square of $w$ is a cyclic function, which by hypothesis is to be known. We
therefore set, denoting by the small latin letters $a, b, c \ldots$
quantities in $\mathfrak{K}$, that is, rational quantities,
\weq{2}{w = 4\sqrt{c},}
and remark further that every function of the $k$ which changes its sign
under the cyclic permutation $(0, 1, 2, 3)$ is the product of a rational
quantity with $\sqrt{c}$.

It is then only a matter of forming cyclic functions in sufficient number, so
that the four quantities $k_0, k_1, k_2, k_3$ can be computed from them and
expressed algebraically in terms of mutually independent rational quantities.

Now two further cyclic functions are
\[
(k_0 - k_2)^2 + (k_1 - k_3)^2, \qquad \frac{(k_0 - k_2)^2 - (k_1 - k_3)^2}{(k_0 - k_2)(k_1 - k_3)},
\]
and we thus obtain, if $a, b$ are rational quantities,
\weq{3}{\begin{aligned}
(k_0 - k_2)^2 + (k_1 - k_3)^2 &= 8\,b \\
(k_0 - k_2)^2 - (k_1 - k_3)^2 &= 8\,a\sqrt{c}.
\end{aligned}}
The three rational quantities $a, b, c$ are, however, not independent of one
another; rather, by (2) and (3) there subsists between them the relation
\weq{4}{b^2 = c\,(1 + a^2).}
From (3) it follows further that
\weq{5}{\begin{aligned}
k_0 - k_2 &= 2\sqrt{b + a\sqrt{c}} \\
k_1 - k_3 &= 2\sqrt{b - a\sqrt{c}},
\end{aligned}}
and between the two square roots there subsists the relation
\weq{6}{\sqrt{c} = \sqrt{b + a\sqrt{c}}\; \sqrt{b - a\sqrt{c}}.}
If we further denote by $C$ and $B$ again two rational quantities, then
%CONT>>


% ===== scan leaf 675, printed 650 =====
% !PRINTED 650
\weberpage{650}
%<<CONT
\weq{7}{\begin{aligned}k_0+k_1+k_2+k_3 &= 4\,C,\\ k_0-k_1+k_2-k_3 &= 4\,B\sqrt{c},\end{aligned}}
hence
\weq{8}{\begin{aligned}k_0+k_2 &= 2\,(C+B\sqrt{c}),\\ k_1+k_3 &= 2\,(C-B\sqrt{c}),\end{aligned}}
from which by (5) it follows:
\weq{9}{\begin{aligned}k_0 &= C+B\sqrt{c}+\sqrt{b+a\sqrt{c}}\\ k_1 &= C-B\sqrt{c}+\sqrt{b-a\sqrt{c}}\\ k_2 &= C+B\sqrt{c}-\sqrt{b+a\sqrt{c}}\\ k_3 &= C-B\sqrt{c}-\sqrt{b-a\sqrt{c}}.\end{aligned}}
Conversely, it is also easy to show that these quantities
are the roots of a biquadratic cyclic equation.
For if we set, as an abbreviation,
\weq{10}{r = \sqrt{c}, \quad \varrho = \sqrt{b+a\sqrt{c}}, \quad \varrho' = \sqrt{b-a\sqrt{c}},}
and hence by (6)
\[ r = \varrho\varrho', \quad \varrho^2 = b+ar, \quad \varrho'^2 = b-ar, \]
then we can represent every rational function of $k_0, k_1, k_2, k_3$ as a linear
function with rational coefficients of the six radicals
\[ 1, \quad r, \quad \varrho, \quad \varrho', \quad r\varrho, \quad r\varrho' \]
The cyclic permutation $(k_0, k_1, k_2, k_3)$ corresponds
to the permutation
\[
\begin{array}{lll}
r, & \varrho, & \varrho',\\
-r, & \varrho', & -\varrho,
\end{array}
\]
and if one repeats this permutation, one sees
that a cyclic function of the $k$ cannot change
if, in the linear representation by these six quantities,
the following permutations are made:
\[
\begin{array}{llllll}
1, & r, & \varrho, & \varrho', & r\varrho, & r\varrho'\\
1, & -r, & \varrho', & -\varrho, & -r\varrho', & r\varrho\\
1, & r, & -\varrho, & -\varrho', & -r\varrho, & -r\varrho'\\
1, & -r, & -\varrho', & \varrho, & r\varrho', & -r\varrho,
\end{array}
\]
and if one adds the four expressions thus obtained,
one obtains a rational expression for the cyclic function
of the $k$.

If one wishes to form the biquadratic equation whose roots
are the $k$, it is best to introduce $k - C = x$ as the unknown
%CONT>>


% ===== scan leaf 676, printed 651 =====
% !PRINTED 651
\weberpage{651}

%<<CONT
in. One obtains in this way, by a simple calculation, the equation
\weq{11}{x^4 - 2\,(B^2c + b)\,x^2 - 4\,B\,a\,c\,x + B^4c^2 - 2\,B^2bc + c = 0.}
The representation of the $k$ by the formulae (9) is, however, not yet completely satisfactory, because among the rational quantities $a$, $b$, $c$ occurring in them the relation (4) still holds, and we seek a representation by independent quantities.

One special case we must first dispose of, namely $b = 0$, which has $a = i$ as a consequence. Then the formulae (9) give
\[
\begin{aligned}
k_0 &= C + B\sqrt{c} + \frac{1+i}{\sqrt{2}}\sqrt[4]{c}\\
k_1 &= C - B\sqrt{c} + \frac{1-i}{\sqrt{2}}\sqrt[4]{c}\\
k_2 &= C + B\sqrt{c} - \frac{1+i}{\sqrt{2}}\sqrt[4]{c}\\
k_3 &= C - B\sqrt{c} - \frac{1-i}{\sqrt{2}}\sqrt[4]{c},
\end{aligned}
\]
and the biquadratic equation (11) becomes
\[
x^4 - 2\,B^2c\,x^2 - 4\,i\,B\,c\,x + B^4c^2 + c = 0.
\]
But this case belongs here only when $i$ is contained in the field $\mathfrak{R}$, and hence never for real fields.

If, however, $b$ does not vanish, we introduce a new rational quantity $h$ by setting
\[
b = h\,(1 + a^2), \quad c = h^2\,(1 + a^2)
\]
whereby the relation (4) is then identically satisfied, and the formulae (9), if one replaces $Bh$ by $B$, give
\weq{12}{
\begin{aligned}
k_0 &= C + B\sqrt{1+a^2} + \sqrt{h\,\bigl(1+a^2+a\sqrt{1+a^2}\bigr)}\\
k_1 &= C - B\sqrt{1+a^2} + \sqrt{h\,\bigl(1+a^2-a\sqrt{1+a^2}\bigr)}\\
k_2 &= C + B\sqrt{1+a^2} - \sqrt{h\,\bigl(1+a^2+a\sqrt{1+a^2}\bigr)}\\
k_3 &= C - B\sqrt{1+a^2} - \sqrt{h\,\bigl(1+a^2-a\sqrt{1+a^2}\bigr)}
\end{aligned}
}
This expression for $k_0$ has moreover the advantage that, however one may determine the signs of the square roots occurring in it, it represents only four distinct values.


% ===== scan leaf 677, printed 652 =====
% !PRINTED 652
\weberpage{652}

In \gesperrt{Abel} a somewhat different expression is found for $k_0$,
namely
\[
k_0 = C + B\sqrt{1+e^2} + \sqrt{h\left(1+e^2+\sqrt{1+e^2}\right)},
\]
which arises from (12) when one sets $a = 1:e$ and then replaces $B$ and $h$
by $Be$ and $he^2$. The expression (12) is thus more general in so far
as it also embraces the special case $a=0$, in which the
biquadratic equation splits into two quadratic equations.

Now, in order to represent the root of a metacyclic equation of degree
5, these expressions for $k_0, k_1, k_2, k_3$ are to be substituted into
the formula (9) of \S.\,186. If we assume $g=2$, the
exponents $r_0, r_1, r_2, r_3$ become congruent in turn with
1, 2, 4, 8, hence equal to 1, 2, 4, 3, and there results, when, as in
\S.\,185,
\[
\begin{aligned}
\tau_0 &= \sqrt[5]{k_0}, & \tau_1 &= \sqrt[5]{k_1} \\
\tau_2 &= \sqrt[5]{k_2}, & \tau_3 &= \sqrt[5]{k_3}
\end{aligned}
\]
is set,
\weq{13}{
\begin{aligned}
\xi = A &+ K_0\, \tau_0^3\tau_1^4\tau_2^2\tau_3 + K_1\, \tau_1^3\tau_2^4\tau_3^2\tau_0 \\
        &+ K_2\, \tau_2^3\tau_3^4\tau_0^2\tau_1 + K_3\, \tau_3^3\tau_0^4\tau_1^2\tau_2.
\end{aligned}
}
The coefficients in this expression, $K_0, K_1, K_2, K_3$, are
rational functions of the $k_0, k_1, k_2, k_3$, which are themselves
cyclically permuted by cyclic permutation of the $k$. According to
\gesperrt{Abel} one is to set
\[
K_0 = A_1 + A_2 k_0 + A_3 k_2 + A_4 k_0 k_2
\]
wherein $A_1, A_2, A_3, A_4$ are rational; but this is not
general enough, because between $k_0 k_2$, $k_0$ and $k_2$ there subsists a
linear relation easily derivable from (12). Most simply one
expresses the $k$ in terms of the three radicals
\[
\begin{aligned}
r &= \sqrt{1+a^2}, & \varrho &= \sqrt{h\left(1+a^2+a\sqrt{1+a^2}\right)} \\
\varrho' &= \sqrt{h\left(1+a^2-a\sqrt{1+a^2}\right)}, & \varrho\varrho' &= hr
\end{aligned}
\]

One can express the radical $\varrho'$ linearly in terms of $r\varrho$ and
$\varrho$, as one sees from the formula
\[
hr^2\varrho' = r\varrho\varrho'^2
\]
when one substitutes on the right for $\varrho'^2$ its expression in terms of $r$
and bears in mind that $r^2$ is rational. Thus one obtains, when
instead of the cyclic permutation of the $k$ one takes the permutation
\[
\begin{pmatrix}
r, & \varrho, & \varrho' \\
-r, & \varrho', & -\varrho
\end{pmatrix}
\]
%CONT>>


% ===== scan leaf 678, printed 653 =====
% !PRINTED 653
\weberpage{653}
%<<CONT
applies, and denotes by $A_1$, $A_2$, $A_3$, $A_4$ rational quantities:
\weq{14}{\begin{split}
K_0 &= A_1 + A_2 r + A_3 \varrho + A_4 r\varrho\\
K_1 &= A_1 - A_2 r + A_3 \varrho' - A_4 r\varrho'\\
K_2 &= A_1 + A_2 r - A_3 \varrho + A_4 r\varrho\\
% [sic] Print repeats $A_2$ as the coefficient of $\varrho'$ in the $K_3$
% row; the three parallel rows make $A_3$ the structurally expected coefficient.
K_3 &= A_1 - A_2 r - A_2 \varrho' - A_4 r\varrho'.
\end{split}}
These quantities $K_0$, $K_1$, $K_2$, $K_3$ are themselves again the roots of a cyclic biquadratic equation. They have only apparently a more general form than the $k$; for if one puts $K_0$ into the form
\[
K_0 = A_1 + A_2 r + \sqrt{\varrho^2(A_3+A_4 r)^2},
\]
one recognises the form of $k_0$ in (9) again, naturally with altered $a$, $b$, $c$.
% [printed page ends here: a short horizontal rule follows below this sentence, marking the close of \S.\,188; the remainder of the printed page is blank]


% ===== scan leaf 679, printed None =====
% !PRINTED none
% Page carries no running head or folio number. The upper portion of the leaf is
% blank (scan shows only faint dust/scanning artefacts); printed matter appears
% only in the lower third: a centred heading "Berichtigungen." followed by two
% single-line corrections.
\begin{center}
\gesperrt{Corrections.}

\medskip
\rule{1.5cm}{0.4pt}
\end{center}

Page 182, in the formula, line 12, read $x_m$ instead of $x_n$.\\
Page 347, in the formula, line 24, read $(2x^2+1)^2$ instead of $(2x^2-1)^2$.

\begin{center}
\medskip
\rule{4cm}{0.4pt}
\end{center}


% ===== scan leaf 680, printed None =====
% !PRINTED none
% Blank leaf: no printed text, figure, or rule visible in any of the three scan
% bands (top, mid, bot) — only faint scanning dust artefacts.


% ===== scan leaf 683, printed None =====
% !PRINTED none
% This scan is not a printed page of Weber's text. All three bands (top/mid/bot) show
% the library's own front-matter material bound into the volume: a "RETURN TO / Astronomy,
% Mathematics, Statistics, Computer Science Library / 100 Evans Hall / 642-3381" due-date
% loan card with a grid of stamped and struck-through due dates (e.g. "DEC 6 1985",
% "SEP 24 1987" struck, "OCT 26 1988" struck, "JUL 06 1992", "JUL 26 1994" struck,
% "MAR 06 1998"), the printed lines "ALL BOOKS MAY BE RECALLED AFTER 7 DAYS" and
% "DUE AS STAMPED BELOW", and the imprint "UNIVERSITY OF CALIFORNIA, BERKELEY / BERKELEY,
% CA 94720 / FORM NO. DD3, 1/83". No Weber body text, no running head, and no printed
% page number appears on this leaf. The barcode and call number are on the following leaf.
% There is no German text on this page; nothing to translate.


% ===== scan leaf 684, printed None =====
% !PRINTED none
% This scan is not a printed page of Weber's text. The upper area carries a printed
% "U.C. BERKELEY LIBRARIES" barcode label (C037482993) and the handwritten call number
% "QA 153 / W 35 / v. 1". Elsewhere the otherwise blank library-processing leaf has the
% pencil annotations "9-27-83" and a partial "-593" with a small dot beside it. No Weber
% body text, running head, or printed page number appears on the leaf.


% ===== scan leaf 685, printed None =====
% !PRINTED none
% Non-text terminal scan leaf: the registered source page is a dark, uniformly textured
% cover/end-sheet surface with no printed text, handwriting, running head, folio, rule, or
% discernible letterforms. There is therefore no Weber content to translate.

\end{document}
